% Version 2
% 1.  Define \into
% 2.  Define m-chart at a point and subchart at a point
% 3,  Define (local) m-morphism
% 4.  Define local M-atlas morphism
% 5.  Change definition of M-atlas morphism
% 6.  Define classic M-atlas morphism
% 7.  Change definition of Ck-atlas morphism
% 8.  Define classic Ck-atlas morphism
% 9.  Move definitions and proclamations to arXiv:1906.11690}, 2019:
%     i.    semi-maximal atlas
%     ii.   near morphism
%     iiii, Associated model spaces and functors

\documentclass{article}
\usepackage{amsmath}
\usepackage{amsthm}
\usepackage{bm}
\usepackage{enumitem}
\usepackage{ifthen}
\usepackage{mathtools}
\usepackage{scalerel}
\usepackage{stix2}[notext,not1] %does not reqire XeTeX or luaTeX
% thmtools removed 2026-09-01: loaded but nothing from it (\declaretheorem,
% \listoftheorems, a custom theorem style) was ever actually used anywhere
% in this document (confirmed by grep, not assumed) -- and loading it
% alongside amsthm + cleveref v0.21.4 (the version installed here)
% reproducibly breaks \newtheorem{corollary}[theorem]{Corollary} with a
% fatal "Command \c@corollary already defined" error, isolated to a
% 5-line minimal reproduction with just those three packages loaded and
% one shared-counter \newtheorem call; amsthm + cleveref alone, without
% thmtools, compiles that same reproduction cleanly. Zero functional loss.
\usepackage{tikz-cd}
\usepackage{xparse} % loads expl3
%See interface3.pdf
\usepackage{xstring}
\usepackage{semantic-markup}

\usepackage[cmtip,all,barr]{xy}
%Remove when xybarr.tex bug fixed
\newdir_{ (}{{ }*!/-.5em/@_{(}}

% Group action
\newcommand \action [1] []
   { \ifthenelse {\equal{#1}{}}
       {\mathbin \star}
       {\mathbin{\overset {#1} {\star}}}
   }

%Morphisms of category
\newcommand \Ar {\mathrm{Ar}}

\DeclareMathOperator \arin {\stackrel{\Ar}{\in}}

%Category of atlases from E to C
\newcommand \Atl [1] []
   { \ifthenelse {\equal{#1}{}}
     {\mathscr{A\!t\!l}}
     {\mathscr{A\!t\!l}^{\mathrm{#1}}}
   }

\DeclareMathOperator \Atlas {\mathrm{Atlas}}

\newcommand \Bun {\mathrm{Bun}}

\newcommand \BunProd {\mathrm{Bun}-\mathrm{prod}}

\newcommand \C {\mathrm{C}}

% \cat, \catname, \catseqname removed 2026-09-01, superseded by
% semantic-markup's \catName/\catSeqName (auto-detecting single-token/
% decorated/all-upper-case vs. genuine word, matching this exact
% mathbf-vs-mathscr split -- all call sites renamed, verified against
% every real argument via scripts/check_arg_consistency.py in that
% repo). See SEMANTIC-MARKUP-MIGRATION.md.

%Category of model space
\newcommand \Cat {\mathrm{Cat}}

\newcommand \Ck {\mathrm{C^k}}

\newcommand \Classic {\mathrm{Classic}}

\DeclareMathOperator \codomain {\mathrm{codomain}}

%\newcommand \compat [2] [] {\underset{\mathrm{{#1}compat}{#2}}

% Extended composition operators for function sequences -- \morphCompose,
% \morphComposeHead, \morphComposeTail removed 2026-09-01, superseded by
% semantic-markup's \morphCompose, \morphComposeHead, \morphComposeTail
% (all call sites renamed accordingly). See SEMANTIC-MARKUP-MIGRATION.md.

\DeclareMathOperator \defeq {\stackrel{\mathrm{def}}{=}}

\DeclareMathOperator \domain {\mathrm{domain}}

\DeclareMathOperator \Domain {\seqname{domain}}

\newcommand \false {\mathrm{False}}

\newcommand \Fib   {\mathrm{Fib}}

%\newcommand \full [2] [] {\underset{\mathrm {{#1}full}}{#2}}

\newcommand \fullcref [1]
   { \ifthenelse {\equal{\nameref{#1}}{}}
       {\cref{#1}}
       {\cref{#1} (\nameref{#1})}
   }

% Select font for functions and sequences of functions
% \newcommand \funcname [1] {\mathit{#1}}
% \newcommand \funcseqname [1] {\bm{#1}}

\DeclareMathOperator \Functor {\mathop{\mathscr{F}}}

\newcommand \group {\mathrm{group}}

\DeclareMathOperator \head {\mathrm{head}}

\DeclareMathOperator \Hom {\mathrm{Hom}}

%Identity function
%\DeclareMathOperator* \Id {\mathrm{Id}}

%Sequence of identity functions
%\DeclareMathOperator* \ID {\mathbf{Id}}

% \into (this paper's own macro, since renamed to \morphMono below)
% removed 2026-09-01 (second pass): semantic-markup's \morphMono was
% originally \hookrightarrow, visually distinct from \mon (this
% paper's own tail-arrow) -- re-reading this paper's own Conventions
% section (Part II) settled it: the hook arrow is explicitly this
% paper's INCLUSION MAP symbol, and the tail arrow is explicitly its
% MONOMORPHISM symbol, two genuinely different notions. The package's
% \morphMono was corrected to \rightarrowtail (matching \mon) rather
% than this paper's own arrow being changed to match the package. See
% SEMANTIC-MARKUP-MIGRATION.md.

%Propositional function isCk_{f,A}(x,y,...)
\DeclareMathOperator \isCk {\mathrm{isCk}}

%Propositional function isAtl(A,E,C)
\newcommand \isAtl [1] []
   { \ifthenelse {\equal{#1}{}}
     {\mathrm{isAtl}}
     {\mathrm{isAtl}^{\mathrm{#1}}}
   }

%Propositional function isLCS(L, M, A, F, Sigma)
\DeclareMathOperator \isLCS {\mathrm{isLCS}}

\DeclareMathOperator \iso {\stackrel{\sim}{=}}

% Join two tuples
\DeclareMathOperator \join {\mathrm{join}}

%Category of M-Sigma local coordinate spaces
\newcommand \LCS {\mathrm{LCS}}

\DeclareMathOperator \length {\mathrm{length}}

\DeclareMathOperator \lengtho {\mathrm length0}

\newcommand \M {\mathrm{M}}

\newcommand \Man {\mathrm{Man}}

% Adjust : spacing for f maps a to b
\newcommand \maps {\!\colon}

%Set long names upright, short names italic
\newcommand{\mathvarname}[1]{%
  \begingroup\noexpandarg
  \StrLen{#1}[\temp]%
  \ifnum\temp>1
    \mathrm{#1}%
  \else
    #1%
  \fi
  \endgroup
}

\newcommand \minimal [2] [] {\underset{\mathrm{{#1}min}}{#2}}

%Model category
\newcommand \Mod {\mathrm{Mod}}

%Morphisms between two objects
\newcommand \Mor {\mathrm{Mor}}

% Morphism of category
\DeclareMathOperator \morphin {\stackrel{\Ar}{\in}}

%Object of category
\newcommand \Ob {\mathrm{Ob}}
\DeclareMathOperator \objin {\stackrel{\Ob}{\in}}

% \morphEpi removed 2026-09-01, superseded by semantic-markup's \morphEpi
% (all call sites renamed) -- confirmed visually identical to \epi
% (Xy-pic) first, unlike \morphMono/\mon below. See
% SEMANTIC-MARKUP-MIGRATION.md.

% All non-null open sets
\newcommand \op [2] [] {\underset{\mathrm {{#1}op}}{#2}}

\newcommand \optriv [2] [] {\underset{\mathrm {{#1}op-\mathscr{T\!r\!i\!v}}}{#2}}

\newcommand \pagecref [1]
   { \ifthenelse {\equal{\nameref{#1}}{}}
       {\cref{#1} on \cpageref{#1}}
       {\cref{#1} (\nameref{#1}) on \cpageref{#1}\negmedspace}
   }

\newcommand \Pagecref [1]
   { \ifthenelse {\equal{\nameref{#1}}{}}
       {\Cref{#1} on \cpageref{#1}}
       {\Cref{#1} (\nameref{#1}) on \cpageref{#1}}
   }

\DeclareMathOperator \range {\mathrm{range}}

\DeclareMathOperator \Range {\seqname{range}}

\newcommand \restrictto {\negthickspace\restriction}

\DeclareMathOperator \seqeq {\stackrel{()}{=}}

\DeclareMathOperator \seqin {\stackrel{()}{\in}}

%Font for sequence
%\newcommand \seqname [1] {\bm{\mathit{\mathsf{#1}}}}

%Underlying set of object in concrete category
\DeclareMathOperator* \Set {\mathbf{Set}}
%Use \catName{Set} for the standard category Set.
%
%\newcommand \sing [2] [] {\underset{\mathrm{{#1}Sing}{#2}}
%\newcommand \Sing [2] [] {\underset{\bm{\mathrm{{#1}Sing}}{#2}}

\newcommand \singcat [2] [] {\underset{{#1}\mathscr{S\!i\!n\!g}}{#2}}
\newcommand \Singcat [2] [] {\underset{\bm{{#1}\mathscr{S\!i\!n\!g}}}{#2}}
% Too many fonts to use \mathbfscr without XeTeX

%newcommand \strict [2] [] {\underset{\mathrm {{#1}strict}}{#2}}

%Subcategory
\newcommand \subcat [1] [] {\overset{\mathrm{{#1}cat}}{\subseteq}}

%Sequence of subcategories
\newcommand \SUBCAT [1] [] {\overset{\mathbf{{#1}cat}}{\subseteq}}

%Model subspace
\newcommand \submod [1] [] {\overset{\mathrm{{#1}mod}}{\subseteq}}

\DeclareMathOperator \SUBSETEQ {\stackrel{()}{\subseteq}}

%subspace of topological space
\newcommand \subsp [1] [] {\overset{\mathrm{{#1}subsp}}{\subseteq}}

\DeclareMathOperator \tail {\mathrm{tail}}

\newcommand \toiso {\,\to/{>}->>/^{\iso}}

%Topological space of model space
%Underlying topological space of object in concrete category
\newcommand \Top {\mathrm{Top}}
%Use \catName{Top} for the standard category Top.

%Topological category of model space
\newcommand \Topcat {\mathscr{T\!o\!p}}

% Font for topology
\newcommand \topname [1] {\mathfrak{#1}}

%Topology of topological space and relative topology of subset
\newcommand \Topology {\mathfrak{T\!o\!p}}

% \newcommand \triv [2] [] {\underset{{#1}\mathrm {triv}}{#2}}
% \newcommand \Triv [2] [] {\underset{{#1}\mathbf {triv}}{#2}}

\newcommand \trivcat [2] [] {\underset{\mathrm{#1}\mathscr{T\!r\!i\!v}}{#2}}
%\newcommand \Trivcat [2] [] {\underset{\mathrm{#1}\mathbfscr{T\!r\!i\!v}}{#2}}
\newcommand \Trivcat [2] [] {\underset{\mathrm{#1}\bm{\mathscr{T\!r\!i\!v}}}{#2}}
% Too many fonts to use \mathbfscr without XeTeX

\newcommand \true {\mathrm{True}}
\newcommand \truthcat {\mathscr{Truth}}
\newcommand \truthset {\mathbb{T}}
\newcommand \truthspace {\mathrm{Truthspace}}
\newcommand \truthtop {\mathfrak{Truthtop}}

\newcommand \unioncat [1] [] {\overset{\mathrm{{#1}cat}}{\cup}}
\newcommand \UNIONCAT [1] [] {\overset{\mathbf{{#1}cat}}{\cup}}

% Existential and universal quamtifiers
% Set former
% Union and intersection

% --------------------------------------------------------------------
% | From here to closing --- belongs in package                      |
% |                                                                  |

% Copyright 2016 Shmuel (Seymour J.) Metz
% I grant permission to the AMS, arXiv.org, the LATEX3 project and the
% TeX users group to incorporate these commands in any LaTeX package
% that may be freely redistributed, provided that they attribute the
% source.

% These commands are intended to allow semantic markup for some
% common mathermtical constructs:
%   \equant{variables}{proposition}     Existential qunatifier
%   \uquant{variables}{proposition}     Universal   quantifier
%   \union[indices]{set}                Union
%   \intersection[indices]{set}         Intersection
%   \set {elements}[propositions]       Set of
%   \set{elements}[propositions]*       Set of, split
%   \setupquant{optionstring}           Style of \equant, \uquant
%   \setupset{optionstring}             Style of \set, \union, \intersection
%   \full[qualifiers]{name}             Render name underset with some variation of full
%   \maxfull[qualifiers]{name}          Render name underset with some variation of max or full
%   \maximal[qualifiers]{name}          Render name underset with some variation of max
%   \strict[qualifiers]{name}           Render name underset with some variation of strict
%   \seqname{name}                      Render sequence/set/tuple name
%   \Triv[qualifiers]{name}             Render name underset with some variation of bold triv

% Because LaTeX has problems with parameters containg \\, the \set command
% has code to split the line between the elements and the proposition if
% the invocation is \set* and the environment is multline.

% Surround individual indices in \intersection, individual variables in
% quantifiers, individual propositions in \set and individual indices in
% \union with braces, e.g., \equant {{x \in X},{y \in Y}}{P(x,y)},
% \set{x}[{P(x)}, {Q(x)}], \union[{i \in I},{j \in J}]{O(i.j)}.

% Setup keywords for \equant, \uquant
% subscript  =none            Default
%                             \exists var1 ... \exists varn prop
%             stacked         \exists_\substack{var1 \\ ... varn} prop
%             multiple        \exists_{var1,...,varn} prop
% parentheses=none            Default
%             single          (\exists var1 ... \exists varn) prop
%             multiple        (\exists var1) ... (\exists varn) prop
% separator=                  Default {}

% Setup keywords for \intersection, \set, \union
% subscript  =none            \bigcap ix1, ..., ixn set
%             stacked         Default
%                             \bigcap__\substack{ix1 \\ ... ixn} set
%             multiple        \bigcap_{ix1,...,ixn} set
% separator=                  Default {\mid}
%                             {elements separator prop1 ^ ... propn}

\ExplSyntaxOn

\int_gzero_new:N \g_style_quant_parens_int
\int_gzero_new:N \g_style_quant_subscr_int
\int_gzero_new:N \g_style_set_subscr_int

\NewDocumentCommand{\equant}{mm}
  {
    \quant:nnn {\exists} {#1} {#2}
  }

\NewDocumentCommand{\uquant}{mm}
  {
    \quant:nnn {\forall} {#1} {#2}
  }

\NewDocumentCommand \setupquant {m}
  {
    \keys_set:nn {shmuel / quant} {#1}
  }

\keys_define:nn {shmuel / quant}
 {
  subscript            .choices:nn =
    {
      {
        none,
        stacked,
        multiple
      }
      {
        \int_gset:Nn \g_style_quant_subscr_int {\l_keys_choice_int-1}
      }
    },
  subscript            .default:n = multiple,
  subscript            .initial:n = none,
  parentheses          .choices:nn =
    {
      {
        none,
        single,
        multiple
      }
      {
        \int_gset:Nn \g_style_quant_parens_int {\l_keys_choice_int - 1}
      }
    },
  parentheses          .default:n = multiple,
  parentheses          .initial:n = none,
  separater            .tl_set:N = \g_style_quant_sep_tl,
  separater            .default:n = {.},
  separater            .initial:n = {}
 }

\cs_new:Npn \quant:nnn #1 #2 #3
  {
    %\int_show:N \g_style_quant_parens_int
    %\int_show:N \g_style_quant_subscr_int
    % g_style_quant_parens_int \ \int_use:N \g_style_quant_parens_int \
    % g_style_quant_subscr_int \ \int_use:N \g_style_quant_subscr_int \
    \clist_set:Nn \l_tmpa_clist {#2}
    \int_case:nn
      {\g_style_quant_subscr_int}
      {
        {0}
        {
          % No subscript
          % Set separater to ) ( quantifier or just quantifier
          \int_compare:nTF {\g_style_quant_parens_int = 2}
            {\tl_set:Nn \l_tmpa_tl {\right ) \left ( #1}}
            {\tl_set:Nn \l_tmpa_tl {#1}}
          \int_compare:nT {\g_style_quant_parens_int > 0} {\left (}
          #1
          \clist_use:Nn \l_tmpa_clist {\l_tmpa_tl}
          \int_compare:nT {\g_style_quant_parens_int > 0} {\right )}
          \g_style_quant_sep_tl #3
        }
        {1}
        {
          % Stacked subscript on single quantifier
          \fp_set:Nn
            \l_tmpa_fp
            {ceil{\clist_count:N{\l_tmpa_clist} - 1} * .1 + 1}
          \int_compare:nT {\g_style_quant_parens_int > 0} {\left (}
          \scaleobj{\fp_to_decimal:N \l_tmpa_fp}{#1} \sb
             { \substack { \clist_use:Nn \l_tmpa_clist { \\ } } }
          \int_compare:nT {\g_style_quant_parens_int > 0} {\right )}
          #3
        }
        {2}
        {
          % Subscripts on separate quantifiers
          \clist_map_inline:Nn
            \l_tmpa_clist
            {
              % (quantifier \sb predicate) or quantifier \sb predicate
              {
                \int_compare:nT
                  {\g_style_quant_parens_int > 0}
                  {\left (}
                \scaleobj{1.2}{#1} \sb
                {##1}
                \int_compare:nT
                  {\g_style_quant_parens_int > 0}
                  {\right )}
              }
            }
          #3
        }
      }
  }

\NewDocumentCommand{\set}{mos}
  {
    \IfBooleanTF {#3}
    {
%      \msg_term:n {set with star}
%      \msg_term:n{{P1 #1}}
%      \msg_term:n{{P2 #2}}
      \bool_set_true:N \l_tmpa_bool
%      \bool_show:N \l_tmpa_bool
    }
    {
      \bool_set_false:N \l_tmpa_bool
    }
    \set_of:nnn {#1} {#2} {\l_tmpa_bool}
  }

%\tl_new:N \g_style_set_sep_tl
%\tl_gset:Nn \g_style_set_sep_tl {\mid}

\NewDocumentCommand \setupset {m}
  {
    \keys_set:nn {shmuel / set} {#1}
  }

\keys_define:nn {shmuel / set}
  {
    separater .tl_set:N = \g_style_set_sep_tl,
    subscript            .choices:nn =
      {
        {
          stacked,
          multiple
        }
        {
          \int_gset:Nn \g_style_set_subscr_int {\l_keys_choice_int-1}
        }
      },
    separater .initial:n = {\mid},
    subscript .initial:n = stacked
  }

\cs_new:Npn \set_of:nnn #1 #2 #3
  {
    \IfValueTF {#2}
    {
%      \msg_term:n {set_of:nn \ has \ predicates \ #2}
    }
    {
%      \msg_term:n {set_of:nn \ has \ no \ predicates}
    }
    \clist_set:Nn \l_tmpa_clist {#2}
%    \msg_term:n {l_tmpa_clist \ set}
    \tl_gset:Nn \g_tmpa_tl {\clist_use:Nn \l_tmpa_clist {\land}}
%    \msg_term:n {g_tmpa_tl \ set \ to \ \g_tmpa_tl}
    \IfValueTF {#2}
    {
      \bool_if:nTF {#3}
      {
        \scalerel*{\{}{#1\g_tmpa_tl}
        #1
%        \msg_term:n  {scalerel returns \scalerel{\g_style_set_sep_tl}{\g_tmpa_tl}}
        \clist_set:Nn \l_tmpa_clist {#2}
        \scalerel*{\mathbin{\g_style_set_sep_tl}}{#1\g_tmpa_tl}
        \\
        \clist_set:Nn \l_tmpa_clist {#2}
        \g_tmpa_tl
%         \tl_show:N \g_tmpa_tl
        \clist_set:Nn \l_tmpa_clist {#2}
        \scalerel*{\}}{#1\g_tmpa_tl}
      }
      {
        \left \{
        #1
        \clist_set:Nn \l_tmpa_clist {#2}
        \scalerel*{\mathbin{\g_style_set_sep_tl}}{#1\g_tmpa_tl}
        \g_tmpa_tl
        \right \}
      }
    }
    {
      \left \{ #1 \right \}
    }
  }

% Render underset with some variation of full
\NewDocumentCommand{\full}{O{full} m}
  {
    \show_full:nn {#1} {#2}
  }

\cs_new:Npn \show_full:nn #1 #2
  {
    \tl_set:Nn \l_tmpa_tl {#1}
    \str_if_eq:eeTF {\str_item:nn {#1} {-1}} {-}
      {
        \tl_set:Nx \l_tmpa_tl {{#1}max}
      }
      {
        \str_case:nn {#1}
        {
          {*}  {\tl_set:Nn \l_tmpa_tl {(full,S-max-full,max-full)}}
          {**} {\tl_set:Nn \l_tmpa_tl {full~(S-max-full,max-full)}}
        }
      }
    \underset {\textup{\l_tmpa_tl}} {#2}
  }

% \funcname/\funcname:n and \funcseqname/\funcseqname:n removed
% 2026-09-01, superseded by semantic-markup's \morphName/\morphSeqName
% (auto-detecting single-token/decorated/all-upper-case vs. genuine
% multi-letter word, matching this exact mathit-vs-mathrm split -- all
% call sites renamed, verified against every real argument via
% scripts/check_arg_consistency.py in that repo). See
% SEMANTIC-MARKUP-MIGRATION.md.

% Identity or inclusion function
\NewDocumentCommand \Id {O{}}
  {
     \Id:nnn {Id} {\operatorname{Id}} {#1}
  }

% Sequence of Identity or Inclusion functions
\NewDocumentCommand \ID {O{}}
  {
     \Id:nnn {ID} {\operatorname{\mathbf{Id}}} {#1}
  }

\cs_new:Npn \Id:nnn #1 #2 #3
  {
    \clist_set:Nn \l_tmpa_clist {#3}
    \int_case:nnF
      {\clist_count:N \l_tmpa_clist}
      {
        {0}
        {
         #2
        }
        {1}
        {
         #2 \sb {#3}
        }
        {2}
        {
         #2
           \sp {\clist_item:Nn \l_tmpa_clist 1}
           \sb {\clist_item:Nn \l_tmpa_clist 2}
        }
      }
      {
         \msg_error:nnnn {shmuel} {toomany} {two} {#1: [#3]}
      }
  }

    \msg_new:nnn
      {shmuel}
      {toomany}
      {More \ than \ #1  \ items \ for \  \c_backslash_str #2}

% Render underset with some variation of max-full
\NewDocumentCommand{\maxfull}{O{max-full} m}
  {
    \maximal:nnnnnnn {max-full} {#1} {#2}
                     {full,max,max-full}
                     {full,max}
                     {full~(max,max-full)}
                     {full~(max)}
  }

% Render underset with some variation of max
\NewDocumentCommand{\maximal}{O{max} m}
  {
    \maximal:nnnnnnn {max} {#1} {#2}
                     {max}
                     {}
                     {max}
                     {}
  }

\cs_new:Npn \maximal:nnnnnnn #1 #2 #3 #4 #5 #6 #7
  {
    \tl_set:Nn \l_tmpa_tl {#2}
    \str_if_eq:eeTF {\str_item:nn {#2} {-1}} {-}
      {
        \tl_set:Nx \l_tmpa_tl {{#2}#1}
      }
      {
        \str_case:nn {#2}
        {
          {*}   {\tl_set:Nn \l_tmpa_tl {(#4)}}
          {*x}  {\tl_set:Nn \l_tmpa_tl {(#5)}}
          {**}  {\tl_set:Nn \l_tmpa_tl {#6}}
          {**x} {\tl_set:Nn \l_tmpa_tl {#7}}
        }
      }
    \underset {\textup{\l_tmpa_tl}} {#3}
  }

\NewDocumentCommand{\donotbreak} {m}
    {
      \mbox{#1}\discretionary{}{}{}
    }

\NewDocumentCommand{\nobreaks} {>{\SplitList{,}}m}
    {
      \ProcessList{#1}{\donotbreak}
    }

% Render cref (nameref) on cpageref
\NewDocumentCommand \pagecrefn {m}
    {
      \pagecref:nn {\cref} {#1}
    }

\cs_new:Npn \pagecref:nn #1 #2
    {
      \bool_if:nTF
        {
          \str_if_eq_p:nn {\str_range:nnn {#2} {1} {3}} {eq:} ||
          \str_if_eq_p:nn {\nameref{#2}}                {}
        }
        {
          #1{#2} on \cpageref{#2}
        }
        {
          #1{#2} (\nameref{#2}) on \cpageref{#2}
        }
    }

% Render sequence, set or tuple name
\NewDocumentCommand{\seqname}{m}
  {
    \seqname:n {#1}
  }

\cs_new:Npn \seqname:n #1
  {
    %code here \tl_count:n
    \int_compare:nTF {\tl_count:n{#1} > 1}
      {
        {\mathbf{#1}}
      }
      {
        {\mathbfit{#1}}
      }
  }

% Render underset with some variation of strict
\NewDocumentCommand{\strict}{O{strict} m}
  {
    \show_strict:nn {#1} {#2}
  }

\cs_new:Npn \show_strict:nn #1 #2
  {
    \tl_set:Nn \l_tmpa_tl {#1}
    \str_if_eq:eeTF {\str_item:nn {#1} {-1}} {-}
      {
        \tl_set:Nx \l_tmpa_tl {{#1}strict}
      }
      {
        \str_case:nn {#1}
        {
          {*}  {\tl_set:Nn \l_tmpa_tl {(semi-strict,strict)}}
          {**} {\tl_set:Nn \l_tmpa_tl {semi-strict~(strict)}}
        }
      }
    \underset {\textup{\l_tmpa_tl}} {#2}
  }

% Render underset with some variation of bold triv
\NewDocumentCommand{\Triv}{o m}
  {
    \Triv:nn {#1} {#2}
  }

\cs_new:Npn \Triv:nn #1 #2
  {
    \IfValueTF
      {#1}
      {
        \str_case:nnF {#1}
        {
           {full-submod-} {\underset{\mathbf{{#1}triv}}{#2}}
           {full-}        {\underset{\mathbf{{#1}triv}}{#2}}
           {submod-}      {\underset{\mathbf{{#1}triv}}{#2}}
        }
        {
          {\underset{#1\mathbf{triv}}{#2}}
        }
      }
      {
         {\underset{\mathbf{triv}}{#2}}
      }
  }

\NewDocumentCommand{\intersection}{om}
  {
    \unint_of:nnn \bigcap {#1} {#2}
  }

\NewDocumentCommand{\union}{om}
  {
    \unint_of:nnn \bigcup {#1} {#2}
  }

\cs_new:Npn \unint_of:nnn #1 #2 #3
  {
    \IfValueTF {#2}
    {
%      \int_show:N \g_style_set_subscr_int
      \clist_set:Nn \l_tmpa_clist {#2}
      \int_case:nn
        {\g_style_set_subscr_int}
        {
          {0}
          {
            % Stacked subscript
%            \msg_term:n {\clist_count:N{\l_tmpa_clist} \ tokens \ stacked}
%            \clist_show:N \l_tmpa_clist
            \fp_set:Nn
              \l_tmpa_fp
              {ceil{\clist_count:N{\l_tmpa_clist} - 1} * .1 + 1}
            \scaleobj{\fp_to_decimal:N \l_tmpa_fp}{#1} \sb
            {\substack { \clist_use:Nn \l_tmpa_clist { \\ } }}
            #3
          }
          {1}
          {
            % Subscripts comma separated
%            \msg_term:n {\clist_count:N{\l_tmpa_clist} \ tokens \ comma \ separated}
%            \clist_show:N \l_tmpa_clist
            #1 \sb
            {\clist_use:Nn \l_tmpa_clist {,}}
            #3
          }
        }
    }
    {
      % No subscript
%      \msg_term:n {No~subscript}
%      \clist_show:N \l_tmpa_clist
      #1 #3
    }
  }

\ExplSyntaxOff

% |                                                                  |
% | From opening to here --- belongs in package                      |
% --------------------------------------------------------------------



\theoremstyle{definition}

\theoremstyle{remark}

\numberwithin{equation}{section}

\newcommand \Alpha   A
\newcommand \Beta    B
\newcommand \Epsilon E
\newcommand \Zeta    Z
\newcommand \Eta     H
\newcommand \Iota    I
\newcommand \Kappa   K
\newcommand \Mu      M
\newcommand \Nu      N
\newcommand \Omicron O
\newcommand \Rho     P
\newcommand \Tau     T
\newcommand \Chi     S

\usepackage[colorlinks,hidelinks,draft=false]{hyperref}

\usepackage{cleveref}

% \newtheorem must follow cleveref
\newtheorem{theorem}{Theorem}[section]
\newtheorem{corollary}[theorem]{Corollary}
\def\corollaryautorefname{Corollary}         % Needed for \autoref
\newtheorem{definition}[theorem]{Definition}
\def\definitionautorefname{Definition}       % Needed for \autoref
\newtheorem{example}[theorem]{Example}
\newtheorem{lemma}[theorem]{Lemma}
\def\lemmaautorefname{Lemma}                 % Needed for \autoref
\newtheorem{remark}[theorem]{Remark}
\newtheorem{xca}[theorem]{Exercise}

\hypersetup {
   bookmarksnumbered=true,
   colorlinks,
   pdfinfo={
      Author={Shmuel (Seymour J.) Metz},
      Keywords={fiber bundles,manifolds},
      Subject={Topology},
      Title={Local Coordinate Spaces: a proposed unification of manifolds with fiber bundles, and associated machinery}
   }
}

\usepackage[draft]{showlabels}
% \usepackage[final]{showlabels}
\showlabels{cite}
\showlabels{cref}
\showlabels{crefrange}

%\DeclareMathAlphabet{\mathbdit}{T1}{\itdefault}{\mddefault}{\sldefault}
%\SetMathAlphabet{\mathbdit}{bold}{T1}{\itdefault}{\bfdefault}{\sldefault}
%
%\DeclareMathAlphabet{\mathsfbd}{T1}{\sfdefault}{\mddefault}{\sldefault}
%\SetMathAlphabet{\mathsfbd}{bold}{T1}{\sfdefault}{\bfdefault}{\sldefault}
%
%\DeclareMathAlphabet{\mathsfit}{T1}{\sfdefault}{}{\sldefault}
%\SetMathAlphabet{\mathsfit}{normal}{T1}{\sfdefault}{}{\sldefault}
%
%\DeclareMathAlphabet{\mathsfbdit}{T1}{\sfdefault}{\mddefault}{\sldefault}
%\SetMathAlphabet{\mathsfbdit}{bold}{T1}{\sfdefault}{\bfdefault}{\sldefault}

\begin{document}
\hyphenation{near pres-en-ta-tion}

\title%
{%
  Local Coordinate Spaces:
  a proposed unification of manifolds and fiber bundles,
  and associated machinery\thanks%
  {
  I wish to gratefully thank Walter Hoffman (z"l),
  Milton Parnes, Dr. Stanley H. Levy (z"l), the Mathematics department of
  Wayne State University, the Mathematics department of the State
  University of New York at Buffalo and others who guided my education.
  }
}
\author{Shmuel (Seymour J.) Metz
}
\maketitle

% \address{4963 Oriskany Drive\\Annandale, VA  22003-5141}
% \email{smetz3@gmu.edu}
% \urladdr{http://mason.gmu.edu/~smetz3}

% \subjclass[2010]{Primary 18F15; Secondary 55R65,57N99,58A05}
% 18F15  Abstract manifolds and fiber bundles
% 32Qxx  Complex manifolds
% 53-    Differential geometry
% 53Axx  Classical differential geometry
% 53A99  None of the above, but in this section
% 55Rxx  Fiber spaces and bundles
% 55R10  Fiber bundles
% 55R65  Generalizations of fiber spaces and bundles
% 57Nxx  Topological manifolds
% 57N99  None of the above, but in this section
% 57Rxx  Differential topology
% 58Axx  General theory of differentiable manifolds
% 58A05  Differentiable manifolds, foundations
% 58Bxx  Infinite-dimensional manifolds

% Primary 18F15  Abstract manifolds and fiber bundles
% Secondary 55R65,57N99,58A05
% \keywords{fiber bundles,manifolds}

\begin{abstract}
Manifolds and fiber bundles, while superficially different, have strong
parallels; in particular, they are both defined in terms of equivalence
classes of atlases or in terms of maximal atlases, with the atlases
treated as mere adjuncts.  This paper introduces the notions of an
m-atlas and of a local coordinate space, and shows that special cases
are equivalent to fiber bundles and manifolds. It defines categories of,
e.g., atlases, and constructs functors among them.
\end{abstract}

\setupquant {parentheses=multiple,subscript=stacked}
\setupset   {}

\tikzset{
    rot90/.style={anchor=south, rotate=90, inner sep=.2mm}
}

\part {Introduction}
\label{part:intro}

Historically, manifolds were formalized in terms of atlases based on
Euclidean spaces and fiber bundles were formalized in terms of atlases
based on product spaces, using either equivalence classes of atlases or
maximal atlases.  The concept of pseudo-groups allowed unifying
manifolds and manifolds with boundary.  The definitions of fiber bundles
and manifolds have strong parallels, and can be unified in a similar
fashion; there are several ways to do so. The central part of this
paper, \pagecref{part:lcs}\, defines an approach using categories and
commutative diagrams that is designed for easy exposition at the
possible expense of abstractness and generality. In particular, I have
chosen to assume the Axiom of Choice (AOC).

This paper treats atlases as objects of interest in their own right,
although it does not give them primacy. It introduces notions that are
convenient for use here and others that, while not used here, may be
useful for future work. \Cref{sec:new} below lists the more imortant of
these.

Although this paper incidentally defines partial equivalents to
manifolds and fiber bundles using model spaces and model atlases, it
proposes the more general
\hyperref[def:LCS]{Local Coordinate Space (LCS)} in order to explicitly
reflect the role of the group in fiber bundles. \cite{Taxonomy} further
develops some of the notions used here.

A local coordinate space (LCS) is a space (total space) with some
additional structure, including a coordinate model space and an atlas
whose transition functions are restricted to morphisms of the coordinate
model space; one can impose, e.g., differentiability restrictions, by
appropriate choice of the coordinate category. There is an equivalent
paradigm that avoids explicit mention of the total space by imposing a
cocycle (compatibility) condition on the transition functions, but that
approach is beyond the scope of this paper.

This paper defines functors among categories of atlases, categories of
model spaces, categories of local coordinate spaces, categories of
manifolds and categories of fiber bundles; it constructs more machinery
than is customary in order to facilitate the presentation of those
categories and functors.

\begin{remark}
The definitions of morphisms given here are nonstandard, but there are
also definitions of classical morphisms that are much closer to the
standard ones.
\end{remark}

\Crefrange{part:conv}{part:pre} present nomenclature and give basic
results.

\Cref{part:M-charts} defines m-atlases, m-atlas morphisms and categories
of them;
{
  \showlabelsinline
  \cref{lem:M-ATLiscat}
}
proves that the defined categories are indeed categories.

\Cref{part:lcs} defines local coordinate spaces and
categories of them; \pagecref{the:LCSiscat} proves that the defined
categories are indeed categories.
\Pagecref{Examples} gives some examples of structures that can be
represented as local coordinate spaces.

\pagecref{part:man} and
\pagecref{part:bun} present two examples of an LCS in detail, showing
the equivalence of manifolds and fiber bundles with special cases of
local coordinate spaces by explicitly exhibiting functors to and from
local coordinate spaces.
\begin{remark}
The unconventional definitions of manifold and fiber bundle morphisms
are intended to make their relationship to local coordinate spaces more
natural.
\end{remark}

Many of the lemmata, theorems and corollaries in this paper are
substantially identical to results that may be familiar to the reader.
What is novel is the perspective and the material directly related to
local coordinate spaces. The presentation assumes only a basic knowledge
of Category Theory, such as may be found in the first chapter of
\cite{CftWM} or {
  \showlabelsinline
  \cite{JoyCat}.
}

Many of the proofs are obvious and may be omitted from other versions of
this paper.

\section{New concepts and notation}
\label{sec:new}
This paper introduces a significant number of new concepts and some
modifications of the definitions for some conventional concepts. It also
introduces some notation of lesser importance.  The following are the
most important.

\begin{enumerate}
\item \hyperref[def:model]{Model space} and related concepts.
A \hyperref[def:model]{model space}\footnote{
  The phrase has been used before, but with a different meaning.
}
is a topological space with a category of permissible open sets and
mappings satisfying specified conditions.  It is similar to a
pseudo-group, with some important differences.

\item \hyperref[def:ModTop]{Model topology} and
\hyperref[def:M-para]{M-paracompactness}

\item \hyperref[sec:sig]{Signature},
\hyperref[def:Sigmacomm]{$\Sigma$-commutation} and related concepts.
A signature is a means of characterising the domains and range of a
multi-variable function.
$\Sigma$-commutation is a natuality condition for a sequence of functions.

\item
An \hyperref[def:m-atlas]{m-atlas} is a generalization of atlases in
manifolds.

\item
An \hyperref[def:M-ATLmorph]{m-atlas morphism} is a generalization of a
$\Ck$ map between differentiable manifolds. It differs in using both a
function between the manifolds and a function between the coordinate
spaces.

\item
\hyperref[def:LCS]{Local Coordinate Space (LCS)} and related concepts

\item
\hyperref[def:lin]{Linear space and related concepts}
A linear space is a subspace of a Banach space or of a Fr\'echet space
that is suitable for defining coordinate patches in a manifold.

\item
\hyperref[def:trivck]{Trivial $\Ck$ linear model space} and related
concepts.
A trivial $\Ck$ linear model space is a maximal model space
with morphisms $\Ck$ maps between open subsets of a linear space.

\item \hyperref[def:BunAtl]{$G$-$\rho$ bundle atlases}%
\footnote{Similar to coordinate bundles}
and related concepts
\end{enumerate}

\part {Conventions}
\label{part:conv}
An arrow with an Equal-Tilde ($A \toiso_\phi B$) represents an
isomorphism. One with a hook ($A \underset{i}{\hookrightarrow} B$)
represents an inclusion map. One with a tail ($A \morphMono_i B$) represents
a monomorphism. One with a double head ($A \morphEpi_\pi B$) represents a
surjection.

When a superscript ends in $-1$, e.g., $\morphName{f}^{i-1}$, it is to be
taken as function inverse rather than subtraction of 1.

All diagrams shown are commutative; none are exact.

A definition of a base term and several related terms may specify
modifiers in parenthese in the base definition and then give the
definitions for each modifier; usually the modifiers will be
restrictions and the base definition will not be repeated.  E.g.,
{\textquotedblleft}$\morphSeqName{f}$ is a (semi-strict, strict)
prestructure morphism of $\seqname{P}^1$ to $\seqname{P}^2$
iff{\textquotedblright} is followed by the definition of prestructure
morphism, then by the restrictions for strict and semi-strict
prestructure morphisms.

When a definition defines a base propositional function and variant
propositional functions with a qualifier given as a superscript or
underset, then the form $\mathrm{base}^{(qualifiers)}$ will refer to
either $\mathrm{base}^{qualifier}$ or $\mathrm{base}$ and the form
$\underset{(qualifiers)}{base}$ will refer to either
$\underset{qualifier}{base}$ or $base$, e.g.,
$\strict[*]{\isAtl[(classic)]_\Ar}$ refers to either $\isAtl_\Ar$,
$\isAtl[classic]_\Ar$, $\strict[semi-]{\isAtl_\Ar}$,
$\strict[semi-]{\isAtl[classic]_\Ar}$, $\strict{\isAtl_\Ar}$ or
$\strict{\isAtl[classic]_\Ar}$.

The form $\mathrm{base}^{qualifier(,qualifiers)}$ will refer to either
$\mathrm{base}^{\mathrm{qualifier}}$ or to the variant with one or more
of the additional qualifiers in the superscript, e.g.,
$\strict{\isAtl[classic(,proper)]_\Ar}$ refers to either
$\strict{\isAtl[classic]_\Ar}$ or
$\strict{\isAtl[classic,proper]_\Ar}$.

Alternatively, a definition may specify a numbered list of alternatives,
and subsequently specify additional numbered lists with items
corresponding to those in the first list, e.g.,
$\morphSeqName{f}$ is also a
\begin{enumerate}
\item full
\item semi-maximal
\item maximal
\item full semi-maximal
\item full maximal
\end{enumerate}
$E^1$-$E^2$ $\Ck$ morphism of
$\seqname{A}^1$ to $\seqname{A}^2$ in the coordinate spaces $C^1$,
$C^2$, abbreviated as
\begin{enumerate}
\item abbreviation for full
\item abbreviation for semi-maximal
\item abbreviation for maximal
\item abbreviation for full semi-maximal
\item abbreviation for full maximal
\end{enumerate}
iff
\begin{enumerate}
\item definition for full
\item definition for semi-maximal
\item definition for maximal
\item definition for full semi-maximal
\item definition for full maximal
\end{enumerate}

A Corollary, Lemma or Theorem that applies to related terms defined
with the above convention will specify the modifiers in parentheses
to indicate that it applies to the base term and to the variant
terms, e.g., {\textquotedblleft}a (semi-strict, strict)
$\seqname{E}^i$-$\seqname{E}^{i+1}$ m-atlas morphism{\textquotedblright}
If it applies only to more restrictive terms then it will specify the
first relevant modifier followed by the remaining relevant modifiers in
parentheses, e.g., {\textquotedblleft}If
$\catName{S}^i \SUBCAT[full-] \catName{S}'^i$ and $\morphSeqName{f}^i$ is
a semi-strict (strict) prestructure morphism{\textquotedblright}.

Blackboard bold upper case will denote specific sets, e.g., the
Naturals.

Bold lower case italic letters will refer to sets, sequences and tuples
of functions, e.g.,
$\morphSeqName{f} \defeq (\morphName{f}_1, \morphName{f}_2)$.

Bold lower case Latin letters will refer to sequence valued functions of
sequences and tuple valued functions of tuples, e.g., $\Range$ yields
the sequence of ranges of a sequence of functions.

Bold upper case italic letters will refer to sequences or tuples, e.g.,
$\seqname{A}=(x,y,z)$, to sets of them, to sets of topological spaces or
to sets of open sets.

Bold upper case script letters will refer to
sequences of categories, e.g., \\*
\nobreaks
  {{
    $\catSeqName{A} \defeq (\catName{A}_\alpha, \alpha \in \Alpha)$.
  }}

Fraktur will refer to topologies and to topology-valued functions, e.g.,
$\Topology$.

Functions have a range, domain and relation, not just a relation. Unless
otherwise stated, they are assumed or proven to be continuous.

Groups are assumed to be topological groups. The ambiguous notation
$x^{-1}$ will be used when it is obvious from context what the group
operation $\star$ and the group identity $\mathbf{1}_G$ are.

Lower case Greek letters other than $\pi$, $\rho$, $\sigma$, $\phi$ and
$\psi$ will refer to ordinals, possibly transfinite, and to formal
labels.  A letter with a Greek superscript and a letter with a Latin or
numeric superscript always refer to distinct variables.

Lower case $\pi$ will refer to a projection operator

Lower case $\rho$ will refer to a continuous effective group action,
i.e., a continuous representation of a group in a homeomorphism group.

Lower case $\sigma$ will refer to a sequence of ordinals, referred to as
a signature.

Lower case $\phi$ will refer to a coordinate function.

Lower case italic and Latin letters will refer to
\begin{enumerate}
\item elements of a set or sequence
\item formal labels \\*
When the letters $u$ and $v$, with or without a superscript, refer to an
element of a set then they may also be used as formal labels for sets
associated with that element.
\item functions
\item morphisms of a category
\item natural numbers
\item objects of a category
\item ordinal numbers
\end{enumerate}

Upper case Greek letters other than $\Sigma$ may refer to
\begin{enumerate}
\item ordinal used as the limit of a sequence of consecutive ordinals,
e.g., $x_\alpha, \alpha \preceq \Alpha$
\item ordinal used as the order type of a sequence of consecutive
ordinals, e.g., $x_\alpha, \alpha \prec \Alpha$
\end{enumerate}

Upper case $\Sigma$ will refer to a sequence of signatures

Upper case Latin letters will refer to
\begin{enumerate}
\item Natural numbers
\item Topological spaces
\item Open sets
\item Formal labels for elements of a sequence or tuple of functions,
e.g., $\morphName{f}_E$ might be $\morphName{f}_0 \maps E_1 \to E_2$.
\end{enumerate}

Upper case Script Latin letters will refer to categories and functors.

Upright Latin letters will be used for long names.
In particular,
\begin{enumerate}
\item
$\Ar$, $\Hom$, $\Mor$ and $\Ob$ are the conventional set-valued
functions on small categories:

\begin{description}
\item[$\Ar(\catName{C})$]
All morphisms of $\catName{C}$
\item[$\Hom_{\catName{C}}(a,b)$]
All morphisms of $\catName{C}$ from $a$ to $b$
\item[$\Mor(\catName{C})$]
All morphisms of $\catName{C}$.
Synonymous with $\Ar(\catName{C})$.
\item[$\Ob(\catName{C})$]
All objects of $\catName{C}$
\end{description}

They will never be applied to large categories in this paper.

\item
$\catName{Set}$ and $\catName{Top}$ are the conventional categories of sets and
topological spaces.

\item
The letters E, G, X and Y will be used for formal labels.
\end{enumerate}

The term $\Ck$ includes $\mathrm{C}^\infty$ (smooth) and
$\mathrm{C}^\omega$ (analytic).

This paper uses the term morphism in preference to arrow, but uses
the conventional $\Ar$.

This paper uses the terms empty, null and void interchangeably.

The term sequence without an explicit reference to $\mathbb{N}$
will refer to a general ordinal sequence, possibly transfinite.

Sequence numbering, unlike tuple numbering, starts at 0 and the
exposition assumes a von Neumann definition of ordinals, so that
$\alpha \in \beta \equiv \alpha \prec \beta$.

Except where explicitly stated otherwise, all categories mentioned are
small categories with underlying sets, but the morphisms will often not
be set functions between the objects and there will not always be a
forgetful functor to $\catName{Set}$ or $\catName{Top}$. By abuse of language no
distinction will consistently be made between a category $\catName{A}$
of topological spaces and the concrete category
$(\catName{A},\catName{U})$ over $\catName{Top}$.  Similarly, no distinction
will consistently be made among the object $U \in \Ob(\catName{A})$, the
topological space $\catName{U}(U)$ and the underlying set.

The notation $G^V$ will refer only to the set of continuous functions
from $V$ to $G$, never to the set of all functions from $V$ to $G$.

When defining a category, the ordered pair $(\seqname{O}, \seqname{M})$
refers to the smallest concrete category over $\catName{Set}$ or $\catName{Top}$
whose objects are the elements in $\seqname{O}$, whose morphisms include
all of the elements of $\seqname{M}$ and whose morphisms from
$o^1 \in \seqname{O}$ to $o^2 \in \seqname{O}$ are functions
$\morphName{f} \maps o^1 \to o^2$ and whose composition is function
composition.

When defining a category, the ordered triple
$(\seqname{O}, \seqname{M}, C)$ refers to the small category whose
objects are in $\seqname{O}$, whose morphisms are in $\seqname{M}$,
whose $\Hom$ is

\begin{equation}
\Hom_{(\seqname{O}, \seqname{M}, C)}
  (o_1 \in \seqname{O}, o_2 \in \seqname{O})
\defeq
\set
  {
    (
      \morphSeqName{f},
      o_1,
      o_2
    )
    \in \seqname{M}
  }
\end{equation}
and whose composition is C.

By abuse of language I may write
``$\catName{S}$'' for $\Ob(\catName{S})$,
``$A \in \catName{A}$'' for  $A \in\ \Ob(\catName{A})$,
``$A \subset \catName{A}$'' for $A \subset \Ob(\catName{A})$,
``$A \in \catName{A} \subset B \in \catName{B}$'' for
``the underlying set of $A$ is contained in the underlying set of $B$
and the inclusion $\morphName{i} \maps x \in A \hookrightarrow x \in B$
is a morphism'' and
``$\morphName{f} \maps A \to B$'' for
$\morphName{f} \in \Hom_{\catName{C}}(A,B)$, where $\catName{C}$ is
understood by context.

By abuse of language I shall use the same nomenclature for sequences and
tuples, and will use parentheses around a single expression both for
grouping and for a tuple with a single element; the intent should be
clear from context.

By abuse of language I will omit parenthese around the operands of
Functors when they can be assumed by context.

By abuse of language I shall use the $\times$ and $\bigtimes$ symbols
for both Cartesian products of sets and Cartesian products of
functions on those sets.

By abuse of language, and assuming AOC,  I shall refer to some sets as
ordinal sequences, e.g., ``$(C_\alpha, \alpha \in \Alpha)$'' for
``$\{C_\alpha \mid \alpha \in \Alpha\}$'', in cases where the order is
irrelevant.

Where a proof requires proving intermediate propositions, a bold leadin
to a paragraph will indicate such a proposition, e.g.

\begin{description}
\item[Intermediate proposition 1:] Intermediate proof 1.

\item[Intermediate proposition 2:]
\leavevmode
Intermediate proof 2.
\end{description}

By abuse of language, I may omit universal quantifiers in cases where
the intent is clear.

In some cases I define a notion similar to a conventional notion and
also need to refer to the conventional notion. In those cases I prefix
a letter or phrase to the term, e.g., m-paracompact versus paracompact.

\part {General notions}
\label{part:notions}
This part of the paper describes nomenclature used throughout the paper.
In some cases this reflects new nomenclature or notions, in others it
simply makes a choice from among the various conventions in the
literature.

\begin{definition}[Operations on categories]
\label{def:catprop}
If $\catName{C}$ is a category then
\begin{enumerate}
\item
$x \objin \catName{C} \equiv x \in \Ob(\catName{C})$

\item
$x \arin \catName{C} \equiv x \in \Ar(\catName{C})$

\item
If $\catName{C}$ is a concrete category over $\catName{Set}$ or $\catName{Top}$
and $o \objin \catName{C}$, then $\Set_{\catName{C}}(o)$ is the underlying
set of $o$. By abuse of language we will write $\Set(o)$ when the
category is understood from context.

\item
If $\catName{C}$ is a concrete category over $\catName{Top}$ and
$o \objin \catName{C}$, then $\Top_{\catName{C}}(o)$ is the underlying
topological space of $o$. By abuse of language we will write $\Top(o)$
when the category is understood from context.
\end{enumerate}

If $\catName{S}$ and $\catName{T}$ are categories then
$\catName{S} \subcat \catName{T}$ iff $S$ is a subcategory of
$\catName{T}$ and $\catName{S} \subcat[full-] \catName{T}$ iff $S$ is a
full subcategory of $\catName{T}$.

If $\catName{S}$ and $\catName{T}$ are categories then the category
union of $\catName{S}$ and $\catName{T}$, abbreviated \\*
$\catName{S} \unioncat \catName{T}$, is the category whose objects are
in $\catName{S}$ or in $\catName{T}$ and whose morphisms are in
$\catName{S}$ or in $\catName{T}$.
\end{definition}

\begin{definition}[Identity]
\label{def:Id}
$\Id[S]$ is the identity function on the space $S$,

$\Id[o]$ is the
identity morphism for the object $o$\footnote{
The object is often expressed as a tuple, e.g.,
$\Id[{{(\seqname{A}, \seqname{B})}}]$ is the identity morphism for the
object $(\seqname{A}, \seqname{B})$
},

$\Id[U,V]$, for $U \subseteq V$, is the inclusion map\footnote{
$U$ need not have the subspace topology.
}.
$\Id[{{U,V}}]$ is synonymous with $\Id[U,V]$.

$\Id[\catName{C}]$ is the identity functor on the category $\catName{C}$.

$\ID[\seqname{S}^i]$, $i=1,2$, is the sequence of identity functions
for the elements of the sequence
$\seqname{S}^i \defeq (\seqname{S}^i_\alpha,~ \alpha \prec \Alpha)$. Let
$\seqname{S}^1 \SUBSETEQ \seqname{S}^2$. Then
$\ID[\seqname{S}^1,\seqname{S}^2]$ is the sequence of inclusion maps
$
  \biggl (
    \Id[\seqname{S}^1_\alpha,\seqname{S}^2_\alpha],~\alpha \prec \Alpha
  \biggr )
$
for the elements of the sequences $\seqname{S}^i$.
$\ID[{{\seqname{S}^1,\seqname{S}^2}}]$ is synonymous with
$\ID[\seqname{S}^1,\seqname{S}^2]$.  A similar definition applies for
tuples, except in a context where they are treated as objects of a
category.

$\ID[{{(A \dots Z)}}]$ is the tuple of identity maps
$\bigl (\Id[A], \dots, \Id[Z] \bigr)$.
$\ID[{(A^1, \dots, Z^1)},{(A^2, \dots, Z^2)}]$ is the tuple of inclusion maps
$\biggl ( \Id[A^1, A^2], \dots, \Id[Z^1, Z^2] \biggr )$.
\end{definition}

\begin{definition}[Images]
$\morphName{f} [U] \defeq \set { {\morphName{f}(x)} }[x \in U]$ is the
image of $U$ under $\morphName{f}$ and \\*
$\morphName{f}^{-1} [V] \defeq \set {x}[\morphName{f}(x) \in V]$ is the
inverse image of $V$ under $\morphName{f}$.

\begin{remark}
  This notation, adopted from \cite{GenTop}, avoids the ambiguity in
  the traditional $\morphName{f}(U)$ and $\morphName{f}^{-1}(V)$.
\end{remark}
\end{definition}

\begin{definition}[Projections]
\label{projections}
$\pi_\alpha$ is the projection function that maps a sequence into
element $\alpha$ of the sequence.  $\pi_i$ is also the projection
function that maps a tuple into element $i$ of the tuple.
\end{definition}

\begin{definition}[Topological category]
\label{def:topcat}
A topological category is a small subcategory of $\catName{Top}$ or its
concrete category over $\catName{Set}$.

$\catName{T}$ is a full topological category iff it is a topological
category and whenever $U^i,V^i \objin \catName{T}$, $i=1,2$,
$V^i \subsp U^i$,
$\morphName{f} \maps U^1 \to U^2 \arin \catName{T}$ and
$\morphName{f}[V^1] \subseteq V^2$ then \\
$\morphName{f} \maps V^1 \to V^2 \arin \catName{T}$.
\end{definition}

\begin{lemma}[Inclusions in topological categories are morphisms]
\label{lem:topInc}
Let $\catName{T}$ be a full topological category,
$S^i \objin \catName{T}$, $i=1,2$, and $S^1 \subsp S^2$. Then
$\Id[S^1,S^2]$ is a morphism of $\catName{T}$

\begin{proof}
\begin{enumerate}
\item
$U^i \defeq S^2 \objin \catName{T}$, $i=1,2$, by hypothesis.

\item
$V^1 \defeq S^1 \objin \catName{T}$ and
$V^2 \defeq S^2 \objin \catName{T}$ by hypothesis.

\item
$V^i \subsp U^i$,

\item
$\morphName{f} \defeq \Id[S^2] \maps U^1 \to U^2 \arin \catName{T}$.

\item
$\morphName{f}[V^1] = V^1 \subsp V^2$.
\end{enumerate}

Thus $\Id[S^1,S^2] = \morphName{f} \maps V^1 \to V^2 \arin \catName{T}$
by \cref{def:topcat}.
\end{proof}
\end{lemma}

\begin{definition}[Local morphisms]
\label{def:topLocal}
Let $\catName{S}^i$, $i=1,2$, be a full topological category and
$S^i \objin \catName{S}^i$, A continuous function
$\morphName{f} \maps S^1 \to S^2$ is a local
$\catName{S}^1$-$\catName{S}^2$ morphism of $S^1$ to $S^2$ iff
$\catName{S}^1 \subcat[full-] \catName{S}^2$ and for every
$u \in S^1$ there is an open neighborhood $U_u$ for $u$ and an open
neighborhood $V_u$ for $v \defeq \morphName{f}(u)$ such that
$\morphName{f}[U_u] \subseteq V_u$ and $\morphName{f} \maps U_u \to V_u$
is a morphism of $\catName{S}^2$.

\begin{remark}
The condition $\morphName{f} \maps U_u \to V_u \arin \catName{S}^2$
ensures that $U_u \objin \catName{S}^1$ and $V_u \objin \catName{S}^2$
\end{remark}

Let $\catName{T}$ be a full topological category and
$S^i \objin \catName{T}$, $i=1,2$. A continuous function
$\morphName{f} \maps S^1 \to S^2$ is a local $\catName{T}$ morphism of
$S^1$ to $S^2$ iff it is a local $\catName{T}$-$\catName{T}$ morphism
of $S^1$ to $S^2$.
\end{definition}

\begin{lemma}[Local morphisms]
\label{lem:topLocal}
Let $\catName{T}^i$, $i \in [1,3]$, be a full topological category,
$S^i \objin \catName{T}^i$ and
$\catName{T}^i \subcat[full-] \catName{T}^{i+1}$.

If $\morphName{f}^i \maps S^i \to S^{i+1} \arin \catName{T}^{i+1}$ then
$\morphName{f}^i$ is a local $\catName{T}^i$-$\catName{T}^{i+1}$
morphism of $S^i$ to $S^{i+1}$.

\begin{proof}
Let $u \in S^i$ and $v \defeq \morphName{f}^i(u) \in S^{i+1}$. $S^i$ is
an open for $u$, $S^{i+1}$ is an open neighborhood of $v$ and
$\morphName{f}^i \maps S^i \to S^{i+1} \arin \catName{T}^{i+1}$ by
hypothesis.
\end{proof}

If each $\morphName{f}^i \maps S^i \to S^{i+1}$, is a local
$\catName{T}^i$-$\catName{T}^{i+1}$ morphism of $S^i$ to
$S^{i+1}$ then
$\morphName{f}^2 \morphCompose \morphName{f}^1 \maps S^1 \to S^3$ is a local
$\catName{T}^1$-$\catName{T}^3$ morphism of $S^1$ to $S^3$.

\begin{proof}
Since $\catName{T}^1 \subcat[full-] \catName{T}^2$ and
$\catName{T}^2 \subcat[full-] \catName{T}^3$,
$\catName{T}^1 \subcat[full-] \catName{T}^3$.
Let $u \in S^1$, $v \defeq \morphName{f}^1(u)$ and
$w \defeq \morphName{f}^2(v)$. There exist an open neighborhood $U_u$
for $u$, open neighborhoods $V_u$, $V'_v$ for $v$ and an open
neighborhood $W_v$ of $w$ such that
$\morphName{f}^1[U_u] \subseteq V_u$,
$\morphName{f}^1 \maps U_u \to V_u$ is a morphism of $\catName{T}^2$,
$\morphName{f}^2[V'_v] \subseteq W_v$ and
$\morphName{f}^2 \maps V'_v \to W_v$ is a morphism of $\catName{T}^3$.
Then $\hat{V_u} \defeq V_u \cap V'_v \neq \emptyset$, $\hat{V_u}$ is an
open neighborhood of $v$ and
$\hat{U_u} \defeq \morphName{f}^{i-1}_1[\hat{V_u}]$ is an open
neighborhood of $u$.  $\morphName{f}^1 \maps \hat{U_u} \to \hat{V_u}$
and $\morphName{f}^2 \maps \hat{V_u} \to W_v$ are morphisms of
$\catName{T}^3$ by \pagecref{def:topcat} and thus
$\morphName{f}^2 \morphCompose \morphName{f}^1 \maps \hat{U_u} \to W_v$ is a
morphism of $\catName{T}^3$.
\end{proof}
\end{lemma}

\begin{corollary}[Local morphisms]
\label{cor:topLocal}
Let $\catName{T}^i$, $i=1,2$, be a full topological category,
$\catName{T}^i \subcat[full-] \catName{T}^{i+1}$,
$S^i \objin \catName{T}^i$ and $S^1 \subseteq S^2$. Then $\Id[S^1,S^2]$
is a local $\catName{T}^1$-$\catName{T}^2$ morphism of $S^1$ to $S^2$
and $\Id[{S^i}]$ is a local $\catName{T}$ morphism
of $S^i$ to $S^i$.

\begin{proof}
$S^1 \objin \catName{T}^2$ because $S^1 \objin \catName{T}^1$ and
$\catName{T}^1 \subcat \catName{T}^2$,
$S^2 \objin \catName{T}^2$ by hypothesis and \\*
$S^1 \subseteq S^2$ by hypothesis, so
$\Id[S^1,S^2] \arin \catName{T}^2$ by \cref{lem:topInc}.

$\Id[S^i] \defeq \Id[S^i,S^i]$.
\end{proof}
\end{corollary}

\begin{definition}[Sequence functions]
Let $\seqname{S} \defeq (s_\alpha, \alpha \prec \Alpha)$ be a sequence
of functions. Then
\begin{equation}
\Domain(\seqname{S}) \defeq \bigl ( \domain(s_\alpha), \alpha \prec \Alpha \bigr )
\end{equation}
\begin{equation}
\Range(\seqname{S}) \defeq \bigl ( \range(s_\alpha), \alpha \prec \Alpha \bigr )
\end{equation}

Let $\seqname{T} \defeq (t_\alpha, \alpha \prec \Alpha)$ be a sequence
of functions with $\Range(\seqname{S}) = \Domain(\seqname{T})$. Then
their composition is the sequence
$
  \seqname{T} \morphCompose[()] \seqname{S} \defeq
  (t_\alpha \morphCompose s_\alpha, \alpha \prec \Alpha)
$,

Let $\seqname{S} \defeq (s_\gamma, \gamma \preceq \Gamma)$, then these
functions extract information about the sequence:
\begin{equation}
\head(\seqname{S},\Omega) \defeq (s_\gamma, \gamma \prec \Omega)
\end{equation}
\begin{equation}
\head(\seqname{S}) \defeq \head(\seqname{S},\Gamma)
\end{equation}
\begin{equation}
\lengtho(\seqname{S}) \defeq \Gamma
\end{equation}
\begin{equation}
\tail(\seqname{S}) \defeq S_\Gamma
\end{equation}

Let $\seqname{S} \defeq (s_\gamma, \gamma \prec \Gamma)$, then
\begin{equation}
\length(\seqname{S}) \defeq \Gamma
\end{equation}

\begin{remark}
If $\lengtho(\seqname{S})$ is defined then
$\length(\seqname{S}) = \lengtho(\seqname{S}) + 1$.
$\lengtho(\seqname{S})$ is the ordinal type of
$\head(\seqname{S})$, not the ordinal type of $\seqname{S}$.
\end{remark}

Let $\catSeqName{S} \defeq (\catName{S}_\alpha, \alpha \prec \Alpha)$
and $\catSeqName{T} \defeq (\catName{T}_\alpha, \alpha \prec \Alpha)$ be
sequences of categories.  Then $\catSeqName{S}$ is a subcategory
sequence of $\catSeqName{T}$, abbreviated
$\catSeqName{S} \SUBCAT \catSeqName{T}$, iff every category in
$\catSeqName{S}$ is a subcategory of the corresponding category in
$\catSeqName{T}$, i.e.,
$
  \uquant%
    {\alpha \prec \Alpha}
    {\catName{S}_\alpha \subcat \catName{T}_\alpha}
$,
and $\catSeqName{S}$ is a full subcategory sequence of $\catSeqName{T}$,
abbreviated $\catSeqName{S} \SUBCAT[full-] \catSeqName{T}$, iff every
category in $\catSeqName{S}$ is a full subcategory of the corresponding
category in $\catSeqName{T}$, i.e.,
$
  \uquant%
    {\alpha \prec \Alpha}
    {\catName{S}_\alpha \subcat[full-] \catName{T}_\alpha}
$.

The category sequence union of $\catSeqName{S}$ and $\catSeqName{T}$,
abbreviated $\catSeqName{S} \UNIONCAT \catSeqName{T}$, is the sequence
of category unions of corresponding categories in $\catSeqName{S}$ and
$\catSeqName{T}$, i.e.,
$(\catSeqName{S}_\alpha \unioncat \catSeqName{T}_\alpha)$.

\end{definition}

\begin{lemma}[Sequence functions]
\label{lem:seqfunc}
Let $\morphSeqName{f}^i \defeq (\morphName{f}^i_\alpha, \alpha \prec \Alpha)$,
$i \in [1,3]$, be sequences of functions with
\donotbreak
  {
    $
      \Domain(\morphSeqName{f}^2) = \Range(\morphSeqName{f}^1)
    $
  }
and
\donotbreak
  {
    $
      \Domain(\morphSeqName{f}^3) = \Range(\morphSeqName{f}^2)
    $.
  }
Then
$
(\morphSeqName{f}^3 \morphCompose[()] \morphSeqName{f}^2) \morphCompose[()] \morphSeqName{f}^1 =
\morphSeqName{f}^3 \morphCompose[()] (\morphSeqName{f}^2 \morphCompose[()] \morphSeqName{f}^1)
$.
\begin{proof}
\begin{equation*}
\begin{split}
  (\morphSeqName{f}^3 \morphCompose[()] \morphSeqName{f}^2) \morphCompose[()]
  \morphSeqName{f}^1
  & =
  \bigl (
    (\morphName{f}^3_\alpha \morphCompose \morphName{f}^2_\alpha) \morphCompose
    \morphName{f}^1_\alpha,
    \alpha \prec \Alpha
  \bigr )
\\*
  & =
  \bigl (
    \morphName{f}^3_\alpha \morphCompose
    (\morphName{f}^2_\alpha \morphCompose \morphName{f}^1_\alpha),
    \alpha \prec \Alpha
  \bigr )
\\*
  & =
  \morphSeqName{f}^3 \morphCompose[()] (\morphSeqName{f}^2 \morphCompose[()] \morphSeqName{f}^1)
\end{split}
\end{equation*}
\end{proof}

Let $\morphSeqName{f} \defeq (\morphName{f}_\alpha, \alpha \prec \Alpha)$
be a sequence of functions, $\seqname{D} = \Domain(\morphSeqName{f})$ and \\*
$\seqname{R} = \Range(\morphSeqName{f})$. Then $\ID[\seqname{R}]$ is a
left $\morphCompose[()]$ identity for $\morphSeqName{f}$ and
$\ID[\seqname{D}]$ is a right $\morphCompose[()]$ identity for
$\morphSeqName{f}$.

\begin{proof}
\begin{equation*}
\begin{split}
  \ID[\seqname{R}] \morphCompose[()] \morphSeqName{f}
  & =
  (\ID[\range(\morphName{f}_\alpha)] \morphCompose \morphName{f}_\alpha,
  \alpha \prec \Alpha)
\\*
  & =
  (\morphName{f}_\alpha, \alpha \prec \Alpha)
\\*
  & =
  \morphSeqName{f}
\end{split}
\end{equation*}
\begin{equation*}
\begin{split}
  \morphSeqName{f} \morphCompose[()] \ID[\seqname{D}]
  & =
  (\morphName{f}_\alpha \morphCompose \Id[\domain(\morphName{f}_\alpha)],
  \alpha \prec \Alpha)
\\*
  & =
  (\morphName{f}_\alpha, \alpha \prec \Alpha)
\\*
  & =
  \morphSeqName{f}
\end{split}
\end{equation*}
\end{proof}
\end{lemma}

\begin{definition}[Tuple functions]
Let $\seqname{S} \defeq (s_n, n \in [1,N])$ be a tuple of functions. Then
\begin{equation}
\Domain(\seqname{S}) \defeq \bigl ( \domain(s_n), n \in [1,N] \bigr )
\end{equation}
\begin{equation}
\Range(\seqname{S}) \defeq \bigl ( \range(s_n), n \in [1,N] \bigr )
\end{equation}

Let $\seqname{T} \defeq (t_n, n \in [1,N])$ be a tuple of functions with
$\Range(\seqname{S})=\Domain(\seqname{T})$,
Then their composition is the tuple
$
  \seqname{T} \morphCompose[()] \seqname{S} \defeq
  (t_n \morphCompose s_n, n \in [1,N])
$

Let $\seqname{S} \defeq (s_m, m \in [1,M])$ and
$\seqname{T} \defeq (t_n, n \in [1,N])$ be tuples.
Then the following are tuple functioms
\begin{equation}
\head(S,I) \defeq (s_m, m \in [1,I])
\end{equation}
\begin{equation}
\head(S) \defeq \head(S,M-1)
\end{equation}
\begin{equation}
\tail(T,I) \defeq (t_n, n \in [I,N])
\end{equation}
\begin{equation}
\tail(T) \defeq t_N
\end{equation}
\begin{equation}
\join(S,T) \defeq (s_1, \dots, s_M, t_1, \dots, t_N)
\end{equation}
\end{definition}

\begin{lemma}[Tuple functions]
\label{lem:tupfunc}
Let $\morphSeqName{f} \defeq (\morphName{f}_n, n \in [1,N])$ be a tuple of
functions,
$\seqname{D} \defeq \Domain(\morphSeqName{f})$ and
$\seqname{R} \defeq \Range(\morphSeqName{f})$. Then $\ID[\seqname{R}]$ is a
left $\morphCompose[()]$ identity for $\morphSeqName{f}$ and $\ID[\seqname{D}]$
is a right $\morphCompose[()]$ identity for $\morphSeqName{f}$.

\begin{proof}
\begin{equation*}
\begin{split}
  \ID[\seqname{R}] \morphCompose[()] \morphSeqName{f}
  & =
  (\Id[\range(\morphName{f}_\alpha)] \morphCompose \morphName{f}_n, n \in [1,N])
\\*
  & =
  (\morphName{f}_n, n \in [1,N])
\\*
  & =
  \morphSeqName{f}
\end{split}
\end{equation*}
\begin{equation*}
\begin{split}
  \morphSeqName{f} \morphCompose[()] \ID[\seqname{D}]
  & =
  (\morphName{f}_n \morphCompose \Id[\domain(\morphName{f}_n)], n \in [1,N])
\\*
  & =
  (\morphName{f}_n, n \in [1,N])
\\*
  & =
  \morphSeqName{f}
\end{split}
\end{equation*}
\end{proof}
\end{lemma}

\begin{lemma}[Tuple composition for unlabeled morphisms]
\label{lem:tupcomp}
Let \\*
$
  \morphSeqName{f}^i \defeq (\morphName{f}^i_n,~ n \in [1,N]),~
  i \in [1,3]
$,
be tuples of functions with
\donotbreak
  {
    $
      \Domain(\morphSeqName{f}^2) = \Range(\morphSeqName{f}^1)
    $
  }
and
\donotbreak
  {
    $
      \Domain(\morphSeqName{f}^3) = \Range(\morphSeqName{f}^2)
    $.
  }
Then
$
(\morphSeqName{f}^3 \morphCompose[()] \morphSeqName{f}^2) \morphCompose[()] \morphSeqName{f}^1 =
\morphSeqName{f}^3 \morphCompose[()] (\morphSeqName{f}^2 \morphCompose[()] \morphSeqName{f}^1)
$.

\begin{proof}
\begin{equation*}
\begin{split}
  (\morphSeqName{f}^3 \morphCompose[()] \morphSeqName{f}^2) \morphCompose[()] \morphSeqName{f}^1
  & =
  \bigl ( (\morphName{f}^3_n \morphCompose \morphName{f}^2_n) \morphCompose \morphName{f}^1_n, n \in [1,N] \bigr )
\\*
  & =
  \bigl ( \morphName{f}^3_n \morphCompose (\morphName{f}^2_n \morphCompose \morphName{f}^1_n), n \in [1,N] \bigr )
\\*
  & =
  \morphSeqName{f}^3 \morphCompose[()] (\morphSeqName{f}^2 \morphCompose[()] \morphSeqName{f}^1)
\end{split}
\end{equation*}
\end{proof}
\end{lemma}

\begin{definition}[Tuple composition for labeled morphisms]
\label{def:labcomp}
Let $\seqname{M}^i \defeq (\morphSeqName{f}^i, o^i_1, o^i_2)$, $i=1,2$,
be tuples such that each $\morphSeqName{f}^i$ is a sequence of functions
or each $\morphSeqName{f}^i$ is a tuple of functions,
$\Range(\morphSeqName{f}^1) = \Domain(\morphSeqName{f}^2)$ and
$0^1_2=o^2_1$. Then

\begin{equation}
\seqname{M}^2 \morphCompose[A] \seqname{M}^1 \defeq
\bigl (
  \morphSeqName{f}^2 \morphCompose[()] \morphSeqName{f}^1,
  o^1_1,
  o^2_2
\bigr )
\end{equation}
\end{definition}

\begin{lemma}[Tuple composition for labeled morphisms]
\label{lem:labcomp}
Let $\seqname{M}^i \defeq (\morphSeqName{f}^i, o^i_1, o^i_2)$,
\donotbreak{$i \in [1,3]$,}
be a tuple such that $\morphSeqName{f}^i$ is a sequence or tuple of
functions,
\donotbreak
  {
    $
      \Range(\morphSeqName{f}^i) =
      \Domain(\morphSeqName{f}^{i+1})
    $
  }
and $o^i_2 = o^{i+1}_1$, $i=1,2$.
Then
\begin{equation}
\seqname{M}^3 \morphCompose[A] \bigl ( \seqname{M}^2 \morphCompose[A] \seqname{M}^1 \bigr )
=
\bigl ( \seqname{M}^3 \morphCompose[A] \seqname{M}^2 \bigr ) \morphCompose[A] \seqname{M}^1
\end{equation}

\begin{proof}
From \fullcref{def:labcomp}, \\
{
  \showlabelsinline
  \pagecref{lem:seqfunc}
}
and \fullcref{lem:tupfunc}, we have
\begin{equation*}
\begin{split}
  \seqname{M}^3 \morphCompose[A] (\seqname{M}^2 \morphCompose[A] \seqname{M}^1)
  & =
  \seqname{M}^3 \morphCompose[A] (\morphSeqName{f}^2 \morphCompose[()] \morphSeqName{f}^1, o^1_1, o^2_2)
\\*
  & =
  (\morphSeqName{f}^3 \morphCompose[()] \morphSeqName{f}^2 \morphCompose[()] \morphSeqName{f}^1, o^1_1, o^3_2)
\\*
  & =
  (\morphSeqName{f}^3 \morphCompose[()] \morphSeqName{f}^2, o^2_1, o^3_2) \morphCompose[A] \seqname{M}^1
\\*
  & =
  (\seqname{M}^3 \morphCompose[A] \seqname{M}^2) \morphCompose[A] \seqname{M}^1
\end{split}
\end{equation*}
\end{proof}

Let $\seqname{D}^i \defeq \Domain(\morphSeqName{f}^i)$ and
$\seqname{R}^i \defeq \Range(\morphSeqName{f}^i)$. Then
\begin{equation}
(\ID[\seqname{R}^i], o^i_2, o^i_2) \morphCompose[A] \seqname{M}^i =
\seqname{M}^i
\end{equation}
\begin{equation}
\seqname{M}^i \morphCompose[A] (\ID[\seqname{D}^i], o^i_1, o^i_1) =
\seqname{M}^i
\end{equation}

\begin{proof}
\begin{equation*}
\begin{split}
  (\ID[\seqname{R}^i], o^i_2, o^i_2) \morphCompose[A] \seqname{M}^i
  & =
  (\ID[\seqname{R}^i] \morphCompose[()] \morphSeqName{f}^i, o^i_1, o^i_2)
\\*
  & =
  (\morphSeqName{f}^i, o^i_1, o^i_2)
\\*
  & =
 \seqname{M}^i
\end{split}
\end{equation*}
\begin{equation*}
\begin{split}
  \seqname{M}^i \morphCompose[A] (\ID[\seqname{D}^i], o^i_1, o^i_1)
  & =
  (\morphSeqName{f}^i \morphCompose[()] \ID[\seqname{D}^i], o^i_1, o^i_2)
\\*
  & =
  (\morphSeqName{f}^i, o^i_1, o^i_2)
\\*
  & =
 \seqname{M}^i
\end{split}
\end{equation*}
\end{proof}
\end{lemma}

\begin{definition}[Cartesian product of sequence]
\label{def:Cart}
Let $\seqname{S}^i \defeq (S^i_\alpha, \alpha \prec \Alpha)$, $i=1,2$,
be a sequence and
$\morphSeqName{f} \defeq (\morphName{f}_\alpha \maps S^1_\alpha \to
S^2_\alpha, \alpha \prec \Alpha)$
be a sequence of functions, then
$
  \bigtimes \seqname{S}^i \defeq
  \bigtimes_{\alpha \prec \Alpha} S^i_\alpha
$
is the generalized Cartesian product of the sequence $\seqname{S}^i$ and \\*
\donotbreak
  {
    $
      \bigtimes \morphSeqName{f} \maps \seqname{S}^1 \to \seqname{S}^2
      \defeq
      \bigtimes_{\alpha \prec \Alpha}\morphName{f}_\alpha
    $
  }
is the generalized Cartesian product of the function sequence
$\morphSeqName{f}$.
\end{definition}

\begin{definition}[underline]
Let $\seqname{S}^1 \defeq (S^1_\alpha, \alpha \preceq \Alpha)$,
$\seqname{S}^2 \defeq (S^2_\alpha, \alpha \preceq \Alpha)$ be
sequences and
$
  \morphSeqName{f} \defeq
  (
    \morphName{f}_\alpha \maps S^1_\alpha \to S^2_\alpha,
    \alpha \preceq \Alpha
  )
$
be a sequence of functions, then
\begin{equation}
\underline{\morphName{f}} \maps \head(S_1) \to \head(S_2) \defeq
\bigtimes \head(\morphSeqName{f}) =
\bigtimes_{\alpha \prec \Alpha} \morphName{f}_\alpha
\end{equation}
is the function mapping
$(s_\alpha \in S^1_\alpha, \alpha \prec \Alpha)$
into $(\morphName{f}_\alpha(s_\alpha), \alpha \prec \Alpha)$.
\end{definition}

\begin{definition}[Head and tail compositions]
Let \\
$\morphName{f}^1 \maps (S^1_\alpha, \alpha \prec \Alpha) \to S^1_\Alpha$,
$\morphName{f}^2 \maps (S^2_\alpha, \alpha \prec \Alpha) \to S^2_\Alpha$
and
$
  \morphSeqName{g} \defeq
  (
    \morphName{g}_\alpha \maps S^1_\alpha \to S^2_\alpha,
    \alpha \preceq \Alpha
  )
$.
Then (see \cref{fig:fg,fig:gf})

\begin{subequations}
\begin{equation}
\morphName{f}^2 \morphComposeHead \morphSeqName{g} \defeq
\morphName{f}^2 \morphCompose \underline{\morphSeqName{g}}
\end{equation}
\begin{equation}
\morphSeqName{g} \morphComposeTail \morphName{f}^1 \defeq
\tail(\morphSeqName{g}) \morphCompose \morphName{f}^1
\end{equation}
\end{subequations}
\end{definition}

\begin{figure}
\[ \bfig
\node s1(0,0)[{\head(S^1)}]
\node s2(0,-1000)[{\head(S^2)}]
\node s2a(1000,-1000)[{S^2_\Alpha}]
\arrow |l|[s1`s2;{\underline{\morphSeqName{g}}}]
\arrow |r|[s1`s2a;{\morphSeqName{f}^2 \morphComposeHead \morphSeqName{g}}]
\arrow |b|[s2`s2a;{\morphSeqName{f^2}}]
\efig \]
\caption{$\morphName{f}^2 \morphComposeHead \morphSeqName{g}$}
\label{fig:fg}
\end{figure}

\begin{figure}
\[ \bfig
\node s1(0,0)[{\head(S^1)}]
\node s1a(0,-1000)[{S^1_\Alpha}]
\node s2a(1000,-1000)[{S^2_\Alpha}]
\arrow |l|[s1`s1a;{\morphName{f}^1}]
\arrow |r|[s1`s2a;{\morphSeqName{g} \morphComposeTail \morphSeqName{f}^1}]
\arrow |b|[s1a`s2a;{\tail(\morphSeqName{g}) = \morphSeqName{g}_\Alpha}]
\efig \]
\caption{$\morphSeqName{g} \morphComposeTail \morphName{f}^1$}
\label{fig:gf}
\end{figure}

\begin{definition}[Topology functions]
Let $S$ and $T$ be a topological spaces and $Y$ a subset of $S$. Then

\begin{enumerate}
\item
$S \subsp T$ iff $S$ is a subspace of $T$.
\item
$S \subsp[open-] T$ iff $S$ is an open subspace of $T$.
\item
$\Set(S)$ is the underlying set of $S$.
\item
$\Topology(S)$ is the topology of $S$.
\item
$\Topology(Y,S) \defeq \set{U \cap Y}[{U \in \Topology(S)}]$
is the relative topology of $Y$.
\item
$\Top(Y,S) \defeq \bigl ( Y, \Topology(Y,S) \bigr )$ is $Y$ with the
relative topology.
\item
$
  \op{S} \defeq
  \set
    {(U, \Topology(U,S))}%
    [{{ U \in \Topology(S) \setminus \{ \emptyset \} }}]
$
is the set of all non-null open subspaces of $S$.
\end{enumerate}

Let $\seqname{S}$ be a set of topological spaces. Then
$\op{\seqname{S}} \defeq \union [{S \in \seqname{S}}] { {\op{S}}}$ is the set
of open subspaces in $\seqname{S}$.

Let $S$ and $T'$ be spaces, $T \subseteq T'$ be a subspace and
$\morphName{f} \maps S \to T$ a function. Then
$\morphName{f} \maps S \to T' \defeq \Id[{T},{T'}] \morphCompose \morphName{f}$
is $\morphName{f}$ considered as a function from $S$ to $T'$.

Let $S'$ and $T'$ be spaces, $S \subseteq S'$, $T \subseteq T'$ be
subspaces and $\morphName{f}' \maps S' \to T'$ a function such that
$\morphName{f}'[S] \subseteq T$. Then $\morphName{f}' \maps S \to T$, also
written $\morphName{f}' \restrictto_{S,T}$, is
$\morphName{f}' \restrictto_S$ considered as a function from $S$ to
$T$.

Let $\seqname{S}^i \defeq (S^i_\alpha, \alpha \preceq \Alpha)$, $i-1,2$,
be a sequence of spaces, $\seqname{S}^1 \SUBSETEQ \seqname{S}^2$ and \linebreak
$\morphName{f}^2 \maps \head(\seqname{S}^2) \to \tail(\seqname{S}^2)$ a
function.
$
  \morphName{f}^2 \restrictto_{\head(\seqname{S}^1)} \defeq
  \morphName{f}^2 \restrictto_{\bigtimes \head(\seqname{S}^1)}
$.
If \linebreak
$
  \morphName{f}^2 \restrictto_{\head(\seqname{S}^1)}
    [\bigtimes \head(\seqname{S}^1)]
  \subseteq \tail(\seqname{S}^1)
$
then
$ \morphName{f}^2 \maps \head(\seqname{S}^1) \to \tail(\seqname{S}^1)$,
also written
$
  \morphName{f}^2 \restrictto_%
    {\head(\seqname{S}^1),\tail(\seqname{S}^1)}
$,
is $\morphName{f}^2 \restrictto_{\head(\seqname{S}^1)}$ considered as a
function from $\bigtimes \head(\seqname{S}^1)$ to $\tail(\seqname{S}^1)$.
\end{definition}

\begin{definition}[Truth space]
$\false \defeq \emptyset$, $\true \defeq \set{\emptyset}$,
the truth set is \linebreak
$\truthset \defeq \set{\false, \true}$,
$\truthtop \defeq \set {\emptyset, \set{\true}, \truthset}$ and
the truth space $\truthspace \defeq (\truthset, \truthtop)$ is
$\truthset$ with $\set{\true}$ and $\truthset$ as the only nonvoid open
sets
\end{definition}

\begin{definition}[Truth category]
\label{def:truthcat}
Let $\morphName{f}_0 = \Id[\truthspace]$ and
$\morphName{f}_1 \maps \truthspace \to \truthspace$ be the constant true
function, i.e., value $\true$.  The truth category is
$
   \truthcat \defeq \\
   \bigl (
     \truthspace,
     \set{\morphName{f}_0, \morphName{f}_1}
   \bigr )
$.
The truth model space is \\*
$\seqname{Truthspace} \defeq (\truthspace, \truthcat)$.
\end{definition}

\begin{definition}[Constraint functions]
A constraint function is a continuous function with range
$\truthspace$ or a model
function with range $\seqname{Truthspace}$.
\end{definition}

\begin{definition}[Sequence inclusion]
\label{def:seqin}
Let $\seqname{S} \defeq (S_\alpha, \alpha \prec \Alpha)$ and
$\seqname{T} \defeq (T_\alpha, \alpha \prec \Alpha)$ be sequences.
$\seqname{S} \seqin \seqname{T}$ iff
$\uquant{\alpha \prec \Alpha} {S_\alpha \in T_\alpha}$ or
$\uquant{\alpha \prec \Alpha} {S_\alpha \objin T_\alpha}$.
\end{definition}

\begin{lemma}[Sequence inclusion]
\label{lem:seqin}
Let $\seqname{S} \defeq (S_\alpha, \alpha \prec \Alpha)$
be a sequence, \\*
\nobreaks
  {{
    $
      \catSeqName{T}^i \defeq
      (\catName{T}^i_\alpha, \alpha \prec \Alpha)
    $,
  }}
$i=1,2$, a sequence of
categories, $\catSeqName{T}^1 \SUBCAT \catSeqName{T}^2$ and
$\seqname{S} \seqin \catSeqName{T}^1$ Then
$\seqname{S} \seqin \catSeqName{T}^2$.

\begin{proof}
If
$\uquant{\alpha \prec \Alpha} {S_\alpha \objin \catName{T}^1_\alpha}$
then
$\uquant{\alpha \prec \Alpha} {S_\alpha \objin \catName{T}^2_\alpha}$.
\end{proof}
\end{lemma}

\part{Model spaces and allied notions}
\label{part:modeldef}
Let $S$ be a topological space. We need to formalize the notions of an
open cover by sets that are ``well behaved'' in some sense, e.g.,
convex, sufficiently small, and of {\textquotedblleft}well
behaved{\textquotedblright} functions among those sets, e.g., preserving
fibers, smooth. We do this by associating a category of acceptable sets
and functions.

\begin{remark}
Using pseudo-groups, as in \cite[p.~1]{FoundDiffGeo1}, would not allow
restricting model neighborhoods to, e.g., convex sets.
\end{remark}

\section{Model spaces}
\begin{definition}[Model spaces]
\label{def:model}
Let $S$ be a topological space and $\catName{S}$ a small category whose
objects are open subsets of $S$ and whose morphisms are continuous
functions. $\seqname{S} \defeq (S, \catName{S})$ is a model space for
$S$ iff

\begin{enumerate}
\item $\Ob(\catName{S})$ is an open cover for $S$. Note that it need
not be a basis for $S$.
\label{mod:cover}
\item $\Ob(\catName{S})$ is closed under finite intersections.
\label{mod:closed}
\item The morphisms of $\catName{S}$ are continuous functions in $S$.
\label{mod:continuous}
\item If $\morphName{f} \maps A \to B$ is a morphism,
$A' \objin \catName{S} \subseteq A \objin \catName{S}$,
$B' \objin \catName{S} \subseteq B \objin \catName{S}$ and
\nobreaks
  {{
    $\morphName{f}[A'] \subseteq B'$
  }}
then
$\morphName{f} \restrictto_{A'} \maps A' \to B'$ is a morphism.
\label{mod:restriction}
\item If $A' \objin \catName{S} \subseteq A \objin \catName{S}$ then
the inclusion map $\Id[{A'},{A}] \maps A' \hookrightarrow A$ is a
morphism.
\label{mod:inclusion}
\begin{remark}
This is actually a consequence of \cref{mod:restriction},
but it is convenient to give it here.
\end{remark}
\item Restricted sheaf condition:
\label{mod:sheaf}
informally, consistent morphisms can be glued together. Whenever
\begin{enumerate}
\item $U_\alpha$ and $V_\alpha$, $\alpha \prec \Alpha$,
are objects of $\catName{S}$.
\item $\morphName{f}_\alpha \maps U_\alpha \to V_\alpha$ are morphisms of
$\catName{S}$.
\item
$U = \union[\alpha \prec \Alpha]{U_\alpha}$ is an object of $\catName{S}$.
\item
$V = \union[\alpha \prec \Alpha]{V_\alpha}$ is an object of $\catName{S}$.
\item
$\morphName{f} \maps U \to V$ is a continuous function and
for every $\alpha \prec \Alpha$,
$\morphName{f}$ agrees with $\morphName{f}_\alpha$ on $U_\alpha$
\end{enumerate}
then $\morphName{f}$ is a morphism of $\catName{S}$.
\end{enumerate}

By abuse of language we write
{\textquotedblleft}$U \objin \seqname{S}${\textquotedblright} for
{\textquotedblleft}$U \objin \Cat{\seqname{S}}${\textquotedblright} and
{\textquotedblleft}$\morphName{f} \arin \seqname{S}${\textquotedblright}
for
{\textquotedblleft}$\morphName{f} \arin \Cat{\seqname{S}}${\textquotedblright}.

\begin{remark}
A model space is similar to a pseudo-group, but has some important
differences:
\begin{enumerate}
\item
Open subsets of objects of $\catName{S}$ need not be objects of
$\catName{S}$.

\item
The morphisms of $\catName{S}$ need not be homeomorphisms or even 1-1,
thus they need not have inverses.

\item
The morphisms of $\catName{S}$ need not satisfy the sheaf condition,
but only the restricted sheaf condition in \cref{mod:sheaf}.
\end{enumerate}
\end{remark}

$\seqname{S}$ is fine grained iff every open subset of an object of
$\catName{S}$ is an object of $ \catName{S}$. Note that it is not
sufficient that $\Ob(\catName{S})$ be a basis for $S$.

$\Top(\seqname{S}) \defeq S = \pi_1(\seqname{S})$ is the topological
space of $\seqname{S}$.

Let $\seqname{S}^i$, $i=1,2$, be a model space.
$\Cat(\seqname{S}^i) \defeq \catName{S}^i = \pi_2(\seqname{S}^i)$ is the
category of model space $\seqname{S}^i$.
By abuse of language we write $\seqname{S}^1 \subsp \seqname{S}^2$
$\bigl ( \seqname{S}^1 \subsp[open-] \seqname{S}^2 \bigr )$ for
$\Top(\seqname{S}^1) \subsp \Top(\seqname{S}^2)$
$\bigl ( \Top(\seqname{S}^1) \subsp[open-] \Top(\seqname{S}^2) \bigr )$.

Let $U \objin \catName{S}$. Then $\Top(U,\seqname{S}) \defeq \Top(U,S)$
is $U$ with the relative topology.

$
  \Topcat(\seqname{S}) \defeq
  \Biggl (
    \set
      {{\Top(U,S)}}%
      [{U \objin \catName{S}}],
    \Ar(\catName{S})
  \Biggr )
$
is the topological category of $\seqname{S}$.

By abuse of language we write $U \subseteq \seqname{S}$ for
$U \objin \catName{S}$, $\Topology(\seqname{S})$ for
$\Topology \bigl ( \Top(\seqname{S}) \bigr )$, $\morphName{f}$ is a
morphism of $\seqname{S}$ for $\morphName{f}$ is a morphism of
$\catName{S}$ and $\morphName{f}$ is an isomorphism of $\seqname{S}$ for
$\morphName{f}$ is an isomorphism of $\catName{S}$.

Let $\seqname{C}$ be a set of model spaces. $\seqname{C}$ is fine
grained iff every $(C,\catName{C}) \in \seqname{C}$ is fine grained.
\end{definition}

\begin{lemma}[The topological category of a model space is a full topological category]
\label{lem:modtopcat}
Let $\seqname{M} \defeq (S, \catName{S})$ be a model space for $S$. Then
$\Topcat(\seqname{M})$ is a full topological category.
\begin{proof}
$\Topcat(\seqname{M})$ is a small subcategory of $\catName{Top}$ by
construction. $\Topcat(\seqname{M})$ is a full topological category by
\cref{mod:restriction} of
{
  \showlabelsinline
  \cref{def:model}.
}
\end{proof}
\end{lemma}

\begin{definition}[Model neighborhoods]
Let $\seqname{S} \defeq (S, \catName{S})$ be a model space for $S$.
Then the objects of $\catName{S}$ are model neighborhoods of
$\seqname{S}$. If $u \in U \objin \catName{S}$ then $U$ is a model
neighborhood of $\seqname{S}$ for $u$. If $U_i \objin \catName{S}$,
$i=1,2$, and $U_1 \subseteq U_2$ then $U_1$ is a model subneighborhood
of $U_2$; if $u^1 \in U_1$ then $U_1$ is also a model subneighborhood
for $u^1$.

If $\emptyset \objin \catName{S}$ then $\emptyset$ is a degenerate model
neighborhood of $\seqname{S}$; any nonvoid model neighborhood of
$\seqname{S}$ is a nondegenerate model neighborhood of $\seqname{S}$.
\end{definition}

\begin{definition}[Model subspaces]
\label{def:modsub}
$\seqname{S} \defeq (S, \catName{S})$ is a model subspace of
$\seqname{T} \defeq (T, \catName{T})$, abbreviated
$\seqname{S} \submod \seqname{T}$, iff

\begin{enumerate}
\item $\seqname{S}$ and $\seqname{T}$ are model spaces

\item $S$ is a subspace of $T$
\label{modsub:sub}

\item $\catName{S} \subcat \catName{T}$

\label{modsub:subcat}

\item every intersection of $S$ with an object of $\catName{T}$ is an object of $\catName{S}$
\label{modsub:rest}
\begin{equation}
\uquant%
  {U \objin \catName{T}}
  {U \cap S \objin \catName{S}}
\end{equation}

\item
\label{modsub:full}
Every object of $\catName{T}$ contained in $S$ is an object of
$\catName{S}$
\begin{equation}
\uquant%
  {U \objin \catName{T}: U \subseteq S}
  {U \objin \catName{S}}
\end{equation}

\begin{remark}
This is actually a consequence of \cref{modsub:rest}, but it is
convenient to state it here.
\end{remark}
\end{enumerate}

$\seqname{S}$ is also a full model subspace of $\seqname{T}$,
abbreviated $\seqname{S} \submod[full-] \seqname{T}$, iff
$S = T$.

$\seqname{S}$ is also a strict model subspace of $\seqname{T}$,
abbreviated $\seqname{S} \submod[strict-] \seqname{T}$, iff
$S$ is a model neighborhood of $\seqname{T}$.
\begin{remark}
$S$ is also a model neighborhood of $\seqname{S}$ by \cref{modsub:full}.
\end{remark}

Let $\seqname{S} \defeq (S, \catName{S})$ be a model space and $T$ a
model neighborhood of $\seqname{S}$. Then $\Mod(T,\seqname{S})$, the
relative model space of $T$, is $(T, \catName{T})$, where $\catName{T}$
is the full subcategory of $\catName{S}$ containing all model
subneighborhoods of $T$.

By abuse of language we write $T$ for $\Mod(T,\seqname{S})$ when
$\seqname{S}$ is understood by context.
\end{definition}

\begin{lemma}[Model subspaces]
\label{lem:modsub}

\begin{enumerate}

\item
$\submod$, $\submod[full-]$ and $\submod[strict-]$ are transitive.

\begin{proof}
Let $\seqname{S}^i \defeq (S^i, \catName{S}^i)$ $i=1,2,3$, be a model
space. If $\seqname{S}^i \submod \seqname{S}^{i+1}$, $i=1,2$, then

\begin{enumerate}
\item
$\seqname{S}^1$ and $\seqname{S}^3$ are model spaces by hypothesis.

\item Since $S^i$ is a subspace of $S^{i+1}$, $i=1,2$, then $S^1$ is a
subspace of $S^3$.

\item
Each $\catName{S}^i \subcat \catName{S}^{i+1}$ by
{
  \showlabelsinline
  \cref{modsub:subcat}
}
of \fullcref{def:modsub} above.
Then $\catName{S}^1 \subcat \catName{S}^3$.

\item
Each intersection of $S^i$ with an object of $\catName{S}^{i+1}$
is an object of $\catName{S}^i$ by
{
  \showlabelsinline
  \cref{modsub:rest}
}
of \cref{def:modsub}. Then every intersection of $S^1$ with an object
of $\catName{S}^3$ is an object of $\catName{S}^1$.

\item
Let $U \objin \seqname{S}^3$ and $U \subseteq S^1$.
$U \cap S^1 = U \objin \seqname{S}^1$.
\end{enumerate}

If each $\seqname{S}^i \submod[full-] \seqname{S}^{i+1}$, then each
$S^i = S^{i+1}$ by \cref{def:modsub}; then $S^1 = S^3$ and thus
$\seqname{S}^1 \submod[full-] \seqname{S}^3$.

If each $\seqname{S}^i \submod[strict-] \seqname{S}^{i+1}$, then each
$S^i$ is a model neighborhood of $\seqname{S}^{i+1}$ by
\cref{def:modsub}, and $S^1$ is a model neighborhood of
$\seqname{S}^3$ by item
{
  \showlabelsinline
  \cref{modsub:full}
}
of \cref{def:modsub}.
\end{proof}

\item
Let $\seqname{S} \defeq (S.\catName{S})$ be a model space and
$T \subseteq \seqname{S}$. Then
$\Mod(T,\seqname{S}) \submod \seqname{S}$

\begin{proof}
Let
$
  \seqname{T} = \bigl ( (T.\topname{T}), \catName{T} \bigr )
  \defeq \Mod(T,\seqname{S})
$.

\begin{enumerate}
\item
$\seqname{S}$ and $\seqname{T}$ are model spaces.

\item
$T$ is a subspace of $S$.

\item
$\catName{T} \subcat \catName{S}$ by construction.

\item
Let $U \objin \seqname{S}$. $U \cap T \objin \seqname{S}$. Since
$U \cap T \subseteq T$ is a model neighborhood of $\seqname{S}$, it is a
model neighborhood of $\Mod(T,\seqname{S})$ by construction.

\item
Let $U \objin \seqname{S}$ be contained in $T$. Then $U$ is a model
neighborhood of $\Mod(T,\seqname{S})$ by construction.
\end{enumerate}
\end{proof}

\item
Let $\seqname{S}^i \defeq (S^i,\catName{S}^i)$, $i=1,2$, and
$\seqname{S}^1 \submod \seqname{S}^2$.
Then $S^1$ is an open subspace of $S^2$.

\begin{proof}
The model neighborhoods of $\seqname{S}^1$ are model neighborhoods of
$\seqname{S}^2$, hence open in $S^2$. Their union, which is open in
$S^2$, is $S^1$ by \cref{mod:cover} of
{
  \showlabelsinline
  \pagecref{def:modsub}
}.
\end{proof}

\item
Let $S^i$, $i=1,2$, be a topological space and $S^1 \subsp[open-] S^2$.
Then $\Id[S^1,S^2]$ is open and continuous.

\begin{proof}
Let $U$ be open in $S^1$.
Since $S^1 \subsp[open-] S^2$, there is a $V$ open in $S^2$ such that
$U = S^1 \cap V$.  Since $V$ and $S^1$ are open in $S^2$, so is
$U = S^1 \cap V$.

Let $V$ be open in $S^2$. Since $S^1 \subsp S^2$,
${\Id[S^1,S^2]}^{-1}[V] = V \cap S^1$ is open in $S^1$.
\end{proof}
\end{enumerate}
\end{lemma}

\begin{corollary}[Model subspaces]
\label{cor:modsub}
Let $\seqname{S}^i \defeq (S^i,\catName{S}^i)$, $i=1,2$,
$\seqname{S}'^2 \defeq (S'^2,\catName{S}'^2)$ be model spaces, and
$\seqname{S}^1 \submod \seqname{S}^2 \submod \seqname{S}'^2$.
\begin{enumerate}

\item
$S^1 \subsp[open-] S'^2$.

\begin{proof}
$S^1 \subsp[open-] S^2 \subsp[open-] S'^2$.
\end{proof}

\item
Let $\morphName{f} \maps S^1 \to S^2$ be continuous. Then
$\morphName{f} \maps S^1 \to S'^2$ is continuous.
\begin{proof}
$S^1 \subsp[open-] S'^2$, Let $V \subseteq S'^2$ be an arbitrary open
set.  Since $\seqname{S}^2 \submod \seqname{S}'^2$ by hypothesis,
$V \cap S^2$ is open in $S^2$.  Since $\morphName{f} \maps S^1 \to S^2$
is continuous by hypothesis,
$\morphName{f}^{-1}[V] = \morphName{f}^{-1}[V \cap S^2]$
is open in $S^1$.
\end{proof}
\end{enumerate}
\end{corollary}

\section{Trivial model spaces}
\label{sec:triv}
Informally, a trivial model space of a specific type is one that does
not restrict the potential objects and morphisms of its type.

\begin{definition}[Trivial model spaces]
\label{def:trivmod}
Let $S$ be a topological space and $\catName{S}$ the category of all
continuous functions between open sets of $S$.

$\Triv{S} \defeq (S, \catName{S})$ is the trivial model space
of $S$ and $\Triv{S}$ is a trivial model space.

Let $\seqname{S}$ be a set of topological spaces.

$
  \Triv{\seqname{S}}
  \defeq
  \set
    {\Triv{S'}}%
    [S' \in \seqname{S}]
$
is the set of all trivial model spaces in $\seqname{S}$.

$\Trivcat{\seqname{S}}$, the category of all continuous functions
between elements of $\Triv{\seqname{S}}$, is the trivial model category
of $\seqname{S}$.

$\optriv{\seqname{S}}$, the category of open trivial model spaces
in $\seqname{S}$, is the category whose objects are
$\set {\Triv{U}}[U \in \op{\seqname{S}}]$, the trivial
model spaces of non-null open sets of spaces in $\seqname{S}$,
and whose morphisms are all the continuous functions among them.
\end{definition}

\begin{lemma}[Trivial model spaces]
\label{lem:trivmod}
Let $S^i$, $i=1,2$, be a topological space.

\begin{enumerate}
\item
$\Triv{S^i}$ is a model space.

\begin{proof}
Let $\catName{S}^i \defeq \Cat \biggl ( \Triv{S^i} \biggr )$,
$\morphName{f} \maps A \to B$ be a morphism of $\catName{S}^i$,
$A' \objin \catName{S}^i \subseteq A$ and
$B' \objin \catName{S}^i \subseteq B$. Then the definitions of
continuity and of $\Triv{S^i}$ imply each of the following.
\begin{enumerate}
\item $\Ob(\catName{S}^i)=\Topology(S^i)$ and thus is an open cover of
$S^i$.
\item $\Ob(\catName{S}^i)=\Topology(S^i)$ and thus is closed under
finite intersections.
\item All morphisms in $\Ar(\catName{S}^i)$ are continuous.
\item Since $\morphName{f} \maps A \to B$ is a morphism of
$\catName{S}^i$, $\morphName{f}$ is continuous.
\item
Since $A' \objin \catName{S}^i$ then $A'$ is open.
\item
Since $B' \objin \catName{S}^i$ then $B'$ is open.
\item
Since $\morphName{f} \maps A \to B$ is continuous then
$\morphName{f} \restrictto_{A'} \maps A' \to B$ is continuous.
\item
Since $\morphName{f}[A'] \subseteq B'$ then
$\morphName{f} \restrictto_{A',B'} \maps A' \to B'$ is well defined and
continuous, hence a morphism.
\item If $A' \objin \catName{S}^i \subseteq A \objin \catName{S}^i$ then
\begin{enumerate}
\item $A$ and $A'$ are open.
\item
The inclusion map $\Id[A',A]$ is continuous.
\item The inclusion map $\Id[A',A]$ is a
morphism of $\catName{S}^i$ by the definition of $\Triv{S^i}$.

\item $\Triv{S^i}$ trivially satisfies the restricted sheaf condition
because every continuous function between open sets of $S^i$ is a
morphism of $\catName{S}^i$.
\end{enumerate}
\end{enumerate}
\end{proof}

\item
Let $S^1$ be an open subspace of $S^2$. Then
$\Triv{S^1} \submod[strict-] \Triv{S^2}$.

\begin{proof}
Let $\catName{S}^i \defeq \Cat \bigl ( \Triv{S^i} \bigr )$, $i=1,2$.
By \pagecref{def:trivmod}, $\catName{S}^i$ is the category of all
continuous functions between open sets of $S^i$,

$\Triv{S^1}$ and $\Triv{S^2}$ satisfy the six conditions of
\pagecref{def:modsub}:

\begin{description}
\item[$\Triv{S^1}$ and $\Triv{S^2}$ are model spaces:]
Proved above.

\item[$S^1$ is a subspace of $S^2$:]
\leavevmode \\*
$\Top \biggl ( \Triv{S^i} \biggr )= S^i$ by \pagecref{def:trivmod}.
$S^1$ is an open subspace of $S^2$ by hypothesis.

\item[\ensuremath{\catName{S}^1 \subcat[full-] \catName{S}^2}:]
\leavevmode \\*
If $U \objin \catName{S}^1$ then $U$ is open in $S^1$. Since $S^1$ is
an open subspace of $S^2$, $U$ is open in $S^1$ and thus
$U \objin \catName{S}^2$.

If $\morphName{f} \maps U \to V \arin \catName{S}^1$ then $\morphName{f}$
is a continuous function between open sets of $S^1$. Since $S^1$ is an
open subspace of $S^2$, $\morphName{f}$ is a continuous function between
open sets of $S^2$ and thus a morphism of $\catName{S}^2$.

If $\morphName{f} \maps U \to V \arin \catName{S}^2$ then $\morphName{f}$
is a continuous function between open sets of $S^2$. If
$U \objin \catName{S}^1$ and $U \objin \catName{S}^1$, then $U$ and
$V$ are open sets of $S^1$. Since $S^1$ is an open subspace of $S^2$,
$\morphName{f}$ is a continuous function between open sets of $S^1$ and
thus a morphism of $\catName{S}^1$.

\item[Every intersection of $S^1$ with an object of $\catName{S}^2$ is an object of $\catName{S}^1$:]
\leavevmode \\*
Let $U \objin \catName{S}^2$. $U$ and $S^1$
are open in $S^2$. $U /cup \Set(S^1)$ is open in $S^2$ and contained in
$S^1$. Since $S^1$ is an open subspace of $S^2$, $U /cup \Set(S^1)$ is
open in $S^1$ and thus an object of $\catName{S}^1$.

\item[Every object of $\catName{S}^2$ contained in $S^1$ is an object of $\catName{S}^1$:]
\leavevmode \\*
Let $U \objin \catName{S}^2$. $U$ is open in $S^2$. If $U$ is
contained in $S^1$ then it is open in $S^1$ and thus an object of
$\catName{S}^1$.
\end{description}

Since $S^1$ is an open subspace of $S^2$, $\Set(S^1)$ is a model
neighborhood of $\Triv{S^2}$ by \cref{def:trivmod} above
\end{proof}

$\Triv{S^1} \submod \Triv{S^2}$.

\begin{proof}
$S^1 \subsp[open-] S^2$ by \pagecref{lem:modsub}.
\end{proof}

\item
Let $\morphName{f} \maps S^1 \to S^2$ be continuous. Then $\morphName{f}$
is a model function from $\Triv{S^1}$ to $\Triv{S^2}$.

\begin{proof}
Let $U^i$ be a model neighborhood of $\Triv{S^i}$, $i=1,2$.
\begin{enumerate}
\item
Since $U^2$ is a model neighborhood of $\Triv{S^2}$, $U^2$ is open,
$\morphName{f}^{-1}[U^2]$ is open and thus a model neighborhood of
$\Triv{S^1}$.
\item
$\morphName{f}[U^1] \subseteq S^2$. Since $S^2$ is open, it is a model
neighborhood of $\Triv{S^2}$.
\end{enumerate}
\end{proof}
\end{enumerate}
\end{lemma}

\section{Minimal model spaces}
\label{sec:min}
Even if a set of open sets fails one of \cref{mod:cover,mod:closed} or a
set of functions among them fails one of
\crefrange{mod:continuous}{mod:sheaf} in the definition of a model
space, there is a minimal model space containing them.

\begin{definition}[Minimal model spaces]
\label{def:minmod}
Let $S$ be a topological space, $\seqname{O}$ a set of open sets in $S$,
$\morphSeqName{f}$ a set of continuous functions between elements of
$\seqname{O}$ and $\catName{S}$ the smallest concrete category  over
$\catName{Top}$ having all sets in $\seqname{O}$ as objects, having all
functions in $\morphSeqName{f}$ as morphisms and satisfying
\crefrange{mod:continuous}{mod:sheaf} of
{
  \showlabelsinline
  \pagecref{def:model}
}.

Then
\begin{equation}
\minimal{\Mod}
  (
    S,
    \seqname{O},
    \morphSeqName{f}
  )
\defeq
  \left (
    \Top \bigl ( \union{O}, S \bigr ),
    \catName{S}
  \right )
\end{equation}
is the minimal model space of $S$ with neighborhoods $\seqname{O}$ and
neighborhood mappings $\morphSeqName{f}$.
\begin{remark}
The trivial model space $\Triv{S}$ is a special case.
\end{remark}

Let $\seqname{S} \defeq (S, \catName{S})$ be a model space,
$\seqname{O}$ a set of model neighborhoods of $\seqname{S}$ and
$\morphSeqName{f}$ a set of morphisms between elements of $\seqname{O}$.

Then
\begin{equation}
\minimal{\Mod}
  (
    \seqname{S},
    \seqname{O},
    \morphSeqName{f}
  )
\defeq
\minimal{\Mod}
  (
    S,
    \seqname{O},
    \morphSeqName{f}
  )
\end{equation}
is the minimal model space of $\seqname{S}$ with neighborhoods
$\seqname{O}$ and neighborhood mappings $\morphSeqName{f}$.

\begin{remark}
$
  \minimal{\Mod}
    (
      \seqname{S},
      \seqname{O},
      \morphSeqName{f}
    )
$
is not, in general, a model subspace of $\seqname{S}$, since
{
  \showlabelsinline
  \cref{modsub:rest,modsub:full}
}
of \pagecref{def:modsub} may fail.

While every model neighborhood of
$\minimal{\Mod}(\seqname{S}, \seqname{O}, \morphSeqName{f})$
is a model neighborhood of
$\minimal{\Mod}(S, \seqname{O}, \morphSeqName{f})$, the converse need not
be true.

\end{remark}
\end{definition}

\begin{lemma}[Minimal model spaces are model spaces]
\label{lem:minmod}
Let $S$ be a topological space, $\seqname{O}$ a set of open sets in $S$
and $\morphSeqName{f}$ a set of continuous functions between elements of
$\seqname{O}$. Then
$(C,\catName{C}) \defeq \minimal{\Mod}(S, \seqname{O}, \morphSeqName{f})$
is a model space.

\begin{proof}
\leavevmode
\begin{enumerate}
\item Finite intersections of open sets are open and
$C = \union[U \in \seqname{O}]{U}$ by construction
\item $\Ob(\catName{C})$ is closed under finite intersections by
construction
\item Compositions of continuous functions, inclusion maps and
restrictions of continuous functions are continuous
\item Restrictions of morphisms are morphisms by construction
\item Inclusion maps are morphisms by construction
\item The restricted sheaf condition holds by construction
\end{enumerate}
\end{proof}

Let $\seqname{S} \defeq (S, \catName{S})$ be a model space,
$\seqname{O}$ a set of model neighborhoods of $\seqname{S}$ and
$\morphSeqName{f}$ a set of morphisms between elements of $\seqname{O}$.
Then
$\minimal{\Mod}(\seqname{S}, \seqname{O}, \morphSeqName{f})$
is a model space.

\begin{proof}
$
  \minimal{\Mod}(\seqname{S}, \seqname{O}, \morphSeqName{f}) =
  \minimal{\Mod}(S, \seqname{O}, \morphSeqName{f})
$
and $\minimal{\Mod}(S, \seqname{O}, \morphSeqName{f})$ is a model space.
\end{proof}
\end{lemma}

\section{M-paracompact model spaces}
\label{sec:m-para}
Paracompactness is an important property for topological spaces because
of partitions of unity. There is an analogous property for model
spaces.

\begin{definition}[Model topology]
\label{def:ModTop}
Let $\seqname{M} \defeq (S, \catName{S})$ be a model space for $S$.
Then the model topology $\topname{M}^*$ for $M$ is the topology
generated by $\Ob(\catName{S})$.
\begin{remark}
$\topname{M}^*$ is not guarantied to be T0 even if $S$ is T4. However,
$\topname{M}^*$ may be normal or regular even if $S$ is not.
\end{remark}
\end{definition}

\begin{definition}[m-paracompactness]
\label{def:M-para}
A model space $\seqname{M} \defeq (S, \catName{S})$ is \\
m-paracompact iff $\topname{M}^*$ is regular and every cover of $S$ by
model neighborhoods has a locally finite refinement by model neighborhoods.
\begin{remark}
This is a stronger condition than merely requiring $\topname{M}^*$ to
be paracompact.
\end{remark}
\end{definition}

\begin{theorem}[m-paracompactness and paracompactness]
If $\seqname{M} = (S, \catName{S})$ is m-paracompact and the model
neighborhoods form a basis for $S$ then $S$ is paracompact.
\begin{proof}
Let $R$ be an open cover of $S$.  Since the model neighborhoods form a
basis, every set in $R$ is a union of model neighborhoods and thus there
is a refinement $R_1$ by model neighborhoods.  Since by hypothesis
$\seqname{M}$ is m-paracompact, $R_1$ has a locally finite refinement
$R_2$ by model neighborhoods.  Since model neighborhoods are open sets,
$R_2$ is a locally finite refinement of $R$ in the conventional sense.
\end{proof}
\end{theorem}

\section{Model functions and model categories}
\label{sec:modcat}
It is convenient to have a notion of mappings between model spaces that
are well behaved in some sense, e.g., fiber preserving; using that
notion it is then possible to group model spaces into categories.

\subsection{Model functions}
\begin{definition}[Model functions]
\label{def:modfunc}
Let $\seqname{S}^i \defeq (S^i, \catName{S}^i)$, $i=1,2$, be a model
space and $\morphName{f} \maps S^1 \to S^2$ be a continuous function.
$\morphName{f}$ is a model function of $\seqname{S}^1$ to $\seqname{S}^2$
iff the inverse images of model neighborhoods are model neighborhoods.

\begin{equation}
\uquant%
  {V \objin \catName{S}^2}
  {\morphName{f}^{-1}[V] \objin \catName{S}^1}
\end{equation}

$\morphName{f}$ is also a constrained model function of $\seqname{S}^1$
to $\seqname{S}^2$ iff $\morphName{f}[S^1]$ is contained in a model
neighborhood of $\seqname{S}^2$.

\begin{equation}
\equant%
  {V \objin \catName{S}^2}
  {\morphName{f}[S^1] \subseteq V}
\end{equation}

\begin{remark}
When $\morphName{f} \maps \seqname{S}^1 \to \seqname{S}^2$ is
constrained, $S^1$ is a model neighborhood of $\seqname{S}^1$.
\end{remark}

By abuse of language we write
$\morphName{f} \maps \seqname{S}^1 \to \seqname{S}^2$ both for
$\morphName{f}$ considered as a model function and for $\morphName{f}$
considered as a continuous function.

Let $\seqname{S}^1 \defeq (S^1, \catName{S}^1)$ and
$\seqname{S}'^2 \defeq (S'^2, \catName{S}'^2)$ be model spaces,
$\seqname{S}^2 \defeq (S^2, \catName{S}^2) \submod \seqname{S}'^2$ be a
model subspace and $\morphName{f} \maps S^1 \to S^2$ be a model function
of $\seqname{S}^1$ to $\seqname{S}^2$. Then
$
  \morphName{f} \maps \seqname{S}^1 \to \seqname{S}'^2 \defeq
  \Id[{\seqname{S}^2},{\seqname{S}'^2}] \morphCompose \morphName{f}
$
is $\morphName{f}$ considered as a model function from $\seqname{S}^1$ to
$\seqname{S}'^2$.

Let $\seqname{S}'^i \defeq (S'^i, \catName{S}'^i)$, $i=1,2$, be a model
space,
$\seqname{S}^i \defeq (S^i, \catName{S}^i) \submod \seqname{S}'^i$ be a
model subspace and $\morphName{f}'$ be a model function of
$\seqname{S}'^1$ to $\seqname{S}'^2$ such that
$\morphName{f}'[S^1] \subseteq S^2$. Then
$\morphName{f} \maps \seqname{S}^1 \to \seqname{S}^2$, also written
$\morphName{f}' \restrictto_{\seqname{S}^1,\seqname{S}^2}$, is
$\morphName{f}' \restrictto_{\seqname{S}^1}$ considered as a model
function of $\seqname{S}^1$ to $\seqname{S}^2$.
\end{definition}

\begin{lemma}[Model functions]
\label{lem:modfunc}
Let $\seqname{S}^i \defeq (S^i, \catName{S}^i)$ and
$\seqname{S}'^i \defeq (S'^i, \catName{S}'^i)$, $i=1,2$, be model
spaces, $\seqname{S}^i \submod \seqname{S}'^i$ and
$\morphName{f}' \maps S'^1 \to S'^2$ be a model function of
$\seqname{S}'^1$ to $\seqname{S}'^2$.

$\morphName{f} \defeq \morphName{f}' \restrictto_{S^1,S^2}$ is a model
function of $\seqname{S}^1$ to $\seqname{S}^2$.

\begin{proof}
$\morphName{f}'$ is a model function by hypothesis.
$\morphName{f}'$ is continuous by \cref{def:modfunc} above, so
$\morphName{f} \defeq \morphName{f}' \restrictto_{S^1,S^2}$ is continuous.

Let $U$ be a model neighborhood of $\seqname{S}^2$.
$\seqname{S}^2 \submod \seqname{S}'^2$ by hypothesis, so
$U$ is a model neighborhood of $\seqname{S}'^2$.
$\morphName{f}'$ is a model function of $\seqname{S}'^1$ to
$\seqname{S}'^2$ by hypothesis, so
$\morphName{f}'^{-1}[U] \objin \seqname{S}'^1$ by \cref{def:modfunc}.
Then
$
  \morphName{f}^{-1}[U] = \morphName{f}'^{-1}[U] \cap S^1 \objin
  \seqname{S}^1
$
by \cref{def:modfunc}.
\end{proof}
\end{lemma}

\begin{lemma}[Model functions and trivial model spaces]
\label{lem:modfunctriv}
Let $S^i$, $i=1,2$, be a topological space and
$\morphName{f} \maps S^1 \to S^2$ be a continuous function.
Then $\morphName{f} \maps \Triv{S^1} \to \Triv{S^2}$ is a model function.

\begin{proof}
Let $U^i$ be a model neighborhood of $\Triv{S^i}$, $i=1,2$.

$\morphName{f}^{-1}[U^2]$ is open, hence a model neighborhood of
$\Triv{S^1}$.

$S^2$ is a model neighborhood of $\Triv{S^2}$ and
$\morphName{f}[U^1] \subseteq U^2$.
\end{proof}

Let $\seqname{S}^i \defeq (S^i,\catName{S}^i)$, $i=1,2$, be a model
space and $\morphName{f} \maps \seqname{S}^1 \to \seqname{S}^2$ be a
model function. Then $\morphName{f} \maps \Triv{S^1} \to \Triv{S^2}$ is a
model function.

\begin{proof}
$\morphName{f} \maps S^1 \to S^2$ is continuous.
\end{proof}

Let $\seqname{S}^i \defeq (S^i,\catName{S}^i)$, $i=1,2$, be a model
space and $\morphName{f} \maps \seqname{S}^1 \to \seqname{S}^2$ be a
model function. Then $\morphName{f} \maps \Triv{S^1} \to \seqname{S}^2$
is a model function iff
$\morphName{f} \maps \seqname{S}^1 \to \seqname{S}^2$  is a constrained
model function.

\begin{proof}
let $U^2$ be a model neighborhood of $\seqname{S}^2$.
$\morphName{f}^{-1}[U^2]$ is open, hence a model neighborhood of
$\Triv{S^1}$.

$S^1$ is a model neighborhood of $\Triv{S^1}$. If $\morphName{f} \maps
\Triv{S^1} \to \seqname{S}^2$ is a model function then
$\morphName{f}[S^1]$ is contained in a model neighbothood of
$\seqname{S}^2$ and thus $\morphName{f}$ is constrained.

If $\morphName{f}$ is constrained, let $U^1$ be a model neighborhood of
$\Triv{S^1}$. Then $\morphName{f}[U^1] \subseteq \morphName{f}[S^1]$ is
contained in a model neighbothood of $\seqname{S}^2$.
\end{proof}
\end{lemma}

\begin{lemma}[Composition of model functions]
\label{lem:modcomp}
Let $\seqname{S}_i \defeq (S_i,\catName{S}_i)$, $i \in [1,3]$, be a
model space and
$\morphName{f}_i \maps \seqname{S}_i \to \seqname{S}_{i+1}$, $i=1,2$, be
a model function.

$\morphName{f}_2 \morphCompose \morphName{f}_1$ is a model function.

\begin{proof}
Let $U_i \objin \catName{S}_i,$ $i=1,2,3$.

Since $\morphName{f}_2$ is a model function,
\nobreaks
  {{
    $\morphName{f}^{-1}_2[U_3] \objin \catName{S}_2$.
  }}
Since $\morphName{f}_1$ is a model function,
$
  (\morphName{f}_2 \morphCompose \morphName{f}_1)^{-1}[U_3]
  =
  \morphName{f}^{-1}_1[\morphName{f}^{-1}_2[U_3]]
  \objin
  \catName{S}_1
$.
Since $\morphName{f}_1$ is a model functions, $\morphName{f}_1[U_1]$
is contained in a model neighborhood $V_2$.
Since $\morphName{f}_2$ is a model functions, $\morphName{f}_2[V_2]$
is contained in a model neighborhood $V_3$.
Then $(\morphName{f}_2 \morphCompose \morphName{f}_1)[U_1] \subseteq V_3$.
\end{proof}

If each $\morphName{f}_i$ is constrained then
$\morphName{f}_2 \morphCompose \morphName{f}_1$ is constrained.

\begin{proof}
$\morphName{f}_1[S^1] \subseteq S^2$ and
$\morphName{f}_2[S^2] \subseteq S^3$, hence
$\morphName{f}_2 \morphCompose \morphName{f}_1[S^1] \subseteq S^3$.
\end{proof}
\end{lemma}

\begin{definition}[Model homeomorphisms]
\label{def:modelhomeo}
Let $\seqname{S}^i \defeq (S^i, \catName{S}^i)$, $i=1,2$, be a model
space and $\morphName{f} \maps S^1 \to S^2$ be a model function.
$\morphName{f}$ is a model homeomorphism iff it is also invertible and
its inverse is a model function.
\end{definition}

\subsection{Model categories}
\begin{definition}[Model categories]
\label{def:modcat}
A category $\catName{M}$ is a model category iff

\begin{enumerate}
\item the objects of $\catName{M}$ are model spaces

\item the morphisms of $\catName{M}$ are model functions.

\item composition is functional composition.

\item
\label{modcat:inc}
If $\seqname{S} \defeq (S, \catName{S}) \objin \catName{M}$ and
$S^i \objin \catName{S}$, $i=1,2$, then
$\seqname{S}^i \defeq \Mod(S^i,\seqname{S}) \objin \catName{M}$.
If further $S^1 \subseteq S^2$ then the inclusion map
$\Id[S^1,S^2] \maps \seqname{S}^1 \morphMono \seqname{S}^2$ is a morphism
of $\catName{S}$.

\item
\label{modcat:rest}
If $\seqname{S}^i \defeq (S^i, \catName{S}^i) \objin \catName{M}$,
$U^i \objin \catName{S}^i$,
$i=1,2$, $\morphName{f} \maps \seqname{S}^1 \to \seqname{S}^2$ is a
morphism of $\catName{M}$ and $\morphName{f}[U^1] \subseteq U^2$ then
$\morphName{f} \maps \Mod(U^1,\seqname{S}^1) \to \Mod(U^2,\seqname{S}^2)$
is a morphism of $\catName{M}$.

\item
\label{modcat:morph}
If $\seqname{S}^i \defeq ({S}^i, \catName{S}^i)$, $i=1,2$, is a
subspace of $\seqname{S} \defeq (S, \catName{S}) \objin \catName{M}$
and $\morphName{f} \maps S^1 \to S^2$ is a morphism of $\catName{S}$ then
$\morphName{f} \maps \seqname{S}^1 \to \seqname{S}^2$ is a morphism of
$\catName{M}$.
\end{enumerate}

A model category $\catName{M}$ satisfies the restricted
sheaf condition iff whenever
\begin{enumerate}
\item $U$ and $V$ are objects of $\catName{S}$.

\item $U_\alpha \submod U$ and $V_\alpha \submod V$,
$\alpha \prec \Alpha$, are objects of $\catName{M}$.

\item
$\Set(U) = \union[\alpha \prec \Alpha]{\Set(U_\alpha)}$

\item
$\Set(V) = \union[\alpha \prec \Alpha]{\Set(V_\alpha)}$

\item $\morphName{f}_\alpha \maps U_\alpha \to V_\alpha$ are morphisms of
$\catName{M}$.

\item
$\morphName{f} \maps U \to V$ is a model function and
for every $\alpha \prec \Alpha$,
$\morphName{f}$ agrees with $\morphName{f}_\alpha$ on $U_\alpha$
\end{enumerate}
then $\morphName{f}$ is a morphism of $\catName{M}$.

Let $\catName{M}^i$, $i=1,2$, be a model category. $\catName{M}^1$ is a
strict subcategory of $\catName{M}^2$, abbreviated
$\catName{M}^1 \subcat[strict-] \catName{M}^2$, iff
$\catName{M}^1 \subcat[fullt-] \catName{M}^2$ and $\catName{M}^2$
satisfies the restricted sheaf condition.
\end{definition}

\begin{lemma}[Model categories]
\label{lem:modcat}
Let $\catName{M}$ be a model category,
$\seqname{S}^i \defeq (S^i, \catName{S}^i) \objin \catName{M}$. \\*
$i=1,2$,
$U^i \objin \catName{S}^i$,
$\seqname{U}^i \defeq \Mod(U^i,\seqname{S}^i)$ and
$\morphName{f} \maps \seqname{U}^1 \to \seqname{U}^2 \arin \catName{M}$.

$\morphName{f} \maps \seqname{U}^1 \to \seqname{S}^2 \arin \catName{M}$.

\begin{proof}
Since $\Id[\seqname{U}^2,\seqname{S}^2] \arin \catName{M}$ by
\cref{modcat:inc} of
{
  \showlabelsinline
  \cref{def:modcat}
}
above and
$\morphName{f} \maps \seqname{U}^1 \to \seqname{U}^2 \arin \catName{M}$
by hypothesis,
$
  \morphName{f} \maps \seqname{U}^1 \to \seqname{S}^2 =
    \Id[\seqname{U}^2,\seqname{S}^2] \morphCompose
    \morphName{f} \maps \seqname{U}^1 \to \seqname{U}^2
  \arin \catName{M}
$
\end{proof}
\end{lemma}

\begin{definition}[Trivial model categories]
\label{def:modtrivcat}
Let $\seqname{S} \defeq (S, \catName{S})$ be a model space.

$
  \Triv[submod-]{\seqname{S}}
  \defeq
  \set
    {\seqname{M}}%
    [\seqname{M} \submod \seqname{S}]
$
is the set of all subspaces of $\seqname{S}$.

$\Trivcat[submod-]{\seqname{S}}$, the category of all model functions
between elements of $\Triv[submod-]{\seqname{S}}$, is the trivial
subspace model category of $\seqname{S}$.

$
  \Triv[full-submod-]{\seqname{S}}
  \defeq
  \set
    {\seqname{M}}%
    [\seqname{M} \submod[full-] \seqname{S}]
$
is the set of all full subspaces of $\seqname{S}$.

$\Trivcat[full-submod-]{\seqname{S}}$, the category of all full model
functions between elements of $\Triv[full-submod-]{\seqname{S}}$, is the
trivial full subspace model category of $\seqname{S}$.

Let $\seqname{S}$ be a set of model spaces.

$
  \Triv[submod-]{\seqname{S}}
  \defeq
  \set
    {\seqname{M}}%
    [
      \equant
        {\seqname{M}' \in \seqname{S}}
        {\seqname{M} \submod \seqname{M}'}
    ]
$
is the set of all model subspaces of model spaces in $\seqname{S}$.

$\Trivcat[submod-]{\seqname{S}}$, the category of all model functions
between elements of $\Triv[submod-]{\seqname{S}}$, is the trivial
subspace model category of $\seqname{S}$.

$
  \Triv[full-submod-]{\seqname{S}}
  \defeq
  \set
    {\seqname{M}}%
    [
      \equant
        {\seqname{M}' \in \seqname{S}}
        {\seqname{M} \submod[full-] \seqname{M}'}
    ]
$
is the set of all full subspaces in $\seqname{S}$.

$\Trivcat[full-submod-]{\seqname{S}}$, the category of all full model
functions between elements of $\Triv[full-submod-]{\seqname{S}}$, is the
trivial full subspace model category of $\seqname{S}$.
\end{definition}

\begin{lemma}[Trivial model categories]
\label{lem:modtrivcat}

Let $\seqname{S} \defeq (S, \catName{S})$ be a model space. \\*
\nobreaks
  {{
    $
      \Trivcat[full-submod-]{\seqname{S}} \subcat[full-]
      \Trivcat[submod-]{\seqname{S}}
    $.
  }}

\begin{proof}
$
  \uquant
    {\seqname{M} \submod[full-] \seqname{S}}
    {\seqname{M} \submod \seqname{S}}
$,
so
$
  \Ob \biggl ( \Trivcat[full-submod-]{\seqname{S}} \biggr ) \subseteq
  \Ob \biggl ( \Trivcat[submod-]{\seqname{S}} \biggr )
$.

It remains to show that for a pair of objects in both categories, the $\Hom$ sets are the same.
Let $\seqname{M}^i \submod[full-] \seqname{S}$, $i=1.2$.
\begin{enumerate}
\item
Let
$
  \morphName{f} \maps \seqname{M}^1 \to \seqname{M}^2 \arin
  \Trivcat[full-submod-]{\seqname{S}}
$.
Then $\morphName{f}$ is a model function and thus
$
  \morphName{f} \maps \seqname{M}^1 \to \seqname{M}^2 \arin
  \Triv[submod-]{\seqname{S}}
$.
\item
Let
$
  \morphName{f} \maps \seqname{M}^1 \to \seqname{M}^2 \arin
  \Triv[submod-]{\seqname{S}}
$.
Then $\morphName{f}$ is a model function and thus \\*
$
  \morphName{f} \maps \seqname{M}^1 \to \seqname{M}^2 \arin
  \Trivcat[full-submod-]{\seqname{S}}
$.
\end{enumerate}

\end{proof}
\end{lemma}

\subsection{(Local) m-morphisms}

\begin{definition}[Local m-morphisms]
\label{def:modlocalMmorph}
Let $\catName{M}^i$ be a model category, $i=1,2$, and
$\seqname{S}^i \defeq (S^i, \catName{S}^i) \objin \catName{M}^i$.

Let $U^i \subseteq \seqname{S}^i$, $i=1,2$. be open. A
function $\morphName{f} \maps U^1 \to U^2$ is a local
$\seqname{S}^1$-$\seqname{S}^2$ m-morphism of $U^1$ to $U^2$ iff
$\seqname{S}^1 \submod \seqname{S}^2$ and for every $u \in U^1$ there is
a model neighborhood $U_u \subseteq U^1$ for $u$ and a model
neighborhood $V_u \subseteq U^2$ for $v \defeq \morphName{f}(u)$ such
that $\morphName{f}[U_u] \subseteq V_u$ and
$\morphName{f} \maps U_u \to V_u$ is a morphism of $\seqname{S}^2$.

Let $U^1 \subseteq \seqname{S}^i$ and $U^2 \subseteq \seqname{S}^i$
be open.  A function
$\morphName{f} \maps U^1 \to U^2$ is a local $\seqname{S}^i$
m-morphism of $U^1$ to $U^2$ iff it is a local
$\seqname{S}^i$-$\seqname{S}^i$ m-morphism of $U^1$ to $U^2$.

A model function $\morphName{f} \maps \seqname{S}^1 \to \seqname{S}^2$ is
a local m-morphism of $\seqname{S}^1$ to $\seqname{S}^2$ iff
it is a local $\seqname{S}^1$-$\seqname{S}^2$ m-morphism of
$S^1$ to $S^2$.\footnote{
  This definition does not require that $S^i$ be a model neighborhod of
  $S^i$.
}
\newcounter{fn:nonmodel}
\setcounter{fn:nonmodel} {\value{footnote}}

A model function $\morphName{f} \maps \seqname{S}^i \to \seqname{S}^i$ is
a local m-morphism of $\seqname{S}^i$ iff it is a
local m-morphism of $\seqname{S}^i$ to $\seqname{S}^i$.

Let $U^i \objin \catName{S}^i$, $i=1,2$. A function
$\morphName{f} \maps U^1 \to U^2$ is a local
$\seqname{S}^1$-$\seqname{S}^2$-$\catName{M}^1$-$\catName{M}^2$
m-morphism of $U^1$ to $U^2$ iff
$\catName{M}^1 \subcat[full-] \catName{M}^2$ and for every $u \in U^1$
there is a model neighborhood $U_u \subseteq U^1$ for $u$ and a model
neighborhood $V_u \subseteq U^2$ for $v \defeq \morphName{f}(u)$ such
that $\morphName{f}[U_u] \subseteq V_u$ and
$
  \morphName{f} \maps \Mod(U_u,\seqname{S}^1)
  \to \Mod(V_u,\seqname{S}^2)
$
is a morphism of $\catName{M}^2$.

Let $U^1 \subseteq \seqname{S}^i$ and $U^2 \subseteq \seqname{S}^i$, be
open. A model function $\morphName{f} \maps U^1 \to U^2$ is a local
$\seqname{S}^i$-$\catName{M}^i$ m-morphism of $U^1$ to $U^2$ iff it is a
local $\seqname{S}^i$-$\seqname{S}^i$-$\catName{M}^i$-$\catName{M}^i$
m-morphism of $U^1$ to $U^2$.

A model function $\morphName{f} \maps \seqname{S}^1 \to \seqname{S}^2$ is
a local $\catName{M}^1$-$\catName{M}^2$ m-morphism of $\seqname{S}^1$ to
$\seqname{S}^2$ iff it is a local
$\seqname{S}^1$-$\seqname{S}^2$-$\catName{M}^1$-$\catName{M}^2$
m-morphism of $S^1$ to $S^2$.\footnotemark[\value{fn:nonmodel}]

A model function $\morphName{f} \maps \seqname{S}^i \to \seqname{S}^i$ is
a local $\catName{M}^i$ m-morphism iff it is a local
$\catName{M}^i$-$\catName{M}^i$ m-morphism of $\seqname{S}^i$ to
$\seqname{S}^i$.

The phrases {\textquotedblleft}locally a(n) ...
m-morphism{\textquotedblright} and {\textquotedblleft}a local ...
m-morphism{\textquotedblright} are equivalent.
\end{definition}

\begin{definition}[M-morphisms]
\label{def:modMmorph}
Let $\catName{M}^i$ be a model category, $i=1,2$, and \\*
$\seqname{S}^i \defeq (S^i, \catName{S}^i) \objin \catName{M}^i$.

Let $U^i \subseteq \seqname{S}^i$, $i=1,2$. be open. A model
function $\morphName{f} \maps U^1 \to U^2$ is an
$\seqname{S}^1$-$\seqname{S}^2$ m-morphism of $U^1$ to $U^2$ iff
$\seqname{S}^1 \submod \seqname{S}^2$ and $\morphName{f}$ is a morphism
of $\seqname{S}^2$.

A model function $\morphName{f} \maps \seqname{S}^1 \to \seqname{S}^2$ is
an m-morphism of $\seqname{S}^1$ to $\seqname{S}^2$ iff it is an
$\seqname{S}^1$-$\seqname{S}^2$ m-morphism of $S^1$ to
$S^2$.\footnotemark[\value{fn:nonmodel}]

A model function $\morphName{f} \maps \seqname{S}^i \to \seqname{S}^i$ is
an m-morphism of $\seqname{S}^i$ iff it is an m-morphism
of $\seqname{S}^i$ to $\seqname{S}^i$.

Let $U^i \subseteq \seqname{S}^i$, $i=1,2$. be model neighborhoods. A
model function $\morphName{f} \maps U^1 \to U^2$ is an
$\seqname{S}^1$-$\seqname{S}^2$-$\catName{M}^1$-$\catName{M}^2$
m-morphism of $U^1$ to $U^2$ iff
$\catName{M}^1 \subcat[full-] \catName{M}^2$ and \\*
$
  \morphName{f} \maps \Mod(U^1,\seqname{S}^1)
  \to \Mod(U^2,\seqname{S}^2)
$
is a morphism of $\catName{M}^2$.

Let $U^1 \subseteq \seqname{S}^i$ and $U^2 \subseteq \seqname{S}^i$,
$i=1,2$, be model neighborhoods. A model function
$\morphName{f} \maps U^1 \to U^2$ is an $\seqname{S}^i$-$\catName{M}^i$
m-morphism of $U^1$ to $U^2$ iff it is an
$\seqname{S}^i$-$\seqname{S}^i$-$\catName{M}^i$-$\catName{M}^i$
m-morphism of $U^1$ to $U^2$.

A model function $\morphName{f} \maps \seqname{S}^1 \to \seqname{S}^2$ is
an $\catName{M}^1$-$\catName{M}^2$ m-morphism of $\seqname{S}^1$ to
$\seqname{S}^2$ iff
\nobreaks{{$\catName{M}^1 \subcat[full-] \catName{M}^2$}} and
$\morphName{f}$ is a morphism of
$\catName{M}^2$.\footnotemark[\value{fn:nonmodel}]

A model function $\morphName{f} \maps \seqname{S}^i \to \seqname{S}^i$ is
an $\catName{M}^i$ m-morphism of $\seqname{S}^i$ iff it is an
$\catName{M}^i$-$\catName{M}^i$ m-morphism of $\seqname{S}^i$ to
$\seqname{S}^i$.
\end{definition}

\begin{lemma}[(Local) m-morphisms]
\label{lem:modMmorph}
Let $\catName{M}'^i$, $\catName{M}^i \subcat[full-] \catName{M}'^i$,
$i=1,2$, be model categories,
$\seqname{S}^i \defeq (S^i, \catName{S}^i) \objin \catName{M}^i$,
$\seqname{S}'^i \defeq (S'^i, \catName{S}'^i) \objin \catName{M}'^i$
and $\seqname{S}^i \submod \seqname{S}'^i$.

\begin{enumerate}
\item
\label{M-morph:openmorphislocal}
Let $U^i \subseteq \seqname{S}^i$, $i=1,2$, be open.  If
$\morphName{f} \maps U^1 \to U^2$ is a
$\seqname{S}^1$-$\seqname{S}^2$ m-morphism of $U^1$ to $U^2$ then
$\morphName{f} \maps U^1 \to U^2$ is a local
$\seqname{S}^1$-$\seqname{S}^2$ m-morphism of $U^1$ to $U^2$.

\begin{proof}
\leavevmode
\begin{enumerate}
\item
Since $\morphName{f}$ is an $\seqname{S}^1$-$\seqname{S}^2$ m-morphism of
$U^1$ to $U^2$, $\seqname{S}^1 \submod \seqname{S}^2$ and
$\morphName{f}$ is a morphism of $\seqname{S}^2$,

\item
$U^1$ and $U^2$ are model neighborhoods of $\seqname{S}^2$. $U^1$ is a
model neighborhood of $\seqname{S}^1$ by \cref{modsub:full} of
{
  \showlabelsinline
  \pagecref{def:modsub}.
}

\item
Let $u^1 \in U^1$, $u^2 \defeq \morphName{f}(u^1) \in U^2$, Then $U^1$ is
a model neighborhood for $u^1$ and $U^2$ is a model neighborhood for
$u^2$.  $\morphName{f}[U^1] \subseteq U^2$ by hypothesis.
$\morphName{f} \maps U^1 \to U^2$ is a morphism of $\seqname{S}^2$ as
shown above.
\end{enumerate}
\end{proof}

\item
\label{M-morph:openlocalismorph}
Let $U^i \subseteq \seqname{S}^i$ be a model neighborhood of
$\seqname{S}^i$, $i=1,2$. If $\morphName{f} \maps U^1 \to U^2$ is a
$\seqname{S}^1$-$\seqname{S}^2$ local m-morphism of $U^1$ to
$U^2$ then $\morphName{f} \maps U^1 \to U^2$ is a
$\seqname{S}^1$-$\seqname{S}^2$ m-morphism of $U^1$ to $U^2$.

\begin{proof}
$\seqname{S}^1 \submod \seqname{S}^2$ by hypothesis.
$U^1$ is a model neighborhood of $\seqname{S}^1$ by hypothesis, thus a
model neighborhood of $\seqname{S}^2$
\item
For each $u \in U^1$, let $U_u$ be a model subneighborhood of $U^1$  for
$u$ and $V_u$ be a model subneighborhood of $U^2$ for
$v \defeq \morphName{f}(u)$ such that $\morphName{f}[U_u] \subseteq V_u$
and $\morphName{f} \maps U_u \to  V_u$ is a morphism of $\seqname{S}^2$.
Then $\morphName{f} \maps S^1 \to S^2$ is a morphism of $\seqname{S}^2$
by
{
  \showlabelsinline
  \cref{mod:sheaf}
}
of \pagecref{def:model}.
\end{proof}

\item
\label{M-morph:opencatmorphislocal}
Let $U^i \subseteq \seqname{S}^i$ be a model neighborhood, $i=1,2$, If
$\morphName{f} \maps U^1 \to U^2$ is a
$\seqname{S}^1$-$\seqname{S}^2$-$\catName{M}^1$-$\catName{M}^2$
m-morphism of $U^1$ to $U^2$ then $\morphName{f} \maps U^1 \to U^2$ is a
local
$\seqname{S}^1$-$\seqname{S}^2$-$\catName{M}^1$-$\catName{M}^2$
m-morphism of $U^1$ to $U^2$.

\begin{proof}
\leavevmode
\begin{enumerate}
\item
Each $U^i$ is a model neighborhood of $\seqname{S}^i$ by hypothesis.

\item
Since $\morphName{f} \maps U^1 \to U^2$ is an
$\seqname{S}^1$-$\seqname{S}^2$-$\catName{M}^1$-$\catName{M}^2$
m-morphism of $U^1$ to $U^2$, \\*
\nobreaks
  {
    $\catName{M}^1 \subcat[full-] \catName{M}^2$
  }
and
$
  \morphName{f} \maps \Mod(U^1,\seqname{S}^1)
  \to \Mod(U^2,\seqname{S}^2)
$
is a morphism of $\catName{M}^2$ by \pagecref{def:modMmorph}.

\item
Let $u^1 \in U^1$, $u^2 \defeq \morphName{f}(u^1) \in U^2$.

\begin{enumerate}
\item
Each $U^i$ is a model neighborhood of $\seqname{S}^i$ for $u^i$,

\item
Each $U^i \subseteq U^i$.

\item
$\morphName{f}[U^1] \subseteq U^2$.
\end{enumerate}
\end{enumerate}
\end{proof}

\item
\label{M-morph:opencatlocalismorph}
Let $\catName{M}^2$ satisfy the restricted sheaf condition and each
$U^i \subseteq \seqname{S}^i$ be a model neighborhood, $i=1,2$, Every
$\morphName{f} \maps U^1 \to U^2$ that is a
$\seqname{S}^1$-$\seqname{S}^2$-$\catName{M}^1$-$\catName{M}^2$ local
m-morphism of $U^1$ to $U^2$ is a
$\seqname{S}^1$-$\seqname{S}^2$-$\catName{M}^1$-$\catName{M}^2$
m-morphism of $U^1$ to $U^2$.

\begin{proof}
\leavevmode
\begin{enumerate}
\item
Since each $U^i$ is a model neighborhood of $\seqname{S}^i$ by hypothesis,
each \\*
$\Mod(U^i,\seqname{S}^i) \objin \catName{M}^i$ by \cref{modcat:inc}
of
{
  \showlabelsinline
  \cref{def:modcat}
}

\item
$
  \morphName{f} \maps
  \Mod(U^1,\seqname{S}^1) \to
  \Mod(U^2,\seqname{S}^2)
$
is a model function by \cref{def:modlocalMmorph}

\item
Each $\seqname{S}^i \objin \catName{M}^i$ by hypothesis

\item
$\catName{M}^1 \subcat \catName{M}^2$ by hypothesis

\item
$\catName{M}^2$ satisfies the restricted sheaf condition by hypothesis

\item
For every $u^1 \in U^1$ there is a model neighborhood $U^1_{u^1}$ for
$u^1$ and a model neighborhood $U^2_{u^1}$ for
$u^2 \defeq \morphName{f}(u^1)$ such that
$\morphName{f}[U^1_{u^1}] \subseteq U^2_{u^1}$ and
$
  \morphName{f} \maps
  \Mod(U^1_{u^1},\seqname{S}^i) \to
  \Mod(U^2_{u^1},\seqname{S}^i)
$
is a morphism of $\catName{M}^i$.

\begin{enumerate}
\item
$\Mod(U^1_{u^1},\seqname{S}^1) \submod \seqname{S}^1$.

\item
Each
$
  \morphName{f} \maps
  \Mod(U^1_{u^1},\seqname{S}^1) \to
  \Mod(U^2_{u^1},\seqname{S}^2)
$
agrees with $\morphName{f}$ on $U^1_{u^1}$.

\item
Since $\union[{u^1 \in U^1}]{U^1_{u^1}} = U^1$ and
$\union[{u^1 \in U^1}]{U^2} = U^2$, then
$\morphName{f} \arin \catName{M}^2$ by \pagecref{def:modcat}.
\end{enumerate}
\end{enumerate}
\end{proof}

\item
\label{M-morph:localislocaltosup}
If $\morphName{f} \maps S^1 \to S^2$ is a local m-morphism of
$\seqname{S}^1$ to $\seqname{S}^2$ then $\morphName{f}$ is a local
m-morphism of $\seqname{S}^1$ to $\seqname{S}'^2$.

\begin{proof}
$\morphName{f} \maps S^1 \to S'^2$ is continuous by
\pagecref{cor:modsub}.

Let $u \in S^1$, $U_u$ be a model neighborhood of $\seqname{S}^1$ for
$u$ and $V_u$ be a model neighborhood of $\seqname{S}^2$ for
$v \defeq \morphName{f}(u)$ such that $\morphName{f}[U_u] \subseteq V_u$
and $\morphName{f} \maps U_u \to  V_u$ is a morphism of $\seqname{S}^2$.

Since $\seqname{S}^2 \submod \seqname{S}'^2$ by hypothesis,
$\morphName{f} \maps U_u \to  V_u$ is a morphism of $\seqname{S}'^2$.
\end{proof}

\item
\label{M-morph:morphismorphtosup}
If $\morphName{f} \maps S^1 \to S^2$ is an m-morphism of
$\seqname{S}^1$ to $\seqname{S}^2$, then $\morphName{f}$ is an m-morphism
of $\seqname{S}^1$ to $\seqname{S}'^2$.

\begin{proof}
$\seqname{S}^1 \submod \seqname{S}^2$ by \pagecref{def:modMmorph} and
$\seqname{S}^2 \submod \seqname{S}'^2$ by hypothesis, so
$\seqname{S}^1 \submod \seqname{S}'^2$.

$\morphName{f}$ is a morphism of $\seqname{S}^2$ by
\pagecref{def:modMmorph}, hence a morphism of $\seqname{S}'^2$.
\end{proof}

\item
\label{M-morph:catlocalislocaltosup}
If $\morphName{f} \maps S^1 \to S^2$ is a local
$\catName{M}^1$-$\catName{M}^2$ m-morphism of $\seqname{S}^1$ to
$\seqname{S}^2$, then $\morphName{f}$ is a local
$\catName{M}'^1$-$\catName{M}'^2$ m-morphism of $\seqname{S}^1$ to
$\seqname{S}'^2$.

\begin{proof}
Let $u \in S^1$, $U_u$ be a model neighborhood of $\seqname{S}^1$ for
$u$ and $V_u$ be a model neighborhood of $\seqname{S}^2$ for
$v \defeq \morphName{f}(u)$ such that $\morphName{f}[U_u] \subseteq V_u$
and $\morphName{f} \maps U_u \to  V_u$ is a morphism of $\catName{S}^2$.

Since $\catName{S}^2 \subcat[full-] \catName{S}'^2$ by hypothesis,
$\morphName{f} \maps U_u \to  V_u$ is a morphism of $\catName{S}'^2$.
\end{proof}

\item
\label{M-morph:catmorphismorphtosup}
If $\morphName{f} \maps S^1 \to S^2$ is an
$\catName{M}^1$-$\catName{M}^2$ m-morphism of $\seqname{S}^1$ to
$\seqname{S}^2$, then $\morphName{f}$ is an
$\catName{M}'^1$-$\catName{M}'^2$ m-morphism of $\seqname{S}^1$ to
$\seqname{S}'^2$.

\begin{proof}
$\catName{S}^1 \subcat[full-] \catName{S}^2$ by \pagecref{def:modMmorph} and \\*
$\catName{S}^2 \subcat[full-] \catName{S}'^2$ by hypothesis, so
$\catName{S}^1 \subcat[full-] \catName{S}'^2$.

$\morphName{f}$ is a morphism of $\catName{S}^1$ by
\cref{def:modMmorph}, hence a morphism of $\catName{S}'^2$.
\end{proof}

\item
\label{M-morph:contisstrict}
If $\seqname{S}^1 \submod \seqname{S}^2$ and
$\morphName{f} \maps S^1 \to S^2$ is continuous, then $\morphName{f}$ is
an m-morphism of $\Triv{S^1}$ to $\Triv{S^2}$.

\begin{proof}
\leavevmode \\*
\begin{description}
\item[$\morphName{f}$ is continuous:]
\leavevmode \\*
$\morphName{f} \maps S^1 \to S^2$ is continuous.
$
  \Top \biggl ( \Triv{S^i} \biggr ) =
  \Top(\seqname{S}^i) =
  S^i
$,
$i=1,2$.

\item [$\morphName{f}$ maps model neighborhoods into model neighborhoods:]
\leavevmode \\*
$S^2$ is open in $S^2$, hence a model neighborhood of $\Triv{S^2}$.
Let $U$ be a model neighborhood of $\Triv{S^1}$. Then
$\morphName{f}[U] \subseteq S^2$.

\item [Inverse images of model neighborhoods are model neighborhoods:]
\leavevmode \\*
Let $V$ be a model neighborhood of $\Triv{S^2}$. $V$ is open in $S^2$
and $\morphName{f}$ is continuous, so $\morphName{f}^{-1}[V]$ is open in
$S^1$ and hence a model neighborhood of $\Triv{S^1}$.
\end{description}
\end{proof}

\item
\label{M-morph:Idislocal}
If $\seqname{S}^1 \submod \seqname{S}^2$ then $\Id[S^1,S^2]$ is a local
m-morphism of $\seqname{S}^1$ to $\seqname{S}^2$. If each ${S}^i$ is
a model neighborhood of $\seqname{S}^i$ then $\Id[S^1,S^2]$ is a
morphism of $\catName{S}^2$.

\begin{proof}
Let $u \in S^1$ and $U_u$ a model neighborhood of $\seqname{S}^1$ for
$u$. Since $\seqname{S}^1 \submod \seqname{S}^2$, $U_u$ is also a
model neighborhood of $\seqname{S}^2$ for $u$, and hence
$\Id[U_u] \arin \catName{S}^2$.

If each $S^i$ is a model neighborhood of $\seqname{S}^i$, then since
$\seqname{S}^1 \submod \seqname{S}^2$ by hypothesis. $S^1$ is a model
neighborhood of $\seqname{S}^2$.  The result follows by
\cref{mod:inclusion} of
{
  \showlabelsinline
  \pagecref{def:model}.
}
\end{proof}

\item
\label{M-morph:Idcatislocal}
If $\catName{M}^1 \subcat \catName{M}^2$ and
$\seqname{S}^1 \submod \seqname{S}^2$ then $\Id[S^1,S^2]$ is a
morphism of $\catName{M}^2$.

\begin{proof}
$\seqname{S}^1 \objin \catName{M}^2$ because
$\catName{M}^1 \subcat \catName{M}^2$.
$\Id[S^2] \arin \catName{M}^2$. Then
$\Id[S^1,S^2] \arin \catName{M}^2$ by
{
  \showlabelsinline
  \cref{modcat:inc}
}
of \pagecref{def:modcat}.
\end{proof}

\item
\label{M-morph:openmorphistrivmorph}
Let $U^i \subseteq \seqname{S}^i$, $i=1,2$, be open.  If
$\morphName{f} \maps U^1 \to U^2$ is a (local)
$\seqname{S}^1$-$\seqname{S}^2$ m-morphism of $U^1$ to $U^2$ then
$\morphName{f} \maps U^1 \to U^2$ is a $\Triv{S^1}$-$\Triv{S^2}$
m-morphism of $U^1$ to $U^2$.

\begin{proof}
Since $\seqname{S}^1 \submod \seqname{S}^2$, $\Triv{S^1} \subsp \Triv{^2}$.
Since $\morphName{f} \maps U^1 \to U^2$ is a model function, it is
continuous and hence a morphism of $\Triv{S^2}$
\end{proof}
\end{enumerate}
\end{lemma}

\begin{corollary}[(Local) m-morphisms]
\label{cor:modMmorph}
Let $\catName{M}'^i$, $\catName{M}^i \subcat[full-] \catName{M}'^i$,
$i=1,2$, be model categories,
$\seqname{S}^i \defeq (S^i, \catName{S}^i) \objin \catName{M}^i$,
$\seqname{S}'^i \defeq (S'^i, \catName{S}'^i) \objin \catName{M}'^i$
and $\seqname{S}^i \submod \seqname{S}'^i$.

\begin{enumerate}
\item
\label{M-morph:morphislocal}
If $\morphName{f} \maps S^1 \to S^2$ is an m-morphism of
$\seqname{S}^1$ to $\seqname{S}^2$ then $\morphName{f} \maps S^1 \to S^2$
is a local m-morphism of $\seqname{S}^1$ to $\seqname{S}^2$.

\begin{proof}
Since $\morphName{f} \maps S^1 \to S^2$ is a
$\seqname{S}^1$-$\seqname{S}^2$ m-morphism of $S^1$ to $S^2$ then
$\morphName{f} \maps S^1 \to S^2$ is a local
$\seqname{S}^1$-$\seqname{S}^2$ m-morphism of $S^1$ to $S^2$.
\end{proof}

\item
\label{M-morph:sublocalismorph}
If $U^1$ and $U^2$ are model neighborhoods of $\seqname{S}^i$ then every
function $\morphName{f} \maps U^1 \to U^2$ that is a local
$\seqname{S}^i$ m-morphism of $U^1$ to $U^2$ is a $\seqname{S}^i$
m-morphism of $U^1$ to $U^2$.

\begin{proof}
Since $\morphName{f} \maps U^1 \to U^2$ is a local
$\seqname{S}^i$-$\seqname{S}^i$ m-morphism of $U^1$ to $U^2$ then it is
an $\seqname{S}^i$-$\seqname{S}^i$ m-morphism of $U^1$ to $U^2$.
\end{proof}

\item
\label{M-morph:localismorph}
If $\morphName{f} \maps S^1 \to S^2$ is a local m-morphism of
$\seqname{S}^1$ to $\seqname{S}^2$ and each $S^i$ is a model
neighborhood of $\seqname{S}^i$, then $\morphName{f} \maps S^1 \to S^2$
is an m-morphism of $\seqname{S}^1$ to $\seqname{S}^2$.

\begin{proof}
Since $\morphName{f} \maps S^1 \to S^2$ is a
$\seqname{S}^1$-$\seqname{S}^2$ local m-morphism of $S^1$ to $S^2$ then
it is an $\seqname{S}^1$-$\seqname{S}^2$ m-morphism of $S^1$ to
$S^2$.
\end{proof}

\item
\label{M-morph:catmorphislocal}
If $\morphName{f} \maps S^1 \to S^2$ is a $\catName{M}^1$-$\catName{M}^2$
m-morphism of $\seqname{S}^1$ to $\seqname{S}^2$ then
$\morphName{f} \maps S^1 \to S^2$ is a local
$\catName{M}^1$-$\catName{M}^2$ m-morphism of $\seqname{S}^1$ to
$\seqname{S}^2$.

\begin{proof}
Since $\morphName{f}$ is a
$\seqname{S}^1$-$\seqname{S}^2$-$\catName{M}^1$-$\catName{M}^2$
m-morphism of $S^1$ to $S^2$ then $\morphName{f}$ is a local
$\seqname{S}^1$-$\seqname{S}^2$-$\catName{M}^1$-$\catName{M}^2$
m-morphism of $S^1$ to $S^2$.
\end{proof}

\item
\label{M-morph:subcatlocalismorph}
If $\catName{M}^i$ satisfies the restricted sheaf condition and $U^j$ is
a model neighborhood of $\seqname{S}^i$, $j=1,2$, then every function
$\morphName{f} \maps U^1 \to U^2$ that is a local
$\seqname{S}^i$-$\catName{M}^i$ m-morphism of $U^1$ to $U^2$ is a
$\seqname{S}^i$-$\catName{M}^i$ m-morphism of $U^1$ to $U^2$.

\begin{proof}
\leavevmode
\begin{enumerate}
\item
$
  \morphName{f} \maps
  \Mod(U^1,\seqname{S}^i) \to
  \Mod(U^2,\seqname{S}^i)
$
is a local
\nobreaks%
  {
    {$\seqname{S}^i$-$\seqname{S}^i$-},
    {$\catName{M}^i$-$\catName{M}^i$}
  }
m-morphism of $U^1$ to $U^2$ by \pagecref{def:modlocalMmorph}

\item
$
  \morphName{f} \maps
  \Mod(U^1,\seqname{S}^i) \to
  \Mod(U^2,\seqname{S}^i)
$
is a
\nobreaks%
  {
    {$\seqname{S}^i$-$\seqname{S}^i$-},
    {$\catName{M}^i$-$\catName{M}^i$}
  }
m-morphism of $U^1$ to $U^2$ by \cref{M-morph:opencatlocalismorph} of
{
  \showlabelsinline
  \cref{lem:modMmorph}.
}

\item
The result follows from \cref{def:modlocalMmorph}.
\end{enumerate}
\end{proof}

\item
\label{M-morph:catlocalismorph}
Let each $S^i$ be a model neighborhood of $\seqname{S}^i$, Then every
$\morphName{f} \maps S^1 \to S^2$ that is a
$\catName{M}^1$-$\catName{M}^2$ local m-morphism of $\seqname{S}^1$ to
$\seqname{S}^2$ is a $\catName{M}^1$-$\catName{M}^2$ m-morphism
of $\seqname{S}^1$ to $\seqname{S}^2$.

\end{enumerate}

\begin{enumerate}
\item
If $\morphName{f} \maps S^1 \to S^2$ is a local m-morphism of
$\seqname{S}^1$ to $\seqname{S}^2$ and
$\seqname{S}^2 \submod \seqname{S}'^2$ then
$\morphName{f}$ is a local m-morphism of $\seqname{S}^1$ to
$\seqname{S}'^2$.

\begin{proof}
$\morphName{f} \maps S^1 \to S'^2$ is continuous by
\pagecref{cor:modsub}.
\end{proof}

\item
If $\morphName{f} \maps S^1 \to S^2$ is a(n) (local) m-morphism of
$\seqname{S}^1$ to $\seqname{S}^2$ then $\morphName{f}$ is an m-morphism
of $\Triv{S^1}$ to $\Triv{S^2}$.

\begin{proof}
$\morphName{f}$ is continuous by \pagecref{def:modlocalMmorph}.
$
  \Top \biggl ( \Triv{S^i} \biggr ) =
  \Top(\seqname{S}^i) =
  S^i
$,
$i=1,2$.
\end{proof}

\item
If $S^i$ is a model neighborhood of $\seqname{S}^i$ then
$\Id[{\seqname{S}^i}]$ is a morphism of $\catName{S}^i$.

\begin{proof}
$\seqname{S}^i \submod \seqname{S}^i$.
\end{proof}

\item
If $\morphName{f} \maps S^1 \to S^2$ is a local m-morphism of
$\seqname{S}^1$ to $\seqname{S}^2$ and
$\morphName{f}^2 \maps S^2 \to S^3$ is an m-morphism of $\seqname{S}^2$
to $\seqname{S}^3$ then $\morphName{f}^2 \morphCompose \morphName{f}$ is a
local m-morphism of $\seqname{S}^1$ to $\seqname{S}^3$.

\begin{proof}
Since $\morphName{f}^2 \maps S^2 \to S^3$ is an m-morphism of
$\seqname{S}^2$ to $\seqname{S}^3$ then $\morphName{f}^2$ is a local
m-morphism of $\seqname{S}^2$ to $\seqname{S}^3$.
\end{proof}

\item
If $\morphName{f} \maps S^1 \to S^2$ is an m-morphism of
$\seqname{S}^1$ to $\seqname{S}^2$ and
$\morphName{f}^2 \maps S^2 \to S^3$ is a local m-morphism of
$\seqname{S}^2$ to $\seqname{S}^3$ then
$\morphName{f}^2 \morphCompose \morphName{f}$ is a local m-morphism of
$\seqname{S}^1$ to $\seqname{S}^3$.

\begin{proof}
Since $\morphName{f} \maps S^1 \to S^2$ is an m-morphism of
$\seqname{S}^1$ to $\seqname{S}^2$ then $\morphName{f}$ is a local
m-morphism of $\seqname{S}^2$ to $\seqname{S}^2$.
\end{proof}

\item
If $\morphName{f} \maps S^1 \to S^2$ is a local
$\catName{M}^1$-$\catName{M}^2$ m-morphism of $\seqname{S}^1$ to
$\seqname{S}^2$ and $\morphName{f}^2 \maps S^2 \to S^3$ is an
$\catName{M}^2$-$\catName{M}^3$ m-morphism of $\seqname{S}^2$ to
$\seqname{S}^3$ then $\morphName{f}^2 \morphCompose \morphName{f}$ is a local
$\catName{M}^1$-$\catName{M}^3$ m-morphism of $\seqname{S}^1$ to
$\seqname{S}^3$.

\begin{proof}
Since $\morphName{f}^2 \maps S^2 \to S^3$ is an
$\catName{M}^2$-$\catName{M}^3$ m-morphism of $\seqname{S}^2$ to
$\seqname{S}^3$ then $\morphName{f}^2$ is a local
$\catName{M}^2$-$\catName{M}^3$ m-morphism of $\seqname{S}^2$ to
$\seqname{S}^3$.
\end{proof}

\item
If $\morphName{f} \maps S^1 \to S^2$ is an
$\catName{M}^1$-$\catName{M}^2$ m-morphism of $\seqname{S}^1$ to
$\seqname{S}^2$ and $\morphName{f}^2 \maps S^2 \to S^3$ is a local
$\catName{M}^2$-$\catName{M}^3$ m-morphism of $\seqname{S}^2$ to
$\seqname{S}^3$ then $\morphName{f}^2 \morphCompose \morphName{f}$ is a local
$\catName{M}^1$-$\catName{M}^3$ m-morphism of $\seqname{S}^1$ to
$\seqname{S}^3$.

\begin{proof}
Since $\morphName{f} \maps S^1 \to S^2$ is an
$\catName{M}^1$-$\catName{M}^2$ m-morphism of $\seqname{S}^1$ to
$\seqname{S}^2$ then $\morphName{f}$ is a local
$\catName{M}^1$-$\catName{M}^2$ m-morphism of $\seqname{S}^2$ to
$\seqname{S}^2$.
\end{proof}

\item
\label{M-morph:morphistrivmorph}
If $\morphName{f} \maps \seqname{S}^1 \to \seqname{S}^2$ is a (local)
m-morphism of $\seqname{S}^1$ to $\seqname{S}^2$ then
$\morphName{f} \maps S^1 \to S^2$ is an m-morphism of $\Triv{S^1}$ to
$\Triv{S^2}$.

\begin{proof}
Each $S^i$ is open, so $\morphName{f} \maps S^1 \to S^2$ is a
$\Triv{S^1}$-$\Triv{S^2}$ m-morphism of $S^1$ to $S^2$.
\end{proof}
\end{enumerate}
\end{corollary}

\begin{lemma}[Composition of local m-morphisms]
\label{lem:modlocalMmorphcomp}
Let $\catName{M}^i$, $i=1,2,3$, be model categories,
$\seqname{S}^i \defeq (S^i, \catName{S}^i) \objin \catName{M}^i$ and
$U^i \subseteq S^i$ be open.

\begin{enumerate}
\item
If $\morphName{f}^i \maps U^i \to U^{i+1}$, $i=1,2$, is a local
$\seqname{S}^i$-$\seqname{S}^{i+1}$ m-morphism of $U^i$ to $U^{i+1}$ then
$\morphName{f}^2 \morphCompose \morphName{f}^1$ a local
$\seqname{S}^1$-$\seqname{S}^3$ m-morphism of $U^i$ to $U^{i+1}$.

\begin{proof}
Since $\seqname{S}^1 \submod \seqname{S}^2$ and
$\seqname{S}^2 \submod \seqname{S}^3$ then
$\seqname{S}^1 \submod \seqname{S}^3$, Let $u^1 \in U^1$,
$u^2 \defeq \morphName{f}^1(u^1)$ and $u^3 \defeq \morphName{f}^2(u^2)$.
Let $U_{u^i}$ and $V_{u^i}$ be model neighborhoods of $\seqname{S}^i$
and $\seqname{S}^{i+1}$ such that that
$\morphName{f}^i[U_{u^i}] \subseteq V_{u^i}$ and
$
  \morphName{f}^i \maps \Mod(U_{u^i},\seqname{S}^i)
  \to \Mod(V_{u^i},\seqname{S}^{i+1})
$
is a morphism of $\seqname{S}^{i+1}$.

Let
$
  U'_{u^1} \defeq
  U_{u^1} \cap (\morphName{f}^2 \morphCompose \morphName{f}^1)^{-1}[V_{u^3}]
$.
Then $\morphName{f}^2 \morphCompose \morphName{f}^1 \maps U'_{u^1} \to V_{u^3}$
is a morphism of $\seqname{S}^3$.
\end{proof}

\item
If $\morphName{f}^i \maps S^i \to S^{i+1}$, $i=1,2$, is a local
$\catName{M}^i$-$\catName{M}^{i+1}$ m-morphism of $\seqname{S}^i$ to
$\seqname{S}^{i+1}$ then $\morphName{f}^2 \morphCompose \morphName{f}^1$ a
local $\catName{M}^i$-$\catName{M}^{i+2}$ m-morphism of
$\seqname{S}^1$ to $\seqname{S}^3$.

\begin{proof}
Since $\catName{M}^1 \subcat[full-] \catName{M}^2$ and
$\seqname{M}^2 \subcat[full-] \catName{M}^3$ then
$\seqname{M}^1 \subcat[full-] \catName{M}^3$.
Let
\nobreaks
  {{
    $u^1 \in \seqname{S}^1$,
  }}
$u^2 \defeq \morphName{f}^1(u^)$ and $u^3 \defeq \morphName{f}^2(u^2)$.
Let $U_{u^i}$ and $V_{u^i}$ be model neighborhoods of $\seqname{S}^i$
and $\seqname{S}^{i+1}$ such that
$\morphName{f}^i \maps U_{u^i} \to V_{u^i}$ is a morphism of
$\catName{M}^{i+1}$. Let
\nobreaks
  {{
    $
      U'_{u^1} \defeq
      U_{u^1} \cap (\morphName{f}^3 \morphCompose \morphName{f}^2)^{-1}[V_{u^3}]
    $.
  }}
Then $\morphName{f}^2 \morphCompose \morphName{f}^2 \maps U'_{u^1} \to V_{u^3}$
is a morphism of $\catName{M}^3$ by
{
  \showlabelsinline
  \cref{modcat:rest}
}
of \pagecref{def:modcat}.
\end{proof}
\end{enumerate}

If each $\morphName{f}^i \maps \seqname{S}^i \to \seqname{S}^{i+1}$,
$i=1,2$, is a local m-morphism of $\seqname{S}^i$ to
$\seqname{S}^{i+1}$, then
\nobreaks
  {{
    $\morphName{f}^2 \morphCompose \morphName{f}^1 \maps S^1 \to S^3$
  }}
is a local m-morphism of $\seqname{S}^1 \maps S^1 \to S^3$.

\begin{proof}
Since $\seqname{M}^1 \submod \seqname{M}^2$ and
$\seqname{M}^2 \submod \seqname{M}^3$,
$\seqname{M}^1 \submod \seqname{M}^3$.
Let $u \in S^i$, $v \defeq \morphName{f}^1(u)$ and
$w \defeq \morphName{f}^2(v)$. There exist a model neighborhood $U_u$
for $u$, model neighborhoods $V_u$, $V'_u$ for $v$ and a model
neighborhood $W_v$ of $w$ such that
\nobreaks
  {{
    $\morphName{f}_1[U_u] \subseteq V'_u$,
  }}
$\morphName{f}_2[V_v] \subseteq W_v$,
$\morphName{f}_1 \maps U_u \to V'_u$ is a morphism of $\seqname{S}^2$
and $\morphName{f}_2 \maps V_v \to W_v$ is a morphism of
$\catName{S}^3$. Then $\hat{V_u} \defeq V_v \cap V'_u \neq \emptyset$,
$\hat{V_u}$ is a model neighborhood of $v$ and
$\hat{U_u} \defeq \morphName{f}^{i-1}_1[\hat{V_u}]$ is a model
neighborhood for $u$.  $\morphName{f}_1 \maps \hat{U_u} \to \hat{V_u}$
and $\morphName{f}_2 \maps \hat{V_u} \to W_v$ are morphisms of
$\catName{S}^3$ by
{
  \showlabelsinline
  \cref{mod:restriction}
}
of \pagecref{def:model} and thus
$\morphName{f}_2 \morphCompose \morphName{f}_1 \maps \hat{U_u} \to W_v$ is a
morphism of $\seqname{S}^3$.
\end{proof}

If each $\morphName{f}^i \maps \seqname{S}^i \to \seqname{S}^{i+1}$ is
a local $\catName{M}^i$-$\catName{M}^{i+1}$ m-morphism of
$\seqname{S}^i$ to $\seqname{S}^{i+1}$ then \\*
\nobreaks
  {{
    $
      \morphName{f}^2 \morphCompose \morphName{f}^1 \maps
      \seqname{S}^1 \to \seqname{S}^3
    $
  }}
is a local $\catName{M}^1$-$\catName{M}^3$ m-morphism of
$\seqname{S}^1$ to $\seqname{S}^3$.

\begin{proof}
Since $\catName{M}^1 \subcat[full-] \catName{M}^2$ and
$\catName{M}^2 \subcat[full-] \catName{M}^3$ then
$\catName{M}^1 \subcat[full-] \catName{M}^3$.
Let \nobreaks{{$u \in \seqname{S}^1$,}} $v \defeq \morphName{f}^1(u)$ and
$w \defeq \morphName{f}^2(v)$. There exist a model neighborhood $U_u$
for $u$, model neighborhoods $V_v$, $V'_u$ for $v$ and a model
neighborhood $W_v$ of $w$ such that
$\morphName{f}^1[U_u] \subseteq V'_u$,
$\morphName{f}^1 \maps U_u \to V'_u$ is a morphism of $\catName{M}^2$,
$\morphName{f}^2[V_v] \subseteq W_v$ and
$\morphName{f}^2 \maps V_v \to W_v$ is a morphism of $\catName{M}^3$.
Then $\hat{V_u} \defeq V_v \cap V'_u \neq \emptyset$, $\hat{V_u}$ is a
model neighborhood of $v$ and
$\hat{U_u} \defeq \morphName{f}^{i-1}_1[\hat{V_u}]$ is a model
neighborhood for $u$.  $\morphName{f}^1 \maps \hat{U_u} \to \hat{V_u}$
and $\morphName{f}^2 \maps \hat{V_u} \to W_v$ are morphisms of
$\catName{M}^3$ by
{
  \showlabelsinline
  \cref{modcat:inc}
}
of \pagecref{def:modcat} and thus
$\morphName{f}^2 \morphCompose \morphName{f}^1 \maps \hat{U_u} \to W_v$ is a
morphism of $\catName{M}^3$.
\end{proof}
\end{lemma}

\begin{corollary}[[Composition of local m-morphisms]
\label{cor:modlocalMmorphcomp}
Let $\catName{M}^i$, $i=1,2,3$, be model categories,
$\seqname{S}^i \defeq (S^i, \catName{S}^i) \objin \catName{M}^i$.

\begin{enumerate}
\item
Let $U^j \subseteq S^i$, $j=1,2,3$, be open.
If $\morphName{f}^j \maps U^j \to U^{j+1}$, $j=1,2$, is a local
$\seqname{S}^i$ m-morphism of $U^j$ to $U^{j+1}$ then
$\morphName{f}^2 \morphCompose \morphName{f}^1$ is a local
$\seqname{S}^i$ m-morphism of $U^1$ to $U^3$.

\begin{proof}
Since each  $\morphName{f}^j \maps U^j \to U^{j+1}$ is a local
$\seqname{S}^i$-$\seqname{S}^i$ m-morphism of $U^j$ to $U^{j+1}$ then
$\morphName{f}^2 \morphCompose \morphName{f}^1$ is a local
$\seqname{S}^i$-$\seqname{S}^i$ m-morphism of $U^j$ to $U^{j+1}$.
\end{proof}

\item
If $\morphName{f}^i \maps S^i \to S^{i+1}$, $i=1,2$, is a local
m-morphism of $\seqname{S}^i$ to $\seqname{S}^{i+1}$ then
$\morphName{f}^2 \morphCompose \morphName{f}^1$ is a local m-morphism of
$\seqname{S}^1$ to $\seqname{S}^3$.

\begin{proof}
Since $\seqname{S}^1 \submod \seqname{S}^2$ and
$\seqname{S}^2 \submod \seqname{S}^3$ then
$\seqname{S}^1 \submod \seqname{S}^3$.
Let $u^1 \in \seqname{S}^1$, $u^2 \defeq \morphName{f}^1(u^1)$ and
$u^3 \defeq \morphName{f}^2(u^2)$.
Let $U_{u^i}$ and $V_{u^i}$ be model neighborhoods of $\seqname{S}^i$ and $\seqname{S}^{i+1}$
such that$\morphName{f}^i \maps U_{u^i} \to V_{u^i}$ is a morphism of $\seqname{S}^{i+1}$.
Let
\nobreaks
  {{
    $
      U'_{u^1} \defeq
      U_{u^1} \cap (\morphName{f}^2 \morphCompose \morphName{f}^2)^{-1}[V_{u^3}]
    $.
  }}
Then $\morphName{f}^2 \morphCompose \morphName{f}^2 \maps U'_{u^1} \to V_{u^3}$
is a morphism of $\seqname{S}^3$.
\end{proof}
\end{enumerate}

\end{corollary}

\begin{lemma}[Composition of m-morphisms]
\label{lem:modMmorphcomp}
Let $\catName{M}^i$, $i=1,2,3$, be model categories and
$\seqname{S}^i \defeq (S^i, \catName{S}^i) \objin \catName{M}^i$.

\begin{enumerate}

\item
If $\morphName{f}^i \maps S^i \to S^{i+1}$, $i=1,2$, is an m-morphism of
$\seqname{S}^i$ to $\seqname{S}^{i+1}$ then
$\morphName{f}^2 \morphCompose \morphName{f}^1 \maps S^1 \to S^3$ is an
m-morphism of $\seqname{S}^1 \maps S^1 \to S^3$.

\begin{proof}
Since each $\morphName{f}^i$ is a morphism of $\seqname{S}^{i+1}$ by
\pagecref{def:modMmorph} and
$\seqname{S}^1 \submod \seqname{S}^2 \submod \seqname{S}^3$, then
each $\morphName{f}^i$ is a morphism of $\seqname{S}^3$ by
\pagecref{def:modMmorph}, and thus
$\morphName{f}^2 \morphCompose \morphName{f}^1 \maps S^1 \to S^3$ is a morphism
of $\seqname{S}^3$.
\end{proof}

\item
If $\morphName{f}^i \maps S^i \to S^{i+1}$, $i=1,2$, is an
$\catName{M}^i$-$\catName{M}^{i+1}$ m-morphism of $\seqname{S}^i$ to
$\seqname{S}^{i+1}$ then $\morphName{f}^2 \morphCompose \morphName{f}^1$ an
$\catName{M}^i$-$\catName{M}^{i+2}$ m-morphism of $\seqname{S}^1$ to
$\seqname{S}^3$.

\begin{proof}
Since $\catName{M}^1 \subcat[full-] \catName{M}^2$ and
$\seqname{M}^2 \subcat[full-] \catName{M}^3$ then
$\seqname{M}^1 \subcat[full-] \catName{M}^3$.
Since each $\morphName{f}^i \maps S^i \to S^{i+1}$ is a morphism of
$\catName{S}^3$, then so is $\morphName{f}^2 \morphCompose \morphName{f}^1$
\end{proof}

\item
If each $\morphName{f}^i \maps \seqname{S}^i \to \seqname{S}^{i+1}$ is
an $\catName{M}^i$-$\catName{M}^{i+1}$ m-morphism of
$\seqname{S}^i$ to $\seqname{S}^{i+1}$ then \\*
\nobreaks
  {{
  $
    \morphName{f}^2 \morphCompose \morphName{f}^1 \maps
    \seqname{S}^1 \to \seqname{S}^3
  $
  }}
is an $\catName{M}^1$-$\catName{M}^3$ m-morphism of
$\seqname{S}^1$ to $\seqname{S}^3$.

\begin{proof}
Since $\catName{M}^1 \subcat[full-] \catName{M}^2$ and
$\catName{M}^2 \subcat[full-] \catName{M}^3$ then
$\catName{M}^1 \subcat[full-] \catName{M}^3$ and
$\morphName{f}^i$ is a morphism of $\catName{S}^{i+1}$ by
\pagecref{def:modMmorph}.
Since
\nobreaks
  {{
    $\catName{M}^2 \subcat[full-] \catName{M}^3$,
  }}
$\morphName{f}^1$ is a
morphism of $\catName{S}^3$. and thus
$\morphName{f}^2 \morphCompose \morphName{f}^1 \maps S^1 \to S^3$ is a morphism
of $\catName{S}^3$.
\end{proof}
\end{enumerate}
\end{lemma}

\section{Spaces and proper functions}
\label{sec:spaces}
Several of the following definitions involve spaces of a restrictive
character and specific types of mappings among them. Except where
otherwise qualified, the word \textit{space} will have this restricted
meaning.

\begin{definition}[Spaces]
\label{def:space}
A space is a topological space, a model space or either with an
additional associated structure.
\end{definition}

\begin{definition}[Proper functions]
\label{def:proper}
Let $S^1$, $S^2$ be spaces and $\morphName{f} \maps S^1 \to S^2$ a
continuous function; $\morphName{f}$ need not preserve any associated
algebraic structure\footnote{
  However, in practice commutation relations will often enforce the
  preservation of algebraic structures.
}.
$\morphName{f}$ is a proper\footnote{
  In the spirit of \cite[footnote, p.~112] {GenTop}
  \begin{quote}
    This nomenclature is an excellent example of the time-honored
    custom of referring to a problem we cannot handle as abnormal,
    irregular, improper, degenerate, inadmissible, and otherwise
    undesirable.
  \end{quote}
}
function iff

\begin{enumerate}
\item $S^1$ and $S^2$ are both $\truthspace$ or both
$\seqname{Truthspace}$, and $\morphName{f}(\true)=\true$.

\item $S^1$, $S^2$ are topological spaces other than
$\truthspace$.

\item $S^1$ is a topological space other than
$\truthspace$, $S^2$ is a model space other than
$\seqname{Truthspace}$ and the images of open sets are contained in
model neighborhoods.

\item $S^1$ is a model space other than $\seqname{Truthspace}$, $S^2$ is
a topological space other than $\truthspace$ and the
inverse images of open sets are model neighborhoods.

\item $S^1$, $S^2$ are both model spaces other than
$\seqname{Truthspace}$ and $\morphName{f}$ is a model function.
\end{enumerate}
\end{definition}

\begin{definition}[Singleton categories and sequences of singleton categories]
\label {def:singcat}
Let $S$ be a topological space. Then the singleton category of $S$,
abbreviated $\singcat{S}$, is the category whose sole object is $S$ and
whose sole morphism is the identity morphism of $S$ to itself.

Let $\seqname{S}$ be a model space. Then the singleton category of
$\seqname{S}$, abbreviated $\singcat{\seqname{S}}$, is the category whose
sole object is $\seqname{S}$ and whose sole morphism is the identity
morphism of $\seqname{S}$ to itself.

Let $\seqname{S} \defeq (S^\alpha, \alpha \prec \Alpha)$ be a sequence
of spaces. Then the singleton category sequence of $\seqname{S}$,
abbreviated $\Singcat{\seqname{S}}$, is
$(\singcat{S^\alpha}, \alpha \prec \Alpha)$.

Let $\seqname{S}$ be a set of spaces. Then the singleton category of
$\seqname{S}$, abbreviated $\Singcat{\seqname{S}}$, is
$\union[S \in \seqname{S}]{\singcat{S}}$.

\begin{remark}
By abuse of language the notation $\Singcat{\catSeqName{S}}$ will be used
to name sequences of categories constructed with this and similar
functions.
\end{remark}
\end{definition}

% Let $\seqname{M}=(S, \catName{S})$ be a model space for $S$ and $P\subset \catName{S}$. We say that $P$ is a
% patching for $\catName{S}$, the elements of $P$ are patches and that $\seqname P \defeq
% (S, \catName{S},P)$ is a patched space iff
%
% \begin{enumerate}
% \item $P$ is a cover for  $S$.
% \item $P$ is a basis for the finite intersections of elements of  $P$.
% \end{enumerate}

\part{Signatures}
\label{part:sig}
Given a sequence of spaces, we need a way to characterise a function
that takes arguments in those spaces and has a value in them. For this
purpose we use a sequence of ordinals where the last ordinal in the
sequence identifies the space containing the function's value.

\section{Definitions of signatures}
\begin{definition}[Signature]
Let $\seqname{S} \defeq (S_\alpha, \alpha \prec \Alpha)$ be a sequence
and $\sigma \defeq (\alpha_\beta \prec \Alpha, \beta \preceq \Beta)$ be
a sequence of ordinals in $\Alpha$. Then $\sigma$ is a signature over
$\seqname{S}$.
\end{definition}

\begin{definition}[$\seqname{S}$-signature of function]
Let $\seqname{S} \defeq (S_\alpha, \alpha \prec \Alpha)$ be a
sequence,
$\sigma \defeq (\alpha_\beta \prec \Alpha, \beta \preceq \Beta)$ be a
signature over $\seqname{S}$ and
$\morphName{f} \maps (S_{\alpha_\beta}, \beta \prec \Beta) \to
S_{\alpha_\Beta}$ continuous.
Then $\sigma$ is the $\seqname{S}$-signature for $\morphName{f}$ and
$\morphName{f}$ has $\seqname{S}$-signature $\sigma$.

\begin{remark}
If $n \in \mathbb{N}$ and $\sigma \defeq (\alpha, \beta \preceq n)$
then $\morphName{f}$ is an n-ary operation over $S_\alpha$.
\end{remark}
$\morphName{f}$ is proper iff either
\begin{enumerate}
\item
$S_{\alpha_\beta}$, $\beta \preceq \Beta)$ are all topological spaces.
\item $S_{\alpha_\beta}$, $\beta \preceq \Beta)$ are all model spaces.
\end{enumerate}
\end{definition}

\begin{definition}[$\catSeqName{S}$-signature of function]
Let $\seqname{S} \defeq (S_\alpha, \alpha \prec \Alpha)$ be a
sequence,
$\catSeqName{S} \defeq (\catName{S}_\alpha, \alpha \prec \Alpha)$ be a
sequence of categories with
$
  \uquant%
    {\alpha \prec \Alpha}
    {S_\alpha \objin \catName{S}_\alpha}
$,
$\sigma \defeq (\alpha_\beta \prec \Alpha, \beta \preceq \Beta)$ be a
signature over $\catSeqName{S}$ and
$\morphName{f} \maps (S_{\alpha_\beta}, \beta \prec \Beta) \to
S_{\alpha_\Beta}$ continuous.
Then $\sigma$ is the $\catSeqName{S}$-signature for $\morphName{f}$ and
$\morphName{f}$ has $\catSeqName{S}$-signature $\sigma$.
\end{definition}

\begin{definition}[$\seqname{S}$-signature of function sequence]
Let \\
$\Sigma \defeq (\sigma_\gamma, \gamma \prec \Gamma) \defeq
\bigl ( (\sigma_{\gamma,{\alpha_\beta}}, \alpha_\beta \prec \Alpha,
\beta \preceq \Beta_\gamma), \gamma \prec \Gamma \bigr )$
be a sequence of signatures over
$\seqname{S} \defeq (S_\alpha, \alpha \prec \Alpha)$,
$\morphSeqName{F}=(\morphName{F}_\gamma, \gamma \prec \Gamma)$ a
sequence of functions where $\morphName{F}_\gamma$ has
$\seqname{S}$-signature $\sigma_\gamma$, then $\Sigma$ is the
$\seqname{S}$-signature for $\morphSeqName{F}$ and
$\morphSeqName{F}$ has $\seqname{S}$-signature $\Sigma$.
\end{definition}

\begin{definition}[$\catSeqName{S}$-signature of function sequence]
Let \\
$\Sigma \defeq (\sigma_\gamma, \gamma \prec \Gamma) \defeq
\bigl ( (\sigma_{\gamma,{\alpha_\beta}},
\alpha_\beta \prec \Alpha,
\beta \preceq \Beta_\gamma), \gamma \prec \Gamma \bigr )$
be a sequence of signatures over
$\catSeqName{S} \defeq (\catName{S}_\alpha, \alpha \prec \Alpha)$,
$\morphSeqName{F}=(\morphName{F}_\gamma, \gamma \prec \Gamma)$
a sequence of functions where $\morphName{F}_\gamma$ has
$\catSeqName{S}$-signature $\sigma_\gamma$,
then $\Sigma$ is the
$\catSeqName{S}$-signature for $\morphSeqName{F}$ and
$\morphSeqName{F}$ has $\catSeqName{S}$-signature $\Sigma$.
\end{definition}

\section{Functions associated with signatures and function sequences}
\label{sec:sig}
When characterizing functions with signatures, several auxillary
functions are helpful.

\begin{definition}[Signature operators]
\label{def:sigop}
Let $\catSeqName{S}^i \defeq (\catName{S}^i_\alpha, \alpha \prec \Alpha)$, $i=1,2$,
be sequences of categories with $\seqname{S}^i \defeq (S^i_\alpha, \alpha \prec
\Alpha)$ sequences of sets in those categories, $\morphName{f}^i,i=1,2$ functions with
$\catSeqName{S}^i$-signatures $\sigma = (\alpha_\beta \prec \Alpha, \beta \preceq \Beta)$,
$\morphSeqName{g} \defeq (\morphName{g}_\alpha \maps S^1_\alpha
\to S^2_\alpha, \alpha \prec \Alpha)$
a sequence of functions then

\begin{equation}
\langle \sigma | \seqname{S}^i \rangle
\defeq \overset{\sigma}{\seqname{S}^i} \defeq (S^i_{\alpha_\beta}, \beta \preceq \Beta)
\end{equation}
selects sets from $\seqname{S}^i$ according to signature $\sigma$.

\begin{equation}
\langle \sigma | \morphSeqName{g} \rangle
\defeq \overset{\sigma}{\morphSeqName{g}} \defeq (\morphName{g}_{\alpha_\beta}, \beta \preceq \Beta)
\end{equation}
selects functions from $\morphSeqName{g}$ according to signature $\sigma$.

\begin{equation}
\overset{\sigma}{\underline{\morphSeqName{g}}} \maps
\bigtimes_{\beta \prec \Beta} S^1_{\alpha_\beta} \to
\bigtimes_{\beta \prec \Beta} S^2_{\alpha_\beta}
= \underline{\langle \sigma | \morphSeqName{g} \rangle} =
\bigtimes_{\beta \prec \Beta} \morphName{g}_{\alpha_\beta}
\end{equation}
is the product of the functions selected by all but the last ordinal in signature $\sigma$.

\begin{equation}
\morphSeqName{g} \morphComposeTail[\sigma] \morphName{f}^1 \defeq
\tail(\langle \sigma | \morphSeqName{g} \rangle) \morphCompose \morphName{f}^1
\end{equation}
composes the last function selected from $\morphSeqName{g}$ by signature
$\sigma$ with $\morphName{f}^1$.

\begin{equation}
  \morphName{f}^2 \morphComposeHead[\sigma] \morphSeqName{g} \defeq
  \morphName{f}^2 \morphCompose \underline{\overset{\sigma}{\morphSeqName{g}}}
\end{equation}
composes $\morphName{f}$ with all but the last function in
$\morphSeqName{g}$ selected by signature $\sigma$.
\end{definition}

\begin{figure}
\[ \bfig
\node s1ab(0,0)[\underline{\overset{\sigma}{\seqname{S}^1}}]
\node s2ab(0,-1000)[\underline{\overset{\sigma}{\seqname{S}^2}}]
\node s2aB(1000,-1000)[\tail(\overset{\sigma}{\seqname{S}^2})]
\arrow |l|[s1ab`s2ab;%
  {%
    \underline{\overset{\sigma}{\morphSeqName{g}}} \defeq%
    \bigtimes \head(\overset{\sigma}{\morphSeqName{g}})
  }]
\arrow |r|[s1ab`s2aB;{{\morphName{f}_2 \morphComposeHead[\sigma] \morphSeqName{g}}}]
\arrow |r|[s2ab`s2aB;\morphName{f}_2]
\efig \]
% \[\begin{tikzcd}
%   {\underline{\overset{\sigma}{\seqname{S}^1}}}
%     \ar[d,"{\underline{\overset{\sigma}{\morphSeqName{g}}}}"']
%     \ar[dr, "{\morphName{f}_2 \morphComposeHead[\sigma] \morphSeqName{g}}"]
% \\
%    {\underline{\overset{\sigma}{\seqname{S}^2}}}
%     \ar[r,"\morphName{f}_2"']
%    &
%    {\tail(\overset{\sigma}{\seqname{S}^2})}
% \end{tikzcd}\]
\caption{$\morphName{f}_2 \morphComposeHead[\sigma] \morphSeqName{g}$}
\label{fig:fgsigma}
\end{figure}

\begin{figure}
\[ \bfig
\node s1ab(0,0)[\underline{\overset{\sigma}{\seqname{S}^1}}]
\node s1aB(0,-1000)[\tail(\overset{\sigma}{\seqname{S}^1})]
\node s2aB(1000,-1000)[\tail(\overset{\sigma}{\seqname{S}^2})]
\arrow |l|[s1ab`s1aB;\morphName{f}_1]
\arrow |r|[s1ab`s2aB;{{\morphSeqName{g} \morphComposeTail[\sigma] \morphName{f}_1}}]
\arrow |b|[s1aB`s2aB;\tail(\overset{\sigma}{\morphSeqName{g}})]
\efig \]
% \[\begin{tikzcd}
%   {\underline{\overset{\sigma}{\seqname{S}^1}}}
%     \ar[d,"{\morphName{f}_1}"']
%     \ar[dr, "{\morphSeqName{g} \morphComposeTail[\sigma] \morphName{f}_1}"]
% \\
%   {\tail(\overset{\sigma}{\seqname{S}^1})}
%     \ar[r, "{\tail(\overset{\sigma}{\morphSeqName{g}})}"']
%   &
%   {\tail(\overset{\sigma}{\seqname{S}^2})}
% \end{tikzcd}\]
\caption{$\morphSeqName{g} \morphComposeTail[\sigma] \morphName{f}_1$}
\label{fig:gsigmaf}
\end{figure}

\section{Commutative diagrams with signatures}
\label{sec:cdsig}
The conventional usage of the word commutative diagram is for a
naturality condition on functions of a single arguement. It is
convenient to extend the nomenclature to functions of multiple arguments
and to sequences of such functions. This is convenient, e.g., as a means
of specifying that functions preserve algebraic structures.

\begin{definition}[$\Sigma$-commutation]
\label{def:Sigmacomm}
Let
$\catSeqName{S}^i \defeq (\catName{S}^i_\alpha, \alpha \prec \Alpha)$,
$i=1,2$, be a sequence of categories, with
$
  \seqname{S}^i \defeq (S^i_\alpha, \alpha \prec \Alpha)
  \seqin \catSeqName{S}^i
$
a sequence of spaces,
$\morphSeqName{F}^i \defeq (\morphName{F}^i_\gamma, \gamma \prec \Gamma)$
be a sequence of functions with
$\catSeqName{S}^i$-signature
$
  \Sigma \defeq (\sigma_{\gamma \prec \Gamma}) \defeq \\
    \bigl (
      (
        \sigma_{\gamma,\alpha_\beta}, \alpha_\beta \prec \Alpha,
        \beta \preceq \Beta_\gamma
      ),
      \gamma \prec \Gamma
    \bigr )
$,
$
  \morphSeqName{f} \defeq
  (
    \morphName{f}_\alpha \maps S^1_\alpha \to S^2_\alpha,
    \alpha \prec \Alpha
  )
$
a sequence of functions satisfying the naturality condition
$
  \morphSeqName{f} \morphComposeTail[\sigma_\gamma] \morphName{F}^1_\gamma
  = \morphName{F}^2_\gamma \morphComposeHead[\sigma_\gamma] \morphSeqName{f},
  \gamma \prec \Gamma
$,
then
$\morphSeqName{f}$ $\Sigma$-commutes with the function sequences
$\morphSeqName{F}^1, \morphSeqName{F}^2$.
\end{definition}

\begin{remark}
If all of the $\sigma_{\gamma,\alpha_\beta}$ are equal to the same
ordinal $\hat{\alpha}$ and each $\Beta_\gamma$ is finite then
$\set{{(S^i_{\hat{\alpha}},\morphName{F}^i_\gamma)}}[{\gamma \prec \Gamma}]$ is
an $\Omega$-algebra and $\morphSeqName{f}$ is an $\Omega$-homomorphism.
\end{remark}

\begin{figure}
\[ \bfig
\node s1b(0,0)[\underline{\overset{\sigma_\gamma}{\seqname{S}^1_\gamma}}]
\node s1B(1000,0)[\tail \left ( \overset{\sigma_\gamma}{\seqname{S}^1_\gamma} \right )]
\node s2b(0,-1000)[\underline{\overset{\sigma_\gamma}{\seqname{S}^2_\gamma}}]
\node s2B(1000,-1000)[\tail \left ( \overset{\sigma_\gamma}{\seqname{S}^2_\gamma} \right )]
\arrow |a|[s1b`s1B;F^1_\gamma]
\arrow |l|[s1b`s2b;\underline{\overset{\sigma_\gamma}{\morphSeqName{f}}} ]
\arrow |r|[s1B`s2B;\tail \left ( \overset{\sigma_\gamma}{\morphSeqName{f}} \right )]
\arrow |a|[s2b`s2B;F^2_\gamma]
\efig \]
\caption{$\morphSeqName{f}$ $\Sigma$-commutes with the function sequences $\morphSeqName{F}^1, \morphSeqName{F}^2$}
\label{fig:fF1F2}
\end{figure}

\begin{lemma}[$\Sigma$-commutation]
\label{lem:Sigmacomm}
Let
$
  \catSeqName{S}^i \defeq (\catName{S}^i_\alpha, \alpha \prec \Alpha)
$,
$i \in [1,3]$, be a sequence of categories, with
$
  \seqname{S}^i \defeq (S^i_\alpha, \alpha \prec \Alpha)
  \seqin \catSeqName{S}^i
$
a sequence of spaces and
$\morphSeqName{F}^i \defeq (\morphName{F}^i_\gamma, \gamma \prec \Gamma)$
be a sequence of functions with $\catSeqName{S}^i$-signature
$
  \Sigma \defeq (\sigma_\gamma, \gamma \prec \Gamma) \defeq
  \bigl (
    (
      \sigma_{\gamma,\alpha_\beta}, \alpha_\beta \prec \Alpha,
      \beta \preceq \Beta_\gamma
    ),
    \gamma \prec \Gamma
  \bigr )
$.

Let
$
  \morphSeqName{f}^i \defeq
  (
    \morphName{f}^i_\alpha \maps S^i_\alpha \to S^{i+1}_\alpha,
    \alpha \prec \Alpha
  )
$,
$i=1,2$, $\Sigma$-commute with $\morphSeqName{F}^i,\morphSeqName{F}^{i+1}$.
Then $\morphSeqName{f}^2 \morphCompose[()] \morphSeqName{f}^1$ $\Sigma$-commutes
with $\morphSeqName{F}^1,\morphSeqName{F}^3$.

\begin{figure}
\[ \bfig
\node s1b(0,2000)[\underline{\overset{\sigma_\gamma}{\seqname{S}^1_\gamma}}]
\node s1B(1000,2000)[\tail \left ( \overset{\sigma_\gamma}{\seqname{S}^1_\gamma} \right )]
\node s2b(0,1000)[\underline{\overset{\sigma_\gamma}{\seqname{S}^2_\gamma}}]
\node s2B(1000,1000)[\tail \left ( \overset{\sigma_\gamma}{\seqname{S}^2_\gamma} \right )]
\node s3b(0,0)[\underline{\overset{\sigma_\gamma}{\seqname{S}^3_\gamma}}]
\node s3B(1000,0)[\tail \left ( \overset{\sigma_\gamma}{\seqname{S}^3_\gamma} \right )]
\arrow |a|[s1b`s1B;F^1_\gamma]
\arrow |l|[s1b`s2b;\underline{\overset{\sigma_\gamma}{\morphSeqName{f}^1}} ]
\arrow |r|[s1B`s2B;\tail \left ( \overset{\sigma_\gamma}{\morphSeqName{f}^1} \right )]
\arrow |a|[s2b`s2B;F^2_\gamma]
\arrow |l|[s2b`s3b;\underline{\overset{\sigma_\gamma}{\morphSeqName{f}^2}} ]
\arrow |r|[s2B`s3B;\tail \left ( \overset{\sigma_\gamma}{\morphSeqName{f}^2} \right )]
\arrow |a|[s3b`s3B;F^3_\gamma]
\efig \]
\caption{$\morphSeqName{f}^2 \morphCompose[()] \morphSeqName{f}^1$ $\Sigma$-commute with $\morphSeqName{F}^1,\morphSeqName{F}^3$}
\label{fig:fF1F2F3}
\end{figure}

\begin{proof}
The result follows by diagram chasing in \cref{fig:fF1F2F3}.
\end{proof}

Let
$\seqname{S}^1 \SUBSETEQ \seqname{S}^2$
and
$
  \morphName{F}^1_\gamma =
  \morphName{F}^2_\gamma \maps
    \bigtimes \underline{\overset{\sigma_\gamma}{\seqname{S}^1}} \to
    \tail \left ( \overset{\sigma_\gamma}{\seqname{S}^1} \right )
$.
Then $\ID[{{\seqname{S}^1},{\seqname{S}^2}}]$ $\Sigma$-commutes with
$\morphSeqName{F}^1,\morphSeqName{F}^2$.

\begin{proof}
The result follows from \pagecref{def:sigop} and
{
  \showlabelsinline
  \pagecref{def:Sigmacomm}
}:
\begin{equation*}
\begin{split}
    \ID[{\seqname{S}^1},{\seqname{S}^2}] \morphComposeTail[\sigma_\gamma]
    \morphName{F}^1_\gamma
  & =
    \Id%
      [
        S^1_{\gamma,\alpha_{\Beta_\gamma}},
        S^2_{\gamma,\alpha_{\Beta_\gamma}}
      ]
    \morphCompose
    \morphName{F}^1_\gamma
\\*
  & =
    \morphName{F}^2_\gamma \restrictto_%
      {
        \bigtimes \underline{\overset{\sigma_\gamma}{\seqname{S}^1}},
      }
\\*
  & =
    \morphName{F}^2_\gamma \morphCompose
    \bigtimes_{\beta \prec \Beta_\gamma}
      \Id[{S^1_{\alpha_\beta}},{S^2_{\alpha_\beta}}]
\\*
  & =
    \morphName{F}^2_\gamma \morphComposeHead[\sigma_\gamma]
    \ID[{{\seqname{S}^1},{\seqname{S}^2}}]
\end{split}
\end{equation*}
\end{proof}
\end{lemma}

\begin{corollary}[$\Sigma$-commutation]
\label{cor:Sigmacomm}
Let
$
  \catSeqName{S} \defeq (\catName{S}_\alpha, \alpha \prec \Alpha)
$,
be a sequence of categories, with
$
  \seqname{S} \defeq (S_\alpha, \alpha \prec \Alpha)
  \seqin \catSeqName{S}
$
a sequence of spaces and
$\morphSeqName{F} \defeq (\morphName{F}_\gamma, \gamma \prec \Gamma)$
be a sequence of functions with $\catSeqName{S}$-signature
$
  \Sigma \defeq (\sigma_\gamma, \gamma \prec \Gamma) \defeq
  \bigl (
    (
      \sigma_{\gamma,\alpha_\beta}, \alpha_\beta \prec \Alpha,
      \beta \preceq \Beta_\gamma
    ),
    \gamma \prec \Gamma
  \bigr )
$.
Then $\ID[{\seqname{S}}]$ $\Sigma$-commutes with
$\morphSeqName{F},\morphSeqName{F}$.

\begin{proof}
$\seqname{S} \SUBSETEQ \seqname{S}$,
$
  \morphName{F}_\gamma =
  \morphName{F}_\gamma \maps
    \bigtimes \underline{\overset{\sigma_\gamma}{\seqname{S}}} \to
    \tail \left ( \overset{\sigma_\gamma}{\seqname{S}} \right )
$
and $\ID[{\seqname{S}}] = \ID[{\seqname{S}},{\seqname{S}}]$.
\end{proof}
\end{corollary}

\part{Prestructures}
\label{part:pre}
Prestructures and prestructure morphisms are generalizations of
$\Omega$-algebras and $\Omega$-homomorphisms, and are notational
conveniences to simplify imposing commutation relations on multiple
unrelated functions of multiple variables.

\begin{definition}[Prestructures]
\label{def:pre}
Let $\catSeqName{S} \defeq (\catName{S}_\alpha, \alpha \prec \Alpha)$ be
a sequence of categories,
$
  \seqname{S} \defeq (S_\alpha, \alpha \prec \Alpha)
$
a sequence of spaces,
$\morphSeqName{F}=(F_\gamma, \gamma \prec \Gamma)$
a sequence of continuous functions and
$
  \Sigma \defeq
  (\sigma_\gamma, \gamma \prec \Gamma) \defeq
  \bigl (
    (
      \sigma_{\gamma.{\alpha_\beta}},
      \beta \preceq \Beta_\gamma
    ),
    \gamma \prec \Gamma
  \bigr )
$
a sequence of signatures.
Then $\seqname{P} \defeq (\catSeqName{S}, \seqname{S}, \Sigma, F)$ is a
$\catSeqName{S}-\Sigma$ prestructure iff
\begin{enumerate}
\item $\catName{S}_\alpha$ is either a topological category or a model category.
\item $\seqname{S} \seqin \catSeqName{S}$.
\item $\morphSeqName{F}$ has $\catSeqName{S}$-signature $\Sigma$.
\end{enumerate}
\end{definition}

\begin{lemma}[Prestructures]
\label{lem:pre}
Let

\begin{enumerate}
\item
$\catSeqName{S}^i$, $i=1,2$, be a sequence of
categories
\item
\nobreaks
  {{
    $\catSeqName{S}^1 \SUBCAT \catSeqName{S}^2$,
  }}
\item
$\seqname{S} \seqin \catSeqName{S}^1$
\item
$\morphSeqName{F}$ a sequence of continuous functions with
$\catSeqName{S}^1$-signatures $\Sigma$
\item
$\seqname{P}^1 \defeq (\catSeqName{S}^1, \seqname{S}, \Sigma, F)$ a
$\catSeqName{S}^1-\Sigma$ prestructure
\end{enumerate}

Then
$\seqname{P}^2 \defeq (\catSeqName{S}^2, \seqname{S}, \Sigma, F)$ is a
$\catSeqName{S}^2-\Sigma$ prestructure.

\begin{proof}
$\seqname{S} \seqin \catSeqName{S}^2$ by \pagecref{lem:seqin}.
The other conditions do not depend on the category sequence.
\end{proof}

Let

\begin{enumerate}
\item
$\catSeqName{S}^i$, $i=1,2$, be a sequence of categories
\item
$\catSeqName{S}^1 \SUBCAT \catSeqName{S}^2$
\item
$\seqname{S}^i \seqin \catSeqName{S}^i$
\item
$\seqname{S}^1 \SUBSETEQ \seqname{S}^2$
\item
$
  \morphSeqName{F}^2 \defeq
  (\morphName{F}^2_\gamma, \gamma \prec \Gamma)
$
be a sequence of continuous functions with $\catSeqName{S}^2$-signatures
$
  \Sigma \defeq
  (\sigma_\gamma, \gamma \prec \Gamma) \defeq
  \bigl (
    (
      \sigma_{\gamma.{\alpha_\beta}},
      \beta \preceq \Beta_\gamma
    ),
    \gamma \prec \Gamma
  \bigr )
$
\item
$
  \morphName{F}^2_\gamma
    \left [
      \bigtimes \underline{\overset{\sigma_\gamma}{\seqname{S}^1}}
    \right ]
  \subseteq
  \tail \left ( \overset{\sigma_\gamma}{\seqname{S}^1} \right )
$
\item
$
  \morphSeqName{F}^1 \defeq
  \left (
    \morphName{F}^1_\gamma \defeq
        \morphName{F}^2_\gamma \maps
          \bigtimes \underline{\overset{\sigma_\gamma}{\seqname{S}^1}}
          \to
          \tail \left ( \overset{\sigma_\gamma}{\seqname{S}^1} \right ),
        \gamma  \prec \Gamma
  \right )
$
\item
$
  \seqname{P}^2 \defeq
  (\catSeqName{S}^2, \seqname{S}^2, \Sigma, \morphSeqName{F}^2)
$
a $\catSeqName{S}^2-\Sigma$ prestructure.
\end{enumerate}

Then
$
  \seqname{P}^1 \defeq
  (\catSeqName{S}^1, \seqname{S}^1, \Sigma, \morphSeqName{F}^1)
$
is a $\catSeqName{S}^1-\Sigma$ prestructure.

\begin{proof}
$\seqname{S}^1 \seqin \catSeqName{S}^1$ by hypothesis.
$\morphSeqName{F}^1$ is continuous and has $\catSeqName{S}^1$-signatures
$\Sigma$ by construction.
\end{proof}
\end{lemma}


\begin{definition}[Morphisms of prestructures]
\label{def:premorph}
Let $\catSeqName{S}^i=(\catName{S}^i_\alpha, \alpha \prec \Alpha)$,
$i=1,2$, be a sequence of categories,
$
  \seqname{S}^i \defeq (S^i_\alpha \alpha \prec \Alpha)
  \seqin \catSeqName{S}^i
$
be a sequence of spaces, \\*
$\morphSeqName{F}^i=(\morphName{F}^i_\gamma, \gamma \prec \Gamma)$
be a sequence of continuous functions with
$\catSeqName{S}^i$-signatures $\Sigma$,
$
\morphSeqName{f} \defeq
  (
    \morphName{f}_\alpha \maps S^1_\alpha \to S^2_\alpha,
    \alpha \prec \Alpha
  )
$
be a sequence of proper functions and
$
  \seqname{P}^i \defeq
  (\catSeqName{S}^i, \seqname{S}^i, \Sigma, \morphSeqName{F}^i)
$
be a $\catSeqName{S}^i-\Sigma$ prestructure. $\morphSeqName{f}$ is a
(sem-strict,strict) prestructure morphism of
$\seqname{P}^1$ to $\seqname{P}^2$ iff it $\Sigma$-commutes with
$\morphSeqName{F}^1,\morphSeqName{F}^2$, as shown in
\pagecref{fig:fF1F2}.

\begin{remark}
If $\morphSeqName{F}^1 = \morphSeqName{F}^2 = \emptyset$ then
$\morphSeqName{f}$ is a morphism since it vacuously $\Sigma$-commutes
with $\emptyset,\emptyset$.
\end{remark}

It is a semi-strict prestructure morphism iff
$\catSeqName{S}^1 \SUBCAT[full-] \catSeqName{S}^2$ and for each $\alpha$

\begin{description}
\item[If $\catName{S}^i_\alpha$ is a model category:]
$\morphName{f}_\alpha$ is a local
$\catName{S}^1_\alpha$-$\catName{S}^2_\alpha$ m-morphism of $S^1_\alpha$
to $S^2_\alpha$.

\item[If $\catName{S}^i_\alpha$ is a topological category:]
$\morphName{f}_\alpha$ is a local
$\catName{S}^1_\alpha$-$\catName{S}^2_\alpha$ morphism of $S^1_\alpha$
to $S^2_\alpha$.

\item[Otherwise:]
$\morphName{f}_\alpha$ is a morphism of $\catName{S}^2_\alpha$.
\end{description}

It is a strict prestructure morphism iff
$\catSeqName{S}^1 \SUBCAT[full-] \catSeqName{S}^2$ and each function
$\morphName{f}_\alpha$ is a morphism of $\catName{S}^2_\alpha$.

\begin{remark}
It is not sufficient to require that each $\morphName{f}_\alpha$ be a
morphism of $\catName{S}^1_\alpha \unioncat \catName{S}^2_\alpha$
because that would not ensure that the composition of strict
prestructure morphisms is strict.
\end{remark}

If $\seqname{S}^1 \SUBSETEQ \seqname{S}^2$ and
$
  \morphName{F}^1_\gamma =
  \morphName{F}^2_\gamma \maps
    \bigtimes \underline{\overset{\sigma_\gamma}{\seqname{S}^1}} \to
    \tail \left ( \overset{\sigma_\gamma}{\seqname{S}^1} \right )
$
then
$
  \Id[{{\seqname{P}^1},{\seqname{P}^2}}] \defeq
  \ID[{{\seqname{S}^1},{\seqname{S}^2}}]
$.

$
  \Id[{\seqname{P}^i}] \defeq
  \ID[{\seqname{S}^i}] = \Id[{{\seqname{P}^i},{\seqname{P}^i}}]
$.
\end{definition}

For some application, certain spaces of a prestructure are of primary
interest while the rest serve only in an auxilliary role. For those
applications, morphisms that are identical on the primary spaces may
be considered equivalent.

\begin{definition}[$\seqname{K}$-equivalence of prestructure morphisms]
\label{def:premorphKequiv}
Let $\catSeqName{S}^i=(\catName{S}^i_\alpha, \alpha \prec \Alpha)$,
$i=1,2$, be a sequence of categories,
$
  \seqname{P}^i \defeq
  (\catSeqName{S}^i, \seqname{S}^i, \Sigma, \morphSeqName{F}^i)
$
be a $\catSeqName{S}^i-\Sigma$ prestructure,
$
  \morphSeqName{f} \defeq
    (
      \morphName{f}_\alpha \maps S^1_\alpha \to S^2_\alpha,
      \alpha \prec \Alpha
    )
$
and
$
  \morphSeqName{g} \defeq
    (
      \morphName{g}_\alpha \maps S^1_\alpha \to S^2_\alpha,
      \alpha \prec \Alpha
    )
$
be prestructure morphisms of $\seqname{P}^1$ to $\seqname{P}^2$ and
$\seqname{K}$ be a set of ordinals less than $\Alpha$.

$\morphSeqName{f}$ is $\seqname{K}$-equivalent to $\morphSeqName{g}$ as a
prestructure morphism from $\seqname{P}^1$ to $\seqname{P}^2$,
abbreviated $\morphSeqName{f} \maps \seqname{P}^1 \to \seqname{P}^2$ is
$\seqname{K}$-equivalent to
$\morphSeqName{f} \maps \seqname{P}^1 \to \seqname{P}^2$, iff
$
  \uquant
    {\alpha \in \seqname{K}}
    {\morphName{f}_\alpha = \morphName{g}_\alpha}
$.

By abuse of language we write $\morphSeqName{f}$ is
$\seqname{K}$-equivalent to $\morphSeqName{g}$ when the prestructures
$\seqname{P}^1$ and $\seqname{P}^2$ are understood by context.
\end{definition}

\begin{lemma}[Prestructure morphisms]
\label{lem:premorph}
Let $\catSeqName{S}^i=(\catName{S}^i_\alpha, \alpha \prec \Alpha)$,
$i = 1,2$, be a sequence of categories,
$
  \seqname{S}^i \defeq (S^i_\alpha, \alpha \prec \Alpha)
  \seqin \catSeqName{S}^i
$
a sequence of spaces,
$\morphSeqName{F}^i=(\morphName{F}^i_\gamma, \gamma \prec \Gamma)$,
be a sequence of continuous functions with
$\seqname{S}^i$-signatures $\Sigma$ and
$
  \seqname{P}^i \defeq
  (\catSeqName{S}^i, \seqname{S}^i, \Sigma, \morphSeqName{F}^i)
$
be a $\catSeqName{S}^i-\Sigma$ prestructure.

If $\seqname{S}^1 \SUBSETEQ \seqname{S}^2$ and each
$
  \morphName{F}^1_\gamma =
  \morphName{F}^2_\gamma \maps
    \bigtimes \underline{\overset{\sigma_\gamma}{\seqname{S}^1}} \to
    \tail \left ( \overset{\sigma_\gamma}{\seqname{S}^1} \right )
$
then $\Id[{{\seqname{P}^1},{\seqname{P}^2}}]$ is a prestructure morphism from
$\seqname{P}^1$ to $\seqname{P}^2$.

$\ID[{{\seqname{S}^1},{\seqname{S}^2}}]$ is a semi-strict prestructure
morphism from $\seqname{P}^1$ to $\seqname{P}^2$ iff
$\catSeqName{S}^1 \SUBCAT[full-] \catSeqName{S}^2$.

$\ID[{{\seqname{S}^1},{\seqname{S}^2}}]$ is a strict prestructure
morphism from $\seqname{P}^1$ to $\seqname{P}^2$ iff each \\
$\Id[{S^1_\alpha},{S^2_\alpha}] \arin \catName{S}^2_\alpha$ and
$\catSeqName{S}^1 \SUBCAT[full-] \catSeqName{S}^2$.

\begin{proof}
$\ID[{{\seqname{S}^1},{\seqname{S}^2}}]$ $\Sigma$-commutes with
$\morphSeqName{F}^1,\morphSeqName{F}^2$ by
\pagecref{lem:Sigmacomm}.

The hypotheses for strictness are precisely those of
\fullcref{def:premorph}.
\end{proof}

Let
$
  \morphSeqName{f} \defeq
  (
    \morphName{f}_\alpha \maps S^1_\alpha \to S^2_\alpha,
    \alpha \prec \Alpha
  )
$
be a sequence of proper functions,

If $\catSeqName{S}'^i=(\catName{S}'^i_\alpha, \alpha \prec \Alpha)$,
$i = 1,2$, is a sequence of categories,
$\catSeqName{S}^i \SUBCAT \catName{S}'^i$ and
$
  \seqname{P}'^i \defeq
  (\catSeqName{S}'^i, \seqname{S}^i, \Sigma, \morphSeqName{F}^i)
$,
then $\morphSeqName{f}$ is a prestructure morphism of $\seqname{P}'^1$
to $\seqname{P}'^2$ iff $\morphSeqName{f}$ is a
prestructure morphism of $\seqname{P}^1$ to
$\seqname{P}^2$.

If $\catSeqName{S}^i \SUBCAT[full-] \catName{S}'^i$ and
$\morphSeqName{f}$ is a semi-strict (strict) prestructure morphism of
$\seqname{P}^1$ to $\seqname{P}^2$ then $\morphSeqName{f}$ is a
semi-strict (strict) prestructure morphism of $\seqname{P}'^1$ to
$\seqname{P}'^2$.

\begin{proof}
The commutation relations do not depend on the choice of categories.

If $\catSeqName{S}^i \SUBCAT[full-] \catName{S}'^i$ and
$\morphSeqName{f}$ is a semi-strict prestructure morphism of
$\seqname{P}^1$ to $\seqname{P}^2$ then for each $\alpha$, if

\begin{enumerate}
\item Each $\catName{S}^i_\alpha$ is a model category:
$\morphName{f}_\alpha$ is a local
$\catName{S}^1_\alpha$-$\catName{S}^2_\alpha$ m-morphism of
$S^1_\alpha$ to $S^2_\alpha$ and thus a local
$\catName{S}'^1_\alpha$-$\catName{S}'^2_\alpha$ m-morphism of
$S^1_\alpha$ to $S^2_\alpha$.

\item Each $\catName{S}^i_\alpha$ is a topological category:
$\morphName{f}_\alpha$ is a local
$\catName{S}^1_\alpha$-$\catName{S}^2_\alpha$ morphism of
$S^1_\alpha$ to $S^2_\alpha$ and thus a local
$\catName{S}'^1_\alpha$-$\catName{S}'^2_\alpha$ morphism of
$S^1_\alpha$ to $S^2_\alpha$.

\item Otherwise
$\morphName{f}_\alpha$ is a morphism of $\catName{S}^2_\alpha$ and
thus a morphism of $\catName{S}'^2_\alpha$.
\end{enumerate}

If $\catSeqName{S}^i \SUBCAT[full-] \catName{S}'^i$ and
$\morphSeqName{f}$ is a strict prestructure morphism of $\seqname{P}^1$
to $\seqname{P}^2$ then each
$
  \morphName{f}_\alpha \arin \catSeqName{S}^2_\alpha
  \subcat[full-] \catName{S}'^2_\alpha
$,
$\alpha \prec \Alpha$, and thus $\morphSeqName{f}$ is a strict
prestructure morphism of $\seqname{P}'^1$ to $\seqname{P}'^2$.
\end{proof}

$\Id[{{\seqname{P}^i}}]$ is an identity morphism of $\seqname{P}^i$.

\begin{proof}
$\morphSeqName{f} \morphCompose[()] \ID[{\seqname{S}^1}] = \morphSeqName{f}$
and
$\ID[{\seqname{S}^2}] \morphCompose[()] \morphSeqName{f} = \morphSeqName{f}$
\end{proof}
\end{lemma}

\begin{corollary}[Prestructure morphisms]
\label{cor:premorph}
Let $\catSeqName{S}=(\catName{S}_\alpha, \alpha \prec \Alpha)$
be a sequence of categories,
$
  \seqname{S} \defeq (S_\alpha \alpha \prec \Alpha)
  \seqin \catSeqName{S}
$
be a sequence of spaces,
$\morphSeqName{F}=(\morphName{F}_\gamma, \gamma \prec \Gamma)$
be a sequence of continuous functions with $\catSeqName{S}$-signatures
$\Sigma$ and
$
  \seqname{P} \defeq
  (\catSeqName{S}, \seqname{S}, \Sigma, \morphSeqName{F})
$
be a $\catSeqName{S}-\Sigma$ prestructure. Then $\Id[{\seqname{P}}]$ is a
strict prestructure morphism from $\seqname{P}$ to $\seqname{P}$.

\begin{proof}
$\seqname{S} \SUBSETEQ \seqname{S}$ and
$
  \morphName{F}_\gamma =
  \morphName{F}_\gamma \restrictto_%
    {
      \bigtimes \underline{\overset{\sigma_\gamma}{\seqname{S}}},
      \tail \left ( \overset{\sigma_\gamma}{\seqname{S}} \right )
    }
$.

Each
$\Id[{S_\alpha},{S_\alpha}] = \Id[{S_\alpha}] \arin \catName{S}_\alpha$ and
$\catSeqName{S} \SUBCAT[full-] \catSeqName{S}$.
\end{proof}
\end{corollary}

\begin{lemma}[Composition of prestructure morphisms]
\label{lem:premorphcomp}
Let
\begin{enumerate}
\item
$\catSeqName{S}^i=(\catName{S}^i_\alpha, \alpha \prec \Alpha)$,
$i \in [1,4]$, be a sequence of categories
\item
$
  \seqname{S}^i \defeq (S^i_\alpha, \alpha \prec \Alpha)
  \seqin \catSeqName{S}^i
$
be a sequence of spaces
\item
$\morphSeqName{F}^i=(\morphName{F}^i_\gamma, \gamma \prec \Gamma)$,
be a sequence of continuous functions with
$\seqname{S}^i$-signatures $\Sigma$
\item
$
  \seqname{P}^i \defeq
  (\catSeqName{S}^i, \seqname{S}^i, \Sigma, \morphSeqName{F}^i)
$
be a $\catSeqName{S}^i-\Sigma$ prestructure
\item
$
  \morphSeqName{f}^i \defeq
  (
    \morphName{f}^i_\alpha \maps S^i_\alpha \to S^{i+1}_\alpha,
    \alpha \prec \Alpha
  )
$
and
$
  \morphSeqName{g}^i \defeq
  (
    \morphName{g}^i_\alpha \maps S^i_\alpha \to S^{i+1}_\alpha,
    \alpha \prec \Alpha
  )
$
$i=1,2,3$, be
(sem-strict,strict) prestructure morphisms of $\seqname{P}^i$ to
$\seqname{P}^{i+1}$.
\end{enumerate}

$\morphSeqName{f}^{i+1} \morphCompose[()] \morphSeqName{f}^i$, $i \in [1,3]$,
is a prestructure morphism and
$
    \morphSeqName{f}^3 \morphCompose[()]
    \bigl ( \morphSeqName{f}^2 \morphCompose[()] \morphSeqName{f}^1 \bigr )
  = \bigr ( \morphSeqName{f}^3 \morphCompose[()] \morphSeqName{f}^2 \bigr )
    \morphCompose[()] \morphSeqName{f}^1
$.

\begin{proof}
If
$\morphSeqName{f}^1$ $\Sigma$-commutes with $\morphSeqName{F}^1,\morphSeqName{F}^2$ and
$\morphSeqName{f}^2$ $\Sigma$-commutes with $\morphSeqName{F}^2,\morphSeqName{F}^3$
then by \pagecref{lem:Sigmacomm}
$\morphSeqName{f}^2 \morphCompose[()] \morphSeqName{f}^1$ $\Sigma$-commutes with $\morphSeqName{F}^1,\morphSeqName{F}^3$.
Thus $\morphSeqName{f}^2 \morphCompose[()] \morphSeqName{f}^1$ is a morphism.

Associativity is just \pagecref{lem:seqfunc}.
\end{proof}

Let each $\morphSeqName{f}^i$ be a semi-strict prestructure morphism.
Then $\morphSeqName{f}^2 \morphCompose[()] \morphSeqName{f}^1$ is a semi-strict
prestructure morphism.

\begin{proof}
If $\catSeqName{S}^1 \SUBCAT[full-] \catSeqName{S}^2$ and
$\catSeqName{S}^2 \SUBCAT[full-] \catSeqName{S}^3$ then
$\catSeqName{S}^1 \SUBCAT[full-] \catSeqName{S}^3$.
For each $\alpha \prec \Alpha$:

\begin{enumerate}
\item Each $\catName{S}^i_\alpha$, $i=1,2,3$, is a model category:
\newline
Each $\morphName{f}^i_\alpha$ is a local
$\catName{S}^1_\alpha$-$\catName{S}^{i+1}_\alpha$ m-morphism of
$S^i_\alpha$ to $S^{i+1}_\alpha$.
Then $\morphName{f}^2_\alpha \morphCompose \morphName{f}^1_\alpha$ is a local
$\catName{S}^i_\alpha$-$\catName{S}^{i+2}_\alpha$ m-morphism of
$S^i_\alpha$ to $S^{i+2}_\alpha$ by
\pagecref{lem:modMmorphcomp}.

\item Each $\catName{S}^i_\alpha$, $i=1,2,3$, is a topological category:
\newline
Each $\morphName{f}^i_\alpha$ is a local
$\catName{S}^i_\alpha$-$\catName{S}^{i+1}_\alpha$ morphism of
$S^i_\alpha$ to $S^{i+1}_\alpha$.
Then $\morphName{f}^2_\alpha \morphCompose \morphName{f}^1_\alpha$ is a local
$\catName{S}^i_\alpha$-$\catName{S}^{i+2}_\alpha$ morphism of
$S^i_\alpha$ to $S^{i+2}_\alpha$ by \pagecref{lem:topLocal}.

\item Otherwise
\newline
Each $\morphName{f}^i_\alpha$ is a morphism of $\catName{S}^{i+1}_\alpha$.
Then $\morphName{f}^2_\alpha \morphCompose \morphName{f}^1_\alpha$ is a
morphism of $\catName{S}^{i+2}_\alpha$.
\end{enumerate}
\end{proof}

Let each $\morphSeqName{f}^i$ be a strict prestructure morphism. Then
$\morphSeqName{f}^2 \morphCompose[()] \morphSeqName{f}^1$ is a strict
prestructure morphism.

\begin{proof}
If $\catSeqName{S}^1 \SUBCAT[full-] \catSeqName{S}^2$ and
$\catSeqName{S}^2 \SUBCAT[full-] \catSeqName{S}^3$ then
$\catSeqName{S}^1 \SUBCAT[full-] \catSeqName{S}^3$.

If $\morphName{f}^1_\alpha$ is a morphism of
$\catName{S}^2_\alpha \subcat[full-] \catName{S}^3_\alpha$ and
$\morphName{f}^2_\alpha$ is a morphism of $\catName{S}^3_\alpha$ then
$\morphName{f}^2_\alpha \morphCompose \morphName{f}^1_\alpha$ is a morphism of
$\catName{S}^3_\alpha$.
\end{proof}

If each $\morphSeqName{f}^i$ is $\seqname{K}$-equivalent to
$\morphSeqName{g}^i$ then each
$\morphSeqName{f}^{i+1} \morphCompose[()] \morphSeqName{f}^i$ is
$\seqname{K}$-equivalent to
$\morphSeqName{g}^{i+1} \morphCompose[()] \morphSeqName{g}^i$.

\begin{proof}
$
  \uquant
    {{i \in {[1-4]}},{\alpha \in \seqname{K}}}
    {\morphName{f}^i_\alpha = \morphName{g}^i_\alpha}
$
by \pagecref{def:premorphKequiv}.
Then
$
  \uquant
    {{i \in {[[1-3]}},{\alpha \in \seqname{K}}}
    {
      \morphName{f}^2_\alpha \morphCompose \morphName{f}^1_\alpha
      =
      \morphName{g}^2_\alpha \morphCompose \morphName{g}^1_\alpha
    }
$
\end{proof}
\end{lemma}

\part{M-charts and m-atlases}
\label{part:M-charts}
The literature defines fiber bundles and manifolds using the language of
charts, atlases and transition functions; it has multiple equivalent
definitions. Some authors start with topological spaces and define
atlases over them. Some start with abstract sets and define atlases over
them, deriving the topology from the atlas. Some start with an indexed
set of open patches in a coordinate space and transition functions among
them satisfying a cocycle (compatibility) constraint, and then derive
the total space as a quotient space of the disjoin union of the patches.
Some authors use maximal atlases while others use equivalence classes of
atlases.

This paper uses the first approach, explicitly making the relevant
spaces topological spaces, but modifies the definitions of charts in
order to make them fit more natuarally into the context of local
coordinate spaces. It adds a prefix, e.g., $\Ck$, M, in order to avoid
confusion with the conventional definitions. An m-atlas based on a
topological spaces is closer to the conventional definitions of an atlas
for a manifold while an m-atlas based on a model space is suitable for
defining both manifolds and fiber bundles. Although simple manifolds and
fiber bundles could both be defined directly in terms of maximal
m-atlases, this paper has a different perspective, and uses the
m-atlases as part of the more general Local Coordinate Space (LCS),
presented in \pagecref{part:lcs}.

\Crefrange{sec:M-charts}{sec:M-atlas} define m-charts and m-atlasses.
Maximal m-atlasses are essentially the same as manifolds.

\Cref{sec:M-ATLmorph} defines morphisms between m-atlasses, referred to
as m-atlas morphisms, in a fashion tailored to use in defining local
coordinate spaces. It constructs categories of M-atlases and functors,
and proves some basic results.

\Cref{sec:M-ATLcmorph} defines morphisms between m-atlasses in a
fashion similar to that conventionally used for $\Ck$ maps between
differentiable manifolds, in order to provide context. It constructs
categories of M-atlases and functors, and proves some basic results.

\section{M-charts}
{
  \showlabelsinline
  \label{sec:M-charts}
}
\begin{definition}[M-charts]
\label{def:M-chart}
Let $\seqname{E} \defeq (E, \catName{E})$ and
$\seqname{C} \defeq (C,\catName{C})$ be model spaces. An m-chart
$(U, V, \phi)$ of $\seqname{E}$ in the coordinate space $\seqname{C}$
consists of
\begin{enumerate}
\item A nonvoid model neighborhood $U \objin \catName{E}$, known as a
coordinate patch
\item A model neighborhood $V\objin \catName{C}$
\item A model homeomorphism $\phi \maps U \toiso V$, known as a
coordinate function
\end{enumerate}
\begin{remark}
I consider it clearer to explicate the range, rather than the
conventional usage of specifying only the domain and function or the
minimalist usage of specifying only the function.
\end{remark}
For evey $u \in U$, $(U, V, \phi)$ is an m-chart of $\seqname{E}$ in the
coordinate space $\catName{C}$ at $u$ and covers $u$.
$\seqname{E}$ is the total model space or total space for the chart and
$\seqname{C}$ is the coordinate model space or coordinate space for the
chart.

Let $E$ be a topological space and $\seqname{C} \defeq (C,\catName{C})$
be a model space. $(U, V, \phi)$ is an m-chart of $E$ in the coordinate
space $\catName{C}$ iff it is an m-chart of $\Triv{E}$ in the coordinate
space $\catName{C}$.

For evey $x \in U$, $(U, V, \phi)$ is an m-chart of $E$ in the
coordinate space $\catName{C}$ at $x$.
$E$ is the total space for the chart and $\seqname{C}$ is the coordinate
model space or coordinate space for the chart.
\end{definition}

\begin{lemma}[M-charts]
\label{lem:M-chart}
\leavevmode
\begin{enumerate}
\item
\label{M-chart:openC}
Let $\seqname{E} \defeq (E, \catName{E})$ and
$\seqname{C} \defeq (C,\catName{C})$ be model spaces, and $(U, V, \phi)$
be an m-chart of $\seqname{E}$ in the coordinate space $\seqname{C}$
such that every open subset of $V$ is a model neighborhood of
$\seqname{C}$. Then $(U, V, \phi)$ is an m-chart of $E$ in the
coordinate space $\seqname{C}$.

\begin{proof}
$U$, $V$ and $\phi$ satisfy the conditions of the definition:
\begin{description}
\item%
  [
    $U$ is a model neighborhood of \ensuremath{\Triv{E}}:
  ]
$U$ is a model neighborhood of $\seqname{E}$, hence open.

\item[$V$ is a model neighborhood of $\seqname{C}$:]
By hypothesis.

\item%
  [
    $\phi$ is a model homeomorphism from
    \ensuremath{\Mod \biggl ( U,\Triv{E} \biggr )} to
    $\Mod(V,\seqname{C})$:
  ]
\leavevmode
\begin{enumerate}
\item
Since $\phi$ is a model homeomorphism from $\Mod(U,\seqname{E})$ to
$\Mod(V,\seqname{C})$, it is a homeomorphism.

\item
If $V'$ is a model subneighborhood of $V$, then $\phi^{-1}[V']$ is open
and hence a model neighborhood of $\Triv{E}$.

\item
If $U'$ is a model neighborhood of $\Triv{E}$ then it is open, hence
$\phi[U]$ is an open subset of $V$ and thus by hypothesis a model
neighborhood of $\seqname{C}$.
\end{enumerate}
\end{description}
\end{proof}

\item
\label{M-chart:openE}
Let $(U, V, \phi)$ be an m-chart of $E$ in the coordinate space
$\seqname{C}$. $(U, V, \phi)$ is an m-chart of $\seqname{E}$ in the
coordinate space $\seqname{C}$ iff every open subset of $U$ is a model
neighborhood of $\seqname{E}$.

\begin{proof}
If $(U, V, \phi)$ is an m-chart of $E$ in the coordinate space
$\seqname{C}$ and also an m-chart of $\seqname{E}$ in the coordinate
space $\seqname{C}$ then let $U' \subseteq U$ be open. Since
$(U, V, \phi)$ is an m-chart of $E$ in the coordinate space
$\seqname{C}$, $\phi^{-1}$ is a model function from
$\Mod(V,\seqname{C})$ to $\Triv{E}$ and $\phi[U']$ is a model
neighborhood of $\seqname{C}$.  Since $(U, V, \phi)$ is an m-chart of
$\seqname{E}$ in the coordinate space $\seqname{C}$, $\phi$ is a model
function from $\Mod(U,\seqname{E})$ to $\seqname{C}$ and thus
$\phi^{-1} \bigl [ \phi[U'] \bigr ]$ is a model neighborhood of $\seqname{E}$.

If $(U, V, \phi)$ is an m-chart of $E$ in the coordinate space
$\seqname{C}$ and every open
\nobreaks
  {{
    $U'\subseteq U$
  }}
is a model neighborhood of $\seqname{E}$ then
\begin{enumerate}
\item Since $U \subseteq U$ is open, $U$ is a model neighborhood of
$\seqname{E}$ by hypothesis.

\item $V$ is a model neighborhood of $\seqname{C}$ by hypothesis.

\item If $V' \subseteq V$ is a model neighborhood of $\seqname{C}$ then
$\phi^{-1}[V']$ is open and hence a model neighborhod of $\seqname{E}$
by hypothesis.

$\phi^{-1}$ is a model function from $\Mod \biggl ( U,\Triv{E} \biggr )$ to
$\Mod(V,\seqname{C})$. if $U' \subseteq U$ is a model neighborhood of
$\seqname{E}$ then $U'$ is open and hence a model neighborhood of
$\Triv{E}$. Then $\phi[U']$ is a model neighborhood of $\seqname{C}$.
\end{enumerate}
\end{proof}
\end{enumerate}
\end{lemma}

\begin{corollary}[M-charts]
\label{cor:M-chart}
Let $\seqname{E} \defeq (E, \catName{E})$ and
$\seqname{C} \defeq (C,\catName{C})$ be model spaces, $\seqname{C}$ be
fine grained and $(U, V, \phi)$ be an m-chart of $\seqname{E}$ in the
coordinate space $\seqname{C}$.  Then $(U, V, \phi)$ is an m-chart of
$E$ in the coordinate space $\seqname{C}$.

\begin{proof}
Since $V$ is open and $\seqname{C}$ is fine grained, every open subset of
$V$ is a model neighborhood of $\seqname{C}$. The result follows by
\cref{M-chart:openC} of
{
  \showlabelsinline
  \fullcref{lem:M-chart}
  on \cpageref{M-chart:openC}.
}
\end{proof}
\end{corollary}

\begin{definition}[Subcharts]
\label{def:M-subchart}
Let $(U, V, \phi)$ be an m-chart of
$\seqname{E} \defeq (E, \catName{E})$ in the coordinate space
$(C,\catName{C})$ and $U' \subseteq U$ a nonvoid model neighborhood of
$\seqname{E}$. Then
$(U', V', \phi') \defeq (U', \phi[U'], \phi \restrictto_{U',V'})$ is a
subchart of $(U, V, \phi)$.

Let $(U, V, \phi)$ be an m-chart of $E$ in the coordinate
space $(C,\catName{C})$ and $U' \subseteq U$ an open set of
$E$. Then
$(U', V', \phi') \defeq (U', \phi[U'], \phi \restrictto_{U',V'})$
is a subchart of $(U, V, \phi)$.

$(U', V', \phi')$ is a subchart of $(U, V, \phi)$ at  $x$ iff
$(U', V', \phi')$ is a subchart of $(U, V, \phi)$ and $x \in U'$.
\end{definition}

\begin{lemma}[Subcharts]
\label{lem:M-subchart}
Let $\seqname{E} \defeq (E, \catName{E})$ and
$\seqname{C} \defeq (C,\catName{C})$ be model spaces.

Let $(U, V, \phi)$ be an m-chart of $\seqname{E}$ in the coordinate
space $\seqname{C}$ and $(U', V', \phi')$ a subchart of $(U, V, \phi)$.
Then $(U', V', \phi')$ is an m-chart of $\seqname{E}$ in the coordinate
space $\seqname{C}$.

\begin{proof}
$U'$, $V'$ and $\phi'$ satisfy the conditions of the definition of an m-chart:
\begin{enumerate}
\item $U'$ is a model neighborhood of $\seqname{E}$ by the definition of
subchart.
\item Since $U'$ is a model neighborhood and $\phi$ is a model
homeomorphism,
\nobreaks
  {{
    $V'=\phi[U']$
  }}
is a model neighborhood.
\item Since $\phi$ and $\phi^{-1}$ are model functions, so are
$\phi' = \phi \restrictto_{U',V'}$ and
$\phi'^{-1} = \phi^{-1} \restrictto_{V',U'}$. Thus $\phi'$ is a model
homeomorphism.
\end{enumerate}
\end{proof}

Let $(U, V, \phi)$ be an m-chart of $E$ in the coordinate space
$\seqname{C}$ and $(U', V', \phi')$ a subchart of $(U, V, \phi)$. Then
$(U', V', \phi')$ is an m-chart of $E$ in the coordinate space
$\seqname{C}$.

\begin{proof}
$U'$, $V'$ and $\phi'$ satisfy the conditions of the definition of an m-chart:
\begin{enumerate}
\item $U'$ is open by the definition of subchart.
\item Since $\phi$ is a homeomorphism and $U'$ is open, $V'=\phi[U']$ is open.
\item Since $\phi$ is a homeomorphism that maps open sets into model
neighborhoods, $\phi \restrictto_{U',V'}$ is a homeomorphism that maps
open sets into model neighborhoods.
\end{enumerate}
\end{proof}
\end{lemma}

\begin{definition}[M-compatibility]
\label{def:M-compat}
Let $\seqname{E} \defeq (E, \catName{E})$ and
$\seqname{C} \defeq (C, \catName{C})$ be model spaces and
$(U_i, V_i, \phi_i)$, $i=1,2$, be an m-chart of $\seqname{E}$ in the
coordinate space $\seqname{C}$. Then $(U_1, V_1, \phi_1)$ is
m-compatible with $(U_2, V_2, \phi_2)$ in the coordinate space
$\seqname{C}$ iff either

\begin{enumerate}
\item $U_1$ and $U_2$ are disjoint

\item The transition function
$
  \morphName{t}=\phi_2 \morphCompose \phi^{-1}_1
  \maps \phi_1[U_1 \cap U_2] \to \phi_2[U_1 \cap U_2]
$
is an isomorphism of $\catName{C}$.
\end{enumerate}

$(U_1, V_1, \phi_1)$ is also properly m-compatible with
\nobreaks
  {{
    $(U_2, V_2, \phi_2)$
  }}
in the space $\seqname{E}$ and the coordinate space $\seqname{C}$, and
properly m-compatible in the space $\seqname{E}$ and the coordinate
space $\seqname{C}$ with
\nobreaks
  {{
    $(U_2, V_2, \phi_2)$,
  }}
iff
for every model neighborhoods $U'_i \subseteq U_i$

\begin{enumerate}
\item
For every morphism $\morphName{f} \maps U'_1 \to U'_2$ of $\seqname{E}$,
the composition \\*
$
  \phi_2 \morphCompose \morphName{f} \morphCompose \phi^{-1}_1
  \maps \phi_1[U'_1] \to \phi_2[U'_2]
$
is a morphism of $\seqname{C}$.

\item
For every morphism $\morphName{g} \maps U'_2 \to U'_1$ of $\seqname{E}$,
the composition
$
  \phi_1 \morphCompose \morphName{g} \morphCompose \phi^{-1}_2
  \maps \phi_2[U'_2] \to \phi_1[U'_1]
$
is a morphism of $\seqname{C}$.
\end{enumerate}

\begin{remark}
This is a stronger condition than requiring that for every
\nobreaks{{$\morphName{f} \maps U_1 \to U_2$}} and
\nobreaks{{$\morphName{g} \maps U_2 \to U_1$}} the compositions be
morphisms of $\seqname{C}$.
\end{remark}

Let $E$ be a topological space, $\seqname{C} \defeq (C, \catName{C})$ be
a model space and $(U_i, V_i, \phi_i)$, $i=1,2$, be an m-chart of $E$ in
the coordinate space $\seqname{C}$. Then $(U_1, V_1, \phi_1)$ is
m-compatible with $(U_2, V_2, \phi_2)$ iff either

\begin{enumerate}
\item $U_1$ and $U_2$ are disjoint
\item The transition function
$
  \morphName{t}=\phi_2 \morphCompose \phi^{-1}_1
  \maps \phi_1[U_1 \cap U_2] \to \phi_2[U_1 \cap U_2]
$
is an isomorphism of $\catName{C}$
\end{enumerate}

$(U_1, V_1, \phi_1)$ is also properly m-compatible with
\nobreaks
  {{
    $(U_2, V_2, \phi_2)$
  }}
in the space $E$ and the coordinate space $\seqname{C}$, and properly
m-compatible in the space $E$ and the coordinate space $\seqname{C}$ with
\nobreaks
  {{
    $(U_2, V_2, \phi_2)$,
  }}
iff
for every open $U'_i \subseteq U_i$

\begin{enumerate}
\item
For every continuous $\morphName{f} \maps U'_1 \to U'_2$,
the composition
$
  \phi_2 \morphCompose \morphName{f} \morphCompose \phi^{-1}_1
  \maps \phi_1[U'_1] \to \phi_2[U'_2]
$
is a morphism of $\seqname{C}$.

\item
For every continuous $\morphName{g} \maps U'_2 \to U'_1$,
the composition
$
  \phi_1 \morphCompose \morphName{g} \morphCompose \phi^{-1}_2
  \maps \phi_2[U'_2] \to \phi_1[U'_1]
$
is a morphism of $\seqname{C}$.
\end{enumerate}

In both cases, the space and coordinate space may be omitted when they
are clear from context.
\end{definition}

\begin{lemma}[M-compatibility]
\label{lem:M-compat}
\mbox{}
\begin{enumerate}
\item
Let $(U_i,V_i,\phi_i)$, $i=1,2$. be a chart of $E$ in the coordinate
space $\seqname{C}$. Then $(U_1, V_1, \phi_1)$ is (properly)
m-compatible with
\nobreaks
  {{
    $(U_2, V_2, \phi_2)$
  }}
in the space $E$ and the coordinate space $\seqname{C}$ iff
$(U_1, V_1, \phi_1)$ is (properly) m-compatible with
\nobreaks
  {{
    $(U_2, V_2, \phi_2)$
  }}
in the space $\Triv{E}$ and the coordinate space $\seqname{C}$.

\begin{proof}
$U_1 \cap U_2$ and the transition function
\nobreaks
  {{
    $
      \morphName{t}=\phi_2 \morphCompose \phi^{-1}_1
      \maps \phi_1[U_1 \cap U_2] \to \phi_2[U_1 \cap U_2]
$
}}
is the same for the case of $E$ and the case of $\Triv{E}$.

Any
$U'_i \subseteq U_i$ is a model neighborhood of $\Triv{E}$ iff it is open.

Any $\morphName{f} \maps U'_1 \to U'_2$ is a morphism of $\Triv{E}$ iff
it is continuous.

Any $\morphName{g} \maps U'_2 \to U'_1$ is a morphism of $\Triv{E}$ iff
it is continuous.

The compositions
$
  \phi_2 \morphCompose \morphName{f} \morphCompose \phi^{-1}_1
  \maps \phi_1[U'_1] \to \phi_2[U'_2]
$
and
$
  \phi_1 \morphCompose \morphName{g} \morphCompose \phi^{-1}_2
  \maps \phi_2[U'_2] \to \phi_1[U'_1]
$
are the same for the case of $E$ and the case of $\Triv{E}$.
\end{proof}
\end{enumerate}
\end{lemma}

\begin{corollary}
\label{cor:M-compat}
Let $(U_i,V_i,\phi_i)$, $i=1,2$. be a chart of $E$ in the coordinate
space $\seqname{C}$ snd let $(U'_i,V'_i,\phi'_i)$ be a subchart. If
$(U_1,V_1,\phi_1)$ is (properly) m-compatible with $(U_2,V_2,\phi_2)$
in the space $E$ and the coordinate space $\seqname{C}$ then
$(U'_1,V'_1,\phi'_1)$ is (properly) m-compatible with
$(U'_2,V'_2,\phi'_2)$.

\begin{proof}
By \cref{lem:M-compat} above, $(U'_1,V'_1,\phi'_1)$ is (properly)
m-compatible with $(U'_2,V'_2,\phi'_2)$ in the space $E$ and the
coordinate space $\seqname{C}$ iff $(U'_1,V'_1,\phi'_1)$ is (properly)
m-compatible with $(U'_2,V'_2,\phi'_2)$ in the space $\Triv{E}$ and the
coordinate space $\seqname{C}$.

The result follows from \cref{lem:M-compat}.
\end{proof}
\end{corollary}

\begin{lemma}[Symmetry of m-compatibility]
\label{lem:M-compatsym}
Let $(U_i, V_i, \phi_i)$, $i=1,2$, be an m-chart of
$(E, \catName{E})$ in the coordinate space $(C,\catName{C})$. Then
\nobreaks
  {{
    $(U_1, V_1, \phi_1)$
  }}
is (properly) m-compatible with
\nobreaks
  {{
    $(U_2, V_2, \phi_2)$
  }}
in the space $\seqname{E}$ and the coordinate space $\seqname{C}$ iff
\nobreaks
  {{
    $(U_2, V_2, \phi_2)$
  }}
is (properly) m-compatible with
\nobreaks
  {{
    $(U_1, V_1, \phi_1)$.
  }}

\begin{proof}
It suffices to prove the implication in only one direction.
\begin{enumerate}
\item $U_1 \cap U_2 = U_2 \cap U_1$.
\item Since the transition function
\nobreaks
  {{
    $
      \morphName{t}=\phi_2 \morphCompose \phi^{-1}_1
      \maps \phi_1[U_1 \cap U_2] \to \phi_2[U_1 \cap U_2]
    $
  }}
is an isomorphism of $\catName{C}$, so is
\nobreaks
  {{
    $
      \morphName{t}^{-1} = \phi_1 \morphCompose \phi_2^{-1}
      \maps \phi_2[U_1 \cap U_2] \to \phi_1[U_1 \cap U_2]
    $.
  }}
\item The constraint for proper maximality is intrinsically symmetrical.
\end{enumerate}
\end{proof}

Let $(U_i, V_i, \phi_i)$, $i=1,2$, be an m-chart of $E$ in the
coordinate space $(C,\catName{C})$. Then $(U_1, V_1, \phi_1)$ is
(properly) m-compatible with $(U_2, V_2, \phi_2)$ in the space $E$ and
the coordinate space $\seqname{C}$ iff iff $(U_2, V_2, \phi_2)$ is
(properly) m-compatible with $(U_1, V_1, \phi_1)$.

\begin{proof}
The same proof applies as above.
\end{proof}
\end{lemma}

\begin{lemma}[M-compatibility of subcharts]
\label{lem:M-compatsub}
Let $(U_i, V_i, \phi_i)$, $i=1,2$, be an m-chart of $(E, \catName{E})$
in the coordinate space $(C,\catName{C})$,
\nobreaks
  {{
    $(U_1, V_1, \phi_1)$
  }}
be (properly) m-compatible with
\nobreaks
  {{
    $(U_2, V_2, \phi_2)$
  }}
in the space $\seqname{E}$ and the coordinate space $\seqname{C}$ and
\nobreaks
  {{
    $(U'_i, V'_i, \phi'_i)$
  }}
be a subchart. Then
\nobreaks
  {{
    $(U'_1, V'_1, \phi'_1)$
  }}
is (properly) m-compatible with
\nobreaks
  {{
    $(U'_2, V'_2, \phi'_2)$.
  }}

\begin{proof}
Since subcharts are charts, $U'_i$ and $V'_i$ are model neighborhoods.
If \nobreaks{{$U_1 \cap U_2 = \emptyset$}} then
\nobreaks{{$U'_1 \cap U'_2 = \emptyset$}} and thus
\nobreaks{{$(U'_1, V'_1, \phi'_1)$}} is m-compatible with
\nobreaks{{$(U'_2, V'_2, \phi'_2)$}}.  Otherwise, the transition
function
\nobreaks
  {{
    $
      t^1_2 \defeq
      \phi_2 \morphCompose \phi^{-1}_1
      \maps \phi_1[U_1 \cap U_2] \to \phi_2[U_1 \cap U_2]
    $
  }}
is an isomorphism of $\catName{C}$ and hence
\nobreaks
  {{
    $
      t^1_2 \maps \phi_1[U'_1 \cap U'_2] \toiso \phi_2[U'_1 \cap U'_2]
    $
 }}
is an isomorphism of $\catName{C}$.

If $(U_1, V_1, \phi_1)$ is properly m-compatible with
$(U_2, V_2, \phi_2)$,
$U''_i \subseteq U'_i$, $i=1,2$, is a model neighbothood and
$\morphName{f} \maps U''_1 \to U''_2$ is a morphism of $\seqname{E}$,
then each $U''_i \subseteq U_i$ and thus by \cref{def:M-compat} the
composition
$
  \phi_2 \morphCompose \morphName{f} \morphCompose \phi^{-1}_1
  \maps \phi_1[U''_1] \to \phi_2[U''_2]
$
is a morphism of $\seqname{C}$.
\end{proof}

Let $(U_i, V_i, \phi_i)$, $i=1,2$, be an m-chart of $E$ in the
coordinate space $(C,\catName{C})$, $(U'_i, V'_i, \phi'_i)$ be a
subchart and $(U_1, V_1, \phi_1)$ be (properly) m-compatible with
\nobreaks
  {{
    $(U_2, V_2, \phi_2)$
  }}
in the space $E$ and the coordinate space $\seqname{C}$. Then
\nobreaks
  {{
    $(U'_1, V'_1, \phi'_1)$
  }}
is (properly) m-compatible with
\nobreaks
  {{
    $(U'_2, V'_2, \phi'_2)$.
  }}

\begin{proof}
By \pagecref{lem:M-compat}, $(U_1, V_1, \phi_1)$ is (properly)
m-compatible with $(U_2, V_2, \phi_2)$ in the space $\Triv{E}$ and the
coordinate space $\seqname{C}$. Then by \cref{lem:M-compatsub} above,
$(U'_1,V'_1,\phi'_1)$ is (properly) m-compatible with
$(U'_2,V'_2,\phi'_2)$ in the space $\Triv{E}$ and the coordinate space
$\seqname{C}$.

The result follows from \cref{lem:M-compat}.
\end{proof}
\end{lemma}

\begin{corollary}[M-compatibility of subcharts]
\label{cor:M-compatsub}
Let $(U, V, \phi)$ be an m-chart of $(E, \catName{E})$ in
the coordinate space $(C,\catName{C})$ and $(U', V', \phi')$ a
subchart. Then $(U', V', \phi')$ is properly m-compatible with
\nobreaks
  {{
    $(U, V, \phi)$
  }}
in the space $\seqname{E}$ and the coordinate space $\seqname{C}$.

\begin{proof}
$(U, V, \phi)$ is m-compatible with itself and is a subchart of
itself,
\end{proof}

Let $(U, V, \phi)$ be an m-chart of $E$ in the coordinate
space $(C,\catName{C})$ and $(U', V', \phi')$ a subchart. Then
$(U', V', \phi')$ is m-compatible with
\nobreaks
  {{
    $(U,V,\phi)$
  }}
in the space $E$ and the coordinate space $\seqname{C}$.

\begin{proof}
$(U, V, \phi)$ is m-compatible with itself and is a subchart of
itself,
\end{proof}
\end{corollary}

\begin{definition}[Covering by m-charts]
Let $\seqname{A}$ be a set of charts of the topological space $E$ in
the coordinate space
$\seqname{C}=(C,\catName{C})$.
$\seqname{A}$ covers $E$ iff $\pi_1[\seqname{A}]$ covers $E$.

Let $\seqname{A}$ be a set of charts of the model space
$\seqname{E}=(E, \catName{E})$ in the coordinate space
$\seqname{C}=(C,\catName{C})$.  $\seqname{A}$ covers $\seqname{E}$ iff
$\pi_1[\seqname{A}]$ covers $E$.
\end{definition}

\section{M-atlases}
\label{sec:M-atlas}
A set of charts can be atlases for different coordinate model spaces
even if it is for the same total model space. In order to aggregate
atlases into categories, there must be a way to distinguish them.
Including the two\footnote{
The spaces are redundant, but convenient.}
spaces in the definitions of the categories serves the purpose.

\begin{definition}[M-atlases]
\label{def:M-atlas}
Let $\seqname{A}$ be a set of mutually (properly) m-compatible m-charts
of $\seqname{E}=(E, \catName{E})$ in the coordinate space
$\seqname{C}=(C,\catName{C})$. Then

\begin{enumerate}
\item
$\seqname{A}$ is an (a proper) m-atlas of $\seqname{E}$ in the
coordinate space $\seqname{C}$, abbreviated \\*
$\isAtl[(proper)]_\Ob(\seqname{A}, \seqname{E}, \seqname{C})$, iff
$\seqname{A}$ covers $E$.

\item
$\seqname{A}$ is a full (proper) atlas of $\seqname{E}$ in the
coordinate space $\seqname{C}$, abbreviated \\*
$\full{\isAtl[(proper)]_\Ob}(\seqname{A}, E, \seqname{C})$, iff
$\seqname{A}$ covers $E$ and $\pi_2[\seqname{A}]$ covers $C$.

\item
The triple $(\seqname{A}, \seqname{E}, \seqname{C})$ refers to
$\seqname{A}$ considered as an m-atlas of $\seqname{E}$ in the
coordinate space $\seqname{C}$.

\item
$\seqname{E}$ is the total model space or total space for the atlas

\item
$\seqname{C}$ is the coordinate model space or coordinate space for the
atlas.
\end{enumerate}

Let $E$ be a topological space, $\seqname{C}=(C,\catName{C})$ a model
space and $\seqname{A}$ be a set of mutually m-compatible m-charts of
$E$ in the coordinate space $\seqname{C}$. Then

\begin{enumerate}[resume]
\item
$\seqname{A}$ is an m-atlas of $E$ in the coordinate space
$\seqname{C}$, abbreviated as
$\isAtl_\Ob(\seqname{A}, E, \seqname{C})$, iff $\seqname{A}$ is an
m-atlas of $\Triv{E}$ in the coordinate space $\seqname{C}$.

\item
$\seqname{A}$ is a full atlas of $E$ in the coordinate space
$\seqname{C}$, abbreviated
$\full{\isAtl_\Ob}(\seqname{A}, E, \seqname{C})$, iff $\seqname{A}$ is a
full atlas of $\Triv{E}$ in the coordinate space $\seqname{C}$.

\item
The triple $(\seqname{A}, E, \seqname{C})$ refers to $\seqname{A}$
considered as an m-atlas of $E$ in the coordinate space $\seqname{C}$.

\item
$E$ is the total space for the atlas.

\item
$\seqname{C}$ is the coordinate model space or coordinate space for the
atlas.
\end{enumerate}

Let $\seqname{A}$ be an (a proper) m-atlas of $\seqname{E}$ in the
coordinate space $\seqname{C}$.  $\seqname{A}$ is fine grained iff
whenever $(U,V,\phi) \in \seqname{A}$, every subchart of
$(U,V,\phi)$ is also an atlas of $\seqname{A}$.

\begin{remark}
$\seqname{A}$ can be a fine graind (proper) atlas of $\seqname{E}$ in
the coordinate space $\seqname{C}$ even if $\seqname{C}$ is not fine
grained.
\end{remark}

Let $\seqname{A}$ be an (a proper) m-atlas of $E$ in the coordinate
space $\seqname{C}$.  $\seqname{A}$ is fine grained iff whenever
$(U,V,\phi) \in \seqname{A}$, every subchart of $(U,V,\phi)$ is also an
atlas of $\seqname{A}$.

By abuse of language we write $U \in \seqname{A}$ for
$U \in \pi_1[\seqname{A}]$.
\end{definition}

\begin{lemma}[M-atlases]
\label{lem:M-atlas}
Let $\seqname{A}$ be an m-atlas of $\seqname{E} \defeq (E, \catName{E})$
in the coordinate space $\seqname{C} \defeq (C,\catName{C})$. If
$\seqname{C}$ is fine grained then $\seqname{A}$ is an m-atlas of $E$ in
the coordinate space $\seqname{C}$.

\begin{proof}
Let $(U,V,\phi) \in \seqname{A}$. Then $(U,V,\phi)$ is an m-chart of $E$
in the coordinate space $\seqname{C}$ by \pagecref{cor:M-chart}.
\end{proof}
\end{lemma}

\begin{definition}[M-compatibility with m-atlases]
\label{def:M-CompatATL}
An m-chart $(U, V, \phi)$ of $(E, \catName{E})$ is (properly)
m-compatible with an m-atlas $\seqname{A}$ in the space $\seqname{E}$
and the coordinate space $\seqname{C}$ iff it is (properly) m-compatible
in the space $\seqname{E}$ and the coordinate space $\seqname{C}$
with every chart in $\seqname{A}$.

An m-chart $(U, V, \phi)$ of $E$ is (properly) m-compatible with an
m-atlas $\seqname{A}$ in the space $\seqname{E}$ and the coordinate
space $\seqname{C}$ iff it is (properly) m-compatible in the space
$\seqname{E}$ and the coordinate space $\seqname{C}$ with every chart in
$\seqname{A}$.
\end{definition}

\begin{lemma}[M-compatibility with m-atlases]
\label{lem:M-CompatATL}
Let $\seqname{A}$ be an m-atlas of $\seqname{E} \defeq (E, \catName{E})$
in the coordinate space $\seqname{C} \defeq (C,\catName{C})$. If
$\seqname{C}$ is fine grained then $\seqname{A}$ is an m-atlas of $E$ in
the coordinate space $\seqname{C}$.

\begin{proof}
\end{proof}
\end{lemma}

\begin{lemma}[M-compatibility of subcharts with atlases]
\label{lem:M-CompatsubATL}
Let $\seqname{A}$ be an (a proper) m-atlas of $(E, \catName{E})$ in the
coordinate space $(C,\catName{C})$ and $(U_1, V_1, \phi_1)$ an m-chart
in $\seqname{A}$. Then any subchart of $(U_1, V_1, \phi_1)$ is
properly m-compatible with $\seqname{A}$ in the space $\seqname{E}$
and the coordinate space $\seqname{C}$.

\begin{proof}
Let $(U'^1, V'^1, \phi'^1)$ be a subchart of $(U^1, V^1, \phi^1)$ and
$(U_2, V_2, \phi_2)$ be another chart in $\seqname{A}$.  Then
$(U_1, V_1, \phi_1)$ is properly compatible with $(U_2, V_2, \phi_2)$
by \pagecref{def:M-atlas}, $(U_2, V_2, \phi_2)$ is properly
m-compatible with itself and $(U_2, V_2, \phi_2)$ is a subchart of
itself. $(U'^1, V'^1, \phi'^1)$ is properly m-compatible with
$(U_2, V_2, \phi_2)$ by \pagecref{lem:M-compatsub}. Since
$(U_2, V_2, \phi_2)$ is arbitrary, $(U'^1, V'^1, \phi'^1)$ is properly
m-compatible with $\seqname{A}$ by \pagecref{def:M-CompatATL}.
\end{proof}

Let $\seqname{A}$ be an m-atlas of $E$ in the coordinate space
$(C,\catName{C})$ and $(U_1, V_1, \phi_1)$ an m-chart in $\seqname{A}$.
Then any subchart of $(U_1, V_1, \phi_1)$ is properly m-compatible with
$\seqname{A}$ in the space $E$ and the coordinate space $\seqname{C}$.

\begin{proof}
The same proof applies in as above.
\end{proof}
\end{lemma}

\begin{lemma}[Extensions of m-atlases]
\label{lem:M-ATLextensions}
Let $\seqname{E}=(E,\catName{E})$ and $\seqname{C}=(C,\catName{C})$ be
model spaces, $\seqname{A}$ an m-atlas of $\seqname{E}$ in the
coordinate space $\seqname{C}$, and $(U_i,V_i,\phi_i)$, $i=1,2$, an
m-chart of $\seqname{E}$ in the coordinate space $\seqname{C}$
(properly) m-compatible with $\seqname{A}$ in the space $\seqname{E}$
and the coordinate space $\seqname{C}$.  Then $(U_1, V_1, \phi_1)$ is
(properly) m-compatible with $(U_2, V_2, \phi_2)$.

\begin{proof}
If $U_1 \cap U_2 = \emptyset$ then $(U_1, V_1, \phi_1)$ is m-compatible
with $(U_2 ,V_2, \phi_2)$. Otherwise, $\phi_1$ is a model homeomorphism,
$U_1 \cap U_2$ is a model neighborhood of $\seqname{E}$,
$\phi_1[U_1 \cap U_2]$ is a model neighborhood of $\seqname{C}$,
$\phi_2[U_1 \cap U_2]$ is a model neighborhood of $\seqname{C}$ and
$
  \phi_2 \morphCompose \phi^{-1}_1
  \maps \phi_1[U_1 \cap U_2] \toiso \phi_2[U_1 \cap U_2]
$
is a homeomorphism.  It remains to show that
$\phi_2 \morphCompose \phi^{-1}_1 \restrictto_{\phi_1[U_1 \cap U_2]}$ is a
model homeomorphism, Let
$(U_\alpha,V_\alpha,\phi_\alpha)$, $\alpha \prec \Alpha$, be charts in
$\seqname{A}$ such that
$U_1 \cap U_2 \subseteq \union[\alpha \prec \Alpha]{U_\alpha}$
and
\nobreaks
  {{
    $
      \uquant
        {
          \alpha \prec \Alpha
        }
        {
          U_1 \cap U_2 \cap U_\alpha \neq \emptyset
        }
    $.
  }}
Since the charts are m-compatible with
\nobreaks
  {{
    $(U_\alpha, V_\alpha, \phi_\alpha)$,
  }}
the compositions \\*
\nobreaks
  {{
    $
      \phi_2 \morphCompose \phi^{-1}_\alpha
      \maps \phi_\alpha[U_1 \cap U_2 \cap U_\alpha]
      \to \phi_2[U_1 \cap U_2 \cap U_\alpha]
    $
  }}
and \\*
\nobreaks
  {{
    $
      \phi_\alpha \morphCompose \phi^{-1}_1
      \maps \phi_1[U_1 \cap U_2 \cap U_\alpha]
      \to \phi_\alpha[U_1 \cap U_2 \cap U_\alpha]
    $
  }}
are isomorphisms of $\seqname{C}$ and thus the composition
\nobreaks
  {{
    $
      \phi_2 \morphCompose \phi^{-1}_1
      \maps \phi_1[U_1 \cap U_2 \cap U_\alpha]
      \to \phi_2[U_1 \cap U_2 \cap U_\alpha]
      =
    $
  }} \\*
\nobreaks
  {{
    $
      \phi_2 \morphCompose \phi^{-1}_\alpha \morphCompose
      \phi_\alpha \morphCompose \phi^{-1}_1
      \maps \phi_1[U_1 \cap U_2 \cap U_\alpha]
      \to \phi_2[U_1 \cap U_2 \cap U_\alpha]
    $
  }}
is an isomorphisms of $\seqname{C}$. Then by
{
  \showlabelsinline
  \cref{mod:sheaf}
},
the restricted sheaf property, of \pagecref{def:model},
$\phi_2 \morphCompose \phi^{-1}_1$
is an isomorphism of $\seqname{C}$.

If each $(U_i,V_i,\phi_i)$ is properly m-compatible with $\seqname{A}$
then let
\nobreaks
  {{
    $\morphName{f} \maps U'_1 \to U'_2$
  }}
and
\nobreaks
  {{
    $\morphName{g} \maps U'_2 \to U'_1$
  }}
be morphisms of $\seqname{E}$ and $U'_i \subseteq U_i$, $i=1,2$.
Let $u$ be an arbitary point in $U_1$ and
\nobreaks
  {{
    $(U_u, V_u, \phi_u) \in \seqname{A}$
  }}
be an arbitray chart at $u$. Since
\nobreaks
  {{
    $(U_1,V_1,\phi_1)$
  }}
is properly m-compatible with $\seqname{A}$,
it is properly m-compatible with
\nobreaks
  {{
    $(U_u, V_u, \phi_u) \in \seqname{A}$,
  }}
thus the compositions
\nobreaks
  {{
    $
      \phi_2 \morphCompose \morphName{f} \morphCompose \phi^{-1}_u
      \maps \phi^{-1}_u[U_u] \to \phi^{-1}_2[U'_2]
    $
  }}
and
\nobreaks
  {{
    $
      \phi_u \morphCompose \phi^{-1}_1
      \maps \phi^{-1}_1[U_u] \to \phi^{-1}_u[U_u]
    $
  }}
are morphisms of $\seqname{C}$, their composition
\nobreaks
  {{
    $
      \phi_2 \morphCompose \morphName{f} \morphCompose \phi^{-1}_u \morphCompose
      \phi_u \morphCompose \phi^{-1}_1
      \maps \phi^{-1}_1[U_u] \to \phi^{-1}_2[U'_2] =
    $
  }}
\nobreaks
  {{
    $
      \phi_2 \morphCompose \morphName{f} \morphCompose \phi^{-1}_1
      \maps \phi^{-1}_1[U_u] \to \phi^{-1}_2[U'_2]
    $
  }}
is a morphism of $\seqname{C}$ and by
{
  \showlabelsinline
  \cref{mod:sheaf}
},
the restricted sheaf property, of \pagecref{def:model}, \\*
\nobreaks
  {{
    $
      \phi_2 \morphCompose \morphName{f} \morphCompose \phi^{-1}_1
      \maps \phi^{-1}_1[U'_1] \to \phi^{-1}_2[U'_2]
    $
  }}
is a morphism of $\seqname{C}$.

The same proof applies in the opposite direction.
\end{proof}

Let $E$ be a topological space, $\seqname{C}=(C,\catName{C})$ a model
space, $\seqname{A}$ an (a proper) m-atlas of $E$ in the coordinate
space $\seqname{C}$ and $(U_i, V_i, \phi_i)$, $i=1,2$, a chart
(properly) m-compatible with $\seqname{A}$ in the space $E$ and the
coordinate space $\seqname{C}$.  Then $(U_1, V_1, \phi_1)$ is (properly)
m-compatible with $(U_2, V_2, \phi_2)$.

\begin{proof}
The same proof applies as above.
\end{proof}
\end{lemma}

\begin{definition}[Maximal m-atlases]
\label{def:M-ATLmax}
Let $\seqname{E}$ and $\seqname{C}$ be model spaces and $\seqname{A}$ an
m-atlas of $\seqname{E}$ in the coordinate space $\seqname{C}$.
$\seqname{A}$ is a (proper) maximal m-atlas of $\seqname{E}$ in the
coordinate space $\seqname{C}$, abbreviated
\nobreaks
  {{
    $
      \maximal{\isAtl[(proper)]_\Ob}
      (\seqname{A}, \seqname{E}, \seqname{C})
    $,
  }}
iff $\seqname{A}$ cannot be extended by adding an additional (properly)
m-compatible chart.

Let $E$ be a topological space, $\seqname{C}=(C,\catName{C})$ be a model
space and $\seqname{A}$ be an m-atlas of $E$ in the coordinate space
$\seqname{C}$. $\seqname{A}$ is a (properly) maximal m-atlas of $E$ in
the coordinate space $\seqname{C}$, abbreviated
$\maximal{\isAtl[(proper)]_\Ob}(\seqname{A}, E, \seqname{C})$, iff
$\seqname{A}$ is an m-atlas that cannot be extended by adding an
additional (properly) m-compatible chart.

\begin{remark}
There is no sheaf condition\footnote{
  However, note \cref{mod:sheaf} (restricted sheaf condition) of
  {
    \showlabelsinline
    \pagecref{def:model}.
  }
};
the union of a set of coordinate patches in the maximal atlas whose
coordinate functions match on the intersections need not be a coordinate
patch in the atlas.
\end{remark}

\begin{equation}
  \maxfull{\isAtl[(proper)]_\Ob}(\seqname{A}, E, \seqname{C}) \defeq
    \full{\isAtl[(proper)]_\Ob}(\seqname{A}, E, \seqname{C}) \land
    \maximal{\isAtl[(proper)]_\Ob}(\seqname{A}, E, \seqname{C})
\end{equation}
\begin{equation}
  \maxfull[S-]{\isAtl[(proper)]_\Ob}(\seqname{A}, E, \seqname{C}) \defeq
    \full{\isAtl[(proper)]_\Ob}(\seqname{A}, E, \seqname{C}) \land
    \maximal[S-]{\isAtl[(proper)]_\Ob}(\seqname{A}, E, \seqname{C})
\end{equation}
\end{definition}

\begin{theorem}[Existence and uniqueness of maximal m-atlases]
\label{the:M-ATLmax}
Let $\seqname{A}$ be an (a proper) m-atlas of
$\seqname{E} \defeq (E, \catName{E})$ in the coordinate space
$\seqname{C} \defeq (C,\catName{C})$. Then there exists a unique
(proper) maximal m-atlas
\nobreaks
  {{
    $\maximal{\Atlas[(proper)]}(\seqname{A}, \seqname{E}, \seqname{C})$
  }}
of $\seqname{E}$ in the coordinate space $\seqname{C}$ m-compatible with
$\seqname{A}$.

\begin{proof}
Let $\seqname{P}$ be the set of all m-atlases of $\seqname{E}$ in the
coordinate space $\seqname{C}$ containing $\seqname{A}$ and (properly)
m-compatible in the coordinate space $\seqname{C}$ with all of the
m-charts in $\seqname{A}$. Let $\maximal{\seqname{P}}$ be a maximal
chain of $\seqname{A}$. Then
\nobreaks
  {{
    $
      \maximal{\Atlas[(proper)]}(\seqname{A}, \seqname{E}, \seqname{C})
      \defeq \union{\maximal{\seqname{P}}}
    $
  }}
is a (proper) maximal m-atlas of $\seqname{E}$ in the coordinate space
$\seqname{C}$ m-compatible with $\seqname{A}$. Uniqueness follows from
\pagecref{lem:M-ATLextensions}.
\end{proof}
\end{theorem}

\begin{corollary}[Existence and uniqueness of maximal m-atlases]
\label{cor:M-ATLmax}
Let $\seqname{A}$ be an (a proper) m-atlas of $E$
in the coordinate space $\seqname{C} \defeq (C,\catName{C})$. Then there
exists a unique (proper) maximal m-atlas
$\maximal{\Atlas[(proper)]}(\seqname{A}, E, \seqname{C})$ of $E$
in the coordinate space $\seqname{C}$ m-compatible with $\seqname{A}$.

\begin{proof}
$
  \maximal{\Atlas[(proper)]}(\seqname{A}, E, \seqname{C}) \defeq
  \maximal{\Atlas[(proper)]}(\seqname{A}, \Triv{E}, \seqname{C})
$.
\end{proof}
\end{corollary}

\begin{definition}[Sets of m-atlases]
\label{def:M-ATLsets}

Let $\seqname{E}$ and $\seqname{C}$ be model spaces.

\begin{multline}
\maxfull[*]{\Atl[(proper)]_\Ob(\seqname{E}, \seqname{C})} \defeq
\set
{{
   (\seqname{A}, \seqname{E}, \seqname{C})
}}%
[
  {
    \maxfull[*]
    {
      \isAtl[(proper)]_\Ob
      (\seqname{A}, \seqname{E}, \seqname{C})
    }
  }
]*
\end{multline}

Let $\seqname{E}$ and $\seqname{C}$ be sets of model spaces.

\begin{multline}
\maxfull[*]{\Atl[(proper)]_\Ob(\seqname{E}, \seqname{C})} \defeq
\set
{{
   \bigl ( \seqname{A}, (E,\catName{E}), (C,\catName{C}) \bigr )
}}%
[
  {(E,\catName{E}) \in \seqname{E}},
  {(C,\catName{C}) \in \seqname{C}},
  {
    \maxfull[*]
    {
      \isAtl[(proper)]_\Ob
      \bigl ( \seqname{A}, (E,\catName{E}), (C,\catName{C}) \bigr )
    }
  }
]*
\end{multline}

Let $\seqname{E}$ be a set of topological spaces and $\seqname{C}$ a
set of model spaces. Then
\begin{multline}
\maxfull[*]{\Atl[(proper)]_\Ob (\seqname{E}, \seqname{C})} \defeq
\set
{{
   \bigl ( \seqname{A}, E, (C,\catName{C}) \bigr )
}}%
[
  {E \in \seqname{E}},
  {(C,\catName{C}) \in \seqname{C}},
  {
    \maxfull[*]
    {
      \isAtl[(proper)]_\Ob
      \bigl ( \seqname{A}, E, (C,\catName{C}) \bigr )
    }
  }
]*
\end{multline}
\end{definition}

\begin{lemma}[Sets of m-atlases]
\label{lem:M-ATLsets}

Let $\seqname{E}$ and $\seqname{C}$ be sets of model spaces.

\begin{equation}
  \maxfull{\Atl[(proper)]_\Ob(\seqname{E}, \seqname{C})} \subseteq
  \maxfull[*x]{\Atl[(proper)]_\Ob(\seqname{E}, \seqname{C})}
\end{equation}

\begin{equation}
  \maxfull[**]{\Atl[(proper)]_\Ob(\seqname{E}, \seqname{C})} \subseteq
  \Atl[(proper)]_\Ob(\seqname{E}, \seqname{C})
\end{equation}

\begin{proof}
The result follows from
{
  \showlabelsinline
  \pagecref{def:M-atlas},
}
\pagecref{def:M-ATLmax}
and \cref{def:M-ATLsets} above.
\end{proof}

\begin{equation}
  \maxfull{\Atl[(proper)]_\Ob(\seqname{E}, \seqname{C})} \subseteq
  \full{\Atl[(proper)]_\Ob(\seqname{E}, \seqname{C})}
\end{equation}

\begin{proof}
The result follows from \pagecref{def:M-ATLmax} and \cref{def:M-ATLsets}
above.
\end{proof}
\end{lemma}

\begin{corollary}[Sets of m-atlases]
\label{cor:M-ATLsets}

Let $\seqname{E}$ be a set of topological spaces and $\seqname{C}$ a
set of model spaces. Then

\begin{equation}
  \maxfull{\Atl[(proper)]_\Ob(\seqname{E}, \seqname{C})} \subseteq
  \maxfull[*x]{\Atl[(proper)]_\Ob(\seqname{E}, \seqname{C})}
\end{equation}

\begin{equation}
  \maxfull[**]{\Atl[(proper)]_\Ob(\seqname{E}, \seqname{C})} \subseteq
  \Atl[(proper)]_\Ob(\seqname{E}, \seqname{C})
\end{equation}

\begin{equation}
  \maximal[max-full]{\Atl[(proper)]_\Ob(\seqname{E}, \seqname{C})}
  \subseteq
  \maximal[max]{\Atl[(proper)]_\Ob(\seqname{E}, \seqname{C})}
\end{equation}

\begin{equation}
  \maxfull{\Atl[(proper)]_\Ob(\seqname{E}, \seqname{C})} \subseteq
  \full{\Atl[(proper)]_\Ob(\seqname{E}, \seqname{C})}
\end{equation}

\begin{proof}
The result follows from \pagecref{def:M-atlas},  \pagecref{def:M-ATLmax}
and \cref{def:M-ATLsets} above.
\end{proof}
\end{corollary}

\section{M-atlas morphisms and functors}
\label{sec:M-ATLmorph}
This section introduces a taxonomy for morphisms between m-atlases,
defines categories of m-atlases, defines functors amomg them and proves
some basic reasults.

\subsection{M-atlas morphisms}
\label{sub:M-ATLmorph}
This subsection introduces the notions of m-atlas morphisms, strict
m-atlas morphisms and semi-strict m-atlas morphisms.

\subsubsection{Definitions of m-atlas morphisms}
\begin{definition}[M-atlas morphisms for model spaces]
\label{def:M-ATLmorph}
Let $\catName{E}^i$, $\catName{C}^i$, $i=1,2$, be model categories,
$\seqname{E}^i \objin \catName{E}^i$,
$\seqname{C}^i \objin \catName{C}^i$, $\seqname{A}^i$ be an m-atlas of
$\seqname{E}^i$ in the coordinate space $\seqname{C}^i$ and
$
  \morphSeqName{f} \defeq
  (
    \morphName{f}_0 \maps \seqname{E}^1 \to \seqname{E}^2,
    \morphName{f}_1 \maps \seqname{C}^1 \to \seqname{C}^2
  )
$
a pair of model functions.

$\morphSeqName{f}$ is an
$\seqname{E}^1$-$\seqname{E}^2$ m-atlas morphism of $\seqname{A}^1$ to
$\seqname{A}^2$ in the coordinate spaces $\seqname{C}^1$,
$\seqname{C}^2$, abbreviated as
$
  \isAtl_\Ar
    (
      \seqname{A}^1,
      \seqname{E}^1,
      \seqname{C}^1,
      \seqname{A}^2,
      \seqname{E}^2,
      \seqname{C}^2,
      \morphName{f}_0,
      \morphName{f}_1
    )
$
and an $\catName{E}^1$-$\catName{E}^2$ m-atlas morphism of
$\seqname{A}^1$ to $\seqname{A}^2$ in the coordinate model categories
$\catName{C}^1$, $\catName{C}^2$, abbreviated as \\*
$
  \isAtl_\Ar
  (
    \seqname{A}^1,
    \catName{E}^1,
    \catName{C}^1,
    \seqname{A}^2,
    \catName{E}^2,
    \catName{C}^2,
    \morphName{f}_0,
    \morphName{f}_1
  )
$,
iff for any charts
$(U^i, V^i, \phi^i \maps U^i \toiso V^i) \in \seqname{A}^i$, $i=1,2$,
with $I \defeq U^1 \cap \morphName{f}^{-1}_0[U^2] \neq \emptyset$ and any
$u^1 \in I$, there exists a subchart
$
  (U'^1, V'^1, \phi^1 \maps U'^1 \toiso V'^1)
  \in \seqname{A}^1
$
at $u^1$, a model neighborhood $U'^2 \subseteq U^2$ of
$u^2 \defeq \morphName{f}_0(u^1)$ and a chart
$
  (U'^2 ,\hat{V}'^2, \phi'^2 \maps U'^2 \toiso \hat{V}'^2)
  \in \seqname{A}^2
$
at $u^2$ such that
$\morphName{f}_0[U'^1] \subseteq U'^2$ and diagram \eqref{eq:AtlM1}
below is commutative, as shown in \pagecref{fig:AtlM1}.

\begin{equation}
\begin{array}{lll}
\label{eq:AtlM1}
D \defeq & \Biggl \{
\\
  &&
  \set{u^1} \to^{\morphName{i}} U'^1
  \to/{>}->>/^{\phi^1}_\iso V'^1
  \to^{\morphName{f}_1} \hat{V}'^2,
\\
  &&
  U'^1 \to^{\morphName{f}_0}  U'^2 \to/{>}->>/^{\phi'^2}_\iso \hat{V}'^2
\\
& \Biggr \}
\end{array}
\end{equation}

\begin{figure}[ht]
\[ \bfig
\node e1(500,0)[u^1]
\node e2(1500,0)[u^2]
\node u1(-300,-500)[U^1]
\node up1(500,-500)[U'^1]
\node up2(1500,-500)[U'^2]
\node u2(2000,-500)[U^2]
\node v1(-300,-1000)[V^1]
\node vp1(500,-1000)[V'^1]
\node vhp2(1500,-1000)[\hat{V}'^2]
\node v2(2000,-1000)[V^2]
\arrow |a|[e1`e2;\morphName{f}_0]
\arrow |r|/^{ (}->/[e1`up1;\morphName{i}]
\arrow |r|/^{ (}->/[e2`up2;\morphName{i}]
\arrow |r|/>->/[u1`v1;\phi^1]
\arrow |l|//[u1`v1;\iso]
\arrow |a|/_{ (}->/[up1`u1;\morphName{i}]
\arrow |a|[up1`up2;\morphName{f}_0]
\arrow |r|[up1`vp1;{\phi^1}]
\arrow |l|//[up1`vp1;\iso]
\arrow |a|/^{ (}->/[up2`u2;i]
\arrow |r|[up2`vhp2;\phi'^2]
\arrow |l|//[up2`vhp2;\iso]
\arrow |r|[u2`v2;\phi^2]
\arrow |l|//[u2`v2;\iso]
\arrow |a|/_{ (}->/[vp1`v1;\morphName{i}]
\arrow |a|[vp1`vhp2;\morphName{f}_1]
\efig \]
\caption{Completed m-atlas morphism}
\label{fig:AtlM1}
\end{figure}

The triple
$
  \Bigl (
    \morphSeqName{f},
    (\seqname{A}^1, \seqname{E}^1, \seqname{C}^1),
    (\seqname{A}^2, \seqname{E}^2, \seqname{C}^2)
  \Bigr )
$
will refer to $\morphSeqName{f}$ considered as an
$\seqname{E}^1$-$\seqname{E}^2$ m-atlas morphism of
$\seqname{A}^1$ to $\seqname{A}^2$ in the coordinate spaces
$\seqname{C}^1$, $\seqname{C}^2$.

$\morphSeqName{f}$ is also a constrained
$\seqname{E}^1$-$\seqname{E}^2$ m-atlas morphism of $\seqname{A}^1$ to
$\seqname{A}^2$ in the coordinate spaces $\seqname{C}^1$,
$\seqname{C}^2$, abbreviated as
$
  \isAtl[const]_\Ar
    (
      \seqname{A}^1,
      \seqname{E}^1,
      \seqname{C}^1,
      \seqname{A}^2,
      \seqname{E}^2,
      \seqname{C}^2,
      \morphName{f}_0,
      \morphName{f}_1
    )
$
and a constrained $\catName{E}^1$-$\catName{E}^2$ m-atlas morphism of
$\seqname{A}^1$ to $\seqname{A}^2$ in the coordinate model categories
$\catName{C}^1$, $\catName{C}^2$, abbreviated as
$
  \isAtl[const]_\Ar
  (
    \seqname{A}^1,
    \catName{E}^1,
    \catName{C}^1,
    \seqname{A}^2,
    \catName{E}^2,
    \catName{C}^2,
    \morphName{f}_0,
    \morphName{f}_1
  )
$,
iff $\morphName{f}_0$ and $\morphName{f}_1$ are constrained, i.e.,
$\morphName{f}_0[\seqname{E}^1]$ is contained in a model neighborhood of
$\seqname{E}^2$ and $\morphName{f}_1[\seqname{C}^1]$ is contained in a
model neighborhood of $\seqname{C}^2$.

\begin{remark}
If $\morphSeqName{f}$ is constrained then $\Top(\seqname{E}^1)$ is a
model neighborhood of $\seqname{E}^1$ and $\Top(\seqname{C}^1)$ is a
model neighborhood of $\seqname{C}^1$.
\end{remark}
\end{definition}

\begin{definition}[M-atlas morphisms for topological spaces]
\label{def:M-ATLmorphtop}
Let $\catName{C}^i$, $i=1,2$, be a model category,
$\seqname{C}^i \objin \catName{C}^i$, $E^i$ be a topological space and
$\seqname{A}^i$ be an m-atlas of $E^i$ in the coordinate space
$\seqname{C}^i$.

A pair of functions
$
  \morphSeqName{f} \defeq
  (
    \morphName{f}_0 \maps E^1 \to E^2,
    \morphName{f}_1 \maps \seqname{C}^1 \to \seqname{C}^2
  )
$
is an $E^1$-$E^2$ m-atlas morphism of
$\seqname{A}^1$ to $\seqname{A}^2$ in the coordinate spaces
$\seqname{C}^i$, abbreviated as \\*
$
  \isAtl_\Ar
    (
      \seqname{A}^1,
      E^1,
      \seqname{C}^1,
      \seqname{A}^2,
      E^2,
      \seqname{C}^2,
      \morphName{f}_0,
      \morphName{f}_1
    )
$,
and an $E^1$-$E^2$ m-atlas morphism of $\seqname{A}^1$ to $\seqname{A}^2$
in the coordinate model categories $\catName{C}^i$, abbreviated as \\*
$
  \isAtl_\Ar
    (
      \seqname{A}^1,
      E^1,
      \catName{C}^1,
      \seqname{A}^2,
      E^2,
      \catName{C}^2,
      \morphName{f}_0,
      \morphName{f}_1
    )
$,
iff it is an $\Triv{E^1}$-$\Triv{E^2}$ m-atlas morphism of
$\seqname{A}^1$ to $\seqname{A}^2$ in the coordinate model categories
$\catName{C}^1$, $\catName{C}^2$.

The triple
$
  \bigl (
    \morphSeqName{f},
    (\seqname{A}^1, E^1, \seqname{C}^1),
    (\seqname{A}^2, E^2, \seqname{C}^2)
  \bigr )
$
will refer to $\morphSeqName{f}$ considered as an $E^1$-$E^2$ m-atlas
morphism of $\seqname{A}^1$ to $\seqname{A}^2$ in the coordinate
spaces $\seqname{C}^1$, $\seqname{C}^2$.

$\morphSeqName{f}$ is also a constrained $E^1$-$E^2$ m-atlas morphism of
$\seqname{A}^1$ to $\seqname{A}^2$ in the coordinate spaces
$\seqname{C}^1$, $\seqname{C}^2$, abbreviated as
$
  \isAtl[const]_\Ar
    (
      \seqname{A}^1,
      E^1,
      \seqname{C}^1,
      \seqname{A}^2,
      E^2,
      \seqname{C}^2,
      \morphName{f}_0,
      \morphName{f}_1
    )
$,
and a constrained $E^1$-$E^2$ m-atlas morphism of $\seqname{A}^1$ to $\seqname{A}^2$
in the coordinate model categories $\catName{C}^1$, $\catName{C}^2$, abbreviated as
$
  \isAtl[const]_\Ar
    (
      \seqname{A}^1,
      E^1,
      \catName{C}^1,
      \seqname{A}^2,
      E^2,
      \catName{C}^2,
      \morphName{f}_0,
      \morphName{f}_1
    )
$,
iff $\morphName{f}_1$ is constrained, i.e.,
$\morphName{f}_1[\seqname{C}^1]$ is contained in a model neighborhood of
$\seqname{C}^2$.
\end{definition}

\begin{definition}[M-atlas inclusion and identity morphisms]
\label{def:M-ATLident}
Let
\begin{enumerate}
\item
$\catName{E}^i$, $\catName{C}^i$, $i=1,2$, be model categories
\item
$\seqname{E}^i \objin \catName{E}^i$
\item
$\seqname{E}^1 \submod \seqname{E}^2$
\item
$\seqname{C}^i \objin \catName{C}^i$
\item
$\seqname{C}^1 \submod \seqname{C}^2$
\item
$\seqname{A}^i$ be an m-atlas of $\seqname{E}^i$ in the coordinate
space $\seqname{C}^i$
\item
$\seqname{A}^1 \subseteq \seqname{A}^2$
\end{enumerate}

The $\seqname{E}^1$-$\seqname{E}^2$ m-atlas inclusion morphism of
$\seqname{A}^1$ to $\seqname{A}^2$ in the coordinate spaces
$\seqname{C}^1$, $\seqname{C}^2$ is
$
 \ID%
   [
     {(\seqname{E}^1, \seqname{C}^1)},
     {(\seqname{E}^2, \seqname{C}^2)}
   ]
$
considered as an $\seqname{E}^1$-$\seqname{E}^2$ m-atlas morphism of
$\seqname{A}^1$ to $\seqname{A}^2$ in the coordinate spaces
$\seqname{C}^1$, $\seqname{C}^2$.

The m-atlas inclusion morphism of
$(\seqname{A}^1, \seqname{E}^1, \seqname{C}^1)$ to
$(\seqname{A}^2, \seqname{E}^2, \seqname{C}^2)$ is \\*
$
  \Id%
    [
        {(\seqname{A}^1, \seqname{E}^1, \seqname{C}^1)},
        {(\seqname{A}^2, \seqname{E}^2, \seqname{C}^2)}
    ]
  \defeq
  \biggl (
    \ID%
      [
        {(\seqname{E}^1,\seqname{C}^1)},
        {(\seqname{E}^2,\seqname{C}^2)}
      ],
    (\seqname{A}^1, \seqname{E}^1, \seqname{C}^1),
    (\seqname{A}^2, \seqname{E}^2, \seqname{C}^2)
  \biggr )
$,\hspace{1pt}
$
 \ID%
   [
     {(\seqname{E}^1, \seqname{C}^1)},
     {(\seqname{E}^2, \seqname{C}^2)}
   ]
$
considered as an $\seqname{E}^1$-$\seqname{E}^2$ m-atlas morphism of
$\seqname{A}^1$ to $\seqname{A}^2$ in the coordinate spaces
$\seqname{C}^1$, $\seqname{C}^2$.

The $\seqname{E}^i$ m-atlas identity morphism of $\seqname{A}^i$ in the
coordinate space $\seqname{C}^1$, $\seqname{C}^2$ is
$
 \ID%
   [
     {(\seqname{E}^i, \seqname{C}^i)}
   ]
$
considered as an $\seqname{E}^1$-$\seqname{E}^2$ m-atlas morphism of
$\seqname{A}^1$ to $\seqname{A}^2$ in the coordinate spaces
$\seqname{C}^1$, $\seqname{C}^2$.

The m-atlas identity morphism of
$(\seqname{A}^i, \seqname{E}^i, \seqname{C}^i)$ is \\*
$
  \Id[{{(\seqname{A}^i, \seqname{E}^i, \seqname{C}^i)}}] \defeq
  \bigl (
    \ID[{{(\seqname{E}^i,\seqname{C}^i)}}],
    (\seqname{A}^i, \seqname{E}^i, \seqname{C}^i),
    (\seqname{A}^i, \seqname{E}^i, \seqname{C}^i)
  \bigr )
$,\hspace{1pt}
$
 \ID%
   [
     {(\seqname{E}^i, \seqname{C}^i)}
   ]
$\hspace{1pt}
considered as an $\seqname{E}^i$ m-atlas morphism of $\seqname{A}^i$ in
the coordinate space $\seqname{C}^i$.

Let
\begin{enumerate}
\item
$E^i$ be a topological space, $i=1,2$
\item
$E^1 \subsp E^2$
\item
$\catName{C}^i$ be a model category
\item
$\seqname{C}^i \objin \catName{C}^i$
\item
$\seqname{C}^1 \submod \seqname{C}^2$
\item
$\seqname{A}^i$ be an m-atlas of $E^i$ in the coordinate space
$\seqname{C}^i$
\item
$\seqname{A}^1 \subseteq \seqname{A}^2$
\end{enumerate}

The $E^1$-$E^2$ m-atlas inclusion morphism of
$\seqname{A}^1$ to $\seqname{A}^2$ in the coordinate spaces
$\seqname{C}^1$, $\seqname{C}^2$ is
$
 \ID%
   [
     {(E^1, \seqname{C}^1)},
     {(E^2, \seqname{C}^2)}
   ]
$.

The m-atlas inclusion morphism of
$(\seqname{A}^1, E^1, \seqname{C}^1)$ to
$(\seqname{A}^2, E^2, \seqname{C}^2)$ is \\*
$
  \Id%
    [
        {(\seqname{A}^1, E^1, \seqname{C}^1)},
        {(\seqname{A}^2, E^2, \seqname{C}^2)}
    ]
  \defeq
  \biggl (
    \ID%
    [
      {(E^1,\seqname{C}^1)},
      {(E^2,\seqname{C}^2)}
    ],
    (\seqname{A}^1, E^1, \seqname{C}^1),
    (\seqname{A}^2, E^2, \seqname{C}^2)
  \biggr )
$,\hspace{1pt}
$
 \ID%
   [
     {(E^1, \seqname{C}^1)},
     {(E^2, \seqname{C}^2)}
   ]
$
considered as an $E^1$-$E^2$ m-atlas morphism of
$\seqname{A}^1$ to $\seqname{A}^2$ in the coordinate spaces
$\seqname{C}^1$, $\seqname{C}^2$.

The $E^i$ m-atlas identity morphism of $\seqname{A}^i$ in the
coordinate space $\seqname{C}^i$ is
$
 \ID%
   [
     {(E^i, \seqname{C}^i)}
   ]
$.

The m-atlas identity morphism of
$(\seqname{A}^i, E^i, \seqname{C}^i)$ is \\*
$
\Id[{{(\seqname{A}^i, E^i, \seqname{C}^i)}}] \defeq
\bigl (
  \ID[{{(E^i,\seqname{C}^i)}}],
  (\seqname{A}^i, E^i, \seqname{C}^i),
  (\seqname{A}^i, E^i, \seqname{C}^i)
\bigr )
$,\hspace{1pt}
$
 \ID%
   [
     {(E^i, \seqname{C}^i)}
   ]
$
considered as an $E^i$ m-atlas morphism of $\seqname{A}^i$ in
the coordinate space $\seqname{C}^i$.
\end{definition}

\begin{definition}[Equivalence of m-atlas morphisms]
\label{def:M-ATLmorphequiv}
Let $\seqname{E}^i$, $\seqname{C}^i$, $i=1,2$, be model spaces,
$\seqname{A}^i$ be an m-atlas of $\seqname{E}^i$ in the coordinate space
$\seqname{C}^i$, and \\*
$
  \morphSeqName{f} \defeq
  (
    \morphName{f}_0 \maps \seqname{E}^1 \to \seqname{E}^2,
    \morphName{f}_1 \maps \seqname{C}^1 \to \seqname{C}^2
  )
$ and
$
  \morphSeqName{g} \defeq
  (
    \morphName{g}_0 \maps \seqname{E}^1 \to \seqname{E}^2,
    \morphName{g}_1 \maps \seqname{C}^1 \to \seqname{C}^2
  )
$
be m-atlas morphisms of $\seqname{A}^1$ to $\seqname{A}^2$ in the
coordinate spaces $\seqname{C}^1$, $\seqname{C}^2$.  Then
$\morphSeqName{f}$ is equivalent to $\morphSeqName{g}$ as an
$\seqname{E}^1$-$\seqname{E}^2$ m-atlas morphism of
$\seqname{A}^1$ to $\seqname{A}^2$ in the coordinate spaces
$\seqname{C}^1$, $\seqname{C}^2$, abbreviated
$
  \morphSeqName{f} \maps
  (\seqname{A}^1, \seqname{E}^1, \seqname{C}^1) \to
  (\seqname{A}^2, \seqname{E}^2, \seqname{C}^2)
$
is equivalent to
$
  \morphSeqName{g} \maps
  (\seqname{A}^1, \seqname{E}^1, \seqname{C}^1) \to
  (\seqname{A}^2, \seqname{E}^2, \seqname{C}^2)
$
and
$
  \bigl (
    \morphSeqName{f},
    (\seqname{A}^1, \seqname{E}^1, \seqname{C}^1),
    (\seqname{A}^2, \seqname{E}^2, \seqname{C}^2)
  \bigr )
$
is equivalent to
$
  \bigl (
    \morphSeqName{g},
    (\seqname{A}^1, \seqname{E}^1, \seqname{C}^1),
    (\seqname{A}^2, \seqname{E}^2, \seqname{C}^2)
  \bigr )
$
iff $\morphName{f}_0 = \morphName{g}_0$.

By abuse of language we write $\morphSeqName{f}$ is equivalent to
$\morphSeqName{g}$ when the spaces and atlases are understood by context.

Let $E^i$, $i=1,2$, be a topological space, $\seqname{C}^i$ be a model
space, $\seqname{A}^i$ be an m-atlas of $E^i$ in the coordinate space
$\seqname{C}^i$ and
$
  \morphSeqName{f} \defeq
  (
    \morphName{f}_0 \maps E^1 \to E^2,
    \morphName{f}_1 \maps \seqname{C}^1 \to \seqname{C}^2
  )
$, and \\*
$
  \morphSeqName{g} \defeq
  (
    \morphName{g}_0 \maps E^1 \to E^2,
    \morphName{g}_1 \maps \seqname{C}^1 \to \seqname{C}^2
  )
$
be m-atlas morphisms of $\seqname{A}^1$ to $\seqname{A}^2$ in the
coordinate spaces $\seqname{C}^1$, $\seqname{C}^2$.
Then $\morphSeqName{f}$ is equivalent to $\morphSeqName{g}$ as an
$E^1$-$E^2$ m-atlas morphism of $\seqname{A}^1$ to
$\seqname{A}^2$ in the coordinate spaces $\seqname{C}^1$,
$\seqname{C}^2$, abbreviated
$
  \morphSeqName{f} \maps
  (\seqname{A}^1, E^1, \seqname{C}^1) \to
  (\seqname{A}^2, E^2, \seqname{C}^2)
$
is equivalent to
$
  \morphSeqName{g} \maps
  (\seqname{A}^1, E^1, \seqname{C}^1) \to
  (\seqname{A}^2, E^2, \seqname{C}^2)
$,
iff $\morphName{f}_0 = \morphName{g}_0$.

By abuse of language we write $\morphSeqName{f}$ is equivalent to
$\morphSeqName{g}$ when the spaces and atlases are understood by context.
\end{definition}

\subsubsection{Definitions of semi-strict and strict}
\begin{definition}[Semi-strict m-atlas morphisms]
\label{def:M-ATLmorphsemi}
Let $\catName{E}^i$, $\catName{C}^i$, $i=1,2$, be model categories,
$\seqname{E}^i \objin \catName{E}^i$,
$\seqname{C}^i \objin \catName{C}^i$, $\seqname{A}^i$ be an m-atlas of
$\seqname{E}^i$ in the coordinate space $\seqname{C}^i$ and
$
  \morphSeqName{f} \defeq
  (
    \morphName{f}_0 \maps \seqname{E}^1 \to \seqname{E}^2,
    \morphName{f}_1 \maps \seqname{C}^1 \to \seqname{C}^2
  )
$
an $\seqname{E}^1$-$\seqname{E}^2$ m-atlas morphism of
$\seqname{A}^1$ to $\seqname{A}^2$ in the coordinate spaces
$\seqname{C}^1$, $\seqname{C}^2$.

$\morphSeqName{f}$ is a semi-strict $\seqname{E}^1$-$\seqname{E}^2$
m-atlas morphism of $\seqname{A}^1$ to $\seqname{A}^2$ in the
coordinate spaces $\seqname{C}^1$, $\seqname{C}^2$, abbreviated as
$
  \strict[semi-]{\isAtl_\Ar}
  (
    \seqname{A}^1,
    \seqname{E}^1,
    \seqname{C}^1,
    \seqname{A}^2,
    \seqname{E}^2,
    \seqname{C}^2,
    \morphName{f}_0,
    \morphName{f}_1
  )
$,
iff $\morphName{f}_0$ is a local m-morphism of $\seqname{E}^1$
to $\seqname{E}^2$, $\morphName{f}_1$ is a local m-morphism of
$\seqname{C}^1$ to $\seqname{C}^2$ and for every \\*
$(U^i, V^i, \phi^i \maps U^i \toiso V^i) \in \seqname{A}^i$, $i=1,2$,
with $I \defeq U^1 \cap \morphName{f}^{-1}_0[U^2] \neq \emptyset$,
$
  \phi^2 \morphCompose \morphName{f}_0 \morphCompose \phi^{1-1} \maps
  \phi^1[I] \to V^2
$
is a local $\seqname{C}^1$-$\seqname{C}^2$ m-morphism of $\phi^1[I]$ to
$V^2$.

It is also a coherent semi-strict
\nobreaks{{$\seqname{E}^1$-$\seqname{E}^2$}}
m-atlas morphism of $\seqname{A}^1$ to $\seqname{A}^2$ in the
coordinate spaces $\seqname{C}^1$, $\seqname{C}^2$ iff
each $E^i \objin \seqname{E}^i$,
each $C^i \objin \seqname{C}^i$,
$\seqname{E}^1 \submod[strict-] \seqname{E}^2$ and
$\seqname{C}^1 \submod[strict-] \seqname{C}^2$.

$\morphSeqName{f}$ is a semi-strict
$\seqname{E}^1$-$\seqname{E}^2$-$\catName{C}^1$-$\catName{C}^2$ m-atlas
morphism of $\seqname{A}^1$ to $\seqname{A}^2$ in the coordinate
spaces $\seqname{C}^1$, $\seqname{C}^2$, abbreviated as
$
  \strict[semi-]{\isAtl_\Ar}
  (
    \seqname{A}^1,
    \seqname{E}^1,
    \seqname{C}^1,
    \seqname{A}^2,
    \seqname{E}^2,
    \seqname{C}^2,
    \morphName{f}_0,
    \morphName{f}_1,
    \catName{C}^1,
    \catName{C}^2
  )
$,
iff $\morphName{f}_0$ is a local
m-morphism of $\seqname{E}^1$ to $\seqname{E}^2$, $\morphName{f}_1$ is a
local $\catName{C}^1$-$\catName{C}^2$ m-morphism of $\seqname{C}^1$ to
$\seqname{C}^2$ and for every
$(U^i,V^i, \phi^i \maps U^i \toiso V^i) \in \seqname{A}^i$, $i=1,2$,
with $I \defeq U^1 \cap \morphName{f}^{-1}_0[U^2] \neq \emptyset$,
$
  \phi^2 \morphCompose \morphName{f}_0 \morphCompose \phi^{1-1} \maps
  \phi^1[I] \to V^2
$
is a local $\catName{C}^1$-$\catName{C}^2$ m-morphism of $\phi^1[I]$ to
$V^2$.

It is also a coherent semi-strict
\nobreaks
  {
    {$\seqname{E}^1$-$\seqname{E}^2$-},
    {$\catName{C}^1$-$\catName{C}^2$}
  }
m-atlas morphism of $\seqname{A}^1$ to $\seqname{A}^2$ in the
coordinate spaces $\seqname{C}^1$, $\seqname{C}^2$ iff
each $E^i \objin \seqname{E}^i$,
$\seqname{E}^1 \submod[strict-] \seqname{E}^2$ and
$\catName{C}^1 \subcat[strict-] \catName{C}^2$.

$\morphSeqName{f}$ is a semi-strict
$\catName{E}^1$-$\catName{E}^2$-$\seqname{E}^1$-$\seqname{E}^2$ m-atlas
morphism of $\seqname{A}^1$ to $\seqname{A}^2$ in the coordinate
spaces $\seqname{C}^1$, $\seqname{C}^2$, abbreviated as
$
  \strict[semi-]{\isAtl_\Ar}
  (
    \seqname{A}^1,
    \seqname{E}^1,
    \seqname{C}^1,
    \seqname{A}^2,
    \seqname{E}^2,
    \seqname{C}^2,
    \morphName{f}_0,
    \morphName{f}_1,
    \catName{E}^1,
    \catName{E}^2
  )
$,
iff $\morphName{f}_0$ is a local $\catName{E}^1$-$\catName{E}^2$
m-morphism of $\seqname{E}^1$ to $\seqname{E}^2$, $\morphName{f}_1$ is
a local m-morphism of $\seqname{C}^1$ to $\seqname{C}^2$ and for
every $(U^i, V^i, \phi^i \maps U^i \toiso V^i) \in \seqname{A}^i$,
$i=1,2$, with
$I \defeq U^1 \cap \morphName{f}^{-1}_0[U^2] \neq \emptyset$,
$
  \phi^2 \morphCompose \morphName{f}_0 \morphCompose \phi^{1-1} \maps
  \phi^1[I] \to U^2
$
is a local $\seqname{C}^1$-$\seqname{C}^2$ m-morphism of $\phi^1[I]$ to $V^2$.

It is also a coherent semi-strict
\nobreaks
  {
    {$\catName{E}^1$-$\catName{E}^2$-},
    {$\seqname{E}^1$-$\seqname{E}^2$}
  }
m-atlas morphism of $\seqname{A}^1$ to $\seqname{A}^2$ in the
coordinate spaces $\seqname{C}^1$, $\seqname{C}^2$ iff
each $C^i \objin \seqname{C}^i$,
$\catName{E}^1 \subcat[strict-] \catName{E}^2$ and
$\seqname{C}^1 \submod[strict-] \seqname{C}^2$.

$\morphSeqName{f}$ is a semi-strict
\nobreaks
  {
    {$\catName{E}^1$-$\catName{E}^2$-},
    {$\seqname{E}^1$-$\seqname{E}^2$-},
    {$\catName{C}^1$-$\catName{C}^2$}
  }
m-atlas morphism of $\seqname{A}^1$ to $\seqname{A}^2$ in the
coordinate spaces $\seqname{C}^1$, $\seqname{C}^2$, abbreviated as \\*
$
  \strict[semi-]{\isAtl_\Ar}
  (
    \seqname{A}^1,
    \seqname{E}^1,
    \seqname{C}^1,
    \seqname{A}^2,
    \seqname{E}^2,
    \seqname{C}^2,
    \morphName{f}_0,
    \morphName{f}_1,
    \catName{E}^1,
    \catName{E}^2,
    \catName{C}^1,
    \catName{C}^2
  )
$,
iff $\morphName{f}_0$ is a local $\catName{E}^1$-$\catName{E}^2$
m-morphism of $\seqname{E}^1$ to $\seqname{E}^2$ and $\morphName{f}_1$ is
a local $\catName{C}^1$-$\catName{C}^2$ m-morphism of $\seqname{C}^1$ to
$\seqname{C}^2$ and for every \\*
$(U^i,V^i, \phi^i \maps U^i \toiso V^i) \in \seqname{A}^i$, $i=1,2$,
with $I \defeq U^1 \cap \morphName{f}^{-1}_0[U^2] \neq \emptyset$,
$
  \phi^2 \morphCompose \morphName{f}_0 \morphCompose \phi^{1-1} \maps
  \phi^1[I] \to V^2
$
is a local \nobreaks{{$\catName{C}^1$-$\catName{C}^2$}} m-morphism of
$\phi^1[I]$ to $V^2$.

It is also a coherent semi-strict
\nobreaks
  {
    {$\catName{E}^1$-$\catName{E}^2$-},
    {$\seqname{E}^1$-$\seqname{E}^2$-},
    {$\catName{C}^1$-$\catName{C}^2$}
  }
m-atlas morphism of $\seqname{A}^1$ to $\seqname{A}^2$ in the
coordinate spaces $\seqname{C}^1$, $\seqname{C}^2$ iff
$\catName{E}^1 \subcat[strict-] \catName{E}^2$ and
$\catName{C}^1 \subcat[strict-] \catName{C}^2$.

By abuse of language we will write $\morphSeqName{f}$ is coherent when it
is a
\nobreaks
  {
    {($\catName{E}^1$-$\catName{E}^2$-)},
    {$\seqname{E}^1$-$\seqname{E}^2$},
    {(-$\catName{C}^1$-$\catName{C}^2)$}
  }
m-atlas morphism of $\seqname{A}^1$ to $\seqname{A}^2$ in the
coordinate spaces $\seqname{C}^1$, $\seqname{C}^2$ where the model
spaces, model categories and atlases are understood by context.
\end{definition}

\begin{definition}[Strict m-atlas morphisms]
\label{def:M-ATLmorphstrict}
Let
\begin{enumerate}
\item
$\catName{E}^i$, $\catName{C}^i$, $i=1,2$, be model categories
\item
$\seqname{E}^i \objin \catName{E}^i$
\item
$\seqname{C}^i \objin \catName{C}^i$
\item
$\seqname{A}^i$ be an m-atlas of
$\seqname{E}^i$ in the coordinate space $\seqname{C}^i$
\item
$
  \morphSeqName{f} \defeq
  (
    \morphName{f}_0 \maps \seqname{E}^1 \to \seqname{E}^2,
    \morphName{f}_1 \maps \seqname{C}^1 \to \seqname{C}^2
  )
$
be an m-atlas morphism
\end{enumerate}

$\morphSeqName{f}$ is a strict \nobreaks{$\seqname{E}^1-\seqname{E}^2$}
m-atlas morphism of $\seqname{A}^1$ to $\seqname{A}^2$ in the
coordinate spaces $\seqname{C}^1$, $\seqname{C}^2$, abbreviated as
$
   \strict{\isAtl_\Ar}
   (
     \seqname{A}^1,
     \seqname{E}^1,
     \seqname{C}^1,
     \seqname{A}^2,
     \seqname{E}^2,
     \seqname{C}^2,
     \morphName{f}_0,
     \morphName{f}_1
   )
$,
iff $\morphName{f}_0$ is an m-morphism of $\seqname{E}^1$ to
$\seqname{E}^2$, $\morphName{f}_1$ is an m-morphism of
$\seqname{C}^1$ to $\seqname{C}^2$ and for every \\*
$(U^i, V^i, \phi^i \maps U^i \toiso V^i) \in \seqname{A}^i$, $i=1,2$,
with $I \defeq U^1 \cap \morphName{f}^{-1}_0[U^2] \neq \emptyset$,
$
   \phi^2 \morphCompose \morphName{f}_0 \morphCompose \phi^{1-1} \maps
   \phi^1[I] \to U^2
$
is a morphism of $\Mod(V^2,\seqname{C}^2)$.

It is also a coherent strict
\nobreaks{{$\seqname{E}^1$-$\seqname{E}^2$}}
m-atlas morphism of $\seqname{A}^1$ to $\seqname{A}^2$ in the
coordinate spaces $\seqname{C}^1$, $\seqname{C}^2$ iff
each $E^i \objin \seqname{E}^i$,
each $C^i \objin \seqname{C}^i$,
$\seqname{E}^1 \submod[strict-] \seqname{E}^2$ and
$\seqname{C}^1 \submod[strict-] \seqname{C}^2$.

$\morphSeqName{f}$ is a strict
\nobreaks
  {
    {$\seqname{E}^1$-$\seqname{E}^2$-}
    {$\catName{C}^1$-$\catName{C}^2$}
  }
m-atlas morphism of $\seqname{A}^1$ to $\seqname{A}^2$ in the coordinate
spaces $\seqname{C}^1$, $\seqname{C}^2$, abbreviated as
$
  \strict{\isAtl_\Ar}
  (
    \seqname{A}^1,
    \seqname{E}^1,
    \seqname{C}^1,
    \seqname{A}^2,
    \seqname{E}^2,
    \seqname{C}^2,
    \morphName{f}_0,
    \morphName{f}_1,
    \catName{C}^1,
    \catName{C}^2
  )
$,
iff $\morphName{f}_0$ is an m-morphism of $\seqname{E}^1$ to
$\seqname{E}^2$, $\morphName{f}_1$ is an m-morphism of
$\catName{C}^1$ to $\catName{C}^2$ and for every \\*
$(U^i,V^i, \phi^i \maps U^i \toiso V^i) \in \seqname{A}^i$, $i=1,2$,
with $I \defeq U^1 \cap \morphName{f}^{-1}_0[U^2] \neq \emptyset$,
$
  \phi^2 \morphCompose \morphName{f}_0 \morphCompose \phi^{1-1} \maps
  \phi^1[I] \to V^2
$
is a morphism of $\catName{C}^2$.

It is also a coherent strict
\nobreaks
  {
    {$\seqname{E}^1$-$\seqname{E}^2$-},
    {$\catName{C}^1$-$\catName{C}^2$}
  }
m-atlas morphism of $\seqname{A}^1$ to $\seqname{A}^2$ in the
coordinate spaces $\seqname{C}^1$, $\seqname{C}^2$ iff
each $E^i \objin \seqname{E}^i$,
$\seqname{E}^1 \submod[strict-] \seqname{E}^2$ and
$\catName{C}^1 \subcat[strict-] \catName{C}^2$.

$\morphSeqName{f}$ is a strict
\nobreaks
  {
    {$\catName{E}^1$-$\catName{E}^2$-},
    {$\seqname{E}^1$-$\seqname{E}^2$}
  }
m-atlas morphism of $\seqname{A}^1$ to $\seqname{A}^2$ in the coordinate
spaces $\seqname{C}^1$, $\seqname{C}^2$, abbreviated as
$
  \strict{\isAtl_\Ar}
  (
    \seqname{A}^1,
    \seqname{E}^1,
    \seqname{C}^1,
    \seqname{A}^2,
    \seqname{E}^2,
    \seqname{C}^2,
    \morphName{f}_0,
    \morphName{f}_1,
    \catName{E}^1,
    \catName{E}^2
  )
$,
iff $\morphName{f}_0$ is an m-morphism of $\catName{E}^1$ to
$\catName{E}^2$, $\morphName{f}_1$ is an m-morphism of $\seqname{C}^1$ to
$\seqname{C}^2$ and for every \\*
$(U^i, V^i, \phi^i \maps U^i \toiso V^i) \in \seqname{A}^i$,
$i=1,2$, with
$I \defeq U^1 \cap \morphName{f}^{-1}_0[U^2] \neq \emptyset$,
$
  \phi^2 \morphCompose \morphName{f}_0 \morphCompose \phi^{1-1} \maps
  \phi^1[I] \to U^2
$
is a morphism of $\Mod(\phi^1[I],\seqname{C}^2)$.

It is also a coherent strict
\nobreaks
  {
    {$\catName{E}^1$-$\catName{E}^2$-},
    {$\seqname{E}^1$-$\seqname{E}^2$}
  }
m-atlas morphism of $\seqname{A}^1$ to $\seqname{A}^2$ in the
coordinate spaces $\seqname{C}^1$, $\seqname{C}^2$ iff
$\seqname{C}^1 \submod[strict-] \seqname{C}^2$.
$\catName{E}^1 \subcat[strict-] \catName{E}^2$ and
$\seqname{C}^1 \submod[strict-] \seqname{C}^2$.

$\morphSeqName{f}$ is a strict
\nobreaks
  {
    {$\catName{E}^1$-$\catName{E}^2$-},
    {$\seqname{E}^1$-$\seqname{E}^2$-},
    {$\catName{C}^1$-$\catName{C}^2$}
  }
m-atlas morphism of $\seqname{A}^1$ to $\seqname{A}^2$ in the
coordinate spaces $\seqname{C}^1$, $\seqname{C}^2$, abbreviated as \\*
$
  \strict{\isAtl_\Ar}
  (
    \seqname{A}^1,
    \seqname{E}^1,
    \seqname{C}^1,
    \seqname{A}^2,
    \seqname{E}^2,
    \seqname{C}^2,
    \morphName{f}_0,
    \morphName{f}_1,
    \catName{E}^1,
    \catName{E}^2,
    \catName{C}^1,
    \catName{C}^2
  )
$,
iff $\morphName{f}_0$ is an m-morphism of $\catName{E}^1$ to
$\catName{E}^2$, $\morphName{f}_1$ is an m-morphism of $\catName{C}^1$ to
$\catName{C}^2$ and for every
$(U^i,V^i, \phi^i \maps U^i \toiso V^i) \in \seqname{A}^i$, $i=1,2$,
with $I \defeq U^1 \cap \morphName{f}^{-1}_0[U^2] \neq \emptyset$,
$
  \phi^2 \morphCompose \morphName{f}_0 \morphCompose \phi^{1-1} \maps
  \phi^1[I] \to V^2
$
is a morphism of $\catName{C}^2$.

It is also a coherent strict
\nobreaks
  {
    {$\catName{E}^1$-$\catName{E}^2$-},
    {$\seqname{E}^1$-$\seqname{E}^2$-},
    {$\catName{C}^1$-$\catName{C}^2$}
  }
m-atlas morphism of $\seqname{A}^1$ to $\seqname{A}^2$ in the
coordinate spaces $\seqname{C}^1$, $\seqname{C}^2$ iff
$\catName{E}^1 \subcat[strict-] \catName{E}^2$ and
$\catName{C}^1 \subcat[strict-] \catName{C}^2$.

By abuse of language we will write $\morphSeqName{f}$ is coherent when it
is a
\nobreaks
  {
    {($\catName{E}^1$-$\catName{E}^2$-)},
    {$\seqname{E}^1$-$\seqname{E}^2$},
    {(-$\catName{C}^1$-$\catName{C}^2$)}
  }
m-atlas morphism of $\seqname{A}^1$ to $\seqname{A}^2$ in the
coordinate spaces $\seqname{C}^1$, $\seqname{C}^2$ where the model
spaces, model categories and atlases are understood by context.
\end{definition}

\subsubsection{Abbreviated nomenclature for m-atlas morphisms}
\begin{definition}[Abbreviated nomenclature for m-atlas morphisms]
\label{def:M-ATLmorphabbrev}

Let $\catName{E}^i$, $\catName{C}^i$, $i=1,2$, be model categories,
$\seqname{E}^i \objin \catName{E}^i$,
$\seqname{C}^i \objin \catName{C}^i$, $\seqname{A}^i$ be an m-atlas of
$\seqname{E}^i$ in the coordinate space $\seqname{C}^i$ and
$
  \morphSeqName{f} \defeq
  (
    \morphName{f}_0 \maps \seqname{E}^1 \to \seqname{E}^2,
    \morphName{f}_1 \maps \seqname{C}^1 \to \seqname{C}^2
  )
$
a pair of model functions.

$\morphSeqName{f}$ is a constrained semi-strict~(strict)
$\seqname{E}^1$-$\seqname{E}^2$ m-atlas morphism of
$\seqname{A}^1$ to $\seqname{A}^2$ in the coordinate spaces
$\seqname{C}^1$, $\seqname{C}^2$, abbreviated as \\*
$
  \strict[**]{\isAtl[const]_\Ar}
  (
    \seqname{A}^1,
    \seqname{E}^1,
    \seqname{C}^1,
    \seqname{A}^2,
    \seqname{E}^2,
    \seqname{C}^2,
    \morphName{f}_0,
    \morphName{f}_1
  )
$,
iff it is a constrained $\seqname{E}^1$-$\seqname{E}^2$ m-atlas
morphism of $\seqname{A}^1$ to $\seqname{A}^2$ in the coordinate spaces
$\seqname{C}^1$, $\seqname{C}^1$ and it is a semi-strict~(strict)
$\seqname{E}^1$-$\seqname{E}^2$ m-atlas morphism of
$\seqname{A}^1$ to $\seqname{A}^2$ in the coordinate spaces
$\seqname{C}^1$, $\seqname{C}^1$.

$\morphSeqName{f}$ is a (constrained) (semi-strict, strict)
$\seqname{E}^1$-$\seqname{E}^2$ m-atlas morphism of
$\seqname{A}^1$ to $\seqname{A}^2$ in the coordinate space
$\seqname{C}^1$, abbreviated as \\*
$
  \strict[*]{\isAtl[(const)]_\Ar}
  (
    \seqname{A}^1,
    \seqname{E}^1,
    \seqname{C}^1,
    \seqname{A}^2,
    \seqname{E}^2,
    \morphName{f}_0,
    \morphName{f}_1
  )
$,
iff it is a (constrained) (semi-strict, strict)
$\seqname{E}^1$-$\seqname{E}^2$ m-atlas morphism of
$\seqname{A}^1$ to $\seqname{A}^2$ in the coordinate spaces
$\seqname{C}^1$, $\seqname{C}^1$.

$\morphSeqName{f}$ is a (constrained) (semi-strict, strict)
\nobreaks
  {
    $\catName{E}^1$-$\catName{E}^2$-,
    $\seqname{E}^1$-$\seqname{E}^2$
  }
m-atlas morphism of $\seqname{A}^1$ to $\seqname{A}^2$ in the
coordinate space $\seqname{C}^1$, abbreviated as \\*
$
  \strict[*]{\isAtl[(const)]_\Ar}
  (
    \seqname{A}^1,
    \seqname{E}^1,
    \seqname{C}^1,
    \seqname{A}^2,
    \seqname{E}^2,
    \morphName{f}_0,
    \morphName{f}_1,
    \catName{E}^1,
    \catName{E}^2
  )
$,
iff it is a (constrained) (semi-strict, strict)
\nobreaks
  {
    {$\catName{E}^1$-$\catName{E}^2$-},
    {$\seqname{E}^1$-$\seqname{E}^2$}
  }
m-atlas morphism of $\seqname{A}^1$ to $\seqname{A}^2$ in
the coordinate spaces $\seqname{C}^1$, $\seqname{C}^1$.

$\morphSeqName{f}$ is a (constrained) (semi-strict, strict)
$\seqname{E}^1$ m-atlas morphism of $\seqname{A}^1$ to
$\seqname{A}^2$ in the coordinate spaces $\seqname{C}^1$,
$\seqname{C}^2$, abbreviated as \\*
$
  \strict[*]{\isAtl[(const)]_\Ar}
  (
    \seqname{A}^1,
    \seqname{E}^1,
    \seqname{C}^1,
    \seqname{A}^2,
    \seqname{C}^2,
    \morphName{f}_0,
    \morphName{f}_1
    )
$,
iff it is a (constrained) (semi-strict, strict)
$\seqname{E}^1$-$\seqname{E}^1$ m-atlas morphism of
$\seqname{A}^1$ to $\seqname{A}^2$ in the coordinate spaces
$\seqname{C}^1$, $\seqname{C}^2$.

$\morphSeqName{f}$ is a (constrained) (semi-strict, strict)
$\seqname{E}^1$ m-atlas morphism of $\seqname{A}^1$ to
$\seqname{A}^2$ in the coordinate space $\seqname{C}^1$, abbreviated as
\\*
$
  \strict[*]{\isAtl[(const)]_\Ar}
  (
    \seqname{A}^1,
    \seqname{E}^1,
    \seqname{C}^1,
    \seqname{A}^2,
    \morphName{f}_0,
    \morphName{f}_1
  )
$,
iff it is a (constrained) (semi-strict, strict) $\seqname{E}^1$-$\seqname{E}^1$
m-atlas morphism of
$\seqname{A}^1$ to $\seqname{A}^2$ in the coordinate spaces
$\seqname{C}^1$, $\seqname{C}^1$.

$\morphSeqName{f}$ is a (constrained) (semi-strict, strict)
$\seqname{E}^1$-$\catName{C}^1$-$\catName{C}^2$ m-atlas
morphism of $\seqname{A}^1$ to $\seqname{A}^2$ in the coordinate
spaces $\seqname{C}^1$, $\seqname{C}^2$, abbreviated as \\*
$
  \strict[*]{\isAtl[(const)]_\Ar}
    (
      \seqname{A}^1,
      \seqname{E}^1,
      \seqname{C}^1,
      \seqname{A}^2,
      \seqname{C}^2,
      \morphName{f}_0,
      \morphName{f}_1,
      \catName{C}^1,
      \catName{C}^2
    )
$,
iff it is a (constrained) (semi-strict, strict)
$\seqname{E}^1$-$\seqname{E}^1$-$\catName{C}^1$-$\catName{C}^2$ m-atlas
morphism of $\seqname{A}^1$ to $\seqname{A}^2$ in the coordinate
spaces $\seqname{C}^1$, $\seqname{C}^2$.

$\morphSeqName{f}$ is a (constrained) (semi-strict, strict)
$\seqname{E}^1$-$\catName{C}^1$ m-atlas morphism of
$\seqname{A}^1$ to $\seqname{A}^2$ in the coordinate space
$\seqname{C}^1$, abbreviated as \\*
$
  \strict[*]{\isAtl[(const)]_\Ar}
  (
    \seqname{A}^1,
    \seqname{E}^1,
    \seqname{C}^1,
    \seqname{A}^2,
    \morphName{f}_0,
    \morphName{f}_1,
    \catName{C}^1
  )
$,
iff it is a (constrained) (semi-strict, strict)
$\seqname{E}^1$-$\seqname{E}^1$-$\catName{C}^1$-$\catName{C}^1$ m-atlas
morphism of $\seqname{A}^1$ to $\seqname{A}^2$ in the coordinate
spaces $\seqname{C}^1$, $\seqname{C}^1$.

Let $\catName{C}^i$, $i=1,2$, be a model category,
$\seqname{C}^i \objin \catName{C}^i$, $E^i$ be a topological space,
$\seqname{A}^i$ be an m-atlas of $E^i$ in the coordinate space
$\seqname{C}^i$, $\morphName{f}_0 \maps E^1 \to E^2$ be continuous,
$\morphName{f}_1 \maps \seqname{C}^1 \to \seqname{C}^2$ be a model
function and
$
  \morphSeqName{f} \defeq
  (
    \morphName{f}_0 \maps E^1 \to E^2,
    \morphName{f}_1 \maps \seqname{C}^1 \to \seqname{C}^2
  )
$.

$\morphSeqName{f}$ is a (constrained) (semi-strict, strict) $E^1$-$E^2$
m-atlas morphism of $\seqname{A}^1$ to $\seqname{A}^2$ in the
coordinate spaces $\seqname{C}^1$, $\seqname{C}^2$, abbreviated as \\*
$
  \strict[*]{\isAtl[(const)]_\Ar}
    (
      \seqname{A}^1,
      E^1,
      \seqname{C}^1,
      \seqname{A}^2,
      E^2,
      \seqname{C}^2,
      \morphName{f}_0,
      \morphName{f}_1
    )
$,
and a (constrained) (semi-strict, strict) $E^1$-$E^2$ m-atlas
morphism of $\seqname{A}^1$ to $\seqname{A}^2$ in the coordinate model
categories $\catName{C}^1$, $\catName{C}^2$, abbreviated as \\*
$
  \strict[*]{\isAtl[(const)]_\Ar}
    (
      \seqname{A}^1,
      E^1,
      \catName{C}^1,
      \seqname{A}^2,
      E^2,
      \catName{C}^2,
      \morphName{f}_0,
      \morphName{f}_1
    )
$,
iff it is a (constrained) (semi-strict, strict)
$\Triv{E^1}$-$\Triv{E^2}$ m-atlas morphism of $\seqname{A}^1$ to
$\seqname{A}^2$ in the coordinate model spaces $\seqname{C}^1$,
$\seqname{C}^2$.

$\morphSeqName{f}$ is a (constrained) (semi-strict, strict) $E^1$-$E^2$
m-atlas (constrained) morphism of $\seqname{A}^1$ to
$\seqname{A}^2$ in the coordinate space $\seqname{C}^1$, abbreviated as
\\*
$
  \strict[*]{\isAtl[(const)]_\Ar}
  (
    \seqname{A}^1,
    E^1,
    \seqname{C}^1,
    \seqname{A}^2,
    E^2,
    \morphName{f}_0,
    \morphName{f}_1
  )
$,
iff it is a (constrained) (semi-strict, strict) $\Triv{E^1}$-$\Triv{E^2}$
m-atlas morphism of $\seqname{A}^1$ to $\seqname{A}^2$ in the coordinate
spaces $\seqname{C}^1$, $\seqname{C}^1$.

$\morphSeqName{f}$ is a (constrained) (semi-strict, strict)
\nobreaks
  {
    {$E^1$-$E^2$-},
    {$\catName{C}^1$-$\catName{C}^2$}
  }
m-atlas morphism of $\seqname{A}^1$ to $\seqname{A}^2$ in the
coordinate spaces $\seqname{C}^1$, $\seqname{C}^2$, abbreviated as \\*
$
  \strict[*]{\isAtl[(const)]_\Ar}
  (
    \seqname{A}^1,
    E^1,
    \seqname{C}^1,
    \seqname{A}^2,
    E^2,
    \seqname{C}^2,
    \morphName{f}_0,
    \morphName{f}_1,
    \catName{C}^1,
    \catName{C}^2
  )
$,
iff it is a (constrained) (semi-strict, strict)
\nobreaks
  {
    {$\Triv{E^1}$-$\Triv{E^2}$-},
    {$\catName{C}^1$-$\catName{C}^2$}
  }
m-atlas morphism of $\seqname{A}^1$ to $\seqname{A}^2$ in the
coordinate model spaces $\seqname{C}^1$, $\seqname{C}^2$.

$\morphSeqName{f}$ is a (constrained) (semi-strict, strict)
$\Triv{E^1}$-$\Triv{E^2}$-$\catName{C}^1$
m-atlas morphism of $\seqname{A}^1$
to $\seqname{A}^2$ in the coordinate model category $\catName{C}^1$,
abbreviated as \\*
$
  \strict[*]{\isAtl[(const)]_\Ar}
  (
    \seqname{A}^1,
    E^1,
    \seqname{C}^1,
    \seqname{A}^2,
    E^2,
    \morphName{f}_0,
    \morphName{f}_1,
    \catName{C}^1
  )
$,
iff it is a (constrained) (semi-strict, strict)
\nobreaks
  {
    {$\Triv{E^1}$-$\Triv{E^2}$-},
    {$\catName{C}^1$-$\catName{C}^1$}
  }
m-atlas morphism of $\seqname{A}^1$ to $\seqname{A}^2$ in the
coordinate model categories $\catName{C}^1$, $\catName{C}^1$.

Similar definitions apply with restrictions on the admissible
atlases, i.e.,
$\morphSeqName{f}$ is a (constrained) (semi-strict, strict)
full (maximal, full maximal)
\nobreaks
  {
    {$\catName{E}^1$-$\catName{E}^2$-},
    {$\seqname{E}^1$-$\seqname{E}^2$-},
    {$\catName{C}^1$-$\catName{C}^2$}
  }
m-atlas morphism of
$\seqname{A}^1$ to $\seqname{A}^2$ in the coordinate spaces
$\seqname{C}^1$, abbreviated as \\*

\begin{equation*}
\begin{array}{lll}
\strict[*]
  {
    \maxfull[**]
      {
        \isAtl[(const)]_\Ar
          (
            \seqname{A}^1,
            \seqname{E}^1,
            \seqname{C}^1,
            \seqname{A}^2,
            \seqname{E}^2,
            \seqname{C}^2,
            \morphName{f}_0,
            \morphName{f}_1,
            \catName{E}^1,
            \catName{E}^2,
            \catName{C}^1,
            \catName{C}^2
          )
      }
  }
  & & \iff
\\*
\strict[**]
  {
    \isAtl[(const)]_\Ar
    (
      \seqname{A}^1,
      \seqname{E}^1,
      \seqname{C}^1,
      \seqname{A}^2,
      \seqname{E}^2,
      \seqname{C}^2,
      \morphName{f}_0,
      \morphName{f}_1,
      \catName{E}^1,
      \catName{E}^2,
      \catName{C}^1,
      \catName{C}^2
    )
  }
  & \land
\\*
\maxfull[**]{\isAtl_\Ob}
  (
    \seqname{A}^1,
    \seqname{E}^1,
    \seqname{C}^1
  )
  & \land
\\*
\maxfull[**]{\isAtl_\Ob}
  (
    \seqname{A}^2,
    \seqname{E}^2,
    \seqname{C}^2
  )
\end{array}
\end{equation*}

\end{definition}

\subsubsection{Proclamations on m-atlas morphisms}
\begin{lemma}[M-atlas morphisms]
\label{lem:M-ATLmorph}
Let
\begin{enumerate}
\item
$\catName{E}^i$, $\catName{C}^i$,
$\catName{E}'^i$, $\catName{C}'^i$, $i=1,2$, be model categories

\item
$\catName{E}^i \subcat[full-] \catName{E}'^i$,
$\catName{C}^i \subcat[full-] \catName{C}'^i$

\item
$\seqname{E}^i \objin \catName{E}^i$,
$\seqname{C}^i \objin \catName{C}^i$,
$\seqname{C}'^i \objin \catName{C}'^i$,
$\seqname{C}^i \submod \seqname{C}'^i$

\item
$\morphName{f}'_1 \maps \seqname{C}'^1 \to \seqname{C}'^2$ be a model
function such that
$\morphName{f}'_1[\seqname{C}^1] \subseteq \seqname{C}^2$ and
$
  \morphName{f}_1 \defeq
  \morphName{f}'_1 \restrictto_{\seqname{C}^1,\seqname{C}^2}
$
is a model function

\item
$\seqname{A}^i$ be an M-atlas
of $\seqname{E}^i$ in the coordinate space $\seqname{C}^i$

\item
$
  \morphSeqName{f} \defeq
  (
    \morphName{f}_0 \maps \seqname{E}^1 \to \seqname{E}^2,
    \morphName{f}_1 \maps \seqname{C}^1 \to \seqname{C}^2
  )
$
be a (semi-strict, strict)
\nobreaks
  {
    {($\catName{E}^1$-$\catName{E}^2$-)},
    {$\seqname{E}^1$-$\seqname{E}^2$},
    {(-$\catName{C}^1$-$\catName{C}^2$)}
  }
m-atlas morphism of $\seqname{A}^1$ to $\seqname{A}^2$ in the
coordinate spaces $\seqname{C}^1$, $\seqname{C}^2$

\item
$
  \morphSeqName{f}' \defeq
  (
    \morphName{f}_0 \maps \seqname{E}^1 \to \seqname{E}^2,
    \morphName{f}'_1 \maps \seqname{C}^1 \to \seqname{C}'^2
  )
$, \\*
$
  \morphSeqName{f}'' \defeq
  (
    \morphName{f}_0 \maps \seqname{E}^1 \to \seqname{E}^2,
    \morphName{f}'_1 \maps \seqname{C}'^1 \to \seqname{C}'^2
  )
$
be pairs of model functions
\end{enumerate}

Then

\begin{enumerate}
\item
\label{M-ATLmorph:fineistopmorph}
Let $E^i \defeq \Top(\seqname{E}^i)$, $i=1,2$. If each $\seqname{C}^i$
is fine grained then \\*
$
   \hat{\morphName{f}} \defeq
   (
      \morphName{f}_0 \maps E^1 \to E^2,
      \morphName{f}_1 \maps \seqname{C}^1 \to \seqname{C}^2
   )
$
is a (semi-strict, strict) $E^1$-$E^2$
m-atlas morphism of $\seqname{A}^1$ to $\seqname{A}^2$ in the
coordinate spaces $\seqname{C}^1$, $\seqname{C}^2$.

\begin{proof}
\leavevmode
\begin{enumerate}
\item
$\seqname{A}^i$ is an m-atlas of $E^i$ in the coordinate space
$\seqname{C}^i$ by \pagecref{lem:M-atlas}.

\item
$\morphName{f}_0 \maps \seqname{E}^1 \to \seqname{E}^2$ is a model
function, hence continuous, hence by \pagecref{lem:modfunctriv},
$\morphName{f}_0 \maps \Triv{E^1} \to \Triv{E^2}$ is a model function.

\item
$\morphName{f}_1 \maps \seqname{C}^1 \to \seqname{C}^2$ is a model
function by hypothesis.

\item
Any model neighborhood of $\seqname{E}^i$ is open, hence a model
neighborhood of $\Triv{E^i}$.

\item
Let $(U^i,V^i,\phi^i) \in \seqname{A}^i$, $i=1,2$.  All of the model
neighborhoods of $\seqname{E}^i$ mentioned in
\pagecref{def:M-ATLmorph} are also
model neighborhoods of $\Triv{E^i}$.
\end{enumerate}
\end{proof}

\item
\label{M-ATLmorph:istosup}
$\morphSeqName{f}'$ is an $\seqname{E}^1$-$\seqname{E}^2$ m-atlas
morphism of $\seqname{A}^1$ to $\seqname{A}^2$ in the coordinate spaces
$\seqname{C}^1$, $\seqname{C}'^2$ and $\morphSeqName{f}''$ is an
$\seqname{E}^1$-$\seqname{E}^2$ m-atlas morphism of
$\seqname{A}^1$ to $\seqname{A}^2$ in the coordinate spaces
$\seqname{C}'^1$, $\seqname{C}'^2$.

\begin{proof}
$\morphName{f}_0 \maps \seqname{E}^1 \to \seqname{E}^2$ is a model
function by \pagecref{def:M-ATLmorph}.

$\morphName{f}'_1 \maps \seqname{C}'^1 \to \seqname{C}'^2$ and
$\morphName{f}'_1 \maps \seqname{C}^1 \to \seqname{C}'^2$ are model
functions by hypothesis.

$\seqname{C}^i \submod \seqname{C}'^i$ by hypothesis,

$C^i$ is an open subset of $C'^i$ by \pagecref{lem:modsub}.

The expanded model space does not change the atlases or the commutation.
\end{proof}

\item
\label{M-ATLmorph:strictissemi}
If $\morphSeqName{f}$ is a strict
\nobreaks
  {
    {($\catName{E}^1$-$\catName{E}^2$-)},
    {$\seqname{E}^1$-$\seqname{E}^2$},
    {(-$\catName{C}^1$-$\catName{C}^2$)}
  }
m-atlas morphism of $\seqname{A}^1$ to $\seqname{A}^2$ in the
coordinate spaces $\seqname{C}^1$, $\seqname{C}^2$ then
$\morphSeqName{f}$ is a semi-strict
\nobreaks
  {
    {($\catName{E}^1$-$\catName{E}^2$-)},
    {$\seqname{E}^1$-$\seqname{E}^2$},
    {(-$\catName{C}^1$-$\catName{C}^2$)}
  }
m-atlas morphism of $\seqname{A}^1$ to $\seqname{A}^2$ in the
coordinate spaces $\seqname{C}^1$, $\seqname{C}^2$

\begin{proof}
There are four overlapping cases:
\begin{description}
\item[$\seqname{E}^i$ without category:]
$\morphName{f}_0$ is an m-morphism of $\seqname{E}^1$ to $\seqname{E}^2$
by \pagecref{def:M-ATLmorphstrict}. $\morphName{f}_0$ is a local
m-morphism of $\seqname{E}^1$ to $\seqname{E}^2$ by
\pagecref{lem:modMmorph}.

\item[$\seqname{C}^i$ without category:]
$\morphName{f}_1$ is an m-morphism of $\seqname{C}^1$ to $\seqname{C}^2$
by \cref{def:M-ATLmorphstrict}. $\morphName{f}_1$ is a local m-morphism
of $\seqname{C}^1$ to $\seqname{C}^2$ by \cref{lem:modMmorph}.

\item[$\seqname{E}^i$ with category:]
$\morphName{f}_0$ is an $\catName{E}^1$-$\catName{E}^2$ m-morphism of
$\seqname{E}^1$ to $\seqname{E}^2$ by \cref{def:M-ATLmorphstrict}.
$\morphName{f}_0$ is a local $\catName{E}^1$-$\catName{E}^2$ m-morphism
of $\seqname{E}^1$ to $\seqname{E}^2$ by \cref{lem:modMmorph}.

\item[$\seqname{C}^i$ with category:]
$\morphName{f}_1$ is a $\catName{C}^1$-$\catName{C}^2$ m-morphism of
$\seqname{C}^1$ to $\seqname{C}^2$ by \cref{def:M-ATLmorphstrict}.
$\morphName{f}_1$ is a local $\catName{C}^1$-$\catName{C}^2$ m-morphism
of $\seqname{C}^1$ to $\seqname{C}^2$ by \cref{lem:modMmorph}.

\end{description}
\end{proof}

\item
\label{M-ATLmorph:strictSstrictisstrict}
If $\morphSeqName{f}$ is a coherent semi-strict
\nobreaks
  {
    {($\catName{E}^1$-$\catName{E}^2$-)},
    {$\seqname{E}^1$-$\seqname{E}^2$},
    {(-$\catName{C}^1$-$\catName{C}^2$)}
  }
m-atlas morphism of $\seqname{A}^1$ to $\seqname{A}^2$ in the
coordinate spaces $\seqname{C}^1$, $\seqname{C}^2$, then
$\morphSeqName{f}$ is a strict
\nobreaks
  {
    {($\catName{E}^1$-$\catName{E}^2$-)},
    {$\seqname{E}^1$-$\seqname{E}^2$},
    {(-$\catName{C}^1$-$\catName{C}^2$)}
  }
m-atlas morphism of $\seqname{A}^1$ to $\seqname{A}^2$ in the
coordinate spaces $\seqname{C}^1$, $\seqname{C}^2$.

\begin{proof}
There are four overlapping cases.
\begin{description}
\item[$\seqname{E}^i$ without category:]
\leavevmode
$\morphName{f}_0$ is a local m-morphism of $\seqname{E}^1$ to
$\seqname{E}^2$ by \pagecref{def:M-ATLmorphsemi}.
$\seqname{E}^1 \submod[strict-] \seqname{E}^2$ by \pagecref{def:M-ATLmorphsemi}.
$\Set(\seqname{E}^2)$ is a model neighborhood of $\seqname{E}^2$ by \pagecref{def:modsub}.
The result follows by \cref{mod:sheaf} (Restricted sheaf condition) of
{
  \showlabelsinline
  \cref{def:modsub}
}.

\item[$\seqname{E}^i$ with category:]
\leavevmode
$\morphName{f}_0$ is a local $\catName{E}^1$-$\catName{E}^2$ m-morphism
of $\seqname{E}^1$ to $\seqname{E}^2$ by \pagecref{def:M-ATLmorphsemi}.
$\catName{E}^1 \subcat[strict-] \catName{E}^2$ by
\cref{def:M-ATLmorphsemi}.  $\catName{E}^2$ satisfies the restricted
sheaf condition by \pagecref{def:modcat} and the result follows.

\item[$\seqname{C}^i$ without category:]
\leavevmode
$\morphName{f}_0$ is a local m-morphism of $\seqname{C}^1$ to
$\seqname{C}^2$ by \pagecref{def:M-ATLmorphsemi}.
$\seqname{C}^1 \submod[strict-] \seqname{C}^2$ by \pagecref{def:M-ATLmorphsemi}.
$\Set(\seqname{C}^2)$ is a model neighborhood of $\seqname{C}^2$ by \pagecref{def:modsub}.
The result follows by \cref{mod:sheaf} (Restricted sheaf condition) of
{
  \showlabelsinline
  \cref{def:modsub}
}.

\item[$\seqname{C}^i$ with category:]
\leavevmode
$\morphName{f}_0$ is a local $\catName{C}^1$-$\catName{C}^2$ m-morphism
of $\seqname{C}^1$ to $\seqname{C}^2$ by \pagecref{def:M-ATLmorphsemi}.
$\catName{C}^1 \subcat[strict-] \catName{C}^2$ by
\cref{def:M-ATLmorphsemi}.  $\catName{E}^2$ satisfies the restricted
sheaf condition by \pagecref{def:modcat} and the result follows.
\end{description}
\end{proof}

\item
\label{M-ATLmorph:SstrictisSstricttosup}
If $\morphSeqName{f}$ is a semi-strict
\nobreaks
  {
    {($\catName{E}^1$-$\catName{E}^2$-)},
    {$\seqname{E}^1$-$\seqname{E}^2$},
    {(-$\catName{C}^1$-$\catName{C}^2$)}
  }
m-atlas morphism of $\seqname{A}^1$ to $\seqname{A}^2$ in the
coordinate spaces $\seqname{C}^1$, $\seqname{C}^2$ then
$\morphSeqName{f}'$ is a semi-strict
\nobreaks
  {
    {($\catName{E}'^1$-$\catName{E}'^2$-)},
    {$\seqname{E}^1$-$\seqname{E}^2$},
    {(-$\catName{C}'^1$-$\catName{C}'^2$)}
  }
m-atlas morphism of $\seqname{A}^1$ to $\seqname{A}^2$ in the
coordinate spaces $\seqname{C}^1$, $\seqname{C}'^2$.

\begin{proof}
There are four overlapping cases.
\begin{description}
\item[$\seqname{E}^i$ without category:]
$\morphName{f}_0$ is a (strict) local m-morphism of $\seqname{E}^1$ to
$\seqname{E}^2$ by \pagecref{def:M-ATLmorphsemi}.

\item[$\seqname{E}^i$ with category:]
$\morphName{f}_0$ is a (strict) local $\catName{E}^1$-$\catName{E}^2$
m-morphism of $\seqname{E}^1$ to $\seqname{E}^2$ by
\pagecref{def:M-ATLmorphsemi}.
$\catName{E}^i \subcat[full-] \catName{E}'^i$ by hypothesis.
$\morphName{f}_0$ is a (strict) local $\catName{E}'^1$-$\catName{E}'^2$
m-morphism of $\seqname{E}^1$ to $\seqname{E}^2$ by
\pagecref{lem:modMmorph}.

\item[$\seqname{C}^i$ without category:]
$\morphName{f}_1$ is a (strict) local m-morphism of $\seqname{C}^1$ to
$\seqname{C}^2$ by \pagecref{def:M-ATLmorphsemi}.
$\seqname{C}^2 \submod \seqname{C}'^2$ by hypothesis.
$\morphName{f}_1$ is a (strict) local m-morphism of $\seqname{C}^1$ to
$\seqname{C}'^2$ by \pagecref{lem:modMmorph}.

\item[$\seqname{C}^i$ with category:]
$\morphName{f}_0$ is a local $\catName{C}^1$-$\catName{C}^2$
m-morphism of $\seqname{C}^1$ to $\seqname{C}^2$ by
\pagecref{def:M-ATLmorphsemi}.
$\catName{C}^i \subcat[full-] \catName{C}'^i$ by hypothesis.
$\morphName{f}_0$ is a (strict) local $\catName{C}'^1$-$\catName{C}'^2$
m-morphism of $\seqname{C}^1$ to $\seqname{C}^2$ by
\pagecref{lem:modMmorph}.
\end{description}
\end{proof}

\item
\label{M-ATLmorph:strictisstricttosup}
If $\morphSeqName{f}$ is a strict
\nobreaks
  {
    {($\catName{E}^1$-$\catName{E}^2$-)},
    {$\seqname{E}^1$-$\seqname{E}^2$},
    {(-$\catName{C}^1$-$\catName{C}^2$)}
  }
m-atlas morphism of $\seqname{A}^1$ to $\seqname{A}^2$ in the
coordinate spaces $\seqname{C}^1$, $\seqname{C}^2$ then
$\morphSeqName{f}'$ is a strict
\nobreaks
  {
    {($\catName{E}'^1$-$\catName{E}'^2$-)},
    {$\seqname{E}^1$-$\seqname{E}^2$},
    {(-$\catName{C}'^1$-$\catName{C}'^2$)}
  }
m-atlas morphism of $\seqname{A}^1$ to $\seqname{A}^2$ in the
coordinate spaces $\seqname{C}^1$, $\seqname{C}'^2$.

\begin{proof}
There are four overlapping cases.
\begin{description}
\item[$\seqname{E}^i$ without category:]
$\morphName{f}_0$ is an m-morphism of $\seqname{E}^1$ to
$\seqname{E}^2$ by \pagecref{def:M-ATLmorphsemi}.

\item[$\seqname{E}^i$ with category:]
$\morphName{f}_0$ is an $\catName{E}^1$-$\catName{E}^2$
m-morphism of $\seqname{E}^1$ to $\seqname{E}^2$ by
\pagecref{def:M-ATLmorphsemi}.
$\catName{E}^i \subcat[full-] \catName{E}'^i$ by hypothesis.
$\morphName{f}_0$ is an $\catName{E}'^1$-$\catName{E}'^2$
m-morphism of $\seqname{E}^1$ to $\seqname{E}^2$ by
\pagecref{lem:modMmorph}.

\item[$\seqname{C}^i$ without category:]
$\morphName{f}_1$ is an m-morphism of $\seqname{C}^1$ to
$\seqname{C}^2$ by \pagecref{def:M-ATLmorphsemi}.
$\seqname{C}^2 \submod \seqname{C}'^2$ by hypothesis.
$\morphName{f}_1$ is an m-morphism of $\seqname{C}^1$ to
$\seqname{C}'^2$ by \pagecref{lem:modMmorph}.

\item[$\seqname{C}^i$ with category:]
$\morphName{f}_0$ is an $\catName{C}^1$-$\catName{C}^2$
m-morphism of $\seqname{C}^1$ to $\seqname{C}^2$ by
\pagecref{def:M-ATLmorphsemi}.
$\catName{C}^i \subcat[full-] \catName{C}'^i$ by hypothesis.
$\morphName{f}_0$ is a $\catName{C}'^1$-$\catName{C}'^2$
m-morphism of $\seqname{C}^1$ to $\seqname{C}^2$ by
\pagecref{lem:modMmorph}.
\end{description}
\end{proof}

\item
\label{M-ATLmorph:SstrictlocalisSstricttosup}
If $\morphSeqName{f}$ is a semi-strict
\nobreaks
  {
    {($\catName{E}^1$-$\catName{E}^2$-)},
    {$\seqname{E}^1$-$\seqname{E}^2$},
    {(-$\catName{C}^1$-$\catName{C}^2$)}
  }
m-atlas morphism of $\seqname{A}^1$ to $\seqname{A}^2$ in the
coordinate spaces $\seqname{C}^1$, $\seqname{C}^2$ and $\morphName{f}'_1$
is a local m-morphism of $\seqname{C}'^1$ to $\seqname{C}'^2$ (a
local $\catName{C}'^1$-$\catName{C}'^2$ m-morphism) of $\seqname{C}'^1$ to
$\seqname{C}'^2$) then $\morphSeqName{f}''$ is a semi-strict
\nobreaks
  {
    {($\catName{E}'^1$-$\catName{E}'^2$-)},
    {$\seqname{E}^1$-$\seqname{E}^2$},
    {(-$\catName{C}'^1$-$\catName{C}'^2$)}
  }
m-atlas morphism of $\seqname{A}^1$ to $\seqname{A}^2$ in the
coordinate spaces $\seqname{C}'^1$, $\seqname{C}'^2$.

\begin{proof}
There are four overlapping cases:

\begin{description}
\item[$\seqname{E}^i$ without categories:]
$\morphName{f}_0$ is a local m-morphism of $\seqname{E}^1$ to
$\seqname{E}^2$ by \pagecref{def:M-ATLmorphsemi}.

\item[$\seqname{E}^i$ with categories:]
$\morphName{f}_0$ is a local $\catName{E}^1$-$\catName{E}^2$
m-morphism of $\seqname{E}^1$ to $\seqname{E}^2$ by
\cref{def:M-ATLmorphsemi}.
$\catName{E}^i \subcat[full-] \catName{E}'^i$ by hypothesis.
$\morphName{f}_0$ is a local $\catName{E}'^1$-$\catName{E}'^2$
m-morphism of $\seqname{E}^1$ to $\seqname{E}^2$ by
\pagecref{lem:modMmorph}.

\item[$\seqname{C}'^i$ without categories:]
$\morphName{f}'_1$ is a local m-morphism of $\seqname{C}'^1$ to
$\seqname{C}'^2$ by hypothesis.

\item[$\seqname{C}'^i$ with categories:]
$\morphName{f}'_1$ is a local $\catName{C}'^1$-$\catName{C}'^2$
m-morphism of $\seqname{C}^1$ to $\seqname{C}^2$ by hypothesis.
\end{description}
\end{proof}

\item
\label{M-ATLmorph:strictmorphisstricttosup}
If $\morphSeqName{f}$ is a strict
\nobreaks
  {
    {($\catName{E}^1$-$\catName{E}^2$-)},
    {$\seqname{E}^1$-$\seqname{E}^2$},
    {(-$\catName{C}^1$-$\catName{C}^2$)}
  }
m-atlas morphism of $\seqname{A}^1$ to $\seqname{A}^2$ in the
coordinate spaces $\seqname{C}^1$, $\seqname{C}^2$ and $\morphName{f}'_1$
is an m-morphism of $\seqname{C}'^1$ to $\seqname{C}'^2$ (a
$\catName{C}'^1$-$\catName{C}'^2$ m-morphism) of $\seqname{C}'^1$ to
$\seqname{C}'^2$) then $\morphSeqName{f}''$ is a strict
\nobreaks
  {
    {($\catName{E}'^1$-$\catName{E}'^2$-)},
    {$\seqname{E}^1$-$\seqname{E}^2$},
    {(-$\catName{C}'^1$-$\catName{C}'^2$)}
  }
m-atlas morphism of $\seqname{A}^1$ to $\seqname{A}^2$ in the
coordinate spaces $\seqname{C}'^1$, $\seqname{C}'^2$.

\begin{proof}
There are four overlapping cases:

\begin{description}
\item[$\seqname{E}^i$ without categories:]
$\morphName{f}_0$ is an m-morphism of $\seqname{E}^1$ to
$\seqname{E}^2$ by \pagecref{def:M-ATLmorph}.

\item[$\seqname{E}^i$ with categories:]
$\morphName{f}_0$ is an $\catName{E}^1$-$\catName{E}^2$
m-morphism of $\seqname{E}^1$ to $\seqname{E}^2$ by
\cref{def:M-ATLmorph}.
$\catName{E}^i \subcat[full-] \catName{E}'^i$ by hypothesis.
$\morphName{f}_0$ is an $\catName{E}'^1$-$\catName{E}'^2$
m-morphism of $\seqname{E}^1$ to $\seqname{E}^2$ by
\pagecref{lem:modMmorph}.

\item[$\seqname{C}'^i$ without categories:]
$\morphName{f}'_1$ is an m-morphism of $\seqname{C}'^1$ to
$\seqname{C}'^2$ by hypothesis.

\item[$\seqname{C}'^i$ with categories:]
$\morphName{f}'_1$ is a $\catName{C}'^1$-$\catName{C}'^2$
m-morphism of $\seqname{C}^1$ to $\seqname{C}^2$ by hypothesis.
\end{description}
\end{proof}
\end{enumerate}
\end{lemma}

\begin{corollary}[M-atlas morphisms]
\label{cor:M-ATLmorph}
Let
\begin{enumerate}
\item
$\catName{C}^i$, $\catName{C}'^i$, $i=1,2$, be model categories

\item
$\catName{C}^i \subcat[full-] \catName{C}'^i$

\item
$\seqname{C}^i \objin \catName{C}^i$

\item
$\seqname{C}'^i \objin \catName{C}'^i$

\item
$\seqname{C}^i \submod \seqname{C}'^i$

\item
$\morphName{f}'_1 \maps \seqname{C}'^1 \to \seqname{C}'^2$ be a model
function such that
$\morphName{f}'_1[\seqname{C}^1] \subseteq \seqname{C}^2$ and
$
  \morphName{f}_1 \defeq
  \morphName{f}'_1 \restrictto_{\seqname{C}^1,\seqname{C}^2}
$
is a model function

\item
$E^i$ be a topological space

\item
$\seqname{A}^i$ be an m-atlas of $E^i$ in the coordinate space
$\seqname{C}^i$

\item
$
  \morphSeqName{f} \defeq
  (
    \morphName{f}_0 \maps E^1 \to E^2,
    \morphName{f}_1 \maps \seqname{C}^1 \to \seqname{C}^2
  )
$
be a (semi-strict, strict)
\nobreaks
  {
    {($\catName{E}^1$-$\catName{E}^2$-)},
    {$E^1$-$E^2$},
    {(-$\catName{C}^1$-$\catName{C}^2$)}
  }
m-atlas morphism of $\seqname{A}^1$ to $\seqname{A}^2$ in the
coordinate spaces $\seqname{C}^1$, $\seqname{C}^2$
\end{enumerate}

If $\morphSeqName{f}$ is a semi-strict (strict)
$E^1$-$E^2$-$\catName{C}^1$-$\catName{C}^2$ m-atlas morphism of
$\seqname{A}^1$ to $\seqname{A}^2$ in the coordinate spaces
$\seqname{C}^1$, $\seqname{C}^2$ then $\morphSeqName{f}$ is a semi-strict
(strict) $E^1$-$E^2$-$\catName{C}'^1$-$\catName{C}'^2$ m-atlas
morphism of $\seqname{A}^1$ to $\seqname{A}^2$ in the coordinate spaces
$\seqname{C}^1$, $\seqname{C}^2$.

\begin{proof}
The result follows from \pagecref{def:M-ATLmorphabbrev}.
\end{proof}
\end{corollary}

\begin{lemma}[Equivalence of m-atlas morphisms]
\label{lem:M-ATLmorphequiv}
Let
\begin{enumerate}
\item
$\catName{E}^i$, $\catName{C}^i$, $i=1,2$, be model categories
\item
$\seqname{E}^i \objin \catName{E}^i$,
$\seqname{C}^i \objin \catName{C}^i$,
\item
$\seqname{A}^i$ be an m-atlas of
$\seqname{E}^i$ in the coordinate space $\seqname{C}^i$,
\item
$
  \morphSeqName{f} \defeq
  (
    \morphName{f}_0 \maps \seqname{E}^1 \to \seqname{E}^2,
    \morphName{f}_1 \maps \seqname{C}^1 \to \seqname{C}^2
  )
$ and
$
  \morphSeqName{g} \defeq
  (
    \morphName{g}_0 \maps \seqname{E}^1 \to \seqname{E}^2,
    \morphName{g}_1 \maps \seqname{C}^1 \to \seqname{C}^2
  )
$
be equivalent m-atlas morphisms of $\seqname{A}^1$ to
$\seqname{A}^2$ in the coordinate spaces $\seqname{C}^1$,
$\seqname{C}^2$.
\end{enumerate}

If $\morphSeqName{f}$ is a (coherent) semi-strict
\nobreaks
  {
    {($\catName{E}^1$-$\catName{E}^2$-)},
    {$\seqname{E}^1$-$\seqname{E}^2$},
    {(-$\catName{C}^1$-$\catName{C}^2$)}
  }
m-atlas morphism of $\seqname{A}^1$ to $\seqname{A}^2$ in the
coordinate spaces $\seqname{C}^1$, $\seqname{C}^2$, then
$\morphSeqName{g}$ is a (coherent) semi-strict
\nobreaks
  {
    {($\catName{E}^1$-$\catName{E}^2$-)},
    {$\seqname{E}^1$-$\seqname{E}^2$},
    {(-$\catName{C}^1$-$\catName{C}^2$)}
  }
m-atlas morphism of $\seqname{A}^1$ to $\seqname{A}^2$ in the
coordinate spaces $\seqname{C}^1$, $\seqname{C}^2$.

\begin{proof}
Let $v^1 \in \seqname{C}^1$ be an arbitrary point. Let
$(U^1,V^1,\phi^1) \in \seqname{A}^1$ be an arbitrary chart such that
$v^1 \in V^1$, $u^1 \defeq \phi^{1-1}(v^1)$,
$(U^2,V^2,\phi^2) \in \seqname{A}^2$ be an arbitrary chart at
$u^2 \defeq \morphName{f}_0(u^1)$, $v^2 \defeq \phi^2(u^2)$ and
$I \defeq U^1 \cap \morphName{f}_0^{-1}[U^2]$, as shown in
\fullcref{fig:AtlE1}.

\begin{figure}[ht]
\[ \bfig
\node v1f(0,0)[v^1]
\node V1f(1000,0)[V^1]
\node C2f(1500,0)[C^2]
\node V2f(2000,0)[V^2]
\node U1f(1000,-500)[U^1]
\node u1(0,-1000)[u^1]
\node Uint(1000,-1000)[{I=U^1 \cap \morphName{f}^{-1}_0[U^2]}]
\node U2(2000,-1000)[U^2]
\node U1g(1000,-1500)[U^1]
\node v1g(0,-2000)[v^1]
\node V1g(1000,-2000)[V^1]
\node C2g(1500,-2000)[C^2]
\node V2g(2000,-2000)[V^2]
%
% y=0
\arrow |a|/^{ (}->/[v1f`V1f;\morphName{i}]
\arrow |a|[V1f`C2f;\morphName{f}_1]
%
% y=-500
\arrow |r|[U1f`V1f;\phi^1]
\arrow |l|//[U1f`V1f;\iso]
\arrow |r|[U1f`U2;\morphName{f}_0]
\arrow |r|[U2`V2f;\phi^2]
\arrow |l|//[U2`V2f;\iso]
%
% y=-1000
\arrow |r|[u1`v1f;\phi^1]
\arrow |a|/^{ (}->/[u1`U1f;\morphName{i}]
\arrow |a|/^{ (}->/[u1`Uint;\morphName{i}]
\arrow |a|/^{ (}->/[u1`U1g;\morphName{i}]
\arrow |r|[u1`v1g;\phi^1]
\arrow |r|/_{ (}->/[Uint`U1f;\morphName{i}]
\arrow |a|[Uint`U2;\morphName{f}_0]
\arrow |r|/^{ (}->/[Uint`U1g;\morphName{i}]
%
% y=-1500
\arrow |a|[U1g`U2;\morphName{f}_0]
\arrow |r|[U1g`V1g;\phi^1]
\arrow |l|//[U1g`V1g;\iso]
\arrow |r|[U2`V2g;\phi^2]
\arrow |l|//[U2`V2g;\iso]
%
% y=-2000
\arrow |a|/^{ (}->/[v1g`V1g;\morphName{i}]
\arrow |a|[V1g`C2g;\morphName{g}_1]
\efig \]
\caption{Uncompleted equivalent m-atlas morphisms}
\label{fig:AtlE1}
\end{figure}

Let
$
  ({^f}U'^1, {^f}V'^1, \phi^1 \maps {^f}U'^1 \toiso {^f}V'^1)
  \in \seqname{A}^1
$
be a subchart of $(U^1,V^1,\phi^1)$ at $u^1$,
${^f}U'^2 \subseteq U^2$ a model neighborhood of $\seqname{E}^2$ for
$u^2$ and \\*
$
  ({^f}U'^2 ,{^f}\hat{V}'^2, {^f}\phi'^2 \maps {^f}U'^2 \toiso {^f}\hat{V}'^2)
  \in \seqname{A}^2
$
be a chart at $u^2$  as in \pagecref{def:M-ATLmorph}, i.e.,

\begin{equation}
\morphName{f}_0[U'^1] \subseteq {^f}U'^2
\end{equation}

\begin{equation}
{^f}\phi'^2 \morphCompose \morphName{f}_0 \maps {^f}U'^1 \to {^f}\hat{V}'^2 =
\morphName{f}_1 \morphCompose \phi^1 \maps {^f}U'^1 \to {^f}\hat{V}'^2
\end{equation}

${^f}\hat{V}'^2$ is a model neighborhood of $\seqname{C}^2$ by
\pagecref{def:M-chart}.

$
  \morphName{f}_1 \maps {^f}V'^1 \to {^f}\hat{V}'^2 =
  {^f}\phi'^2 \morphCompose \morphName{f}_0 \morphCompose \phi^{1-1}
$

Let
$
  ({^g}U'^1, {^g}V'^1, {^g}\phi^1 \maps {^g}U'^1 \toiso {^g}V'^1)
  \in \seqname{A}^1
$
be a subchart of $(U^1,V^1,\phi^1)$ at $u^1$,
${^g}U'^2 \subseteq U^2$ a model neighborhood of $\seqname{E}^2$ for
$u^2$ and \\*
$
  ({^g}U'^2 ,{^g}\hat{V}'^2, {^g}\phi'^2 \maps {^g}U'^2 \toiso {^g}\hat{V}'^2)
  \in \seqname{A}^2
$
be a chart at $u^2$  as in \pagecref{def:M-ATLmorph}, i.e.,

\begin{equation}
\morphName{g}_0[U'^1] \subseteq {^g}U'^2
\end{equation}

\begin{equation}
\begin{split}
  {^g}\phi'^2 \morphCompose \morphName{f}_0 \maps {^g}U'^1 \to {^g}\hat{V}'^2
  & =
  {^g}\phi'^2 \morphCompose \morphName{g}_0 \maps {^g}U'^1 \to {^g}\hat{V}'^2
\\*
  & =
  \morphName{g}_1 \morphCompose \phi^1 \maps {^g}U'^1 \to {^g}\hat{V}'^2
\end{split}
\end{equation}

${^g}\hat{V}'^2$ is a model neighborhood of
$\seqname{C}^2$ by \cref{def:M-chart}.

\nobreaks
  {{
    $U'^1 \defeq {^f}U'^1 \cap {^g}U'^1$
  }}
is a model neighborhood of $\seqname{E}^1$,
\nobreaks
  {{
    $V'^1 \defeq {^f}U'^1 \cap {^g}U'^1$
  }}
is a model neighborhood of $\seqname{C}^1$,
\nobreaks
  {{
    $U'^2 \defeq {^f}U'^2 \cap {^g}U'^2$
  }}
is a model neighborhood of $\seqname{E}^2$ and
\nobreaks
  {{
    $V'^2 \defeq \hat{^{f}V'^2} \cap \hat{^{g}V'^2}$
  }}
is a model neighborhood of $\seqname{C}^2$ by
{
  \showlabelsinline
  \cref{mod:closed}
}
of \pagecref{def:model}.
$\morphName{f}_1[V'^1] \subseteq V'^2$ and
$\morphName{g}_1[V'^1] \subseteq V'^2$.

$
  \morphName{f}_0 \restrictto{U'^1,U'^2} =
  \morphName{g}_0 \restrictto{U'^1,U'^2}
$
is a morphism of $\seqname{E}^2$ by \pagecref{lem:modMmorph}.
$
  \morphName{f}_1 \restrictto{V'^1,V'^2} =
$
is a morphism of $\seqname{C}^2$ by \cref{lem:modMmorph}.

${^g}\phi'^2 \morphCompose {^f}\phi'^{2-} \maps V'^2 \to V'^2$ is an
isomorphism of $\seqname{C}^2$ by \fullcref{def:M-compat}.

\begin{equation}
\begin{split}
  \morphName{g}_1
   & =
   {^g}\phi'^2 \morphCompose \morphName{f}_0 \morphCompose \phi^{1-1}
\\*
  & =
    {^g}\phi'^2 \morphCompose \morphName{f}_0 \morphCompose \phi^{1-1}
\\*
  & =
    {^g}\phi'^2 \morphCompose {^f}\phi'^{2-1} \morphCompose
    {^f}\phi'^2 \morphCompose \morphName{f}_0 \morphCompose \phi^{1-1}
\\*
  & =
    {^g}\phi'^2 \morphCompose {^f}\phi'^{2-1} \morphCompose
    \morphName{f}_1
\end{split}
\end{equation}

There are four overlapping cases.
\begin{description}
\item[$\seqname{E}^i$ without category:]
$\morphName{f}_0 = \morphName{g}_0$ by \pagecref{def:M-ATLmorphequiv}.
$\morphName{f}_0$ is a local m-morphism of $\seqname{E}^1$ to
$\seqname{E}^2$ by \pagecref{def:M-ATLmorphsemi}.

If $\morphSeqName{f}$ is coherent then
$\seqname{E}^1 \submod[strict-] \seqname{E}^2$ by
\pagecref{def:M-ATLmorphsemi}.

\item[$\seqname{C}^i$ without category:]
$\morphName{f}_1$ is a local m-morphism of $\seqname{C}^1$ to
$\seqname{C}^2$. \\*
$
  \phi^2 \morphCompose \morphName{f}_0 \morphCompose \phi^{1-1} \maps
  \phi^1[I] \to V^2
$
is a local m-morphism of $\phi^1[I]$ to $\seqname{C}^2$ by
\pagecref{def:M-ATLmorphsemi}.

Let $\hat{V}'^2 \subseteq \hat{^{f}V'^2} \cap \hat{^{g}V'^2}$ be a model
neighborhod of $\seqname{C}^2$ at $v^2$ such that \\*
$
  \morphName{f}_1 \maps V'^1 \defeq \morphName{f}_1^{-1}[\hat{V}'^2)
  \to \hat{V}'^2
$
is a morphism of $\seqname{C}^2$.

The following, shown in
{
  \showlabelsinline
  \crefrange{fig:AtlE2}{fig:AtlE3},
}
are model
neighborhoods of $\seqname{C}^2$ by \pagecref{def:modfunc}:
\begin{align*}
V'^2 & \defeq \hat{f}^{-1}[\hat{V}'^2]
\\*
U'^2 & \defeq \phi^{2-1}[V'^2]
\\*
U'^1 & \defeq \morphName{f}^{-1}_0[U'^2]
\end{align*}

\begin{figure}[ht!]
\[ \bfig
\node v1f(1000,500)[v^1]
%
\node V1f(0,0)[V^1]
\node fVp1(1000,0)[{^f}V'^1]
\node fVhp2(2000,0)[\hat{^{f}V'^2}]
\node V2f(2500,0)[V^2]
\node u2(2000,0)[u^2]
%
\node U1f(0,-500)[U^1]
\node fUp1(1000,-500)[{^f}U'^1]
\node fUp2(2000,-500)[{^f}U'^2]
%
\node u1f(600,-1100)[u^1]
%
\node Uint(0,-1500)[{I=U^1 \cap \morphName{f}^{-1}_0[U^2]}]
\node Up1(1000,-1500)[U'^1]
\node Up2(2000,-1500)[U'^2]
\node U2(2500,-1500)[U^2]
%
\node U1g(0,-2000)[U^1]
\node gUp1(1000,-2000)[^{g}U'^1]
\node gUp2(2000,-2000)[^{g}U'^2]
%
\node V1g(0,-2500)[V^1]
\node gVp1(1000,-2500)[^{g}V'^1]
\node gVhp2(2000,-2500)[\hat{^{g}V'^2}]
\node V2g(2500,-2500)[V^2]
%
\node v1g(1000,-3000)[v^1]
%
% y=500
\arrow |r|/^{ (}->/[v1f`fVp1;\morphName{i}]
%
% y=0
\arrow |a|/_{ (}->/[fVp1`V1f;\morphName{i}]
\arrow |a|[fVp1`fVhp2;\morphName{f}_1]
\arrow |a|[fVhp2`fVhp2;\hat{\morphName{f}}]
\arrow |b|[fVhp2`fVhp2;\iso]
\arrow |a|/^{ (}->/[fVhp2`V2f;\morphName{i}]
%
% y=-500
\arrow |r|[U1f`V1f;\phi^1]
\arrow |l|//[U1f`V1f;\iso]
\arrow |r|[fUp1`fVp1;\phi^1]
\arrow |l|//[fUp1`fVp1;\iso]
\arrow |a|/_{ (}->/[fUp1`U1f;\morphName{i}]
\arrow |a|[fUp1`fUp2;\morphName{f}_0]
\arrow |a|/_{ (}->/[fUp1`Uint;\morphName{i}]
\arrow |r|[fUp2`fVhp2;\phi^2]
\arrow |l|//[fUp2`fVhp2;\iso]
\arrow |a|/^{ (}->/[fUp2`U2;\morphName{i}]
%
% y=-1000
\arrow |a|/_{ (}->/[u1f`Up1;\morphName{i}]
%
% y=-1500
\arrow |a|/>->/[Uint`U1f;\morphName{i}]
\arrow |a|/>->/[Uint`U1g;\morphName{i}]
\arrow |a|/_{ (}->/[Up1`fUp1;\morphName{i}]
\arrow |a|/_{ (}->/[Up1`Uint;\morphName{i}]
\arrow |r|/^{ (}->/[Up1`gUp1;\morphName{i}]
\arrow |a|[Up1`Up2;\morphName{f}_0]
\arrow |a|/^{ (}->/[Up2`U2;\morphName{i}]
\arrow |a|/^{ (}->/[Up2`fUp2;\morphName{i}]
\arrow |r|/^{ (}->/[gUp2`U2;\morphName{i}]
\arrow |r|/^{ (}->/[Up2`gUp2;\morphName{i}]
%
% y=-2000
\arrow |r|[U1g`V1g;\phi^ 1]
\arrow |a|/_{ (}->/[gUp1`Uint;\morphName{i}]
\arrow |a|/_{ (}->/[gUp1`U1g;\morphName{i}]
\arrow |a|[gUp1`gUp2;\morphName{g}_0]
\arrow |r|[gUp1`gVp1;\phi^1]
\arrow |l|//[gUp1`gVp1;\iso]
\arrow |a|/^{ (}->/[gUp2`U2;\morphName{i}]
\arrow |r|[gUp2`gVhp2;\phi^2]
\arrow |l|//[gUp2`gVhp2;\iso]
\arrow |r|[U2`V2f;\phi^2]
\arrow |r|[U2`V2g;\phi^2]
\arrow |l|//[U2`V2g;\iso]
%
% y=-2500
\arrow |a|/_{ (}->/[gVp1`V1g;\morphName{i}]
\arrow |a|[gVp1`gVhp2;\morphName{g}_1]
\arrow |a|[gVhp2`gVhp2;\hat{\morphName{g}}]
\arrow |b|[gVhp2`gVhp2;\iso]
\arrow |a|/^{ (}->/[gVhp2`V2g;\morphName{i}]
%
% y=-3000
\arrow |r|/_{ (}->/[v1g`gVp1;\morphName{i}]
\efig \]
\caption{Completed equivalent m-atlas morphisms}
\label{fig:AtlE2}
\end{figure}

\begin{figure}[ht!]
\[ \bfig
% y=0
\node Up1f(1000,0)[U'^1]
\node Up2f(2000,0)[U'^2]
%
% y=-500
\node u1f(0,-500)[u^1]
\node fUp1(1000,-500)[{^f}U'^1]
\node fUp2(2000,-500)[{^f}U'^2]
%
% y=-1000
\node fVp1(1000,-1000)[{^f}V'^1]
\node fVhp2(2000,-1000)[\hat{^{f}V'^2}]
%
% y=-1500
\node v1(0,-1500)[v^1]
\node Vp1(1000,-1500)[V'^1]
\node Vhp2(1500,-1500)[\hat{V}'^2]
\node Vp2(2000,-1500)[V'^2]
%
% y=-2000
\node gVp1(1000,-2000)[^{g}V'^1]
\node gVhp2(2000,-2000)[\hat{^{g}V'^2}]
%
% y=-2500
\node u1g(0,-2500)[u^1]
\node gUp1(1000,-2500)[{^g}U'^1]
\node gUp2(2000,-2500)[{^g}U'^2]
%
% y=-3000
\node Up1g(1000,-3000)[U'^1]
\node Up2g(2000,-3000)[U'^2]
%
% y=0
\arrow |a|[Up1f`Up2f;\morphName{f}_0]
\arrow |a|/^{ (}->/[Up1f`fUp1;\morphName{i}]
\arrow |r|/^{ (}->/[Up2f`fUp2;\morphName{i}]
%
% y=-500
\arrow |a|/^{ (}->/[u1f`Up1f;\morphName{i}]
\arrow |a|/^{ (}->/[u1f`fUp1;\morphName{i}]
\arrow |a|[u1f`fVp1;\phi^1]
\arrow |b|//[u1f`fVp1;\iso]
\arrow |a|[fUp1`fUp2;\morphName{f}_0]
\arrow |r|[fUp1`fVp1;\phi^1]
\arrow |l|//[fUp1`fVp1;\iso]
\arrow |r|[fUp2`fVhp2;\phi^2]
\arrow |l|//[fUp2`fVhp2;\iso]
%
% y=-1000
\arrow |a|[fVp1`fVhp2;\morphName{f}_1]
%
% y=-1500
\arrow |a|/^{ (}->/[v1`fVp1;\morphName{i}]
\arrow |a|/^{ (}->/[v1`Vp1;\morphName{i}]
\arrow |a|/^{ (}->/[v1`gVp1;\morphName{i}]
\arrow |r|/_{ (}->/[Vp1`fVp1;\morphName{i}]
\arrow |a|[Vp1`fVhp2;\morphName{f}_1]
\arrow |a|[Vp1`gVhp2;\morphName{g}_1]
\arrow |r|/^{ (}->/[Vp1`gVp1;\morphName{i}]
\arrow |r|/_{ (}->/[Vp2`fVhp2;\morphName{i}]
\arrow |r|/^{ (}->/[Vp2`gVhp2;\morphName{i}]
%
% y=-2000
\arrow |a|[gVp1`gVhp2;\morphName{g}_1]
%
% y=-2500
\arrow |a|[u1g`gVp1;\phi^1]
\arrow |b|//[u1g`gVp1;\iso]
\arrow |a|/^{ (}->/[u1g`gUp1;\morphName{i}]
\arrow |a|/^{ (}->/[u1g`Up1g;\morphName{i}]
\arrow |r|[gUp1`gVp1;\phi^1]
\arrow |a|[gUp1`gUp2;\morphName{f}_0]
\arrow |r|[gUp2`gVhp2;\phi^2]
\arrow |l|//[gUp2`gVhp2;\iso]
%
% y=-3000
\arrow |a|[Up1g`Up2g;\morphName{f}_0]
\arrow |a|/_{ (}->/[Up1g`gUp1;\morphName{i}]
\arrow |a|/_{ (}->/[Up2g`gUp2;\morphName{i}]
\efig \]
\caption{Completed equivalent m-atlas morphisms}
\label{fig:AtlE3}
\end{figure}

Then
\begin{equation*}
\begin{array} {llll}
  \morphName{g}_1 \maps V'^1 \to \hat{V}'^2
  & =
  & \hat{\morphName{g}} \maps V'^2 \to \hat{V}'^2
  & \morphCompose
\\*
  &
  & \phi^2 \maps U'^2 \to V'^2
  & \morphCompose
\\*
  &
  & \morphName{f}_0 \maps U'^1 \to U'^2
  & \morphCompose
\\*
  &
  & \phi^{1-1} \maps V'^1 \to U'^1
\\*
  & =
  & \hat{\morphName{g}} \maps  V'^2 \to \hat{V}'^2
  & \morphCompose
\\*
  &
  & \hat{\morphName{f}}^{-1} \maps  \hat{V}'^2 \to  V'^2
  & \morphCompose
\\*
  &
  & \hat{\morphName{f}} \maps V'^2 \to \hat{V}'^2
  & \morphCompose
\\*
  &
  & \phi^2 \maps  U'^2 \to '^2
  & \morphCompose
\\*
  &
  & \morphName{f}_0 \maps U'^1 \to U'^2
  & \morphCompose
\\*
  &
  & \phi^{1-1} \maps V'^1 \to U'^1
\\*
  & =
  & \hat{\morphName{g}} \maps  V'^2 \to \hat{V}'^2
  & \morphCompose
\\*
  &
  & \hat{\morphName{f}}^{-1} \maps  \hat{V}'^2 \to  V'^2
  & \morphCompose
\\*
  &
  & \morphName{f}_1 \maps V'^1 \to \hat{V}'^2
\end{array}
\end{equation*}
is a composition of morphisms of $\seqname{C}^2$, hence a morphism.

If $\morphSeqName{f}$ is coherent then
$\seqname{C}^1 \submod[strict-] \seqname{C}^2$ by
\pagecref{def:M-ATLmorphsemi}.

\item[$\seqname{E}^i$ with category:]
$\morphName{f}_0 = \morphName{g}_0$ by \cref{def:M-ATLmorphequiv}.
$\morphName{f}_0$ is a local $\catName{E}^1$-$\catName{E}^2$
m-morphism of $\seqname{E}^1$ to $\seqname{E}^2$ by
\cref{def:M-ATLmorphsemi}.

If $\morphSeqName{f}$ is coherent then
$\catName{E}^1 \subcat[strict-] \catName{E}^2$ by
\pagecref{def:M-ATLmorphsemi}.

\item[$\seqname{C}^i$ with category:]
$\morphName{f}_1$ is a local $\catName{C}^1$-$\catName{C}^2$ m-morphism
of $\seqname{C}^1$ to $\seqname{C}^2$ and \\*
$
  \phi^2 \morphCompose \morphName{f}_0 \morphCompose \phi^{1-1} \maps
  \phi^1[I] \to V^2
$
is a local $\catName{C}^1$-$\catName{C}^2$ m-morphism of
$\Mod(\phi^1[I],\seqname{C}^2)$ to $\Mod(V^2,\seqname{C}^2)$ by
\pagecref{def:M-ATLmorphsemi}.

Let $U'^2 \subseteq {}^f\!U'^2 \cap ^{g}U'^2$ be a model neighborhod of
$\seqname{C}^2$ at $u^2$ such that \\*
$\morphName{f}_0 \maps U'^1 \defeq \morphName{f}_0^{-1}[U'^2] \to U'^2$ is
a morphism of $\catName{C}^2$.

Let $V'^2 \defeq \phi^2[U'^2]$,
$\hat{^{f}V'^2} \defeq \hat{\morphName{f}}[^{f}V'^2]$,
$\hat{^{g}V'^2} \defeq \hat{\morphName{g}}[^{g}V'^2]$,
$^{f}V'^1 \defeq \morphName{f}_1^{-1}[\hat{^{f}V'^2}]$ and
$^{g}V'^1 \defeq \morphName{g}_1^{-1}[\hat{^{g}V'^2}]$.
All of these are model neighborhoods by \pagecref{def:modfunc} and the
corresponding relative model subspaces of $\seqname{V}^i$ are objects of
$\catName{V}^i$ by
{
  \showlabelsinline
  \cref{modcat:rest}
}
of \pagecref{def:modcat}. The restrictions of morphisms to them are
morphisms by
{
  \showlabelsinline
  \cref{modcat:rest}
}
of \cref{def:modcat}.

Then
\begin{equation*}
\begin{array} {llll}
  \morphName{g}_1 \maps V'^1 \to \hat{V}'^2
  & =
  & \hat{\morphName{g}} \maps V'^2 \to \hat{V}'^2
  & \morphCompose
\\*
  &
  & \phi^2 \maps U'^2 \to V'^2
  & \morphCompose
\\*
  &
  & \morphName{f}_0 \maps U'^1 \to U'^2
  & \morphCompose
\\*
  &
  & \phi^{1-1} \maps V'^1 \to U'^1
\\*
  & =
  & \hat{\morphName{g}} \maps  V'^2 \to \hat{V}'^2
  & \morphCompose
\\*
  &
  & \hat{\morphName{f}}^{-1} \maps  \hat{V}'^2 \to  V'^2
  & \morphCompose
\\*
  &
  & \hat{\morphName{f}} \maps V'^2 \to \hat{V}'^2
  & \morphCompose
\\*
  &
  & \phi^2 \maps  U'^2 \to '^2
  & \morphCompose
\\*
  &
  & \morphName{f}_0 \maps U'^1 \to U'^2
  & \morphCompose
\\*
  &
  & \phi^{1-1} \maps V'^1 \to U'^1
\\*
  & =
  & \hat{\morphName{g}} \maps  V'^2 \to \hat{V}'^2
  & \morphCompose
\\*
  &
  & \hat{\morphName{f}}^{-1} \maps  \hat{V}'^2 \to  V'^2
  & \morphCompose
\\*
  &
  & \morphName{f}_1 \maps V'^1 \to \hat{V}'^2
\end{array}
\end{equation*}
is a composition of morphisms of $\catName{C}^2$, hence a morphism.

If $\morphSeqName{f}$ is coherent then
$\catName{C}^1 \subcat[strict-] \catName{C}^2$ by
\pagecref{def:M-ATLmorphsemi}.
\end{description}
\end{proof}
\end{lemma}

\begin{corollary}[Equivalence of m-atlas morphisms]
\label{cor:M-ATLmorphequiv}
Let
\begin{enumerate}
\item
$\catName{E}^i$, $\catName{C}^i$, $i=1,2$, be model categories
\item
$\seqname{E}^i \objin \catName{E}^i$
\item
$\seqname{C}^i \objin \catName{C}^i$
\item
$\seqname{A}^i$ be an m-atlas of $\seqname{E}^i$ in the coordinate space
$\seqname{C}^i$
\item
$\seqname{A}^1$ be full
\item
$
  \morphSeqName{f} \defeq
  (
    \morphName{f}_0 \maps \seqname{E}^1 \to \seqname{E}^2,
    \morphName{f}_1 \maps \seqname{C}^1 \to \seqname{C}^2
  )
$
be a coherent semi-strict
\nobreaks
  {
    {($\catName{E}^1$-$\catName{E}^2$-)},
    {$\seqname{E}^1$-$\seqname{E}^2$},
    {(-$\catName{C}^1$-$\catName{C}^2$)}
  }
m-atlas morphism of $\seqname{A}^1$ to $\seqname{A}^2$ in the
coordinate spaces $\seqname{C}^1$, $\seqname{C}^2$.
\item
$
  \morphSeqName{g} \defeq
  (
    \morphName{g}_0 \maps \seqname{E}^1 \to \seqname{E}^2,
    \morphName{g}_1 \maps \seqname{C}^1 \to \seqname{C}^2
  )
$
be an m-atlas morphisms of $\seqname{A}^1$ to $\seqname{A}^2$ in
the coordinate spaces $\seqname{C}^1$, $\seqname{C}^2$ equivalent to
$\morphSeqName{f}$
\end{enumerate}

Then $\morphSeqName{g}$ is a strict
\nobreaks
  {
    {($\catName{E}^1$-$\catName{E}^2$-)},
    {$\seqname{E}^1$-$\seqname{E}^2$},
    {(-$\catName{C}^1$-$\catName{C}^2$)}
  }
m-atlas morphism of $\seqname{A}^1$ to $\seqname{A}^2$ in the
coordinate spaces $\seqname{C}^1$, $\seqname{C}^2$.

\begin{proof}
$\morphSeqName{g}$ is a coherent semi-strict
\nobreaks
  {
    {($\catName{E}^1$-$\catName{E}^2$-)},
    {$\seqname{E}^1$-$\seqname{E}^2$},
    {(-$\catName{C}^1$-$\catName{C}^2$)}
  }
m-atlas morphism of $\seqname{A}^1$ to $\seqname{A}^2$ in the
coordinate spaces $\seqname{C}^1$, $\seqname{C}^2$ by
\pagecref{lem:M-ATLmorphequiv}.
The result follows from \cref{M-ATLmorph:strictSstrictisstrict} of
{
  \showlabelsinline
  \pagecref{lem:M-ATLmorph}.
}
\end{proof}
\end{corollary}

\begin{lemma}[Composition of m-atlas morphisms]
\label{lem:M-ATLmorphcomp}
Let
\begin{enumerate}
\item
$\catName{E}^i$, $\catName{C}^i$, $i=1,2,3$, be a model category,
$\seqname{E}^i \objin \catName{E}^i$,
$\seqname{C}^i \objin \catName{C}^i$,
\item
$\seqname{A}^i$ an m-atlas of
$\seqname{E}^i$ in the coordinate space $\seqname{C}^i$
\item
$
  \morphSeqName{f}^i \defeq
  (
    \morphName{f}^i_0 \maps \seqname{E}^i \to \seqname{E}^{i+1},
    \morphName{f}^i_1 \maps \seqname{C}^i \to \seqname{C}^{i+1}
  )
$
and
$
  \morphSeqName{g}^i \defeq
  (
    \morphName{g}^i_0 \maps \seqname{E}^i \to \seqname{E}^{i+1},
    \morphName{g}^i_1 \maps \seqname{C}^i \to \seqname{C}^{i+1}
  )
$
$i=1,2$,
be equivalent (semi-strict, strict)
\nobreaks
  {
    {($\catName{E}^i$-$\catName{E}^{i+1}$-)},
    {$\seqname{E}^i$-$\seqname{E}^{i+1}$},
    {(-$\catName{C}^i$-$\catName{C}^{i+1}$)}
  }
m-atlas morphisms of $\seqname{A}^i$ to $\seqname{A}^{i+1}$ in
the coordinate spaces $\seqname{C}^i$, $\seqname{C}^{i+1}$.
\end{enumerate}

Then:
\begin{enumerate}
\item
\label{M-morphcomp:morphcompmorph}
$\morphSeqName{f}^2 \morphCompose[()] \morphSeqName{f}^1$ is an
$\seqname{E}^1$-$\seqname{E}^3$ m-atlas morphism of $\seqname{A}^1$ to
$\seqname{A}^3$ in the coordinate spaces $\seqname{C}^1$,
$\seqname{C}^3$.

\begin{proof}
$\morphName{f}^2_0 \morphCompose \morphName{f}^1_0$ and
$\morphName{f}^2_1 \morphCompose \morphName{f}^1_1$ are model functions by
\pagecref{lem:modcomp}.

Let $(U^i, V^i, \phi^i) \in \seqname{A}^i$, $i=1,3$, be a chart,
$
  I \defeq
  U^1 \cap (\morphName{f}^2_0 \morphCompose \morphName{f}^1_0)^{-1}[U^3]
  \neq \emptyset
$.

Let $u^1 \in I$, $u^2 \defeq \morphName{f}^1_0(u^1)$,
$u^3 \defeq \morphName{f}^2_0(u^2)$ and
$(U^2, V^2, \phi^2) \in \seqname{A}^2$ be a chart at $u^2$.
Since $\morphSeqName{f}^1$ is a morphism by hypothesis, there exists a
subchart \\*
\nobreaks
  {{
    $
      (U'^1, V'^1, \phi^1 \maps U'^1 \toiso V'^1)
      \in \seqname{A}^1
    $
  }}
at $u^1$, a model neighborhood $U'^2 \subseteq U^2$ of $u^2$ and a chart
$
  (U'^2 ,\hat{V}'^2, \phi'^2 \maps U'^2 \toiso \hat{V}'^2)
  \in \seqname{A}^2
$
at $u^2$ such that
\nobreaks
  {{
    $\morphName{f}^1_0[U'^1] \subseteq U'^2$
  }}
and diagram \eqref{eq:AtlM1} in \fullcref{def:M-ATLmorph} on \cpageref{eq:AtlM1} is
commutative, as shown in \pagecref{fig:AtlM1}.

\begin{figure}[ht]

\[ \bfig
\node x(0,0)[u^1]
\node up1(500,0)[U'^1]
\node upp1(500,-1000)[U''^1]
\node vpp1(500,-2000)[V''^1]
\node vp1(500,-2500)[V'^1]
\node up2(1500,0)[U'^2]
\node uint(1500,-500)[{I^{2,3} = U'^2 \cap \morphName{f}^{-1}_0[U^3]}]
\node upp2(1500,-1000)[U''^2]
\node vpp2(1500,-2000)[V''^2]
\node vhp2(1500,-2500)[\hat{V}'^2]
\node E3(2500,0)[E^3]
\node u3(2500,-500)[U^3]
\node up3(2500,-1000)[U'^3]
%\node vp3(2500,-1500)[V'^3]
\node vhp3(2500,-2000)[\hat{V}'^3]
\node v3(2500,-1500)[V^3]
\node c3(2500,-2500)[C^3]
\arrow |a|/^{ (}->/[x`up1;\morphName{i}]
\arrow |a|[up1`up2;\morphName{f}^1_0]
\arrow  |r|/>->>/[up1`vp1;\phi^1]
%\arrow |r|/^{ (}->/[upp1`up1;\morphName{i}]
%\arrow |a|[upp1`upp2;\morphName{f}^1_0]
%\arrow |r|/>->>/[upp1`vpp1;\phi^1]
%\arrow |l|//[upp1`vpp1;\iso]
%\arrow |a|[vpp1`vpp2;\morphName{f}^1_1]
%\arrow |r|/^{ (}->/[vpp1`vp1;\morphName{i}]
\arrow |a|[vp1`vhp2;\morphName{f}^1_1]
\arrow |a|[up2`E3;\morphName{f}^2_0]
\arrow |r|/^{ (}->/[uint`up2;\morphName{i}]
\arrow |a|[uint`u3;\morphName{f}^2_0]
\arrow |a|[uint`vhp2;\phi'^2]
%\arrow |r|/^{ (}->/[upp2`uint;\morphName{i}]
%\arrow |a|[upp2`up3;\morphName{f}^2_0]
%\arrow |r|/>->>/[upp2`vpp2;\phi'^2]
%\arrow |l|//[upp2`vpp2;\iso]
%\arrow |a|[vpp2`vhp3;\morphName{f}^2_1]
%\arrow |a|/^{ (}->/[vpp2`vhp2;\morphName{i}]
\arrow |a|[vhp2`c3;\morphName{f}^2_1]
\arrow |r|/^{ (}->/[u3`E3;\morphName{i}]
\arrow |r|/^{ (}->/[u3`v3;\phi'^3]
%\arrow |r|/^{ (}->/[up3`u3;\morphName{i}]
%\arrow |r|/>->>/[up3`vhp3;\phi'^3]
%\arrow |l|//[up3`vhp3;\iso]
%\arrow |r|/>->>/[vp3`vhp3;\hat{\morphName{f}}]
%\arrow |l|//[vp3`vhp3;\iso]
%\arrow |r|/^{ (}->/[vhp3`c3;\morphName{i}]
\efig \]
\caption{Uncompleted composition of m-atlas morphisms}
\label{fig:morphcompM1}
\end{figure}

Since $\morphSeqName{f}^2$ is a morphism by hypothesis, there exists a
subchart \\
$
  (U''^2, V''^2, \phi^2 \maps U''^2 \toiso V''^2)
  \in \seqname{A}^2
$
at $u^2$, a model neighborhood $U'^3 \subseteq U^3$ of $u^3$ and a chart
$
  (U'^3 ,\hat{V}'^3, \phi'^3 \maps U'^3 \toiso \hat{V}'^3)
  \in \seqname{A}^3
$
at $u^2$ such that $\morphName{f}^2_0[U''^2] \subseteq U'^3$ and diagram
\eqref{eq:AtlM2} below is commutative, as shown in
\cref{fig:morphcompM1,fig:morphcompM2}.

\begin{equation}
\begin{array}{lll}
\label{eq:AtlM2}
D \defeq & \bigl \{
\\
  &&
  \morphName{i} \maps \set{u^2} \morphMono U''^2,
  \phi^2 \maps U''^2 \toiso V'^2,
  \morphName{f}_1 \maps  V''^2 \to \hat{V}'^3,
\\
  &&
  \morphName{f}_0 \maps  U'^2 \to U'^3,
  \phi'^3 \maps U'^3 \toiso \hat{V}'^3
\\
& \bigr \}
\end{array}
\end{equation}

\begin{figure}[ht]
\[ \bfig
\node x(0,0)[u^1]
\node up1(500,0)[U'^1]
\node upp1(500,-1000)[U''^1]
\node vpp1(500,-2000)[V''^1]
\node vp1(500,-2500)[V'^1]
\node up2(1500,0)[U'^2]
\node uint(1500,-500)[{I^{2,3} = U'^2 \cap \morphName{f}^{-1}_0[U^3]}]
\node upp2(1500,-1000)[U''^2]
\node vpp2(1500,-2000)[V''^2]
\node vhp2(1500,-2500)[\hat{V}'^2]
\node E3(2500,0)[E^3]
\node u3(2500,-500)[U^3]
\node up3(2500,-1000)[U'^3]
%\node vp3(2500,-1500)[V'^3]
\node vhp3(2500,-2000)[\hat{V}'^3]
\node v3(2500,-1500)[V^3]
\node c3(2500,-2500)[C^3]
\arrow |a|/^{ (}->/[x`up1;\morphName{i}]
\arrow |a|[up1`up2;\morphName{f}^1_0]
\arrow |r|/>->>/[up1`vp1;\phi^1]
\arrow |l|//[up1`vp1;\iso]
% \arrow |r|/^{ (}->/[upp1`up1;\morphName{i}]
% \arrow |a|[upp1`upp2;\morphName{f}^1_0]
% \arrow |r|/>->>/[upp1`vpp1;\phi^1]
% \arrow |l|//[upp1`vpp1;\iso]
% \arrow |a|[vpp1`vpp2;\morphName{f}^1_1]
% \arrow |r|/^{ (}->/[vpp1`vp1;\morphName{i}]
\arrow |a|[vp1`vhp2;\morphName{f}^1_1]
\arrow |a|[up2`E3;\morphName{f}^2_0]
\arrow |r|/^{ (}->/[uint`up2;\morphName{i}]
\arrow |a|[uint`u3;\morphName{f}^2_0]
\arrow |r|/^{ (}->/[upp2`uint;\morphName{i}]
\arrow |a|[upp2`up3;\morphName{f}^2_0]
\arrow |r|/>->>/[upp2`vpp2;\phi'^2]
\arrow |l|//[upp2`vpp2;\iso]
\arrow |a|[vpp2`vhp3;\morphName{f}^2_1]
\arrow |a|/^{ (}->/[vpp2`vhp2;\morphName{i}]
\arrow |a|[vhp2`c3;\morphName{f}^2_1]
\arrow |r|/^{ (}->/[u3`E3;\morphName{i}]
\arrow |r|/^{ (}->/[up3`u3;\morphName{i}]
\arrow |r|/>->>/[up3`vhp3;\phi'^3]
\arrow |l|//[up3`vhp3;\iso]
%\arrow |r|/>->>/[vp3`vhp3;\hat{\morphName{f}}]
%\arrow |l|//[vp3`vhp3;\iso]
\arrow |r|/^{ (}->/[vhp3`c3;\morphName{i}]
\efig \]
\caption{Partially completed composition of m-atlas morphisms}
\label{fig:morphcompM2}
\end{figure}

\begin{figure}[ht]
\[ \bfig
\node x(0,-1000)[u^1]
\node up1(500,0)[U'^1]
\node upp1(500,-1000)[U''^1]
\node vpp1(500,-2000)[V''^1]
\node vp1(500,-2500)[V'^1]
\node up2(1500,0)[U'^2]
\node uint(1500,-500)[{I^{2,3} = U'^2 \cap \morphName{f}^{-1}_0[U^3]}]
\node upp2(1500,-1000)[U''^2]
\node vpp2(1500,-2000)[V''^2]
\node vhp2(1500,-2500)[\hat{V}'^2]
\node E3(2500,0)[E^3]
\node u3(2500,-500)[U^3]
\node up3(2500,-1000)[U'^3]
%\node vp3(2500,-1500)[V'^3]
\node vhp3(2500,-2000)[\hat{V}'^3]
\node v3(2500,-1500)[V^3]
\node c3(2500,-2500)[C^3]
\arrow |a|/^{ (}->/[x`upp1;\morphName{i}]
\arrow |a|[up1`up2;\morphName{f}^1_0]
\arrow |r|/^{ (}->/[upp1`up1;\morphName{i}]
\arrow |a|[upp1`upp2;\morphName{f}^1_0]
\arrow |r|/>->>/[upp1`vpp1;\phi^1]
\arrow |l|//[upp1`vpp1;\iso]
\arrow |a|[vpp1`vpp2;\morphName{f}^1_1]
\arrow |r|/^{ (}->/[vpp1`vp1;\morphName{i}]
\arrow |a|[vp1`vhp2;\morphName{f}^1_1]
\arrow |a|[up2`E3;\morphName{f}^2_0]
\arrow |r|/^{ (}->/[uint`up2;\morphName{i}]
\arrow |a|[uint`u3;\morphName{f}^2_0]
\arrow |r|/^{ (}->/[upp2`uint;\morphName{i}]
\arrow |a|[upp2`up3;\morphName{f}^2_0]
\arrow |r|/>->>/[upp2`vpp2;\phi'^2]
\arrow |l|//[upp2`vpp2;\iso]
\arrow |a|[vpp2`vhp3;\morphName{f}^2_1]
\arrow |a|/^{ (}->/[vpp2`vhp2;\morphName{i}]
\arrow |a|[vhp2`c3;\morphName{f}^2_1]
\arrow |r|/^{ (}->/[u3`E3;\morphName{i}]
\arrow |r|/^{ (}->/[up3`u3;\morphName{i}]
\arrow |r|/>->>/[up3`vhp3;\phi'^3]
\arrow |l|//[up3`vhp3;\iso]
%\arrow |r|/>->>/[vp3`vhp3;\hat{\morphName{f}}]
%\arrow |l|//[vp3`vhp3;\iso]
\arrow |r|/^{ (}->/[vhp3`c3;\morphName{i}]
\efig \]
\caption{Completed composition of M-atlas morphisms}
\label{fig:morphcompM3}
\end{figure}

Define $U''^1 \defeq \morphName{f}^{1-1}_0[U''^2]$,
$
  V''^1 \defeq
  \morphName{f}^{1-1}_1[V''^2] =
  \phi^1[U''^1]
$,
\Cref{fig:morphcompM3} is commutative.
\end{proof}

\item
\label{M-morphcomp:semicompsemi}
If each $\morphName{f}^i$, $i=1,2$, is a semi-strict
\nobreaks
  {
    {($\catName{E}^1$-$\catName{E}^2$-)},
    {$\seqname{E}^1$-$\seqname{E}^2$},
    {(-$\catName{C}^1$-$\catName{C}^2$)}
  }
m-atlas morphism
then $\morphSeqName{f}^2 \morphCompose[()] \morphSeqName{f}^1$ is a semi-strict
\nobreaks
  {
    {($\catName{E}^1$-$\catName{E}^2$-)},
    {$\seqname{E}^1$-$\seqname{E}^2$},
    {(-$\catName{C}^1$-$\catName{C}^2$)}
  }
m-atlas morphism.

\begin{proof}
There are four overlapping cases.
\begin{description}
\item[$\seqname{E}^i$ without category:]
If each $\seqname{E}^i \submod \seqname{E}^{i+1}$ and each
$\morphSeqName{f}^i_0$ is a local m-morphism of $\seqname{E}^i$ to
$\seqname{E}^{i+1}$, then by \pagecref{lem:modMmorphcomp}
$\morphSeqName{f}^2_0 \morphCompose \morphSeqName{f}^1_0$ is a local
m-morphism of $\seqname{E}^1$ to $\seqname{E}^3$.

\item[$\seqname{C}^i$ without category:]
If each $\seqname{C}^i \submod \seqname{C}^{i+1}$ and each
$\morphSeqName{f}^i_1$ is a local m-morphism of $\seqname{C}^i$ to
$\seqname{C}^{i+1}$, then by \cref{lem:modMmorphcomp}
$\morphSeqName{f}^2_1 \morphCompose \morphSeqName{f}^1_1$ is a local
morphism of $\seqname{C}^1$ to $\seqname{C}^3$.
Let $(U^i, V^i, \phi^i \maps U^i \toiso V^i) \in \seqname{A}^i$, $i=1,3$,
with
$
  I \defeq
  U^1 \cap (\morphSeqName{f}^2_0 \morphCompose \morphSeqName{f}^1_0)^{-1}[U^3]
  \neq \emptyset
$.
For every $u^1 \in I$ let $u^2 \defeq \morphSeqName{f}^1_0(u^1)$,
$u^3 \defeq \morphSeqName{f}^2_0(u^2)$ and
$(U^2, V^2, \phi^2 \maps U^2 \toiso V^2) \in \seqname{A}^2$ be a chart
at $u^2$,
$
  I^1 \defeq
  U^1 \cap \morphSeqName{f}^{1-1}_0[U^2]
  \neq \emptyset
$
and
$
  I^2 \defeq
  U^2 \cap \morphSeqName{f}^{2-1}_0 [U^3]
  \neq \emptyset
$.
If
$
  \phi^2 \morphCompose \morphName{f}^1_0 \morphCompose \phi^{1-1} \maps
  \phi^1[I^1] \to V^2
$
is a local m-morphism of $\Mod(\phi^1[I^1],\seqname{C}^1)$
to $\Mod(V^2,\seqname{C}^2)$ and
$
  \phi^3 \morphCompose \morphName{f}^2_0 \morphCompose \phi^{2-1} \maps
  \phi^2[I^2] \to V^3
$
is a local m-morphism of $\Mod(\phi^1[I^2],\seqname{C}^1)$
to $\Mod(V^2,\seqname{C}^2)$ then by \cref{def:modlocalMmorph} there
exist model neighborhoods $V'^i \subseteq V^i$, $i \in [1,3]$, such that
$\morphName{f}^i[V'^i] \subseteq V'^{i+1}$ and \\*
$
  \phi^{i+1} \morphCompose \morphName{f}^i_0 \morphCompose \phi^{i-1} \maps
  V'^i \to V'^{i+1}
$
is an m-morphism of $\seqname{C}^{i+1}$, $i=1,2$.
Then
\begin{equation*}
\begin{array}{lll}
  \phi^3 \morphCompose \morphName{f}^2_0 \morphCompose \morphName{f}^1_0 \morphCompose
  \phi^{1-1} \maps V'^1 \to V'^3
  &
  =
  &
  \phi^3 \morphCompose \morphName{f}^2_0 \morphCompose \phi^{2-1} \maps
  V'^2 \to V'^3 \morphCompose
\\
  &&
  \phi^2 \morphCompose \morphName{f}^1_0 \morphCompose \phi^{1-1} \maps
  V'^1 \to V'^2
\end{array}
\end{equation*}
is an m-morphism of $\seqname{C}^3$. Since $u^1 \in I$ is arbitray,
$
  \phi^3 \morphCompose
  \morphName{f}^2_0 \morphCompose \morphName{f}^1_0 \morphCompose
  \phi^{1-1} \maps \phi^1[I] \to V^3
$
is a local m-morphism of $\Mod(\phi^1[I],\seqname{C}^3)$
to $\Mod(V^3,\seqname{C}^3)$.

\item[$\seqname{E}^i$ with category:]
If each $\catName{E}^i \subcat[full-] \catName{E}^{i+1}$ and each
$\morphSeqName{f}^i_0$ is a local $\catName{E}^i$-$\catName{E}^{i+1}$
m-morphism of $\seqname{E}^i$ to $\seqname{E}^{i+1}$, then by
\pagecref{lem:modMmorphcomp}
$\morphSeqName{f}^2_0 \morphCompose \morphSeqName{f}^1_0$ is a local
$\catName{E}^1$-$\catName{E}^3$ m-morphism of $\seqname{E}^1$ to
$\seqname{E}^3$.

\item[$\seqname{C}^i$ with category:]
If each $\catName{C}^i \subcat[full-] \catName{C}^{i+1}$ and each
$\morphSeqName{f}^i_1$ is a local $\catName{C}^i$-$\catName{C}^{i+1}$
m-morphism of $\seqname{C}^i$ to $\seqname{C}^{i+1}$, then by
\pagecref{lem:modMmorphcomp}
$\morphSeqName{f}^2_1 \morphCompose \morphSeqName{f}^1_1$ is a local
$\catName{C}^1$-$\catName{C}^3$ m-morphism of $\seqname{C}^1$ to
$\seqname{C}^3$.
\end{description}
\end{proof}

\item
\label{M-morphcomp:strictcompstrict}
If each $\morphName{f}^i$, $i=1,2$, is strict then
$\morphSeqName{f}^2 \morphCompose[()] \morphSeqName{f}^1$ is strict.

\begin{proof}
There are four overlapping cases.
\begin{description}
\item[$\seqname{E}^i$ without category:]
If each $\seqname{E}^i \submod \seqname{E}^{i+1}$ and each
$\morphSeqName{f}^i_0$ is an m-morphism of $\seqname{E}^i$ to
$\seqname{E}^{i+1}$, then by \pagecref{lem:modMmorphcomp}
$\morphSeqName{f}^2_0 \morphCompose \morphSeqName{f}^1_0$ is an
m-morphism of $\seqname{E}^1$ to $\seqname{E}^3$.

\item[$\seqname{C}^i$ without category:]
If each $\seqname{C}^i \submod \seqname{C}^{i+1}$ and each
$\morphSeqName{f}^i_1$ is an m-morphism of $\seqname{C}^i$ to
$\seqname{C}^{i+1}$, then by \pagecref{lem:modMmorphcomp}
$\morphSeqName{f}^2_1 \morphCompose \morphSeqName{f}^1_1$ is an
m-morphism of $\seqname{C}^1$ to $\seqname{C}^3$.

\item[$\seqname{E}^i$ with category:]
If each $\catName{E}^i \subcat[full-] \catName{E}^{i+1}$ and each
$\morphSeqName{f}^i_0$ is a morphism of $\catName{E}^{i+1}$, then
$\morphSeqName{f}^2_0 \morphCompose \morphSeqName{f}^1_0$ is a morphism of
$\catName{E}^3$.

\item[$\seqname{C}^i$ with category:]
If each $\catName{C}^i \subcat[full-] \catName{C}^{i+1}$ and each
$\morphSeqName{f}^i_1$ is a morphism of $\catName{C}^i$, then
$\morphSeqName{f}^2_1 \morphCompose \morphSeqName{f}^1_1$ is a
morphism of $\catName{C}^3$.
\end{description}
\end{proof}

\item
\label{M-morphcomp:constcompconst}
If each $\morphSeqName{f}^i$ is constrained then
$\morphSeqName{f}^2 \morphCompose[()] \morphSeqName{f}^1$ is constrained.

\begin{proof}
This is a corollary of \pagecref{lem:modcomp}.
\end{proof}

\item
\label{M-morphcomp:compequivcomp}
$\morphSeqName{f}^2 \morphCompose[()] \morphName{f}^1$ is equivalent to
$\morphSeqName{g}^2 \morphCompose[()] \morphName{g}^1$.

\begin{proof}
$\morphName{f}^1_0 = \morphName{g}^1_0$ and
$\morphName{f}^2_0 = \morphName{g}^2_0$, hence
$
  \morphName{f}^2_0 \morphCompose \morphName{f}^1_0 =
  \morphName{g}^2_0 \morphCompose \morphName{g}^1_0
$.
\end{proof}
\end{enumerate}
\end{lemma}

\begin{corollary}[Composition of m-atlas morphisms]
\label{cor:M-ATLmorphcomp}
Let $\catName{E}^1$ be a model category,
$\seqname{E}^1 \objin \catName{E}^1$,
$\catName{C}^i$, $i=1,2,3$, be a model category,
$\seqname{C}^i \objin \catName{C}^i$, $\seqname{A}^i$ be an m-atlas of
$\seqname{E}^1$ in the coordinate space $\seqname{C}^i$ and
$
  \morphSeqName{f}^i \defeq
    (
      \morphName{f}^i_0 \maps \seqname{E}^1 \to \seqname{E}^1,
      \morphName{f}^i_1 \maps \seqname{C}^i \to \seqname{C}^{i+1}
    )
$, $i=1,2$,
be a (semi-strict, strict)
\nobreaks
  {
    {($\catName{E}^1$-)},
    {$\seqname{E}^1$},
    {(-$\catName{C}^i$-$\catName{C}^{i+1}$)}
  }
m-atlas morphism of
$\seqname{A}^i$ to $\seqname{A}^{i+1}$ in the coordinate spaces
$\seqname{C}^i$, $\seqname{C}^{i+1}$.

Then

\begin{enumerate}
\item
$
  \morphSeqName{f}^2 \morphCompose[()] \morphSeqName{f}^1 =
    (
      \morphName{f}^2_0 \morphCompose \morphName{f}^1_0,
      \morphName{f}^2_1 \morphCompose \morphName{f}^1_1
    )
$
is a (semi-strict, strict)
\nobreaks
  {
    {($\catName{E}^1$-)},
    {$\seqname{E}^1$},
    {(-$\catName{C}^1$-$\catName{C}^3$)}
  }
m-atlas morphism of $\seqname{A}^1$ to $\seqname{A}^3$ in the
coordinate spaces $\seqname{C}^1$, $\seqname{C}^3$.

\begin{proof}
Since each $\morphSeqName{f}^i$ is a (semi-strict, strict)
\nobreaks
  {
    {($\catName{E}^1$-$\catName{E}^1$-)},
    {$\seqname{E}^1$-$\seqname{E}^1$},
    {(-$\catName{C}^i$-$\catName{C}^{i+1}$)}
  }
m-atlas morphism of
$\seqname{A}^i$ to $\seqname{A}^{i+1}$ in the coordinate spaces
$\seqname{C}^i$, $\seqname{C}^{i+1}$ by
\pagecref{def:M-ATLmorphabbrev},
$\morphSeqName{f}^2 \morphCompose[()] \morphSeqName{f}^1$ is a (semi-strict,
strict)
\nobreaks
  {
    {($\catName{E}^1$-$\catName{E}^1$-)},
    {$\seqname{E}^1$-$\seqname{E}^1$},
    {(-$\catName{C}^1$-$\catName{C}^3$)}
  }
m-atlas morphism of
$\seqname{A}^1$ to $\seqname{A}^3$ in the coordinate spaces
$\seqname{C}^1$, $\seqname{C}^3$ by \cref{M-morphcomp:morphcompmorph} of
{
  \showlabelsinline
  \pagecref{lem:M-ATLmorphcomp}.
}
Apply \cref{def:M-ATLmorphabbrev}.
\end{proof}
\end{enumerate}

Let $\catName{E}^i$ $i=1,2,3$, be a model category,
$\seqname{E}^i \objin \catName{E}^i$,
$\catName{C}^1$ be a model category,
$\seqname{C}^1 \objin \catName{C}^1$, $\seqname{A}^i$ be an m-atlas of
$\seqname{E}^i$ in the coordinate space $\seqname{C}^1$ and \\*
\nobreaks
  {{
    $
      \morphSeqName{f}^i \defeq
        (
          \morphName{f}^i_0 \maps \seqname{E}^i \to \seqname{E}^{i+1}
          \morphName{f}^i_1 \maps \seqname{C}^1 \to \seqname{C}^1
        )
    $,
  }}
$i=1,2$,
be a (semi-strict, strict)
\nobreaks
  {
    {($\catName{E}^i$-$\catName{E}^{i+1}$-)},
    {$\seqname{E}^i$-$\seqname{E}^{i+1}$},
    {(-$\catName{C}^1$)}
  }
m-atlas morphism of
$\seqname{A}^i$ to $\seqname{A}^{i+1}$ in the coordinate space
$\seqname{C}^1$.

\begin{enumerate}[resume]
\item
$
  \morphSeqName{f}^2 \morphCompose[()] \morphSeqName{f}^1 =
    (
      \morphName{f}^2_0 \morphCompose \morphName{f}^1_0,
      \morphName{f}^2_1 \morphCompose \morphName{f}^1_1
    )
$
is a (semi-strict, strict)
\nobreaks
 {
   {($\catName{E}^1$-$\catName{E}^3$-)},
   {$\seqname{E}^1$-$\seqname{E}^3$},
   {(-$\catName{C}^1$)}
  }
m-atlas morphism of
$\seqname{A}^1$ to $\seqname{A}^3$ in the coordinate space
$\seqname{C}^1$.

\begin{proof}
Apply \cref{def:M-ATLmorphabbrev} as above.
\end{proof}
\end{enumerate}

Let $\catName{E}^1$, $\catName{C}^1$ be model categories,
$\seqname{E}^1 \objin \catName{E}^1$,
$\seqname{C}^1 \objin \catName{C}^1$, $\seqname{A}^i$, $i=1,2,3$, be an
m-atlas of $\seqname{E}^1$ in the coordinate space $\seqname{C}^1$ and
$
  \morphSeqName{f}^i \defeq
    (
      \morphName{f}^i_0 \maps \seqname{E}^1 \to \seqname{E}^1,
      \morphName{f}^i_1 \maps \seqname{C}^1 \to \seqname{C}^1
    )
$, $i=1,2$,
be a (semi-strict, strict) $\seqname{E}^1$ m-atlas morphism of
$\seqname{A}^i$ to $\seqname{A}^{i+1}$ in the coordinate space
$\seqname{C}^1$.

\begin{enumerate}[resume]
\item
$
  \morphSeqName{f}^2 \morphCompose[()] \morphSeqName{f}^1 =
    (
      \morphName{f}^2_0 \morphCompose \morphName{f}^1_0,
      \morphName{f}^2_1 \morphCompose \morphName{f}^1_1
    )
$
is a (semi-strict, strict) $\seqname{E}^1$ m-atlas m-atlas morphism
of $\seqname{A}^1$ to $\seqname{A}^3$ in the coordinate space
$\seqname{C}^1$.

\begin{proof}
Apply \cref{def:M-ATLmorphabbrev} as above.
\end{proof}
\end{enumerate}

Let $\catName{C}^i$, $i=1,2,3$, be a model category,
$\seqname{C}^i \objin \catName{C}^i$, $E^i$ be a topological space,
$\seqname{A}^i$ be an m-atlas of $E^i$ in the coordinate space
$\seqname{C}^i$ and
\nobreaks
  {{
    $
      \morphSeqName{f}^i \defeq
        (
          \morphName{f}^i_0 \maps E^i \to E^{i+1},
          \morphName{f}^i_1 \maps \seqname{C}^i \to \seqname{C}^{i+1}
        )
    $,
  }}
$i=1,2$,
be a (semi-strict, strict)
\nobreaks
 {
    {$E^i$-$E^{i+1}$},
    {(-$\catName{C}^i$-$\catName{C}^{i+1}$)}
  }
m-atlas morphism of $\seqname{A}^i$ to $\seqname{A}^{i+1}$ in the
coordinate spaces $\seqname{C}^i$, $\seqname{C}^{i+1}$. Then

\begin{enumerate}[resume]
\item
$
  \morphSeqName{f}^2 \morphCompose[()] \morphSeqName{f}^1 =
    (
      \morphName{f}^2_0 \morphCompose \morphName{f}^1_0,
      \morphName{f}^2_1 \morphCompose \morphName{f}^1_1
    )
$
is a (semi-strict, strict)
$E^1$-$E^3$
\nobreaks
 {
    {$E^1$-$E^3$},
    {(-$\catName{C}^1$-$\catName{C}^3$)}
  }
m-atlas morphisms of $\seqname{A}^1$ to $\seqname{A}^3$ in the
coordinate spaces $\seqname{C}^1$, $\seqname{C}^3$.

\begin{proof}
Apply \cref{def:M-ATLmorphabbrev} as above.
\end{proof}
\end{enumerate}

Let $\catName{E}^i$, $\catName{C}^i$, $i \in [1,3]$, be model categories,
$\seqname{E}^i \objin \catName{E}^i$,
$\seqname{C}^i \objin \catName{C}^i$ $\seqname{A}^i$ an m-atlas
of $\seqname{E}^i$ in the coordinate space $\seqname{C}^i$ and
\nobreaks
  {{
    $
      \morphSeqName{f}^i \defeq
      (
        \morphName{f}^i_0 \maps \seqname{E}^i \to \seqname{E}^{i+1},
        \morphName{f}^i_1 \maps \seqname{C}^i \to \seqname{C}^{i+1}
      )
    $
  }}
a (semi-strict, strict) $\seqname{E}^i$-$\seqname{E}^{i+1}$ m-atlas
morphism of $\seqname{A}^i$ to $\seqname{A}^{i+1}$ in the coordinate
spaces $\seqname{C}^i$, $\seqname{C}^{i+1}$. Then

\begin{enumerate}[resume]
\item
$
  \morphSeqName{f}^2 \morphCompose[()] \morphSeqName{f}^1 =
  (
    \morphName{f}^2_0 \morphCompose \morphName{f}^1_0,
    \morphName{f}^2_1 \morphCompose \morphName{f}^1_1
  )
$
is a (semi-strict, strict) $\seqname{E}^1$-$\seqname{E}^3$ m-atlas
morphism of $\seqname{A}^1$ to $\seqname{A}^3$ in the coordinate spaces
$\seqname{C}^1$, $\seqname{C}^3$.

\begin{proof}
Apply \cref{def:M-ATLmorphabbrev} as above.
\end{proof}
\end{enumerate}

Let $\catName{C}^i$, $i=1,2,3$, be a model category,
$\seqname{C}^i \objin \catName{C}^i$, $E^i$ be a topological space,
$\seqname{A}^i$ be an m-atlase of $E^i$ in the coordinate space
$\seqname{C}^i$ and
\nobreaks
  {{
    $
      \morphSeqName{f}^i \defeq
      (
        \morphName{f}^i_0 \maps E^i \to E^{i+1},
        \morphName{f}^i_1 \maps \seqname{C}^i \to \seqname{C}^{i+1}
      )
    $
  }}
be a (semi-strict, strict)
$E^i$-$E^{i+1}$
m-atlas morphism of
$\seqname{A}^i$ to $\seqname{A}^{i+1}$ in the coordinate spaces
$\seqname{C}^i$, $\seqname{C}^{i+1}$. Then

\begin{enumerate}[resume]
\item
$
  \morphSeqName{f}^2 \morphCompose[()] \morphSeqName{f}^1 =
  (
    \morphName{f}^2_0 \morphCompose \morphName{f}^1_0,
    \morphName{f}^2_1 \morphCompose \morphName{f}^1_1
  )
$
is a (semi-strict, strict) $E^1$-$E^3$ m-atlas morphism of
$\seqname{A}^1$ to $\seqname{A}^3$ in the coordinate spaces
$\seqname{C}^1$, $\seqname{C}^3$.

\begin{proof}
Apply \cref{def:M-ATLmorphabbrev} as above.
\end{proof}
\end{enumerate}

Let $\catName{C}^i$, $i=1,2,3$, be a model category,
$\seqname{C}^i \objin \catName{C}^i$, $E^i$ be a topological space,
$\seqname{A}^i$ be an m-atlas of $E^i$ in the coordinate space
$\seqname{C}^i$ and
$
  \morphSeqName{f}^i \defeq
    (
      \morphName{f}^i_0 \maps E^i \to E^{i+1},
      \morphName{f}^i_1 \maps \seqname{C}^i \to \seqname{C}^{i+1}
    )
$, $i=1,2$,
be a (semi-strict, strict) $E^i$-$E^{i+1}$ m-atlas morphism of
$\seqname{A}^i$ to $\seqname{A}^{i+1}$ in the coordinate spaces
$\seqname{C}^i$, $\seqname{C}^{i+1}$. Then

\begin{enumerate}[resume]
\item
$
  \morphSeqName{f}^2 \morphCompose[()] \morphSeqName{f}^1 =
    (
      \morphName{f}^2_0 \morphCompose \morphName{f}^1_0,
      \morphName{f}^2_1 \morphCompose \morphName{f}^1_1
    )
$
is a (semi-strict, strict) $E^1$-$E^3$ m-atlas morphism of
$\seqname{A}^1$ to $\seqname{A}^3$ in the coordinate spaces
$\seqname{C}^1$, $\seqname{C}^3$.

\begin{proof}
Apply \cref{def:M-ATLmorphabbrev} as above.
\end{proof}
\end{enumerate}
\end{corollary}

\begin{lemma}[M-atlas identity]
\label{lem:M-ATLid}
Let
\begin{enumerate}
\item
$\catName{E}^i$, $\catName{C}^i$, $i=1,2$, be model categories
\item
$\catName{E}^1 \subcat[full-] \catName{E}^2$
\item
$\catName{C}^1 \subcat[full-] \catName{C}^2$
\item
$\seqname{E}^i \objin \catName{E}^i$
\item
$\seqname{E}^1 \submod \seqname{E}^2$
\item
$\seqname{C}^i \objin \catName{C}^i$
\item
$\seqname{C}^1 \submod \seqname{C}^2$
\item
$\seqname{A}^i$ be a maximal m-atlas of $\seqname{E}^i$ in the
coordinate space $\seqname{C}^i$
\item
$\seqname{A}^1 \subseteq \seqname{A}^2$
\end{enumerate}

The inclusion morphism
$
  \ID%
    [
      {(\seqname{E}^1, \seqname{C}^1)},
      {(\seqname{E}^2, \seqname{C}^2)}
    ]
$
is a semi-strict
\nobreaks
  {
    {($\catName{E}^1$-$\catName{E}^2$-)},
    {$\seqname{E}^1$-$\seqname{E}^2$-},
    {(-$\catName{C}^1$-$\catName{C}^2$)}
  }
m-atlas morphism of
$\seqname{A}^1$ to $\seqname{A}^2$ in the coordinate spaces
$\seqname{C}^1$, $\seqname{C}^2$.

\begin{proof}
\begin{description}

\item%
  [
    \ensuremath
      {
        \ID%
          [
            {(\seqname{E}^1, \seqname{C}^1)},
            {(\seqname{E}^2, \seqname{C}^2)}
          ]
      }
    is an  m-atlas morphism:
  ]

Let $(U^i,V^i,\phi^i) \in \seqname{A}^i$, $i=1,2$, be charts such that
$I \defeq U^1 \cap U^2 \neq \emptyset$. Since each $U^i$ is a model
neighborhood of $\seqname{E}^i$ and
$\seqname{E}^1 \submod \seqname{E}^2$ by hypothesis,
$U^2 \cap \Top(\seqname{E}^1)$ is a model neighborhood of
$\seqname{E}^1$ by \cref{modsub:rest} of
{
  \showlabelsinline
  \pagecref{def:modsub}
}
and \\*
$I = U^1 \cap U^2 \cap \Top(\seqname{E}^1)$ is a model neighborhood
of $\seqname{E}^1 \submod \seqname{E}^2$.

Define
\begin{enumerate}
\item
$
  \seqname{f} \defeq
  \ID%
    [
      {(\seqname{E}^1,\seqname{C}^1)},
      {(\seqname{E}^2,\seqname{C}^2)}
    ]
$
\item
$U'^1 \defeq I$
\item
$V'^1 \defeq \phi^1[I]$
\item
$U'^2 \defeq U'^1 = I$,
\item
$V'^2 \defeq \phi^2[U'^2]$
\item
$\hat{V}'^2 \defeq V'^1$
\item
$\phi'^2 \defeq \phi^1 \maps U'^2 = U'^1 \toiso \hat{V}'^2 = V'^1$.
\end{enumerate}

Then
\begin{enumerate}
\item
$U'^1 = I$ is a model neighborhood of $\seqname{E}^1$.
\item
Since $\phi^1$ is a model isomorphism, $V'^1$ is a model neighborhood of
$\seqname{C}^1$.
\item
$\phi^1 \restrictto_{U'^1,V'^1}$ is a model isomorphism of $\seqname{C}^1$.
\item
Since $\seqname{A}^i$ is maximal,
$(U'^1,V'^1,\phi^1 \restrictto_{U'^1,V'^1})$ is an m-chart of
$\seqname{E}^1$ in the coordinate space $\seqname{C}^1$.

\item
Since $\seqname{A}^1$ is maximal,
$(U'^1,V'^1, \phi^1 \maps U'^1 \toiso V'^1) \in \seqname{A}^1$.
\item
Since $\seqname{E}^1 \submod \seqname{E}^2$, $U'^2 = U'^1$ is a
model neighborhood of $\seqname{E}^2$.
\item
Since $\phi^2$ is a model isomorphism, $V'^2$ is a model neighborhood of
$\seqname{C}^2$.
\item
Since $\seqname{C}^1 \submod \seqname{C}^2$, $\hat{V}'^2 = V'^1$ is a
model neighborhood of $\seqname{C}^2$.
\item
$\phi^1 \restrictto_{U'^2,\hat{V}'^2}$ is a model isomorphism of
$\seqname{C}^2$.
\item
$(U'^2,\hat{V}'^2,\phi^1 \restrictto_{U'^2,\hat{V}'^2})$ is an m-chart
of $\seqname{E}^2$ in the coordinate space $\seqname{C}^2$.
\item
$
  \morphName{f}_0[U'^1] =
  \Id[{\seqname{E}^1},{\seqname{E}^2}][U'^1] =
  U'^1 \subseteq U'^2
$
\item
Diagram \eqref{eq:AtlM1} in
{
  \showlabelsinline
  \namecref{def:M-ATLmorph}
}
is commutative, as shown in \pagecref{fig:AtlM1}.
\end{enumerate}

\item%
  [
    \ensuremath
      {
        \ID%
          [
            {(\seqname{E}^1, \seqname{C}^1)},
            {(\seqname{E}^2, \seqname{C}^2)}
          ]
      }
    is semi-strict:
  ]

Let $u^1 \in \seqname{E}^1$. Since the model neighborhoods of
$\seqname{E}^1$ form an open cover, there is a model neighborhood $U$ of
$\seqname{E}^1$ at $u^1$. $\Id[{U}]$ is a morphism of $\seqname{E}^1$
and of $\catName{E}^1$; since $\seqname{E}^1 \submod \seqname{E}^2$ and
$\catName{E}^1 \subcat[full-] \catName{E}^2$, $\Id[{U}]$ is also a
morphism of $\seqname{E}^2$ and of $\catName{E}^2$. Thus
$
\Id%
  [
    {\seqname{E}^1},
    {\seqname{E}^2}
  ]
$
is a local m-morphism of $\seqname{E}^1$ to $\seqname{E}^12$ and a local
$\catName{E}^1$-$\catName{E}^2$ m-morphism $\seqname{E}^1$ to $\seqname{E}^12$.

Let $v^1 \in \seqname{C}^1$. Since the model neighborhoods of
$\seqname{C}^1$ form an open cover, there is a model neighborhood $V$ of
$\seqname{C}^1$ at $v^1$. $\Id[{V}]$ is a morphism of $\seqname{C}^1$
and of $\catName{C}^1$; since $\seqname{C}^1 \submod \seqname{C}^2$ and
$\catName{C}^1 \subcat[full-] \catName{C}^2$, $\Id[{V}]$ is also a
morphism of $\seqname{C}^2$ and of $\catName{C}^2$. Thus
$
  \Id%
    [
      {\seqname{C}^1},
      {\seqname{C}^2}
    ]
$
is a local m-morphism of $\seqname{C}^1$ to $\seqname{C}^12$ and a local
$\catName{C}^1$-$\catName{C}^2$ m-morphism $\seqname{C}^1$ to $\seqname{C}^12$.

\end{description}
\end{proof}

If each $\Top(\seqname{E}^i) \objin \seqname{E}^i$
($\seqname{E}^i \objin \catName{E}^i$)
and each $\Top(\seqname{C}^i) \objin \seqname{C}^i$
($\seqname{C}^i \objin \catName{C}^i$) then
$
  \ID%
    [
      {(\seqname{E}^1, \seqname{C}^1)},
      {(\seqname{E}^2, \seqname{C}^2)}
    ]
$
is a strict
\nobreaks
  {
    {($\catName{E}^1$-$\catName{E}^2$-)},
    {$\seqname{E}^1$-$\seqname{E}^2$-},
    {(-$\catName{C}^1$-$\catName{C}^2$)}
  }
m-atlas morphism of
$\seqname{A}^1$ to $\seqname{A}^2$ in the coordinate spaces
$\seqname{C}^1$, $\seqname{C}^2$.

\begin{proof}
There are four overlapping cases
\begin{description}
\item%
  [
    Each
    \ensuremath
      {
        \Top{\seqname{E}^i} \objin \seqname{E}^i
      }:
  ]
$\Id[{\seqname{E}^1},{\seqname{E}^2}]$ is a morphism of
$\seqname{E}^2$ by \cref{mod:inclusion} of
{
  \showlabelsinline
  \pagecref{def:model}
},
hence an m-morphism of $\seqname{E}^1$ to $\seqname{E}^2$.

\item%
  [
    Each
    \ensuremath
      {
        \seqname{E}^i \objin \catName{E}^2
      }:
  ]
$\Id[{\seqname{E}^1},{\seqname{E}^2}]$ is a morphism of
$\catName{E}^2$ by \cref{modcat:morph} of
{
  \showlabelsinline
  \pagecref{def:modcat}
},
hence an $\catName{E}^1$-$\catName{E}^2$ m-morphism of $\seqname{E}^1$
to $\seqname{E}^2$.

\item%
  [
    Each
    \ensuremath
      {
        \Top{\seqname{C}^i} \objin \seqname{C}^i
      }:
  ]
$\Id[{\seqname{C}^1},{\seqname{C}^2}]$ is a morphism of
$\seqname{C}^2$ by \cref{mod:inclusion} of
{
  \showlabelsinline
  \cref{def:model}
},
hence an m-morphism of $\seqname{C}^1$ to $\seqname{C}^2$.

\item%
  [
    Each
    \ensuremath
      {
        \seqname{C}^i \objin \catName{C}^2
      }:
  ]
$\Id[{\seqname{C}^1},{\seqname{C}^2}]$ is a morphism of
$\catName{C}^2$ by \cref{modcat:morph} of
{
  \showlabelsinline
  \cref{def:modcat}
},
hence a $\catName{C}^1$-$\catName{C^2}$ m-morphism of $\seqname{C}^1$ to
$\seqname{C}^2$.
\end{description}
\end{proof}
\end{lemma}

\subsection{Categories of m-atlases and functors}
\label{sub:M-cat}
This subsection defines categories of m-atlases and functors among them,
and proves some basic results.

\subsubsection{M-atlas categories}
In the following, \textbf{...} will refer to one of
\begin{equation*}
  \seqname{A}^1, \seqname{E}^1, \seqname{C}^1,
  \seqname{A}^2, \seqname{E}^2, \seqname{C}^2,
  \morphName{f}_0, \morphName{f}_1
\end{equation*}
\begin{equation*}
  \seqname{A}^1, E^1, \seqname{C}^1,
  \seqname{A}^2, E^2, \seqname{C}^2,
  \morphName{f}_0, \morphName{f}_1
\end{equation*}
depending on context.

\begin{definition}[Sets of m-atlas morphisms]
\label{def:M-ATLmorphsets}
Let $\seqname{E}^i$ and $\seqname{C}^i$. $i=1,2$, be model spaces. Then
the set of (constrained) (semi-strict, strict)
(full, maximal, full maximal)
$\seqname{E}^1$-$\seqname{E}^2$ m-atlas morphisms in the
coordinate spaces $\seqname{C}^1$,  $\seqname{C}^2$ is

\begin{multline}
  \strict[*]
    {
      \maxfull[*]
        {
          \Atl[(const)]_\Ar
            (\seqname{E}^1, \seqname{C}^1, \seqname{E}^2, \seqname{C}^2)
        }
    }
  \defeq \\*
  \set
  {{
    \bigl (
      (
        \morphName{f}_0,
        \morphName{f}_1
      ),
      (\seqname{A}^1, \seqname{E}^1, \seqname{C}^1),
      (\seqname{A}^2, \seqname{E}^2, \seqname{C}^2)
    \bigr )
  }}%
  [
    {
      \strict[*]
      {
        \maxfull[*]
          {
            \isAtl[(const)]
              (
                \seqname{A}^1, \seqname{E}^1, \seqname{C}^1,
                \seqname{A}^2, \seqname{E}^2, \seqname{C}^2,
                \morphName{f}_0, \morphName{f}_1
              )
          }
      }
    }
  ]*
\end{multline}

Let $E^i$, $i=1,2$, be a topological space and
$\seqname{C}^1=(C^i, \catName{C}^i)$, $i=1,2$, be a model space.
Then
Then the set of (constrained) (semi-strict, strict)
(full, maximal, full maximal)
$E^1$-$E^2$ m-atlas morphisms in the
coordinate spaces $\seqname{C}^1$,  $\seqname{C}^2$ is

\begin{multline}
  \strict[*]
    {
      \maxfull[*]
        {
          \Atl[(const)]_\Ar
            (E^1, \seqname{C}^1, E^2, \seqname{C}^2)
        }
    }
  \defeq \\*
  \set
  {{
    \bigl (
      (
        \morphName{f}_0,
        \morphName{f}_1
      ),
      (\seqname{A}^1, E^1, \seqname{C}^1),
      (\seqname{A}^2, E^2, \seqname{C}^2)
    \bigr )
  }}%
  [
    {
      \strict[*]
      {
        \maxfull[*]
          {
            \isAtl[(const)]
              (
                \seqname{A}^1, E^1, \seqname{C}^1,
                \seqname{A}^2, E^2, \seqname{C}^2,
                \morphName{f}_0, \morphName{f}_1
              )
          }
      }
    }
  ]*
\end{multline}

Let
\begin{enumerate}
\item
$\seqname{E}$ and $\seqname{C}$ be sets of model spaces
\item
$\seqname{P} \defeq \seqname{E} \times \seqname{C}$
\end{enumerate}

\begin{multline}
  \strict[*]
    {
      \maxfull[*]
        {
          \Atl[(const)]_\Ar
            (\seqname{E}, \seqname{C})
        }
    }
  \defeq \\*
  \set
  {{
    \bigl (
      (
        \morphName{f}_0,
        \morphName{f}_1
      ),
      (\seqname{A}^1, \seqname{E}^1, \seqname{C}^1),
      (\seqname{A}^2, \seqname{E}^2, \seqname{C}^2)
    \bigr )
  }}%
  [
    {(\seqname{E}^i, \seqname{C}^i) \in \seqname{P}},
    {
      \strict[*]
      {
        \maxfull[*]
          {
            \isAtl[(const)]
              (
                ...
              )
          }
      }
    }
  ]*
\end{multline}

Let
\begin{enumerate}
\item
$\catName{E}$ be a model category
\item
$\seqname{C}$ be a set of model spaces
\item
$\seqname{P} \defeq \Ob(\catName{E}) \times \seqname{C}$
\end{enumerate}

\begin{multline}
  \strict[*]
    {
      \maxfull[*]
        {
          \Atl[(const)]_\Ar
            (\seqname{E}, \catName{C})
        }
    }
  \defeq \\*
  \set
  {{
    \bigl (
      (
        \morphName{f}_0,
        \morphName{f}_1
      ),
      (\seqname{A}^1, \seqname{E}^1, \seqname{C}^1),
      (\seqname{A}^2, \seqname{E}^2, \seqname{C}^2)
    \bigr )
  }}%
  [
    {(\seqname{E}^i, \seqname{C}^i) \in \seqname{P}},
    {
      \strict[*]
      {
        \maxfull[*]
          {
            \isAtl[(const)]
              (
                ...,
                \catName{E},
                \catName{E}
              )
          }
      }
    }
  ]*
\end{multline}

Let
\begin{enumerate}
\item
$\seqname{E}$ be a set of model spaces
\item
$\seqname{C}$ be a model category
\item
$\seqname{P} \defeq \seqname{E} \times Ob(\catName{C})$
\end{enumerate}

\begin{multline}
  \strict[*]
    {
      \maxfull[*]
        {
          \Atl[(const)]_\Ar
            (\catName{E}, \catName{C})
        }
    }
  \defeq \\*
  \set
  {{
    \bigl (
      (
        \morphName{f}_0,
        \morphName{f}_1
      ),
      (\seqname{A}^1, \seqname{E}^1, \seqname{C}^1),
      (\seqname{A}^2, \seqname{E}^2, \seqname{C}^2)
    \bigr )
  }}%
  [
    {(\seqname{E}^i, \seqname{C}^i) \in \seqname{P}},
    {
      \strict[*]
      {
        \maxfull[*]
          {
            \isAtl[(const)]
              (
                ...,
                \catName{C},
                \catName{C}
              )
          }
      }
    }
  ]*
\end{multline}

Let
\begin{enumerate}
\item
$\catName{E}$ and $\catName{C}$ be model categories
\item
$\seqname{P} \defeq \Ob(\catName{E}) \times Ob(\catName{C})$
\end{enumerate}

\begin{multline}
  \strict[*]
    {
      \maxfull[*]
        {
          \Atl[(const)]_\Ar
            (\catName{E}, \catName{C})
        }
    }
  \defeq \\*
  \set
  {{
    \bigl (
      (
        \morphName{f}_0,
        \morphName{f}_1
      ),
      (\seqname{A}^1, \seqname{E}^1, \seqname{C}^1),
      (\seqname{A}^2, \seqname{E}^2, \seqname{C}^2)
    \bigr )
  }}%
  [
    {(\seqname{E}^i, \seqname{C}^i) \in \seqname{P}},
    {
      \strict[*]
      {
        \maxfull[*]
          {
            \isAtl[(const)]
              (
                ...,
                \catName{E},
                \catName{E},
                \catName{C},
                \catName{C}
              )
          }
      }
    }
  ]*
\end{multline}
\end{definition}

\begin{lemma}[Sets of m-atlas morphisms]
\label{lem:M-ATLmorphsets}
Let $\seqname{E}$ and $\seqname{C}$ be sets of model spaces.

\begin{align*}
  \strict[*]
    {
      \maxfull[*]
        {
          \Atl[(const)]_\Ar
            (\seqname{E}, \seqname{C})
        }
    }
  & \subseteq
  \Atl_\Ar
    (\seqname{E}, \seqname{C})
\\
  \strict[*]
    {
      \maxfull
        {
          \Atl[(const)]_\Ar
            (\seqname{E}, \seqname{C})
        }
    }
  & \subseteq
  \strict[*]
    {
      \full
        {
          \Atl[(const)]_\Ar
            (\seqname{E}, \seqname{C})
        }
    }
\end{align*}

\begin{proof}
The result follows from
{
  \showlabelsinline
  \pagecref{def:M-ATLmorphabbrev}
}
and \cref{def:M-ATLmorphsets} above.
\end{proof}
\end{lemma}

\begin{corollary}[Sets of m-atlas morphisms]
\label{cor:M-ATLmorphsets}
Let $\seqname{E}$ and $\seqname{C}$ be sets of topological spaces.

\begin{align*}
  \strict[*]
    {
      \maxfull[*]
        {
          \Atl[(const)]_\Ar
            (\seqname{E}, \seqname{C})
        }
    }
  & \subseteq
  \Atl_\Ar
    (\seqname{E}, \seqname{C})
\\
  \strict[*]
    {
      \maxfull
        {
          \Atl[(const)]_\Ar
            (\seqname{E}, \seqname{C})
        }
    }
  & \subseteq
  \strict[*]
    {
      \full
        {
          \Atl[(const)]_\Ar
            (\seqname{E}, \seqname{C})
        }
    }
\end{align*}

\begin{proof}
The result follows from
{
  \showlabelsinline
  \pagecref{def:M-ATLmorphabbrev}
}
and \cref{def:M-ATLmorphsets} above.
\end{proof}
\end{corollary}

\begin{definition}[Categories $\Atl(\seqname{E}, \seqname{C})$]
\label{def:M-cat}
Let $\seqname{E}$ and $\seqname{C}$ be sets of model spaces.
Let $P \defeq \seqname{E} \times \seqname{C}$.

\begin{multline}
  \strict[*]
    {
      \maxfull[*]
        {
          \Atl[(const)]
            (\seqname{E}, \seqname{C})
        }
    }
  \defeq \\*
  \biggl (
    \maxfull[*]
      {
        \Atl_\Ob
          (\seqname{E}, \seqname{C})
      },
    \strict[*]
      {
        \maxfull[*]
          {
            \Atl[(const)]_\Ar
              (\seqname{E}, \seqname{C})
          }
      },
    \morphCompose[A]
  \biggr )
\end{multline}

Let $\seqname{E}^i$ and $\seqname{C}^i$, $i=1,2$, be model spaces,
$\seqname{A}^i$ be an atlas of $\seqname{E}^i$ in the coordinate
space $\seqname{C}^i$ and
$
\morphSeqName{f} \defeq
(
  \morphName{f}_0 \maps \seqname{E}^1 \to \seqname{E}^2,
  \morphName{f}_0 \maps \seqname{C}^1 \to \seqname{C}^2
)
$
be a (strict, semi-strict) $\seqname{E}^1$-$\seqname{E}^2$ M-atlas
morphism of $\seqname{A}^1$ to $\seqname{A}^2$ in the coordinate
spaces $\seqname{C}^1$, $\seqname{C}^2$.

The identity functor $\Functor_{\M,\Id}$ is

\begin{equation}
\Functor_{\M,\Id} (\seqname{A}^i, \seqname{E}^i, \seqname{C}^i)
\defeq
(\seqname{A}^i, \seqname{E}^i, \seqname{C}^i)
\end{equation}

This nomenclature will be justified below.

\begin{multline}
\Functor_{\M,\Id}
  \biggl (
    (
      \morphName{f}_0 \maps \seqname{E}^1 \to \seqname{E}^2,
      \morphName{f}_1 \maps \seqname{C}^1 \to \seqname{C}^2,
    ),
    (\seqname{A}^1, \seqname{E}^1, \seqname{C}^1),
    (\seqname{A}^2, \seqname{E}^2, \seqname{C}^2)
  \biggr )
\defeq
\\*
  \bigl (
    (
      \morphName{f}_0 \maps \seqname{E}^1 \to \seqname{E}^2,
      \morphName{f}_1 \maps \seqname{C}^1 \to \seqname{C}^2,
    ),
    (\seqname{A}^1, \seqname{E}^1, \seqname{C}^1),
    (\seqname{A}^2, \seqname{E}^2, \seqname{C}^2)
  \bigr )
\end{multline}

These are fine grained iff $\seqname{C}$ is fine grained.
\end{definition}

\begin{lemma}[$\Atl(\seqname{E}, \seqname{C})$ is a category]
\label{lem:M-ATLiscat}
Let $\seqname{E}$ and $\seqname{C}$ be sets of model spaces. Then
\begin{enumerate}
\item
\label{Mcat:morph}
$
  \strict[*]
    {
      \maxfull[*]
        {
          \Atl[(const)]
            (\seqname{E}, \seqname{C})
        }
    }
$
is a category and the identity
morphism for an object $(\seqname{A}^i, \seqname{E}^i, \seqname{C}^i)$
of
$
  \strict[*]
    {
      \maxfull[*]
        {
          \Atl[(const)]_\Ob
            (\seqname{E}, \seqname{C})
        }
    }
$
is
$\Id[{{(\seqname{A}^i, \seqname{E}^i, \seqname{C}^i)}}]$.

\begin{proof}
Let $(\seqname{A}^i,\seqname{E}^i,\seqname{C}^i)$, $i \in [1,3]$,
be an object of $\Atl(\seqname{E}, \seqname{C})$ and
let \\
$
  \morphName{m}^i \defeq
  \bigl (
    (\morphName{f}_0^i,\morphName{f}_1^i),
    (\seqname{A}^i,\seqname{E}^i,\seqname{C}^i),
    (\seqname{A}^{i+1},\seqname{E}^{i+1},\seqname{C}^{i+1})
  \bigr )
$
be a morphism of $\Atl(\seqname{E}, \seqname{C})$. Then
\begin{description}
\item [Composition:]
\leavevmode \\*
$
    \morphSeqName{m}^2 \morphCompose[A] \morphSeqName{m}^1 =
  \bigl (
    (
      \morphName{f}_0^2 \morphCompose \morphName{f}_0^1,
      \morphName{f}_1^2 \morphCompose \morphName{f}_1^1
    ),
    (\seqname{A}^1,\seqname{E}^1,\seqname{C}^1),
    (\seqname{A}^3,\seqname{E}^3,\seqname{C}^3)
  \bigr )
$
is a morphism of $\Atl(\seqname{E}, \seqname{C})$ by
\cref{M-morphcomp:morphcompmorph} of
{
  \showlabelsinline
  \fullcref{lem:M-ATLmorphcomp} on
  \cpageref{M-morphcomp:morphcompmorph}.
}

\item [Associativity:]
\leavevmode \\*
Composition is associative by \pagecref{lem:labcomp}.

\item [Identity:]
\leavevmode \\*
$\Id[{{(\seqname{A}^i, \seqname{E}^i, \seqname{C}^i)}}]$ is an identity
morphism by \cref{lem:labcomp}.
\end{description}
\end{proof}
\end{enumerate}
\end{lemma}

\begin{corollary}[Subcategories of $\Atl(\seqname{E}, \seqname{C})$]
\label{cor:M-ATLiscat}
Let $\seqname{E}$ and $\seqname{C}$ be sets of model spaces. Then
$
  \strict[*]
    {
      \maxfull[*]
        {
          \Atl[(const)]
            (\seqname{E}, \seqname{C})
        }
    }
$
is a subcategory of
$\Atl[(const)] (\seqname{E}, \seqname{C})$.

\begin{remark}
They are not, in general, full subcategories.
\end{remark}

The identity functor $\Functor_{\M,\Id}$ is a functor from each of the
subcategories to each of the containing categories.
\end{corollary}

\begin{definition}[Categories $\Atl(\catName{E}, \catName{C})$]
\label{def:M-cat-cat}
Let $\catName{E}$ and $\catName{C}$ be model categories.
Let $\seqname{P} \defeq Ob(\catName{E}) \times \Ob(\catName{C})$.
Then

\begin{multline}
  \maxfull[*]
    {
      \Atl_\Ob(\catName{E}, \catName{C})
    }
  \defeq \\*
  \union
    [
        {\seqname{E} \objin \catName{E}},
        {\seqname{C} \objin \catName{C}}
    ]
    {{
      \maxfull[*]
        {
          \Atl_\Ob(\seqname{E}, \seqname{C})
        }
    }}
\end{multline}
\begin{multline}
  \strict[*]
    {
      \maxfull[*]
        {
          \Atl[(const)]_\Ar
            (\catName{E}, \catName{C})
        }
    }
  \defeq \\*
    \set
    {{
      \bigl (
        (
          \morphName{f}_0,
          \morphName{f}_1
        ),
        (\seqname{A}^1, \seqname{E}^1, \seqname{C}^1),
        (\seqname{A}^2, \seqname{E}^2, \seqname{C}^2)
      \bigr )
    }}%
    [
      {
        (\seqname{E}^i, \seqname{C}^i) \in \seqname{P}
      },
      {
        \strict[*]
        {
          \maxfull[*]
          {
            \isAtl[(const)]_\Ar
            (
              \seqname{A}^1,
              \seqname{E}^1,
              \seqname{C}^1,
              \seqname{A}^2,
              \seqname{E}^2,
              \seqname{C}^2,
              \morphName{f}_0,
              \morphName{f}_1,
              \catName{E},
              \catName{C}
            )
          }
        }
      }
    ]*
\end{multline}
\begin{multline}
  \strict[*]
  {
    \maxfull[*]
    {
      \Atl[(const)](\catName{E}, \catName{C})
    }
  }
  \defeq \\*
  \bigl (
    \Atl_\Ob(\catName{E}, \catName{C}),
    \Atl[(const)]_\Ar(\catName{E}, \catName{C}),
    \morphCompose[A]
  \bigr )
\end{multline}
\end{definition}

\begin{lemma}[$\Atl(\catName{E}, \catName{C})$ is a category]
\label{lem:M-cat-ATLiscat}
Let $\catName{E}$ and $\catName{C}$ be model categories. Then

$
  \strict[*]
  {
    \maxfull[*]
    {
      \Atl[(const)](\catName{E}, \catName{C})
    }
  }
$
is a category and the identity morphism for an object
$(\seqname{A}^i, \seqname{E}^i, \seqname{C}^i)$ of
$
  \strict[*]
  {
    \maxfull[*]
    {
      \Atl[(const)](\catName{E}, \catName{C})
    }
  }
$
is $\Id[{{(\seqname{A}^i, \seqname{E}^i, \seqname{C}^i)}}]$.

\begin{proof}
Let $(\seqname{A}^i,\seqname{E}^i,\seqname{C}^i)$, $i \in [1,3]$,
be an object of $\Atl(\catName{E}, \catName{C})$ and
let \\
$
  \morphName{m}^i \defeq
  \bigl (
    (\morphName{f}_0^i,\morphName{f}_1^i),
    (\seqname{A}^i,\seqname{E}^i,\seqname{C}^i),
    (\seqname{A}^{i+1},\seqname{E}^{i+1},\seqname{C}^{i+1})
  \bigr )
$
be a morphism of $\Atl(\catName{E}, \catName{C})$. Then
\begin{description}
\item[Composition:]
\leavevmode \\*
$\morphName{f}_0^2 \morphCompose \morphName{f}_0^1$ is a morphism of
$\catName{E}$, $\morphName{f}_1^2 \morphCompose \morphName{f}_1^1$ is a
morphism of $\catName{C}$ and
$
  \morphSeqName{f}^2 \morphCompose[()] \morphSeqName{f}^1 =
  (
    \morphName{f}_0^2 \morphCompose \morphName{f}_0^1,
    \morphName{f}_1^2 \morphCompose \morphName{f}_1^1
  )
$
is an $E^1$-$E^3$ m-atlas morphism of $\seqname{A}^1$ to $\seqname{A}^3$
in the coordinate spaces $\seqname{C}^1$, $\seqname{C}^2$ by
\cref{M-morphcomp:morphcompmorph} of
{
  \showlabelsinline
  \fullcref{lem:M-ATLmorphcomp} on
  \cpageref{M-morphcomp:morphcompmorph}.
}

\item[Associativity:]
\leavevmode \\*
Composition is associative by \pagecref{lem:labcomp}.

\item[Identity:]
$\Id[{{(\seqname{A}^i, \seqname{E}^i, \seqname{C}^i)}}]$ is an identity
morphism by \cref{lem:labcomp}.
\end{description}
\end{proof}
\end{lemma}

\begin{definition}[Categories $\Atl(\seqname{E}, \seqname{C})$ of topological spaces]
\label{def:T-cat}
Let $\seqname{E}$ be a set of topological spaces and $\seqname{C}$ a
set of model spaces.
Let $\seqname{P} \defeq \seqname{E} \times \seqname{C}$.
Then
\begin{equation}
  \Atl_\Ob(\seqname{E}, \seqname{C})
  \defeq
  \union
    [
        {E^1 \in \seqname{E}},
        {\seqname{C}^1 \in \seqname{C}}
    ]
    {{
      \Atl_\Ob(E^1, \seqname{C}^1)
    }}
\end{equation}
\begin{equation}
  \Atl_\Ar(\seqname{E}, \seqname{C})
  \defeq
  \union
    [
        {(E^1, \seqname{C}^1) \in \seqname{P}},
        {(E^2, \seqname{C}^2) \in \seqname{P}}
    ]
    {{
      \Atl_{\Ar}
      (
        E^1,
        \seqname{C}^1,
        E^2,
        \seqname{C}^2
      )
    }}
\end{equation}
\begin{equation}
  \Atl(\seqname{E}, \seqname{C})
  \defeq
  \bigl (
    \Atl_\Ob(\seqname{E}, \seqname{C}),
    \Atl_\Ar(\seqname{E}, \seqname{C}),
    \morphCompose[A]
  \bigr )
\end{equation}

Similar definitions apply with restrictions on the admissible atlases,
the admissible morphisms, or both:
\begin{enumerate}
\item $\full{\Atl}$
\item $\maximal[S-]{\Atl}$
\item $\maximal{\Atl}$
\item $\maxfull[S-]{\Atl}$
\item $\maxfull{\Atl}$
\item $\strict[semi-]{\full{\Atl}}$
\item $\strict[semi-]{\maximal[S-]{\Atl}}$
\item $\strict[semi-]{\maximal{\Atl}}$
\item $\strict[semi-]{\maxfull[S-]{\Atl}}$
\item $\strict[semi-]{\maxfull{\Atl}}$
\item $\strict{\full{\Atl}}$
\item $\strict{\maximal[S-]{\Atl}}$
\item $\strict{\maximal{\Atl}}$
\item $\strict{\maxfull[S-]{\Atl}}$
\item $\strict{\maxfull{\Atl}}$
\end{enumerate}
e.g.,

\begin{equation}
  \maxfull[S-]{\Atl_\Ob}(\seqname{E}, \seqname{C})
  \defeq
  \union
    [
      {(E^1, \seqname{C}^1) \in \seqname{P}},
    ]
    {{
      \maxfull[S-]{\Atl_\Ob}(E^1, \seqname{C}^1)
    }}
\end{equation}
\begin{equation}
  \strict[semi-]{\maxfull[S-]{\Atl_\Ar}}(\seqname{E}, \seqname{C})
  \defeq
  \union
    [
      {(E^1, \seqname{C}^1) \in \seqname{P}},
      {(E^2, \seqname{C}^2) \in \seqname{P}}
    ]
    {{
      \strict[semi-]{\maxfull[S-]{\Atl_\Ar}}
      (
        E^1,
        \seqname{C}^1,
        E^2,
        \seqname{C}^2
      )
    }}
\end{equation}
\begin{equation}
  \strict[semi-]{\maxfull[S-]{\Atl}}(\seqname{E}, \seqname{C})
  \defeq
  \bigl (
    \maxfull[S-]{\Atl_\Ob}(\seqname{E}, \seqname{C}),
    \strict[semi-]{\maxfull[S-]{\Atl_\Ar}}(\seqname{E}, \seqname{C}),
    \morphCompose[A]
  \bigr )
\end{equation}

Let
$
  (\seqname{A}^1, \seqname{E}^1, \seqname{C}^1) \in
  \Atl_\Ob(\seqname{E}, \seqname{C})
$.

\begin{equation}
\Id[{{(\seqname{A}^1, \seqname{E}^1, \seqname{C}^1)}}] \defeq
\bigl (
  (\Id[{{\seqname{E}^1}}], \Id[{{\seqname{C}^1}}]),
  (\seqname{A}^1, \seqname{E}^1, \seqname{C}^1),
  (\seqname{A}^1, \seqname{E}^1, \seqname{C}^1)
\bigr )
\end{equation}
\end{definition}

\begin{lemma}[$\Atl(\seqname{E}, \seqname{C})$ of topological spaces is a category]
\label{lem:T-ATLiscat}
Let $\seqname{E}$ be a set of topological spaces and $\seqname{C}$ a
set of model spaces. Then $\Atl(\seqname{E}, \seqname{C})$ is a category
and the identity morphism for an object
$(\seqname{A}^1, E^1, \seqname{C}^1)$ is
$\Id[{{(\seqname{A}^1, E^1, \seqname{C}^1)}}]$.

\begin{proof}
Let $(\seqname{A}^i,E^i,\seqname{C}^i)$, $i \in [1,3]$,
be an object of $\Atl(E, \seqname{C})$. and let
$
  \morphName{m}^i \defeq \\
  \bigl (
    (\morphName{f}_0^i,\morphName{f}_1^i),
    (\seqname{A}^i,E^i,\seqname{C}^i),
    (\seqname{A}^{i+1},E^{i+1},\seqname{C}^{i+1})
  \bigr )
$
be a morphism of $\Atl(\seqname{E}, \seqname{C})$. Then
\begin{description}
\item[Composition:]
\leavevmode \\*
$
  \bigl (
    (
      \morphName{f}_0^2 \morphCompose \morphName{f}_0^1,
      \morphName{f}_1^2 \morphCompose \morphName{f}_1^1
    ),
    (\seqname{A}^1,E^1,\seqname{C}^1),
    (\seqname{A}^3,E^3,\seqname{C}^3)
  \bigr )
$
is a morphism of $\Atl(\seqname{E}, \seqname{C})$ by
\pagecref{cor:M-ATLmorphcomp}.

\item[Associativity:]
\leavevmode \\*
Composition is associative by \pagecref{lem:labcomp}.

\item[Identity:]
$\Id[{{(\seqname{A}^i, \seqname{E}^i, \seqname{C}^i)}}]$ is an identity
morphism by \cref{lem:labcomp}.
\end{description}
\end{proof}

Similar results follow with restrictions on the admissible atlases,
the admissible morphisms, or both:
\begin{enumerate}
\item $\full{\Atl}(\catName{E}, \catName{C})$
\item $\maximal[S-]{\Atl}(\catName{E}, \catName{C})$
\item $\maximal{\Atl}(\catName{E}, \catName{C})$
\item $\maxfull[S-]{\Atl}(\catName{E}, \catName{C})$
\item $\maxfull{\Atl}(\catName{E}, \catName{C})$
\item $\strict[semi-]{\full{\Atl}}(\catName{E}, \catName{C})$
\item $\strict[semi-]{\maximal[S-]{\Atl}}(\catName{E}, \catName{C})$
\item $\strict[semi-]{\maximal{\Atl}}(\catName{E}, \catName{C})$
\item $\strict[semi-]{\maxfull[S-]{\Atl}}(\catName{E}, \catName{C})$
\item $\strict[semi-]{\maxfull{\Atl}}(\catName{E}, \catName{C})$
\item $\strict{\full{\Atl}}(\catName{E}, \catName{C})$
\item $\strict{\maximal[S-]{\Atl}}(\catName{E}, \catName{C})$
\item $\strict{\maximal{\Atl}}(\catName{E}, \catName{C})$
\item $\strict{\maxfull[S-]{\Atl}}(\catName{E}, \catName{C})$
\item $\strict{\maxfull{\Atl}}(\catName{E}, \catName{C})$
\end{enumerate}
\end{lemma}

\begin{definition}[$\Triv{\seqname{E}}$ and $\Triv{\Atl}(\seqname{E}, \seqname{C})$]
\label{def:M-triv}
Let $\seqname{E}$ be a set of topological spaces and $\seqname{C}$ a
set of model spaces. Then
\begin{equation}
\Triv{\seqname{E}} \defeq \set {\Triv{E^\mu}}[E^\mu \in \seqname{E}]
\end{equation}
\begin{equation}
\Triv{\Atl}(\seqname{E}, \seqname{C}) \defeq
\Atl \Big ( \Triv{\seqname{E}}, \seqname{C} \Big )
\end{equation}
\end{definition}

\subsubsection{M-atlas functors}
\begin{definition}[$\Functor_{\M,\Top}$]
\label{def:Func-M-Top}
Let $\seqname{E}^i \defeq (E^i,\catName{E}^i)$, $i=1,2$, and
$\seqname{C}^i \defeq (C^i,\catName{C}^i)$
be model spaces and $\seqname{A}^i$ be an m-atlas of
$\seqname{E}^i$ in the coordinate space $\seqname{C}^i$. Then
\begin{equation}
\Functor_{\M,\Top} (\seqname{A}^i, \seqname{E}^i, \seqname{C}^i)
\defeq
(\seqname{A}^i, E^i, \seqname{C}^i)
\end{equation}

Let $\morphName{f}_0 \maps \seqname{E}^1 \to \seqname{E}^2$ and
$\morphName{f}_1 \maps \seqname{C}^1 \to \seqname{C}^2$ be model
functions. Then
\begin{multline}
\Functor_{\M,\Top}
  \bigl (
    (
      \morphName{f}_0 \maps \seqname{E}^1 \to \seqname{E}^2,
      \morphName{f}_1 \maps \seqname{C}^1 \to \seqname{C}^2,
    ),
    (\seqname{A}^1, \seqname{E}^1, \seqname{C}^1),
    (\seqname{A}^2, \seqname{E}^2, \seqname{C}^2)
  \bigr )
\defeq \\
  \bigl (
    (
      \morphName{f}_0 \maps E^1 \to E^2,
      \morphName{f}_1 \maps \seqname{C}^1 \to \seqname{C}^2
    ),
    ( \seqname{A}^1, E^1, \seqname{C}^1 ),
    ( \seqname{A}^2, E^2, \seqname{C}^2 )
  \bigr )
\end{multline}
\end{definition}

\begin{lemma}[$\Functor_{\M,\Top}$ maps M-atlases to M-atlases if the coordinate model space is fine grained.]
\label{lem:Func-M-Top}
Let $\seqname{E} \defeq (E,\catName{E})$ and $\seqname{C}$ be model
spaces, $\seqname{C}$ be fine grained and $\seqname{A}$ be an M-atlas of
$\seqname{E}$ in the model space $\seqname{C}$. Then
$\seqname{A}$ is an M-atlas of $E$ in the model space $\seqname{C}$.

\begin{proof}
Let $(U, V, \phi)$ be an M-chart in $\seqname{A}$, Every open subset of
$V$ is a model neighborhood of $\seqname{C}$ since it is fine grained by
hypothesis. Then $(U, V, \phi)$ is an M-chart of $E$ in the model space
$\seqname{C}$ by \pagecref{lem:M-chart}.

The definition of mutually compatible charts
neighborhood of $\seqname{C}$. $(U, V, \phi)$ is also an M-chart of $E$
in the coordinate space $\seqname{C}$:

Since $U$ is a model neighborhood of $\seqname{E}$, it is open and hence
a model neighborhood of $\Triv{E}$.

$V$ is a model neighborhood of $\seqname{C}$ by hypothesis.

Since $U'$ is a model neighborhood of $\Triv{E}$ and hence open, and
$\phi$ is a homeomorphism, $\phi[U']$ is open.
Since $V$ is a model neighborhood of $\seqname{C}$,
$\seqname{C}$ is fine grained and $\phi[U'] \subseteq V$ is open,
$\phi[U']$ is a model neighborhood of $\seqname{C}$,

Since $V'$ is a model neighborhood of $\seqname{C}$, it is open,
$\phi^{-1}[V']$ is open and hence a model neighborhood of $\Triv{E}$.
\end{proof}
\end{lemma}

\begin{theorem}[$\Functor_{\M,\Top}$ is a functor if $\seqname{C}$ is fine grained.]
\label{the:Func-M-Top}
Let $\seqname{E}$ and $\seqname{C}$ be sets of model spaces and
$\seqname{C}$ be fine grained. Then $\Functor_{\M,\Top}$ is a functor
from $\Atl(\seqname{E}, \seqname{C})$ to
$\Atl \bigl ( \Top(\seqname{E}), \seqname{C} \bigr )$

\begin{proof}
Let
$
  o^i \defeq
  \bigl (
    \seqname{A}^i, \seqname{E}^i \defeq (E^i,\catName{E^i}), \seqname{C}^i
  \bigr )
$,
$i \in [1,3]$, be an object of $\Atl(\seqname{E}, \seqname{C})$,
$
  o'^i \defeq \Functor_{\M,\Top} o^i =
  (\seqname{A}^i, E^i, \seqname{C}^i)
$
the corresponding object of $\Atl \bigl ( \Top(\seqname{E}), \seqname{C} \bigr )$,
$
  m^i \defeq \\*
  \bigl (
    (
      \morphName{f}_0^i \maps E^i \to E^{i+1},
      \morphName{f}_1^i \maps \seqname{C}^i \to \seqname{C}^{i+1},
    ),
    o^i,
    o^{i+1}
  \bigr )
$, $i=1,2$,
be a morphism of $\Atl(\seqname{E}, \seqname{C})$
and
$
  m'^i \defeq \Functor_{\M,\Top} m^i =
  \bigl (
    (
      \morphName{f}_0^i \maps E^i \to E^{i+1},
      \morphName{f}_1^i \maps \seqname{C}^i \to \seqname{C}^{i+1}
    ),
    o'^i,
    o'^{i+1}
  \bigr )
$
the corresponding morphism of $\Atl \bigl ( \Top(\seqname{E}), \seqname{C} \bigr )$.

$\Functor_{\M,\Top}$ satisfies these criteria:
\begin{description}
\item[Preservation of objects:]
\leavevmode \\*
$\seqname{A}^i$ is an M-atlas of $E^i$ in the coordinate space
$\seqname{C}^i$ by \fullcref{lem:Func-M-Top} above.

\item[Preservation of morphisms:]
\leavevmode \\*
$\morphName{f}_0^i \maps E^i \to E^{i+1}$ is continuous and
$\morphName{f}_1^i \maps \seqname{C}^i \to \seqname{C}^{i+1}$ is a model
function.

\item[Preservation of endpoints:]
\begin{align*}
  \Functor_{\M,\Top} o^i
  & =
  o'^i \objin \Atl \bigl ( \Top(\seqname{E}), \seqname{C} \bigr )
\\*
  \Functor_{\M,\Top} m^i
  & =
  \bigl (
    (
      \morphName{f}_0^i \maps E^i \to E^{i+1},
      \morphName{f}_1^i \maps \seqname{C}^i \to \seqname{C}^{i+1}
    ),
    o'^i,
    o'^{i+1}
  \bigr )
  \arin \Atl \bigl ( \Top(\seqname{E}), \seqname{C} \bigr )
\\*
  & \quad
\text{is a morphism from $o'^i$ to $o'^{i+1}$.}
\end{align*}

\item[Identity:]
\leavevmode \\*
Let
$
  \bigl (
    (\Id[{\seqname{E}^i}], \Id[{\seqname{C}^i}]),
   o^i,
   o^i
  \bigr )
$
be an identity morphism of
$o^i$ in $\Atl(\seqname{E}, \seqname{C})$.
Then
$
  \bigl (
    (
      \Id[E^i] \maps \Top(\seqname{E}^i) \to \Top(\seqname{E}^i),
      \Id[{\seqname{C}^i}]
    ),
    o'^i,
    o'^i
  \bigr )
$
is an identity morphism of $o'^i$ in $\Atl \bigl ( \Top(\seqname{E}), \seqname{C} \bigr )$.

\item[Composition:]
\begin{equation*}
\begin{array}{lll}
  \Functor_{\M,\Top} (m^2 \morphCompose[A] m^1)
  & =
  &
  \Functor_{\M,\Top}
    \bigl (
      (
        \morphName{f}_0^2 \morphCompose \morphName{f}_0^1,
        \morphName{f}_1^2 \morphCompose \morphName{f}_1^1
      ),
      o^1,
      o^3
    \bigr )
\\*
  & =
  &
    \bigl (
      (
        \morphName{f}_0^2 \morphCompose \morphName{f}_0^1,
        \morphName{f}_1^2 \morphCompose \morphName{f}_1^1
      ),
      o'^1,
      o'^3
    \bigr )
\\*
  & =
  &
    \bigl (
      (\morphName{f}_0^2, \morphName{f}_1^2),
      o'^1,
      o'^2
    \bigr )
    \morphCompose[A]
    \bigl (
      (\morphName{f}_0^1, \morphName{f}_1^1),
      o'^2,
      o'^3
    \bigr )
\\*
  & =
  &
  \Functor_{\M,\Top}
    \bigl (
      (\morphName{f}_0^2, \morphName{f}_1^2),
      o^2,
      o^3
    \bigr )
  \morphCompose[A]
\\*
  &
  &
    \Functor_{\M,\Top}
    \bigl (
      (\morphName{f}_0^1, \morphName{f}_1^1),
      o^1,
      o^2
    \bigr )
\\*
  & =
  &
  \Functor_{\M,\Top} m^2
  \morphCompose[A]
  \Functor_{\M,\Top} m^1
\end{array}
\end{equation*}
\end{description}
\end{proof}
\end{theorem}

\begin{definition}[$\Functor_{\Top,\M}$]
\label{def:Func-Top-M}
Let $E^i$, $i=1,2$, be a topological space, $\seqname{C}^i$ a model
space and $\seqname{A}^i$ an m-atlas of $E^i$
and $\seqname{A}$ an m-atlas of $E$ in the coordinate space
$\seqname{C}^i$. Then
\begin{equation}
\Functor_{\Top,\M} (\seqname{A}^i, E^i, \seqname{C}^i)
\defeq
\Big ( \seqname{A}, \Triv{E^i}, \seqname{C}^i \Big )
\end{equation}

Let $\morphName{f}_0 \maps E^1 \to E^2$ be continuous and
$\morphName{f}_1 \maps \seqname{C}^1 \to \seqname{C}^2$ be a model
function. Then
\begin{multline}
\Functor_{\Top,\M}
  \bigl (
    (
      \morphName{f}_0 \maps E^1 \to E^2,
      \morphName{f}_1 \maps \seqname{C}^1 \to \seqname{C}^2,
    ),
    (\seqname{A}^1, E^1, \seqname{C}^1),
    (\seqname{A}^2, E^2, \seqname{C}^2)
  \bigr )
\defeq \\
  \biggl (
    \Bigl (
      \morphName{f}_0 \maps \Triv{E^1} \to \Triv{E^2},
      \morphName{f}_1 \maps \seqname{C}^1 \to \seqname{C}^2
     \Bigr ),
    \Bigl ( \seqname{A}^1, \Triv{E^1}, \seqname{C}^1 \Bigr ),
    \Bigl ( \seqname{A}^2, \Triv{E^2}, \seqname{C}^2 \Bigr )
  \biggr )
\end{multline}
\end{definition}

\begin{lemma}[$\Functor_{\Top,\M}$ maps M-atlases to M-atlases]
\label{lem:Func-Top-M}
Let $E$ be a topological space, $C$ be a model space and $\seqname{A}$
be an M-atlas of $E$ in the model space $\seqname{C}$. Then
$\seqname{A}$ is an M-atlas of $\Triv{E}$ in the model space
$\seqname{C}$.

\begin{proof}
Let $(U, V, \phi)$ be a M-chart in $\seqname{A}$, $U' \subseteq U$ a
model neighborhood of $\Triv{E}$ and $V' \subseteq V$ a model
neighborhood of $\seqname{C}$. $(U, V, \phi)$ is also an M-chart of $E$
in the coordinate space $\seqname{C}$:
\end{proof}
\end{lemma}

\begin{theorem}[$\Functor_{\Top,\M}$ is a functor]
\label{the:Func-Top-M}
Let $\seqname{E}$ be a set of topological spaces and $\seqname{C}$ a set
of model spaces, Then $\Functor_{\Top,\M}$ is a functor from
$\Atl(\seqname{E}, \seqname{C})$ to
$\Atl \big ( \Triv{\seqname{E}}, \seqname{C} \big )$.

\begin{proof}
Let $o^i \defeq (\seqname{A}^i, E^i, \seqname{C}^i)$, $i \in [1,3]$,
be an object of $\Atl(\seqname{E}, \seqname{C})$,
$o'^i \defeq (\seqname{A}^i, \Triv{E^i}, \seqname{C}^i)$
the corresponding object of
$\Atl \big ( \Triv{\seqname{E}}, \seqname{C} \big )$,
$
 m^i \defeq
  \bigl (
    (f_0^i \maps  E^i \to E^{i+1}, f_1^i \maps  \seqname{C}^i \to  \seqname{C}^{i+1}),
   o^i ,
   o^{i+1}
 \bigr )
$,
$i=1,2$, be a morphism of $\Atl(\seqname{E}, \seqname{C})$ and
$
  m'^i \defeq
  \bigl (
    (f_0^i, f_1^i),
    o'^i ,
   o'^{i+1}
  \bigr )
$
the corresponding morphism of
$\Atl \big ( \Triv{\seqname{E}}, \seqname{C} \big )$.

$\Functor_{\Top,\M}$ satisfies these criteria:
\begin{description}
\item[Preservation of endpoints:]
$\seqname{A}^i$ is an M-atlas of $\Triv{E^i}$ in the coordinate space
$\seqname{C}^i$ by \fullcref{lem:Func-M-Top} above.
$\morphName{f}_0^i \maps E^i \to E^{i+1}$ is continuous and
$\morphName{f}_1^i \maps \seqname{C}^i \to \seqname{C}^{i+1}$ is a model
function.

\begin{align*}
  \Functor_{\Top,\M} o^i
  & =
  o'^i \objin \Atl \big ( \Triv{\seqname{E}}, \seqname{C} \big )
\\*
  \Functor_{\Top,\M} m^i
  & =
  m'^i \arin \Atl \big ( \Triv{\seqname{E}}, \seqname{C} \big )
\\*
  & \quad
\text{is a morphism from $o'^i$ to $o'^{i+1}$.}
\end{align*}

\item[Identity:]
$
  \Id[o^i] \defeq
  \bigl (
    (\Id[E^i], \Id[\seqname{C}^i]),
    o^i,
    o^i
  \bigr )
$.
Then
$
  \Functor_{\Top,\M} \Id[o^i] =
  \bigl (
    (\Id[{\Triv{E^i}}], \Id[{\seqname{C}^i}]),
    o'^i,
    o'^i
  \bigr )
$
is an identity morphism of
$o'^i$.

\item[Composition:]
\begin{equation*}
\begin{array}{lll}
  \Functor_{\Top,\M} (m^2 \morphCompose[A] m^1)
  & =
  &
    \Functor_{\Top,\M}
      \bigl (
        (f_0^2 \morphCompose f_0^1, f_1^2 \morphCompose f_1^1),
        o^1,
        o^3
      \bigr )
\\*
  & =
  &
      \bigl (
        (f_0^2 \morphCompose f_0^1, f_1^2 \morphCompose f_1^1),
        o'^1,
        o'^3
      \bigr )
\\*
  & =
  &
      \bigl (
        (f_0^2, f_1^2),
        o'^1,
        o'^2
      \bigr )
      \morphCompose[A]
      \bigl (
        (f_0^1, f_1^1),
        o'^2,
        o'^3
      \bigr )
\\*
  & =
  &
    \Functor_{\Top,\M}
      \bigl (
        (f_0^2, f_1^2),
        o^2,
        o^3
      \bigr )
    \morphCompose[A]
\\*
  &
  &
    \Functor_{\Top,\M}
      \bigl (
        (f_0^1, f_1^1),
        o^1,
        o^2
      \bigr )
\\*
  & =
  &
  \Functor_{\Top,\M} m^2
  \morphCompose[A]
  \Functor_{\Top,\M} m^1
\end{array}
\end{equation*}
\end{description}
\end{proof}

$\Functor_{\M,\Top} \morphCompose \Functor_{\Top,\M}$ is the identity functor
on $\Atl(\seqname{E}, \seqname{C})$.

\begin{proof}
\begin{align*}
  \Functor_{\M,\Top} \morphCompose \Functor_{\Top,\M} o^i
  & =
  \Functor_{\M,\Top} o'^i
\\
  & =
   o^i
\end{align*}

\begin{align*}
  \Functor_{\M,\Top} \morphCompose \Functor_{\Top,\M} m^i
  & =
  \Functor_{\M,\Top} m'^i
\\
  & =
  m^i
\end{align*}

\end{proof}
\end{theorem}

\section{Classic m-atlas morphisms and functors}
\label{sec:M-ATLcmorph}
This subsection introduces an alternate definition of and taxonomy for
morphisms between m-atlases, defines categories of m-atlases, defines
functors amomg them and proves some basic reasults.  It introduces the
notion of classic m-atlas morphisms, which is very similar to
the conventional definitions for a manifold, although it does not
require the space to be locally Euclidean. The reader can safely skip
it, but it does help to establish a context for other notions that are
used in the exposition of local coordinate spaces.

\subsection{Classic m-atlas morphisms}
\label{sub:M-ATLcmorph}
This subsection introduces the notion of classic m-atlas morphisms, and
proves some basic results.

\subsubsection{Definitions of classic m-atlas morphisms}
\begin{definition}[Classic m-atlas morphisms for model spaces]
\label{def:M-ATLcmorph}
Let $\catName{E}^i$, $\catName{C}^i$, $i=1,2$, be model categories,
$\seqname{E}^i \objin \catName{E}^i$,
$\seqname{C}^i \objin \catName{C}^i$
and $\seqname{A}^i$ be an m-atlas
of $\seqname{E}^i$ in the coordinate space $\seqname{C}^i$.

$\morphName{f} \maps \seqname{E}^1 \to \seqname{E}^2$ is a classic
\nobreaks
  {
    {$\seqname{E}^1$-$\seqname{E}^2$},
    {(-$\catName{C}^1$-$\catName{C}^2$)}
  }
m-atlas morphism of $\seqname{A}^1$ to $\seqname{A}^2$ in the coordinate
spaces $\seqname{C}^1$, $\seqname{C}^2$, abbreviated as
$
  \isAtl[classic]_\Ar
  (
    \seqname{A}^1,
    \seqname{E}^1,
    \seqname{C}^1,
    \seqname{A}^2,
    \seqname{E}^2,
    \seqname{C}^2,
    \morphName{f}
  )
$.
and a classic
\nobreaks
 {
   {$\catName{E}^1$-$\catName{E}^2$},
   {(-$\catName{C}^1$-$\catName{C}^2$)}
  }
m-atlas morphism of $\seqname{A}^1$ to $\seqname{A}^2$ in the coordinate
model categories $\catName{C}^1$, $\catName{C}^2$, abbreviated as
$
  \isAtl[classic]_\Ar
  (
    \seqname{A}^1,
    \catName{E}^1,
    \catName{C}^1,
    \seqname{A}^2,
    \catName{E}^2,
    \catName{C}^2,
    \morphName{f}
  )
$,
iff
\nobreaks
  {{
    $\seqname{C}^1 \submod \seqname{C}^2$
  }}
($\catName{C}^1 \subcat \catName{C}^2)$
and for any $(U^i, V^i, \phi^i \maps U^i \toiso V^i) \in \seqname{A}^i$,
$i=1,2$, with $I \defeq U^1 \cap \morphName{f}^{-1}[U^2] \neq \emptyset$,
$
  \phi^2 \morphCompose \morphName{f} \morphCompose \phi^{1-1} \maps
    \phi^1[I] \to V^2
$
is a morphism of $\seqname{C}^2$ ($\catName{C}^2$).

$\morphName{f}$ is also a constrained classic
\nobreaks
  {
    {$\seqname{E}^1$-$\seqname{E}^2$},
    {(-$\catName{C}^1$-$\catName{C}^2$)}
  }
m-atlas
morphism of $\seqname{A}^1$ to $\seqname{A}^2$ in the coordinate spaces
$\seqname{C}^1$, $\seqname{C}^2$, abbreviated as \\*
$
  \isAtl[constrained,classic]_\Ar
    (
      \seqname{A}^1,
      E^1,
      \seqname{C}^1,
      \seqname{A}^2,
      E^2,
      \seqname{C}^2,
      \morphName{f}_0,
      \morphName{f}_1
    )
$,
and a constrained classic
\nobreaks
  {
    {$\catName{E}^1$-$\catName{E}^2$},
    {(-$\catName{C}^1$-$\catName{C}^2$)}
  }
m-atlas morphism of
$\seqname{A}^1$ to $\seqname{A}^2$ in the coordinate model categories
$\catName{C}^1$, $\catName{C}^2$, abbreviated as
$
  \isAtl[constrained,classic]_\Ar
    (
      \seqname{A}^1,
      E^1,
      \catName{C}^1,
      \seqname{A}^2,
      E^2,
      \catName{C}^2,
      \morphName{f}_0,
      \morphName{f}_1
    )
$,
iff $\morphName{f}$ is constrained, i.e.,
$\morphName{f}[\seqname{E}]$ is contained in a model neighborhood of
$\seqname{E}^2$.

Let $\morphName{f}$ be a classic
\nobreaks
  {
    {$\seqname{E}^1$-$\seqname{E}^2$},
    {(-$\catName{C}^1$-$\catName{C}^2$)}
  }
m-atlas morphism of $\seqname{A}^1$ to $\seqname{A}^2$ in the
coordinate spaces $\seqname{C}^1$, $\seqname{C}^2$.

The triple
$
  \Bigl (
    (\morphName{f}),
    (\seqname{A}^1, \seqname{E}^1, \seqname{C}^1),
    (\seqname{A}^2, \seqname{E}^2, \seqname{C}^2)
  \Bigr )
$
will refer to $\morphName{f}$ considered as a classic
\nobreaks
  {
    {$\seqname{E}^1$-$\seqname{E}^2$},
    {(-$\catName{C}^1$-$\catName{C}^2$)}
  }
and
\nobreaks
  {
    {$\catName{E}^1$-$\catName{E}^2$},
    {(-$\catName{C}^1$-$\catName{C}^2$)}
  }
m-atlas morphism of $\seqname{A}^1$ to $\seqname{A}^2$ in the coordinate
spaces $\seqname{C}^1$, $\seqname{C}^2$.
\end{definition}

\begin{definition}[Classic m-atlas morphisms for topological spaces]
\label{def:M-ATLcmorphtop}
Let $\catName{C}^i$, $i=1,2$, be a model category, $\catName{E}^i$ be a
topological category, $E^i \objin \catName{E}^i$,
$\seqname{C}^i \objin \catName{C}^i$ and $\seqname{A}^i$ be an m-atlas
of $E^i$ in the coordinate space $\seqname{C}^i$.

$\morphName{f} \maps E^1 \to E^2$ is a classic
\nobreaks
  {
    {$E^1$-$E^2$},
    {(-$\catName{C}^1$-$\catName{C}^2$)}
  }
m-atlas morphism of $\seqname{A}^1$ to $\seqname{A}^2$ in the coordinate
spaces $\seqname{C}^1$, $\seqname{C}^2$, abbreviated as
$
  \isAtl[classic]_\Ar
  (
    \seqname{A}^1,
    E^1,
    \seqname{C}^1,
    \seqname{A}^2,
    E^2,
    \seqname{C}^2,
    \morphName{f}
  )
$.
and a classic
\nobreaks
  {
    {$\catName{E}^1$-$\catName{E}^2$},
    {(-$\catName{C}^1$-$\catName{C}^2$)}
  }
m-atlas morphism of $\seqname{A}^1$ to $\seqname{A}^2$ in the coordinate
model categories $\catName{C}^1$, $\catName{C}^2$, abbreviated as
$
  \isAtl[classic]_\Ar
  (
    \seqname{A}^1,
    \catName{E}^1,
    \catName{C}^1,
    \seqname{A}^2,
    \catName{E}^2,
    \catName{C}^2,
    \morphName{f}
  )
$,
iff $\morphName{f}$ is a classic
\nobreaks
  {
    {$\Triv{E^1}$-$\Triv{E^2}$},
    {(-$\catName{C}^1$-$\catName{C}^2$)}
  }
m-atlas
morphism of $\seqname{A}^1$ to $\seqname{A}^2$ in the coordinate spaces
$\seqname{C}^1$, $\seqname{C}^2$.

The triple
$
  \Bigl (
    (\morphName{f}),
    (\seqname{A}^1, E^1, \seqname{C}^1),
    (\seqname{A}^2, E^2, \seqname{C}^2)
  \Bigr )
$
will refer to $\morphName{f}$ considered as a classic
\nobreaks
  {
    {$E^1$-$E^2$}%
    {(-$\catName{C}^1$-$\catName{C}^2$)}
  }
and
\nobreaks
  {
    {$\catName{E}^1$-$\catName{E}^2$}%
    {(-$\catName{C}^1$-$\catName{C}^2$)}
  }
m-atlas morphism of $\seqname{A}^1$ to $\seqname{A}^2$ in the coordinate
spaces $\seqname{C}^1$, $\seqname{C}^2$.
\end{definition}

\begin{definition}[Classic m-atlas inclusion and identity morphisms]
\label{def:M-ATLcid}
Let
\begin{enumerate}
\item
$\catName{E}^i$, $\catName{C}^i$, $i=1,2$, be model categories
\item
$\seqname{E}^i \objin \catName{E}^i$
\item
$\seqname{E}^1 \submod \seqname{E}^2$
\item
$\seqname{C}^i \objin \catName{C}^i$
\item
\nobreaks
  {
    {$\seqname{C}^1 \submod \seqname{C}^2$},
    {($\catName{C}^1 \subcat \catName{C}^2$)}
  }
\item
$\seqname{A}^i$ be an m-atlas of $\seqname{E}^i$ in the coordinate
space $\seqname{C}^i$
\item
$\seqname{A}^1 \subseteq \seqname{A}^2$
\end{enumerate}

The classic
\nobreaks
  {
    {$\seqname{E}^1$-$\seqname{E}^2$},
    {(-$\catName{C}^1$-$\catName{C}^2$)}
  }
m-atlas inclusion morphism of $\seqname{A}^1$ to $\seqname{A}^2$ in the
coordinate spaces $\seqname{C}^1$, $\seqname{C}^2$ is
$
  \Id%
    [
      {(\seqname{E}^1, \seqname{C}^1)},
      {(\seqname{E}^2, \seqname{C}^2)}
    ]
  \defeq
  \Id%
    [
      {\seqname{E}^1},
      {\seqname{E}^2}
    ]
$.

The classic m-atlas inclusion morphism of
$(\seqname{A}^1, \seqname{E}^1, \seqname{C}^1)$ to
$(\seqname{A}^2, \seqname{E}^2, \seqname{C}^2)$ is \\*
$
  \Id%
    [
      {(\seqname{A}^1, \seqname{E}^1, \seqname{C}^1)},
      {(\seqname{A}^2, \seqname{E}^2, \seqname{C}^2)}
    ]
  \defeq
  \biggl (
    \ID%
    [
      {(\seqname{E}^1)},
      {(\seqname{E}^2)}
    ],
    (\seqname{A}^1, \seqname{E}^1, \seqname{C}^1),
    (\seqname{A}^2, \seqname{E}^2, \seqname{C}^2)
  \biggr )
$,
$
 \Id%
   [
     {\seqname{E}^1},
     {\seqname{E}^2}
   ]
$
considered as a classic
\nobreaks
  {
    {$\seqname{E}^1$-$\seqname{E}^2$},
    {(-$\catName{C}^1$-$\catName{C}^2$)}
  }
m-atlas morphism of $\seqname{A}^1$ to $\seqname{A}^2$ in the coordinate
spaces $\seqname{C}^1$, $\seqname{C}^2$.

The classic
\nobreaks
  {
    {$\seqname{E}^i$},
    {(-$\catName{C}^i$)}
  }
m-atlas identity morphism of $\seqname{A}^i$
in the coordinate space $\seqname{C}^i$ is
$
  \Id[{{(\seqname{E}^i, \seqname{C}^i)}}] \defeq
  \Id%
    [
      {\seqname{E}^i}
    ]
$.

The classic m-atlas identity morphism of
$(\seqname{A}^i, \seqname{E}^i, \seqname{C}^i)$ is \\*
$
  \Id[{{(\seqname{A}^i, \seqname{E}^i, \seqname{C}^i)}}] \defeq
  \bigl (
    \ID[{{(\seqname{E}^i)}}],
    (\seqname{A}^i, \seqname{E}^i, \seqname{C}^i),
    (\seqname{A}^i, \seqname{E}^i, \seqname{C}^i)
  \bigr )
$,\hspace{1pt}
$
 \Id%
   [
     {\seqname{E}^i}
   ]
$
considered as a classic
\nobreaks
  {
    {$\seqname{E}^i$},
    {(-$\catName{C}^i$)}
  }
m-atlas morphism of
$\seqname{A}^i$ in the coordinate space $\seqname{C}^i$.

Let
\begin{enumerate}
\item
$E^i$ be a topological space, $i=1,2$,
\item
$E^1 \subsp E^2$
\item
$\catName{C}^i$ be a model category
\item
$\seqname{C}^i \objin \catName{C}^i$
\item
$\catName{C}^1 \subcat \catName{C}^2$
\item
$\seqname{A}^i$ be an m-atlas of $E^i$ in the coordinate space
$\seqname{C}^i$
\item
$\seqname{A}^1 \subseteq \seqname{A}^2$
\end{enumerate}

The classic
\nobreaks
  {
    {$E^1$-$E^2$},
    {(-$\catName{C}^1$-$\catName{C}^2$)}
  }
m-atlas inclusion morphism of
$\seqname{A}^1$ to $\seqname{A}^2$ in the coordinate spaces
$\seqname{C}^1$, $\seqname{C}^2$ is
$
 \Id%
   [
     {E^1},
     {E^2}
   ]
$.

The classic m-atlas inclusion morphism of
$(\seqname{A}^1, E^1, \seqname{C}^1)$ to
$(\seqname{A}^2, E^2, \seqname{C}^2)$ is \\*
$
  \Id%
    [
        {(\seqname{A}^1, E^1, \seqname{C}^1)},
        {(\seqname{A}^2, E^2, \seqname{C}^2)}
    ]
  \defeq
  \biggl (
    \ID%
    [
      {(E^1)},
      {(E^2)}
    ],
    (\seqname{A}^1, E^1, \seqname{C}^1),
    (\seqname{A}^2, E^2, \seqname{C}^2)
  \biggr )
$,
$
 \Id%
   [
     {E^1},
     {E^2}
   ]
$
considered as a classic $E^1$-$E^2$ m-atlas morphism of
$\seqname{A}^1$ to $\seqname{A}^2$ in the coordinate spaces
$\seqname{C}^1$, $\seqname{C}^2$.

The classic $E^i$ m-atlas identity morphism of $\seqname{A}^i$ in the
coordinate space $\seqname{C}^i$ is
$
 \Id%
   [
     {(E^i, \seqname{C}^i)}
   ]
$.

The classic m-atlas identity morphism of $(\seqname{A}^i, E^i, \seqname{C}^i)$ is \\*
$
  \Id[{{(\seqname{A}^i, E^i, \seqname{C}^i)}}] \defeq
  \bigl (
    \ID[{{(E^i)}}],
    (\seqname{A}^i, E^i, \seqname{C}^i),
    (\seqname{A}^i, E^i, \seqname{C}^i)
  \bigr )
$,
$
 \Id%
   [
     {E^i}
   ]
$
considered as an $E^i$ m-atlas morphism of $\seqname{A}^i$ in
the coordinate space $\seqname{C}^i$.
\end{definition}

\subsubsection{Definitions of semi-strict and strict}
\begin{definition}[Semi-strict classic m-atlas morphisms]
\label{def:M-ATLcmorphsemi}
Let $\catName{E}^i$, $\catName{C}^i$, $i=1,2$, be model categories,
$\seqname{E}^i \objin \catName{E}^i$,
$\seqname{C}^i \objin \catName{C}^i$ and $\seqname{A}^i$ be an m-atlas
of $\seqname{E}^i$ in the coordinate space $\seqname{C}^i$.

Let $\morphName{f} \maps \seqname{E}^1 \to \seqname{E}^2$ be a classic
$\seqname{E}^1$-$\seqname{E}^2$ m-atlas morphism of
$\seqname{A}^1$ to $\seqname{A}^2$ in the coordinate spaces
$\seqname{C}^1$, $\seqname{C}^2$.

$\morphName{f}$ is a semi-strict classic $\seqname{E}^1$-$\seqname{E}^2$
m-atlas morphism of $\seqname{A}^1$ to $\seqname{A}^2$ in the
coordinate spaces $\seqname{C}^1$, $\seqname{C}^2$, abbreviated as
$
  \strict[semi-]{\isAtl[classic]_\Ar}
  (
    \seqname{A}^1,
    \seqname{E}^1,
    \seqname{C}^1,
    \seqname{A}^2,
    \seqname{E}^2,
    \seqname{C}^2,
    \morphName{f}
  )
$,
iff $\morphName{f}$ is a local m-morphism of $\seqname{E}^1$ to
$\seqname{E}^2$.

\begin{remark}
Definitions of semi-strict classic m-atlas morphisms for
topological spaces would be pointless, as any semi-strict classic
m-atlas morphism would be strict.
\end{remark}

It is a semi-strict classic
$\catName{E}^1$-$\catName{E}^2$-$\seqname{E}^1$-$\seqname{E}^2$ m-atlas
morphism of $\seqname{A}^1$ to $\seqname{A}^2$ in the coordinate
spaces $\seqname{C}^1$, $\seqname{C}^2$, abbreviated as \\*
$
  \strict[semi-]{\isAtl[classic]_\Ar}
  (
    \seqname{A}^1,
    \seqname{E}^1,
    \seqname{C}^1,
    \seqname{A}^2,
    \seqname{E}^2,
    \seqname{C}^2,
    \morphName{f},
    \catName{E}^1,
    \catName{E}^2
  )
$,
iff $\morphName{f}$ is a local $\catName{E}^1$-$\catName{E}^2$
m-morphism of $\seqname{E}^1$ to $\seqname{E}^2$.

Let $\morphName{f} \maps \seqname{E}^1 \to \seqname{E}^2$ be a classic
$\seqname{E}^1$-$\seqname{E}^2$-$\catName{C}^1$-$\catName{C}^2$ m-atlas
morphism of $\seqname{A}^1$ to $\seqname{A}^2$ in the coordinate
spaces $\seqname{C}^1$, $\seqname{C}^2$.

$\morphName{f}$ is a semi-strict classic
$\seqname{E}^1$-$\seqname{E}^2$-$\catName{C}^1$-$\catName{C}^2$
m-atlas morphism of $\seqname{A}^1$ to $\seqname{A}^2$ in the
coordinate spaces $\seqname{C}^1$, $\seqname{C}^2$, abbreviated as \\*
$
  \strict[semi-]{\isAtl[classic]_\Ar}
  (
    \seqname{A}^1,
    \seqname{E}^1,
    \seqname{C}^1,
    \seqname{A}^2,
    \seqname{E}^2,
    \seqname{C}^2,
    \morphName{f},
    \catName{C}^1,
    \catName{C}^2
  )
$,
iff $\morphName{f}$ is a local m-morphism of $\seqname{E}^1$ to
$\seqname{E}^2$.

It is a semi-strict classic
$\catName{E}^1$-$\catName{E}^2$-%
$\seqname{E}^1$-$\seqname{E}^2$-%
$\catName{C}^1$-$\catName{C}^2$
m-atlas morphism of $\seqname{A}^1$ to $\seqname{A}^2$ in the coordinate
spaces $\seqname{C}^1$, $\seqname{C}^2$, abbreviated as \\*
$
  \strict[semi-]{\isAtl[classic]_\Ar}
  (
    \seqname{A}^1,
    \seqname{E}^1,
    \seqname{C}^1,
    \seqname{A}^2,
    \seqname{E}^2,
    \seqname{C}^2,
    \morphName{f},
    \catName{E}^1,
    \catName{E}^2,
    \catName{C}^1,
    \catName{C}^2
  )
$,
iff $\morphName{f}$ is a local $\catName{E}^1$-$\catName{E}^2$
m-morphism of $\seqname{E}^1$ to $\seqname{E}^2$.
\end{definition}

\begin{definition}[Strict classic m-atlas morphisms]
\label{def:M-ATLcmorphstrict}
Let $\catName{E}^i$, $\catName{C}^i$, $i=1,2$, be model categories,
$\seqname{E}^i \objin \catName{E}^i$,
$\seqname{C}^i \objin \catName{C}^i$, $\seqname{A}^i$ be an m-atlas of
$\seqname{E}^i$ in the coordinate space $\seqname{C}^i$.

Let $\morphName{f} \maps \seqname{E}^1 \to \seqname{E}^2$ be a classic
$\seqname{E}^1$-$\seqname{E}^2$ m-atlas morphism of
$\seqname{A}^1$ to $\seqname{A}^2$ in the coordinate spaces
$\seqname{C}^1$, $\seqname{C}^2$.

$\morphName{f}$ is a strict classic $\seqname{E}^1$-$\seqname{E}^2$
m-atlas morphism of $\seqname{A}^1$ to $\seqname{A}^2$ in the
coordinate spaces $\seqname{C}^1$, $\seqname{C}^2$, abbreviated as
$
  \strict{\isAtl[classic]_\Ar}
  (
    \seqname{A}^1,
    \seqname{E}^1,
    \seqname{C}^1,
    \seqname{A}^2,
    \seqname{E}^2,
    \seqname{C}^2,
    \morphName{f}
  )
$,
iff $\morphName{f}$ is an m-morphism of $\seqname{E}^1$ to $\seqname{E}^2$.

It is a strict classic
$\catName{E}^1$-$\catName{E}^2$-$\seqname{E}^1$-$\seqname{E}^2$ m-atlas
morphism of $\seqname{A}^1$ to $\seqname{A}^2$ in the coordinate
spaces $\seqname{C}^1$, $\seqname{C}^2$, abbreviated as \\*
$
  \strict{\isAtl[classic]_\Ar}
  (
    \seqname{A}^1,
    \seqname{E}^1,
    \seqname{C}^1,
    \seqname{A}^2,
    \seqname{E}^2,
    \seqname{C}^2,
    \morphName{f},
    \catName{E}^1,
    \catName{E}^2
  )
$,
iff $\morphName{f}$ is an $\catName{E}^1$-$\catName{E}^2$ m-morphism of
$\seqname{E}^1$ to $\seqname{E}^2$.

Let $\morphName{f} \maps \seqname{E}^1 \to \seqname{E}^2$ be a classic
$\seqname{E}^1$-$\seqname{E}^2$-$\catName{C}^1$-$\catName{C}^2$ m-atlas
morphism of $\seqname{A}^1$ to $\seqname{A}^2$ in the coordinate
spaces $\seqname{C}^1$, $\seqname{C}^2$.

$\morphName{f}$ is a strict classic
$\seqname{E}^1$-$\seqname{E}^2$-$\catName{C}^1$-$\catName{C}^2$
m-atlas morphism of $\seqname{A}^1$ to $\seqname{A}^2$ in the coordinate
spaces $\seqname{C}^1$, $\seqname{C}^2$, abbreviated as
$
  \strict{\isAtl[classic]_\Ar}
  (
    \seqname{A}^1,
    \seqname{E}^1,
    \seqname{C}^1,
    \seqname{A}^2,
    \seqname{E}^2,
    \seqname{C}^2,
    \morphName{f},
    \catName{C}^1,
    \catName{C}^2
  )
$,
iff $\morphName{f}$ is an m-morphism of $\seqname{E}^1$ to $\seqname{E}^2$.

It is a strict classic
$\catName{E}^1$-$\catName{E}^2$-%
$\seqname{E}^1$-$\seqname{E}^2$-%
$\catName{C}^1$-$\catName{C}^2$
m-atlas morphism of $\seqname{A}^1$ to $\seqname{A}^2$ in the coordinate
spaces $\seqname{C}^1$, $\seqname{C}^2$, abbreviated as \\*
$
  \strict{\isAtl[classic]_\Ar}
  (
    \seqname{A}^1,
    \seqname{E}^1,
    \seqname{C}^1,
    \seqname{A}^2,
    \seqname{E}^2,
    \seqname{C}^2,
    \morphName{f},
    \catName{E}^1,
    \catName{E}^2,
    \catName{C}^1,
    \catName{C}^2
  )
$,
iff $\morphName{f}$ is an $\catName{E}^1$-$\catName{E}^2$
m-morphism of $\seqname{E}^1$ to $\seqname{E}^2$.

\end{definition}

\subsubsection{Abbreviated nomenclature for classic m-atlas morphisms}
\begin{definition}[Abbreviated nomenclature for classic m-atlas morphisms]
\label{def:M-ATLcmorphabbrev}
Let $\catName{E}^i$, $\catName{C}^i$, $i=1,2$, be model categories,
$\seqname{E}^i \objin \catName{E}^i$,
$\seqname{C}^i \objin \catName{C}^i$, $\seqname{A}^i$ be an m-atlas of
$\seqname{E}^i$ in the coordinate space $\seqname{C}^i$ and
$\morphName{f} \maps \seqname{E}^1 \to \seqname{E}^2$ be a classic
$\seqname{E}^1$-$\seqname{E}^2$ m-atlas morphism of
$\seqname{A}^1$ to $\seqname{A}^2$ in the coordinate spaces
$\seqname{C}^1$,
$\seqname{C}^2$.

$\morphName{f}$ is a (full) (maximal) (constrained) (semi-strict, strict)
classic $\seqname{E}^1$-$\seqname{E}^2$ m-atlas morphism of
$\seqname{A}^1$ to $\seqname{A}^2$ in the coordinate space
$\seqname{C}^1$, abbreviated as \\*
$
  \strict[*]{\maxfull[*]{\isAtl[(constrained,)classic]_\Ar}}
  (
    \seqname{A}^1,
    \seqname{E}^1,
    \seqname{C}^1,
    \seqname{A}^2,
    \seqname{E}^2,
    \morphName{f}
  )
$,
iff it is a (full) (maximal) (constrained) (semi-strict, strict) classic
$\seqname{E}^1$-$\seqname{E}^2$ m-atlas morphism of
$\seqname{A}^1$ to $\seqname{A}^2$ in the coordinate spaces
$\seqname{C}^1$, $\seqname{C}^1$.

$\morphName{f}$ is a (full) (maximal) (constrained) (semi-strict, strict)
classic
$\catName{E}^1$-$\catName{E}^2$-$\seqname{E}^1$-$\seqname{E}^2$
m-atlas morphism of $\seqname{A}^1$ to $\seqname{A}^2$ in the
coordinate space $\seqname{C}^1$, abbreviated as \\*
$
  \strict[*]{\maxfull[*]{\isAtl[(constrained,)classic]_\Ar}}
  (
    \seqname{A}^1,
    \seqname{E}^1,
    \seqname{C}^1,
    \seqname{A}^2,
    \seqname{E}^2,
    \morphName{f},
    \catName{E}^1,
    \catName{E}^2
  )
$,
iff it is a (full) (maximal) (constrained) (semi-strict, strict) classic
$\catName{E}^1$-$\catName{E}^2$-$\seqname{E}^1$-$\seqname{E}^2$
m-atlas morphism of $\seqname{A}^1$ to $\seqname{A}^2$ in the
coordinate spaces $\seqname{C}^1$, $\seqname{C}^1$.

Let $\catName{C}^i$, $i=1,2$, be a model category,
$\seqname{C}^i \objin \catName{C}^i$,
$E^i$ be a topological space,
$\seqname{A}^i$ be an m-atlas of $E^i$ in the coordinate space
$\seqname{C}^i$ and $\morphName{f} \maps E^1 \to E^2$ be a classic
$E^1$-$E^2$ m-atlas morphism of $\seqname{A}^1$ to
$\seqname{A}^2$ in the coordinate spaces $\seqname{C}^1$, $\seqname{C}^2$.

$\morphName{f}$ is a (full) (maximal) (constrained) (semi-strict, strict)
classic $E^1$-$E^2$ m-atlas morphism
of $\seqname{A}^1$ to $\seqname{A}^2$ in the coordinate spaces
$\seqname{C}^1$, $\seqname{C}^2$, abbreviated as \\*
$
  \strict[*]{\maxfull[*]{\isAtl[(constrained,)classic]_\Ar}}
  (
    \seqname{A}^1,
    E^1,
    \seqname{C}^1,
    \seqname{A}^2,
    E^2,
    \seqname{C}^2,
    \morphName{f}
  )
$,
iff it is a (full) (maximal) (constrained) (semi-strict, strict) classic
$\Triv{E}^1$-$\Triv{E}^2$ m-atlas
morphism of $\seqname{A}^1$ to $\seqname{A}^2$ in the coordinate spaces
$\seqname{C}^1$, $\seqname{C}^2$.

$\morphName{f}$ is a (full) (maximal) (constrained) (semi-strict, strict)
classic $E^1$-$E^2$ m-atlas morphism of $\seqname{A}^1$ to
$\seqname{A}^2$ in the coordinate space $\seqname{C}^1$, abbreviated
as \\*
$
  \strict[*]{\maxfull[*]{\isAtl[(constrained,)classic]_\Ar}}
  (
    \seqname{A}^1,
    E^1,
    \seqname{C}^1,
    \seqname{A}^2,
    E^2,
    \morphName{f}
  )
$,
iff it is a (full) (maximal) (constrained) (semi-strict, strict) classic
$E^1$-$E^2$ m-atlas morphism of $\seqname{A}^1$ to $\seqname{A}^2$ in
the coordinate spaces $\seqname{C}^1$, $\seqname{C}^1$.
\end{definition}

\subsubsection{Proclamations on classic m-atlas morphisms}
\begin{lemma}[Classic m-atlas morphisms]
\label{lem:M-ATLcmorph}
Let $\seqname{E}$ be a set of topological spaces,
$E^i \in \seqname{E}$, $i=1,2$,
$\seqname{C}^i$ be a model space,
$\seqname{A}^i$ be an m-atlas of $\Triv{E^i}$ in the coordinate space
$\seqname{C}^i$.

\begin{enumerate}
\item
Let $\morphName{f} \maps \Triv{E^1} \to \Triv{E^2}$ be a classic
$\Triv{E^1}$-$\Triv{E^2}$ m-atlas morphism of $\seqname{A}^1$ to
$\seqname{A}^2$ in the coordinate spaces $\seqname{C}^1$,
$\seqname{C}^2$ and $\seqname{C}^1 \submod \seqname{C}^2$. Then
$\morphName{f}$ is a strict constrained classic
\nobreaks
  {
    {$\Trivcat{\seqname{E}}$-$\Trivcat{\seqname{E}}$-},
    {$\Triv{E^1}$-$\Triv{E^2}$}
  }
m-atlas morphism of $\seqname{A}^1$ to $\seqname{A}^2$ in the coordinate
spaces $\seqname{C}^1$, $\seqname{C}^2$.

\begin{proof}
\leavevmode
\begin{description}
\item[$\morphName{f}$ is constrained:]
\leavevmode
$\morphName{f}[E^1] \subseteq E^2$ by hypothesis.
$E^2$ is a model neighborhood of $\Triv{E^2}$ by \pagecref{def:trivmod}.
\item[$\morphName{f}$ is strict:]
\leavevmode
$\Trivcat{\seqname{E}} \subcat[full-] \Trivcat{\seqname{E}}$.
$\morphName{f}$ is continuous by \pagecref{def:modfunc}.
$E^i \in \seqname{E}$ by hypothesis,
$E^i \objin  \Trivcat{\seqname{E}}$ and
$\morphName{f} \arin \Trivcat{\seqname{E}}$ by \cref{def:trivmod}.
$\seqname{C}^1 \submod \seqname{C}^2$
by \pagecref{def:M-ATLcmorph}.
\end{description}
\end{proof}

\item
If $E^1$ is an open subspace of $E^2$ then
$\morphName{f} \maps \Triv{E^1} \to \Triv{E^2}$ is also a strict
constrained classic $\Triv{E^1}$-$\Triv{E^2}$ m-atlas morphism of
$\seqname{A}^1$ to $\seqname{A}^2$ in the coordinate spaces
$\seqname{C}^1$, $\seqname{C}^2$.

\begin{proof}
Each $E^i \in \seqname{E}$ by hypothesis and each
$\Triv{E^i} \objin \Trivcat{\seqname{E}}$ by \cref{def:trivmod}.

$E^1 \subsp[open-] E^2$ by hypothesis.

$\Triv{E^1} \submod \Triv{E^2}$ by \pagecref{lem:trivmod}.

$\morphName{f} \arin \Trivcat{\seqname{E}}$ by \cref{def:trivmod}.
\end{proof}
\end{enumerate}

Let
\begin{enumerate}
\item
$\catName{E}^i$, $\catName{C}^i$,
$\catName{E}'^i$, $\catName{C}'^i$, $i=1,2$, be model categories
\item
$\catName{E}^i \subcat[full-] \catName{E}'^i$
\item
$\catName{C}^i \subcat[full-] \catName{C}'^i$
\item
$\catName{C}'^1 \subcat[full-] \catName{C}'^2$
\item
$\seqname{E}^i \objin \catName{E}^i$
\item
$\seqname{C}^i \objin \catName{C}^i$
\item
$\seqname{C}'^i \objin \catName{C}'^i$
\item
$\seqname{C}^i \submod \seqname{C}'^i$
\item
$\seqname{A}^i$ be an m-atlas of $\seqname{E}^i$ in the coordinate space
$\seqname{C}^i$
\item
$\morphName{f} \maps \seqname{E}^1 \to \seqname{E}^2$ be a
(semi-strict, strict)
\nobreaks
  {
    {($\catName{E}^1$-$\catName{E}^2$-)},
    {$\seqname{E}^1$-$\seqname{E}^2$},
    {(-$\catName{C}^1$-$\catName{C}^2$)}
  }
m-atlas morphism of $\seqname{A}^1$ to $\seqname{A}^2$ in the
coordinate spaces $\seqname{C}^1$, $\seqname{C}^2$
(model function?)
\end{enumerate}

If $\morphName{f}$ is a (semi-strict, strict)
$\catName{E}^1$-%
$\catName{E}^2$-%
$\seqname{E}^1$-%
$\seqname{E}^2$
(%
$\seqname{E}^1$-%
$\seqname{E}^2$%
)
classic m-atlas morphism of
$\seqname{A}^1$ to $\seqname{A}^2$ in the coordinate spaces
$\seqname{C}^1$, $\seqname{C}^2$, then $\morphName{f}$ is a
(semi-strict, strict)
$\catName{E}'^1$-%
$\catName{E}'^2$-%
$\seqname{E}^1$-%
$\seqname{E}^2$
(
$\seqname{E}^1$-%
$\seqname{E}^2$%
)
classic m-atlas morphism of $\seqname{A}^1$ to $\seqname{A}^2$ in the
coordinate spaces $\seqname{C}'^1$, $\seqname{C}'^2$.

\begin{proof}
\leavevmode
\begin{description}
\item[$\morphName{f}$ remains a classic m-atlas morphism with the expanded categories.]
$\catName{C}'^1 \subcat[full-] \catName{C}'^2$ by hypothesis.

Let $(U^i, V^i, \phi^i \maps U^i \toiso V^i) \in \seqname{A}^i$, $i=1,2$,
with $I \defeq U^1 \cap \morphName{f}^{-1}[U^2] \neq \emptyset$.
$\phi^2 \morphCompose \morphName{f} \morphCompose \phi^{-1}$ is a morphism of
$\catName{C}^2 \subcat[full-] \catName{C}'^2$ by \cref{def:M-ATLcmorph}.
\item[$\morphName{f}$ remains semi-strict (strict).]
If $\morphName{f}$ is a local morphism of $\catName{E}^2$ then
$\morphName{f}$ is a local morphism of $\catName{E}'^2$.    If
$\morphName{f}$ is a local morphism of $\seqname{E}^2$ then
$\morphName{f}$ is a local morphism of $\seqname{E}'^2$.    If
$\morphName{f} \arin \catName{E}^2$ then
$\morphName{f} \arin \catName{E}'^2$.  If $\morphName{f}$ is a
morphism of $\seqname{E}^2$ then $\morphName{f}$ is a morphism of
$\seqname{E}'^2$.
\end{description}
\end{proof}
\end{lemma}

\begin{lemma}[Composition of classic m-atlas morphisms]
\label{lem:M-ATLcmorphcomp}
Let $\catName{E}^i$, $\catName{C}^i$, $i \in [1,3]$, be a model
category,
$\seqname{E}^i \objin \catName{E}^i$,
$\seqname{C}^i \objin \catName{C}^i$, $\seqname{A}^i$ be an m-atlas
of $\seqname{E}^i$ in the coordinate space $\seqname{C}^i$ and
$\morphName{f}^i \maps \seqname{E}^i \to \seqname{E}^{i+1}$, $i=1,2$, be
a (semi-strict, strict) classic $\seqname{E}^i$-$\seqname{E}^{i+1}$
m-atlas morphism of $\seqname{A}^i$ to $\seqname{A}^{i+1}$ in the
coordinate spaces $\seqname{C}^i$, $\seqname{C}^{i+1}$.

$
  \morphName{f}^2 \morphCompose \morphName{f}^1
  \maps \seqname{E}^1 \to \seqname{E}^3
$
is a (semi-strict, strict) classic $\seqname{E}^1$-$\seqname{E}^3$
m-atlas morphism of $\seqname{A}^1$ to $\seqname{A}^3$ in the coordinate
spaces $\seqname{C}^1$, $\seqname{C}^3$.

\begin{proof}
\leavevmode
\begin{description}
\item[\ensuremath{\morphName{f}^2 \morphCompose \morphName{f}^1} is a classic $\seqname{E}^1$-$\seqname{E}^3$ m-atlas morphism.]
If $\catName{C}^1 \subcat \catName{C}^2$ and
$\catName{C}^2 \subcat \catName{C}^3$ then
$\catName{C}^1 \subcat \catName{C}^3$.
Let
$(U^1, V^1, \phi^1 \maps U^1 \toiso V^1) \in \seqname{A}^1$,
$(U^3 ,V^3, \phi^3 \maps U^3 \toiso V^3) \in \seqname{A}^3$
with
$
  I \defeq
  U^1 \cap (\morphName{f}^2 \morphCompose \morphName{f}^1)^{-1}[U^3]
  \neq \emptyset
$,
$(\morphName{f}^2 \morphCompose \morphName{f}^1)^{-1}[U^3]$ is a model
neighborhood by \pagecref{def:modfunc} and
$I$ is a model neighborhood by \cref{mod:closed} of
{
  \showlabelsinline
  \pagecref{def:model}
}.

For any $u^1 \in I$ and any chart
$(U^2 ,V^2, \phi^2 \maps U^2 \toiso V^2) \in \seqname{A}^2$ at
$\morphName{f}^1(u^1)$, define
$I^{1,2} \defeq U^1 \cap \morphName{f}^{1-1}[U^2]$,
$I^{2,3} \defeq U^2 \cap \morphName{f}^{2-1}[U^3]$ and
$I^{1,3}_{u^1} \defeq U^1 \cap \morphName{f}^{1-1}[I^{2,3}] \subseteq I$.
$I^{1,3}_{u^1}$ is a model neighborhood by
{
  \showlabelsinline
  \cref{def:modfunc}
}
 and
\cref{mod:closed} of
{
  \showlabelsinline
  \cref{def:model}
}.

$
  \phi^2 \morphCompose \morphName{f}^1 \morphCompose \phi^{1-1} \maps
    \phi^1[I^{1,2}] \to V^2
$
is a morphism of $\catName{C}^2 \subcat \catName{C}^3$ by hypothesis and
$
  \phi^2 \morphCompose \morphName{f}^1 \morphCompose \phi^{1-1} \maps
    \phi^1[I^{1,3}_{u^1}] \to V^2
$
is a morphism of $\catName{C}^2 \subcat \catName{C}^3$ by
{
  \showlabelsinline
  \cref{mod:restriction}.
}
of \cref{def:model}.
$
  \phi^3 \morphCompose \morphName{f}^2 \morphCompose \phi^{2-1} \maps
    \phi^2[I^{2,3}] \to V^3
$
is a morphism of $\catName{C}^3$ by hypothesis and
and thus
$
  (
    \phi^3 \morphCompose
    \morphName{f}^2 \morphCompose
    \morphName{f}^1 \morphCompose
    \phi^{1-1}
    \maps \phi^1[I^{1,3}_{u^1}] \to V^3
  )
  =
  (
    \phi^3 \morphCompose \morphName{f}^2 \morphCompose \phi^{2-1} \maps
      \phi^2[I^{2,3}] \to V^3
  )
  \morphCompose
  (
    \phi^2 \morphCompose \morphName{f}^1 \morphCompose \phi^{1-1} \maps
      \phi^1[I^{1,3}_{u^1}] \to V^2
  )
$
is a morphism of $\catName{C}^3$.

$
  (
    \phi^3 \morphCompose
    \morphName{f}^2 \morphCompose
    \morphName{f}^1 \morphCompose
    \phi^{1-1}
    \maps \phi^1[I^{1,3}_{u^1}] \to V^3
  )
$
agrees with
$
  (
    \phi^3 \morphCompose
    \morphName{f}^2 \morphCompose
    \morphName{f}^1 \morphCompose
    \phi^{1-1}
    \maps \phi^1[I] \to V^3
  )
$
on $\phi^1[I^{1,3}_{u^1}]$ and
$\phi^1[I] = \union[u^1 \in I]{\phi^1[I^{1,3}_{u^1}]}$,
thus
$
  (
    \phi^3 \morphCompose
    \morphName{f}^2 \morphCompose
    \morphName{f}^1 \morphCompose
    \phi^{1-1}
    \maps \phi^1[I] \to V^3
  )
$
is a morphism of $\catName{C}^3$ by
{
  \showlabelsinline
  \cref{mod:sheaf}
}
of \cref{def:model}.

\item [\ensuremath{\morphName{f}^2 \morphCompose \morphName{f}^1} is (semi-strict, strict).]
If $\seqname{E}^1 \submod \seqname{E}^2$ and
$\seqname{E}^2 \submod \seqname{E}^3$ then
$\seqname{E}^1 \submod \seqname{E}^3$. If
$\seqname{C}^1 \submod \seqname{C}^2$ and
$\seqname{C}^2 \submod \seqname{C}^3$ then
$\seqname{C}^1 \submod \seqname{C}^3$.
If $\catName{C}^1 \subcat \catName{C}^2$ and
$\catName{C}^2 \subcat \catName{C}^3$ then
$\catName{C}^1 \subcat \catName{C}^3$.

If $\morphName{f}^1$ is a local morphism of $\seqname{E}^2$ and
$\morphName{f}^2$ is a local morphism of $\seqname{E}^3$ then
$\morphName{f}^2 \morphCompose \morphName{f}^1$ is a local morphism of
$\seqname{E}^3$ by \pagecref{lem:modlocalMmorphcomp}.

If $\morphName{f}^1$ is a local $\catName{E}^1$-$\catName{E}^2$
morphism of $\seqname{E}^2$ and $\morphName{f}^2$ is a local
$\catName{E}^2$-$\catName{E}^3$ morphism of $\seqname{E}^3$ then
$\morphName{f}^2 \morphCompose \morphName{f}^1$ is a local
$\catName{E}^1$-$\catName{E}^3$ morphism of $\seqname{E}^3$ by
\cref{lem:modlocalMmorphcomp}.

If $\morphName{f}^1$ is a morphism of $\seqname{E}^2 \seqname{E}^3$ and
$\morphName{f}^2$ is a morphism of $\seqname{E}^3$ then
$\morphName{f}^2 \morphCompose \morphName{f}^1$ is a morphism of
$\seqname{E}^3$.
\end{description}
\end{proof}
Similar results apply with restrictions on the admissible atlases.
\end{lemma}

\begin{lemma}[Classic m-atlas inclusion and identity]
\label{lem:M-ATLcid}
Let
\begin{enumerate}
\item
$\catName{E}^i$, $\catName{C}^i$, $i=1,2$, be model categories
\item
$\catName{E}^1 \subcat[full-] \catName{E}^2$
\item
$\seqname{E}^i \objin \catName{E}^i$
\item
$\seqname{E}^1 \submod \seqname{E}^2$
\item
$\seqname{C}^i \objin \catName{C}^i$
\item
$\seqname{C}^1 \submod \seqname{C}^2$
($\catName{C}^1 \subcat[full-] \catName{C}^2$)
\item
$\seqname{A}^i$ be a maximal m-atlas of $\seqname{E}^i$ in the
coordinate space $\seqname{C}^i$
\item
$\seqname{A}^1 \subseteq \seqname{A}^2$
\end{enumerate}

The classic m-atlas inclusion morphism
$
 \Id%
   [
     {(\seqname{E}^1,\seqname{C}^1)},
     {(\seqname{E}^2,\seqname{C}^2)}
   ]
  =
  \ID%
    [
      {(\seqname{E}^1)},
      {(\seqname{E}^2)}
    ]
$
is a semi-strict classic
\nobreaks
  {
    {$\seqname{E}^1$-$\seqname{E}^2$},
    {(-$\catName{C}^1$-$\catName{C}^2$)}
  }
m-atlas morphism of $\seqname{A}^1$ to $\seqname{A}^2$ in the coordinate
spaces $\seqname{C}^1$, $\seqname{C}^2$ and a classic
\nobreaks
  {
    {$\catName{E}^1$-$\catName{E}^2$},
    {(-$\catName{C}^1$-$\catName{C}^2$)}
  }
m-atlas morphism of $\seqname{A}^1$ to $\seqname{A}^2$ in the coordinate
spaces $\seqname{C}^1$, $\seqname{C}^2$.

\begin{proof}
\mbox{}
\begin{description}
\item%
[
  \ensuremath
  {
    \Id%
      [
        {\seqname{E}^1},
        {\seqname{E}^2}
      ]
  }
  is a classic m-atlas morphism:
]
\mbox{}

Let $(U^i, V^i, \phi^i) \in \seqname{A}^i$, $i=1,2$, with
$I \defeq U^1 \cap U^2 \neq \emptyset$.

Since $\seqname{A}^1 \subseteq \seqname{A}^2$, $(U^1, V^1, \phi^1)$ is
compatible with $(U^2, V^2, \phi^2)$, \\*
$
  \phi^2 \morphCompose
  \Id%
    [
      {\seqname{E}^1},
      {\seqname{E}^2}
    ]
  \morphCompose
  \phi^{1-1}
  \maps \phi^1[I] \to V^2
  =
  \phi^2 \morphCompose
  \phi^{1-1}
  \maps \phi^1[I] \to V^2
$
is a morphism of $\seqname{C}^2$ and
a morphism of $\catName{C}^2$.

\item%
[
  \ensuremath
  {
    \Id%
      [
        {\seqname{E}^1},
        {\seqname{E}^2}
      ]
  }
  is semi-striict:
]
\end{description}
\end{proof}

If each $\Top(\seqname{E}^i) \objin \seqname{E}^i$
($\seqname{E}^i \objin \catName{E}^i$)
and
\nobreaks{{$\seqname{C}^1 \submod \seqname{C}^2$}}
\nobreaks{{($\catName{C}^1 \subcat \catName{C}^2$)}}
then the classic m-atlas inclusion morphism
$
  \Id%
    [
     {(\seqname{E}^1,\seqname{C}^1)},
     {(\seqname{E}^2,\seqname{C}^2)}
    ]
$
is a strict classic
\nobreaks
  {
    {($\catName{E}^1$-$\catName{E}^2$-)},
    {$\seqname{E}^1$-$\seqname{E}^2$},
    {(-$\catName{C}^1$-$\catName{C}^2$)}
  }
m-atlas morphism of
$\seqname{A}^1$ to $\seqname{A}^2$ in the coordinate spaces
$\seqname{C}^1$, $\seqname{C}^2$.

\begin{proof}
\mbox{}
\begin{description}
\item%
[
  \ensuremath
  {
    \Id%
      [
        {(\seqname{E}^1,\seqname{C}^1)},
        {(\seqname{E}^2,\seqname{C}^2)}
      ]
  }
  is a classic m-atlas morphism
]

$
 \Id%
   [
     {(\seqname{E}^1,\seqname{C}^1)},
     {(\seqname{E}^2,\seqname{C}^2)}
   ]
$
is a classic
\nobreaks{$\seqname{E}^1$-$\seqname{E}^2$}%
\nobreaks{(-$\catName{C}^1$-$\catName{C}^2$)}
m-atlas morphism of $\seqname{A}^1$ to $\seqname{A}^2$ in the coordinate
spaces $\seqname{C}^1$, $\seqname{C}^2$ and a classic
\nobreaks{$\catName{E}^1$-$\catName{E}^2$}%
\nobreaks{(-$\catName{C}^1$-$\catName{C}^2$)}
m-atlas morphism of $\seqname{A}^1$ to $\seqname{A}^2$ in the coordinate
spaces $\seqname{C}^1$, $\seqname{C}^2$.
by the above proof.

\item%
  [
    \ensuremath
      {
        \Id%
          [
            {(\seqname{E}^1,\seqname{C}^1)},
            {(\seqname{E}^2,\seqname{C}^2)}
          ]
      }
    is strict:
  ]


If each $\Top(\seqname{E}^i) \objin \seqname{E}^i$ then
$\Id[{\seqname{E}^1},{\seqname{E}^2}]$ is a morphism of
$\seqname{E}^2$ by \cref{mod:inclusion} of
{
  \showlabelsinline
  \pagecref{def:model}
},
hence an m-morphism of $\seqname{E}^1$ to $\seqname{E}^2$.

If each
$\seqname{E}^i \objin \catName{E}^i$
$\Id[{\seqname{E}^1},{\seqname{E}^2}]$ is a morphism of
$\catName{E}^2$ by \cref{modcat:morph} of
{
  \showlabelsinline
  \pagecref{def:modcat}
},
hence an $\catName{E}^1$-$\catName{E}^2$ m-morphism of $\seqname{E}^1$
to $\seqname{E}^2$.
\end{description}
\end{proof}
\end{lemma}

\begin{corollary}[Classic m-atlas identity]
\label{cor:M-ATLcid}
Let $\catName{E}$, $\catName{C}$, be model categories,
$\seqname{E} \objin \catName{E}$,
$\seqname{C} \objin \catName{C}$,
$\seqname{A}$ an m-atlas of
$\seqname{E}$ in the coordinate spaces $\seqname{C}$.

The identity morphism
$\seqname{f} \defeq \Id[{{(\seqname{A}, \seqname{E}, \seqname{C})}}]$ is a
strict m-atlas morphism of $\seqname{A}$ to $\seqname{A}$ in the
coordinate spaces $\seqname{C}$, $\seqname{C}$.
\end{corollary}

\subsection{Categories of classic m-atlases and functors}
\label{sub:C-cat}
This subsection defines categories of m-atlases with classic m-atlas
morphisms and functors among them, and proves some basic results.

\subsubsection{Classic m-atlas categories}
\begin{definition}[Sets of classic m-atlas morphisms]
Let $\seqname{E}^i$ and $\seqname{C}^i$. $i=1,2$, be model spaces.
Then
\begin{multline}
  \strict[*]
    {
      \maxfull[*]
        {
          \Atl[classic]_\Ar
            (\seqname{E}^1, \seqname{C}^1, \seqname{E}^2, \seqname{C}^2)
        }
    }
  \defeq \\*
  \set
  {{
    \bigl (
      (\morphName{f}),
      (\seqname{A}^1, \seqname{E}^1, \seqname{C}^1),
      (\seqname{A}^2, \seqname{E}^2, \seqname{C}^2)
    \bigr )
  }}%
  [
    {
      \strict[*]
      {
        \maxfull[*]
          {
            \isAtl[classic]_\Ar
            (
              \seqname{A}^1, \seqname{E}^1, \seqname{C}^1,
              \seqname{A}^2, \seqname{E}^2, \seqname{C}^2,
              \morphName{f}
            )
          }
      }
    }
  ]
\end{multline}
\end{definition}

\begin{definition}[Categories $\Atl^\Classic(\seqname{E}, \seqname{C})$]
\label{def:C-cat}
Let $\seqname{E}$ and $\seqname{C}$ be sets of model spaces.
Let $P \defeq \seqname{E} \times \seqname{C}$.
Then
\begin{equation}
  \strict[*]
    {
      \maxfull[*]
        {
          \Atl[classic]_\Ob(\seqname{E}, \seqname{C})
        }
    }
  \defeq
  \union
    [
      {(\seqname{E}^\mu, \seqname{C}^\mu) \in \seqname{P}},
    ]
    {{
      \strict[*]
        {
          \maxfull[*]
            {
              \Atl[classic]_\Ob
              (
                \seqname{E}^\mu,
                \seqname{C}^\mu
              )
            }
        }
    }}
\end{equation}
\begin{equation}
  \strict[*]
    {
      \maxfull[*]
        {
          \Atl[classic]_\Ar(\seqname{E}, \seqname{C})
        }
    }
  \defeq
  \union
    [
      {(\seqname{E}^\mu, \seqname{C}^\mu) \in \seqname{P}},
      {(\seqname{E}^\nu, \seqname{C}^\nu) \in \seqname{P}}
    ]
    {{
      \strict[*]
        {
          \maxfull[*]
            {
              \Atl[classic]_\Ar
              (
                \seqname{E}^\mu,
                \seqname{C}^\mu,
                \seqname{E}^\nu,
                \seqname{C}^\nu
              )
            }
        }
    }}
\end{equation}
\begin{equation}
  \strict[*]
    {
      \maxfull[*]
        {
          \Atl[classic](\seqname{E}, \seqname{C})
        }
    }
  \defeq
  \bigl (
    \maxfull[*]
      {
        \Atl[classic]_\Ob(\seqname{E}, \seqname{C})
      },
    \maxfull[*]
      {
        \Atl[classic]_\Ar(\seqname{E}, \seqname{C})
      },
    \morphCompose[A]
  \bigr )
\end{equation}

Let $\seqname{E}^i$ and $\seqname{C}^i$, $i=1,2$, be model spaces,
$\seqname{A}^i$ be an atlas of $\seqname{E}^i$ in the coordinate
space $\seqname{C}^i$ and
$\morphName{f} \maps \seqname{E}^1 \to \seqname{E}^2$ be a
(strict, semi-strict) $\seqname{E}^1$-$\seqname{E}^2$ classic M-atlas
 morphism of $\seqname{A}^1$ to $\seqname{A}^2$ in the coordinate
spaces $\seqname{C}^1$, $\seqname{C}^2$.

The identity functor $\Functor_{\Classic,\Id}$ is

\begin{equation}
\Functor_{\Classic,\Id} (\seqname{A}^i, \seqname{E}^i, \seqname{C}^i)
\defeq
(\seqname{A}^i, \seqname{E}^i, \seqname{C}^i)
\end{equation}
\begin{multline}
\Functor_{\Classic,\Id}
  \bigl (
    (\morphName{f} \maps \seqname{E}^1 \to \seqname{E}^),
    (\seqname{A}^1, \seqname{E}^1, \seqname{C}^1),
    (\seqname{A}^2, \seqname{E}^2, \seqname{C}^2)
  \bigr )
\defeq
\\*
  \bigl (
    (\morphName{f} \maps \seqname{E}^1 \to \seqname{E}^2 ),
    (\seqname{A}^1, \seqname{E}^1, \seqname{C}^1),
    (\seqname{A}^2, \seqname{E}^2, \seqname{C}^2)
  \bigr )
\end{multline}

This nomenclature will be justified below.
\end{definition}

\begin{lemma}[$\Atl^\Classic(\seqname{E}, \seqname{C})$ is a category]
\label{lem:M-ATLCiscat}
Let $\seqname{E}$ and $\seqname{C}$ be sets of model spaces. Then
$\Atl[Classic](\seqname{E}, \seqname{C})$ is a category and the identity
morphism for an object $(\seqname{A}^i, \seqname{E}^i, \seqname{C}^i)$
of $\Atl[Classic]_\Ob(\seqname{E}, \seqname{C})$ is
$\Id[{{(\seqname{A}^i, \seqname{E}^i, \seqname{C}^i)}}]$.

\begin{proof}
Let $(\seqname{A}^i,\seqname{E}^i,\seqname{C}^i)$, $i \in [1,3]$,
be an object of $\Atl[Classic](\seqname{E}, \seqname{C})$ and
let \\
$
  \morphName{m}^i \defeq
  \bigl (
    (\morphName{f}^i),
    (\seqname{A}^i,\seqname{E}^i,\seqname{C}^i),
    (\seqname{A}^{i+1},\seqname{E}^{i+1},\seqname{C}^{i+1})
  \bigr )
$
be a morphism of $\Atl[Classic](\seqname{E}, \seqname{C})$. Then
\begin{description}
\item [Composition:]
$\morphName{f}^2 \morphCompose \morphName{f}^1$ is a classic
$\seqname{E}^1$-$\seqname{E}^3$ m-atlas morphism of $\seqname{A}^1$ to
$\seqname{A}^3$ in the coordinate spaces $\seqname{C}^1$,
$\seqname{C}^3$ by \pagecref{lem:M-ATLcmorphcomp} and
$
  (\morphName{f}^2) \morphCompose[()] (\morphName{f}^1) =
  (\morphName{f}^2 \morphCompose \morphName{f}^1)
$.
\item [Associativity:]
Composition is associative by \pagecref{lem:labcomp}.
\item [Identity:]
$\Id[{{(\seqname{A}^i, \seqname{E}^i, \seqname{C}^i)}}]$ is an identity
morphism by \cref{lem:labcomp}.
\end{description}
\end{proof}
\end{lemma}


\begin{corollary}[Subcategories of $\Atl^\Classic(\seqname{E}, \seqname{C})$]
\label{cor:ATLCiscat}
Let $\seqname{E}$ and $\seqname{C}$ be sets of model spaces. Then
all of the following are subcategories of
$\Atl[Classic](\seqname{E}, \seqname{C})$.

\begin{enumerate}
\item $\full{\Atl[Classic]}(\seqname{E}, \seqname{C})$
\item $\maximal[S-]{\Atl[Classic]}(\seqname{E}, \seqname{C})$
\item $\maximal{\Atl[Classic]}(\seqname{E}, \seqname{C})$
\item $\maxfull[S-]{\Atl[Classic]}(\seqname{E}, \seqname{C})$
\item $\maxfull{\Atl[Classic]}(\seqname{E}, \seqname{C})$
\item $\strict[semi-]{\Atl[Classic]}(\seqname{E}, \seqname{C})$
\item $\strict[semi-]{\full{\Atl[Classic]}}(\seqname{E}, \seqname{C})$
\item $\strict[semi-]{\maximal[S-]{\Atl[Classic]}}(\seqname{E}, \seqname{C})$
\item $\strict[semi-]{\maximal{\Atl[Classic]}}(\seqname{E}, \seqname{C})$
\item $\strict[semi-]{\maxfull[S-]{\Atl[Classic]}}](\seqname{E}, \seqname{C})$
\item $\strict[semi-]{\maxfull{\Atl[Classic]}}(\seqname{E}, \seqname{C})$
\item $\strict{\Atl[Classic]}(\seqname{E}, \seqname{C})$
\item $\strict{\full{\Atl[Classic]}}(\seqname{E}, \seqname{C})$
\item $\strict{\maximal[S-]{\Atl[Classic]}}(\seqname{E}, \seqname{C})$
\item $\strict{\maximal{\Atl[Classic]}}(\seqname{E}, \seqname{C})$
\item $\strict{\maxfull[S-]{\Atl[Classic]}}(\seqname{E}, \seqname{C})$
\item $\strict{\maxfull{\Atl[Classic]}}(\seqname{E}, \seqname{C})$
\end{enumerate}

\begin{remark}
They are not, in general, full subcategories.
\end{remark}

The identity functor $\Functor_{\Classic,\Id}$ is a functor from each of the
subcategories to each of the containing categories.
\end{corollary}

\begin{definition}[Categories $\Atl^\Classic(\catName{E}, \catName{C})$]
\label{def:C-cat-cat}
Let $\catName{E}$ and $\catName{C}$ be model categories.
Let $\seqname{P} \defeq Ob(\catName{E}) \times \Ob(\catName{C})$.
Then
\begin{equation}
  \Atl[\Classic]_\Ob(\catName{E}, \catName{C})
  \defeq
  \union
    [
        {\seqname{E} \objin \catName{E}},
        {\seqname{C} \objin \catName{C}}
    ]
    {{
      \Atl[\Classic]_\Ob(\seqname{E}, \seqname{C})
    }}
\end{equation}
\begin{multline}
\Atl[Classic]_\Ar(\catName{E}, \catName{C})
\defeq
  \set
  {{
    \bigl (
      (\morphName{f}),
      (\seqname{A}^1, \seqname{E}^1, \seqname{C}^1),
      (\seqname{A}^2, \seqname{E}^2, \seqname{C}^2)
    \bigr )
  }}%
  [
    {
      (\seqname{E}^i, \seqname{C}^i) \in \seqname{P}
    },
    {
      \strict{\isAtl_\Ar}
      (
        \seqname{A}^1,
        \seqname{E}^1,
        \seqname{C}^1,
        \seqname{A}^2,
        \seqname{E}^2,
        \seqname{C}^2,
        \morphName{f},
        \catName{E},
        \catName{C}
      )
    }
  ]*
\end{multline}
\begin{equation}
  \Atl[\Classic](\catName{E}, \catName{C})
  \defeq
  \bigl (
    \Atl[\Classic]_\Ob(\catName{E}, \catName{C}),
    \Atl[\Classic]_\Ar(\catName{E}, \catName{C}),
    \morphCompose[A]
  \bigr )
\end{equation}

Similar definitions apply with restrictions on the admissible atlases,
the admissible morphisms, or both:
\begin{enumerate}
\item $\full{\Atl[\Classic]}$
\item $\maximal[S-]{\Atl[\Classic]}$
\item $\maximal{\Atl[\Classic]}$
\item $\maxfull[S-]{\Atl[\Classic]}$
\item $\maxfull{\Atl[\Classic]}$
\item $\strict[semi-]{\full{\Atl[\Classic]}}$
\item $\strict[semi-]{\maximal[S-]{\Atl[\Classic]}}$
\item $\strict[semi-]{\maximal{\Atl[\Classic]}}$
\item $\strict[semi-]{\maxfull[S-]{\Atl[\Classic]}}$
\item $\strict[semi-]{\maxfull{\Atl[\Classic]}}$
\item $\strict{\full{\Atl[\Classic]}}$
\item $\strict{\maximal[S-]{\Atl[\Classic]}}$
\item $\strict{\maximal{\Atl[\Classic]}}$
\item $\strict{\maxfull[S-]{\Atl[\Classic]}}$
\item $\strict{\maxfull{\Atl[\Classic]}}$
\end{enumerate}
e.g.,

\begin{equation}
  \maxfull[S-]{\Atl[\Classic]_\Ob}l(\catName{E}, \catName{C})
  \defeq
  \union
    [
      {(\seqname{E}^\mu, \seqname{C}^\mu) \in \seqname{P}},
    ]
    {{
      \maxfull[S-]{\Atl[\Classic]_\Ob}(\seqname{E}^\mu, \seqname{C}^\mu)
    }}
\end{equation}
\begin{equation}
  \strict[semi-]{\maxfull[S-]{\Atl[\Classic]_\Ar}}(\catName{E}, \catName{C})
  \defeq
  \union
    [
      {(\seqname{E}^\mu, \seqname{C}^\mu) \in \seqname{P}},
      {(\seqname{E}^\nu, \seqname{C}^\nu) \in \seqname{P}}
    ]
    {{
      \strict[semi-]{\maxfull[S-]{\Atl[\Classic]_\Ar}}
      (
        \seqname{E}^\mu,
        \seqname{C}^\mu,
        \seqname{E}^\nu,
        \seqname{C}^\nu
      )
    }}
\end{equation}
\begin{equation}
  \strict[semi-]{\maxfull[S-]{\Atl[\Classic]}}(\catName{E}, \catName{C})
  \defeq
  \biggl (
    \maxfull[S-]{\Atl[\Classic]_\Ob}l(\catName{E}, \catName{C}),
    \strict[semi-]{\maxfull[S-]{\Atl[\Classic]_\Ar}}(\seqname{E}, \seqname{C}),
    \morphCompose[A]
  \biggr )
\end{equation}

Let $(\seqname{A}^1, \seqname{E}^1, \seqname{C}^1) \in \Atl[\Classic]_\Ob(\catName{E}, \catName{C})$.
\begin{equation}
\Id[{{(\seqname{A}^1, \seqname{E}^1, \seqname{C}^1)}}] \defeq
\bigl (
  (\Id[{\seqname{E}^1}], \Id[{\seqname{C}^1}]),
  (\seqname{A}^1, \seqname{E}^1, \seqname{C}^1),
  (\seqname{A}^1, \seqname{E}^1, \seqname{C}^1)
\bigr )
\end{equation}
\end{definition}

\begin{lemma}[$\Atl^\Classic(\catName{E}, \catName{C})$ is a category]
\label{lem:C-cat-ATLCiscat}
Let $\catName{E}$ and $\catName{C}$ be model categories. Then
$\Atl[\Classic](\catName{E}, \catName{C})$ is a category and the identity morphism
for an object $(\seqname{A}^i, \seqname{E}^i, \seqname{C}^i)$ of
$\Atl[\Classic]_\Ob(\catName{E}, \catName{C})$ is
$\Id[{{(\seqname{A}^i, \seqname{E}^i, \seqname{C}^i)}}]$.

\begin{proof}
Let $(\seqname{A}^i,\seqname{E}^i,\seqname{C}^i)$, $i \in [1,3]$,
be an object of $\Atl[\Classic](\catName{E}, \catName{C})$ and
let \\
$
  \morphName{m}^i \defeq
  \bigl (
    (\morphName{f}^i),
    (\seqname{A}^i,\seqname{E}^i,\seqname{C}^i),
    (\seqname{A}^{i+1},\seqname{E}^{i+1},\seqname{C}^{i+1})
  \bigr )
$, $i=1,2$,
be a morphism of $\Atl[\Classic](\catName{E}, \catName{C})$. Then
\begin{description}
\item[Composition:]
$\morphName{f}^2 \morphCompose \morphName{f}^1$ is a morphism of $\catName{E}$ and
is an $E^1$-$E^3$ m-atlas morphism of $\seqname{A}^1$ to $\seqname{A}^3$
in the coordinate spaces $\seqname{C}^1$, $\seqname{C}^3$ by
\pagecref{lem:M-ATLcmorphcomp}.

\item[Associativity:]
Composition is associative by \pagecref{lem:labcomp}.

\item[Identity:]
$\Id[{{(\seqname{A}^i, \seqname{E}^i, \seqname{C}^i)}}]$ is an identity
morphism by \cref{lem:labcomp}.
\end{description}
\end{proof}

Similar results follow with restrictions on the admissible atlases,
the admissible morphisms, or both.
\end{lemma}

\subsubsection{Classic m-atlas functors}
\begin{definition}[$\Functor_{\M,\Classic}$]
\label{def:Func-M-C}
Let $\seqname{E}^i$, $i=1,2$, $\seqname{C}^i$ be model spaces and
$\seqname{A}^i$ an m-atlas of $\seqname{E}^i$ in the coordinate model
space $\seqname{C}^i$. Then
\begin{equation}
\Functor_{\M,\Classic} (\seqname{A}^i, \seqname{E}^i, \seqname{C}^i)
\defeq (\seqname{A}^i, \seqname{E}^i, \seqname{C}^i)
\end{equation}

Let
$
  \morphSeqName{f} \defeq
    (
    \morphName{f}_0 \maps \seqname{E}^1 \to \seqname{E}^2,
    \morphName{f}_1 \maps \seqname{C}^1 \to \seqname{C}^2,
    )
$
be an M-atlas morphism of $\seqname{A}^1$ to $\seqname{A}^2$ in the
coordinate spaces $\seqname{C}^1$, $\seqname{C}^2$. Then
\begin{multline}
\Functor_{\M,\Classic}
  \bigl (
    (
      \morphName{f}_0 \maps \seqname{E}^1 \to \seqname{E}^2,
      \morphName{f}_1 \maps \seqname{C}^1 \to \seqname{C}^2,
    ),
    (\seqname{A}^1, \seqname{E}^1, \seqname{C}^1),
    (\seqname{A}^2, \seqname{E}^2, \seqname{C}^2)
  \bigr )
\defeq \\
  \bigl (
    (
      \morphName{f}_0 \maps \seqname{E}^1 \to \seqname{E}^2
    ),
    (\seqname{A}^1, \seqname{E}^1, \seqname{C}^1),
    (\seqname{A}^2, \seqname{E}^2, \seqname{C}^2)
  \bigr )
\end{multline}
\end{definition}

\begin{theorem}[$\Functor_{\M,\Classic}$ is a functor]
\label{the:Func-M-C}
Let $\seqname{E}$ and $\seqname{C}$ be sets of model spaces, Then
$\Functor_{\M,\Classic}$ is a functor from
$\Atl(\seqname{E}, \seqname{C})$ to
$\Atl[classic](\seqname{E}, \seqname{C})$.

\begin{proof}
Let $o^i \defeq  (\seqname{A}^i, \seqname{E}^i, \seqname{C}^i)$,
$i \in [1,3]$, be an object of $\Atl(\seqname{E}, \seqname{C})$ and \\*
$
  m^i \defeq
  \bigl (
    (f_0^i, f_1^i),
    o^i,
    o^{i+1}
  \bigr )
$, $i=1,2$,
be a morphism of $\Atl(\seqname{E}, \seqname{C})$.

$\Functor_{\M,\Classic}$ satisfies these criteria:
\begin{description}
\item[Preservation of endpoints:]
\begin{align*}
  \Functor_{\M,\Classic} o^i
  & =
  o^i \objin \Atl[classic](\seqname{E}, \seqname{C})
\\*
  \Functor_{\Top,\M} m^i
  & =
  \bigl (
    (
      \morphName{f}_0 \maps \seqname{E}^i \to \seqname{E}^{i+1}
    ),
    (\seqname{A}^i, \seqname{E}^i, \seqname{C}^i),
    (\seqname{A}^{i+1}, \seqname{E}^{i+1}, \seqname{C}^{i+1})
  \bigr )
\\*
  & \quad
  \text{is a morphism from $o^i$ to $o^{i+1}$.}
\end{align*}

\item[Identity:]
Let
$
  \bigl (
    (\Id[{\seqname{E}^i}], \Id[{\seqname{C}^i}]),
   o^i,
   o^i
  \bigr )
$
be an identity morphism of
$o^i$ in $\Atl(\seqname{E}, \seqname{C})$.
Then
$
  \bigl (
    (
     \Id[\seqname{E}^i] \maps \seqname{E}^i \to \seqname{E}^i
    ),
    o^i,
    o^i
  \bigr )
$
is an identity morphism of $o^i$ in $\Atl[classic](\seqname{E}, \seqname{C})$.

\item[Composition:]
\begin{equation*}
\begin{array}{lll}
  \Functor_{\M,\Classic} (m^2 \morphCompose[A] m^i)
  & =
  &
  \Functor_{\M,\Classic}
    \bigl (
      (f_0^2 \morphCompose f_0^1, f_1^2 \morphCompose f_1^1),
      o^1,
      o^3
    \bigr )
\\*
  & =
  &
    \bigl (
      (f_0^2 \morphCompose f_0^1),
      o^1,
      o^3
    \bigr )
\\*
  & =
  &
    \bigl (
      (f_0^2),
      o^1,
      o^2
    \bigr )
    \morphCompose[A]
    \bigl (
      (f_0^1),
      o^2,
      o^3
    \bigr )
\\*
  & =
  &
  \Functor_{\M,\Classic}
    \bigl (
      (f_0^2, f_1^2),
      o^2,
      o^3
    \bigr )
  \morphCompose[A]
\\*
  &
  &
    \Functor_{\M,\Classic}
    \bigl (
      (f_0^1, f_1^1),
      o^1,
      o^2
    \bigr )
\\*
  & =
  &
  \Functor_{\M,\Classic} m^2
  \morphCompose[A]
  \Functor_{\M,\Classic} m^1
\end{array}
\end{equation*}
\end{description}
\end{proof}
\end{theorem}

\part{Local Coordinate Spaces}
\label{part:lcs}

This part of the paper defines local coordinate spaces, morphisms among
them and categories of them.  The definitions given here are based on
m-atlas morphisms, but they could as well have been based on classic
m-atlas morphisms.

\section{Local $\seqname{M}$-Sigma and \ensuremath{\catName{M}}-Sigma coordinate spaces}
\label{sec:lcs}
\begin{definition}[Local $\seqname{M}-\Sigma$ and \ensuremath{\catName{M}-\Sigma} coordinate spaces]
{
  \showlabelsinline
  \label{def:LCS}
}
Let $\catSeqName{M} \defeq (\catName{M}_\alpha, \alpha \prec \Alpha)$ be
a sequence of categories,
$\seqname{M} \defeq (M_\alpha, \alpha \prec \Alpha)$ be a
sequence of spaces and
$\morphSeqName{F} \defeq (\morphName{F}_\gamma, \gamma \prec \Gamma)$ a
sequence of functions.
$
  \seqname{L} \defeq
  (\catSeqName{M}, \seqname{M}, \seqname{A}, \morphSeqName{F}, \Sigma)
$
is a (maximal) local $\seqname{M}-\Sigma$ coordinate space, abbreviated
$\max[*]{isLCS_\Ob} \seqname{L}$, a (maximal) local
$\catSeqName{M}-\Sigma$ coordinate space, abbreviated
$\max[*]{isLCS_\Ob} \seqname{L}$, and a (maximal) local $\Alpha-\Sigma$
coordinate space, abbreviated $\max[*]{isLCS_\Ob}(\seqname{L},\Alpha)$,
iff

\begin{enumerate}
\item $\catName{M}_1$ is a model category

\item $(\catSeqName{M}, \seqname{M}, \Sigma, \morphSeqName{F})$ is a
$\catSeqName{M}$-$\Sigma$ prestructure

\item $\seqname{A}$ is a (maximal) m-atlas of $M_0$ in $M_1$.

\item $\forall \gamma \prec \Gamma$, if
$\morphName{F}_\gamma \maps (M_{\sigma_{\gamma,\alpha_\beta}}, \beta
\prec \Beta_\gamma) \to \truthset$ is a constraint function then \\
$\morphName{F}_\gamma
  [
    \bigtimes_{\beta \prec \Beta_\gamma}
    M_{\sigma_{\gamma,\alpha_\beta}}
  ]
= \{ \true \}$, i.e., \\
$\uquant
  {
    {
      (
        s_\beta \in M_{\sigma_{\gamma,\alpha_\beta}},
        \beta \prec \Beta_\gamma
      )
    }
  }
  {{\morphName{F}_\gamma(s_\beta, \beta \prec \Beta_\gamma) = \true}}
  $.
\end{enumerate}

Define the total space $(E, \catName{E})$ to be $M_0$, the Coordinate
space $(C,\catName{C})$ to be $M_1$, the adjunct spaces, if any, to be
$M_\alpha$, $\alpha >1$, and the adjunct functions if any, to be
$F_\alpha$, $\alpha \prec \Alpha$.
\end{definition}

\begin{remark}
There are alternative approaches that are beyond the scope of this
paper, e.g.,
\begin{enumerate}
\item Requiring the atlas to be full would eliminate certain pathologies.
\item A more complex definition would more easily accommodate structures
with more than one atlas, e.g.,
\begin{enumerate}
\item A fiber bundle with a differential structure
\item Multiple fiber bundles on the same base space.
\end{enumerate}
\end{enumerate}
\end{remark}

\begin{lemma}[Local $*-\Sigma$ coordinate spaces]
\label{lem:LCS}
let $\catSeqName{M}^i$, $i=1,2$, be a sequence of categories,
$\catSeqName{M}^2_j$, $j=0,1$, model categories,
$\catSeqName{M}^1 \SUBCAT[full-] \catSeqName{M}^2$ and
\nobreaks
    {{
     $
        \seqname{L}^1 \defeq
        (\catSeqName{M}^1, \seqname{M}, \seqname{A}, \morphSeqName{F}, \Sigma)
    $
  }}
a local $\catSeqName{M}^1-\Sigma$ coordinate space, Then
$
  \seqname{L}^2 \defeq
  (\catSeqName{M}^2, \seqname{M}, \seqname{A}, \morphSeqName{F}, \Sigma)
$
is a local $\catSeqName{M}^2-\Sigma$ coordinate space,

\begin{proof}
$\seqname{L}^2$ satisfies the conditions of \fullcref{def:LCS}:
\begin{enumerate}
\item $\catName{M}^2_0$ and $\catName{M}^2_1$ are model categories by
hypothesis.
\item $(\catSeqName{M}^2, \seqname{M}, \Sigma, \morphSeqName{F})$ is a
$\catSeqName{M}^2$-$\Sigma$ prestructure by
\pagecref{lem:pre}.
\item $\seqname{A}$ is a maximal m-atlas of $M_0$ in $M_1$ by hypothesis.
\item $\forall \gamma \prec \Gamma$, if
$\morphName{F}_\gamma \maps (M_{\sigma_{\gamma,\alpha_\beta}}, \beta
\prec \Beta_\gamma) \to \truthset$ is a constraint function then \\
$
  \morphName{F}_\gamma
    [
      \bigtimes_{\beta \prec \Beta_\gamma}
      M_{\sigma_{\gamma,\alpha_\beta}}
    ]
  = \{ \true \}
$,
i.e., \\
$
  \uquant
  {
    {
      (
        s_\beta \in M_{\sigma_{\gamma,\alpha_\beta}},
        \beta \prec \Beta_\gamma
      )
    }
  }
  {{\morphName{F}_\gamma(s_\beta, \beta \prec \Beta_\gamma) = \true}}
$.
\end{enumerate}
\end{proof}
\end{lemma}

\begin{definition}[$\LCS_\Ob$]
\begin{equation}
  \LCS_\Ob(\catSeqName{M}, \Sigma) \defeq
  \set
    {{
      \seqname{L} \defeq
      (\catSeqName{M}, \seqname{M}, \seqname{A}, \morphSeqName{F}, \Sigma)
    }}%
    [
      {{
        \seqname{M} \seqin \catSeqName{M} \land \isLCS_\Ob \; \seqname{L}
      }}
    ]
\end{equation}
\end{definition}

A variety of conditions can be imposed on the transition functions
$\morphName{t}_\beta^\alpha=\phi_\beta \morphCompose \phi_\alpha^{-1}$ by
appropriate choice of category.

\begin{remark}
While some of cases in \pagecref{Examples} require constraint functions,
the specific examples worked out in \cref{part:man,part:bun} below
do not.
\end{remark}

\section {Morphisms of local coordinate spaces}
\label{sec:lcsmorph}
Informally, a morphism between two local coordinate spaces is a sequence
of functions that is compatible with the functions and m-atlases of the
two local coordinate spaces. The definition implicitly uses commutation
relations, which allows a variety of properties to fall out
automatically, e.g., preservation of fibers, being a homomorphism.

\begin{definition}[Morphisms of local coordinate spaces]
\label{def:LCSmorph}
Let \\*
$\catSeqName{M}^i \defeq (\catName{M}^i_\alpha, \alpha \prec \Alpha)$,
$i=1,2$, be a sequence of categories,
$
  \seqname{M}^i \defeq (\seqname{M}^i_\alpha, \alpha \prec \Alpha)
  \seqin \catSeqName{M}^i
$
a sequence of spaces,
$
  \seqname{L}^i \defeq
  (
    \catSeqName{M}^i,
    \seqname{M}^i,
    \seqname{A}^i,
    \morphSeqName{F}^i,
    \Sigma
  )
$
a local $\seqname{M}^i$-$\Sigma$ coordinate space and
$
  \morphSeqName{f} \defeq
    (
      \morphName{f}_\alpha \maps
        \seqname{M}^1_\alpha \to \seqname{M}^2_\alpha,
      \alpha \prec \Alpha
    )
$.

$\morphSeqName{f}$ is a morphism from $\seqname{L}^1$ to
$\seqname{L}^2$, abbreviated \\*
$
  \isLCS_\Ar
    (
      \catSeqName{M}^1,
      \seqname{M}^1,
      \seqname{A}^1,
      \morphSeqName{F}^1,
      \Sigma,
      \catSeqName{M}^2,
      \seqname{M}^2,
      \seqname{A}^2,
      \morphSeqName{F}^2,
      \morphSeqName{f}
    )
$,
iff

\begin{enumerate}
\item $\morphSeqName{f}$ is a morphism from the prestructure
$
  \seqname{P}^1 \defeq
  (\catSeqName{M}^1, \seqname{M}^1, \Sigma, \morphSeqName{F}^1)
$
to the prestructure
$
  \seqname{P}^2 \defeq
  (\catSeqName{M}^2, \seqname{M}^2, \Sigma, \morphSeqName{F}^2)
$.
\item $(\morphName{f}_0, \morphName{f}_1)$ is a
$\catName{M}^1_0$-$\catName{M}^2_0$ m-atlas morphism from
$\seqname{A}^1$ to $\seqname{A}^2$ in the coordinate model categories
$\catName{M}^1_1$-$\catName{M}^2_1$.
\end{enumerate}

$\morphSeqName{f} \maps \seqname{L}^1 \to \seqname{L}^2$ refers to
$\morphSeqName{f}$ regarded as a morphism from $\seqname{L}^1$ to
$\seqname{L}^2$.

$\morphSeqName{f}$ is a semi-strict morphism from $\seqname{L}^1$
to $\seqname{L}^2$, abbreviated \\*
$
  \strict[semi-]{\isLCS_\Ar}
    (
      \catSeqName{M}^1,
      \seqname{M}^1,
      \seqname{A}^1,
      \morphSeqName{F}^1,
      \Sigma,
      \catSeqName{M}^2,
      \seqname{M}^2,
      \seqname{A}^2,
      \morphSeqName{F}^2,
      \morphSeqName{f}
    )
$,
iff

\begin{enumerate}
\item $\morphSeqName{f}$ is a semi-strict morphism from the prestructure
$
  \seqname{P}^1 \defeq
  (\catSeqName{M}^1, \seqname{M}^1, \Sigma, \morphSeqName{F}^1)
$
to the prestructure
$
  \seqname{P}^2 \defeq
  (\catSeqName{M}^2, \seqname{M}^2, \Sigma, \morphSeqName{F}^2)
$.

\item $(\morphName{f}_0, \morphName{f}_1)$ is a semi-strict
$\catName{M}^1_0$-%
$\catName{M}^2_0$-%
$\seqname{M}^1_0$-%
$\seqname{M}^2_0$-%
$\catName{M}^1_1$-%
$\catName{M}^2_1$
m-atlas morphism from $\seqname{A}^1$ to $\seqname{A}^2$ in the
coordinate model categories $\catName{M}^1_1$-$\catName{M}^2_1$.
\end{enumerate}

$\morphSeqName{f}$ is a strict morphism from
$
  (
    \catSeqName{M}^1,
    \seqname{M}^1,
    \seqname{A}^1,
    \morphSeqName{F}^1,
    \Sigma
  )
$ to
$
  (
    \catSeqName{M}^2,
    \seqname{M}^2,
    \seqname{A}^2,
    \morphSeqName{F}^2,
    \Sigma
  )
$,
abbreviated
$
  \strict{\isLCS_\Ar}
  (
    \catSeqName{M}^1,
    \seqname{M}^1,
    \seqname{A}^1,
    \morphSeqName{F}^1,
    \Sigma,
    \catSeqName{M}^2,
    \seqname{M}^2,
    \seqname{A}^2,
    \morphSeqName{F}^2,
    \morphSeqName{f}
  )
$,
iff

\begin{enumerate}
\item $\morphSeqName{f}$ is a strict morphism from the prestructure
$
  \seqname{P}^1 \defeq
  (\catSeqName{M}^1, \seqname{M}^1, \Sigma, \morphSeqName{F}^1)
$
to the prestructure
$
  \seqname{P}^2 \defeq
  (\catSeqName{M}^2, \seqname{M}^2, \Sigma, \morphSeqName{F}^2)
$.
\item $(\morphName{f}_0, \morphName{f}_1)$ is a strict
$\catName{M}^1_0$-%
$\catName{M}^2_0$-%
$\seqname{M}^1_0$-%
$\seqname{M}^2_0$-%
$\catName{M}^1_1$-%
$\catName{M}^2_1$
m-atlas morphism from $\seqname{A}^1$ to $\seqname{A}^2$ in the
coordinate spaces $\seqname{M}^1_1$, $\seqname{M}^2_1$.
\end{enumerate}

\end{definition}

\begin{definition}[$\seqname{K}$-equivalence of local coordinate space morphisms]
\label{def:LCSmorphKequiv}
Let \\*
$\catSeqName{M}^i \defeq (\catName{M}^i_\alpha, \alpha \prec \Alpha)$,
$i =1,2$, be a sequence of categories,
$
  \seqname{M}^i \defeq (\seqname{M}^i_\alpha, \alpha \prec \Alpha)
  \seqin \catSeqName{M}^i
$
a sequence of spaces,
$
  \seqname{P}^i \defeq
  (
    \catSeqName{M}^i,
    \seqname{M}^i,
    \morphSeqName{F}^i,
    \Sigma
  )
$,
$
  \seqname{L}^i \defeq
  (
    \catSeqName{M}^i,
    \seqname{M}^i,
    \seqname{A}^i,
    \morphSeqName{F}^i,
    \Sigma
  )
$
a local $\seqname{M}^i$-$\Sigma$ coordinate space
$
  \morphSeqName{f} \defeq
    (
      \morphName{f}_\alpha \maps
        \seqname{M}^1_\alpha \to \seqname{M}^2_\alpha,
      \alpha \prec \Alpha
    )
$
and
$
  \morphSeqName{g} \defeq
    (
      \morphName{g}_\alpha \maps
        \seqname{M}^1_\alpha \to \seqname{M}^2_\alpha,
      \alpha \prec \Alpha
    )
$
morphisms from $\seqname{L}^1$ to $\seqname{L}^2$ and
$\seqname{K}$ be a set of ordinals less than $\Alpha$.

$\morphSeqName{f}$ is $\seqname{K}$-equivalent to $\morphSeqName{g}$ iff
$\morphSeqName{f}$ is $\seqname{K}$-equivalent to $\morphSeqName{g}$ as a
prestructure morphism, i.e.,
$
  \uquant
    {\alpha \in \seqname{K}}
    {\morphName{f}_\alpha = \morphName{g}_\alpha}
$.
\end{definition}

\begin{lemma}[morphisms of local coordinate spaces]
\label{lem:LCSmorph}
Let
$\catSeqName{M}^i \defeq (\catName{M}^i_\alpha, \alpha \prec \Alpha)$,
\nobreaks{{$i=1,2$,}} be a sequence of categories,
$
  \seqname{M}^i \defeq (\seqname{M}^i_\alpha, \alpha \prec \Alpha)
  \seqin \catSeqName{M}^i
$
a sequence of spaces, \\
$
  \seqname{P}^i \defeq
  (
    \catSeqName{M}^i,
    \seqname{M}^i,
    \morphSeqName{F}^i,
    \Sigma
  )
$,
$
  \seqname{L}^i \defeq
  (
    \catSeqName{M}^i,
    \seqname{M}^i,
    \seqname{A}^i,
    \morphSeqName{F}^i,
    \Sigma
  )
$
a local $\seqname{M}^i$-$\Sigma$ coordinate space
and
$
  \morphSeqName{f} \defeq
    (
      \morphName{f}_\alpha \maps
        \seqname{M}^1_\alpha \to \seqname{M}^2_\alpha,
      \alpha \prec \Alpha
    )
$,
$i=1,2,3$,
a (semi-strict, strict) morphism from $\seqname{L}^2$ to
$\seqname{L}^2$.

If $\catSeqName{M}^i \SUBCAT[full-] \catSeqName{M}^2$,
$\seqname{S}^i \SUBSETEQ \seqname{S}^2$,
$\seqname{A}^i = \seqname{A}^2$ and
$
  \morphName{F}^i_\gamma = \\
  \morphName{F}^2_\gamma \maps
    \bigtimes \underline{\overset{\sigma_\gamma}{\seqname{S}^i}} \to
    \tail \left ( \overset{\sigma_\gamma}{\seqname{S}^i} \right )
$
then
$\Id[{\seqname{M}^1,\seqname{M}^2}]$ is a strict morphism from
$\seqname{L}^1$ to $\seqname{L}^2$.

\begin{proof}
$\Id[{\seqname{P}^1,\seqname{P}^2}]$ is a strict morphism from
$\seqname{P}^1$ to $\seqname{P}^2$
by \pagecref{lem:premorph}

$
\ID%
  [
    {(\seqname{M}^1_0,\seqname{M}^1_1)},
    {(\seqname{M}^2_0,\seqname{M}^2_1)}
  ]
$
is a strict
$\catName{M}^1_0$-$\catName{M}^2_0$-%
$\seqname{M}^1_0$-$\seqname{M}^2_0$-%
$\catName{M}^1_1$-$\catName{M}^2_1$ m-atlas morphism of
$\seqname{A}^1$ to $\seqname{A}^2$ in the coordinate spaces
$\seqname{M}^1_1$, $\seqname{M}^2_1$.
by
\cref{M-morphcomp:morphcompmorph} of
{
  \showlabelsinline
  \fullcref{lem:M-ATLmorphcomp} on
  \cpageref{M-morphcomp:morphcompmorph}.
}
\end{proof}

Let
$\catSeqName{M}'^i \defeq (\catName{M}'^i_\alpha, \alpha \prec \Alpha)$,
$i=1,2$, be a sequence of categories and \\
\nobreaks
  {{
    $
      \seqname{L}'^i \defeq
      (
        \catSeqName{M}'^i,
        \seqname{M}^i,
        \seqname{A}^i,
        \morphSeqName{F}^i,
        \Sigma
      )
    $
 }}
be a local coordinate space. Then $\morphSeqName{f}$ is a morphism from
$\seqname{L}'^1$ to $\seqname{L}'^2$

\begin{proof}
The commutation relations do not depend on the categories.
\end{proof}

Let
$\catSeqName{M}'^i \defeq (\catName{M}'^i_\alpha, \alpha \prec \Alpha)$,
$i=1,2$, be a sequence of categories,
$\catName{M}'^i_0$ and $\catName{M}'^i_1$ be model categories,
$\catSeqName{M}^i \SUBCAT[full-] \catSeqName{M}'^i$,
$
  \seqname{P}'^i \defeq
  (
    \catSeqName{M}'^i,
    \seqname{M}^i,
    \morphSeqName{F}^i,
    \Sigma
  )
$ and \\*
$
  \seqname{L}'^i \defeq
  (
    \catSeqName{M}'^i,
    \seqname{M}^i,
    \seqname{A}^i,
    \morphSeqName{F}^i,
    \Sigma
  )
$.
Then
\begin{enumerate}
\item $\seqname{L}'^i$ is a local $\seqname{M}^i$-$\Sigma$ coordinate
space and $\morphSeqName{f}$ is a (semi-strict, strict) morphism from
$\seqname{L}'^1$ to $\seqname{L}'^2$.

\item If $\morphSeqName{f}$ is a semi-strict morphism from
$\seqname{L}^1$ to $\seqname{L}^2$ then $\morphSeqName{f}$ is
a semi-strict morphism from $\seqname{L}'^1$ to $\seqname{L}'^2$.

\item If $\morphSeqName{f}$ is a strict morphism from $\seqname{L}^1$
to $\seqname{L}^2$ then $\morphSeqName{f}$ is a strict
morphism from $\seqname{L}'^1$ to $\seqname{L}'^2$.
\end{enumerate}

\begin{proof}
Each $\seqname{L}'^i$ is a local $\seqname{M}^i$-$\Sigma$
coordinate space:

\begin{enumerate}
\item $\catSeqName{M}'^i_0$ and $\catSeqName{M}'^i_1$ are model
categories by hypothesis.

\item $\seqname{P}'^i$ is a $\catSeqName{M}'^i$-$\Sigma$ prestructure
by \pagecref{lem:pre}.

\item $\seqname{A}^i$ is a maximal m-atlas of $\seqname{M}^i_0$ in
$\seqname{M}^i_1$ by hypothesis.

\item All constraint functions evaluate to $\true$ by hypothesis.
\end{enumerate}

Each $\morphSeqName{f}$ is a (semi-strict, strict) morphism from
$\seqname{L}'^1$ to $\seqname{L}'^2$:

\begin{enumerate}
\item $\morphSeqName{f}$ is a morphism from the prestructure
$\seqname{P}'^1$ to the prestructure $\seqname{P}'^2$ by
\pagecref{lem:premorph}.

If $\morphSeqName{f}$ is a semi-strict (strict) morphism from the
prestructure $\seqname{P}^1$ to the prestructure $\seqname{P}^2$
then $\morphSeqName{f}$ is a semi-strict (strict) morphism from the
prestructure $\seqname{P}'^1$ to the prestructure $\seqname{P}'^2$
by \cref{lem:premorph}.

\item $(\morphName{f}_0, \morphName{f}_1)$ is a
$\seqname{M}^1_0$-$\seqname{M}^2_0$ m-atlas morphism from
$\seqname{A}^1$ to $\seqname{A}^2$ in the coordinate model
spaces $\seqname{M}^1_i$, $\seqname{M}^2_i$ by hypothesis.

If $(\morphName{f}_0, \morphName{f}_1)$ is a semi-strict (strict)
$\catName{M}^1_0$-$\catName{M}^2_0$-%
$\seqname{M}^1_0$-$\seqname{M}^2_0$-
$\catName{M}^1_1$-$\catName{M}^2_1$ m-atlas morphism from
$\seqname{A}^1$ to $\seqname{A}^2$ in the coordinate model spaces
$\seqname{M}^1_i$, $\seqname{M}^2_i$ then
$(\morphName{f}_0, \morphName{f}_1)$ is a semi-strict (strict)
$\catName{M}'^1_0$-$\catName{M}'^2_0$-%
$\seqname{M}^1_0$-$\seqname{M}^2_0$-%
$\catName{M}'^1_1$-$\catName{M}'^2_1$ m-atlas morphism from
$\seqname{A}^1$ to $\seqname{A}^2$ in the coordinate model spaces
$\seqname{M}^1_i$, $\seqname{M}^2_i$ by
\pagecref{lem:M-ATLmorph}.
\end{enumerate}
%\end{enumerate}
\end{proof}
\end{lemma}

\begin{lemma}[Composition of local coordinate space morphisms]
\label{lem:LCSmorphcomp}
Let \\*
$\catSeqName{M}^i \defeq (\catName{M}^i_\alpha, \alpha \prec \Alpha)$,
$i \in [1,4]$, be a sequence of categories,
$
  \seqname{M}^i \defeq (\seqname{M}^i_\alpha, \alpha \prec \Alpha)
  \seqin \catSeqName{M}^i
$
a sequence of spaces,
$
  \seqname{P}^i \defeq
  (
    \catSeqName{M}^i,
    \seqname{M}^i,
    \morphSeqName{F}^i,
    \Sigma
  )
$,
$
  \seqname{L}^i \defeq
  (
    \catSeqName{M}^i,
    \seqname{M}^i,
    \seqname{A}^i,
    \morphSeqName{F}^i,
    \Sigma
  )
$
a local $\seqname{M}^i$-$\Sigma$ coordinate space
and
$
  \morphSeqName{f}^i \defeq
    (
      \morphName{f}^i_\alpha \maps
        \seqname{M}^i_\alpha \to \seqname{M}^{i+1}_\alpha,
      \alpha \prec \Alpha
    )
$,
$i=1,2,3$,
a (semi-strict, strict) morphism from $\seqname{L}^i$ to
$\seqname{L}^{i+1}$.

$\morphSeqName{f}^{i+1} \morphCompose \morphSeqName{f}^i$ is
a (semi-strict, strict) morphism from $\seqname{L}^i$ to
$\seqname{L}^{i+2}$:

\begin{proof}
$\morphSeqName{f}^{i+1} \morphCompose \morphSeqName{f}^i$ satisfies the
criteria of \pagecref{def:LCSmorph}.

\begin{enumerate}
\item If each $\morphSeqName{f}^j$ is a (semi-strict, strict) morphism
from the prestructure $\seqname{P}^j$ to the prestructure
$\seqname{P}^{j+1}$ then
$\morphSeqName{f}^{i+1} \morphCompose \morphSeqName{f}^i$ is a
(semi-strict, strict) morphism from the prestructure $\seqname{P}^i$ to
the prestructure $\seqname{P}^{i+2}$ by
\pagecref{lem:premorph}.

If each $\morphSeqName{f}^j$ is a semi-strict (strict) morphism from the
prestructure $\seqname{P}^j$ to the prestructure $\seqname{P}^{j+1}$
then $\morphSeqName{f}^{i+1} \morphCompose \morphSeqName{f}^i$ is a semi-strict
(strict) morphism from the prestructure $\seqname{P}^i$ to the
prestructure $\seqname{P}^{i+2}$ by \cref{lem:premorph}.

\item
$
  (
    \morphName{f}^{i+1}_0 \morphCompose \morphName{f}^i_0,
    \morphName{f}^{i+1}_1 \morphCompose \morphName{f}^i_1
  )
$
is a semi-strict (strict)
$\catName{M}^i_0$-$\catName{M}^{i+1}_0$-%
$\seqname{M}^i_0$-$\seqname{M}^{i+1}_0$-%
$\catName{M}^i_1$-$\catName{M}^{i+1}_1$ m-atlas morphism from
$\seqname{A}^i$ to $\seqname{A}^{i+2}$ in the coordinate model spaces
$\seqname{M}^i_i$, $\seqname{M}^{i+2}_i$ by
\pagecref{lem:M-ATLmorph}.
If
$
  (
    \morphName{f}^{i+1}_0 \morphCompose \morphName{f}^i_0,
    \morphName{f}^{i+1}_1 \morphCompose \morphName{f}^i_1
  )
$
is a semi-strict (strict)
$\catName{M}^i_0$-$\catName{M}^{i+1}_0$-%
$\seqname{M}^i_0$-$\seqname{M}^{i+1}_0$-%
$\catName{M}^i_1$-$\catName{M}^{i+1}_1$
m-atlas morphism from $\seqname{A}^i$ to $\seqname{A}^{i+2}$ in the
coordinate model spaces $\seqname{M}^i_i$, $\seqname{M}^{i+2}_i$ then
$
  (
    \morphName{f}^{i+1}_0 \morphCompose \morphName{f}^i_0,
    \morphName{f}^{i+1}_1 \morphCompose \morphName{f}^i_1
  )
$
is a semi-strict (strict)
$\catName{M}^i_0$-$\catName{M}^{i+1}_0$-%
$\seqname{M}^i_0$-$\seqname{M}^{i+1}_0$-%
$\catName{M}^i_1$-$\catName{M}^{i+1}_1$
m-atlas morphism from $\seqname{A}^i$ to $\seqname{A}^{i+2}$ in the
coordinate model spaces $\seqname{M}^i_i$, $\seqname{M}^{i+2}_i$ by
\cref{M-morphcomp:morphcompmorph} of
{
  \showlabelsinline
  \fullcref{lem:M-ATLmorphcomp} on
  \cpageref{M-morphcomp:morphcompmorph}.
}
\end{enumerate}
\end{proof}

\end{lemma}

\begin{lemma}[Composition of $\seqname{K}$-equivalent local coordinate space morphisms]
\label{lem:LCSmorphcompKequiv}
Let
\begin{enumerate}
\item
$\catSeqName{M}^i \defeq (\catName{M}^i_\alpha, \alpha \prec \Alpha)$,
$i=1,2,3$, be a sequence of categories,
\item
$
  \seqname{M}^i \defeq (\seqname{M}^i_\alpha, \alpha \prec \Alpha)
  \seqin \catSeqName{M}^i
$
be a sequence of spaces,
\item
$
  \seqname{P}^i \defeq
  (
    \catSeqName{M}^i,
    \seqname{M}^i,
    \morphSeqName{F}^i,
    \Sigma
  )
$
be a prestructure
\item
$
  \seqname{L}^i \defeq
  (
    \catSeqName{M}^i,
    \seqname{M}^i,
    \seqname{A}^i,
    \morphSeqName{F}^i,
    \Sigma
  )
$
be a local $\seqname{M}^i$-$\Sigma$ coordinate space
\item
$
  \morphSeqName{f}^1 \defeq
    (
      \morphName{f}^1_\alpha \maps
        \seqname{M}^1_\alpha \to \seqname{M}^{i+1}_\alpha,
      \alpha \prec \Alpha
    )
$
and
$
  \morphSeqName{g}^1 \defeq
    (
      \morphName{g}^1_\alpha \maps
        \seqname{M}^1_\alpha \to \seqname{M}^{i+1}_\alpha,
      \alpha \prec \Alpha
    )
$, \\*
$i=1,2,3$,
be morphisms from $\seqname{L}^i$ to $\seqname{L}^{i+1}$
\item
$\seqname{K}$ be a set of ordinals less than $\Alpha$
\end{enumerate}

If each $\morphSeqName{f}^i$ is $\seqname{K}$-equivalent to
$\morphSeqName{g}^i$ then
$\morphSeqName{f}^2 \morphCompose[()] \morphSeqName{f}^1$ is
$\seqname{K}$-equivalent to
$\morphSeqName{f}^2 \morphCompose[()] \morphSeqName{f}^1$.

\begin{proof}
The result follows from \pagecref{lem:premorphcomp}.
\end{proof}
\end{lemma}

\begin{definition}[$\LCS_\Ar$]
\label{def:LCSmorphAr}
Let
$\catSeqName{M}^i \defeq (\catName{M}^i_\alpha, \alpha \prec \Alpha)$,
$i=1,2$, be sequences of categories and
$\Sigma \defeq (\sigma_\gamma, \gamma \prec \Gamma) \defeq
\bigl ( (\sigma_{\gamma,{\alpha_\beta}}, \alpha_\beta \prec \Alpha,
\beta \preceq \Beta_\gamma), \gamma \prec \Gamma \bigr )$
a signature over $\catSeqName{M}^i$.

\begin{multline}
\LCS_\Ar(\catSeqName{M}^1, \Sigma, \catSeqName{M}^2) \defeq
\set
  {{
    \bigl (
      \morphSeqName{f},
      (
        \catSeqName{M}^1,
        \seqname{M}^1,
        \seqname{A}^1,
        \morphSeqName{F}^1,
        \Sigma
      ),
      (
        \catSeqName{M}^2,
        \seqname{M}^2,
        \seqname{A}^2,
        \morphSeqName{F}^2,
        \Sigma
      )
    \bigr )
  }}%
  [
    {
      \seqname{M}^i \seqin \catSeqName{M}^i
    },
    {
    \isLCS_\Ar
      (
        \catSeqName{M}^1,
        \seqname{M}^1,
        \seqname{A}^1,
        \morphSeqName{F}^1,
        \Sigma,
        \catSeqName{M}^2,
        \seqname{M}^2,
        \seqname{A}^2,
        \morphSeqName{F}^2,
        \morphSeqName{f}
      )
    }
  ]*
\end{multline}
\begin{equation}
\LCS
  (\catSeqName{M}^i, \Sigma)
\defeq
  \bigl (
    \LCS_\Ob
      (
        \catSeqName{M}^i, \Sigma
      ),
    \LCS_\Ar
      (
        \catSeqName{M}^i,
        \Sigma,
        \catSeqName{M}^i
      ),
    \morphCompose[A]
  \bigr )
\end{equation}

If $\catSeqName{M}^1 \SUBCAT[full-] \catSeqName{M}^2$ then

\begin{multline}
\strict[semi-]{\LCS_\Ar}(\catSeqName{M}^1, \Sigma, \catSeqName{M}^2)
\defeq
\set
  {{
    \bigl (
      \morphSeqName{f},
      (
        \catSeqName{M}^1,
        \seqname{M}^1,
        \seqname{A}^1,
        \morphSeqName{F}^1,
        \Sigma
      ),
      (
        \catSeqName{M}^2,
        \seqname{M}^2,
        \seqname{A}^2,
        \morphSeqName{F}^2,
        \Sigma
      )
    \bigr )
  }}%
  [
    {
      \seqname{M}^i \seqin \catSeqName{M}^i
    },
    {
      \strict[semi-]{\isLCS_\Ar}
        (
          \catSeqName{M}^1,
          \seqname{M}^1,
          \seqname{A}^1,
          \morphSeqName{F}^1,
          \Sigma,
          \seqname{M}^2,
          \seqname{A}^2,
          \morphSeqName{F}^2,
          \morphSeqName{f}
        )
    }
  ]*
\end{multline}

\begin{multline}
\strict{\LCS_\Ar}(\catSeqName{M}^1, \Sigma, \catSeqName{M}^2) \defeq
\set
  {{
    \bigl (
      \morphSeqName{f},
      (
        \catSeqName{M}^1,
        \seqname{M}^1,
        \seqname{A}^1,
        \morphSeqName{F}^1,
        \Sigma
      ),
      (
        \catSeqName{M}^2,
        \seqname{M}^2,
        \seqname{A}^2,
        \morphSeqName{F}^2,
        \Sigma
      )
    \bigr )
  }}%
  [
    {
      \seqname{M}^i \seqin \catSeqName{M}^i
    },
    {
      \strict{\isLCS_\Ar}
        (
          \catSeqName{M}^1,
          \seqname{M}^1,
          \seqname{A}^1,
          \morphSeqName{F}^1,
          \Sigma,
          \seqname{M}^2,
          \seqname{A}^2,
          \morphSeqName{F}^2,
          \morphSeqName{f}
        )
    }
  ]*
\end{multline}

\begin{equation}
\strict[semi-]{\LCS}(\catSeqName{M}^1, \Sigma) \defeq
\Bigl (
  \LCS_\Ob(\catSeqName{M}^1, \Sigma),
  \strict[semi-]{\LCS_\Ar}(\catSeqName{M}^1, \Sigma,\catSeqName{M}^1),
  \morphCompose[A]
\Bigr )
\end{equation}
\begin{equation}
\strict{\LCS}(\catSeqName{M}^1, \Sigma) \defeq
\Bigl (
  \LCS_\Ob(\catSeqName{M}^1, \Sigma),
  \strict{\LCS_\Ar}(\catSeqName{M}^1, \Sigma,\catSeqName{M}^1),
  \morphCompose[A]
\Bigr )
\end{equation}

Let
$
  \seqname{M}^i \defeq (\seqname{M}^i_\alpha, \alpha \prec \Alpha)
  \seqin \catSeqName{M}^i
$,
$i=1,2$,
be a sequence of spaces,
$
  \seqname{P}^i \defeq
  (
    \catSeqName{M}^i,
    \seqname{M}^i,
    \morphSeqName{F}^i,
    \Sigma
  )
$
and
$
  \seqname{L}^i \defeq
  (
    \catSeqName{M}^i,
    \seqname{M}^i,
    \seqname{A}^i,
    \morphSeqName{F}^i,
    \Sigma
  )
$
be a local $\seqname{M}^i$-$\Sigma$ coordinate space.

If $\catSeqName{M}^1 \SUBCAT[full-] \catSeqName{M}^2$ then
$\seqname{S}^1 \SUBSETEQ \seqname{S}^2$,
$\seqname{A}^1 = \seqname{A}^2$ and each
$
  \morphName{F}^1_\gamma =
  \morphName{F}^2_\gamma \maps
    \bigtimes \underline{\overset{\sigma_\gamma}{\seqname{S}^1}} \to
    \tail \left ( \overset{\sigma_\gamma}{\seqname{S}^1} \right )
$
then

\begin{multline}
\Id[{\seqname{L}^1,\seqname{L}^2}] \defeq
\bigl (
  \ID[{{\seqname{M}^1},{\seqname{M}^2}}],
  \seqname{L}^1,
  \seqname{L}^2
\bigr )
\end{multline}
\begin{multline}
\Id[{\seqname{L}^i}] \defeq \Id[{\seqname{L}^i,\seqname{L}^i}]
\end{multline}

This nomenclature is justified below.
\end{definition}

\begin{theorem}[$\LCS(\catSeqName{M}, \Sigma)$ is a category]
\label{the:LCSiscat}
Let $\catSeqName{M} \defeq (\catName{M}_\alpha, \alpha \prec \Alpha)$
be a sequence of categories and
$
\Sigma \defeq
(\sigma_\gamma, \gamma \prec \Gamma) \defeq
\bigl (
  (
    \sigma_{\gamma,{\alpha_\beta}}, \beta \preceq \Beta_\gamma
  ),
\gamma \prec \Gamma
\bigr )
$
a signature over $\catSeqName{M}$. Then
$\LCS(\catSeqName{M}, \Sigma)$,
$\strict[semi-]{\LCS}(\catSeqName{M}, \Sigma)$ and
$\strict{\LCS}(\catSeqName{M}, \Sigma)$ are categories.

Let $\seqname{L}^i \objin \LCS(\catSeqName{M}, \Sigma)$.
Then $\Id[{\seqname{L}^i}]$ is the identity morphism of $\seqname{L}^i$:

\begin{proof}
Let
$
  \seqname{L}^i
  \defeq
  (
    \seqname{M}^i,
    \seqname{A}^i,
    \morphSeqName{F}^i,
    \Sigma
  )
$
be an object of $\LCS(\catSeqName{M}, \Sigma)$, let
$
  \seqname{P}^i \defeq \\
  (
    \catSeqName{M},
    \seqname{M}^i,
    \Sigma,
    \morphSeqName{F}^i
  )
$
and let
$
  m^i \defeq
  \bigl (
    \morphSeqName{f}^i \defeq (\morphName{f}^i_\alpha \maps
    \seqname{M}^i_\alpha \to \seqname{M}^{i+1}_\alpha,
    \alpha \prec \Alpha),
    L^i,
    L^{i+1},
  \bigr )
$
be a morphism.
\begin{description}
\item[Composition:]
$\morphSeqName{f}^{i+1} \morphCompose[()] \morphSeqName{f}^i$ is a
prestructure morphism from $\seqname{M}^i$ to $\seqname{M}^{i+2}$
by \pagecref{lem:premorph}
and
$
\bigl (
  \morphSeqName{f}^{i+1}_0 \morphCompose \morphSeqName{f}^i_0,
  \morphSeqName{f}^{i+1}_1 \morphCompose \morphSeqName{f}^i_1
\bigr )
$
is an m-atlas morphism from $\seqname{A}^i$ to $\seqname{A}^{i+1}$ by
\cref{M-morphcomp:morphcompmorph} of
{
  \showlabelsinline
  \fullcref{lem:M-ATLmorphcomp} on
  \cpageref{M-morphcomp:morphcompmorph}.
}

If each $\morphSeqName{f}^i$ is a strict prestructure morphism from
$\seqname{P}^i$ to $\seqname{P}^{i+1}$ then
$\morphSeqName{f}^{i+1} \morphCompose[()] \morphSeqName{f}^i$ is a strict
prestructure morphism from $\seqname{M}^i$ to $\seqname{M}^{i+2}$
by \cref{lem:premorph}.

If each $(\morphSeqName{f}^i_0, \morphSeqName{f}^i_1)$ is a strict (semi-strict)
$\catName{M}_0$-$\catName{M}_0$-%
$\seqname{M}^i_0$-$\seqname{M}^{i+1}_0$-%
$\catName{M}_1$-$\catName{M}_1$
m-atlas morphism from $\seqname{A}^i$ to $\seqname{A}^{i+1}$ in the
coordinate spaces $\seqname{M}^i_1$, $\seqname{M}^{i+1}_1$ then
$
\bigl (
  \morphSeqName{f}^{i+1}_0 \morphCompose \morphSeqName{f}^i_0,
  \morphSeqName{f}^{i+1}_1 \morphCompose \morphSeqName{f}^i_1
\bigr )
$
is a strict (semi-strict)
$\catName{M}_0$-$\catName{M}_0$-%
$\seqname{M}^i_0$-$\seqname{M}^{i+1}_0$-%
$\catName{M}_1$-$\catName{M}_1$
m-atlas morphism from $\seqname{A}^i$ to $\seqname{A}^{i+1}$ in the
coordinate spaces $\seqname{M}^i_1$, $\seqname{M}^{i+1}_1$ by
\cref{M-morphcomp:morphcompmorph} of
{
  \showlabelsinline
  \fullcref{lem:M-ATLmorphcomp} on
  \cpageref{M-morphcomp:morphcompmorph}.
}

\item[Associativity:]
Composition is associative by \pagecref{lem:labcomp}.

\item[Identity:]
$\Id[{\seqname{L}^1}]$ is the identity morphism of $\seqname{L}^1$
by \cref{lem:labcomp}.

\end{description}
\end{proof}
\end{theorem}

\section{Examples}
{
  \showlabelsinline
  \label{sec:ex}
}
\label{Examples}
Functors can be constructed among special cases of local coordinate
spaces and many mathematical structures, although this paper only gives
details for two of them.

\begin{example}[Manifolds]
Choosing the coordinate category as open subsets of a Banach space or
more generally a Fr\'echet space, yields a manifold;
the definitions and results of
\pagecref{part:man}, also cover manifolds with boundaries and
manifolds with tails.
\end{example}

\begin{example}[Fiber bundles]
Choosing the coordinate category as the space of products of open sets
in a topological space (base) with a fixed fiber and the morphisms as
fiber-preserving maps yields a fiber bundle; \pagecref{part:bun} covers
the more general case of a restricted group of transition functions on
fibers.
\end{example}

\begin{example}[Lie groups]
Let $G$ be a topological group with group operation $\star$ and identity
$\mathbf{1}_G$, $C$ a linear space%
\footnote{
  See \pagecref{def:lin} and
  {
    \showlabelsinline
    \pagecref{def:trivck}.
  }
}
and $\seqname{A}$ a maximal $\Ck$-atlas%
\footnote{See \pagecref{def:Ck-ATLmax}.}
of $G$ in the coordinate space $C$. Define
$\hat{\star} \maps G \times G \to G$,
$\isCk_{\star,\seqname{A}} \maps G \times G \to \truthspace$ and
$\seqname{L}$ by:

\begin{equation}
 \hat{\star}(g_1,g_2)  \defeq  g_1 \star g_2^{-1}
\end{equation}
\begin{equation}
  \isCk_{\star,\mathbf{1}_G,\seqname{A}}(g_1,g_2) \defeq
  \begin{cases}
    \true  & \text{$\hat{\star}$ is $\Ck$ at $(g_1,g_2)$} \\*
    \false & \text{$\hat{\star}$ is not $\Ck$ at $(g_1,g_2)$}
  \end{cases}
\end{equation}
\begin{equation}
  \seqname{M} \defeq \Bigl ( \Triv{G}, \Triv[\Ck-]{C}, \seqname{Truthspace} \Bigr )
\end{equation}
\begin{equation}
  \catSeqName{M} \defeq \Singcat{\seqname{M}}
\end{equation}
\begin{equation}
  \morphSeqName{F} \defeq (\star, \isCk_{\star,\seqname{A}})
\end{equation}
\begin{equation}
  \Sigma \defeq
  \bigl (
    (0,0,0),
    (0,0,2)
  \bigr )
\end{equation}
\begin{equation}
  \seqname{L} \defeq
  (
    \catSeqName{M},
    \seqname{M},
    \seqname{A},
    \morphSeqName{F},
    \Sigma
  )
\end{equation}
Then $\seqname{L}$
is a local coordinate space iff $\hat{\star}$ is $\Ck$; the constrain
function $\isCk_{\star,\seqname{A}}$ expresses this condition.
For $k = \infty$, that LCS is equivalent to a Lie group.
\end{example}

\begin{example}[Fiber bundle with global sections]
Let $\seqname{B} \defeq (E, X, Y, \pi, G, \rho, \seqname{A})$ be a fiber
bundle\footnote{See \pagecref{def:Bun}.},
$\morphName{s} \maps X \to E$ and
\begin{equation}
  \mathrm{isSection}_{E,X,\pi,s} (x) \defeq
  \begin{cases}
    \true  & \pi \morphCompose \morphName{s} (x) = x \\*
    \false & \pi \morphCompose \morphName{s} (x) \ne x
  \end{cases}
\end{equation}
a constraint function. Then a simple modification of
\pagecref{def:BunLCS} gives a local coordinate space equivalent to
$\seqname{B}$ with the global section $\morphName{s}$.
\end{example}

\begin{example}[Minkowsky manifolds with foliations]
Similarly to the other examples, a local coordinate space can represent
a manifold, a global section of the tensor bundle, a constraint function
for the signature, a global time coordinate, constraint functions
enforcing diferentiability and a constraint function enforcing a foliation.
\end{example}

\part{Equivalence of manifolds}
\label{part:man}
For a manifold\footnote{
  The literature has several different definitions of a manifold.
  This paper uses one chosen for ease of exposition.
}
the coordinate category is open subsets of a Banach space or more
generally a Fr\'echet space, with an appropriate choice of morphisms.
Choosing a separating hyperplane and half space, with open sets in the
chosen half space, allows manifolds with boundary. Similarly, choosing a
ball with tails allows a manifold with tails.

For differentiable manifolds the coordinate category is similar, but the
morphisms are limited to those sufficiently differentiable, in order to
impose a diferentiability constraint on the transition functions
$\morphName{t}^\alpha_\beta=\phi_\beta \morphCompose \phi_\alpha^{-1}$.

This part of the paper defines $\Ck$-atlases, $\Ck$-manifolds, local
coordinate spaces equivalent to $\Ck$-manifolds, categories of them and
functors, and gives some basic results.

\section {Linear spaces and linear model spaces}
{
  \showlabelsinline
  \label{sec:lin}
}

This subsection defines spaces and related model spaces suitable for use
as the coordinate spaces of generalized $\Ck$ manifolds.

\begin{definition}[Linear spaces]
{
  \showlabelsinline
  \label{def:lin}
}
Let $C$ be a locally arcwise connected topological subspace of a real
(complex) Banach or Fr\'echet space. Then $C$ is a real (complex) linear
space.  If $C$ has a non-void interior, i.e., contains a ball, then it
is non-degenerate.

\begin{remark}
This paper uses the term \textit{ball} to refer to balls of the
underlying space but uses the terms \textit{open set} and
\textit{neighborhoods} to refer to the relative topology.
\end{remark}

Let $\catName{C}$ be a small category whose objects are real (complex)
linear spaces and whose morphisms are $\Ck$ functions. Then
$\catName{C}$ is a $\Ck$ linear category.  If each object of
$\catName{C}$ is non-degenerate then $\catName{C}$ is non-degenerate.

Let $C$ be a real (complex) linear space. Then $\trivcat[\Ck-]{C}$, the
category of all $\Ck$ functions between nonvoid open subspaces of $C$,
is the trivial $\Ck$ linear category of $C$.
\end{definition}

\begin{lemma}[Open subsets of linear spaces are locally arcwise connected]
Let $U$ be an open subset of the real (complex) linear space $S$. Then
$U$ is locally arcwise connected.

\begin{proof}
An open subset of a locally arcwise connected space is locally arcwise
connected.
\end{proof}
\end{lemma}

\begin{definition}[Linear model spaces]
Let $S$ be a real (complex) linear space and
$\seqname{S} \defeq (S, \catName{S})$ a model space for $S$.
Then $\seqname{S}$ is a real (complex) linear model space.

Let $\seqname{S} \defeq (S, \catName{S})$ be a real (complex) linear
model space such that every morphism of $\catName{S}$ is a $\Ck$
function. Then $\seqname{S}$ is a real (complex) $\Ck$ linear model
space.

Let $\catName{S}$ be a small category whose objects are real (complex)
$\Ck$ linear model spaces and whose morphisms are $\Ck$ model functions.
Then $\catName{S}$ is a $\Ck$ linear model category.
\end{definition}

\begin{definition}[Trivial \ensuremath{\Ck} linear model spaces]
\label{def:trivck}
Let $C$ be a real (complex) linear space and $\catName{C}$ the category
of all $\Ck$ functions between open sets of $C$. Then
$\Triv[\Ck-]{C} \defeq (C, \catName{C})$ is the trivial $\Ck$ linear
model space of $C$ and $\Triv[\Ck-]{C}$ is a real (complex) trivial
$\Ck$ linear model space.

Let $\seqname{C}$ be a set of real (complex) linear spaces.

$
  \Triv[\Ck-]{\seqname{C}}
  \defeq
  \set
    {\Triv[\Ck-]{C'}}%
    [C' \in \seqname{C}]
$
is the set of trivial $\Ck$ linear model spaces of $\seqname{C}$.

The category of trivial $\Ck$ linear model spaces of $\seqname{C}$,
abbreviated $\Trivcat[\Ck-]{\seqname{C}}$, is the category whose
objects are $\Triv[\Ck-]{\seqname{C}}$ and whose morphisms are all the
$\Ck$ functions among them.

$\optriv[\Ck-]{\seqname{C}}$, the category of open trivial $\Ck$ model
spaces in $\seqname{C}$, is the category whose objects are
$\set {\Triv[\Ck-]{U}}[U \in \op{\seqname{C}}]$, the trivial $\Ck$
linear model spaces of non-null open sets of spaces in $\seqname{C}$,
and whose morphisms are all the $\Ck$ functions among them.
\end{definition}

\begin{lemma}[Trivial \ensuremath{\Ck} linear model spaces]
\label{lem:trivck}
Let $C^i$, $i=1,2$, be a real (complex) linear space and
$\catName{C}^i \defeq \Cat \biggl ( \Triv[\Ck-]{C^i} \biggr )$.

\begin{enumerate}
\item
$\Triv[\Ck-]{C^i}$ is a $\Ck$ linear model space.

\begin{proof}
\leavevmode
\begin{enumerate}
\item $\Ob(\catName{C}^i)$ is an open cover for $C^i$.

\item $\Ob(\catName{C}^i)$ is closed under finite intersections.

\item The morphisms of $\catName{C}^i$ are $\Ck$, hence continuous.

\item If $\morphName{f} \maps A \to B$ is a morphism,
$A' \objin \catName{C}^i \subseteq A \objin \catName{C}^i$,
$B' \objin \catName{C}^i \subseteq B \objin \catName{C}^i$
and $\morphName{f}[A'] \subseteq B'$ then
$\morphName{f} \restrictto_{A'} \maps A' \to B'$ is $\Ck$ and thus a
morphism.

\item If $A' \objin \catName{C}^i \subseteq A \objin \catName{C}^i$
then the inclusion map $\morphName{i} \maps A' \hookrightarrow A$ is
$\Ck$ and thus a morphism.

\item Restricted sheaf condition: Whenever
\begin{enumerate}
\item $U_\alpha$ and $V_\alpha$, $\alpha \prec \Alpha$,
are objects of $\catName{C}^i$.

\item $\morphName{f}_\alpha \maps U_\alpha \to V_\alpha$ are morphisms of
$\catName{C}^i$.

\item
$U \defeq \union[\alpha \prec \Alpha]{U_\alpha} \objin \catName{C}^i$,

\item
$V \defeq \union[\alpha \prec \Alpha]{V_\alpha} \objin \catName{C}^i$

\item
$\morphName{f} \maps U \to V$ is a continuous function and
for every $\alpha \prec \Alpha$,
$\morphName{f}$ agrees with $\morphName{f}_\alpha$ on $U_\alpha$
\end{enumerate}
then $\morphName{f}$ is $\Ck$ and thus a morphism of $\catName{C}^i$.
\end{enumerate}
\end{proof}

\item
Let $\morphName{f} \maps C^1 \to C^2$ be $\Ck$. Then $\morphName{f}$
is a model function from $\Triv[\Ck-]{C^1}$ to $\Triv[\Ck-]{C^2}$.

\begin{proof}
Let $U^i$ be a model neighborhood of $\Triv[\Ck-]{C^i}$, $i=1,2$.
\begin{enumerate}
\item
Since $U^2$ is a model neighborhood of $\Triv[\Ck-]{C^2}$, $U^2$ is open,
$\morphName{f}^{-1}[U^2]$ is open and thus a model neighborhood of
$\Triv[\Ck-]{C^1}$.
\item
$\morphName{f}[U^1] \subseteq C^2$. Since $C^2$ is open, it is a model
neighborhood of $\Triv[\Ck-]{C^2}$.
\end{enumerate}
\end{proof}

\item
Let $C^1$ be an open subspace of $C^2$. Then
$\Triv[\Ck-]{C^1} \submod \Triv[\Ck-]{C^2}$.

\begin{proof}
By \cref{def:trivck} above, $\Triv[\Ck-]{C^i} = (C^i, \catName{C}^i)$,
where $\catName{C}^i)$ is the category of all $\Ck$ functions between
open sets of $C^i$.

\begin{description}
\item[\ensuremath{\Triv[\Ck-]{C^1} \submod \Triv[\Ck-]{C^2}}:]
\mbox{}
\begin{enumerate}
\item $C^1$ is an open subspace of $C^2$ by hypothesis.

\item Every object of $\catName{C}^1$ is open in $C^1$, thus open in
$C^2$ and thus an object of $\catName{C}^2$. If
$\morphName{f} \maps U \to V \arin \catName{C}^1$ then $U$
and $V$ are open in $C^1$ and $\morphName{f}$ is $\Ck$ in $C^1$,
thus $U$ and $V$ are open in $C^2$ and $\morphName{f}$ is $\Ck$ in
$C^2$ and $\morphName{f} \maps U \to V \arin \catName{C}^2$.
If $\morphName{f} \maps U \to V$ is a morphism of
$\catName{C}^2$, $U \objin \catName{C}^1$ and
$V \objin \catName{C}^1$, then $U$ and $V$ are open in $C^1$ and
$\morphName{f}$ is $\Ck$ in $C^2$, thus $\Ck$ in $C^1$, and
$\morphName{f} \maps U \to V \arin \catName{C}^1$.

\item If $U \objin \catName{C}^2$ then $U$ is open in $C^2$, hence
$U \cap C^1$ is open in $C^2$, hence $U \cap C^1$ is open in $C^1$ and
$U \objin \catName{C}^1$.

\item If $U \objin \catName{C}^2$ and $U \subseteq C^1$ then U is open
is $C^2$, hence open is $C^1$, hence $U \objin \catName{C}^1$.
\end{enumerate}
\end{description}
\end{proof}
\end{enumerate}
\end{lemma}

\begin{definition}[ \ensuremath{\Ck} singleton categories]
\label {def:singCk}
Let $\seqname{C}$ be a $\Ck$ linear model space. Then the $\Ck$
singleton category of $\seqname{C}$, abbreviated
$\singcat[\Ck-]{\seqname{C}}$, is the category whose sole object is
$\seqname{C}$ and whose morphisms are all of the $\Ck$ model functions
from $\seqname{C}$ to itself.
\end{definition}

\section{$\Ck$ charts}
{
  \showlabelsinline
  \label{sec:ckcharts}
}
\begin{definition}[\ensuremath{\Ck} charts]
\label{def:Ck-chart}
Let $C$ be a linear space and $E$ a topological space. A $\Ck$\footnote{
With $k \in \mathbb{N} \cup \{ \infty, \omega \}$.}
chart $(U, V, \phi)$ of $E$ in the coordinate space $C$ consists of

\begin{enumerate}
\item A nonvoid open subset $U \subseteq E$, known as a coordinate patch
\item An open subset $V \subseteq C$
\item A homeomorphism $\phi \maps U \toiso V$, known as a
coordinate function
\end{enumerate}
\begin{remark}
I consider it clearer to explicate the range, rather than the
conventional usage of specifying only the domain and function or the
minimalist usage of specifying only the function.
\end{remark}

A chart $(U, V, \phi)$ is non-degenerate iff $V$ contains a ball of the
underlying Banach or Fr\'echet space.
\end{definition}

\begin{lemma}[\ensuremath{\Ck} charts]
\label{lem:Ck-chart}
Let $C$ be a linear space and $E$ a topological space. The triple
$(U, V, \phi)$ is a $\Ck$ chart of $E$ in the coordinate space $C$ iff
it is an m-chart of $\Triv{E}$ in the coordinate space $\Triv[\Ck-]{C}$.

\begin{proof}
$U$ is a model neighborhood of $\Triv{E}$ iff it is a nonvoid open set
of $E$.  $V$ is a model neighborhood of $\Triv[\Ck-]{C}$ iff it is a
nonvoid open set of $C$.  $\phi$ is required to be a homeomorphis from
$U$ to $V$ in either case.
\end{proof}
\end{lemma}

\begin{definition}[\ensuremath{\Ck} subcharts]
\label{def:Ck-subchart}
Let $(U, V, \phi)$ be a $\Ck$ chart of $E$ in the coordinate space $C$
and $U'$ be a nonvoid open subset of $U$. Then
$(U', V', \phi') \defeq (U', \phi[U'], \phi \restrictto_{U',V'})$
is a subchart of $(U, V, \phi)$.

By abuse of language we will write $(U', V', \phi)$ for
$(U', V', \phi')$.
\end{definition}

\begin{lemma}[\ensuremath{\Ck} subcharts]
\label{lem:Ck-subchart}
Let $(U, V, \phi)$ be a $\Ck$ chart of $E$ in the coordinate space $C$
and $(U', V', \phi')$ be a triple. Then $(U', V', \phi')$ is a subchart
of $(U, V, \phi)$ iff it is a subchart of $(U, V, \phi)$ considered as a
chart of $\Triv{E}$ in the coordinate space $\Triv[\Ck-]{C}$.

\begin{proof}
$(U',V',\phi')$ satisfies the conditions of \cref{def:Ck-chart};
\begin{enumerate}
\item
$U$ is open by \cref{def:Ck-chart}
\item
$U'$ is required to be a nonvoid open subset of $U$ in either case,
and thus a model neighborhood of $\Triv{E}$.
\item
$\phi$ is a homeomorphism, so $V'=\phi[U']$ is open,
and thus a model neighborhood of $\Triv[\Ck-]{C}$.
\item
$U'$ is a model neighborhood of $\Triv{E}$ iff it is an open set of $E$.
\item
$\phi$ is a homeomorphism, so
$\phi \restrictto_{U'} \maps U' \toiso \phi[U']$ is also.
\end{enumerate}
\end{proof}
\end{lemma}

\begin{corollary}[\ensuremath{\Ck} subcharts]
\label{cor:Ck-subchart}
Let $(U, V, \phi)$ be a $\Ck$ chart of $E$ in the coordinate space $C$
and $(U', V', \phi')$ be a subchart of $(U, V, \phi)$. Then
$(U', V', \phi')$ is a $\Ck$ chart of $E$ in the coordinate space $C$.

\begin{proof}
The result follows from \fullcref{lem:Ck-chart} above,
{
  \showlabelsinline
  \pagecref{lem:M-chart}
}
and \pagecref{def:M-subchart}.
\end{proof}
\end{corollary}

\begin{definition}[\ensuremath{\Ck} compatibility]
\label{Ck-compat}
Let $(U, V, \phi)$ and $(U',V', \phi')$ be $\Ck$ charts of $E$ in the
coordinate space $C$. Then $(U, V, \phi)$ is $\Ck$ compatible with
$(U', V', \phi')$ iff either
\begin{enumerate}
\item $U$ and $U'$ are disjoint
\item The transition function $\morphName{t}=\phi' \morphCompose \phi^{-1} \restrictto_{\phi[U \cap U']}$
is a $\Ck$ diffeomorphism.
\end{enumerate}
\end{definition}

\begin{lemma}[Symmetry of \ensuremath{\Ck} compatibility]
\label{lem:Ck-compatsym}
Let $(U, V, \phi)$ and $(U',V', \phi')$ be $\Ck$ charts of $E$ in the
coordinate space $C$.
Then $(U, V, \phi)$ is $\Ck$ compatible with $(U', V', \phi')$ iff
$(U', V', \phi')$ is $\Ck$ compatible with $(U,V,\phi)$.

\begin{proof}
It suffices to prove the implication in only one direction.
\begin{enumerate}
\item $U \cap U' = U' \cap U$.
\item Since the transition function
$\morphName{t}=\phi' \morphCompose \phi^{-1} \restrictto_{\phi[U \cap U']}$
is a  $\Ck$ diffeomorphism of $C$, so is
$\morphName{t}^{-1} = \phi \morphCompose \phi'^{-1} \restrictto_{\phi'[U \cap U']}$.
\end{enumerate}
\end{proof}
\end{lemma}

\begin{lemma}[\ensuremath{\Ck} compatibility of subcharts]
Let $(U_i, V_i, \phi_i)$, $i=1,2$, be $\Ck$ charts of $E$ in the
coordinate space $C$, $(U'_i, V'_i, \phi'_i)$ be subcharts and
$(U_1, V_1, \phi_1)$ be $\Ck$ compatible with $(U_2, V_2, \phi_2)$. Then
$(U'_1, V'_1, \phi'_1)$ is $\Ck$ compatible with
$(U'_2, V'_2, \phi'_2)$.

\begin{proof}
If $U_1 \cap U_2 = \emptyset$ then $U'_1 \cap U'_2 = \emptyset$. If
$U'_1 \cap U'_2 = \emptyset$ then $(U'_1, V'_1, \phi'_1)$ is $\Ck$
compatible with $(U'_2, V'_2, \phi'_2)$. Otherwise, the transition
function
$
  t^1_2 \defeq
  \phi_2 \morphCompose \phi^{-1}_1 \restrictto_{\phi_1[U_1 \cap U_2]}
$
is a $\Ck$ diffeomorphism and hence
$
  t^1_2 \restrictto_{\phi_1[U'_1 \cap U'_2]} \maps
  \phi_1[U'_1 \cap U'_2] \toiso
  \phi_2[U'_1 \cap U'_2]
$
is a $\Ck$ diffeomorphism.
\end{proof}
\end{lemma}

\begin{corollary}[\ensuremath{\Ck} compatibility with subcharts]
Let $(U, V, \phi)$ be a $\Ck$ charts of $E$ in
the coordinate space $C$ and $(U', V', \phi')$ a
subchart. Then $(U', V', \phi')$ is $\Ck$ compatible with
$(U,V,\phi)$.

\begin{proof}
$(U, V, \phi)$ is $\Ck$ compatible with itself and is a subchart of
itself,
\end{proof}
\end{corollary}

\begin{definition}[Covering by \ensuremath{\Ck} charts]
Let $\seqname{A}$ be a set of charts of $E$ in the coordinate space $C$.
$\seqname{A}$ covers $E$ iff $\pi_1[\seqname{A}]$ covers $E$, i.e.,
$E = \union{{\pi_1[\seqname{A}]}}$.
\end{definition}

\section{$\Ck$-atlases}
\label{sec:ckatlases}
A set of charts can be atlases for different coordinate spaces even if
it is for the same total space.  In order to aggregate them into
categories, there must be a way to distinguish them. Including the
two\footnote{The total space is redundant, but convenient.}
spaces in the definitions of the categories serves the purpose.

\begin{definition}[\ensuremath{\Ck}-atlases]
\label{def:Ck-atlas}
Let $\seqname{A}$ be a set of mutually $\Ck$ compatible charts of $E$ in
the coordinate space $C$. $\seqname{A}$ is a $\Ck$-atlas of $E$ in
the coordinate space $C$, abbreviated
$\isAtl[\Ck]_\Ob(\seqname{A}, E, C)$, iff $\seqname{A}$ covers $E$.

$\seqname{A}$ is a full $\Ck$-atlas of $E$ in
the coordinate space $C$, abbreviated
$\full{\isAtl[\Ck]_\Ob}(\seqname{A}, E, C)$, iff
\begin{enumerate}
\item $\pi_1[\seqname{A}]$ covers $E$
\item $\pi_2[\seqname{A}]$ covers $C$.
\end{enumerate}

$\seqname{A}$ is non-degenerate iff it contains a non-degenerate chart.

By abuse of language we write $U \in \seqname{A}$ for
$U \in \pi_1[\seqname{A}]$.

Let $E$ be a topological space and $C$ a linear space. Then
\begin{equation}
\Atl[\Ck]_\Ob(E,C) \defeq
\set
{{
  (\seqname{A}, E, C)
}}%
[
  {{
    \isAtl[\Ck]_\Ob(\seqname{A}, E, C}
  )}
]
\end{equation}
\begin{equation}
\full{\Atl[\Ck]_\Ob}(E,C) \defeq
\set
{{
  (\seqname{A}, E, C)
}}%
[
  {{
    \full{\isAtl[\Ck]_\Ob}(\seqname{A}, E, C}
  )}
]
\end{equation}

Let $\seqname{E}$ be a set of topological spaces and $\seqname{C}$ a
set of linear spaces. Then
\begin{equation}
\Atl[\Ck]_\Ob(\seqname{E}, \seqname{C}) \defeq
  \union
    [
        {E_\mu \in \seqname{E}},
        {C_\mu \in \seqname{C}}
    ]
    {{
      \Atl_\Ob(E_\mu, C_\mu)
    }}
\end{equation}
\begin{equation}
\full{\Atl[\Ck]_\Ob}(\seqname{E}, C) \defeq
\set
{{
  (\seqname{A}, E, C)
}}%
[
  {{
    \full{\isAtl[\Ck]_\Ob}(\seqname{A}, E, C}
  )}
]
\end{equation}
\end{definition}

\begin{lemma}[\ensuremath{\Ck}-atlases]
\label{lem:Ck-atlas}
Let $E$ be a topological space, $C$ be a linear space and $\seqname{A}$
be a $\Ck$-atlas of $E$ in the coordinate space $C$.

If $\seqname{A}$ is non-degenerate then $C$ is non-degenerate.

\begin{proof}
Let $(U,V,\phi)$ be a non-degenerate chart of $\seqname{A}$. $V$
contains a ball and is contained in $C$.
\end{proof}

If $C$ is non-degenerate and $\seqname{A}$ is full then $\seqname{A}$ is
non-degenerate.

\begin{proof}
Let $B$ be a ball in $C$ with center $v$. Since $\seqname{A}$ is full,
it contains a chart $(U,V,\phi)$ with $v \in V$. Then $V$ contains a
ball with center $v$.
\end{proof}
\end{lemma}

\begin{definition}[Compatibility of charts with \ensuremath{\Ck}-atlases]
\label{def:CkCompat}
A chart $(U, V, \phi)$ of $E$ in the coordinate space $C$ is $\Ck$
compatible with a $\Ck$-atlas $\seqname{A}$ iff it is $\Ck$ compatible with
every chart in the atlas.
\end{definition}

\begin{lemma}[Compatibility of subcharts with \ensuremath{\Ck}-atlases]
\label{lem:CkCompat}
Let $\seqname{A}$ be a $\Ck$-atlas of $E$ in the coordinate space $C$
and $\seqname{C}_1 = (U_1, V_1, \phi_1)$ a $\Ck$ chart in
$\seqname{A}$. Then any subchart of $\seqname{C}_1$ is $\Ck$ compatible
with $\seqname{A}$.

\begin{proof}
Let $\seqname{C}' = (U', V', \phi')$ be a subchart of $\seqname{C}_1$ and
$\seqname{C}_2 = (U_2, V_2, \phi_2)$ another chart in $\seqname{A}$.
\begin{enumerate}
\item If $U_1 \cap U_2 = \emptyset$, then $U' \cap U_2 = \emptyset$.
\item If $U' \cap U_2 = \emptyset$ then $\seqname{C}'$ is $\Ck$
compatible with $\seqname{C}_2$.
\item Otherwise the transition function
$t^1_2 \defeq \phi_2 \morphCompose \phi^{-1}_1 \restrictto_{\phi_1[U_1 \cap U_2]}$
is a $\Ck$ diffeomorphism and thus
$t^1_2 \restrictto_{\phi_1[U' \cap U_2]}$ is a $\Ck$ diffeomorphism.
\end{enumerate}
\end{proof}
\end{lemma}

\begin{lemma}[Extensions of \ensuremath{\Ck}-atlases]
\label{Ck-atl:extensions}
Let $\seqname{A}$ be a $\Ck$ atlas of $\seqname{E}$ in the coordinate
space $C$ and $(U_i,V_i,\phi_i)$, $i=1,2$, be a $\Ck$ chart of $E$ in
the coordinate space $C$ $\Ck$ compatible with $\seqname{A}$ in the
coordinate space $C$. Then $(U_1,V_1,\phi_1)$ is $\Ck$ compatible with
$(U_2,V_2,\phi_2)$ in the coordinate space $C$.

\begin{proof}
If $U_1 \cap U_2 = \emptyset$ then $(U_1,V_1,\phi_1)$ is
$\Ck$ compatible with $(U_2,V_2,\phi_2)$. Otherwise,
$\phi_2 \morphCompose \phi^{-1}_1 \restrictto_{\phi_1[U_1 \cap U_2]} \maps
\phi_1[U_1 \cap U_2] \toiso \phi_2[U_1 \cap U_2]$
is a homeomorphism.  It remains to show that
$\phi_2 \morphCompose \phi^{-1}_1 \restrictto_{\phi_1[U_1 \cap U_2]}$
is a $\Ck$ diffeomorphism.
Let $(U'_\alpha,V'_\alpha,\phi'_\alpha)$, $\alpha \prec \Alpha$, be
charts in $\seqname{A}$ such that
$U_1 \cap U_2 \subseteq \union[\alpha \prec \Alpha]{U'_\alpha}$ and
$U_1 \cap U_2 \cap U'_\alpha \neq \emptyset$, $\alpha \prec \Alpha$.
Since the charts are $\Ck$ compatible with $(U'_\alpha,V'_\alpha,\phi'_\alpha)$,
$\phi_2 \morphCompose \phi'^{-1}_\alpha \restrictto_{U_1 \cap U_2 \cap U'_\alpha}$
and
$\phi'_\alpha \morphCompose \phi^{-1}_1 \restrictto_{U_1 \cap U_2 \cap U'_\alpha}$
are $\Ck$ diffeomorphisms and thus
$\phi_2 \morphCompose \phi^{-1}_1 = \phi_2 \morphCompose \phi'^{-1}_\alpha \morphCompose \phi'_\alpha \morphCompose \phi^{-1}_1$
is a $\Ck$ diffeomorphism.
\end{proof}
\end{lemma}

\begin{definition}[Maximal \ensuremath{\Ck}-atlases]
\label{def:Ck-ATLmax}
Let $E$ be a topological space, $C$ be a linear space and $\seqname{A}$
be a non-degenerate $\Ck$-atlas of $E$ in the coordinate space $C$.

$\seqname{A}$ is a maximal $\Ck$-atlas of $E$ in the coordinate space
$C$, abbreviated \\*
$\maximal{\isAtl[\Ck]_\Ob}(\seqname{A}, E,C)$, iff $\seqname{A}$ is a
$\Ck$-atlas that cannot be extended by adding an additional $\Ck$
compatible chart.

\begin{equation}
  \maxfull{\isAtl[\Ck]_\Ob}(\seqname{A}, E, C) \defeq
    \full{\isAtl[\Ck]_\Ob}(\seqname{A}, E, C) \land
    \maximal{\isAtl[\Ck]_\Ob}(\seqname{A}, E, C)
\end{equation}
\end{definition}

\begin{theorem}[Existence and uniqueness of maximal \ensuremath{\Ck}-atlases]
Let $\seqname{A}$ be a $\Ck$-atlas of $E$ in the coordinate space $C$.
Then there exists a unique maximal $\Ck$-atlas
$\maximal{\Atlas^\Ck}(\seqname{A}, E, C)$ of $E$ in the coordinate space $C$
compatible with $\seqname{A}$.

\begin{proof}
Let $\seqname{P}$ be the set of all $\Ck$-atlases of $\seqname{E}$ in
the coordinate space $C$ containing $\seqname{A}$ and $\Ck$ compatible
in the coordinate space $\seqname{C}$ with all of the $\Ck$ charts in
$\seqname{A}$. Let $\maximal{\seqname{P}}$ be a maximal chain of
$\seqname{A}$. Then $A'=\union{\maximal{\seqname{P}}}$ is a maximal
$\Ck$ atlas of $\seqname{E}$ in the coordinate space $\seqname{C}$
$\Ck$ compatible with $\seqname{A}$.
Uniqueness follows from \pagecref{Ck-atl:extensions}.
\end{proof}
\end{theorem}

\begin{definition}[Sets of $\Ck$-atlases]
\label{def:Ck-ATLsets}

Let $E$ be a topological space and $C$ be a linear space.

\begin{multline}
\maxfull[*]{\Atl[\Ck]_\Ob(E, C)} \defeq
\set
{{
   (\seqname{A}, E, C)
}}%
[
  {
    \maxfull[*]
    {
      \isAtl[\Ck]_\Ob
      (\seqname{A}, E, C)
    }
  }
]
\end{multline}

Let $\seqname{E}$ be a set of topological spaces and $\seqname{C}$ a
set of linear spaces. Then

\begin{multline}
\maxfull[*]{\Atl[\Ck]_\Ob (\seqname{E}, \seqname{C})} \defeq
  \union
    [
        {E \in \seqname{E}},
        {C \in \seqname{C}}
    ]
    {{
      \maxfull[*]{\Atl[\Ck]_\Ob(E, C)}
    }}
\end{multline}

\end{definition}

\section{$\Ck$-atlas morphisms and functors}
\label{sec:ckmorph}
This section introduces a taxonomy for morphisms between $\Ck$-atlases,
defines categories of $\Ck$-atlases, defines functors amomg them,
defines functors between them and categories of m-atlases, defines
inverse functors and proves some basic reasults.

\subsection{$\Ck$-atlas morphisms}
This subsection introduces the notions of $\Ck$-atlas morphisms, and
proves some basic results.

\subsubsection{Definitions of \ensuremath{\Ck}-atlas morphisms}
\begin{definition}[\ensuremath{\Ck}-atlas morphisms]
\label{def:Ck-ATLmorph}
Let $E^i$, $i=1,2$, be a topological space, $C^i$ a linear space,
$\seqname{A}^i$ a $\Ck$-atlases of $E^i$ in the coordinate space $C^i$ and
$\morphSeqName{f} \defeq (\morphName{f}_0, \morphName{f}_1)$ a
pair\footnote{
  The conventional definition uses only the first of the two
  functions and a slightly different compatibility condition.
}
of functions.

$\morphSeqName{f}$ is an $E^1$-$E^2$ $\Ck$-morphism of $\seqname{A}^1$ to
$\seqname{A}^2$ in the coordinate spaces $C^1$, $C^2$, abbreviated as
$\isAtl[\Ck]_\Ar(\seqname{A}^1, E^1, C^1$,
                  $\seqname{A}^2$, $E^2$, $C^2$,
                  $\morphName{f}_0$, $\morphName{f}_1)$,
iff
\begin{enumerate}
\item $\morphName{f}_0 \maps E^1 \to E^2$ is a continuous function.
\item $\morphName{f}_1 \maps C^1 \to C^2$ is a $\Ck$ function.
\item for any $u^1 \in E^1$, any chart
$(U^1, V^1, \phi^1 \maps U^1 \toiso V^1) \in \seqname{A}^1$ at $u^1$
and any chart
$(U^2, V^2, \phi^2 \maps U^2 \toiso V^2) \in \seqname{A}^2$ at
$\morphName{f}_0(u^1)$ there exists a subchart \\*
$
  (U'^1, V'^1, \phi^1 \maps U'^1 \toiso V'^1)
  \in \seqname{A}^1
$
at $u^1$, an open set $U'^2 \subseteq U^2$ and a chart
$
  (U'^2 ,\hat{V}'^2, \phi'^2 \maps U'^2 \toiso \hat{V}'^2)
  \in \seqname{A}^2
$
at $\morphName{f}_0(u^1)$ such that
$\morphName{f}_0[U'^1] \subseteq U'^2$ and diagram \eqref{eq:AtlM1} in
\pagecref{def:M-ATLmorph} is commutative, as shown in \pagecref{fig:AtlM1}.
\end{enumerate}

The identity morphism of $(\seqname{A}^i, E^i, C^i)$ is
\begin{equation}
\Id[{{(\seqname{A}^i, E^i, C^i)}}] \defeq
\bigl (
  (\Id[E^i], \Id[C^i]),
  (\seqname{A}^i, E^i, C^i),
  (\seqname{A}^i, E^i, C^i)
\bigr )
\end{equation}
This nomenclature will be justified later.

Let $E^i$, $i=1,2$, be topological spaces and $C^i$ be linear
spaces. Then
\begin{multline}
\Atl[\Ck]_\Ar(E^1, C^1, E^2, C^2) \defeq
  \set
  {{
    \bigl (
      (
        \morphName{f}_0,
        \morphName{f}_1
      ),
      (\seqname{A}^1, E^1, C^1),
      (\seqname{A}^2, E^2, C^2)
    \bigr )
  }}%
  [{{
    \isAtl[\Ck]_\Ar
    (
      \seqname{A}^1,
      E^1,
      C^1,
      \seqname{A}^2,
      E^2,
      C^2,
      \morphName{f}_0,
      \morphName{f}_1
    )
  }}]*
\end{multline}
\begin{multline}
\full{\Atl[\Ck]_\Ar}(E^1, C^1, E^2, C^2) \defeq
  \set
  {{
    \bigl (
      (
        \morphName{f}_0,
        \morphName{f}_1
      ),
      (\seqname{A}^1, E^1, C^1),
      (\seqname{A}^2, E^2, C^2)
    \bigr )
  }}%
  [{{
    \full{\isAtl[\Ck]_\Ar}
    (
      \seqname{A}^1,
      E^1,
      C^1,
      \seqname{A}^2,
      E^2,
      C^2,
      \morphName{f}_0,
      \morphName{f}_1
    )
  }}]*
\end{multline}
\begin{multline}
\maximal{\Atl[\Ck]_\Ar}(E^1, C^1, E^2, C^2) \defeq
  \set
  {{
    \bigl (
      (
        \morphName{f}_0,
        \morphName{f}_1
      ),
      (\seqname{A}^1, E^1, C^1),
      (\seqname{A}^2, E^2, C^2)
    \bigr )
  }}%
  [{{
    \maximal{\isAtl[\Ck]_\Ar}
    (
      \seqname{A}^1,
      E^1,
      C^1,
      \seqname{A}^2,
      E^2,
      C^2,
      \morphName{f}_0,
      \morphName{f}_1
    )
  }}]*
\end{multline}
\begin{multline}
\maximal[S-]{\Atl[\Ck]_\Ar}(E^1, C^1, E^2, C^2) \defeq
  \set
  {{
    \bigl (
      (
        \morphName{f}_0,
        \morphName{f}_1
      ),
      (\seqname{A}^1, E^1, C^1),
      (\seqname{A}^2, E^2, C^2)
    \bigr )
  }}%
  [{{
    \maximal[S-]{\isAtl[\Ck]_\Ar}
    (
      \seqname{A}^1,
      E^1,
      C^1,
      \seqname{A}^2,
      E^2,
      C^2,
      \morphName{f}_0,
      \morphName{f}_1
    )
  }}]*
\end{multline}
\begin{multline}
\maxfull{\Atl[\Ck]_\Ar}(E^1, C^1, E^2, C^2) \defeq
  \set
  {{
    \bigl (
      (
        \morphName{f}_0,
        \morphName{f}_1
      ),
      (\seqname{A}^1, E^1, C^1),
      (\seqname{A}^2, E^2, C^2)
    \bigr )
  }}%
  [{{
    \maxfull{\isAtl[\Ck]_\Ar}
    (
      \seqname{A}^1,
      E^1,
      C^1,
      \seqname{A}^2,
      E^2,
      C^2,
      \morphName{f}_0,
      \morphName{f}_1
    )
  }}]*
\end{multline}
\begin{multline}
\maxfull[S-]{\Atl[\Ck]_\Ar}(E^1, C^1, E^2, C^2) \defeq
  \set
  {{
    \bigl (
      (
        \morphName{f}_0,
        \morphName{f}_1
      ),
      (\seqname{A}^1, E^1, C^1),
      (\seqname{A}^2, E^2, C^2)
    \bigr )
  }}%
  [{{
    \maxfull[S-]{\isAtl[\Ck]_\Ar}
    (
      \seqname{A}^1,
      E^1,
      C^1,
      \seqname{A}^2,
      E^2,
      C^2,
      \morphName{f}_0,
      \morphName{f}_1
    )
  }}]*
\end{multline}
\end{definition}

\begin{definition}[Equivalence of \ensuremath{\Ck}-atlas morphisms]
\label{def:Ck-ATLmorphequiv}
Let $E^i$, $i=1,2$, be a topological space, $C^i$ a linear space,
$\seqname{A}^i$ a $\Ck$-atlases of $E^i$ in the coordinate space $C^i$ and
$\morphSeqName{f} \defeq (\morphName{f}_0, \morphName{f}_1)$ and
$\morphSeqName{g} \defeq (\morphName{g}_0, \morphName{g}_1)$ $E^1$-$E^2$
$\Ck$ morphisms of $\seqname{A}^1$ to $\seqname{A}^2$ in the
coordinate spaces $C^1$, $C^2$.

$\morphSeqName{f}$ is $\Ck$-equivalent to $\morphSeqName{g}$ iff
$\morphName{f}_0 = \morphName{g}_0$.
\end{definition}

\subsubsection{Proclamations on \ensuremath{\Ck}-atlas morphisms}
\begin{lemma}[\ensuremath{\Ck}-atlas morphisms]
\label{lem:Ck-ATLmorph}
Let $\seqname{E}$ be a set of topological spaces, $\seqname{C}$ a set of
$\Ck$ linear spaces, $E^i \in \seqname{E}$, $i=1,2$,
$C^i \in \seqname{C}$, $\seqname{A}^i$ a $\Ck$-atlas of $E^i$ in the
coordinate space $C^i$ and
$
  \morphSeqName{f} \defeq
  (
    \morphName{f}_0 \maps E^1 \to E^2,
    \morphName{f}_1 \maps C^1 \to C^2
  )
$
a pair of functions.

$\morphSeqName{f}$ is an $E^1$-$E^2$ $\Ck$ morphism of
$\seqname{A}^1$ to $\seqname{A}^2$ in the coordinate spaces $C^1$, $C^2$
iff $\morphSeqName{f}$ is a strict
$\Trivcat{\seqname{E}}$-%
$\Trivcat{\seqname{E}}$-%
$\Triv{E^1}$-$\Triv{E^2}$-%
$\optriv[\Ck-]{\seqname{C}}$-%
$\optriv[\Ck-]{\seqname{C}}$
m-atlas morphism of $\seqname{A}^1$ to $\seqname{A}^2$ in the
coordinate spaces $\Triv[\Ck-]{\seqname{C}^1}$,
$\Triv[\Ck-]{\seqname{C}^2}$.

\begin{proof}
The model neighborhoods of $\Triv[\Ck-]{\seqname{C}^i}$ are the open sets
of $\seqname{C}^i$,
$\Trivcat{\seqname{E}} \subcat[full-] \Trivcat{\seqname{E}}$,
the morphisms of $\Trivcat{\seqname{E}}$ are the continuous
functions between spaces in $\seqname{E}$,
$\Trivcat[\Ck-]{\seqname{C}} \subcat[full-] \Trivcat[\Ck-]{\seqname{C}}$,
and the morphisms of $\Trivcat[\Ck-]{\seqname{C}}$ are the $\Ck$
functions between spaces in $\seqname{C}$.
\end{proof}

$\morphSeqName{f}$ is an $E^1$-$E^2$ $\Ck$-morphism of $\seqname{A}^1$ to
$\seqname{A}^2$ in the coordinate spaces $C^1$, $C^2$ iff
$\morphSeqName{f}$ is a strict
$E^1$-$E^2$-%
$\Trivcat[\Ck-]{\seqname{C}}$-%
$\Trivcat[\Ck-]{\seqname{C}}$
m-atlas morphism of $\seqname{A}^1$ to $\seqname{A}^2$ in the coordinate
spaces $\Triv[\Ck-]{\seqname{C}^1}$, $\Triv[\Ck-]{\seqname{C}^2}$.

\begin{proof}
The model neighborhoods of $\Triv[\Ck-]{\seqname{C}^i}$ are the open sets
of $\seqname{C}^i$,
$\Trivcat{\seqname{E}} \subcat[full-] \Trivcat{\seqname{E}}$,
$\Trivcat[\Ck-]{\seqname{C}} \subcat[full-] \Trivcat[\Ck-]{\seqname{C}}$
and the morphisms of $\Trivcat[\Ck-]{\seqname{C}}$ are the $\Ck$
functions between spaces in $\seqname{C}$.
\end{proof}
\end{lemma}

\begin{corollary}[\ensuremath{\Ck}-atlas morphisms]
\label{cor:Ck-ATLmorph}
Let $E^i$, $i=1,2,3$, be a topological space, $C^i$ a linear space,
$\seqname{A}^i$ a $\Ck$-atlas of $E^i$ in the coordinate space $C^i$ and
$(\morphName{f}^i_0, \morphName{f}^i_1)$ an $E^i$-$E^{i+1}$
$\Ck$-morphism of $\seqname{A}^i$ to $\seqname{A}^{i+1}$ in the
coordinate spaces $C^i$, $C^{i+1}$. Then
$
  (
    \morphName{f}^2_0 \morphCompose \morphName{f}^1_0,
    \morphName{f}^2_1 \morphCompose \morphName{f}^1_1
  )
$
is a $E^1$-$E^3$ $\Ck$-morphism of $\seqname{A}^1$ to $\seqname{A}^3$ in
the coordinate spaces $C^1$, $C^3$.

\begin{proof}
The result follows from
\cref{M-morphcomp:morphcompmorph} of
{
  \showlabelsinline
  \fullcref{lem:M-ATLmorphcomp} on
  \cpageref{M-morphcomp:morphcompmorph}.
}
\end{proof}
\end{corollary}

\begin{lemma}[Composition of equivalent \ensuremath{\Ck}-atlas morphisms]
\label{lem:Ck-ATLmorphequiv}
Let $E^i$, $i=1,2,3$, be a topological space,
$C^i$ be a $\Ck$ linear space,
$\seqname{A}^i$ be a $\Ck$-atlas of $E^i$ in the coordinate space $C^i$ and
$
  \morphSeqName{f}^i \defeq
  (
    \morphName{f}^i_0 \maps \seqname{E}^i \to \seqname{E}^{i+1},
    \morphName{f}^i_1 \maps \seqname{C}^i \to \seqname{C}^{i+1}
  )
$
and \\*
$
  \morphSeqName{g}^i \defeq
  (
    \morphName{g}^i_0 \maps \seqname{E}^i \to \seqname{E}^{i+1},
    \morphName{g}^i_1 \maps \seqname{C}^i \to \seqname{C}^{i+1}
  )
$
be $\Ck$-equivalent $\seqname{E}^i$-$\seqname{E}^{i+1}$ m-atlas
morphisms of $\seqname{A}^i$ to $\seqname{A}^{i+1}$ in the coordinate
spaces $\seqname{C}^i$, $\seqname{C}^{i+1}$.

Then $\morphSeqName{f}^2 \morphCompose[()] \morphName{f}^1$ is $\Ck$-equivalent
to $\morphSeqName{g}^2 \morphCompose[()] \morphName{g}^1$.

\begin{proof}
$\morphName{f}^1_0 = \morphName{g}^1_0$ and
$\morphName{f}^2_0 = \morphName{g}^2_0$, hence
$
  \morphName{f}^2_0 \morphCompose \morphName{f}^1_0 =
  \morphName{g}^2_0 \morphCompose \morphName{g}^1_0
$.
\end{proof}
\end{lemma}

\subsection{Categories of $\mathrm{\Ck}$ atlases and functors}

\subsubsection{Categories of $\mathrm{\Ck}$ atlases}
\begin{definition}[Categories of $\mathrm{\Ck}$ atlases]
\label{def:Ck-ATLcat}
Let $\seqname{E}$ be a set of topological spaces and $\seqname{C}$ a
set of linear spaces.
Let $P \defeq \seqname{E} \times \seqname{C}$. Then
\begin{equation}
  \maximal[*]
    {
      \Atl[\Ck]_\Ar(\seqname{E}, \seqname{C})
    }
  \defeq
    \union
      [
          {(E^\mu, C^\mu) \in \seqname{P}},
          {(E^\nu, C^\nu) \in \seqname{P}}
      ]
      {{
        \maximal[*]
        {
          \Atl[\Ck]_\Ar
          (
            E^\mu,
            C^\mu,
            E^\nu,
            C^\nu
          )
        }
      }}
\end{equation}
\begin{equation}
  \maximal[*]
    {
      \Atl[\Ck](\seqname{E}, \seqname{C})
    }
  \defeq
  \Bigl (
    \maximal[*]
      {
        \Atl[\Ck]_\Ob(\seqname{E}, \seqname{C})
      },
    \maximal[*]
      {
        \Atl[\Ck]_\Ar(\seqname{E}, \seqname{C})
      },
    \morphCompose[A]
  \Bigr )
\end{equation}
\end{definition}

\begin{lemma}[\ensuremath{\Atl[\Ck](\seqname{E}, \seqname{C}} is a category]
\label{lem:Ck-ATLiscat}
Let $\seqname{E}$ be a set of topological spaces and $\seqname{C}$ a set
of linear spaces. Then each of
$\maximal[*]{\Atl[\Ck]}(\seqname{E}, \seqname{C})$ is a category.

Let
$
  (\seqname{A}^i, E^i, C^i)
  \in
  \Atl[\Ck]_\Ob(\seqname{E}, \seqname{C})
$.
Then $\Id[{{(\seqname{A}^i, E^i, C^i)}}]$ is the
identity morphism for $(\seqname{A}^i, E^i, C^i)$.

\begin{proof}
Let $(\seqname{A}^i,E^i,C^i)$, $i \in [1,3]$,
be an object of $\Atl[\Ck](\seqname{E}, \seqname{C})$ and
let
$
  \morphName{m}^i \defeq \\
  \bigl (
    (\morphName{f}_0^i,\morphName{f}_1^i),
    (\seqname{A}^i,E^i,C^i),
    (\seqname{A}^{i+1},E^{i+1},C^{i+1})
  \bigr )
$, $i=1,2$,
be a morphism of $\Atl[\Ck](\seqname{E}, \seqname{C})$. Then
\begin{description}
\item[Composition:]
$
  \bigl (
    (
      \morphName{f}_0^2 \morphCompose \morphName{f}_0^1,
      \morphName{f}_1^2 \morphCompose \morphName{f}_1^1
    ),
    (\seqname{A}^1,E^1,C^1),
    (\seqname{A}^3,E^3,C^3)
  \bigr )
$
is a morphism of \\*
$\Atl[\Ck](\seqname{E}, \seqname{C})$ by
\pagecref{cor:Ck-ATLmorph}.

\item[Associativity:]
Composition is associative by \pagecref{lem:labcomp}.

\item[Identity:]
$\Id[{{(\seqname{A}^i, E^i, C^i)}}]$ is an identity
morphism by \cref{lem:labcomp}.
\end{description}
\end{proof}
\end{lemma}

\subsubsection{\ensuremath{\Ck} atlas functors}
\begin{definition}[Functors from \ensuremath{\Ck} atlases to m-atlases]
\label{def:Func-Ck-M}
Let $E^i$, $i=1,2$, be a topological space, $C^i$ a linear space,
$\seqname{A}^i$ $\Ck$-atlases of $E^i$ in the coordinate space $C^i$,
$\morphName{f}_0 \maps E^1 \to E^2$ continuous and
$\morphName{f}_1 \maps C^1 \to C^2$ $\Ck$. Then
\begin{equation}
\Functor^\Ck_{\Ck,\M}  \bigl ( \seqname{A}^i, E^i, C^i \bigl )
\defeq
  \Bigl (
    \seqname{A}^i,
    \Triv{E^i},
    \Triv[\Ck-]{C^i}
  \Bigr )
\end{equation}
\begin{multline}
\Functor^\Ck_{\Ck,\M}
  \bigl (
    (
      \morphName{f}_0,
      \morphName{f}_1
    ),
    (\seqname{A}^1, E^1, C^1),
    (\seqname{A}^2, E^2, C^2)
  \bigr )
\defeq \\
  \biggl (
    \Bigl (
      \morphName{f}_0,
      \morphName{f}_1
     \Bigr ) ,
    \Bigl ( \seqname{A}^1, \Triv{E^1}, \Triv[\Ck-]{C^1} \Bigr ) ,
    \Bigl ( \seqname{A}^2, \Triv{E^2}, \Triv[\Ck-]{C^2} \Bigr )
  \biggr )
\end{multline}
\end{definition}

\begin{theorem}[Functors from \ensuremath{\Ck} atlases to m-atlases]
\label{the:Func-Ck-M}
Let $\seqname{E}$ be a set of topological spaces and $\seqname{C}$ a
set of linear spaces. Then
$\Functor^\Ck_{\Ck,\M}$ is a functor from
$\Atl[\Ck] \bigl ( \seqname{E}, \seqname{C} \bigl )$
to
$\Atl \bigl ( \Triv{\seqname{E}}, \Triv[\Ck-]{\seqname{C}} \bigl )$

\begin{proof}
Let $o^i \defeq (\seqname{A}^i, E^i, C^i)$,
$i \in [1,3]$, be an object of $\Atl[\Ck](\seqname{E}, \seqname{C})$, \\*
$
  o'^i \defeq
  \Functor^\Ck_{\Ck,\M} o^i
  =
  \bigl (
     \seqname{A}^i,
     \Triv{E^i},
     \Triv[\Ck-]{C^i}
   \bigr )
$,
$
  \morphName{m}^i \defeq
    \bigl (
      (\morphSeqName{f}^i_0, \morphSeqName{f}^i_1),
      o^i,
      o^{i+1}
    \bigr )
$, $i=1,2$,
be a morphism from $o^i$ to $o^{i+1}$ and
$
  \morphName{m}'^i \defeq
  \Functor^\Ck_{\Ck,\M} \morphName{m}^i
  =
    \bigl (
      (\morphSeqName{f}^i_0, \morphSeqName{f}^i_1).
      o'^i,
      o'^{i+1}
    \bigr )
$,
be the corresponding morphism in
$\Atl(\Triv{\seqname{E}}, \Triv[\Ck-]{\seqname{C}})$.

\begin{description}
\item[Preservation of endpoints:]
\begin{align*}
  \Functor^\Ck_{\Ck,\M} o^i
  & =
  \bigl (
     \seqname{A}^i,
     \Triv{E^i},
     \Triv[\Ck-]{C^i}
   \bigr )
\\*
  & =
  o'^i
\\*
  \Functor^\Ck_{\Ck,\M} \morphName{m}^i
  & =
  \morphName{m}'^i
\\*
  & =
    \bigl (
      (
        \morphSeqName{f}^i_0,
        \morphSeqName{f}^i_1
      ),
     o'^i,
     o'^{i+1}
    \bigr )
\end{align*}
$\Functor^\Ck_{\Ck,\M} \morphName{m}^i$ is a morphism from
$\Functor^\Ck_{\Ck,\M} o^i$ to $\Functor^\Ck_{\Ck,\M} o^{i+1}$:

\item[Identity:]
\begin{align*}
  \Functor^\Ck_{\Ck,\M} \Id[{o^i}]
  & =
  \Functor^\Ck_{\Ck,\M}
    \bigl (
      (\Id[{E^1}], \Id[{C^1}]),
      o^i,
      o^i
    \bigr )
\\*
  & =
  \bigl (
    (\Id[{\Triv{E^i}}], \Id[{C^i}]),
    o'^i,
    o'^i
  \bigr )
\\*
  & =
  \bigl (
    (\Id[{\Triv{E^i}}], \Id[{C^i}]),
    \Functor^\Ck_{\Ck,\M} o^i,
    \Functor^\Ck_{\Ck,\M} o^i
  \bigr )
\\*
  & =
  \Id[{\Functor^\Ck_{\Ck,\M} o^i}]
\end{align*}

\item[Composition:]
\begin{align*}
  m^2 \morphCompose[A] m^1
  &=
  \bigl (
    (f_0^2 \morphCompose f_0^1, f_1^2 \morphCompose f_1^1),
    o^1,
    o^3
  \bigr )
\\*
  \Functor^\Ck_{\Ck,\M} (m^2 \morphCompose[A] m^1)
  &=
  \bigl (
    (f_0^2 \morphCompose f_0^1, f_1^2 \morphCompose f_1^1),
    o'^1,
    o'^3
  \bigr )
\\*
  &=
  \bigl (
    (f_0^2, f_1^2),
    o'^2,
    o'^3
  \bigr ) \morphCompose[A]
  \bigl (
    (f_0^1, f_1^1),
    o'^1,
    o'^2
  \bigr )
\\*
  &=
    \Functor^\Ck_{\Ck,\M}
      \bigl (
        (f_0^2, f_1^2),
        o^2,
        o^3
      \bigr ) \morphCompose[A]
    \Functor^\Ck_{\Ck,\M}
      \bigl (
        (f_0^1, f_1^1),
        o^1,
        o^2
      \bigr )
\\*
  &=
  \Functor^\Ck_{\Ck,\M} m^2 \morphCompose[A] \Functor^\Ck_{\Ck,\M} m^1
\end{align*}
\end{description}
\end{proof}
\end{theorem}

\begin{definition}[Functors from m-atlases to \ensuremath{\Ck} atlases]
\label{def:Func-M-Ck}
Let $\seqname{E}^i$, $i=1,2$, be model spaces, $\seqname{C}^i$ be linear
model spaces, $\seqname{A}^i$ a maximal m-atlas of $\seqname{E}^i$ in
the coordinate space $\seqname{C}^i$ and $(\morphName{f}_0,
\seqname{f}_1)$ an $\seqname{E}^1$-$\seqname{E}^2$ m-atlas morphism of
$\seqname{A}^1$ to $\seqname{A}^2$ in the coordinate spaces
$\seqname{C}^1$, $\seqname{C}^2$, Then
\begin{equation}
\Functor^\Ck_{\M,\Ck} (\seqname{A}^i, \seqname{E}^i, \seqname{C}^i)
\defeq
  \bigl (
     \seqname{A}^i,
     \pi_1(\seqname{E}^i),
     \pi_1(\seqname{C}^i)
   \bigr )
\end{equation}
\begin{multline}
\Functor^\Ck_{\M,\Ck}
  \bigl (
    (
      \morphName{f}_0,
      \morphName{f}_1
    ),
    (\seqname{A}^1, \seqname{E}^1, \seqname{C}^1),
    (\seqname{A}^2, \seqname{E}^2, \seqname{C}^2)
  \bigr )
\defeq \\
  \Bigl (
    \bigl (
      \morphName{f}_0,
      \morphName{f}_1
    \bigr ),
    \bigl ( \seqname{A}^1, \pi_1(\seqname{E}^1), \pi_1(\seqname{C}^1) \bigr ),
    \bigl ( \seqname{A}^2, \pi_1(\seqname{E}^2), \pi_1(\seqname{C}^2) \bigr )
  \Bigr )
\end{multline}
\end{definition}

\begin{theorem}[Functors from m-atlases to \ensuremath{\Ck} atlases]
\label{the:Func-M-Ck}
Let $\seqname{E}$ be a set of model spaces and $\seqname{C}$ a
set of $\Ck$ linear model spaces. Then
$\Functor^\Ck_{\M,\Ck}$ is a functor from
$\Atl(\seqname{E}, \seqname{C})$
to
$\Atl[\Ck] \bigl (\pi_1[\seqname{E}], \pi_1[\seqname{C}] \bigr )$.

\begin{proof}
Let
$
o^i \defeq
\bigl (
  \seqname{A}^i, (E^i,\catName{E}^i), (C^i,\catName{C}^i)
\bigr )
$,
$i \in [1,3]$, be an object of $\Atl[\Ck](\seqname{E}, \seqname{C})$ and \\*
$
\morphName{m}^i \defeq
  \bigl (
    (\morphSeqName{f}^i_0, \morphSeqName{f}^i_1).
    o^i,
    o^{i+1}
  \bigr )
$, $i=1,2$,
be a morphism from $o^i$ to $o^{i+1}$.

$\Functor^\Ck_{\M,\Ck} (\morphName{m}^1)$ is a morphism from
$\Functor^\Ck_{\M,\Ck} o^1$ to $\Functor^\Ck_{\M,\Ck} o^2$:

\begin{equation*}
\begin{split}
  &
\Functor^\Ck_{\M,\Ck} (\morphName{m}^1) =
\\*
  &
\Functor^\Ck_{\M,\Ck}
  \bigl (
    (
      \morphSeqName{f}^1_0,
      \morphSeqName{f}^1_1
    ),
    (\seqname{A}^1, \seqname{E}^1, \seqname{C}^1),
    (\seqname{A}^2, \seqname{E}^2, \seqname{C}^2)
  \bigr )
=
\\*
  &
  \Bigl (
    (
      \morphSeqName{f}^1_0,
      \morphSeqName{f}^1_1
    ),
    \bigl (
      \seqname{A}^1, \pi_1(\seqname{E}^1), \pi_1(\seqname{C}^1)
    \bigr ),
    \bigl (
      \seqname{A}^2, \pi_1(\seqname{E}^2), \pi_1(\seqname{C}^2)
    \bigr )
  \Bigr )
\end{split}
\end{equation*}

$\Functor^\Ck_{\M,\Ck}$ maps identity functions to identity functions:

\begin{equation*}
\begin{split}
  &
\Functor^\Ck_{\M,\Ck} \Id[{{(\seqname{A}^i, \seqname{E}^i, \seqname{C}^i)}}]
=
\\*
  &
\Functor^\Ck_{\M,\Ck}
  \bigl (
    (\Id[{\seqname{E}^i}], \Id[{\seqname{C}^i}]),
    (\seqname{A}^i, \seqname{E}^i, \seqname{C}^i),
    (\seqname{A}^i, \seqname{E}^i, \seqname{C}^i)
  \bigr )
=
\\*
  &
  \bigl (
    (\Id[{\pi_1(\seqname{E}^i)}], \Id[{\pi_1(\seqname{C}^i)}]),
    (\seqname{A}^i, \pi_1(\seqname{E}^i), \pi_1(\seqname{C}^i)
    ),
    (\seqname{A}^i, \pi_1(\seqname{E}^i), \pi_1(\seqname{C}^i)
    )
  \bigr )
=
\\*
  &
  \bigl (
    (\Id[{E^i}], \Id[{C^i}]),
    \Functor^\Ck_{\M,\Ck} (\seqname{A}^i, \seqname{E}^i, \seqname{C}^i),
    \Functor^\Ck_{\M,\Ck} (\seqname{A}^i, \seqname{E}^i, \seqname{C}^i)
  \bigr )
=
\\*
  &
\Id[{{\Functor^\Ck_{\M,\Ck} (\seqname{A}^i, \seqname{E}^i, \seqname{C}^i)}}]
\end{split}
\end{equation*}

$
  \Functor^\Ck_{\M,\Ck} m^2 \morphCompose[A] \Functor^\Ck_{\M,\Ck} m^1
  =
  \Functor^\Ck_{\M,\Ck} (m^2 \morphCompose[A] m^1)
$:

\begin{align*}
  m^2 \morphCompose[A] m^1
  & =
  \bigl (
    (f_0^2 \morphCompose f_0^1, f_1^2 \morphCompose f_1^1),
    (\seqname{A}^1, \seqname{E}^1, \seqname{C}^1),
    (\seqname{A}^3, \seqname{E}^3, \seqname{C}^3)
  \bigr )
\\*
  \Functor^\Ck_{\M,\Ck}
    \bigl (
       (\seqname{A}^i, \seqname{E}^i, \seqname{C}^i)
    \bigr )
  & =
  \bigl (
    \seqname{A}^i,
    \pi_1(\seqname{E}^i),
    \pi_1(\seqname{C}^i)
  \bigr )
\\*
  \Functor^\Ck_{\M,\Ck} (m^i)
  & =
    \Bigl (
      (f_0^i, f_1^i),
      \bigl (
        \seqname{A}^i,
        \pi_1(\seqname{E}^i),
        \pi_1(\seqname{C}^i)
      \bigr ),
      \bigl (
        \seqname{A}^{i+1},
        \pi_1(\seqname{E}^{i+1}),
        \pi_1(\seqname{C}^{i+1})
      \bigr )
    \Bigr )
\\*
  \Functor^\Ck_{\M,\Ck} m^2 \morphCompose[A] \Functor^\Ck_{\M,\Ck} m^1
  & =
    \Bigl (
      (f_0^2 \morphCompose f_0^1, f_1^2 \morphCompose f_1^2),
      \bigl (
        \seqname{A}^1,
        \pi_1(\seqname{E}^1),
        \pi_1(\seqname{C}^1)
      \bigr ),
      \bigl (
        \seqname{A}^3,
        \pi_1(\seqname{E}^3),
        \pi_1(\seqname{C}^3)
      \bigr )
    \Bigr )
\\*
  \Functor^\Ck_{\M,\Ck} (m^2 \morphCompose m^1)
  & =
    \Bigl (
      (f_0^2 \morphCompose f_0^1, f_1^2 \morphCompose f_1^2),
      \bigl (
        \seqname{A}^1,
        \pi_1(\seqname{E}^1),
        \pi_1(\seqname{C}^1)
      \bigr ),
      \bigl (
        \seqname{A}^3,
        \pi_1(\seqname{E}^3),
        \pi_1(\seqname{C}^3)
      \bigr )
    \Bigr )
\end{align*}

\end{proof}
\end{theorem}

\section{Classic \ensuremath{\Ck}-atlas morphisms and functors}
\label{sec:Ck-ATLcmorph}
This section introduces an alternate definition of and taxonomy for
morphisms between $\Ck$-atlases, defines categories of $\Ck$-atlases,
defines functors amomg them and proves some basic reasults.  It
introduces the notion of classic $\Ck$-atlas morphisms, which is
equivalent to the conventional definitions for a manifold.  The reader
can safely skip it, but it does help to establish a context for other
notions that are used in the exposition of local coordinate spaces.

\subsection{Classic \ensuremath{\Ck}-atlas morphisms}
\label{sub:Ck-ATLcmorph}
This subsection introduces the notion of classic $\Ck$-atlas morphisms,
and proves some basic results.

\subsubsection{Definitions of classic \ensuremath{\Ck}-atlas morphisms}
\begin{definition}[Classic \ensuremath{\Ck}-atlas morphisms]
\label{def:Ck-ATLcmorph}
Let $E^i$, $i=1,2$, be a topological space, $C^i$ be a linear space,
$\seqname{A}^i$ be a $\Ck$-atlases of $E^i$ in the coordinate space
$C^i$ and $\morphName{f}$ be a continuous function.

$\morphName{f}$ is a $E^1$-$E^2$ classic $\Ck$ morphism of
$\seqname{A}^1$ to $\seqname{A}^2$ in the coordinate spaces $C^1$,
$C^2$, abbreviated as
$
\isAtl[classic,\Ck]_\Ar
  (
    \seqname{A}^1,
    E ^1,
    C^1,
    \seqname{A}^2,
    E^2,
    C^2,
    \morphName{f}
  )
$,
iff for any
$(U^i, V^i, \phi^i \maps U^i \toiso V^i) \in \seqname{A}^i$,
$i=1,2$, with $I \defeq U^1 \cap \morphName{f}^{-1}[U^2] \neq \emptyset$,
$
  \phi^2 \morphCompose \morphName{f} \morphCompose \phi^{1-1} \maps
    \phi^1[I] \to V^2
$
is a $\Ck$ function.
\end{definition}

\begin{remark}
Definitions of constrained, semistrict and strict classic $\Ck$-atlas
morphisms would be pointless, as any classic $\Ck$-atlas
morphism would be constrained and strict.
\end{remark}

\subsubsection{Proclamations on classic \ensuremath{\Ck}-atlas morphisms}
\begin{lemma}[Classic \ensuremath{\Ck}-atlas morphisms]
\label{lem:Ck-ATLcmorph}
Let $\seqname{E}$ be a set of topological spaces, $\seqname{C}$ a set of
$\Ck$ linear spaces, $E^i \in \seqname{E}$, $i=1,2$,
$C^i \in \seqname{C}$, $\seqname{A}^i$ a $\Ck$-atlas of $E^i$ in the
coordinate space $C^i$ and $\morphName{f} \maps E^1 \to E^2$ a function.

Let $E^1$ be an open subspace of $E^2$. Then $\morphName{f}$ is an
$E^1$-$E^2$ classic $\Ck$-atlas morphism of $\seqname{A}^1$ to
$\seqname{A}^2$ in the coordinate spaces $C^1$, $C^2$ iff $\morphName{f}$
is a strict $E^1$-$E^2$ classic m-atlas morphism of $\seqname{A}^1$ to
$\seqname{A}^2$ in the coordinate spaces $\Triv[\Ck-]{\seqname{C}^1}$,
$\Triv[\Ck-]{\seqname{C}^2}$.

$\morphName{f}$ is an $E^1$-$E^2$ classic $\Ck$-atlas morphism of
$\seqname{A}^1$ to $\seqname{A}^2$ in the coordinate spaces $C^1$, $C^2$
iff $\morphName{f}$ is a strict
$E^1$-$E^2$-%
$\optriv[\Ck-]{\seqname{C}}$-%
$\optriv[\Ck-]{\seqname{C}}$
classic m-atlas morphism of $\seqname{A}^1$ to $\seqname{A}^2$ in the
coordinate spaces $\Triv[\Ck-]{\seqname{C}^1}$,
$\Triv[\Ck-]{\seqname{C}^2}$.


\end{lemma}

\section {$\Ck$ manifolds}
\label{sec:ckman}
Conventionally a manifold is different from its atlases, but
$\maximal{\Atl[\Ck]}(\seqname{E}, \seqname{C})$ in
\pagecref{def:Ck-ATLcat} encourages treating them on an equal footing.
All of the results for maximal atlases carry directly over to results
for manifolds.

\begin{remark}
The variations of the $\LCS^\Ck$ constructors in
\pagecref{def:CktoLCS} are very similar, as are the variations of
the $\Functor^\Ck_{\Man,\LCS}$ functors; likewise the variations of the
$\Functor^\Ck_{\LCS,\Man}$ functors in  \pagecref{def:LCStoCk} are very
similar.

The corresponding results and proofs in \pagecref{the:CktoLCS}
are likewise very similar, as are the corresponding results and proofs
in \pagecref{the:LCStoCk}.
\end{remark}

\begin{definition}[\ensuremath{\Ck} manifolds]
\label{def:man}
Let $E$ be a topological space, $C$ a linear space and $\seqname{A}$ a
maximal\footnote{
Requiring that the atlas be full would eliminate some pathologies.}
$\Ck$-atlas of $E$ in the coordinate space $C$. Then $(E, C,
\seqname{A})$ is a $\Ck$ manifold.

Let $\seqname{E}$ be a set of topological spaces and $\seqname{C}$ be a
set of linear spaces. Then
\begin{equation}
\Man^\Ck_\Ob(\seqname{E}, \seqname{C}) \defeq
\maximal{\Atl[\Ck]}(\seqname{E}, \seqname{C})
\end{equation}
\begin{remark}
The manifold $(E, C, \seqname{A})$ corresponds to the object
$(\seqname{A}, E, C)$.
\end{remark}
\end{definition}

\begin{definition}[\ensuremath{\Ck} manifold morphisms]
\label{def:manmorph}
Let $S^i \defeq (E^i,C^i,\seqname{A}^i)$, $i=1,2$, be $\Ck$ manifolds
and $(\morphName{f}_0 \maps E^1 \to E^2, \morphName{f}_1 \maps C^1 \to C^2)$
be an $E^1$-$E^2$ $\Ck$-morphism of $\seqname{A}^1$ to $\seqname{A}^2$
in the coordinate spaces $C^1$, $C^2$. Then
$(\morphName{f}_0, \morphName{f}_1)$ is a $\Ck$ morphism of $S^1$ to
$S^2$.

Let $E^i$, $i=1,2$, be a topological space and $C^i$ be a linear space.
Then
\begin{multline}
\Man^\Ck_\Ar(E^1, C^1, E^2, C^2) \defeq
\Atl[\Ck]_\Ar(E^1, C^1, E^2, C^2)
\end{multline}

Let $\seqname{E}$ be a set of topological spaces and $\seqname{C}$ a
set of linear spaces. Then
\begin{equation}
\Man^\Ck_\Ar(\seqname{E}, \seqname{C}) \defeq
\maximal{\Atl[\Ck]_\Ar}(\seqname{E}, \seqname{C})
\end{equation}
\begin{equation}
\Man^\Ck(\seqname{E}, \seqname{C}) \defeq
\maximal{\Atl[\Ck]}(\seqname{E}, \seqname{C})
\end{equation}
\end{definition}

\begin{theorem}[Categories of \ensuremath{\Ck} manifolds]
Let $\seqname{E}$ be a set of topological spaces and $\seqname{C}$ a set
of linear spaces. Then $\Man^\Ck(\seqname{E}, \seqname{C})$ is a
category and the identity morphism of $(\seqname{A}, E, C)$ is an
identity morphism.

\begin{proof}
{
  \showlabelsinline
  The result follows directly from \cref{def:man,def:manmorph} above and
}
\pagecref{lem:Ck-ATLiscat}.
\end{proof}
\end{theorem}

\begin{definition}[Functors from \ensuremath{\Ck} manifolds to Local 2-$\emptyset$ coordinate spaces]
\label{def:CktoLCS}
Let $\seqname{E}$ be a set of topological spaces,
$\seqname{C}$ be a set of $\Ck$ linear spaces,
$\catName{E}$ and $\hat{\catName{E}}$ model categories,
$\catName{C}$ and $\hat{\catName{C}}$ $\Ck$ linear model categories,
$\catSeqName{M} \defeq (\catName{E}, \catName{C})$,
$
  \hat{\catSeqName{M}} \defeq
  (\hat{\catName{E}}, \hat{\catName{C}})
$, \\*
$
  \Triv{\catSeqName{M}} \defeq
  \biggl ( \Trivcat{\seqname{E}}, \Trivcat[\Ck-]{\seqname{C}} \biggr )
$.
$
  \optriv{\catSeqName{M}} \defeq
  \biggl ( \Trivcat{\seqname{E}}, \optriv[\Ck-]{\seqname{C}} \biggr )
$.
$
  \Triv{\catSeqName{M}} \SUBCAT[full-] \catSeqName{M}
  \SUBCAT[full-] \hat{\catSeqName{M}}
$,
$\optriv{\catSeqName{M}} \SUBCAT \hat{\catSeqName{M}}$.

Let $E^i \in \seqname{E}$, $i=1,2$,
$C^i \in \seqname{C}$,
$\seqname{M}^i \defeq \left ( \Triv{E^i}, \Triv[\Ck-]{C^i} \right )$,
$
  \Singcat{\catSeqName{M}^i} \defeq
  \biggl (
    \singcat{\seqname{M}^i_0},
    \singcat[\Ck-]{\seqname{M}^i_1}
  \biggr )
$,
$\seqname{A}^i$ a maximal $\Ck$-atlas of $E^i$ in $C^i$,
$
  \Functor^{\minimal{\Ck}}_1 (\seqname{A}^i, E^i, C^i) \objin
  \hat{\catName{E}}
$, \\*
$
  \Functor^{\minimal{\Ck}}_2 (\seqname{A}^i, E^i, C^i) \objin
  \hat{\catName{C}}
$,
$(E^i, C^i, \seqname{A}^i)$ a $\Ck$ manifold,
$\seqname{S}^i \defeq (\seqname{A}^i, E^i, C^i)$,
$
  \minimal{\seqname{M}^i} \defeq \\
    \Bigl (
      \Functor^{\minimal{\Ck}}_1 (\seqname{A}^i, E^i, C^i),
      \Functor^{\minimal{\Ck}}_2 (\seqname{A}^i, E^i, C^i)
    \Bigr )
$
and
$
  \Singcat{\Triv{\catSeqName{M}^i}} \defeq
  (
    \singcat{\Triv{\seqname{M}^i_0}},
    \singcat[\Ck-]{\Triv{\seqname{M}^i_1}}
  )
$,
Then
\begin{equation}
\Functor^\Ck_{\Man,\LCS} \seqname{S}^i
\defeq
  \left (
    \Singcat{\catSeqName{M}^i},
    \seqname{M}^i,
    \seqname{A}^i,
    \emptyset,
    \emptyset
  \right )
\end{equation}
\begin{equation}
\Functor^{\Ck,\catSeqName{M}}_{\Man,\LCS} \seqname{S}^i
\defeq
  \left (
    \catSeqName{M},
    \seqname{M}^i,
    \seqname{A}^i,
    \emptyset,
    \emptyset
  \right )
\end{equation}
\begin{equation}
\Functor^{\minimal{\Ck}}_{\Man,\LCS} \seqname{S}^i
\defeq
  \left (
    \Singcat{\Triv{\catSeqName{M}^i}},
    \minimal{\seqname{M}^i},
    \seqname{A}^i,
    \emptyset,
    \emptyset
  \right )
\end{equation}
\begin{equation}
\Functor^{\minimal{\Ck},\hat{\catSeqName{M}}}_{\Man,\LCS} \seqname{S}^i
\defeq
  \left (
    \hat{\catSeqName{M}},
    \minimal{\seqname{M}^i},
    \seqname{A}^i,
    \emptyset,
    \emptyset
  \right )
\end{equation}
\begin{multline}
\Functor^\Ck_{\Man,\LCS}
  \bigl (
    (
      \morphName{f}_0,
      \morphName{f}_1
    ),
    \seqname{S}^1,
    \seqname{S}^2
  \bigr )
\defeq
  \Bigl (
    (
      \morphName{f}_0,
      \morphName{f}_1
    ),
    \Functor^\Ck_{\Man,\LCS} \seqname{S}^1,
    \Functor^\Ck_{\Man,\LCS} \seqname{S}^2
  \Bigr )
\end{multline}
\begin{multline}
\Functor^{\Ck,\catSeqName{M}}_{\Man,\LCS}
  \bigl (
    (
      \morphName{f}_0,
      \morphName{f}_1
    ),
    \seqname{S}^1,
    \seqname{S}^2
  \bigr )
\defeq
  \Bigl (
    (
      \morphName{f}_0,
      \morphName{f}_1
    ),
    \Functor^{\Ck,\catSeqName{M}}_{\Man,\LCS} \seqname{S}^1,
    \Functor^{\Ck,\catSeqName{M}}_{\Man,\LCS} \seqname{S}^2
  \Bigr )
\end{multline}
\begin{multline}
\Functor^{\minimal{\Ck}}_{\Man,\LCS}
  \bigl (
    (
      \morphName{f}_0,
      \morphName{f}_1
    ),
    \seqname{S}^1,
    \seqname{S}^2
  \bigr )
\defeq
  \biggl (
    \bigl (
      \morphName{f}_0,
      \morphName{f}_1
    \bigr ),
    \Functor^{\minimal{\Ck}}_{\Man,\LCS} \seqname{S}^1,
    \Functor^{\minimal{\Ck}}_{\Man,\LCS} \seqname{S}^2
  \biggr )
\end{multline}
\begin{multline}
\Functor^{\minimal{\Ck},\catSeqName{M}}_{\Man,\LCS}
  \bigl (
    (
      \morphName{f}_0,
      \morphName{f}_1
    ),
    \seqname{S}^1,
    \seqname{S}^2
  \bigr )
\defeq
  \Biggl (
    \bigl (
      \morphName{f}_0,
      \morphName{f}_1
    \bigr ),
    \Functor^{\minimal{\Ck},\catSeqName{M}}_{\Man,\LCS} \seqname{S}^1,
    \Functor^{\minimal{\Ck},\catSeqName{M}}_{\Man,\LCS} \seqname{S}^2
  \Biggr )
\end{multline}

\begin{multline}
\LCS^\Ck_\Ob(\seqname{E}, \seqname{C}) \defeq
\set
  {
    \Functor^\Ck_{\Man,\LCS} (\seqname{A}, E, C)
  }%
  [
    E \in \seqname{E},
    C \in \seqname{C},
    {\maximal{\isAtl[\Ck]_\Ob}(\seqname{A}, E, C)}
  ]*
\end{multline}
\begin{multline}
\LCS^\Ck_\Ar(\seqname{E}, \seqname{C}) \defeq
\set
  {
    \Functor^\Ck_{\Man,\LCS}
      \bigl (
        (
          \morphName{f}_0,
          \morphName{f}_1
        ),
        (
          \seqname{A}^1,
          E^1,
          C^1
        ),
        (
          \seqname{A}^2,
          E^2,
          C^2
        )
      \bigr )
  }%
  [
    E^i \in \seqname{E},
    C^i \in \seqname{C},
    {
      \maximal{\isAtl[\Ck]_\Ar}
        (
          \seqname{A}^1, E^1, C^1,
          \seqname{A}^2, E^2, C^2,
          \morphName{f}_0,\morphName{f}_1
        )
    }
  ]*
\end{multline}
\begin{equation}
\LCS^\Ck(\seqname{E}, \seqname{C}) \defeq
\Bigl (
  \LCS^\Ck_\Ob(\seqname{E}, \seqname{C}),
  \LCS^\Ck_\Ar(\seqname{E}, \seqname{C}),
  \morphCompose[A]
\Bigr )
\end{equation}
\begin{multline}
\LCS^{\Ck,\catSeqName{M}}_\Ob(\seqname{E}, \seqname{C}) \defeq
\set
  {
    \Functor^{\Ck,\catSeqName{M}}_{\Man,\LCS} (\seqname{A}, E, C)
  }%
  [
    E \in \seqname{E},
    C \in \seqname{C},
    {\maximal{\isAtl[\Ck]_\Ob}(\seqname{A}, E, C)}
  ]*
\end{multline}
\begin{multline}
\LCS^{\Ck,\catSeqName{M}}_\Ar(\seqname{E}, \seqname{C}) \defeq
\set
  {
    \Functor^{\Ck,\catSeqName{M}}_{\Man,\LCS}
      \bigl (
        (
          \morphName{f}_0,
          \morphName{f}_1
        ),
        (
          \seqname{A}^1,
          E^1,
          C^1
        ),
        (
          \seqname{A}^2,
          E^2,
          C^2
        )
      \bigr )
  }%
  [
    E^i \in \seqname{E},
    C^i \in \seqname{C},
    {
      \maximal{\isAtl[\Ck]_\Ar}
        (
          \seqname{A}^1, E^1, C^1,
          \seqname{A}^2, E^2, C^2,
          \morphName{f}_0,\morphName{f}_1
        )
    }
  ]*
\end{multline}
\begin{equation}
\LCS^{\Ck,\catSeqName{M}}(\seqname{E}, \seqname{C}) \defeq
\Bigl (
  \LCS^{\Ck,\catSeqName{M}}_\Ob(\seqname{E}, \seqname{C}),
  \LCS^{\Ck,\catSeqName{M}}_\Ar(\seqname{E}, \seqname{C}),
  \morphCompose[A]
\Bigr )
\end{equation}
\begin{multline}
\LCS^{\minimal{\Ck}}_\Ob(\seqname{E}, \seqname{C}) \defeq
\set
  {
    \Functor^{\minimal{\Ck}}_{\Man,\LCS} \seqname{S}
  }%
  [
    {\seqname{S} \defeq (\seqname{A}, E, C)},
    E \in \seqname{E},
    C \in \seqname{C},
    {\maximal{\isAtl[\Ck]_\Ob}(\seqname{A}, E, C)}
  ]*
\end{multline}
\begin{multline}
\LCS^{\minimal{\Ck}}_\Ar(\seqname{E}, \seqname{C}) \defeq \\
\set
  {
    {
      \Functor^{\minimal{\Ck}}_{\Man,\LCS}
        \bigl (
          (
            \morphName{f}_0,
            \morphName{f}_1
          ),
          (
            \seqname{A}^1,
            E^1,
            C^1
          ),
          (
              \seqname{A}^2,
              E^2,
              C^2
          )
        \bigr )
    }
  }%
  [
    E^i \in \seqname{E},
    C^i \in \seqname{C},
    {
      \maximal{\isAtl[\Ck]_\Ar}
        (
          \seqname{A}^1, E^1, C^1,
          \seqname{A}^2, E^2, C^2,
          \morphName{f}_0, \morphName{f}_1
        )
    }
  ]*
\end{multline}
\begin{equation}
\LCS^{\minimal{\Ck}}(\seqname{E}, \seqname{C}) \defeq
\Bigl (
  \LCS^{\minimal{\Ck}}_\Ob(\seqname{E}, \seqname{C}),
  \LCS^{\minimal{\Ck}}_\Ar(\seqname{E}, \seqname{C}),
  \morphCompose[A]
\Bigr )
\end{equation}
\begin{multline}
\LCS^{\minimal{\Ck},\hat{\catSeqName{M}}}_\Ob(\seqname{E}, \seqname{C}) \defeq
\set
  {
    \Functor^{\minimal{\Ck}}_{\Man,\LCS} \seqname{S}
  }%
  [
    {\seqname{S} \defeq (\seqname{A}, E, C)},
    E \in \seqname{E},
    C \in \seqname{C},
    {\maximal{\isAtl[\Ck]_\Ob}(\seqname{A}, E, C)}
  ]*
\end{multline}
\begin{multline}
\LCS^{\minimal{\Ck},\hat{\catSeqName{M}}}_\Ar(\seqname{E}, \seqname{C}) \defeq \\
\set
  {
    {
      \Functor^{\minimal{\Ck},\hat{\catSeqName{M}}}_{\Man,\LCS}
        \bigl (
          (
            \morphName{f}_0,
            \morphName{f}_1
          ),
          (
            \seqname{A}^1,
            E^1,
            C^1
          ),
          (
              \seqname{A}^2,
              E^2,
              C^2
          )
        \bigr )
    }
  }%
  [
    E^i \in \seqname{E},
    C^i \in \seqname{C},
    {
      \maximal{\isAtl[\Ck]_\Ar}
        (
          \seqname{A}^1, E^1, C^1,
          \seqname{A}^2, E^2, C^2,
          \morphName{f}_0, \morphName{f}_1
        )
    }
  ]*
\end{multline}
\begin{equation}
\LCS^{\minimal{\Ck},\hat{\catSeqName{M}}}(\seqname{E}, \seqname{C}) \defeq
\Bigl (
  \LCS^{\minimal{\Ck},\hat{\catSeqName{M}}}_\Ob(\seqname{E}, \seqname{C}),
  \LCS^{\minimal{\Ck},\hat{\catSeqName{M}}}_\Ar(\seqname{E}, \seqname{C}),
  \morphCompose[A]
\Bigr )
\end{equation}
\end{definition}

\begin{lemma}[$(0)$-equivalence of local 2-$\emptyset$ coordinate spaces]
\label{lem:CktoLCSequiv}
Let
\begin{enumerate}
\item
$\seqname{E}$ be a set of topological spaces,
$\seqname{C}$ be a set of $\Ck$ linear spaces
\item
$\catName{E}$ and $\hat{\catName{E}} be $ model categories,
$\catName{C}$ and $\hat{\catName{C}}$ be $\Ck$ linear model categories
\item
$\catSeqName{M} \defeq (\catName{E}, \catName{C})$,
$
  \hat{\catSeqName{M}} \defeq
  (\hat{\catName{E}}, \hat{\catName{C}})
$
\item
$
  \Triv{\catSeqName{M}} \defeq
  \biggl ( \Trivcat{\seqname{E}}, \Trivcat[\Ck-]{\seqname{C}} \biggr )
$,
$
  \optriv{\catSeqName{M}} \defeq
  \biggl ( \Trivcat{\seqname{E}}, \optriv[\Ck-]{\seqname{C}} \biggr )
$
\item
$
  \Triv{\catSeqName{M}} \SUBCAT[full-] \catSeqName{M}
  \SUBCAT[full-] \hat{\catSeqName{M}}
$,
$\optriv{\catSeqName{M}} \SUBCAT \hat{\catSeqName{M}}$
\item
$E^i \in \seqname{E}$, $i=1,2$,
$C^i \in \seqname{C}$
\item
$\seqname{M}^i \defeq \left ( \Triv{E^i}, \Triv[\Ck-]{C^i} \right )$,
$
  \Singcat{\catSeqName{M}^i} \defeq
  \biggl (
    \singcat{\seqname{M}^i_0},
    \singcat[\Ck-]{\seqname{M}^i_1}
  \biggr )
$
\item
$\seqname{A}^i$ be a maximal $\Ck$-atlas of $E^i$ in $C^i$
\item
$
  \Functor^{\minimal{\Ck}}_1 (\seqname{A}^i, E^i, C^i) \objin
  \hat{\catName{E}}
$,
$
  \Functor^{\minimal{\Ck}}_2 (\seqname{A}^i, E^i, C^i) \objin
  \hat{\catName{C}}
$
\item
$(E^i, C^i, \seqname{A}^i)$ be a $\Ck$ manifold
\item
$\seqname{S}^i \defeq (\seqname{A}^i, E^i, C^i)$
\item
$
  \minimal{\seqname{M}^i} \defeq
    \Bigl (
      \Functor^{\minimal{\Ck}}_1 (\seqname{A}^i, E^i, C^i),
      \Functor^{\minimal{\Ck}}_2 (\seqname{A}^i, E^i, C^i)
    \Bigr )
$,
$
  \Singcat{\Triv{\catSeqName{M}^i}} \defeq
    \Bigl (
      \singcat{\Triv{\seqname{M}^i_0}},
      \singcat[\Ck-]{\Triv{\seqname{M}^i_1}}
    \Bigr )
$
\item
$
  \morphSeqName{f} \defeq
  (
    \morphName{f}_0 \maps E^1 \to E^2,
    \morphName{f}_1 \maps C^1 \to C^2
  )
$
and \\*
$
  \morphSeqName{g} \defeq
  (
    \morphName{g}_0 \maps E^1 \to E^2,
    \morphName{g}_1 \maps C^1 \to C^2
  )
$
be $E^1$-$E^2$ $\Ck$ morphisms of $\seqname{A}^1$ to
$\seqname{A}^2$ in the coordinate spaces $C^1$, $C^2$
\end{enumerate}

$
  \morphSeqName{f} \maps
  \Functor^\Ck_{\Man,\LCS} \seqname{S}^1 \to
  \Functor^\Ck_{\Man,\LCS} \seqname{S}^2
$
is $(0)$-equivalent to \\*
$
  \morphSeqName{g} \maps
  \Functor^\Ck_{\Man,\LCS} \seqname{S}^1 \to
  \Functor^\Ck_{\Man,\LCS} \seqname{S}^2
$
iff $\morphSeqName{f}$ is $\Ck$-equivalent to $\morphSeqName{g}$.

\begin{proof}
Both equivalences are defined as $\morphName{f}_0 = \morphName{g}_0$.
\end{proof}

\end{lemma}

\begin{theorem}[Functors from \ensuremath{\Ck} manifolds to local 2-$\emptyset$ coordinate spaces]
\label{the:CktoLCS}
Let $\seqname{E}$ be a set of topological spaces,
$\seqname{C}$ be a set of $\Ck$ linear spaces,
$\catName{E}$ and $\hat{\catName{E}}$ model categories,
$\catName{C}$ and $\hat{\catName{C}}$ $\Ck$ linear model categories,
$\catSeqName{M} \defeq (\catName{E}, \catName{C})$,
$
  \hat{\catSeqName{M}} \defeq
  (\hat{\catName{E}}, \hat{\catName{C}})
$,
$
  \Triv{\catSeqName{M}} \defeq
  \Bigl ( \Trivcat{\seqname{E}}, \Trivcat[\Ck-]{\seqname{C}} \Bigr )
$.
$
  \optriv{\catSeqName{M}} \defeq
  \Bigl ( \Trivcat{\seqname{E}}, \optriv[\Ck-]{\seqname{C}} \Bigr )
$.
$
  \Triv{\catSeqName{M}} \SUBCAT \catSeqName{M}
  \SUBCAT \hat{\catSeqName{M}}
$,
$\optriv{\catSeqName{M}} \SUBCAT \hat{\catSeqName{M}}$.

Then
$\LCS^\Ck(\seqname{E}, \seqname{C})$,
$\LCS^{\Ck,\catSeqName{M}}(\seqname{E}, \seqname{C})$,
$\LCS^{\minimal{\Ck}}(\seqname{E}, \seqname{C})$ and
$\LCS^{\minimal{\Ck},\hat{\catSeqName{M}}}(\seqname{E}, \seqname{C})$
are categories and the identity morphism of
$
   \seqname{L}^\alpha \defeq
    \bigl (
      (\catName{E}, \catName{C}),
      (\seqname{E}^\alpha, \seqname{C}^\alpha),
      \seqname{A}^\alpha,
      \emptyset,
      \emptyset
    \bigr )
$
is that given in \pagecref{def:LCSmorphAr}:
$
  \Id[{\seqname{L}^\alpha}]
  \defeq
  \Bigl (
    \bigl ( \Id[{\seqname{E}^\alpha}], \Id[{\seqname{C}^\alpha}] \bigr ),
    \seqname{L}^\alpha,
    \seqname{L}^\alpha
  \Bigr )
$.

\begin{proof}
$\LCS^\Ck(\seqname{E}, \seqname{C})$,
$\LCS^{\Ck,\catSeqName{M}}(\seqname{E}, \seqname{C})$,
$\LCS^{\minimal{\Ck}}(\seqname{E}, \seqname{C})$ and
$\LCS^{\minimal{\Ck},\hat{\catSeqName{M}}}(\seqname{E}, \seqname{C})$
satisfy the definition of a category:
\begin{description}
\item[Composition:]
Let
$
  m^i \defeq
  \Functor^{\Ck}_{\Man,\LCS}
    \bigl (
      (
        \morphName{f}^i_0,
        \morphName{f}^i_1
      ),
      \seqname{S}^i,
      \seqname{S}^{i+1}
    \bigr )
$
be a morphism of $\LCS^{\Ck}(\seqname{E}, \seqname{C})$.
$
  m^2 \morphCompose[A] m^1 =
  \Functor^{\Ck}_{\Man,\LCS}
    \bigl (
      (
        \morphName{f}^2_0 \morphCompose \morphName{f}^1_0,
        \morphName{f}^2_1 \morphCompose \morphName{f}^2_1
      ),
      \seqname{S}^1,
      \seqname{S}^3
    \bigr )
$.
By
\cref{M-morphcomp:morphcompmorph} of
{
  \showlabelsinline
  \fullcref{lem:M-ATLmorphcomp} on
  \cpageref{M-morphcomp:morphcompmorph}
}
$
  (
    \morphName{f}^2_0 \morphCompose \morphName{f}^1_0,
    \morphName{f}^2_1 \morphCompose \morphName{f}^2_1
  )
$
is an m-atlas morphism.
\item[Associativity:]
Composition is associative by \pagecref{lem:labcomp}.
\item[Unit:]
$\Id[{\seqname{L}^\alpha}]$ is an identity morphism by
\cref{lem:labcomp}.
\end{description}

The same proof applies to $\LCS^{\Ck,\catSeqName{M}}(\seqname{E}, \seqname{C})$,
$\LCS^{\minimal{\Ck}}(\seqname{E}, \seqname{C})$ and
$\LCS^{\minimal{\Ck},\hat{\catSeqName{M}}}(\seqname{E}, \seqname{C})$.
\end{proof}

Let
$E^i \in \seqname{E}$, $i \in [1,3]$, $C^i \in \seqname{C}$,
$\seqname{M}^i \defeq \Bigl ( \Triv{E^i}, \Triv[\Ck-]{C^i} \Bigr )$,
$
  \Singcat{\catSeqName{M}^i} \defeq
  (
    \singcat{\seqname{M}^i_0},
    \singcat[\Ck-]{\seqname{M}^i_1}
  )
$,
$\seqname{A}^i$ a maximal $\Ck$-atlas of $E^i$ in $C^i$,
$
  \Functor^{\minimal{\Ck}}_1 (\seqname{A}^i, E^i, C^i) \objin
  \hat{\catName{E}}
$,
$
  \Functor^{\minimal{\Ck}}_2 (\seqname{A}^i, E^i, C^i) \objin
  \hat{\catName{C}}
$,
$(E^i, C^i, \seqname{A}^i)$ a $\Ck$ manifold,
$\seqname{S}^i \defeq (\seqname{A}^i,E^i, C^i)$,
$
  \minimal{\seqname{M}^i} \defeq
    \Bigl (
      \Functor^{\minimal{\Ck}}_1 (\seqname{A}^i, E^i, C^i),
      \Functor^{\minimal{\Ck}}_2 (\seqname{A}^i, E^i, C^i)
    \Bigr )
$, \\
$
  \Singcat{\Triv{\catSeqName{M}^i}} \defeq
  \Bigl (
    \singcat{\Triv{\seqname{M}^i_0}},
    \singcat[\Ck-]{\Triv{\seqname{M}^i_1}}
  \Bigr )
$,
$
\seqname{L}^i \defeq
\Functor^{\Ck}_{\Man,\LCS} \seqname{S}^i =
\Bigl (
  \bigl (
    \Trivcat{E^i}, \Trivcat[\Ck-]{C^i}
  \bigr ),
  \bigl (
    \Triv{E^i}, \Triv[\Ck-]{C^i}
  \bigr ),
  \seqname{A}^i,
  \emptyset,
  \emptyset
\Bigr )
$,
$
\seqname{L}^{i,\catSeqName{M}} \defeq
\Functor^{\Ck,\catSeqName{M}}_{\Man,\LCS} \seqname{S}^i =
\Bigl (
  \catSeqName{M},
  \bigl (
    \Triv{E^i}, \Triv[\Ck-]{C^i}
  \bigr ),
  \seqname{A}^i,
  \emptyset,
  \emptyset
\Bigr )
$,
$
\Triv{\seqname{L}^i} \defeq
\Bigl (
  \singcat{\Triv{\catSeqName{M}^i}},
  \bigl (
    \Triv{E^i}, \optriv[\Ck-]{C^i}
  \bigr ),
  \seqname{A}^i,
  \emptyset,
  \emptyset
\Bigr )
$,
$
\Triv{\seqname{L}^{i,\hat{\catSeqName{M}}}} \defeq
\Bigl (
  \hat{\catSeqName{M}},
  \bigl (
    \Triv{E^i}, \optriv[\Ck-]{C^i}
  \bigr ),
  \seqname{A}^i,
  \emptyset,
  \emptyset
\Bigr )
$,
$\morphSeqName{f}^i \defeq (\morphName{f}^i_0, \morphName{f}^i_1)$,
$\morphName{f}^i_0 \maps E^i \to E^{i+1}$ continuous and
$\morphName{f}^i_1 \maps C^i \to C^{i+1}$ $\Ck$.

$\Functor^\Ck_{\Man,\LCS} \seqname{S}^i$ is a local
$\Singcat{\catSeqName{M}^i}$-$\emptyset$ coordinate space and
$\Functor^{\Ck,\catSeqName{M}}_{\Man,\LCS} \seqname{S}^i$ is a local
$\catSeqName{M}$-$\emptyset$ coordinate space.

\begin{proof}
$\Functor^{\Ck,\catSeqName{M}}_{\Man,\LCS} \seqname{S}^i$
satisfies the conditions of \pagecref{def:LCS}:
\begin{enumerate}
\item Model categories:
$\Triv{E^i}$ and $\Triv[\Ck-]{C^i}$ are model spaces
\item
$
\biggl (
  \catSeqName{M},
  \Bigl ( \Triv{E^i}, \Triv[\Ck-]{C^i} \Bigr )
  \emptyset,
  \emptyset
\biggr )
$
satisfies the conditions of \pagecref{def:pre}
\item Maximal m-atlas:
$\seqname{A}^i$ is a maximal $\Ck$ atlas by hypothesis and a
maximal atlas of $\Triv{E^1}$ in the coordinate space $\Triv[\Ck-]{C^i}$
by construction.
\item Constraint functions:
There are no adjunct functions so there are no constraint functions.
\end{enumerate}

$\Singcat{\catSeqName{M}^i} \SUBCAT \catSeqName{M}$, so
$\Functor^{\Ck,\catSeqName{M}}_{\Man,\LCS} \seqname{S}^i$ is a local
$\catSeqName{M}$-$\emptyset$ coordinate space by
\pagecref{lem:LCS}.
\end{proof}

$\Functor^{\minimal{\Ck}}_{\Man,\LCS} \seqname{S}^i$ is a local
$\Singcat{\Triv{\catSeqName{M}^i}}$-$\emptyset$ coordinate space and
$
  \Functor^{\minimal{\Ck},\hat{\catSeqName{M}}}_{\Man,\LCS}
    \seqname{S}^i
$
is a local $\hat{\catSeqName{M}}$-$\emptyset$ coordinate space.

\begin{proof}
$\Functor^{\minimal{\Ck}}_{\Man,\LCS} \seqname{S}^i$
satisfies the conditions of \pagecref{def:LCS}:
\begin{enumerate}
\item Model categories:
$\Triv{E^1}$ and $\optriv[\Ck-]{C^i}$ are model spaces
\item
$
  \Bigl (
    \bigl ( \Triv{E^i}, \optriv[\Ck-]{C^i} \bigr ), \emptyset, \emptyset
  \Bigr )
$
satisfies the conditions of \pagecref{def:pre}
\item Maximal m-atlas:
$\seqname{A}^i$ is a maximal $\Ck$ atlas by hypothesis and a
maximal atlas of $\Triv{E^1}$ in the coordinate space
$\optriv[\Ck-]{C^i}$ by construction.
\item Constraint functions:
There are no adjunct functions so there are no constraint functions.
\end{enumerate}

$\Singcat{\Triv{\catSeqName{M}^i}} \SUBCAT \hat{\catSeqName{M}}$, so
$\Functor^{\Ck,\hat{\catSeqName{M}}}_{\Man,\LCS} \seqname{S}^i$ is a local
$\hat{\catSeqName{M}}$-$\emptyset$ coordinate space by \cref{lem:LCS}.
\end{proof}

$(\morphName{f}^1_0, \morphName{f}^1_1)$ is a morphism from
$\seqname{L}^1$ to $\seqname{L}^2$ and a strict morphism from
$\seqname{L}^{1,\catSeqName{M}}$ to $\seqname{L}^{2,\catSeqName{M}}$.

\begin{proof}
\leavevmode
\begin{description}
\item[Prestructure morphism:]
$\morphSeqName{f}^1$ $\Sigma$-commutes with $\emptyset,\emptyset$.
$\morphName{f}^1_0$ is continuous, hence a morphism of $\catName{E}$.
$\morphName{f}^1_1$ is $\Ck$, hence a morphism of $\catName{C}$.

\item[m-atlas morphism:]
$(\morphName{f}^1_0, \morphName{f}^1_1)$ is a $\Triv{E^1}$-$\Triv{E^2}$
m-atlas morphism of $\seqname{A}^1$ to $\seqname{A}^2$ in the coordinate
spaces $\Triv[\Ck-]{C^1}$, $\Triv[\Ck-]{C^i}$.

$(\morphName{f}^1_0, \morphName{f}^1_1)$ is a strict
$\singcat{\Triv{E^1}}$-$\singcat{\Triv{E^2}}$-%
$\Triv{E^1}$-$\Triv{E^2}$-%
$\singcat[\Ck-]{\Triv[\Ck-]{C^1}}$-$\singcat[\Ck-]{\Triv[\Ck-]{C^1}}$
m-atlas morphism of $\seqname{A}^1$ to $\seqname{A}^2$ in the coordinate
spaces $\Triv[\Ck-]{C^1}$, $\Triv[\Ck-]{C^i}$. by
\pagecref{lem:Ck-ATLmorph}.
\end{description}
\end{proof}

$
\Functor^{\minimal{\Ck}}_{\Man,\LCS}
  \bigl (
    (\morphName{f}^1_0, \morphName{f}^1_1), \seqname{S}^1, \seqname{S}^2
  \bigr )
$
is a morphism from $\Triv{\seqname{L}^1}$ to
$\Triv{\seqname{L}^{2,\catSeqName{M}}}$ and \\
$
\Functor^{\minimal{\Ck}}_{\Man,\LCS}
  \bigl (
    (\morphName{f}^1_0, \morphName{f}^1_1), \seqname{S}^1, \seqname{S}^2
  \bigr )
$
is a strict morphism from
$\Triv{\seqname{L}^{1,\hat{\catSeqName{M}}}}$ to
$\Triv{\seqname{L}^{2,\hat{\catSeqName{M}}}}$.

\begin{proof}
\leavevmode
\begin{description}
\item[Prestructure morphism:]
$(\morphName{f}^1_0, \morphName{f}^1_1)$ $\Sigma$-commutes with $\emptyset,\emptyset$.

\item[m-atlas morphism:]
$(\morphName{f}^1_0, \morphName{f}^1_1)$ is an m-atlas morphism from
$(\Triv{E^1}, \optriv[\Ck-]{C^1})$ to
$(\Triv{E^2}, \optriv[\Ck-]{C^2})$.
\end{description}
\end{proof}

$\Functor^{\Ck,\catSeqName{M}}_{\Man,\LCS}$ is a functor from
$\Man^\Ck(\seqname{E}, \seqname{C})$ to
$\LCS^{\Ck,\catSeqName{M}}(\seqname{E}. \seqname{C})$.

\begin{proof}
\leavevmode
\begin{description}
\item[Preservation of endpoints:]
\begin{align*}
  \Functor^{\Ck,\catSeqName{M}}_{\Man,\LCS}
    \bigl (
      (
        \morphSeqName{f}^1_0,
        \morphSeqName{f}^1_1
      ),
      \seqname{S}^1,
      \seqname{S}^2
    \bigr )
  & =
    \bigl (
      (
        \morphSeqName{f}^1_0,
        \morphSeqName{f}^1_1
      ),
      \seqname{L}^i,
      \seqname{L}^i
    \bigr )
\\*
  \Functor^{\Ck,\catSeqName{M}}_{\Man,\LCS} \seqname{S}^i
  & =
    \left (
      \Bigl ( \Triv{E^i}, \Triv[\Ck-]{C^i} \Bigr ),
      \seqname{A}^i,
      \emptyset,
      \emptyset
    \right )
\\*
  & =
    \seqname{L}^i
\end{align*}
\item[Composition:]
$
  \Functor^{\Ck,\catSeqName{M}}_{\Man,\LCS}
    \Bigl (
      \bigl (
        (\morphSeqName{f}^2_0, \morphSeqName{f}^2_1),
        \seqname{S}^2,
        \seqname{S}^3
      \bigr )
      \morphCompose[A]
      \bigl (
        (\morphSeqName{f}^1_0, \morphSeqName{f}^1_1),
        \seqname{S}^1,
        \seqname{S}^2
      \bigr )
    \Bigr )
  = \\
  \Functor^{\Ck,\catSeqName{M}}_{\Man,\LCS}
    \Bigl (
      \bigl (
        (
          \morphSeqName{f}^2_0 \morphCompose \morphSeqName{f}^1_0,
          \morphSeqName{f}^2_1 \morphCompose \morphSeqName{f}^1_1
        ),
        \seqname{S}^1,
        \seqname{S}^3
      \bigr )
    \Bigr )
  = \\
  \bigl (
    (
      \morphSeqName{f}^2_0 \morphCompose \morphSeqName{f}^1_0,
      \morphSeqName{f}^2_1 \morphCompose \morphSeqName{f}^1_1
    ),
    \seqname{L}^1,
    \seqname{L}^3
  \bigr )
  = \\
  \Functor^{\Ck,\catSeqName{M}}_{\Man,\LCS}
    \bigl (
      (
        \morphSeqName{f}^2_0,
        \morphSeqName{f}^2_1
      ),
      \seqname{S}^2,
      \seqname{S}^3
    \bigr )
  \morphCompose[A]
  \Functor^{\Ck,\catSeqName{M}}_{\Man,\LCS}
    \bigl (
      (
        \morphSeqName{f}^1_0,
        \morphSeqName{f}^1_1
      ),
      \seqname{S}^1,
      \seqname{S}^2
    \bigr )
$
\item[Identity:]
$
\Functor^{\Ck,\catSeqName{M}}_{\Man,\LCS} (\Id[{\seqname{S}^i})] =
\Functor^{\Ck,\catSeqName{M}}_{\Man,\LCS}
  \Bigl (
    (\Id[E^i], \Id[C^i]),
    \seqname{S}^i,
    \seqname{S}^i
  \Bigr ) =\\
  \biggl (
   \Bigl ( \Id[{\Triv{E^i}}], \Id[{\Triv[\Ck-]{C^i}}] \Bigr ),
    \seqname{L}^1,
    \seqname{L}^1
  \biggr )
= \\
%\Id[{\seqname{L}^1}] = \Id[{\Functor^{\Ck,\catSeqName{M}}_{\Man,\LCS}(\seqname{S}^i)}]
$
\end{description}
\end{proof}

$\Functor^{\minimal{\Ck}}_{\Man,\LCS}$ is a functor from
$\Man^\Ck(\seqname{E}, \seqname{C})$ to
$\LCS^{\minimal{\Ck}}(\seqname{E}. \seqname{C})$.

\begin{proof}
\leavevmode
\begin{description}
\item[Preservation of endpoints:]
\begin{align*}
  \Functor^{\minimal{\Ck}}_{\Man,\LCS}
    \bigl (
      (
        \morphName{f}^1_0,
        \morphName{f}^1_1
      ),
      \seqname{S}^1,
      \seqname{S}^2
    \bigr )
  & =
    \bigl (
      (
        \morphName{f}^1_0,
        \morphName{f}^1_1
      ),
      \seqname{L}^i,
      \seqname{L}^i
    \bigr )
\\*
  \Functor^{\minimal{\Ck}}_{\Man,\LCS} \seqname{S}^i
  & =
    \left (
      \Bigl ( \Triv{E^i}, \optriv[\Ck-]{C^i} \Bigr ),
      \seqname{A}^i,
      \emptyset,
      \emptyset
    \right )
\\*
  & =
    \seqname{L}^i
\end{align*}
\item[Composition:]
$
  \Functor^{\minimal{\Ck}}_{\Man,\LCS}
    \Bigl (
      \bigl (
        (\morphName{f}^2_0, \morphName{f}^2_1),
        \seqname{S}^2,
        \seqname{S}^3
      \bigr )
      \morphCompose[A]
      \bigl (
        (\morphName{f}^1_0, \morphName{f}^1_1),
        \seqname{S}^1,
        \seqname{S}^2
      \bigr )
    \Bigr )
  = \\
  \Functor^{\minimal{\Ck}}_{\Man,\LCS}
    \Bigl (
      \bigl (
        (
          \morphName{f}^2_0 \morphCompose \morphName{f}^1_0,
          \morphName{f}^2_1 \morphCompose \morphName{f}^1_1
        ),
        \seqname{S}^1,
        \seqname{S}^3
      \bigr )
    \Bigr )
  = \\
  \biggl (
    (
      \morphName{f}^2_0 \morphCompose \morphName{f}^1_0,
      \morphName{f}^2_1 \morphCompose \morphName{f}^1_1
    ),
    \hat{\seqname{L}}^1,
    \hat{\seqname{L}}^3
  \biggr )
  = \\
  \Functor^{\minimal{\Ck}}_{\Man,\LCS}
    \bigl (
      (
        \morphName{f}^2_0,
        \morphName{f}^2_1
      ),
      \seqname{S}^2,
      \seqname{S}^3
    \bigr )
  \morphCompose[A]
  \Functor^{\minimal{\Ck}}_{\Man,\LCS}
    \bigl (
      (
        \morphName{f}^1_0,
        \morphName{f}^1_1
      ),
      \seqname{S}^1,
      \seqname{S}^2
    \bigr )
$
\item[Identity:]
$
\Functor^{\minimal{\Ck}}_{\Man,\LCS} (\Id[{\seqname{S}^i})] =
\Functor^{\minimal{\Ck}}_{\Man,\LCS}
  \bigl (
    (\Id[E^i], \Id[C^i]),
    \seqname{S}^i,
    \seqname{S}^i
  \bigr )
= \\
  \biggl (
   \bigl ( \Id[{\Triv{E^i}}], \Id[{\Triv[\Ck-]{C^i}}] \bigr ),
    \Bigl (
      \bigl ( \Triv{E^i}, \Triv[\Ck-]{C^i} \bigr ),
      \seqname{A}^i,
      \emptyset,
      \emptyset
    \Bigr ),
    \Bigl (
      \bigl ( \Triv{E^i}, \Triv[\Ck-]{C^i} \bigr ),
      \seqname{A}^i,
      \emptyset,
      \emptyset
    \Bigr )
  \biggr )
= \\
\Id%
  [{{
    \Bigl (
      \bigl ( \Triv{E^i}, \Triv[\Ck-]{C^i} \bigr ),
      \seqname{A}^i,
      \emptyset,
      \emptyset
    \Bigr )
  }}]
\defeq
\Id_{\Functor^{\minimal{\Ck}}_{\Man,\LCS}(\seqname{S}^i)}
$
\end{description}
\end{proof}
\end{theorem}

\begin{definition}[Functor from Local 2-$\emptyset$ coordinate spaces to \ensuremath{\Ck} manifolds]
\label{def:LCStoCk}
Let $\seqname{C^i} \defeq (C^i, \catName{C}^i)$,
$i=1,2$, be a $\Ck$ linear model space, $\seqname{E}^i$ a
model space,
$\catName{M}^i_0$ a model category,
$\catName{M}^i_1$ a $\Ck$ linear model category,
$\catSeqName{M}^i \defeq (\catName{M}^i_0, \catName{M}^i_1)$,
$\seqname{M}^i \defeq (\seqname{E}^i, \seqname{C}^i)$,
$\seqname{M}^i \seqin \catSeqName{M}^i$,
$\seqname{A}^i$ maximal m-atlases of $\seqname{E}^i$ in $\seqname{C}^i$,
$\morphName{f}_0 \maps \seqname{E}^1 \to \seqname{E}^2$,
$\morphName{f}_1 \maps \seqname{C}^1 \to \seqname{C}^2)$,
$\morphName{f} \defeq (\morphName{f}_0, \morphName{f}_1)$ and
$
  \seqname{L}^i \defeq
  (
    \catSeqName{M}^i,
    \seqname{M}^i,
    \seqname{A}^i,
    \emptyset,
    \emptyset
  )
$
a local 2-$\emptyset$ coordinate spaces. Then
\begin{equation}
\Functor^\Ck_{\LCS,\Man} \seqname{L}^i \defeq
(E^i, C^i, \seqname{A}^i)
\end{equation}
\begin{equation}
\Functor^\Ck_{\LCS,\Man}
  (
    \morphSeqName{f},
    \seqname{L}^1,
    \seqname{L}^2
  )
\defeq
  \bigl (
    \morphSeqName{f},
    (E^1, C^1, \seqname{A}^1),
    (E^2, C^2, \seqname{A}^2)
  \bigr )
\end{equation}
\end{definition}

\begin{theorem}[Functor from Local 2-$\emptyset$ coordinate spaces to \ensuremath{\Ck} manifolds]
\label{the:LCStoCk}
Let $\seqname{C^i} \defeq (C^i, \catName{C}^i)$,
$i=1,2$, be $\Ck$ linear model spaces,
$\seqname{E}^i \defeq (E^i, \catName{E}^i)$ model spaces,
$\seqname{A}^i$ maximal $\Ck$-atlases of $\seqname{E}^i$ in $\seqname{C}^i$,
$\seqname{M}^i \defeq (\seqname{E}^i, \seqname{C}^i)$,
$\catSeqName{M}^i \defeq \Singcat{\seqname{M}^i}$,
$\morphName{f}^i_0 \maps \seqname{E}^i \to \seqname{E}^{i+1}$
$\morphName{f}^i_1 \maps \seqname{C}^i \to \seqname{C}^{i+1})$ and
$
  \seqname{L}^i \defeq
  (
    \catSeqName{M}^i,
    \seqname{M}^i,
    \seqname{A}^i,
    \emptyset,
    \emptyset
  )
$
a local 2-$\emptyset$ coordinate space. Then
\begin{enumerate}
\item
$
  \Functor^\Ck_{\LCS,\Man} \seqname{L}^i
  \defeq (E^i, \seqname{C}^i, \seqname{A}^i)
$
is a $\Ck$ manifold.
\begin{proof}
$(E^i, \seqname{C}^i, \seqname{A}^i)$ satisfies the conditions of
\pagecref{def:man}:
$(E^i$ is a topological space, $\seqname{C}^i$ is a $\Ck$ linear space and
$\seqname{A}^i$ is a maximal $\Ck$-atlas $\seqname{E}^i$ in $\seqname{C}^i$.
\end{proof}
\item
$\Functor^\Ck_{\LCS,\Man}$ is a functor from
$\LCS^{\Ck,\catSeqName{M}}(\seqname{E}, \seqname{C})$ to
$\Man^\Ck(\seqname{E}, \seqname{C})$.
\begin{proof}
\leavevmode
Let $\seqname{S}^i=(\seqname{A}^i, E^i, \seqname{C}^i)$.
\begin{description}
\item[Preservation of endpoints:]
\begin{align*}
  \Functor^\Ck_{\LCS,\Man}
    \bigl
      (
        \morphName{f}^i_0,
        \morphName{f}^i_1
      ),
      \seqname{L}^i,
      \seqname{L}^{i+1}
    \bigr )
  & =
    \bigl
      (
        \morphName{f}^i_0,
        \morphName{f}^i_1
      ),
      \seqname{S}^i,
      \seqname{S}^{i+1}
    \bigr )
\\*
  \Functor^\Ck_{\LCS,\Man} (\seqname{L}^i)
  & =
  \seqname{S}^i
\end{align*}

\item[Composition:]
$
\Functor^\Ck_{\LCS,\Man}
  \Bigl (
    \bigl (
      (
        \morphName{f}^2_0,
        \morphName{f}^2_1
      ),
      \seqname{L}^2,
      \seqname{L}^3
    \bigr )
  \morphCompose[A]
    \bigl (
      (
        \morphName{f}^1_0,
        \morphName{f}^1_1
      ),
      \seqname{L}^1,
      \seqname{L}^2
    \bigr )
  \Bigr )
= \\
\Functor^\Ck_{\LCS,\Man}
  \Bigl (
    \bigl (
      (
        \morphName{f}^2_0 \morphCompose \morphName{f}^1_0,
        \morphName{f}^2_1 \morphCompose \morphName{f}^1_1
      ),
      \seqname{L}^1,
      \seqname{L}^3
    \bigr )
  \Bigr )
= \\
  \Bigl (
    \bigl
      (
        \morphName{f}^2_0 \morphCompose \morphName{f}^1_0,
        \morphName{f}^2_1 \morphCompose \morphName{f}^1_1
      ),
      \seqname{S}^1,
      \seqname{S}^3
    \bigr )
  \Bigr )
= \\
  \Bigl (
    \bigl (
      (
        \morphName{f}^2_0,
        \morphName{f}^2_1
      ),
      \seqname{S}^2,
      \seqname{S}^3
    \bigr )
  \Bigr )
  \morphCompose[A]x
  \Bigl (
    \bigl (
      (
        \morphName{f}^1_0,
        \morphName{f}^1_1
      ),
      \seqname{S}^1,
      \seqname{S}^2
    \bigr )
  \Bigr )
= \\
\Functor^\Ck_{\LCS,\Man}
  \Bigl (
    \bigl (
      (
        \morphName{f}^2_0,
        \morphName{f}^2_1
      ),
      \seqname{L}^2,
      \seqname{L}^3
    \bigr )
  \Bigr )
\morphCompose[A]
\Functor^\Ck_{\LCS,\Man}
  \Bigl (
    \bigl (
      (
        \morphName{f}^1_0,
        \morphName{f}^1_1
      ),
      \seqname{L}^1,
      \seqname{L}^2
    \bigr )
  \Bigr )
$

\item[Identity:]
$
\Functor^\Ck_{\LCS,\Man} (\Id[L^i])
=
\Functor^\Ck_{\LCS,\Man}
  \bigl
    (\Id[E^i], \Id[\seqname{C}^i]), \seqname{L}^i, \seqname{L}^i
  \bigr )
= \\
  \bigl
    (\Id[E^i], \Id[{\seqname{C}^i}]), \seqname{S}^i, \seqname{S}^i
  \bigr )
=
\Id[{\Functor^\Ck_{\LCS,\Man}}] (L^i)
$
\end{description}
\end{proof}
\item
$
  \Functor^\Ck_{\LCS,\Man}
  \morphCompose
  \Functor^{\Ck,\catSeqName{M}}_{\Man,\LCS}
$
is the identity functor on $\Atl[\Ck](\seqname{E}, \seqname{C})$.
\begin{proof}
$
\Functor^\Ck_{\LCS,\Man} \morphCompose \Functor^{\Ck,\catSeqName{M}}_{\Man,\LCS}
  (\seqname{S}^i) =
\Functor^\Ck_{\LCS,\Man} (\seqname{L}^i) =
\seqname{S}^i
$
\end{proof}
\item
$\Functor^{\Ck,\catSeqName{M}}_{\Man,\LCS} \morphCompose \Functor^\Ck_{\LCS,\Man}$
is the identity functor on $\LCS^\Ck(\seqname{E}, \seqname{C})$.
\begin{proof}
$
\Functor^{\Ck,\catSeqName{M}}_{\Man,\LCS} \morphCompose \Functor^\Ck_{\LCS,\Man}
  (\seqname{L}^i) =
\Functor^{\Ck,\catSeqName{M}}_{\Man,\LCS} (\seqname{S}^i) =
\seqname{L}^i
$
\end{proof}
\end{enumerate}
\end{theorem}

\part{Equivalence of fiber bundles}
\label{part:bun}
For fiber bundles\footnote{
  The literature has several definitions of fiber bundle. This paper
  uses one chosen for clarity of exposition. It differs from
  \cite[p.~8]{TopFib} in that, e.g., it uses the machinery of maximal
  atlases rather than equivalence classes of coordinate bundles, the
  nomenclature differs in several minor regards.
},
the adjunct spaces are the base space $\seqname{X} = (X,\catName{X})$,
the fiber $\seqname{Y} = (Y, \catName{Y})$
and the group $G$;
the category of the coordinate space is the category of Cartesian
products $\{U \times Y \mid U \objin \catName{X}\}$ of model
neighborhoods in the base space with the entire fiber, with morphisms
$\morphName{t} \maps U \times Y \to U \times Y$ that preserve the
fibers, i.e., $\pi_1 \morphCompose \morphName{t} = \pi_1$, and are
generated by the group action on the fiber
(Equation~\eqref{eq:transitiongroup}).

The sole adjunct functions are the projection $\pi \maps E \morphEpi X$, the
group operation and the group action on the fiber.

\begin{remark}
While this presentation makes the projection $\pi$ explicit, there is a
less traditional approach that derives the projection and is in some
ways simpler.
\end{remark}

This part of the paper defines bundle atlases, fiber bundles, local
coordinate spaces equivalent to fiber bundles, categories of them and
functors, and gives basic results.

\begin{definition}[Trivial group category of groups]
Let $\seqname{G}$ be a set of topological groups. The trivial group
category of $\seqname{G}$, abbreviated
$\Trivcat[\mathrm{group}-]{\seqname{G}}$, is the category of all
continuous homomorphisms between groups of $\seqname{G}$. By abuse of
language it will be shortened to $\Trivcat{\seqname{G}}$ when the
meaning is clear from context.
\end{definition}

\begin{definition}[Group actions]
Let $Y$ be a topological space, $G$ a topological group, $\rho$ an
effective group action of $G$ on $Y$, $y \in Y$ and $g \in G$. Then
$y \star g \defeq \rho(y,g)$.

Let $X$ be a topological space and $x \in X$. Then
$(x, y) \star g \defeq (x, y \star g)$.

A $\star$ with a subscript, superset, underset or overset refers to the
group action $\rho$ with the same subscript, superset, underset or
overset.

\begin{remark}
This notation is only used when it is clear from context what the group
action is.
\end{remark}
\end{definition}

\begin{definition}[Protobundles]
\label{def:ProtoBun}
Let $\seqname{B} \defeq (E, X, Y, \pi, G, \rho)$, where $E$, $X$ and $Y$
are topological spaces, $G$ a topological group, $\pi \maps E \morphEpi X$ a
continuous surjection and $\rho$ an efective group action of $G$ on $Y$.
Then $\seqname{B}$ is a protobundle.
\begin{remark}
While this definition does not itself require $E$ to have a local product
structure nor require $\pi$
to have the Covering Homotopy Property, only those protobundles having
an atlas are of interest, and for them \pagecref{def:Bun-Atl} imposes
additional constraints.
\end{remark}
\end{definition}

\begin{definition}[$X$-$\pi$-fine grained model spaces]
\label{def:Xpimod}
Let $\seqname{E} \defeq (E,\catName{E})$ be a model space, $X$ be a
topological space and the projection $\pi \maps E \to X$ be a continuous
function. $\seqname{E}$ is a $X$-$\pi$-fine grained model space iff
every $U \objin \seqname{E}$ is of the form $\pi^{-1}[V]$ for
$V \subseteq X$ open and for every open sets
$V' \subseteq V \subseteq X$, if $\pi^{-1}[V] \objin \seqname{E}$ then
$\pi^{-1}[V'] \objin \seqname{E}$.

\begin{remark}
The spaces used for fiber bundles have a local product structure with a
surjective projecttion, although neither is required by this definition.
\end{remark}
\end{definition}

\begin{definition}[$X$-$\pi$-trivial model spaces]
\label{def:Xpitriv}
Let $E$ and $X$ be topological spaces and the projection
$\pi \maps E \to X$ be a continuous function.

Let $\catName{E}$ be the category whose objects are all $\pi^{-1}[V]$
for $V \subseteq X$ open and whose morphisms are all continuous
$\morphName{f} \maps U^1 \to U^2$ that preserve fibers, i.e., \\*
$
  \uquant
    {{x,y \in U^1}}
    {{
      \pi(x) = \pi(y) \vdash
      \pi \bigl ( \morphName{f}(x) \bigr ) =
      \pi \bigl ( \morphName{f}(y) \bigr )
    }}
$

Then the the $X$-$\pi$-trivial model space of $E$ is
$\Triv[X-\pi-]{\seqname{E} \defeq (E,\catName{E})}$ and
$\Triv[X-\pi-]{\seqname{E}}$ is a $X$-$\pi$-trivial model space.
\end{definition}

\begin{lemma}[$X$-$\pi$-trivial model spaces]
\label{lem:Xpitriv}
Let $E$ and $X$ be topological spaces and the projection
$\pi \maps E \to X$ be a continuous function. Then
$\Triv[X-\pi-]{E}$ is a $X$-$\pi$-fine grained model space.

\begin{proof}
\begin{enumerate}
\item
Every $U \objin \seqname{E}$ is of the form $\pi^{-1}[V]$ for
$V \subseteq X$ open and for every open sets by \cref{def:Xpitriv}
above.

\item
For every open $V' \subset X$, $\pi^{-1}[V'] \objin \Triv[X-\pi-]{E}$.
\end{enumerate}
\end{proof}

\end{lemma}

\section{$G$-$rho$ model spaces}
\label{sec:grhomod}
\begin{definition}[$G$-$\rho$-model spaces]
{
  \showlabelsinline
  \label{def:Grhomod}
}
Let $\seqname{XY} \defeq (X \times Y, \catName{XY})$ be a model space,
$G$ a topological group and $\rho$ an effective group action of $G$ on
$Y$ such that the objects of $\catName{XY}$ are products of open sets
with $Y$ and the morphisms are fiber-preserving automorphisms generated
by the group action, i.e.,
\begin{equation}
\uquant
  {
    U \objin \catName{XY}
  }
  {
    \equant
      {
        V \in \op{X}
      }
      {
        U = V \times Y
      }
  }
\end{equation}
\begin{equation}
\label{eq:fGrho}
\uquant
  {
   \morphName{t} \arin \catName{XY} \maps V \times Y \toiso V \times Y
  }
  {
    \equant
    {
      \morphName{g} \in G^V
    }
    {
    \uquant
      {
        {(x,y) \in V \times Y}
      }
      {
        \morphName{t}(x,y) =
          \bigl ( x, y \star \morphName{g}(x) \bigr )
      }
    }
  }
\end{equation}

Then $\seqname{XY}$ is a $G$-$\rho$ model space of $X \times Y$,
abbreviated
$
\mathrm{isG{\rho}}
(
  \seqname{XY},
  Y,
  G,
  \rho
)
$,
\\*
$G_{\seqname{XY},\morphName{t}} \defeq \range(\morphName{g})$
is the set of group elements associated with $\morphName{t}$
and \\*
$
  G_\seqname{XY} \defeq
  \union%
    [{\morphName{t} \arin \catName{XY}}]
    {G_{\seqname{XY},\morphName{t}}}
$
is the set of group elements associated with $\seqname{XY}$ morphisms.

$\seqname{XY}$ is also a fine grained $G$-$\rho$ model space of
$X \times Y$ iff for every $U \times Y \objin \catName{XY}$ and every
open $U' \subseteq U$, the product $U' \times Y$ is an object of
$\catName{XY}$.
\end{definition}

\begin{lemma}[$G$-$\rho$-model spaces]
\label{lem:Grhomod}
Let $\seqname{E}^i \defeq (E^i, \catName{E}^i)$, $i=1,2$, be a model
space, $\seqname{XY}^i$ be a $G^i$-$\rho^i$ model space of
$X^i \times Y^i$ and
$
  \seqname{A}^i \defeq
  \set%
    {{(U^\alpha, V^\alpha \times Y^i, \phi^\alpha)}}%
    [{\alpha \prec \Alpha^i}]
$
be an m-atlas of $\seqname{E}^i$ in the coordinate space $\seqname{XY}^i$.
\begin{enumerate}
\item
Let $\morphName{f} \arin \catName{XY}^i \maps V \times Y \toiso V \times Y$.
The function $\morphName{g}$ in \cref{eq:fGrho} is unique.

\begin{proof}
The group action is effective.
\end{proof}

$G^i_{\seqname{XY}^i}$ is unique.

\begin{proof}
The function $\morphName{g}$ is unique.
\end{proof}

\item
Let $(U^{i,j}, V^{i,j} \times Y^i, \phi^{i.j}) \in \seqname{A}^i$,
$j=1,2$, with $I \defeq U^{i,1} \cap U^{i,2} \neq \emptyset$. Then
$\pi_1 \morphCompose \phi^{i,1} = \pi_1 \morphCompose \phi^{i,2}$ on $I$.

\begin{proof}
By \cref{def:Grhomod} above, the transition function is of the form \\*
$\morphName{t}(x,y) = \bigl ( x, y \star^i \morphName{g}(x) \bigr )$.
Then
$
  \pi_1 \morphCompose \morphName{t}(x,y) =
  \pi_1 \morphCompose \bigl ( x, y \star^i \morphName{g}(x) \bigr ) =
  x
$.
\end{proof}

\item
Let $\pi^{i,\alpha} \defeq \pi_1 \morphCompose \phi^{i,\alpha}$,
$\alpha \prec \Alpha^i$, and let $\pi^i \maps \seqname{E}^i \to X$ be
the function agreeing with each $\pi^{i,\alpha}$ on $U^{i,\alpha}$. Then
$\pi^i$ is continuous and
$\pi^{{i,\alpha}-1}[V^{i,\alpha} \times Y^i] = U^{i,\alpha}$.

\begin{proof}
$\pi^i$ is well defined by the above result. Let
$W \subseteq X^i \times Y^i$ be open. Each $W \cap V^{i.\alpha} \times Y^i$ is
open, so each ${\pi^{i.\alpha}}^{-1}[W \cap V^{i.\alpha} \times Y^i]$ is open,
Then \\*
$
  \pi^{i-1}[W \cap V^{i.\alpha} \times Y^i] =
  \union%
  [{{\alpha \prec \Alpha}}]%
  {{\pi^\alpha}^{-1}[W \cap V^\alpha \times Y^i]}
$
is open.
\end{proof}

\item
$\seqname{A}^i$ is an m-atlas of $\Triv[X^i-\pi^i-]{E^i}$ in the
coordinate space $\seqname{XY}^i$.

\begin{proof}
Let $(U^{i,j},V^{i,j} \times Y^i,\phi^{i,j}) \in \seqname{A}^i$,
$j=1,2$. Then $(U^{i,j} = \pi^{i-1}[V^{i,j}]$ by the above, and hence is
a neighborhood of $\Triv[X^i-\pi^i-]{E^i}$.

The compatibility constraints are the same in both cases.

If $V'^{i,j} \subseteq V^{i,j}$ is open then $\pi^{i,j-1}[V^{i,j}]$ is a
neighborhood of $\Triv[X^i-\pi^i-]{E^i}$ by \pagecref{def:Xpitriv}.

\end{proof}

\item
If $\seqname{XY}^i$ is a fine grained $G^i$-$\rho^i$ model space of
$X^i \times Y^i$ then $\seqname{E}^i$ is an $X^i$-$\pi_1$-fine grained model
space.

\begin{proof}
$\seqname{E}^i$ satisfies the conditions of \pagecref{def:Xpimod}:
\begin{enumerate}
\item
For $V \subseteq X^i$, $\pi^{i-1}_1[V] = V \times Y^i$.

\item
If $\pi^{i-1}_1[V] = V \times Y^i$ is a model neighborhood, $V'$ is open
and $V' \subseteq V$ then $\pi^{i-1}_1[V'] = V' \times Y^i$ is a model
neighborhood,
\end{enumerate}
\end{proof}

\item
Let $\morphSeqName{f} \defeq (\morphName{f}_0, \morphName{f}_1)$ be a
$\seqname{E}^1$-$\seqname{E}^2$ m-atlas morphism of $\seqname{A}^1$ to
$\seqname{A}^2$ in the coordinate spaces $\seqname{XY}^1$,
$\seqname{XY}^2$. Then $\morphSeqName{f}$ is a
$\Triv[X^1-\pi^1-]{E^1}$-$\Triv[X^2-\pi^2-]{E^2}$ m-atlas morphism of
$\seqname{A}^1$ to $\seqname{A}^2$ in the coordinate spaces
$\seqname{XY}^1$, $\seqname{XY}^2$.

\begin{proof}
$\morphSeqName{f}_1$ is a model function by hypothesis, hence continuous.

$\morphSeqName{f}_0 \maps \seqname{E}^1 \to \seqname{E}^2$ is a model
function by hypothesis. Let $U^2$ be a model neighborhood of
$\Triv[X^2-\pi^2-]{E^2}$ and $U^1 \defeq \morphName{f}^{-1}_0[U^2]$. Then
$U^2$ is of the form $\pi^{2-1}[V^2]$ with $V^2$ open. Let
$V^1 \defeq \morphName{f}^{-1}V^2]$. Since $\morphSeqName{f}_1$ is
continuous, $V^1$ is open and thus $\pi{1-1}[V^1] = U^1$ is open and a
model neighborhood of  $\Triv[X^1-\pi^1-]{E^1}$.

Let $(U^i,V^i \times Y^i, \phi^i) \in \seqname{A}^i$, $i=1,2$, and
$I \defeq U^1 \cap \morphName{f}^{-1}_0[U^2] \neq \emptyset$. Each $V^i$
is open, thus $U^i = \pi^{i-1}[V^i]$ is a model neighborhood of
$\Triv[X^i-\pi^i-]{E^i}$.

Let $u^1 \in I$, $(U'^1, V'^1 \times Y^i, \phi^1)$,
$U'^2 \subseteq U^2$, $(U'^2, \hat{V'^2} \times Y^2, \phi'^2)$ be as in
\pagecref{M-ATLmorph}. By the above, $(U'^1, V'^1 \times Y^i, \phi^1)$
and $U'^2 \subseteq U^2$, $(U'^2, \hat{V'^2} \times Y^2, \phi'^2)$ are
M-charts of $\Triv[X^1-\pi^1-]{E^1}$ ($\Triv[X^2-\pi^2-]{E^2}$)
in the coordinate space $\seqname{XY}^1$ ($\seqname{XY}^2$). The
commution relations are not changed by the change of base space.
\end{proof}
\end{enumerate}
\end{lemma}

\begin{definition}[Morphisms of $G$-$\rho$-model spaces]
\label{def:Grhomorph}
Let
$X^i.Y^i$, $i=1,2$, be topological spaces,
$G^i$ a topological group,
$\rho^i$ an effective group action on $Y^i$ and \\
$\seqname{XY}^i \defeq (X^i \times Y^i, \catName{XY}^i)$ a
$G^i$-$\rho^i$ model space of $X^i \times Y^i$.
Then a model function
$\morphName{f}_C \maps \seqname{XY}^1 \to \seqname{XY}^2$ is a
$G^1$-$G^2$-$\rho^1$-$\rho^2$ morphism of $X^1 \times Y^1$ to $X^2 \times Y^2$,
abbreviated
$
\mathrm{isG{\rho}morph}
(
  \seqname{XY}^1,
  G^1,
  \rho^1,
  \seqname{XY}^2,
  G^2,
  \rho^2,
  \morphName{f}_C
)
$
iff it preserves the group action, i.e., there is a continuous
homomorphism $\morphName{f}_G \maps G^1 \to G^2$ such that
\fullcref{fig:Grho} is commutative, i.e.,
{
  \showlabelsinline
  \cref{eq:Grho}
}
below holds.

\begin{figure}
\[ \bfig
\node xyg1(0,1000)[\bigl ( (x^1,y^1), g^1 \bigr )]
\node xys1(0,0)[(x^1, y^1 \star g^1)]
\node xyg2(1000,1000)[\bigl ( (x^2,y^2), g^2 \bigr )]
\node xys2(1000,0)[(x^2, y^2 \star g^2)]
\arrow |r|[xyg1`xyg2;\morphName{f}_C \times \morphName{f}_G]
\arrow |a|[xyg1`xys1;\star^1]
\arrow |r|[xys1`xys2;\morphName{f}_C]
\arrow |a|[xyg2`xys2;\star^2]
\efig \]
\caption{Preserving group action}
\label{fig:Grho}
\end{figure}
\begin{equation}
\showlabelsinline
\label{eq:Grho}
\equant
  {
    \morphName{f}_G \maps G_1 \to G_2
  }
  {
    \uquant
    {
      {
        (x,y) \in X^1 \times Y^1
      },
      {
        g \in G^1
      }
    }
    {
      \morphName{f}_C \bigl ( (x,y) \star^1 g \bigr )
      =
      \morphName{f}_C \bigl ( (x,y) \bigr ) \star^2 \morphName{f}_G (g)
    }
  }
\end{equation}
\end{definition}

\begin{lemma}[Morphisms of $G$-$\rho$-model spaces]
\label{lem:Grhomorph}
Let $X^i.Y^i$, $i=1,2,3$, be topological spaces, $G^i$ a topological
group, $\rho^i$ an effective group action on $Y^i$,
\nobreaks
{{
  $\seqname{XY}^i \defeq (X^i \times Y^i, \catName{XY}^i)$
}}
a $G^i$-$\rho^i$ model space of $X^i \times Y^i$ and
$\morphName{f}^i_C \maps X^i \times Y^i \to X^{i+1} \times Y^{i+1}$ a
$G^i$-$G^{i+1}$-$\rho^i$-$\rho^{i+1}$ morphism of $X^i \times Y^i$ to
$X^{i+1} \times Y^{i+1}$.

\begin{enumerate}
\item
\label{grhomorph:fgunique}
The function $\morphName{f}_G$ in \cref{eq:Grho} is unique.

\begin{proof}
The group action is effective.
\end{proof}

\item
\label{grhomorph:fiber}
$\morphName{f}^i_C$ preserves fibers, i.e.,
$
  \pi_1 \bigl ( \morphName{f}^i_C(x,y) \bigr ) =
  \pi_1 \bigl ( \morphName{f}^i_C(x,y') \bigr )
$
for $x$ in $X^i$ and $y,y'$ in $ Y^i$,
\begin{proof}
Since $\rho^i$ is effective, there is a $g \in G^i$ such that
$y' = y \star^i g$ and thus by \cref{def:Grhomorph}
\begin{equation}
\begin{split}
  \pi_1 \bigl ( \morphName{f}^i_C(x,y') \bigr )
  & =
\\
  \pi_1 \bigl ( \morphName{f}^i_C(x,y \star^i g) \bigr )
  & =
\\
  \pi_1 \bigl ( \morphName{f}^i_C(x,y) \star^{i+1} \morphName{f}^i_G(g) \bigr )
  & =
\\
  \pi_1 \bigl ( \morphName{f}^i_C(x,y) \bigr )
\end{split}
\end{equation}
\end{proof}

\item
\label{grhomorph:fxunique}
There exists a unique function $\morphName{f}^i_X \maps X^i \to X^{i+1}$
such that
$\morphName{f}^i_X \morphCompose \pi_1 = \pi_1 \morphCompose \morphName{f}^i_C$.

\begin{proof}
Define
$\morphName{f}^i_X(x) = \pi_1 \bigl ( \morphName{f}^i_C(x,y) \bigr )$,
where $y$ is an arbitrary point of $Y^i$. It does not depend on the
choice of $y$ because $\morphName{f}^i_C$ preserves fibers.
\end{proof}

$\morphName{f}^i_X$ is continuous.

\begin{proof}
Let $V^{i+1} \subseteq X$ be open. Then
$(\pi_1 \morphCompose \morphName{f}^i_C)^{-1}[V^{i+1}]$ is open and is of the
form $V^i \times Y$ with $V^i \subseteq X^i$ open.
$\morphName{f}^i_X[V^{i+1}] = V^i$ and is thus open.
\end{proof}

\begin{remark}
$\morphName{f}^i_C$ may be a twisted product:  there need not exist
$\morphName{f}^i_Y \maps Y^i \to Y^{i+1}$ such that
$\morphName{f}^i_C = \morphName{f}^i_X \times \morphName{f}^i_Y$.
\end{remark}

\item
\label{grhomorph:comp}
$\morphName{f}^2_C \morphCompose \morphName{f}^1_C$ is a
$G^1$-$G^3$-$\rho^1$-$\rho^3$ morphism of $X^1 \times Y^1$ to $X^3 \times Y^3$.

\begin{proof}

\begin{figure}
\[ \bfig
\node xyg1(0,1500)[\bigl ( (x^1,y^1), g^1 \bigr )]
\node xys1(0,0)[(x^1, y^1 \star g^1)]
\node xyg2(1000,1000)[\bigl ( (x^2,y^2), g^2 \bigr )]
\node xys2(1000,0)[(x^2, y^2 \star g^2)]
\node xyg3(2000,1500)[\bigl ( (x^3,y^3), g^3 \bigr )]
\node xys3(2000,0)[(x^3, y^3 \star g^3)]
\arrow |r|[xyg1`xyg3;\morphName{f}^2_C \morphCompose \morphName{f}^1_C \times%
  \morphName{f}^2_G \morphCompose \morphName{f}^1_G]
\arrow |r|[xyg1`xyg2;\morphName{f}^1_C \times \morphName{f}^1_G]
\arrow |a|[xyg1`xys1;\star^1]
\arrow |r|[xys1`xys2;\morphName{f}^1_C]
\arrow |a|[xyg2`xys2;\star^2]
\arrow |r|[xyg2`xyg3;\morphName{f}^2_C \times \morphName{f}^2_G]
\arrow |r|[xys2`xys3;\morphName{f}^1_C]
\arrow |a|[xyg3`xys3;\star^3]
\efig \]
\caption{Preserving group actions}
\label{fig:Grho2}
\end{figure}

Let $\morphName{f}^i_G \maps G^i \to G^{i+1}$ be a continuous
homomorphism such that
\begin{equation}
\label{eq:Grho2}
    \uquant
    {
      {
        (x^i,y^i) \in X^i \times Y^i
      },
      {
        g^i \in G^i
      }
    }
    {
      \morphName{f}^i_C \bigl ( (x^i,y^i) \star g^i \bigr )
      =
      \morphName{f}^i_C \bigl ( (x^i,y^i) \bigr ) \star \morphName{f}^i_G (g^i)
    }
\end{equation}
Let $(x^1,y^1) \in X^1 \times Y^1$ and $g^1 \in G^1$. Then \fullcref{fig:Grho2} is
commutative and
\begin{align*}
  \morphName{f}^2_C \morphCompose \morphName{f}^1_C
  \bigl ( (x^1,y^1) \star g^1 \bigr )
  & =
  \morphName{f}^2_C
  \Bigl (
    \morphName{f}^1_C \bigl ( (x^1,y^1) \bigr ) \star
    \morphName{f}^1_G (g^1)
  \Bigr )
\\*
  & =
  \morphName{f}^2_C
  \Bigl (
    \morphName{f}^1_C \bigl ( (x^1,y^1) \bigr )
  \Bigr )
  \star\morphName{f}^2_G \morphCompose \morphName{f}^1_G (g^1)
\\*
  & =
  \morphName{f}^2_C \morphCompose \morphName{f}^1_C \bigl ( (x^1,y^1) \bigr )
  \star\morphName{f}^2_G \morphCompose \morphName{f}^1_G (g^1)
\end{align*}
\end{proof}
\end{enumerate}
\end{lemma}

\begin{definition}[Categories of $G$-$\rho$-model spaces]
\label{def:Grhomodcat}
Let $X^\alpha.Y^\alpha$, $\alpha \prec \Alpha$, be topological spaces,
$G^\alpha$ a topological group, $\rho^\alpha$ an effective group action
on $Y^\alpha$, \\
$
  \seqname{XY}^\alpha \defeq
  (X^\alpha \times Y^\alpha, \catName{XY}^\alpha)
$
a $G^\alpha$-$\rho^\alpha$ model space of $X^\alpha \times Y^\alpha$, \\*
$
  \catSeqName{XYG\rho}_\Ob \defeq
  \set
    {\seqname{XY}^\alpha}%
    [\alpha \prec \Alpha]
$, \\*
$
  \catSeqName{XYG\rho}_\Ar \defeq
  \set
    {\morphName{f}_C \maps X^\alpha \times Y^\alpha \to X^\beta \times Y^\beta}%
    [
      {
        \mathrm{isG{\rho}morph}
        (
          \seqname{XY}^\alpha,
          G^\alpha,
          \rho^\alpha,
          \seqname{XY}^\beta,
          G^\beta,
          \rho^\beta,
          \morphName{f}_C
        )
      },
      \alpha \prec \Alpha,
      \beta \prec \Alpha
    ]*
$ \\*
and
$
  \catSeqName{XYG\rho} \defeq
  (
    \catSeqName{XYG\rho}_\Ob,
    \catSeqName{XYG\rho}_\Ar
  )
$.
Then any subcategory of $\catSeqName{XYG\rho}$ is a $G$-$\rho$ model category.
\end{definition}

\begin{definition}[Trivial $G$-$\rho$-model spaces]
\label{def:Grhomodtriv}
Let $X$ and $Y$ be topological spaces, $G$ a topological group, $\rho$
an effective group action of $G$ on $Y$ and $\catName{XY}$ the category
of all products of open subsets of $X$ with $Y$ and all homeomorphisms
induced by the group action, i.e,

\begin{equation}
\Ob(\catName{XY})
\defeq
\set
  {V \times Y}%
  [V \in \op{X}]
\end{equation}
\begin{multline}
\Ar(\catName{XY})
\defeq
\set
  {
   \morphName{t} \maps V \times Y \toiso V \times Y
  }%
  [
    {V \in \op{X}},
    {\pi_1 \morphCompose \morphName{t} = \pi_1},
    {
      \equant
        {
         \morphName{g} \in G^V,
        }
        {
          \uquant
            {
              {x \in V},
              {y \in Y}
            }
            {
               f(x, y) = \bigl ( x, y \star \morphName{g}(x) \bigr )
            }
        }
    }
  ]*
\end{multline}
Then the trivial $G$-$\rho$ model space of $X,Y$, abbreviated
$\Triv[G-\rho-]{X,Y}$, is $(X \times Y, \catName{XY})$ and
$\Triv[G-\rho-]{X,Y}$ is a trivial $G$-$\rho$ model space of $X,Y$.

\begin{remark}
Let $G'$ be a topological group and $\rho'$ an effective group action of
$G'$ on $Y$ such that $\Triv[G-\rho-]{X,Y} = \Triv[G'-\rho'-]{X,Y}$.
Although $G'$ must be isomorphic to $G$. it need not have the same
topology.
\end{remark}

The identity morphism of $\Triv[G-\rho-]{X,Y}$ is
$\Id_{\Triv[G-\rho-]{X,Y}} \defeq \Id[{X \times Y}]$.
\end{definition}

\begin{lemma}[Trivial $G$-$\rho$-model spaces]
\label{lem:Grhomodtriv}
Let $X$ and $Y$ be topological spaces, $G$ a topological group and
$\rho$ an effective group action on $Y$.

$\Triv[G-\rho-]{X,Y}$ is a $G$-$\rho$ model space of $X \times Y$.

\begin{proof}
$\Triv[G-\rho-]{X,Y}$ satisfies the conditions of
\pagecref{def:Grhomod}:
\begin{enumerate}
\item Product with fiber:
\[
\uquant
  {
    U \objin \pi_2 \biggl ( \Triv[G-\rho-]{X,Y} \biggr )
  }
  {
    \equant
      {
        V \in \op{X}
      }
      {
        U = V \times Y
      }
  }
\]
by \cref{def:Grhomodtriv}.
\item Generated by $\rho$:
Let
$
  \morphName{f} \arin
  \pi_2 \biggl ( \Triv[G-\rho-]{X,Y} \biggr )
  \maps V \times Y \toiso V \times Y
$.
Then
\[
  \equant
    {
      \morphName{g} \in G^V
    }
    {
    \uquant
      {
        {(x,y) \in V \times Y}
      }
      {
        \morphName{f}(x,y) =
          \bigl ( x, y \star \morphName{g}(x) \bigr )
      }
    }
\]
by \cref{def:Grhomodtriv}.
\end{enumerate}
\end{proof}

$\Triv[G-\rho-]{X,Y}$ is a fine grained $G$-$\rho$ model space of $X
\times Y$.

\begin{proof}
Let $V \times Y$ be a model neighborhood and $V'$ an open subset.
$V' \times Y$ is open by \cref{def:Grhomodtriv}.
\end{proof}

Let $\seqname{E}^i \defeq (E^i, \catName{E}^i)$, $i=1,2$.  be a model
space, $\seqname{XY}^i$ be a fine grained $G^i$-$\rho^i$ model space of
$X^i \times Y^i$, and
$
  \seqname{A}^i \defeq
  \set%
    {{(U^\alpha, V^\alpha \times Y^i, \phi^\alpha)}}%
    [{\alpha \prec \Alpha}]
$
be an m-atlas of $\seqname{E}$ in the coordinate space $\seqname{XY}^i$.

$\seqname{A}^i$ is an m-atlas of $\Triv[X^i-\pi^i-]{E^i}$ in the coordinate
space $\Triv[G^i-\rho^i-]{X^i,Y^i}$.

\begin{proof}
\mbox{}
$\seqname{XY}^i$ is a fine grained $G^i$-$\rho^i$ model space of $X^i \times Y^i$
by hypothesis.

$\seqname{A}^i$ is an m-atlas of $\seqname{E}$ in the coordinate space
$\seqname{XY}^i$ by hypothesis.

$\seqname{A}^i$ is an m-atlas of $\Triv[X^i-\pi^i-]{E^i}$ in the coordinate
space $\seqname{XY}^i$ by \pagecref{lem:Grhomod}.

Let $(U^i, V^i \times Y^i, \phi^i) \seqname{A}^i$, $i=1,2$. Then
$(U^i, V^i \times Y^i, \phi^i)$ is an m-chart of $\Triv[X^i-\pi^i-]{E^i}$ in the
coordinate space $\seqname{XY}^i$:

\begin{enumerate}
\item
$V^i$ is open by \pagecref{def:Grhomod} and
$\pi^{i-1}[V^i \times Y^i] = U^i$ by \cref{lem:Grhomod}, hence $U^i$ is a
model neighborhood of $\Triv[X^i-\pi^i-]{E^i}$ by \pagecref{def:Grhomodtriv}.
\item
$[V^i \times Y^i]$ is a model neighborhood of $\Triv[X^i-\pi^i-]{E^i}$ by the
above.
\item
The transition function $\morphName{t}^i_j = \phi^j \morphCompose \phi^{i-1}$
is of the form
$
  \morphName{t}^i_j(x^i,y^i) =
  \bigl ( X^i, y^i \star^i \morphName{g}^i_j(x^i) \bigr )
$
\end{enumerate}
\end{proof}

Let $\morphSeqName{f} \defeq (\morphName{f}_0, \morphName{f}_1)$ be a
$\seqname{E}^1$-$\seqname{E}^2$ m-atlas morphism of $\seqname{A}^1$ to
$\seqname{A}^2$ in the coordinate spaces $\seqname{XY}^1$,
$\seqname{XY}^2$. Then $\morphSeqName{f}$ is a
$\Triv[X^1-\pi^1-]{E^1}$-$\Triv[X^2-\pi^2-]{E^2}$ m-atlas morphism of
$\seqname{A}^1$ to $\seqname{A}^2$ in the coordinate spaces spaces
$\Triv[G^1-\rho^1-]{X^1,Y^1}$, $\Triv[G^2-\rho^2-]{X^2,Y^2}$.

\begin{proof}
$\morphSeqName{f}_1$ is a model function by hypothesis, hence continuous.

$\morphSeqName{f}_0 \maps \seqname{E}^1 \to \seqname{E}^2$ is a model
function by hypothesis. Let $U^2$ be a model neighborhood of
$\Triv[X^2-\pi^2-]{E^2}$ and $U^1 \defeq \morphName{f}^{-1}_0[U^2]$. Then
$U^2$ is of the form $\pi^{2-1}[V^2]$ with $V^2$ open. Let
$V^1 \defeq \morphName{f}^{-1}V^2]$. Since $\morphSeqName{f}_1$ is
continuous, $V^1$ is open and thus $\pi{1-1}[V^1] = U^1$ is open and a
model neighborhood of  $\Triv[X^1-\pi^1-]{E^1}$.
Let $(U^i,V^i \times Y^i, \phi^i) \in \seqname{A}^i$, $i=1,2$, and
$I \defeq U^1 \cap \morphName{f}^{-1}_0[U^2] \neq \emptyset$. Each $V^i$ is
open, thus $U^i = \pi^{i-1}[V^i]$ is a model neighborhood of
$\Triv[X^i-\pi^i-]{E^i}$ and each $V^i \times Y^i$ is a model
neighborhood of $\Triv[G^i-\rho^i-]{X^i,Y^i}$.

Let $u^1 \in I$, $(U'^1, V'^1 \times Y^i, \phi^1)$,
$U'^2 \subseteq U^2$, $(U'^2, \hat{V'^2} \times Y^2, \phi'^2)$ be as in
\pagecref{M-ATLmorph}. By the above, $(U'^1, V'^1 \times Y^i, \phi^1)$
and $U'^2 \subseteq U^2$, $(U'^2, \hat{V'^2} \times Y^2, \phi'^2)$ are
M-charts of $\Triv[X^1-\pi^1-]{E^1}$ ($\Triv[X^2-\pi^2-]{E^2}$)
in the coordinate space
$\Triv[G^1-\rho^1-]{X^1,Y^1}$ ($\Triv[G^2-\rho^2-]{X^2,Y^2}$). The
commution relations are not changed by the change of spaces.
\end{proof}
\end{lemma}

\begin{corollary}
\label{cor:Grhomodtriv}
Let $X$ and $Y$ be topological spaces, $G$ a topological group and
$\rho$ an effective group action on $Y$. Then $\Triv[G-\rho-]{X,Y}$ is a
$X$-$\pi_1$-fine grained model space.

\begin{proof}
The proof follows from \pagecref{lem:Grhomod} and
{
   \showlabelsinline
  \cref{lem:Grhomodtriv}
}
above.
\end{proof}
\end{corollary}

\begin{lemma}[Morphisms of trivial $G$-$\rho$-model spaces]
\label{lem:Grhomodtrivmorph}
Let $X^i$ and $Y^i$, $i=1,2$, be topological spaces, $G^i$ a topological
group, $\rho^i$ an effective group action of $G^i$ on $Y^i$ and
\nobreaks
{{
  $(f_X \maps X^1 \to X^2, f_Y \maps Y^1 \to Y^2)$
}}
a pair of continuous functions such that
$
  \morphName{f}_C \defeq f_X \times f_Y \maps
  \Triv[G^1-\rho^1-]{X^1,Y^1} \to \Triv[G^2-\rho^2-]{X^2,Y^2}
$
is a $G^1$-$G^2$-$\rho^1$-$\rho^2$ morphism of $X^1 \times Y^1$ to
$X^2 \times Y^2$,
Then $\morphName{f}_C$ is a model function from
$\Triv[G^1-\rho^1-]{X^1,Y^1}$ to $\Triv[G^2-\rho^2-]{X^2,Y^2}$.
\begin{proof}
\leavevmode
\begin{enumerate}
\item
$\morphName{f}_C$ is a $G^1$-$G^2$-$\rho^1$-$\rho^2$ morphism of
$X^1 \times Y^1$ to $X^2 \times Y^2$ by hypothesis, hence continuous.

\item
Let $m^2 \defeq U^2 \times Y^2$ be a model neighborhood of
$\Triv[G^2-\rho^2-]{X^2,Y^2}$. Then
$U^1 \defeq \morphName{f}_X^{-1}(U^2)$ is open and
$\morphName{f}_C^{-1}(m^2) = U^1 \times Y^1$ is a model function.

\item
Let $m^1 \defeq U^1 \times Y^1$ be a model neighborhood of
$\Triv[G^1-\rho^1-]{X^1,Y^1}$. $X^2 \times Y^2$ is a model neighborhood
of $\Triv[G^2-\rho^2-]{X^2,Y^2}$ and
$\morphName{f}_C(m^1) \subseteq X^2 \times Y^2$.
\end{enumerate}
\end{proof}

\begin{remark}
The hypothesis precludes the possibilty that $\morphName{f}_C$ is
twisted.
\end{remark}
\end{lemma}

\begin{definition}[Categories of trivial $G$-$\rho$-model spaces]
\label{def:Grhomodtrivcat}
Let \\*
$
  \seqname{B}^\alpha \defeq
  (E^\alpha, X^\alpha, Y^\alpha, \pi^\alpha, G^\alpha, \rho^\alpha)
$,
$\alpha \prec \Alpha$, be a protobundle and
$\seqname{B} \defeq \set{{\seqname{B}^\alpha}}[\alpha \prec \Alpha]$
be a set of protobundles.

\begin{remark}
The component $E^\alpha$ plays no role in this definition.
\end{remark}

The trivial coordinate model category of $\seqname{B}$,
$\Trivcat[\Bun-]{\seqname{B}}$, is the category with objects all trivial
$G_\alpha$-$\rho_\alpha$ model spaces of $X_\alpha, Y_\alpha$,
$\alpha \prec \Alpha$ and morphisms all continuous functions compatible
with the group action:

\begin{equation}
\Trivcat[\Bun-]{\seqname{B}}_\Ob \defeq
\set
  {
    \Triv[G-\rho-]{X,Y}
  }%
  [
    {
      (
        E,
        X,
        Y,
        \pi,
        G,
        \rho
      )
      \in \seqname{B}
    }
  ]
\end{equation}
\begin{multline}
\Trivcat[\Bun-]{\seqname{B}}_\Ar \defeq
\set
  {
    \morphName{f}_C \maps
    X^1 \times Y^1 \to
    X^2 \times Y^2
  }%
  [
    {
      \equant
        {
          {
            (
              E^i,
              X^i,
              Y^i,
              \pi^i,
              G^i,
              \rho^i
            )
            \in \seqname{B}
          },
          \morphName{f}_G \maps G^1 \to G^2
        }
        {
          \mathrm{isG{\rho}morph}
          (
            \Triv[G-\rho-]{X^1,Y^1},
            G^1,
            \rho^1,
            \Triv[G-\rho-]{X^2,Y^2},
            G^2,
            \rho^2,
            \morphName{f}^C
          )
        }
    }
  ]*
\end{multline}
\begin{equation}
\Trivcat[\Bun-]{\seqname{B}} \defeq
\bigl (
  \Trivcat[\Bun-]{\seqname{B}}_\Ob,
  \Trivcat[\Bun-]{\seqname{B}}_\Ar
\bigr )
\end{equation}

\begin{remark}
The morphisms $\morphName{f}_C$ may be twisted products:  there need not
exist $\morphName{f}_Y \maps Y^1 \to Y^2$ such that
$\morphName{f}_C = \morphName{f}_X \times \morphName{f}_Y$.
\end{remark}

The trivial product coordinate model category of $\seqname{B}$,
$\Trivcat[\BunProd-]{\seqname{B}}$, is the category with objects all
trivial $G^\alpha$-$\rho^\alpha$ model spaces of $X^\alpha, Y^\alpha$,
$\alpha \prec \Alpha$ and morphisms all products of continuous functions
compatible with the group action:

\begin{equation}
\Trivcat[\BunProd-]{\seqname{B}}_\Ob \defeq \Trivcat[\Bun-]{\seqname{B}}_\Ob
\end{equation}
\begin{multline}
\Trivcat[\BunProd-]{\seqname{B}}_\Ar \defeq
\set
  {
    \morphName{f}_C \defeq  \morphName{f}_X \times \morphName{f}_Y \maps
    X^1 \times Y^1 \to
    X^2 \times Y^2
  }%
  [
    {
      \equant
        {
          {
            (
              E^i,
              X^i,
              Y^i,
              \pi^i,
              G^i,
              \rho^i
            )
            \in \seqname{B}
          },
          \morphName{f}_G \maps G^1 \to G^2
        }
        {
          \mathrm{isG{\rho}morph}
          (
            \Triv[G-\rho-]{X^1,Y^1},
            G^1,
            \rho^1,
            \Triv[G-\rho-]{X^2,Y^2},
            G^2,
            \rho^2,
            \morphName{f}_C
          )
        }
    }
  ]*
\end{multline}
\begin{equation}
\Trivcat[\BunProd-]{\seqname{B}} \defeq
\Bigl (
  \Trivcat[\BunProd-]{\seqname{B}}_\Ob,
  \Trivcat[\BunProd-]{\seqname{B}}_\Ar
\Bigr )
\end{equation}

Any subcategory of $\Trivcat[\Bun-]{\seqname{B}}$ is a trivial coordinate
model category and any subcategory of $\Trivcat[\BunProd-]{\seqname{B}}$ is
a trivial product coordinate model category.
\end{definition}

\begin{lemma}[The trivial coordinate model category of $\seqname{B}$ is a category]
\label{lem:Grhomodtrivcat}
Let \\
$
  \seqname{B}^\alpha \defeq
  (E^\alpha, X^\alpha, Y^\alpha, \pi^\alpha, G^\alpha, \rho^\alpha)
$,
$\alpha \prec \Alpha$, be a protobundle and
$\seqname{B} \defeq \set{{\seqname{B}^\alpha}}[\alpha \prec \Alpha]$
be a set of protobundles.

Then $\Trivcat[\Bun-]{\seqname{B}}$ is a category and the identity
morphism for object $\seqname{B}^\alpha$ is
$\Id[{X^\alpha] \times Y^\alpha}]$.

\begin{proof}
\leavevmode
\begin{description}
\item[Composition:]
The composition of $G$-$\rho$ morphisms is a $G$-$\rho$ morphism by
\pagecref{lem:Grhomorph}.
\item[Associativity:]
Morphisms are simply functions and composition of morphisms is simply
composition of functions.
\item[Unit:]
The identity morphisms are simply identity functions and composition of
morphisms is simply composition of functions.
\end{description}
\end{proof}

$\Trivcat[\BunProd-]{\seqname{B}}$ is a category and the identity
morphism for object $\seqname{B}^\alpha$ is
$\Id[{X^\alpha \times Y^\alpha}]$.

\begin{proof}
\leavevmode
\begin{description}
\item[Composition:]
The composition of $G$-$\rho$ morphisms is a $G$-$\rho$ morphism by
\pagecref{lem:Grhomorph}. Let
$\Triv[G^i-\rho^i-]{X^i,Y^i} \objin\Trivcat[\BunProd-]{\seqname{B}},$
$i=1,2,3$,
$
  \morphName{f}^i_C =
  \morphName{f}^i_X \times \morphName{f}^i_Y \maps
  X^i \times Y^i \to X^{i+1} \times Y^{i+1}
  \arin  \Trivcat[\BunProd-]{\seqname{B}}
$.
Then
$
  \morphName{f}^{i+1}_C \morphCompose \morphName{f}^i_C =
  (\morphName{f}^{i+1}_X \morphCompose \morphName{f}^i_X)
  \times
  (\morphName{f}^{i+1}_Y \morphCompose \morphName{f}^i_Y)
$.

\item[Associativity:]
Morphisms are simply functions and composition of morphisms is simply
composition of functions.

\item[Unit:]
The identity morphisms are simply identity functions and composition of
morphisms is simply composition of functions.
\end{description}
\end{proof}
\end{lemma}

\section{Bundle charts}
{
  \showlabelsinline
  \label{sec:buncharts}
}
\begin{definition}[$Y$-$\pi$-bundle charts]
\label{def:Ypichart}
Let $E,X,Y$ be topological spaces and
\nobreaks
{{
  $\pi \maps E \to X$.
}}
Then $(U, V \times Y, \phi \maps U \toiso V \times Y)$ is a
$Y$-$\pi$-bundle chart of $E$ in the coordinate space $X \times Y$ iff

\begin{figure}[ht]
\[ \bfig
\node u(0,1000)[U]
\node vy(0,0)[V \times Y]
\node x(1000,0)[X]
\arrow |l|[u`vy;\phi]
\arrow |r|[u`vy;\iso]
\arrow |a|[u`x;\pi]
\arrow |r|[vy`x;\pi_1]
\efig \]
\caption{$\phi$ preserves fibers}
\label{fig:Ypichart}
\end{figure}

\begin{enumerate}
\item $U$, known as a coordinate patch, is a novoid open subset of $E$,
\item $V \times Y$ is a novoid open subset of $X \times Y$
\item $\phi$, known as a coordinate function, is a homeomorphism that
preserves fibers. i.e., $\pi_1 \morphCompose \phi = \pi$.
\end{enumerate}

$(U, V \times Y, \phi)$ covers $u$ iff $u \in U$.
\end{definition}

\begin{lemma}[$Y$-$\pi$-bundle charts of m-atlases in $G$-$\rho$ model spaces]
\label{lem:Ypichart}
Let
\begin{enumerate}
\item
$E,X,Y$ be topological spaces
\item
$G$ a topological group
\item
$\rho$ an effective group action of $G$ on $Y$
\item
$\seqname{XY} \defeq (X \times Y, \catName{XY})$
a $G$-$\rho$ model space of $X \times Y$
\item
$\seqname{A}$ an m-atlas of $E$ in the coordinate space $\seqname{XY}$
\end{enumerate}

There is a unique $\pi \maps E \morphEpi X$ such that every chart
$(U, V \times Y, \phi) \in \seqname{A}$ is a
$Y$-$\pi$-bundle chart of $E$ in the coordinate space
$X \times Y$. Further, $\pi$ is continuous.

\begin{proof}
Let $e$ be an arbitrary point of $E$ and let
$(U, V \times Y, \phi) \in \seqname{A}$ be an arbitrary m-chart in
$\seqname{A}$ that covers $e$.  Define $\pi(e) \defeq \pi_1 \morphCompose
\phi(e)$.

To show that $\pi(e)$ is well defined let $(U', V' \times Y, \phi')$ be
another m-chart in $\seqname{A}$ that covers $e$.  By
\pagecref{def:Grhomodtriv}, the transition function is of the form
$\morphName{t}(x, y) = \bigl ( x, y \star \morphName{g}(x) \bigr )$.  Then
$\pi_1 \morphCompose \morphName{t} = \pi_1$ and \\*
$\pi_1 \morphCompose \phi' (e) = \pi_1 \morphCompose \phi (e)$.

Let $V \subseteq X$ be open and let $e \in {\pi_1}^{-1}[V]$. Let
$(U', V' \times Y, \phi')$ be an m-chart of $\seqname{A}$ that covers $e$.
$V \cap V'$ is open, $\pi_1$ is continuous and $\phi'$ is continuous, so \\*
$\pi^{-1}[V'] = \phi'^{-1} \bigl [ \pi_1^{-1}[V'] \bigr ]$ is open.
$\pi^{-1}(V)$ is thus a union of open sets, hence open.
\end{proof}
\end{lemma}

\begin{lemma}[$Y$-$\pi$-bundle charts in \ensuremath{\Triv[G-\rho-]{X,Y}}]
\label{lem:GrhomodtrivYpichart}
Let $E$, $X$ and $Y$ be topological spaces, $G$ a topological group,
$\rho$ an effective group action of $G$ on $Y$ and
$\pi \maps E \to X$.
$(U, V \times Y, \phi)$ is a $Y$-$\pi$-bundle chart of $E$ in
the coordinate space $X \times Y$ iff $(U, V \times Y, \phi)$ is an
m-chart of $E$ in the coordinate space $\Triv[G-\rho-]{X,Y}$ and $\phi$
preserves fibers.

\begin{proof}
If $(U, V \times Y, \phi)$ is a $Y$-$\pi$-bundle chart of $E$
in the coordinate space $X \times Y$ then by \pagecref{def:Ypichart}

\begin{enumerate}
\item
$U$ is a nonvoid open subset of $E$.
\item
$V \times Y$ is an open subset of $X \times Y$, hence a model
neighborhood of $\Triv[G-\rho-]{X,Y}$ by \pagecref{def:Grhomodtriv}.

\item
$\phi$ is a homeomorphism

\item
$(U, V \times Y, \phi)$ is an m-chart of $E$ in the coordinate space
$\Triv[G-\rho-]{X,Y}$ by \cref{def:Grhomodtriv}.

\item
$\phi$ is a homeomorphism that preserves fibers. i.e.,
$\pi_1 \morphCompose \phi = \pi$.
\end{enumerate}

If $(U, V \times Y, \phi)$ is an m-chart of $E$ in the coordinate space
$\Triv[G-\rho-]{X,Y}$ and $\phi$ preserves fibers then by
\cref{def:Grhomodtriv}

\begin{enumerate}
\item
$U$ is a nonvoid open subset of $E$.

\item
$V \times Y$ is an m-chart of $E$ in the coordinate space
$\Triv[G-\rho-]{X,Y}$, hence an open subset of $X \times Y$.

\item
$\phi$ is a homeomorphism

\item
$\phi$ preserves fibers by hypothesis.
\end{enumerate}
\end{proof}
\end{lemma}

\begin{lemma}[Properties of projection]
\label{def:piepi}
Let $(U, V \times Y, \phi)$ be a $Y$-$\pi$-bundle chart of $E$ in the
coordinate space $X \times Y$ and $v \in V$. Then

$\pi \restrictto_U$ is a surjection.
\begin{proof}
Let $v \in V$, $y \in Y$ and
$u \defeq \phi^{-1} \bigl ( (v,y) \bigr ) \in U$. Then $\pi(u) = v$.
\end{proof}
$\pi^{-1}[\set{v}]$ is homeomorphic to $Y$.
\begin{proof}
$\phi$ and $\phi^{-1}$ are homeomorphisms, so their restrictions are
homeomorphisms and thus $\phi^{-1}[\set{v}]$ is homeomorphic to
$\set{v} \times Y$, which is homeomorphic to $Y$.
\end{proof}
\end{lemma}

\begin{definition}[$Y$-$\pi$ subcharts]
\label{def:Ypisub}
Let $(U, V \times Y, \phi)$ be a $Y$-$\pi$-bundle chart of $E$ in the
coordinate space $X \times Y$ and $U'$ be a nonvoid open subset of $U$.
Then \\*
$(U', V' \times Y, \phi') \defeq (U', \phi[U'], \phi \restrictto_{U'})$
is a subchart of $(U, V \times Y, \phi)$.
\end{definition}

\begin{lemma}[$Y$-$\pi$ subcharts]
Let $(U, V \times Y, \phi)$ be a $Y$-$\pi$-bundle chart of $E$ in the
coordinate space $X \times Y$ and $(U', V' \times Y, \phi')$ a subchart
of $(U, V \times Y, \phi)$. Then $(U', V' \times Y, \phi')$ is a
$Y$-$\pi$-bundle chart of $E$ in the coordinate space $X \times Y$.
\begin{proof}
$(U', V' \times Y, \phi')$ satisfies the conditions of \cref{def:Ypichart}:
\begin{enumerate}
\item $U' \subseteq E$ is open.
\item $\phi[U']$ os open since $\phi$ is a homeomorphism.
\item $\phi \restrictto_{U'} \maps U' \to V' \times Y$ is the
restriction of a homeomorphism and thus a homeomorphism.
$\phi \restrictto_{U'}$ preserves fibers because $\phi$ does.
\end{enumerate}
\end{proof}
\end{lemma}

\begin{definition}[$G$-$\rho$-compatibility]
\label{def:Grho-compat}
Let $E,X,Y$ be topological spaces, $G$ a topological group,
$\rho \maps Y \times G \to Y$ an effective right action of
$G$ on $Y$, $\pi \maps E \morphEpi X$ surjective
and $(U^\mu, V^\mu \times Y, \phi^\mu)$,
$(U^\nu, V^\nu \times Y, \phi^\nu)$ $Y$-$\pi$-bundle charts. Then
$(U^\mu, V^\mu \times Y, \phi^\mu)$ and
$(U^\nu, V^\nu \times Y, \phi^\nu)$ are $G$-$\rho$-compatible iff either
\begin{enumerate}
\item $U^\mu$ and $U^\nu$ are disjoint
\item
The transition function
$
  \morphName{t}^\mu_\nu \defeq
    \phi^\mu \morphCompose
    {\phi^\nu}^{-1} \restrictto_{\phi^\nu[U^\mu \cap U^\nu]}
$
is generated by the group action, i.e., there is a continuous function
$
  \morphName{g}^\mu_\nu
  \maps \pi_1 \bigl [ \phi^\nu[U^\mu \cap U^\nu] \bigr ] \to G
$
such that
\begin{equation}
\showlabelsinline
\label{eq:transitiongroup}
\uquant
  {{(x,y) \in \phi^\nu[U^\mu \cap U^\nu]}}
  {
    \morphName{t}^\mu_\nu(x,y) =
    \bigl ( x, y \star \morphName{g}^\mu_\nu(x) \bigr )
  }
\end{equation}
\end{enumerate}
\end{definition}

\begin{lemma}[$G$-$\rho$-compatibility]
\label{lem:Grho-compat}
Let $E,X,Y$ be topological spaces, $G$ a topological group,
$\rho \maps Y \times G \to Y$ an effective right action of
$G$ on $Y$, $\pi \maps E \morphEpi X$ surjective
and $(U^i, V^i \times Y, \phi^i)$, $i=1,2$, $Y$-$\pi$-bundle charts.
Then $(U^1, V^1 \times Y, \phi^1)$ and \\*
$(U^2, V^2 \times Y, \phi^2)$ are $G$-$\rho$-compatible iff
$(U^1, V^1 \times Y, \phi^1)$ and
$(U^2, V^2 \times Y, \phi^2)$ are compatible m-charts of
$\Triv{E}$ in the coordinate space
$\Triv[G-\rho-]{X,Y}$.

\begin{proof}
\begin{enumerate}
\leavevmode
\item
Each $U^i$ is a model neighborhood of $\Triv{E}$ and each $V^i \times Y$
is a model neighborhood of $\Triv[G-\rho-]{X,Y}$.

\item
If $U^1 \cap U^2 = \emptyset$ then the lemma is trivially true.

\item
If the charts are $G$-$\rho$-compatible then the transition function \\*
$
  \morphName{t}^1_2 \defeq
    \phi^1 \morphCompose
    {\phi^2}^{-1} \restrictto_{\phi^2[U^1 \cap U^2]}
$
is generated by the group action, i.e., there is a continuous function
$\morphName{g}^1_2 \maps \pi_1[\phi^2[U^1 \cap U^2]] \to G$
such that \\*
$
  \uquant
    {{(x,y) \in \phi^2[U^1 \cap U^2]}}
    {
      \morphName{t}^1_2(x,y) =
      \bigl ( x, y \star \morphName{g}^1_2(x) \bigr )
    }
$.
Then $\pi_1 \morphCompose \morphName{t}^1_2 = \pi_1$, $\morphName{t}^1_2$ is a
morphism of $\Triv[G-\rho-]{X,Y}$ and the charts are compatible m-charts
of $\Triv[G-\rho-]{X,Y}$.

\item
If the charts are compatible m-charts of $\Triv[G-\rho-]{X,Y}$ then the
transition function is an isomorphism of $\Triv[G-\rho-]{X,Y}$, hence
generated by the group action, and the charts are $G$-$\rho$-compatible.
\end{enumerate}
\end{proof}
\end{lemma}

\begin{corollary}[Symmetry of $G$-$\rho$ compatibility]
\label{cor:Grho-compatsym}
Let $(U^\mu, V^\mu \times Y, \phi^\mu)$ and \linebreak
$(U^\nu, V^\nu \times Y, \phi^\nu)$ be $Y$-$\pi$-bundle charts of $E$ in the
coordinate space $X \times Y$. Then
$(U^\mu, V^\mu \times Y, \phi^\mu)$ is $G$-$\rho$-compatible with
$(U^\nu, V^\nu \times Y, \phi^\nu)$ iff
$(U^\nu, V^\nu \times Y, \phi^\nu)$ is $G$-$\rho$-compatible with
$(U^\mu, V^\mu \times Y, \phi^\mu)$.
\begin{proof}
The result follows from \pagecref{lem:M-compatsym}.
\end{proof}
\end{corollary}

\begin{lemma}[$G$-$\rho$-compatibility of subcharts]
\label{lem:Grho-compatsub}
Let $(U^i, V^i \times Y, \phi^i)$ be a $Y$-$\pi$-bundle chart of $E$ in
the coordinate space $X \times Y$, $(U'^i, V'^i \times Y, \phi'^i)$ a
subchart and $(U^1, V^1, \phi^1)$ be $G$-$\rho$-compatible with
$(U^2, V^2, \phi^2)$.  Then $(U'^1, V'^1 \times Y, \phi'^1)$ is
$G$-$\rho$-compatible with $(U'^2, V'^2 \times Y, \phi'^2)$.

\begin{proof}
If $U^1 \cap U^2 = \emptyset$ then $U'^1 \cap U'^2 = \emptyset$. If
$U'^1 \cap U'^2 = \emptyset$ then
\nobreaks
{{
  $(U'^1, V'^1 \times Y, \phi'^1)$
}}
is $G$-$\rho$-compatible with $(U'^2, V'^2 \times Y, \phi'^2)$. Otherwise,
the transition function \\*
\nobreaks
{{
  $
    t^1_2 \defeq
    \phi^1 \morphCompose \phi^{2-1} \restrictto_{\phi^2[U^1 \cap U^2]}
  $
}}
is generated by the group action and hence \\*
$
  t^1_2 \restrictto_{\phi^2[U'^1 \cap U'^2]}
  \maps
  \phi^2[U'^1 \cap U'^2]
  \toiso
  \phi^1[U'^1 \cap U'^2]
$
is generated by the group action.
\end{proof}
\end{lemma}

\begin{corollary}[$G$-$\rho$-compatibility with subcharts]
Let $(U, V \times Y, \phi)$, be a $Y$-$\pi$-bundle chart of $E$ in the
coordinate space $X \times Y$ and $(U', V' \times Y, \phi')$ a subchart.
Then $(U', V' \times Y, \phi')$ is $G$-$\rho$-compatible with
$(U, V \times Y, \phi)$.

\begin{proof}
$(U, V \times Y, \phi)$ is $G$-$\rho$-compatible with itself and is a
subchart of itself,
\end{proof}
\end{corollary}

\begin{definition}[Covering by $Y$-$\pi$-bundle charts]
Let $\seqname{A}$ be a set of $Y$-$\pi$-bundle
charts of $E$ in the coordinate space $X \times Y$.
$\seqname{A}$ covers $E$ iff $E = \union{{\pi_1[\seqname{A}]}}$.
\end{definition}

\begin{lemma}
Let $\seqname{A}$ be a set of $Y$-$\pi$-bundle charts of $E$ in the
coordinate space $X \times Y$ that covers $E$ and $x \in X$. Then
$\pi^{-1}[\set{x}]$ is homeomorphic to $Y$.
\begin{proof}
Since $\seqname{A}$ covers $E$, there is a chart $(U, V, \phi)$ in
$\seqname{A}$ containing $x$.
Then $\pi^{-1}[\set{x}]$ is homeomorphic to $Y$ by
\pagecref{def:piepi}.
\end{proof}
\end{lemma}

\section{Bundle atlases}
\label{sec:BunAtlases}
A set of charts can be atlases for different fiber bundles even if it is
for the same total model space, base space and fiber. In order to
aggregate atlases into categories, there must be a way to distinguish
them. Including the spaces\footnote{
The spaces are redundant, but convenient.},
group and group action in the definitions of the categories serves the
purpose.

\begin{definition}[Bundle atlases]
\label{def:Bun-Atl}
Let $\seqname{B} \defeq (E, X, Y, \pi, G, \rho)$, be a protobundle. Then
$(\seqname{A},\seqname{B})$ is a bundle atlas, $\seqname{A}$ is a bundle
atlas of $B$, abbreviated \\*
$\isAtl[\Bun]_\Ob(\seqname{A}, \seqname{B})$
and $\seqname{A}$ is a $\pi$-$G$-$\rho$-bundle atlas of $E$ in the
coordinate space $X \times Y$, abbreviated
$\isAtl[\Bun]_\Ob(\seqname{A}, E, X, Y, \pi, G, \rho)$, iff
$\seqname{A}$ consists of a set of mutually $G$-$\rho$-compatible
$Y$-$\pi$-bundle charts of $E$ in the coordinate space $X \times Y$ that
covers $E$\footnote{
  There is no need to introduce the concept of a full
  $\pi$-$G$-$\rho$-bundle atlas because a $\pi$-$G$-$\rho$-bundle
  atlas is automatically full.
}.

By abuse of language we write $U \in \seqname{A}$ for
$U \in \pi_1[\seqname{A}]$.

\begin{remark}
The definition of a $\pi$-$G$-$\rho$-bundle atlas is by design similar
to the definition of a coordinate bundle in \cite[p.~7]{TopFib}, but
there are significant differences. This paper will use the term bundle
atlas to avoid confusion.
\end{remark}

Let
$
  \seqname{B}^\alpha \defeq
  (E^\alpha, X^\alpha, Y^\alpha, \pi^\alpha, G^\alpha, \rho^\alpha)
$,
$\alpha \prec \Alpha$, be a protobundle and
$\seqname{B} \defeq \set{{\seqname{B}^\alpha}}[\alpha \prec \Alpha]$
be a set of protobundles.
Then
\begin{multline}
\Atl[\Bun]_\Ob(\seqname{B}^\alpha) \defeq
  \set
  {{
    (
      \seqname{A},
      \seqname{B}^\alpha
    )
  }}%
  [
    {{
      \isAtl[\Bun]_\Ob
        (
          \seqname{A},
          E^\alpha,
          X^\alpha,
          Y^\alpha,
          G^\alpha,
          \pi^\alpha,
          \rho^\alpha
        )
    }}
  ]
\end{multline}

\begin{equation}
\Atl[\Bun]_\Ob \seqname{B} \defeq
  \union[\alpha \prec \Alpha]{\Atl[\Bun]_\Ob(\seqname{B}^\alpha)}
\end{equation}
\end{definition}

\begin{lemma}[Bundle atlases]
\label{lem:Bun-Atl}
Let $E,X,Y$ be topological spaces, $G$ a topological group,
$\pi \maps E \morphEpi X$ surjective and $\rho \maps Y \times G \to Y$ an
effective right action of $G$ on $Y$.

$\seqname{A}$ is a $\pi$-$G$-$\rho$-bundle atlas of $E$ in the
coordinate space $X \times Y$ iff it is an m-atlas of
$\Triv[X-\pi-]{\seqname{E}}$ in the coordinate space
$\Triv[G-\rho-]{X,Y}$.

\begin{proof}
Let $(U^i, V^i \times Y, \phi^i) \in \seqname{A}$, $i=1,2$, be
$Y$-$\pi$-bundle charts.  Then by \pagecref{lem:Grho-compat},
$(U^1, V^1 \times Y, \phi^1)$ and $(U^2, V^2 \times Y, \phi^2)$ are
$G$-$\rho$-compatible iff $(U^1, V^1 \times Y, \phi^1)$ and
$(U^2, V^2 \times Y, \phi^2)$ are compatible m-charts of $\Triv{E}$ in
the coordinate space $\Triv[G-\rho-]{X,Y}$.
\end{proof}

If $\seqname{A}$ is a $\pi$-$G$-$\rho$-bundle atlas of $E$ in the
coordinate space $X \times Y$ then every chart of $\seqname{A}$ is of
the form $(U, V \times Y, \phi)$.

\begin{proof}
The result follows from \pagecref{def:Grhomodtriv}.
\end{proof}
\end{lemma}

\begin{corollary}[Bundle atlases]
\label{cor:Bun-Atl}
Let $E,X,Y$ be topological spaces, $G$ a topological group,
$\pi \maps E \morphEpi X$ surjective, $\rho \maps Y \times G \to Y$ an
effective right action of $G$ on $Y$ and $\seqname{A}$ an m-atlas of
$\Triv{E}$ in the coordinate space $X \times Y$ such that for every
$e \in E$ there is a $Y$-$\pi$-bundle chart of $E$ in the coordinate
space $X \times Y$ in $\seqname{A}$ that covers $e$.

Then $\seqname{A}$ is a $\pi$-$G$-$\rho$-bundle atlas of $E$ in the
coordinate space $X \times Y$.

\begin{proof}
The result follows from \pagecref{lem:Ypichart}.
\end{proof}
\end{corollary}

\begin{definition}[Compatibility of charts with bundle atlases]
\label{def:piGrho-atl:extensions}
A $Y$-$\pi$-bundle chart $(U, V \times Y, \phi)$ of $E$ in the
coordinate space $X \times Y$ is $G$-$\rho$-compatible with a
$\pi$-$G$-$\rho$-bundle atlas $\seqname{A}$ iff it is
$G$-$\rho$-compatible with every chart in the atlas.
\end{definition}

\begin{lemma}[Compatibility of subcharts with bundle atlases]
\label{lem:BunCompat}
Let $\seqname{A}$ be a $\pi$-$G$-$\rho$-bundle atlas of $E$ in the
coordinate space $X \times Y$ and
$(U, V \times Y, \phi)$ a $Y$-$\pi$-bundle chart in
$\seqname{A}$. Then any subchart of $(U, V \times Y, \phi)$ is
$G$-$\rho$-compatible with $\seqname{A}$.

\begin{proof}
The result follows from \pagecref{lem:M-Compatsub} and
\pagecref{lem:Grho-compat}.
\end{proof}
\end{lemma}

\begin{lemma}[Extensions of bundle atlases]
\label{lem:piGrho-atl:extensions}
Let $\seqname{A}$ be a $\pi$-$G$-$\rho$-atlas of $\seqname{E}$ in the
coordinate space $X \times Y$ and $(U_i,V_i,\phi_i)$, $i=1,2$, be a
$\pi$-$G$-$\rho$ chart of $E$ in the coordinate space $X \times Y$
$G$-$\rho$-compatible with $\seqname{A}$ in the coordinate space $X
\times Y$.  Then $(U_1,V_1,\phi_1)$ is $G$-$\rho$-compatible with
$(U_2,V_2,\phi_2)$ in the coordinate space $X \times Y$.

\begin{proof}
The result follows from \pagecref{lem:M-ATLextensions} and
\pagecref{lem:Grho-compat}.
\end{proof}
\end{lemma}

\begin{definition}[Maximal bundle atlases]
\label{def:Bun-ATLmax}
Let $\seqname{A}$ be a $\pi$-$G$-$\rho$-bundle atlas of $E$ in the
coordinate space $X \times Y$. $\seqname{A}$ is a maximal
$\pi$-$G$-$\rho$-bundle atlas, abbreviated \\*
$
  \maximal{\isAtl[\Bun]_\Ob}
    (
      \seqname{A},
      E,
      X,
      Y,
      G,
      \pi,
      \rho
    )
$,
iff it cannot be extended by adding an additional $G$-$\rho$-compatible
$Y$-$\pi$-bundle chart.

Let
$
  \seqname{B}^\alpha \defeq
  (E^\alpha, X^\alpha, Y^\alpha, \pi^\alpha, G^\alpha, \rho^\alpha)
$,
$\alpha \prec \Alpha$, be a protobundle and
$\seqname{B} \defeq \set{{\seqname{B}^\alpha}}[\alpha \prec \Alpha]$
be a set of protobundles.
Then
\begin{multline}
\maximal[**]{\Atl[\Bun]_\Ob(\seqname{B}^\alpha)} \defeq
  \set
  {{
    (
      \seqname{A},
      \seqname{B}^\alpha
    )
  }}%
  [
    {{
      \maximal[**]
        {
          \isAtl[\Bun]_\Ob
            \biggl (
              \seqname{A},
              E^\alpha,
              X^\alpha,
              Y^\alpha,
              G^\alpha,
              \pi^\alpha,
              \rho^\alpha
            \biggr )
        }
    }}
  ]
\end{multline}

\begin{equation}
\maximal[**]{\Atl[\Bun]_\Ob} \seqname{B} \defeq
  \union
    [\alpha \prec \Alpha]
    {\maximal[**]{\Atl[\Bun]_\Ob}(\seqname{B}^\alpha)}
\end{equation}
\end{definition}

\begin{lemma}[Maximal bundle atlases]
\label{lem:Bun-ATLmax}
Let $\seqname{A}$ be a $\pi$-$G$-$\rho$-bundle atlas of $E$ in the
coordinate space $X \times Y$. $\seqname{A}$ is a maximal
$\pi$-$G$-$\rho$-bundle atlas iff it is a maximal m-atlas of $\Triv{E}$
in the coordinate space $\Triv[G-\rho-]{X,Y}$.

\begin{proof}
The result follows from \pagecref{lem:Grho-compat} and
{
  \showlabelsinline
  \cref{lem:piGrho-atl:extensions}
}
above.
\end{proof}
\end{lemma}

\begin{theorem}[Existence and uniqueness of maximal $\pi$-$G$-$\rho$-bundle atlases]
Let $\seqname{B} \defeq (E, X, Y, \pi, G, \rho)$ be a protobundle and
$\seqname{A}$ $\pi$-$G$-$\rho$-bundle atlas of $E$ in the coordinate
space $X \times Y$. Then there exists a unique maximal
$\pi$-$G$-$\rho$-bundle atlas
$\maximal{\Atlas^\Bun}(\seqname{A},\seqname{B})$ of $E$ in the
coordinate space $X \times Y$ $G$-$\rho$-compatible with $\seqname{A}$.

\begin{proof}
Let $\seqname{P}$ be the set of all $\pi$-$G$-$\rho$-bundle atlases $E$
in the coordinate space $X \times Y$ containing $\seqname{A}$ and
$G$-$\rho$ compatible in the coordinate space $X \times Y$ with
$\seqname{A}$. Let $\maximal{\seqname{P}}$ be a maximal chain of
$\seqname{P}$. Then $A'=\union{\maximal{\seqname{P}}}$ is a maximal
$\pi$-$G$-$\rho$-bundle atlas of $\seqname{E}$ in the coordinate space
$X \times Y$ $G$-$\rho$ compatible with $\seqname{A}$. Uniqueness
follows from \pagecref{def:piGrho-atl:extensions}.
\end{proof}
\end{theorem}

\begin{lemma}[Existence and uniqueness of projection for atlases in $G$-$\rho$-model spaces]
\label{lem:mGrho}
Let $E$, $X$ and $Y$ be topological spaces, $G$ a topological group,
$\rho$ an effective action of $G$ on $Y$,
$\seqname{C} \defeq (X \times Y, \catName{XY})$ a $G$-$\rho$ model space
of $X \times Y$ and $\seqname{A}$ an m-atlas of $E$ in the coordinate
model space $\seqname{C}$.  Then there exists a unique function
$\pi \maps \seqname{E} \to X$ such that for any chart $(U, V, \phi)$ in
$\seqname{A}$, $\pi \restrictto_U = \pi_1 \morphCompose \phi$.   If
$\seqname{A}$ is full then $\pi$ is surjective.

\begin{proof}
Let $(U, V, \phi)$ be an arbitrary chart in $\seqname{A}$ and define
$\pi(e \in U) = \pi_1 \morphCompose \phi(e)$. $\pi(e)$ does not depend on the
choice of chart because the morphisms of a $G$-$\rho$ model space
preserve fibers. $\pi$ is continuous because it is continuous on each
coordinate patch.

Let $x \in X$. If $\seqname{A}$ is full then there exists a chart
$(U, V, \phi)$ in $\seqname{A}$ such that $x \in \pi_1[V]$. Let $u$ be an
arbitrary point in $\phi^{-1}[\set{x} \times Y]$. Then
$\pi(u) = \pi_1(\phi(u)) = x$.
\end{proof}
\end{lemma}
\section{Bundle-atlas morphisms and functors}
\label{sec:bunatlmorph}
This section introduces a taxonomy for morphisms between bundle atlases,
defines categories of bundle atlases and constructs functors from them
to categories of m-atlases.  It only constructs reverse functors for
subcategories of $\Trivcat[\BunProd-]{\seqname{B}}$.

\subsection{Bundle-atlas morphisms}
\subsubsection{Definitions of bundle-atlas morphisms}
\begin{definition}[Bundle-atlas morphisms]
\label{def:Bun-ATLmorph}
Let
$
  \seqname{B}^i
  \defeq
  (
    E^i,
    X^i,
    Y^i,
    \pi^i,
    G^i,
    \rho^i
  )
$,
$i=1,2$, be a protobundle and let $\seqname{A}^i$ be a
$\pi^i$-$G^i$-$\rho^i$-atlas of $E^i$ in the coordinate space
\nobreaks
  {{
    $C^i = X^i \times Y^i$.
  }}
Then
$
  \morphSeqName{f}
  \defeq
  (
    \morphName{f}_E \maps E^1 \to E^2,
    \morphName{f}_X \maps X^1 \to X^2,
    \morphName{f}_Y \maps Y^1 \to Y^2,
    \morphName{f}_G \maps G^1 \to G^2
  )
$
is a $\seqname{B^1}$-$\seqname{B^2}$ bundle-atlas morphism from
$\seqname{A}^1$ to $\seqname{A}^2$, abbreviated
\nobreaks
{{
  $
    \isAtl[\Bun]_\Ar
      (
        \seqname{A^1}, \seqname{B^1},
        \seqname{A^2}, \seqname{B^2},
        \morphSeqName{f}
      )
  $,
}}
iff
\begin{enumerate}
\item all four functions are continuous
\item $\morphName{f}_G$ is a homomorphism
\item $\morphSeqName{f}$ commutes with $\pi^i$ and $\rho^i$, i.e.,
\begin{enumerate}
\item $\pi^2 \morphCompose \morphName{f}_E = \morphName{f}_X \morphCompose \pi^1$
\item
$
  \uquant
    {{y \in Y^1},{g \in G^1}}
    {{
      \morphName{f}_Y(y) \star^2 \morphName{f}_G(g) =
      \morphName{f}_Y(y \star^1 g)
    }}
$
\end{enumerate}
\item For any charts
$(U^i, V^i \times Y, \phi^i \maps U^i \toiso V^ii \times Y) \in \seqname{A}^i$.
$i=1,2$, and any $u^1 \in I \defeq U^1 \cap \morphName{f}^{-1}_E[U^2]$,
there exists a subchart
$
  (U'^1, V'^1 \times Y, \phi^1 \maps U'^1 \toiso V'^1i \times Y)
  \in \seqname{A}^1
$
at $u^1$, an open set $U'^2 \subseteq U^2$ and a chart
$
  (U'^2 ,\hat{V}'^2 \times Y, \phi'^2) \in \seqname{A}^2
$
at $\morphName{f}_E(u^1)$ such that
$\morphName{f}_E[U'^1] \subseteq U'^2$ and diagram \eqref{eq:Bun-AtlM1}
below is commutative, as shown in \cref{fig:Bun-AtlM1}.
\end{enumerate}

\begin{equation}
\begin{array}{lll}
\label{eq:Bun-AtlM1}
D \defeq & \Biggl \{
\\
  &&
  \morphName{i} \maps \set{u^1} \morphMono U'^1,
  \phi^1 \maps U'^1 \toiso V'^1 \times Y,
  \morphName{f}_X \times \morphName{f}_Y \maps V'^1 \times Y \to \hat{V}'^2 \times Y,
\\
  &&
  \morphName{f}_E \maps  U'^1 \to U'^2,
  \phi'^2 \maps U'^2 \toiso \hat{V}'^2 \times Y
\\
& \Biggr \}
\end{array}
\end{equation}

\begin{figure}[ht]
\[ \bfig
\node E1(500,0)[u^1]
\node E2(1500,0)[u^2]
\node U1(-300,-500)[U^1]
\node Up1(500,-500)[U'^1]
\node Up2(1500,-500)[U'^2]
\node U2(2000,-500)[U^2]
\node V1Y(-300,-1000)[V^1i \times Y]
\node Vp1Y(500,-1000)[V'^1 \times Y]
\node Vhp2Y(1500,-1000)[\hat{V}'^2 \times Y]
\node V2Y(2000,-1000)[V^2 \times Y]
\arrow |a|[E1`E2;\morphName{f}_E]
\arrow |r|/^{ (}->/[E1`Up1;\morphName{i}]
\arrow |r|/^{ (}->/[E2`Up2;\morphName{i}]
\arrow |r|/>->/[U1`V1Y;\phi^1]
\arrow |l|//[U1`V1Y;\iso]
\arrow |a|/_{ (}->/[Up1`U1;\morphName{i}]
\arrow |a|[Up1`Up2;\morphName{f}_E]
\arrow |r|[Up1`Vp1Y;{\phi^1}]
\arrow |l|//[Up1`Vp1Y;\iso]
\arrow |a|/^{ (}->/[Up2`U2;i]
\arrow |r|[Up2`Vhp2Y;\phi'^2]
\arrow |l|//[Up2`Vhp2Y;\iso]
\arrow |r|[U2`V2Y;\phi^2]
\arrow |l|//[U2`v2;\iso]
\arrow |a|/_{ (}->/[Vp1Y`V1Y;\morphName{i}]
\arrow |a|[Vp1Y`Vhp2Y;\morphName{f}_X \times \morphName{f}_Y]
\arrow |a|/^{ (}->/[Vhp2Y`V2Y;\morphName{i}]
\efig \]
\caption{Completed Bundle-atlas morphism}
\label{fig:Bun-AtlM1}
\end{figure}

If $\seqname{A}^1$ and $\seqname{A}^2$ are maximal
m-atlases then $\morphSeqName{f}$ is also a maximal
$\seqname{B^1}$-$\seqname{B^2}$ bundle-atlas morphism from
$\seqname{A}^1$ to $\seqname{A}^2$, abbreviated
$
  \maximal[**] {\isAtl[\Bun]_\Ar}
    (
      \seqname{A^1}, \seqname{B^1},
      \seqname{A^2}, \seqname{B^2},
      \morphSeqName{f}
    )
$.

The identity morphism of $(\seqname{A}^i, \seqname{B}^i)$ is
\begin{equation}
\Id[{{(\seqname{A}^i, \seqname{B}^i)}}] \defeq
\Bigl (
  \bigl (
    \Id[E^i],
    \Id[X^i],
    \Id[Y^i],
    \Id[G^i]
  \bigr ),
  \bigl ( \seqname{A}^i, \seqname{B}^i\bigr ),
  \bigl ( \seqname{A}^i, \seqname{B}^i\bigr )
\Bigr )
\end{equation}
This nomenclature will be justified later.
\end{definition}

\begin{definition}[Equivalence of bundle-atlas morphisms]
\label{def:Bun-ATLmorphequiv}
\leavevmode
\\*
$
  \morphSeqName{f}
  \defeq
  (
    \morphName{f}_E \maps E^1 \to E^2,
    \morphName{f}_X \maps X^1 \to X^2,
    \morphName{f}_Y \maps Y^1 \to Y^2,
    \morphName{f}_G \maps G^1 \to G^2
  )
$
is bundle equivalent to
$
  \morphSeqName{g}
  \defeq
  (
    \morphName{g}_E \maps E^1 \to E^2,
    \morphName{g}_X \maps X^1 \to X^2,
    \morphName{g}_Y \maps Y^1 \to Y^2,
    \morphName{g}_G \maps G^1 \to G^2
  )
$
as a $\seqname{B^1}$-$\seqname{B^2}$ bundle-atlas morphism from
$\seqname{A}^1$ to $\seqname{A}^2$, abbreviated $\morphSeqName{f}$ and
$\morphSeqName{g}$ are bundle equivalent $\seqname{B^1}$-$\seqname{B^2}$
bundle-atlas morphisms from $\seqname{A}^1$ to $\seqname{A}^2$, iff
$
  \seqname{B}^i
  \defeq
  (
    E^i,
    X^i,
    Y^i,
    \pi^i,
    G^i,
    \rho^i
  )
$,
$i=1,2$, is a protobundle, $\seqname{A}^i$ is a
$\pi^i$-$G^i$-$\rho^i$-atlas of $E^i$ in the coordinate space
$C^i \defeq X^i \times Y^i$, $\morphSeqName{f}$ and $\morphSeqName{g}$ are
$\seqname{B^1}$-$\seqname{B^2}$ bundle-atlas morphisms from
$\seqname{A}^1$ to $\seqname{A}^2$, $\morphName{f}_E = \morphName{g}_E$,
$\morphName{f}_X = \morphName{g}_X$ and $\morphName{f}_G = \morphName{g}_G$,

By abuse of language we shall write $\morphSeqName{f}$ is bundle
equivalent to $\morphSeqName{g}$ when the protobundles and bundle
atlases are understood by context.
\end{definition}

\subsubsection{Proclamations on bundle-atlas morphisms}
\begin{lemma}[Bundle-atlas morphisms]
\label{lem:Bun-ATLmorph}
Let
$
  \seqname{B}^i
  \defeq
  (
    E^i,
    X^i,
    Y^i,
    \pi^i,
    G^i,
    \rho^i
  )
$, $i=1,2$, be a protobundle and let
$\seqname{A}^i$ be a $\pi^i$-$G^i$-$\rho^i$-atlas of $E^i$ in the
coordinate space
\nobreaks
  {{
    $C^i \defeq X^i \times Y^i$.
  }}
Then
$
  \morphSeqName{f}
  \defeq
  (
    \morphName{f}_E \maps E^1 \to E^2,
    \morphName{f}_X \maps X^1 \to X^2,
    \morphName{f}_Y \maps Y^1 \to Y^2,
    \morphName{f}_G \maps G^1 \to G^2
  )
$
is a $\seqname{B^1}$-$\seqname{B^2}$ bundle-atlas morphism from
$\seqname{A}^1$ to $\seqname{A}^2$ iff
$\morphName{f}_X \times \morphName{f}_Y$ is a
$G^1$-$G^2$-$\rho^1$-$\rho^2$ morphism from
$\Triv[G^1-\rho^1-]{X^1,Y^1}$ to $\Triv[G^2-\rho^2-]{X^2,Y^2}$ and
$(\morphName{f}_E, \morphName{f}_X \times \morphName{f}_Y)$ is a
$\Triv[X^1-\pi^1-]{\seqname{E^1}}$-$\Triv[X^2-\pi^2-]{\seqname{E^2}}$
m-atlas morphism from $\seqname{A}^1$ to $\seqname{A}^2$ in the
coordinate spaces $\Triv[G^1-\rho^1-]{X^1,Y^1}$,
$\Triv[G^2-\rho^2-]{X^2,Y^2}$.

\begin{proof}
\leavevmode
\begin{enumerate}
\item
By \pagecref{lem:Bun-Atl}, $\seqname{A}^i$ is an m-atlas of
$\Triv[X^i-\pi^i-]{E^i}$ in the coordinate space
$\Triv[G^i-\rho^i-]{X^i,Y^i}$, every coordinate function preserves
fibers and every chart of $\seqname{A}$ is of the form
$(U, V \times Y, \phi)$.

\item
If $\morphSeqName{f}$ is a $\seqname{B^1}$-$\seqname{B^2}$ bundle-atlas
morphism then $\morphName{f}_X \times \morphName{f}_Y$ is a model function
from $\Triv[G^1-\rho^1-]{X^1,Y^1}$ to $\Triv[G^2-\rho^2-]{X^2,Y^2}$
by \pagecref{lem:Grhomodtrivmorph}
and $\morphName{f}_G$ is the function asserted to exist in
{
  \showlabelsinline
  \cref{eq:Grho}
}
(preservation of group action) of \pagecref{def:Grhomorph},
so $\morphName{f}_X \times \morphName{f}_Y$ is a
$G^1$-$G^2$-$\rho^1$-$\rho^2$ morphism from
$\Triv[G^1-\rho^1-]{X^1,Y^1}$ to $\Triv[G^2-\rho^2-]{X^2,Y^2}$.

$\morphName{f}_E$ is continuous and preserves fibers, hence a model
function from $\Triv[X^1-\pi^1-]{E^1}$ to $\Triv[X^2-\pi^2-]{E^2}$, and
$\morphName{f}_X \times \morphName{f}_Y$ is continuous and preserves
fibers, hence a model function from $\Triv[G^1-\rho^1-]{X^1,Y^1}$ to
$\Triv[G^2-\rho^2-]{X^2,Y^2}$.

Let $(U^i, V^i \times Y^i, \phi^i) \in \seqname{A}^i$, $i=1,2$,
$u^1 \in I \defeq U^1 \cap \morphName{f}^{-1}_E[U^2]$. Let the chart
$
  (U'^1, V'^1 \times Y, \phi^1 \maps U'^1 \toiso V'^1)
  \in \seqname{A}^1
$
at $u^1$, the open set $U'^2 \subseteq U^2$ and the chart
$
  (U'^2 ,\hat{V}'^2 \times Y, \phi'^2 \maps U'^2 \toiso \hat{V}'^2 \times Y)
  \in \seqname{A}^2
$
at $\morphName{f}_E(u^1)$ be as in \cref{def:Bun-ATLmorph} above, i.e.,
$\morphName{f}_E[U'^1] \subseteq U'^2$ and diagram \eqref{eq:Bun-AtlM1}
is commutative, as shown in
{
  \showlabelsinline
  \cref{fig:Bun-AtlM1}.
}
This is a relabling of the commutative diagram in
\pagecref{def:M-ATLmorph}.

\item
If $\morphName{f}_X \times \morphName{f}_Y$ is a
$G^1$-$G^2$-$\rho^1$-$\rho^2$ morphism from
$\Triv[G^1-\rho^1-]{X^1,Y^1}$ to $\Triv[G^2-\rho^2-]{X^2,Y^2}$ and \\*
$(\morphName{f}_E, \morphName{f}_X \times \morphName{f}_Y)$ is an
$\Triv[X^1-\pi^1-]{E^1}$-$\Triv[X^2-\pi^2-]{E^2}$ m-atlas morphism from
$\seqname{A}^1$ to $\seqname{A}^2$ in the coordinate spaces
$\Triv[G^1-\rho^1-]{X^1,Y^1}$, $\Triv[G^2-\rho^2-]{X^2,Y^2}$, then let
$u^1 \in E^1$ be arbitrary, $u^2 \defeq \morphName{f}_E(u^1)$ and
\nobreaks
  {{
    $(U^i, V^i \times Y, \phi^i) \in \seqname{A}^i$
  }}
be a chart at $u^i$. Let the chart
\nobreaks
  {{
    $
      (U'^1, V'^1 \times Y, \phi^1 \restrictto^{U'^1}_{V'^1 \times Y})
      \in \seqname{A}^1
    $
  }}
at $u^1$, the open set $U'^2 \subseteq U^2$ and the chart
\nobreaks
  {{
    $
      (
        U'^2,
        \hat{V}'^2 \times Y,
        \phi'^2 \restrictto^{U'^2}_{\hat{V}'^2 \times Y}
      )
      \in \seqname{A}^2
    $
  }}
at $\morphName{f}_E(u^1)$ be as in \cref{def:M-ATLmorph}, i.e.,
$\morphName{f}_E[U'^1] \subseteq U'^2$ and diagram \eqref{eq:AtlM1}
is commutative, as shown in
{
  \showlabelsinline
  \pagecref{fig:AtlM1},
}
with relabling as in \cref{fig:Bun-AtlM2} below.

\begin{figure}[ht]
\[ \bfig
\node u1(500,0)[u^1]
\node u2(1500,0)[u^2]
\node U1(-300,-500)[U^1]
\node Up1(500,-500)[U'^1]
\node Up2(1500,-500)[U'^2]
\node U2(2000,-500)[U^2]
\node V1Y(-300,-1000)[V^1 \times Y]
\node Vp1Y(500,-1000)[V'^1 \times Y]
\node Vhp2Y(1500,-1000)[{\hat{V}'^2 \times Y}]
\node V2Y(2000,-1000)[V^2 \times Y]
\node Vp1(500,-1500)[V'^1]
\node Vhp2(1500,-1500)[\hat{V}'^2]
\arrow |a|[u1`u2;\morphName{f}_E]
\arrow |r|/^{ (}->/[u1`Up1;\morphName{i}]
\arrow |r|/^{ (}->/[u2`Up2;\morphName{i}]
\arrow |r|/>->/[U1`V1Y;\phi^1]
\arrow |l|//[U1`V1Y;\iso]
\arrow |a|/_{ (}->/[Up1`U1;\morphName{i}]
\arrow |a|[Up1`Up2;\morphName{f}_E]
\arrow |r|[Up1`Vp1Y;{\phi^1}]
\arrow |l|//[Up1`Vp1Y;\iso]
\arrow |a|/^{ (}->/[Up2`U2;\morphName{i}]
\arrow |r|[Up2`Vhp2Y;\phi'^2]
\arrow |l|//[Up2`Vhp2Y;\iso]
\arrow |r|[U2`V2Y;\phi^2]
\arrow |l|//[U2`V2Y;\iso]
\arrow |a|/_{ (}->/[Vp1Y`V1Y;\morphName{i}]
\arrow |a|[Vp1Y`Vhp2Y;\morphName{f}_X \times \morphName{f}_Y]
\arrow |r|[Vp1Y`Vp1;\morphName{\pi}_1]
\arrow |r|[Vhp2Y`Vhp2;\morphName{\pi}_1]
\arrow |a|/^{ (}->/[Vhp2Y`V2Y;\morphName{i}]
\arrow |a|[Vp1`Vhp2;\morphName{f}_X]
\efig \]
\caption{Completed Bundle-atlas morphism with projection}
\label{fig:Bun-AtlM2}
\end{figure}

\begin{description}
\item[All four functions are continuous:]
$\morphName{f}_E$, $\morphName{f}_X \times \morphName{f}_Y$ and
$\morphName{f}_G$ are continuous by \pagecref{def:M-ATLmorph} and
{
  \showlabelsinline
  \pagecref{def:Grhomorph}.
}
$\morphName{f}_X$ and $\morphName{f}_Y$ are continuous because
$\morphName{f}_X \times \morphName{f}_Y$ is continuous.

\item[$\morphName{f}_G$ is a homomorphism:]
$\morphName{f}_G$ is a homomorphism by \pagecref{def:Grhomorph}.

\item[$\morphSeqName{f}$ commutes with $\pi^i$ and $\rho^i$:]
\begin{enumerate}
\leavevmode
\item
$
  \pi^2 \morphCompose \morphName{f}_E =
  \pi_1 \morphCompose \phi^2 \morphCompose \morphName{f}_E =
  \morphName{f}_X \morphCompose \pi_1 \morphCompose \phi^1 =
  \morphName{f}_X \morphCompose \pi^1
$
on $U'^1$. Since $u$ is arbitray,
$\pi^2 \morphCompose \morphName{f}_E = \morphName{f}_X \morphCompose \pi^1$ on all
of $E^1$.

\item
$
  \uquant
    {{y \in Y^1},{g \in G^1}}
    {{
      \morphName{f}_Y(y) \star^2 \morphName{f}_G(g) =
      \morphName{f}_Y(y \star^1 g)
    }}
$
by \cref{eq:Grho} in
{
  \showlabelsinline
  \pagecref{def:Grhomorph}.
}
\end{enumerate}

\item[Commutation relation of charts:]
The commutativity of \cref{def:Bun-ATLmorph} follows from that of
\cref{fig:Bun-AtlM2}.
\end{description}
\end{enumerate}
\end{proof}
\end{lemma}

\begin{corollary}[Bundle-atlas morphisms]
\label{cor:Bun-ATLmorph}
Let
$
  \seqname{B}^i
  \defeq
  (
    E^i,
    X^i,
    Y^i,
    \pi^i,
    G^i,
    \rho^i
  )
$,
$i=1,2,3$, be a protobundle,
$\seqname{A}^i$ be a $\pi^i$-$G^i$-$\rho^i$-atlas of $E^i$ in the
coordinate space $C^i = X^i \times Y^i$ and
$
  \morphSeqName{f}^i
  \defeq
  (
    \morphName{f}^i_E \maps E^i \to E^{i+1},
    \morphName{f}^i_X \maps X^i \to X^{i+1},
    \morphName{f}^i_Y \maps Y^i \to Y^{i+1},
    \morphName{f}^i_G \maps G^i \to G^{i+1}
  )
$
a $\seqname{B^i}$-$\seqname{B^{i+1}}$ bundle-atlas morphism from
$\seqname{A}^i$ to $\seqname{A}^{i+i}$.

$\morphSeqName{f}^2 \morphCompose[()] \morphSeqName{f}^1$ is a bundle-atlas
morphism from $\seqname{A}^1$ to $\seqname{A}^3$.

\begin{proof}
$
  \morphName{f}^1_X \times \morphName{f} 1_Y \morphCompose
  \morphName{f}^2_X \times \morphName{f}^2_Y
$
is a $G^1$-$G^3$-$\pi^1$-$\pi^3$ morphism from
$\Triv[G^1-\rho^1-]{X^1,Y^1}$ to $\Triv[G^3-\rho^3-]{X^3,Y^3}$ by
\cref{grhomorph:comp} of
{
  \showlabelsinline
  \fullcref{lem:Grhomorph} on
  \cpageref{grhomorph:comp}
}
and
$
  (\morphName{f}^1_E, \morphName{f}^1_X \times \morphName{f}^1_Y)
  \morphCompose[()]
  (\morphName{f}^2_E, \morphName{f}^2_X \times \morphName{f}^2_Y)
$
is a $\Triv[X^1-\pi^1]{E^1}$-$\Triv[X^3-\pi^3]{E^3 }$ m-atlas morphism
from $\seqname{A}^1$ to $\seqname{A}^3$ in the coordinate space
$\Triv[G^3-\rho^3-]{X^3,Y^3}$ by \cref{M-morphcomp:morphcompmorph} of
{
  \showlabelsinline
  \fullcref{lem:M-ATLmorphcomp} on
  \cpageref{M-morphcomp:morphcompmorph}.
}
\end{proof}
\end{corollary}

\begin{lemma}[Composition of equivalent bundle-atlas morphisms]
\label{lem:Bun-ATLmorphequiv}
Let
\begin{enumerate}
\item
$
  \seqname{B}^i
  \defeq
  (
    E^i,
    X^i,
    Y^i,
    \pi^i,
    G^i,
    \rho^i
  )
$, $i=1,2,$, be a protobundle
\item
$\seqname{A}^i$ be a $\pi^i$-$G^i$-$\rho^i$-atlas of $E^i$ in the coordinate space
$C^i = X^i \times Y^i$
\item
$
  \morphSeqName{f}^i \defeq
  (
    \morphName{f}^i_E \maps E^i \to E^{i+1},
    \morphName{f}^i_X \maps X^i \to X^{i+1},
    \morphName{f}^i_Y \maps Y^i \to Y^{i+1},
    \morphName{f}^i_G \maps G^i \to G^{i+1}
  )
$
and
$
  \morphSeqName{g}^i \defeq
  (
    \morphName{g}^i_E \maps E^i \to E^{i+1},
    \morphName{g}^i_X \maps X^i \to X^{i+1},
    \morphName{g}^i_Y \maps Y^i \to Y^{i+1},
    \morphName{g}^i_G \maps G^i \to G^{i+1}
  )
$,
$i=1,2$,
be bundle equivalent $\seqname{B^i}$-$\seqname{B^{i+1}}$ bundle-atlas
morphisms from $\seqname{A}^i$ to $\seqname{A}^{i+1}$
\end{enumerate}

Then $\morphSeqName{f}^2 \morphCompose[()] \morphName{f}^1$ is bundle
equivalent to $\morphSeqName{g}^2 \morphCompose[()] \morphName{g}^1$ as a
$\seqname{B^1}$-$\seqname{B^3}$ bundle-atlas morphism from
$\seqname{A}^1$ to $\seqname{A}^3$.

\begin{proof}
$\morphName{f}^1_E = \morphName{g}^1_E$, and
$\morphName{f}^2_E = \morphName{g}^2_E$, hence
$
  \morphName{f}^2_E \morphCompose \morphName{f}^1_E =
  \morphName{g}^2_E \morphCompose \morphName{g}^1_E
$.
$\morphName{f}^1_X = \morphName{g}^1_X$, and
$\morphName{f}^2_X = \morphName{g}^2_X$, hence
$
  \morphName{f}^2_X \morphCompose \morphName{f}^1_X =
  \morphName{g}^2_X \morphCompose \morphName{g}^1_X
$.
$\morphName{f}^1_G = \morphName{g}^1_G$, and
$\morphName{f}^2_G = \morphName{g}^2_G$, hence
$
  \morphName{f}^2_G \morphCompose \morphName{f}^1_G =
  \morphName{g}^2_G \morphCompose \morphName{g}^1_G
$.
\end{proof}
\end{lemma}

\subsection{Categories of bundle atlases and functors}
\label{sub:Bun-cat}

\begin{definition}[Sets of bundle atlas morphisms]
{
  \showlabelsinline
  \label{def:Bun-ATLsets}
}
Let \\*
$
  \seqname{B}^\alpha \defeq
  (E^\alpha, X^\alpha, Y^\alpha, \pi^\alpha, G^\alpha, \rho^\alpha)
$,
$\alpha \prec \Alpha$, be a protobundle and
$\seqname{B} \defeq \set{{\seqname{B}^\alpha}}[\alpha \prec \Alpha]$
be a set of protobundles.

\begin{multline}
  \maximal[*]
    {
      \Atl[\Bun]_\Ar
        (
          \seqname{B}^\alpha,
          \seqname{B}^\beta
        )
    }
  \defeq
  \set
    {
      \bigl (
        \morphSeqName{f},
        (\seqname{A}^\alpha, \seqname{B}^\alpha),
        (\seqname{A}^\beta, \seqname{B}^\beta)
      \bigr )
    }%
    [
      {
        \maximal[*]
          {
            \isAtl[\Bun]_\Ar
              (
                \seqname{A}^\alpha, \seqname{B}^\alpha,
                \seqname{A}^\beta, \seqname{B}^\beta,
                \morphSeqName{f}
              )
          }
      }
    ]*
\end{multline}

\begin{multline}
  \maximal[*]
    {
      \Atl[\Bun]_\Ar
        (
          \seqname{B}^\alpha
        )
    }
  \defeq
  \set
    {
      \bigl (
        \morphSeqName{f},
        (\seqname{A}^\alpha, \seqname{B}^\alpha),
        (\seqname{A}^\alpha, \seqname{B}^\alpha)
      \bigr )
    }%
    [
      {
        \maximal[*]
          {
            \isAtl[\Bun]_\Ar
              (
                \seqname{A}^\alpha, \seqname{B}^\alpha,
                \seqname{A}^\alpha, \seqname{B}^\alpha,
                \morphSeqName{f}
              )
          }
      }
    ]*
\end{multline}

\begin{equation}
\maximal[*]{\Atl[\Bun]_\Ar} \seqname{B}
\defeq
\union
  [
    \seqname{B^\mu \in \seqname{B}},
    \seqname{B}^\nu \in \seqname{B}
  ]
  {
    \maximal[*]{\Atl[\Bun]_\Ar}
      (
        \seqname{B}_\mu,
        \seqname{B}_\nu
      )
  }
\end{equation}
\end{definition}

\begin{definition}[Categories of bundle atlases]
\label{def:Bun-Atlcat}
Let \\*
$
  \seqname{B}^\alpha \defeq
  (E^\alpha, X^\alpha, Y^\alpha, \pi^\alpha, G^\alpha, \rho^\alpha)
$,
$\alpha \prec \Alpha$, be a protobundle and
$\seqname{B} \defeq \set{{\seqname{B}^\alpha}}[\alpha \prec \Alpha]$
be a set of protobundles. Then

\begin{equation}
\maximal[*]
  {
    \Atl[\Bun]
      (
        \seqname{B}^\alpha
      )
  }
\defeq
  \biggl (
    \maximal[*]
      {
        \Atl[\Bun]_\Ob
          (
            \seqname{B}^\alpha
          )
      },
    \maximal[*]
      {
      \Atl[\Bun]_\Ar
        (
          \seqname{B}^\alpha,
          \seqname{B}^\alpha
        )
      },
    \morphCompose[A]
  \biggr )
\end{equation}

\begin{equation}
\maximal[*]{\Atl[\Bun]} \seqname{B}
\defeq
  \biggl (
    \maximal[*]{\Atl[\Bun]_\Ob}
      (
        \seqname{B}
      ),
    \maximal[*]{\Atl[\Bun]_\Ar}
      (
        \seqname{B}
      ),
    \morphCompose[A]
  \biggr )
\end{equation}
\end{definition}

\begin{lemma}[\ensuremath{\Atl[\Bun] \seqname{B}} is a category]
\label{lem:Bun-ATLiscat}
Let
$
  \seqname{B}^\alpha \defeq
  (E^\alpha, X^\alpha, Y^\alpha, \pi^\alpha, G^\alpha, \rho^\alpha)
$,
$\alpha \prec \Alpha$, be a protobundle and
$\seqname{B} \defeq \set{{\seqname{B}^\alpha}}[\alpha \prec \Alpha]$
be a set of protobundles.
Then $\maximal[*]{\Atl[\Bun] \seqname{B}}$ is a category.

Let
$
  (\seqname{A}^\alpha, \seqname{B}^\alpha)
  \in
  \Atl[\Bun]_\Ob \seqname{B}
$.
Then $\Id[{{(\seqname{A}^\alpha, \seqname{B}^\alpha)}}]$ is the
identity morphism for $(\seqname{A}^\alpha, \seqname{B}^\alpha)$.

\begin{proof}
Let $(\seqname{A}^i,\seqname{B}^i)$, $i \in [1,3]$,
be an object of $\Atl[\Bun] \seqname{B}$ and
let \\
$
  \morphName{m}^i \defeq
  \bigl (
    \morphSeqName{f}^i,
    (\seqname{A}^i,\seqname{B}^i),
    (\seqname{A}^{i+1},\seqname{B}^{i+1})
  \bigr )
$
be a morphism of $\Atl[\Bun] \seqname{B}$. Then
\begin{description}
\item[Composition:]
$
  \bigl (
    \morphSeqName{f}^2 \morphCompose[()] \morphSeqName{f}^1,
    (\seqname{A}^1,\seqname{E}^1,\seqname{C}^1),
    (\seqname{A}^3,\seqname{E}^3,\seqname{C}^3)
  \bigr )
$
is a morphism of $\Atl[\Bun] \seqname{B}$ by
\pagecref{cor:Bun-ATLmorph}.

\item[Associativity:]
Composition is associative by \pagecref{lem:labcomp}.

\item[Identity:]
$\Id[{{(\seqname{A}^i, \seqname{E}^i, \seqname{C}^i)}}]$ is an identity
morphism by \cref{lem:labcomp}.
\end{description}
\end{proof}
\end{lemma}

\begin{definition}[Functor from Bundle atlases to m-atlases]
\label{def:FuncBuntoM}
Let \\
$
  \seqname{B}^i
  \defeq
  (
    E^i,
    X^i,
    Y^i,
    \pi^i,
    G^i,
    \rho^i
  )
$,
$i=1,2$, be a protobundle, let $\seqname{A}^i$ be a
$\pi^i$-$G^i$-$\rho^i$-atlas of $E^i$ in the coordinate space
$C^i = X^i \times Y^i$ and let \\*
$
  \morphSeqName{f}
  \defeq
  (
    \morphName{f}_E \maps E^1 \to E^2,
    \morphName{f}_X \maps X^1 \to X^2,
    \morphName{f}_Y \maps Y^1 \to Y^2,
    \morphName{f}_G \maps G^1 \to G^2
  )
$
be a $\seqname{B^1}$-$\seqname{B^2}$ bundle-atlas morphism from
$\seqname{A}^1$ to $\seqname{A}^2$. Then
\begin{equation}
\Functor^\Bun_{\Bun,\M} (\seqname{A}^i, \seqname{B}^i)
\defeq
\biggl (
  \seqname{A}^i, \Triv[X^i-\pi^i]{E^i}, \Triv[G^i-\rho^i-]{X^i,Y^i}
\biggr )
\end{equation}
\begin{multline}
\Functor^\Bun_{\Bun,\M}
  \bigl (
    \morphSeqName{f},
    (\seqname{A}^1, \seqname{B}^1),
    (\seqname{A}^2, \seqname{B}^2)
  \bigr )
\defeq \\
  \Biggl (
    (\morphSeqName{f}_E, \morphSeqName{f}_X \times \morphSeqName{f}_Y),
    \biggl (
      \seqname{A}^1, \Triv[X^1-\pi^1]{E^1}, \Triv[G^1-\rho^1-]{C^1}
    \biggr ),
    \biggl (
      \seqname{A}^2, \Triv[X^2-\pi^2]{E^2}, \Triv[G^2-\rho^2-]{C^2}
    \biggr )
  \Biggr )
\end{multline}
\end{definition}

\begin{theorem}[Functor from Bundle atlases to m-atlases]
\label{the:FuncBuntoM}
Let $\seqname{E}$ be a set of topological spaces,
$
  \seqname{B}^\alpha \defeq
  (E^\alpha \in \seqname{E}, X^\alpha, Y^\alpha, \pi^\alpha, G^\alpha, \rho^\alpha)
$,
$\alpha \prec \Alpha$, be a protobundle, \\
$\seqname{B} \defeq \set{{\seqname{B}^\alpha}}[\alpha \prec \Alpha]$
be a set of protobundles and
$
  \seqname{C}^\alpha \defeq
  \Triv[G^\alpha-\rho^\alpha-]{(X^\alpha, Y^\alpha)}
$.
Then $\Functor^\Bun_{\Bun,\M}$ is a functor from
$\maximal[*]{\Atl[\Bun]} \seqname{B}$ to
$
  \maximal[*]
    {
      \Atl
        \Bigl (
          \Trivcat{\seqname{E}}, \Trivcat[\Bun-]{\seqname{B}}
        \Bigr )
    }
$
and a functor from \\*
$\maximal[*]{\Atl[\Bun]} \seqname{B}$ to
$
  \maximal[*]
    {
      \Atl
        \Bigl (
          \Trivcat{\seqname{E}}, \Trivcat[\BunProd-]{\seqname{B}}
        \Bigr )
    }
$.

\begin{proof}
Let $o^i \defeq (\seqname{A}^i, B^i)$, $i \in [1,3]$, be an object of
$\Atl[\Bun](\seqname{B})$,
$
  o'^i \defeq \Functor^\Bun_{\Bun,\M} o^i = \\*
  \Bigl (
    \seqname{A}^i,
    \Triv[X^i-\pi^i]{E^i},
    \Triv[G^i-\rho^i-]{C^i}
  \Bigr )
$
be the corresponding object of
$
  \Atl
  \Bigl (
    \Trivcat{\seqname{E}},
    \Trivcat[\Bun-]{\seqname{B}}
  \Bigr )
$, \\
$
\morphName{m}^i \defeq
  \bigl (
    \morphSeqName{f}^i,
    o^i,
    o^{i+1}
  \bigr )
$, $i=1,2$,
be a morphism of $\Atl[\Bun](\seqname{B})$ from $o^i$ to $o^{i+1}$ and
\begin{align*}
  \morphName{m}'^i
  & \defeq
  \Functor^\Bun_{\Bun,\M} \morphName{m}^i
\\*
  & =
  \Biggl (
    (\morphName{f}^i_E, \morphName{f}^i_X \times \morphName{f}^i_Y),
    \biggl (
      \seqname{A}^i,
      \Triv[X^i-\pi^i]{E^i},
      \Triv[G^i-\rho^i-]{C^i}
    \biggl ),
    \biggl (
      \seqname{A}^{i+1},
      \Triv[X^{i+i}-\pi^{i+i}]{E^{i+i}},
      \Triv[G^{i+1}-\rho^{i+1}-]{C^2}
    \biggr )
  \Biggr )
\\*
  & =
  \bigl (
    (\morphName{f}^i_E, \morphName{f}^i_X \times \morphName{f}^i_Y),
    o'^i,
    o'^{i+1}
  \bigr )
\end{align*}
be the corresponding morphism of
$\Atl \biggl ( \Trivcat{\seqname{E}}, \Trivcat[\Bun-]{\seqname{B}} \biggr )$.

\begin{description}
\item[Preservation of endpoints:]
$\Triv[X^i-\pi^i]{E^i} \objin \Trivcat{\seqname{E}}$.
$\Triv[G^i-\rho^i-]{C^i} \arin \Trivcat[\BunProd-]{\seqname{B}}$.

By \pagecref{the:FuncBuntoM},
$(\morphName{f}_E, \morphName{f}_X \times \morphName{f}_Y)$ is a
$\Triv[X^i-\pi^i-]{E^i}$-$\Triv[X^{i+i}-\pi^{i+i}-]{E^{i+i}}$ m-atlas
morphism from $\seqname{A}^i$ to $\seqname{A}^{i+1}$ in the coordinate
spaces $\Triv[G^i-\rho^i-]{X^i,Y^i}$,
$\Triv[G^{i+1}-\rho^{i+1}-]{X^{i+1},Y^{i+1}}$.

\begin{align*}
  \Functor^\Bun_{\Bun,\M} \morphName{m}^i
  & =
  \Functor^\Bun_{\Bun,\M}
    \bigl (
      \morphSeqName{f}^i,
      o^i,
      o^{i+1}
    \bigr )
\\*
  & =
  \bigl (
    (
      \morphName{f}^i_E,
      \morphName{f}^i_X \times \morphName{f}^i_Y
    ),
    o'^i,
    o'^{i+1}
  \bigr )
\end{align*}

\item[Identity:]
\begin{align*}
  \Functor^\Bun_{\Bun,\M} \Id[{{(\seqname{A}^i, \seqname{B}^i)}}]
  & =
  \Functor^\Bun_{\Bun,\M}
    \bigl (
      (\Id[E^i], \Id[X^i], \Id[Y^i], \Id[G^i]),
      o^i,
      o^i
    \bigr )
\\*
  & =
  \bigl (
    (\Id[{\Triv[X^i-\pi^i]{E^i}}], \Id[C^i]),
    o'^i,
    o'^i
  \bigr )
\\*
  & =
  \Bigl (
    (\Id[{\Triv[X^i-\pi^i]{E^i}}], \Id[C^i]),
    \Functor^\Bun_{\Bun,\M} o^i,
    \Functor^\Bun_{\Bun,\M} o^i
  \Bigr )
\\*
  & =
  \Id[{\Functor^\Bun_{\Bun,\M} o^i}]
\end{align*}

\item[Composition:]
\begin{align*}
  m^2 \morphCompose[A] m^1
  & =
  \bigl (
    (
      \morphName{f}^2_0 \morphCompose \morphName{f}^1_0,
      \morphName{f}^2_X \morphCompose \morphName{f}^1_X,
      \morphName{f}^2_Y \morphCompose \morphName{f}^1_Y,
      \morphName{f}^2_G \morphCompose \morphName{f}^1_G
    ),
    o^1,
    o^3
  \bigr )
\\*
  \Functor^\Bun_{\Bun,\M} (m^2 \morphCompose m^1)
  & =
    \bigl (
      (
        \morphName{f}_E^2 \morphCompose \morphName{f}_E^1,
        (\morphName{f}_X^2 \morphCompose \morphName{f}_X^1) \times
        (\morphName{f}_Y^2 \morphCompose \morphName{f}_Y^1)
      ),
      o'^1,
      o'^3
    \bigr )
\\*
  & =
    \bigl (
      (
        \morphName{f}_E^2 \morphCompose \morphName{f}_E^1,
        (\morphName{f}^2_X \times \morphName{f}^2_Y) \morphCompose
        (\morphName{f}^1_X \times \morphName{f}^1_Y)
      ),
      o'^1,
      o'^3
    \bigr )
\\*
  & =
  \Functor^\Bun_{\Bun,\M} m^2 \morphCompose[A] \Functor^\Bun_{\Bun,\M} m^1
\end{align*}
\end{description}
\end{proof}
\end{theorem}

\begin{lemma}[Base space functions derived from bundle-atlas morphisms]
\label{lem:Bun-Atlbasemorph}
Let $\seqname{E}^i$, $i=1,2$, be a model space, $X^i$, $Y^i$ topological
spaces, $G^i$ a topological group, $\rho^i$ an effective action of $G^i$
on $Y^i$, $\seqname{C}^i \defeq (X^i \times Y^i, \catName{XY}^i)$ a
$G^i$-$\rho^i$ model space of $X^i \times Y^i$,
$\morphName{f}_C \maps \seqname{C}^1 \to \seqname{C}^2$ a
$G^1$-$G^2$-$\rho^1$-$\rho^2$ morphism of
$X^1 \times Y^1$ to $X^2 \times Y^2$, i.e., a model function that
preserves group action, $\seqname{A}^i$ an m-atlas of $E^i$ in the
coordinate model space $\seqname{C}^i$ and
$
  \morphSeqName{f} \defeq
  (
    \morphName{f}_E \maps \seqname{E^1} \to \seqname{E}^2,
    \morphName{f}_C \maps \seqname{C}^1 \to \seqname{C}^2
  )
$
an $\seqname{E}^1$-$\seqname{E}^2$ m-atlas morphism of $\seqname{A}^1$
to $\seqname{A}^2$ in the coordinate spaces $\seqname{C}^1$,
$\seqname{C}^2$.

Then there exists a unique function $\morphName{f}_X \maps X^1 \to X^2$
such that $\morphName{f}_X \morphCompose \pi_1 = \pi_1 \morphCompose \morphName{f}_C$.

\begin{proof}
Let $x$ in $X^1$, $y,y'$ in $Y^1$, Then
$\morphName{f}_C$ preserves fibers, i.e., \\*
$
\pi_1 \bigl ( \morphName{f}_C(x,y) \bigr ) =
\pi_1 \bigl ( \morphName{f}_C(x,y') \bigr )
$.
by \pagecref{lem:Grhomorph}. Define
$\morphName{f}_X(x) \defeq \pi_1 \bigl ( \morphName{f}_C(x,y) \bigr )$
\end{proof}

\begin{remark}
There need not exist $\morphName{f}_Y \maps Y^1 \to Y^2$ such that
$\morphName{f}_C = \morphName{f}_X \times \morphName{f}_Y$.
\end{remark}
\end{lemma}

\begin{definition}[Functor from m-atlases to Bundle atlases]
\label{def:FuncMtoBun}
Let $\catSeqName{E}$ be a trivial model category, $\catName{C}$ be a
trivial product coordinate model category with objects \\*
$
\set
  {
    \seqname{C}^\alpha \defeq
    (X^\alpha \times Y^\alpha, \catName{XY}^\alpha)
  }%
  [\alpha \prec \Alpha]
$,
$G$ a group valued function on $\Ob(\seqname{C})$ and $\rho$ a function
valued function on $\Ob(\seqname{C})$ such that for every
$
  \seqname{C}^\alpha \defeq
  (X^\alpha \times Y^\alpha, \catName{XY}^\alpha) \objin \catName{C}
$,
$\rho(\seqname{C}^\alpha)$ is an effective action of
$G(\seqname{C}^\alpha)$ on $Y^\alpha$ and
$
\Triv%
  [
    G(\seqname{C}^\alpha)-\rho(\seqname{C}^\alpha)-
  ]
  {X^\alpha, Y^\alpha} =
  \seqname{C}^\alpha
$.

Let $\seqname{E}^i \defeq (E^i, \catName{E}^i) \objin \catSeqName{E}$,
$i=1,2$, be a trivial model space,
$
  \seqname{C}^i \defeq (X^i \times Y^i, \catName{XY}^i)
  \objin \catName{C}
$
$G^i \defeq G(\seqname{C}^i)$, $\rho^i \defeq \rho(\seqname{C}^i)$,
$\seqname{A}^i$ a full m-atlas of $\seqname{E}^i$ in the coordinate
space $\seqname{C}^i$,
$\pi^i \maps E^i \morphEpi X^i$ the unique function asserted in
\cref{lem:mGrho} and
$
  \seqname{B}^i \defeq
    (
      E^i,
      X^i,
      Y^i,
      \pi^i,
      G^i,
      \rho^i
    )
$.
Then define
\begin{equation}
\Functor^\Bun_{\M-G-\rho,\Bun}
(\seqname{A}^i, \seqname{E}^i, \seqname{C}^i)
\defeq
(\seqname{A}^i, \seqname{B}^i)
\end{equation}

Let
$
  \morphSeqName{f} \defeq
  (
    \morphName{f}_E \maps \seqname{E^1} \to \seqname{E}^2,
    \morphName{f}_C \defeq \morphName{f}_X \times \morphName{f}_Y \maps
    \seqname{C}^1 \to \seqname{C}^2
  )
$
be an $\seqname{E}^1$-$\seqname{E}^2$ m-atlas morphism of
$\seqname{A}^1$ to $\seqname{A}^2$ in the coordinate spaces
$\seqname{C}^1, \seqname{C}^2$ that preserves the group action,
$\morphName{f}_G \maps G^1 \to G^2$ the unique function asserted to exist
in \cref{eq:Grho} (preservation of group action) of
\pagecref{def:Grhomorph}.

Then
\begin{multline}
\Functor^\Bun_{\M-G-\rho,\Bun}
  \bigl (
    \morphSeqName{f},
    (\seqname{A}^1, \seqname{E^1}, \seqname{C}^1),
    (\seqname{A}^2, \seqname{E^2}, \seqname{C}^2)
  \bigr )
\defeq \\
  \bigl (
    (\morphName{f}_E, \morphName{f}_X, \morphName{f}_Y, \morphName{f}_G),
    (\seqname{A}^1, \seqname{B}^1),
    (\seqname{A}^2, \seqname{B}^2)
  \bigr )
\end{multline}
\end{definition}

\begin{theorem}[Functor from m-atlases to bundle atlases]
\label{the:FuncMtoBun}
Let
\begin{enumerate}
\item
$\catName{C}$ be a trivial product coordinate model category,

\item
$\catName{E}$ be a model category,

\item
$G$ be a group valued function on $\Ob(\catName{C})$

\item
$\rho$ be a function valued function on
$\Ob(\seqname{C})$ such that for every \\*
$
  \seqname{C}^\alpha \defeq
  (X^\alpha \times Y^\alpha, \catName{XY}^\alpha) \objin \catName{C}
$,
$\rho(\seqname{C}^\alpha)$ is an effective action of
$G(\seqname{C}^\alpha)$ on $Y^\alpha$ and
$
  \seqname{C}^\alpha =
  \Triv%
    [
      G(\seqname{C}^\alpha)-\rho(\seqname{C}^\alpha)-
    ]
    {X^\alpha, Y^\alpha}
$.

\item
$\pi$ be the unique function valued function on
$\full{\Atl_\Ob}(\catName{E}, \catName{C})$ such that for every
$
\bigl (
  \seqname{A}^\alpha,
  (E^\alpha,\catName{E}^\alpha),
  (C^\alpha, \catName{C}^\alpha)
  \bigr )
  \in \full{\Atl_\Ob}(\catName{E}, \catName{C})
$ and every $(U, V, \phi) \in \seqname{A}^\alpha$,
\nobreaks
  {{
    $
      \pi_1 \morphCompose \phi =
      \pi(\seqname{A}^\alpha) \restrictto_U
    $
  }}

\item
$
  \seqname{B} \defeq
  \set
  {
    \seqname{B}^\alpha \defeq
      \bigl (
        E^\alpha,
        X^\alpha,
        Y^\alpha,
        G(\seqname{C}^\alpha),
        \pi(\seqname{A}^\alpha),
        \rho(\seqname{C}^\alpha)
      \bigr )
  }%
  [
    {
      \Bigl (
        \seqname{A}^\alpha,
        (E^\alpha,\catName{E}^\alpha),
        \bigl (
          C^\alpha \defeq X^\alpha \times Y^\alpha,
          \catName{C}^\alpha
        \bigr )
      \Bigr ) \in \full{\Atl_\Ob}(\catName{E}, \catName{C})
    }
  ]*
$
\end{enumerate}

Then $\Functor^\Bun_{\M-G-\rho,\Bun}$ is a functor from
$
\Atl
  \Bigl (
    \Trivcat{\seqname{E}},
    \Trivcat[\BunProd-]{\seqname{B}}
  \Bigr )
$
to $\Atl[\Bun](\seqname{B})$.

\begin{proof}
Let
$
  o^i \defeq
  \bigl (
    \seqname{A}^i,
    \seqname{E}^i,
    \seqname{C}^i \defeq (X^i \times Y^i, \catName{XY}^i)
  \bigr )
  \objin \catName{M})
$,
$i \in [1,3]$,
$G^i \defeq G(o^i)$, \\
$\pi^i \defeq \pi(o^i)$,
$\rho^i \defeq \rho(o^i)$,
$\seqname{A}^i$ an m-atlas of $E^i$ in the coordinate model space $\seqname{C}^i$, \\
$
  \seqname{B}^i \defeq
    (
      E^i,
      X^i,
      Y^i,
      \pi^i,
      G^i,
      \rho^i
    )
$,
$
  o'^i \defeq
  (
    \seqname{A}^i,
    \seqname{B}^i
  )
$,
$
  \morphName{m}^i \defeq
    \bigl (
      \morphSeqName{f}^i,
      o^i,
      o^{i+1}
    \bigr )
    \arin \catName{M}
$, $i=1,2$,
an $E^i$-$E^{i+1}$ m-atlas morphism of $\seqname{A}^i$ to
$\seqname{A}^{i+1}$ in the coordinate spaces
$\seqname{C}^i, \seqname{C}^{i+1}$ that preserves the group action,
$\morphName{f}_G \maps G^i \to G^{i+1}$ the unique function asserted to
exist in \cref{eq:Grho} (preservation of group action) of
\pagecref{def:Grhomorph} and $\pi^i \maps E^i \morphEpi X^i$ the unique
function asserted in \cref{lem:mGrho}.

Let $\morphName{f}^i_G \maps G^i \to G^{i+1}$ be the function asserted by
\pagecref{eq:Grho}; let
$\morphName{f}^i_1 = \morphName{f}^i_X \times \morphName{f}^i_Y$ be the
decomposition given by \pagecref{def:Grhomodtriv}, and let
$
  \morphName{m}'^i \defeq
  \Functor^\Bun_{\M-G-\rho,\Bun} \morphName{m}^i = \\
    \bigl (
      \morphSeqName{f}^i,
      (
        \morphName{f}^i_0
        \morphName{f}^i_X
        \morphName{f}^i_Y
        \morphName{f}^i_G
      ),
      o'^i,
      o'^{i+1}
    \bigr )
$.

\begin{description}
\item[Preservation of endpoints:]
$E^i$, $X^i$ and $Y^i$ are topological spaces.

$ G^i$ is a topological group,

$\pi^i \maps E^i \morphEpi X^i$ is a continuous surjection.

$\rho^i \defeq \rho(\seqname{C}^i)$ is an effective action of
$G^i \defeq G(\seqname{C}^i)$ on $Y^\alpha$.

Hence $B^i$ is a protobundle.

\item[Identity:]
\begin{align*}
  \Functor^\Bun_{\M-G-\rho,\Bun}  \Id[{{(\seqname{A}^i, \seqname{B}^i)}}]
  & =
   \Id[{{(\seqname{A}^i, \seqname{B}^i)}}]
    \bigl (
      (\Id[E^i], \Id[X^i], \Id[Y^i], \Id[G^i]),
      o^i,
      o^i
    \bigr )
\\*
  & =
  \bigl (
    (\Id[{\Triv[X^i-\pi^i]{E^i}}], \Id[C^i]),
    o'^i,
    o'^i
  \bigr )
\\*
  & =
  \Bigl (
    (\Id[{\Triv[X^i-\pi^i]{E^i}}], \Id[C^i]),
    \Functor^\Bun_{\M-G-\rho,\Bun} o^i,
    \Functor^\Bun_{\M-G-\rho,\Bun} o^i
  \Bigr )
\\*
  & =
  \Id_{\Functor^\Bun_{\M-G-\rho,\Bun} o^i}
\end{align*}

\item[Composition:]
\end{description}
$\Functor^\Bun_{\M-G-\rho,\Bun} \morphName{m}^i$ is a morphism from
$\Functor^\Bun_{\M-G-\rho,\Bun} o^i$ to
$\Functor^\Bun_{\M-G-\rho,\Bun} o^{i+1}$:

\begin{align*}
\Functor^\Bun_{\M-G-\rho,\Bun} (\morphName{m}^i)
  & =
\Functor^\Bun_{\M-G-\rho,\Bun}
  (
    \morphSeqName{f}^i,
    o^i,
    o^{i+1}
  )
\\
  & =
  \bigl (
    (
      \morphName{f}^i_0,
      \morphName{f}^i_X,
      \morphName{f}^i_Y,
      \morphName{f}^i_G
    ),
    (\seqname{A}^i, \seqname{B}^i),
    (\seqname{A}^{i+1}, \seqname{B}^{i+1})
  \bigr )
\\
  & =
  \Bigl (
    (
      \morphName{f}^i_0,
      \morphName{f}^i_X,
      \morphName{f}^i_Y,
      \morphName{f}^i_G
    ),
    \Functor^\Bun_{\M-G-\rho,\Bun} o^i,
    \Functor^\Bun_{\M-G-\rho,\Bun} o^{i+1}
  \Bigr )
\end{align*}

$\Functor^\Bun_{\M-G-\rho,\Bun}$ maps identity functions to identity
functions:

\begin{equation}
\begin{split}
  &
\Functor^\Bun_{\M-G-\rho,\Bun} \Id[o^i] =
\\
  &
\Functor^\Bun_{\M-G-\rho,\Bun}
  \bigl (
    (\Id[E^i], \Id[C^i]),
    o^i,
    o^i
  \bigr ) =
\\
  &
  \bigl (
    (\Id[E^i], \Id[X^i], \Id[Y^i], \Id[G^i]),
    (\seqname{A}^i, \seqname{B}^i),
    (\seqname{A}^i, \seqname{B}^i)
  \bigr ) =
\\
  &
  \bigl (
    (\Id[E^i], \Id[X^i], \Id[Y^i], \Id[G^i]),
    \Functor^\Bun_{\M-G-\rho,\Bun} o^i,
    \Functor^\Bun_{\M-G-\rho,\Bun} o^i
  \bigr ) =
\\
  &
\Id_{\Functor^\Bun_{\M-G-\rho,\Bun} o^i}
\end{split}
\end{equation}

$
  \Functor^\Bun_{\M-G-\rho,\Bun} m^2 \morphCompose[A] \Functor^\Bun_{\M-G-\rho,\Bun} m^1
  =
  \Functor^\Bun_{\M-G-\rho,\Bun} (m^2 \morphCompose[A] m^1)
$:

\begin{enumerate}
\item
$
  m^2 \morphCompose[A] m^1 =
  \bigl (
    (
      \morphName{f}_0^2 \morphCompose \morphName{f}_0^1,
      \morphName{f}_1^2 \morphCompose \morphName{f}_1^1
    ),
    o^1,
    o^3
  \bigr )
  $
\item
$
  \Functor^\Bun_{\M-G-\rho,\Bun}
    \bigl (
       o^i
    \bigr )
  =
  (\seqname{A}^i, \seqname{B}^i)
$
\item
$
  \Functor^\Bun_{\M-G-\rho,\Bun} (m^i)
    = \\
    \bigl (
      (
        \morphName{f}^i_0,
        \morphName{f}^i_X,
        \morphName{f}^i_Y,
        \morphName{f}^i_G
      ),
      (\seqname{A}^i, \seqname{B}^i),
      (\seqname{A}^{i+1}, \seqname{B}^{i+1})
    \bigr )
$
\item
$
  \Functor^\Bun_{\M-G-\rho,\Bun} m^2 \morphCompose[A] \Functor^\Bun_{\M-G-\rho,\Bun} m^1
  = \\
    \bigl (
      (
        \morphName{f}^2_0 \morphCompose \morphName{f}^1_0,
        \morphName{f}^2_X \morphCompose \morphName{f}^1_X,
        \morphName{f}^2_Y \morphCompose \morphName{f}^1_Y,
        \morphName{f}^2_G \morphCompose \morphName{f}^1_G
      ),
      (\seqname{A}^1, \seqname{B}^1),
      (\seqname{A}^3, \seqname{B}^3)
    \bigr )
$
\item
$
  \Functor^\Bun_{\M-G-\rho,\Bun} (m^2 \morphCompose m^1)
  = \\
    \bigl (
      (
        \morphName{f}^2_0 \morphCompose \morphName{f}^1_0,
        \morphName{f}^2_X \morphCompose \morphName{f}^1_X,
        \morphName{f}^2_Y \morphCompose \morphName{f}^1_Y,
        \morphName{f}^2_G \morphCompose \morphName{f}^1_G
      ),
      (\seqname{A}^1, \seqname{B}^1),
      (\seqname{A}^3, \seqname{B}^3)
    \bigr )
$
\end{enumerate}
\end{proof}

$\Functor^\Bun_{\Bun,\M} \morphCompose \Functor^\Bun_{\M-G-\rho,\Bun} = \Id$.
\begin{proof}
Expanding the definitions, we have
\begin{enumerate}
\item
$
  \Functor^\Bun_{\M-G-\rho,\Bun}
  (\seqname{A}^i, \seqname{E}^i, \seqname{C}^i) =
  (\seqname{A}^i, \seqname{B}^i)
$
\item
$
  \Functor^\Bun_{\Bun,\M} (\seqname{A}^i, \seqname{B}^i) =
  (\seqname{A}^i, \Triv[X^i-\pi^i]{E^i}, \Triv[G^i-\rho^i-]{X^i,Y^i})
  (\seqname{A}^i, \seqname{E}^i, \seqname{C}^i)
$,
since by \cref{def:FuncMtoBun}, $\seqname{E}^i$ and
$\seqname{C}^i$ are trivial.
\item
$
  \Functor^\Bun_{\M-G-\rho,\Bun}
    \bigl (
      (
        \morphSeqName{f}^1_E,
        \morphSeqName{f}^1_X \times \morphSeqName{f}^1_Y
      ),
      (\seqname{A}^1, \seqname{E^1}, \seqname{C}^1),
      (\seqname{A}^2, \seqname{E^2}, \seqname{C}^2)
    \bigr ) = \\
    \bigl (
      (
        \morphName{f}^1_E,
        \morphName{f}^1_X,
        \morphName{f}^1_Y,
        \morphName{f}^1_G
      ),
      (\seqname{A}^1, \seqname{B}^1),
      (\seqname{A}^2, \seqname{B}^2)
    \bigr )
$
\item
$
  \Functor^\Bun_{\Bun,\M}
    \bigl (
      (
        \morphName{f}^1_E,
        \morphName{f}^1_X,
        \morphName{f}^1_Y,
        \morphName{f}^1_G
      ),
      (\seqname{A}^1, \seqname{B}^1),
      (\seqname{A}^2, \seqname{B}^2)
    \bigr ) = \\
    \bigl (
      (
        \morphSeqName{f}^1_E,
        \morphSeqName{f}^1_X \times \morphSeqName{f}^1_Y
      ),
      (\seqname{A}^1, \Triv[X^1-\pi^1-]{E^1}, \Triv[G^1-\rho^1-]{C^1}),
      (\seqname{A}^2, \Triv[X^2-\pi^2-]{E^2}, \Triv[G^2-\rho^2-]{C^2})
    \bigr ) = \\
    \bigl (
      (
        \morphSeqName{f}^1_E,
        \morphSeqName{f}^1_X \times \morphSeqName{f}^1_Y
      ),
      (\seqname{A}^1, \seqname{E}^1, \seqname{C}^1),
      (\seqname{A}^2, \seqname{E}^2, \seqname{C}^2)
    \bigr )
$,
since by \cref{def:FuncMtoBun}, $\seqname{E}^i$ and $\seqname{C}^i$ are trivial.
\end{enumerate}
\end{proof}
\end{theorem}

\section{Fiber bundles}
\label{sec:fib}
Conventionally a fiber bundle is different from its atlases, but
\pagecref{def:Bun-Atlcat} encourages treating them on an equal footing. All
of the results for maximal bundle atlases carry directly over to results
for fiber bundles.

\begin{definition}[fiber bundles]
\label{def:Bun}
$(E, X, Y, \pi, G, \rho, \seqname{A})$ is a weak fiber bundle iff

\begin{enumerate}
\item
$E$, $X$ and $Y$ are topological spaces

\item
$\pi \maps E \morphEpi X$ is surjective,

\item
$G$ is a topological group,

\item
$\rho \maps Y \times G \to Y$ is an effective right action of $G$ on $Y$

\item
$\seqname{A}$ is a $\pi$-$G$-$\rho$-bundle atlas of $E$ in the
coordinate space $X \times Y$.
\end{enumerate}

$(E, X, Y, \pi, G, \rho, \seqname{A})$ is also a proper fiber buundle,
or just fiber bundle, iff $\seqname{A}$ is a maximal
$\pi$-$G$-$\rho$-bundle atlas of $E$ in the coordinate space $X \times
Y$.

Let \\
\begin{enumerate}
\item
$E^\alpha, X^\alpha, Y^\alpha$ be topological spaces

\item
$G^\alpha$ be a topological group,

\item
$\pi^\alpha \maps E^\alpha \morphEpi X^\alpha$ be surjective

\item
$\rho^\alpha \maps Y^\alpha \times G^\alpha \to Y^\alpha$ be an
effective right action of $G^\alpha$ on $Y^\alpha$

\item
$
  \seqname{B}
  \defeq
  \set
    {{
      \seqname{B}^\alpha \defeq
      (
        E^\alpha,
        X^\alpha,
        Y^\alpha,
        \pi^\alpha,
        G^\alpha,
        \rho^\alpha
      )
    }}%
    [
      \alpha \prec \Alpha
    ]
$
\end{enumerate}

Then
\begin{multline}
\maximal[*]{\Bun_\Ob} \seqname{B} \defeq
\maximal[*]{\Atl[\Bun]_\Ob} \seqname{B}
\end{multline}
\end{definition}

\begin{lemma}[fiber bundles]
\label{lem:Bun}
$(E, X, Y, \pi, G, \rho, \seqname{A})$ is a weak fiber bundle iff

\begin{enumerate}
\item
$E,X,Y$ are topological spaces

\item
$G$ is a topological group,

\item
$\pi \maps E \morphEpi X$ is surjective

\item
$\rho \maps Y \times G \to Y$ an effective right action of $G$ on $Y$

\item
$\seqname{A}$ is an m-atlas of $\Triv{E}$ in the coordinate space
$\Triv[G-\rho-]{X,Y}$

\item
Every coordinate function in $\seqname{A}$ preserves fibers.
\end{enumerate}

\begin{proof}
The result follows from \cref{def:Bun} and \pagecref{lem:Bun-Atl}.
\end{proof}

$(E, X, Y, \pi, G, \rho, \seqname{A})$ is also a proper fiber bundle iff
$\seqname{A}$ is a maximal m-atlas of $\Triv{E}$ in the coordinate space
$\Triv[G-\rho-]{X,Y}$.

\begin{proof}
The result follows from \pagecref{lem:Bun-ATLmax}.
\end{proof}
\end{lemma}

\section{Bundle maps, categories of fiber bundles and functors}
\label{sec:Bunmap}

\subsection{Bundle maps}
\label{sub:Bunmap}

\subsubsection{Definition of Bundle maps}
\begin{definition}[Bundle maps]
\label{def:Bun-map}
Let \\
\begin{enumerate}
\item
$E^\alpha, X^\alpha, Y^\alpha$ be topological spaces

\item
$G^\alpha$ be a topological group,

\item
$\pi^\alpha \maps E^\alpha \morphEpi X^\alpha$ be surjective

\item
$\rho^\alpha \maps Y^\alpha \times G^\alpha \to Y^\alpha$ be an
effective right action of $G^\alpha$ on $Y^\alpha$

\item
$
  \seqname{B}
  \defeq
  \set
    {{
      \seqname{B}^\alpha \defeq
      (
        E^\alpha,
        X^\alpha,
        Y^\alpha,
        \pi^\alpha,
        G^\alpha,
        \rho^\alpha
      )
    }}%
    [
      \alpha \prec \Alpha
    ]
$
\end{enumerate}

Then
\begin{multline}
\maximal[*]{\Bun_\Ar} \seqname{B} \defeq \maximal[*]{\Atl[\Bun]_\Ar} \seqname{B}
\end{multline}
\begin{multline}
\maximal[*]{\Bun} \seqname{B}
\defeq
  \bigl (
    \maximal[*]{\Bun_\Ob} \seqname{B},
    \maximal[*]{\Bun_\Ar} \seqname{B},
    \morphCompose[A]
  \bigr )
\end{multline}

Let $(\seqname{A}^i,\seqname{B}^i) \in \Bun_\Ob \seqname{B}$, $i=1,2$.
Then  \\*
$
  \morphSeqName{f}
  \defeq
  (
    \morphName{f}_E \maps E^1 \to E^2,
    \morphName{f}_X \maps X^1 \to X^2,
    \morphName{f}_Y \maps Y^1 \to Y^2,
    \morphName{f}_G \maps G^1 \to G^2
  )
$
is a bundle map from $(\seqname{A}^1,\seqname{B}^1)$ to
$(\seqname{A}^2,\seqname{B}^2)$ iff it is a
$\seqname{B^1}$-$\seqname{B^2}$ bundle-atlas morphism from $A^1$ to
$A^2$. The identity morphism for $(\seqname{A}^i,\seqname{B}^i)$ is
\begin{equation}
\Id[{{(\seqname{A}^i,\seqname{B}^i)}}] \defeq
\bigl (
  (\Id[E^i], \Id[X^i], \Id[Y^i], \Id[G^i]),
  (\seqname{A}^i,\seqname{B}^i),
  (\seqname{A}^i,\seqname{B}^i)
\bigr )
\end{equation}
\end{definition}

\subsubsection{Proclamations on Bundle maps}
\begin{lemma}[Bundle maps]
\label{lem:Bun-map}
Let
$
  \seqname{B}^i
  \defeq
  (
    E^i,
    X^i,
    Y^i,
    \pi^i,
    G^i,
    \rho^i
  )
$, $i=1,2$, be a protobundle and let
$\seqname{A}^i$ be a $\pi^i$-$G^i$-$\rho^i$-atlas of $E^i$ in the
coordinate space
\nobreaks
  {{
    $C^i \defeq X^i \times Y^i$.
  }}
Then
$
  \morphSeqName{f}
  \defeq
  (
    \morphName{f}_E \maps E^1 \to E^2,
    \morphName{f}_X \maps X^1 \to X^2,
    \morphName{f}_Y \maps Y^1 \to Y^2,
    \morphName{f}_G \maps G^1 \to G^2
  )
$
is a bundle map from $(\seqname{A}^1,\seqname{B}^1)$ to
$(\seqname{A}^2,\seqname{B}^2)$ iff
$\morphName{f}_X \times \morphName{f}_Y$ is a
$G^1$-$G^2$-$\rho^1$-$\rho^2$ morphism from
$\Triv[G^1-\rho^1-]{X^1,Y^1}$ to $\Triv[G^2-\rho^2-]{X^2,Y^2}$ and
$(\morphName{f}_E,c \morphName{f}_X \times \morphName{f}_Y)$ is a
$\Triv[X^1-\pi^1-]{E^1}$-$\Triv[X^2-\pi^2-]{E^2}$ m-atlas morphism from
$\seqname{A}^1$ to $\seqname{A}^2$ in the coordinate spaces
$\Triv[G^1-\rho^1-]{X^1,Y^1}$, $\Triv[G^2-\rho^2-]{X^2,Y^2}$.

\begin{proof}
The result follows from \cref{def:Bun-map} above and
{
  \showlabelsinline
  \pagecref{lem:Bun-ATLmorph}.
}
\end{proof}
\end{lemma}

\subsection{Categories of fiber bundles and functors}
{
  \showlabelsinline
  \label{sub:Fib-cat}
}
\begin{theorem}[Categories of fiber bundles]
\label{the:Bun}
Let \\
$
  \seqname{B}
  \defeq
  \set
    {{
      \seqname{B}^\alpha \defeq
      (
        E^\alpha,
        X^\alpha,
        Y^\alpha,
        G^\alpha,
        \pi^\alpha,
        \rho^\alpha
      )
    }}%
    [
      \alpha \prec \Alpha
    ]
$,
where $E^\alpha, X^\alpha, Y^\alpha$ are topological spaces,
$G^\alpha$ a topological group,
$\pi^\alpha \maps E^\alpha \morphEpi X^\alpha$ surjective and
$\rho^\alpha \maps Y^\alpha \times G^\alpha \to Y^\alpha$
an effective right action of $G^\alpha$ on $Y^\alpha$. Then
$\Bun \seqname{B}$ is a category and $\Id[{B^\alpha}]$ is the identity
morphism for $B^\alpha$.

\begin{proof}
{
  \showlabelsinline
  The result follows directly from \fullcref{def:Bun},
}
\fullcref{def:Bun-map} and \pagecref{lem:Bun-ATLiscat}.
\end{proof}
\end{theorem}

\subsubsection{Functors for fiber bundles}

\begin{definition}[Functor from fiber bundles to Local Coordinate Spaces]
{
  \showlabelsinline
  \label{def:BunLCS}
}
Let
\begin{enumerate}
\item
$\seqname{E}$, $\seqname{X}$ and $\seqname{Y}$ be sets of
topological spaces

\item
$\seqname{G}$ be a set of topological groups

\item
$
  \seqname{B}^\alpha \defeq
  (
    E^\alpha \in \seqname{E},
    X^\alpha \in \seqname{X},
    Y^\alpha \in \seqname{Y},
    G^\alpha \in \seqname{Y},
    \pi^\alpha,
    \rho^\alpha
  )
$,
$\alpha \prec \Alpha$, be a protobundle

\item
$\seqname{B} \defeq \set{{\seqname{B}^\alpha}}[\alpha \prec \Alpha]$
be a set of protobundles

\item
$C^\alpha \defeq X^\alpha \times Y^\alpha$,

\item
$\catName{E}$, $\catName{X}$, $\catName{Y}$ and $\catName{G}$ be model
categories

\item
$\catSeqName{XYG\rho}$ be a $G$-$\rho$-model category

\item
$
  \catSeqName{M} \defeq
   \Bigl  (
      \catName{E},
      \catSeqName{XYG\rho},
      \catName{X},
      \catName{Y},
      \catName{G}
   \Bigr  )
$,

\item
$
  \Triv{\catSeqName{M}} \defeq
    \Bigl  (
      \Trivcat{\seqname{E}},
      \Trivcat[\Bun-]{\seqname{B}},
      \Trivcat{\seqname{X}},
      \Trivcat{\seqname{Y}},
      \Trivcat{\seqname{G}}
    \Bigr  )
$

\item
$\Triv{\catSeqName{M}} \SUBCAT[full-] \catSeqName{M}$
\item
$
  \seqname{B}^i \defeq
  (
    E^i,
    X^i,
    Y^i,
    \pi^i,
    G^i,
    \rho^i
  )
  \in \seqname{B}
$,
$i=1,2$.
with
\begin{enumerate}
\item
$\star^i \maps G^i \times G^i \morphEpi G^i$ the group operation for $G^i$

\item
$\seqname{A}^i$ a maximal $\pi^i$-$G^i$-$\rho^i$-bundle atlas of $E^i$
in the coordinate space
\nobreaks
  {{
    $
      C^i \defeq X^i \times Y^i
    $
  }}

\item
$\pi^i_X \defeq \pi_1 \maps X^i \times Y^i \morphEpi X^i$

\item
$\pi^i_Y \defeq \pi_2 \maps X^i \times Y^i \morphEpi Y^i$
\item
$o^i \defeq (\seqname{A}^i, \seqname{B}^i)$

\item
$
  \seqname{F}^i \defeq
  (
    \pi^i,
    \pi^i_X,
    \pi^i_Y,
    \star^i,
    \rho^i
  )
$

\item
$
  \Singcat{\Triv{\catSeqName{M}^i}} \defeq
    \Bigl  (
      \singcat{\Triv[X^i-\pi^i]{E^i}},
      \singcat{\Triv[G^i-\rho^i-]{X^i,Y^i}},
      \singcat{X^i},
      \singcat{Y^i},
      \singcat{G^i},
    \Bigr  )
$

\item
$
  \seqname{M}^i \defeq
    \Bigl (
      \Triv[X^i-\pi^i]{E^i},
      \Triv[G^i-\rho^i-]{X^i,Y^i},
      X^i,
      Y^i,
      G^i
    \Bigr )
$

\item
$
  \Sigma \defeq
  \bigl (
      (0,2),
      (1,2),
      (1,3),
      (4,4,4),
      (3,4,3)
  \bigr )
$

\item
$
  \seqname{L}^i \defeq
  \left (
    \Singcat{\Triv{\catSeqName{M}^i}},
    \seqname{M}^i,
    \seqname{A}^i,
    \seqname{F}^i,
    \Sigma
  \right )
$

\item
$
  \seqname{L}^{i,\catSeqName{M}} \defeq
  (
    \catSeqName{M},
    \seqname{M}^i,
    \seqname{A}^i,
    \seqname{F}^i,
    \Sigma
  )
$

\item
$
   \morphSeqName{f} \defeq
     (
       \morphName{f}_E \maps E^1 \to E^2,
       \morphName{f}_X \maps X^1 \to X^2,
       \morphName{f}_Y \maps Y^1 \to Y^2,
       \morphName{f}_G \maps G^1 \to G^2
     )
$
a $\seqname{B^1}$-$\seqname{B^2}$ bundle-atlas morphism from
$\seqname{A}^1$ to $\seqname{A}^2$
\end{enumerate}
\end{enumerate}

Then
\begin{multline}
\Functor^\Bun_{\Fib,\LCS}
  (\seqname{A}^i, \seqname{B}^i) \defeq
\seqname{L}^i
\end{multline}
\begin{multline}
\Functor^\Bun_{\Fib,\LCS}
  \bigl (
    \morphSeqName{f},
    (
      \seqname{A}^1,
      \seqname{B}^1
    ),
    (
      \seqname{A}^2,
      \seqname{B}^2
    )
  \bigr )
\defeq
\Bigl (
  \bigl (
    \morphName{f}_E,
    \morphName{f}_X \times \morphName{f}_Y,
    \morphName{f}_X,
    \morphName{f}_Y,
    \morphName{f}_G
  \bigr ),
  \seqname{L}^1,
  \seqname{L}^2
\Bigr )
\end{multline}

\begin{multline}
\Functor^{\Bun,\catSeqName{M}}_{\Fib,\LCS}
  (\seqname{A}^i, \seqname{B}^i) \defeq
\seqname{L}^{i,\catSeqName{M}}
\end{multline}
\begin{multline}
\Functor^{\Bun,\catSeqName{M}}_{\Fib,\LCS}
  \bigl (
    \morphSeqName{f},
    (
      \seqname{A}^1,
      \seqname{B}^1
    ),
    (
      \seqname{A}^2,
      \seqname{B}^2
    )
  \bigr )
\defeq \\
\Bigl (
  \bigl (
    \morphName{f}_E,
    \morphName{f}_X \times \morphName{f}_Y,
    \morphName{f}_X,
    \morphName{f}_Y,
    \morphName{f}_G
  \bigr ),
  \seqname{L}^{1,\catSeqName{M}},
  \seqname{L}^{2,\catSeqName{M}}
\Bigr )
\end{multline}

\begin{multline}
\LCS^\Bun_\Ob \seqname{B}
\defeq
\set
  {
    \Functor^\Bun_{\Fib,\LCS}
      \bigl (
        \seqname{A},
        (
          E^\alpha,
          X^\alpha,
          Y^\alpha,
          G^\alpha,
          \pi^\alpha,
          \rho^\alpha
        )
      \bigr )
  }%
  [
    {
      (
        E^\alpha,
        X^\alpha,
        Y^\alpha,
        G^\alpha,
        \pi^\alpha,
        \rho^\alpha
      )
      \in
      \seqname{B}
    },
    {
      \maximal{\isAtl[\Bun]_\Ob}
        (
          \seqname{A},
          E^\alpha,
          X^\alpha,
          Y^\alpha,
          G^\alpha,
          \pi^\alpha,
          \rho^\alpha
        )
    }
  ]*
\end{multline}
\begin{multline}
\LCS^\Bun_\Ar \seqname{B}
\defeq
\set
  {
    \Functor^\Bun_{\Fib,\LCS}
      \bigl (
        \morphSeqName{f},
        (\seqname{A}^1, \seqname{B}^1),
        (\seqname{A}^2, \seqname{B}^2)
      \bigr )
  }%
  [
    {\seqname{B}^i \in \seqname{B}},
    {
      \maximal{\isAtl[\Bun]_\Ar}
        (
          \seqname{A}^1,
          \seqname{B}^1,
          \seqname{A}^2,
          \seqname{B}^2,
          \morphSeqName{f}
        )
    }
  ]*
\end{multline}
\begin{multline}
\LCS^\Bun \seqname{B}
\defeq
  \bigl (
    \LCS^\Bun_\Ob \seqname{B},
    \LCS^\Bun_\Ar \seqname{B},
    \morphCompose[A]
  \bigr )
\end{multline}

\begin{multline}
\LCS^{\Bun,\catSeqName{M}}_\Ob \seqname{B}
\defeq
\set
  {
    \Functor^{\Bun,\catSeqName{M}}_{\Fib,\LCS}
      \bigl (
        \seqname{A},
        (
          E^\alpha,
          X^\alpha,
          Y^\alpha,
          G^\alpha,
          \pi^\alpha,
          \rho^\alpha
        )
      \bigr )
  }%
  [
    {
      (
        E^\alpha,
        X^\alpha,
        Y^\alpha,
        G^\alpha,
        \pi^\alpha,
        \rho^\alpha
      )
      \in
      \seqname{B}
    },
    {
      \maximal{\isAtl[\Bun]_\Ob}
        (
          \seqname{A},
          E^\alpha,
          X^\alpha,
          Y^\alpha,
          \pi^\alpha,
          G^\alpha,
          \rho^\alpha
        )
    }
  ]*
\end{multline}
\begin{multline}
\LCS^{\Bun,\catSeqName{M}}_\Ar \seqname{B}
\defeq
\set
  {
    \Functor^{\Bun,\catSeqName{M}}_{\Fib,\LCS}
      \bigl (
        \morphSeqName{f},
        (\seqname{A}^1, \seqname{B}^1),
        (\seqname{A}^2, \seqname{B}^2)
      \bigr )
  }%
  [
    {\seqname{B}^i \in \seqname{B}},
    {
      \maximal{\isAtl[\Bun]_\Ar}
        (
          \seqname{A}^1,
          \seqname{B}^1,
          \seqname{A}^2,
          \seqname{B}^2,
          \morphSeqName{f}
        )
    }
  ]*
\end{multline}
\begin{multline}
\LCS^{\Bun,\catSeqName{M}} \seqname{B}
\defeq
  \bigl (
    \LCS^{\Bun,\catSeqName{M}}_\Ob \seqname{B},
    \LCS^{\Bun,\catSeqName{M}}_\Ar \seqname{B},
    \morphCompose[A]
  \bigr )
\end{multline}
\end{definition}

\begin{theorem}[Functor from fiber bundles to Local Coordinate Spaces]
\label{the:BunLCS}
Let
\begin{enumerate}
\item
$\seqname{E}$, $\seqname{X}$ and $\seqname{Y}$ be sets of
topological spaces, $\seqname{G}$ be a set of topological groups,
\item
$
  \seqname{B}^\alpha \defeq
  (
    E^\alpha \in \seqname{E},
    X^\alpha \in \seqname{X},
    Y^\alpha \in \seqname{Y},
    G^\alpha \in \seqname{Y},
    \pi^\alpha,
    \rho^\alpha
  )
$
$\alpha \prec \Alpha$, be a protobundle with group operation
$\star^\alpha \maps G^\alpha \times G^\alpha \morphEpi G^\alpha$
\item
$\seqname{B} \defeq \set{{\seqname{B}^\alpha}}[\alpha \prec \Alpha]$
be a set of protobundles
\item
$C^\alpha \defeq X^\alpha \times Y^\alpha$,
$\catName{E}$, $\catName{X}$, $\catName{Y}$ and $\catName{G}$ be model
categories
\item
$\catSeqName{XYG\rho}$ be a $G$-$\rho$-model category,
\item
$
  \catSeqName{M} \defeq
   \Bigl  (
      \catName{E},
      \catSeqName{XYG\rho},
      \catName{X},
      \catName{Y},
      \catName{G}
   \Bigr  )
$,
$
  \Triv{\catSeqName{M}} \defeq
    \Bigl  (
      \Trivcat{\seqname{E}},
      \Trivcat[\Bun-]{\seqname{B}},
      \Trivcat{\seqname{X}},
      \Trivcat{\seqname{Y}},
      \Trivcat{\seqname{G}}
    \Bigr  )
$
\item
$\Triv{\catSeqName{M}} \SUBCAT[full-] \catSeqName{M}$,
$
  \seqname{B}^i \defeq
  (
    E^i,
    X^i,
    Y^i,
    \pi^i,
    G^i,
    \rho^i
  )
$
$i=1,2,3$, be a protobundle in $\seqname{B}$
with group operation $\star^i \maps G^i \times G^i \morphEpi G^i$
\item
$\seqname{A}^i$ be a maximal $\pi^i$-$G^i$-$\rho^i$-bundle atlas of
$E^i$ in the coordinate space \\*
$C^i \defeq X^i \times Y^i$, i.e.,
$\maximal{\isAtl[\Bun]_\Ob} (\seqname{A}^i, \seqname{B}^i)$
\item
$o^i \defeq (\seqname{A}^i, \seqname{B}^i)$ be the object corresponding
to the fiber bundle \\*
$
  \seqname{B}^i \defeq
  (
    E^i,
    X^i,
    Y^i,
    \pi^i,
    G^i,
    \rho^i,
    \seqname{A}^i
  )
$
\item
$\pi^i_X \defeq \pi_1 \maps X^i \times Y^i \morphEpi X^i$,
$\pi^i_Y \defeq \pi_2 \maps X^i \times Y^i \morphEpi Y^i$
\item
$
  \seqname{F}^i \defeq
  (
    \pi^i,
    \pi^i_X,
    \pi^i_Y,
    \star^i,
    \rho^i
  )
$
\item
$
  \Singcat{\Triv{\catSeqName{M}^i}} \defeq
    \Bigl  (
      \singcat{\Triv[X^i-\pi^i]{E^i}},
      \singcat{\Triv[G^i-\rho^i-]{X^i,Y^i}},
      \singcat{X^i},
      \singcat{Y^i},
      \singcat{G^i},
    \Bigr  )
$,
$
  \seqname{M}^i \defeq
    \Bigl (
      \Triv[X^i-\pi^i]{E^i},
      \Triv[G^i-\rho^i-]{X^i,Y^i},
      X^i,
      Y^i,
      G^i
    \Bigr )
$
\item
$
  \Sigma \defeq
  \bigl (
      (0,2),
      (1,2),
      (1,3),
      (4,4,4),
      (3,4,3)
  \bigr )
$
be a signature
\item
$
  \seqname{L}^i \defeq
  \Biggl (
    \Singcat{\Triv{\catSeqName{M}^i}},
    \seqname{M}^i,
    \seqname{A}^i,
    \seqname{F}^i,
    \Sigma
  \Biggr )
$,
$
  \seqname{L}^{i,\catSeqName{M}} \defeq
  (
    \catSeqName{M},
    \seqname{M}^i,
    \seqname{A}^i,
    \seqname{F}^i,
    \Sigma
  )
$
be local coordinate spaces
\item
$
   \morphSeqName{f}^i \defeq
     (
       \morphName{f}^i_E \maps E^i \to E^{i+1},
       \morphName{f}^i_X \maps X^i \to X^{i+1},
       \morphName{f}^i_Y \maps Y^i \to Y^{i+1},
       \morphName{f}^i_G \maps G^i \to G^{i+1}
     )
$
and
$
   \morphSeqName{g}^i \defeq
     (
       \morphName{g}^i_E \maps E^i \to E^{i+1},
       \morphName{g}^i_X \maps X^i \to X^{i+1},
       \morphName{g}^i_Y \maps Y^i \to Y^{i+1},
       \morphName{g}^i_G \maps G^i \to G^{i+1}
     )
$,
$i=1,2$, be bundle maps from $(\seqname{A}^i,\seqname{B}^i)$ to
$(\seqname{A}^{i+1},\seqname{B}^{i+1})$
\item
$
  \morphSeqName{f}^{i,LCS} \defeq
    (
      \morphName{f}^i_E,
      \morphName{f}^i_X \times \morphName{f}^i_Y,
      \morphName{f}^i_X,
      \morphName{f}^i_Y,
      \morphName{f}^i_G
    )
$
and
$
  \morphSeqName{g}^{i,LCS} \defeq
    (
      \morphName{g}^i_E,
      \morphName{g}^i_X \times \morphName{f}^i_Y,
      \morphName{g}^i_X,
      \morphName{g}^i_Y,
      \morphName{g}^i_G
    )
$,
$i=1,2$
\end{enumerate}

Then:

$\seqname{L}^i$ is a local $\Singcat{\Triv{\catSeqName{M}^i}}$-$\Sigma$
coordinate space and $\seqname{L}^{i,\catSeqName{M}}$ is a local
$\catSeqName{M}$-$\Sigma$ coordinate space.

\begin{proof}
They satisfy the criteria in \pagecref{def:LCS}:
\begin{enumerate}
\item $\Singcat{\Triv{\catSeqName{M}^i}}$ and $\catSeqName{M}$ are
sequences of categories by construction.

\item
$
  \seqname{M}^i \seqin
  \Singcat{\Triv{\catSeqName{M}^i}} \SUBCAT[full-]
  \Triv{\catSeqName{M}}
$
by construction and $\Triv{\catSeqName{M}} \SUBCAT[full-]\catSeqName{M}$
by hypothesis.

\item $\singcat{\Triv[X^i-\pi^i]{E^i}}$ is a model category by
\pagecref{def:trivmod}, $\singcat{\Triv[G^i-\rho^i-]{X^i,Y^i}}$ is a
model category by \pagecref{def:Grhomodtriv}, $\catName{E}$ is a model
category by hypothesis and $\catSeqName{XYG\rho}$ is a model category by
\pagecref{def:Grhomodcat}.

\item
$
  \Biggl (
    \Singcat{\Triv{\catSeqName{M}^i}},
    \seqname{M}^i,
    \Sigma,
    \seqname{F}^i
  \Biggr )
$
and $(\catSeqName{M}, \seqname{M}^i, \Sigma, \seqname{F}^i)$
are prestructures:
$\seqname{F}^i$ has $\Singcat{\Triv{\catSeqName{M}^i}}$-signatures
$\Sigma$ and $\seqname{F}^i$ has $\catSeqName{M}$-signatures $\Sigma$.

\item $\seqname{A}^i$ is a maximal m-atlas of $\Triv[X^i-\pi^i]{E^i}$ in
$\Triv[G^i-\rho^i-]{X^i,Y^i}$ by \pagecref{lem:Bun-ATLmorph}.
\item There are no constraint functions.
\end{enumerate}
\end{proof}

$\LCS^\Bun \seqname{B}$ and $\LCS^{\Bun,\catSeqName{M}} \seqname{B}$ are
categories and the identity morphism for $\seqname{L}^i$ is \\*
$
  \Id[{\seqname{L}^i}] \defeq
  \Bigl (
    \bigl (
      \Id[E^i]
      \Id[X^i]
      \Id[{X^i \times Y^i}]
      \Id[Y^i]
      \Id[G^i]
    \bigr ),
    \seqname{L}^i,
    \seqname{L}^i
  \Bigr )
$.

\begin{proof}
%$\LCS^\Bun \seqname{B}$ and $\LCS^{\Bun,\catSeqName{M}}\seqname{B}$
satisfy the definition of a category:
\begin{description}
\item[Composition:]
$\morphSeqName{f}^2 \morphCompose \morphSeqName{f}^1$ is a bundle map from
$(\seqname{A}^1,\seqname{B}^1)$ to $(\seqname{A}^1,\seqname{B}^1)$ by
\pagecref{cor:Bun-ATLmorph}.
Then
\begin{align*}
  \Functor^\Bun_{\Fib,\LCS}
    \bigl (
      \morphSeqName{f}^2,
     o^2,
     o^3
    \bigr )
  \morphCompose
  \Functor^\Bun_{\Fib,\LCS}
    \bigl (
      \morphSeqName{f}^1,
     o^1,
     o^2
    \bigr )
  & =
  \Functor^\Bun_{\Fib,\LCS}
    \bigl (
      \morphSeqName{f}^2 \morphCompose \morphSeqName{f}^1,
     o^1,
     o^3
    \bigr )
\\*
  \Functor^{\Bun,\catSeqName{M}}_{\Fib,\LCS}
    \bigl (
      \morphSeqName{f}^2,
     o^2,
     o^3
    \bigr )
  \morphCompose
  \Functor^{\Bun,\catSeqName{M}}_{\Fib,\LCS}
    \bigl (
      \morphSeqName{f}^1,
     o^1,
     o^2
    \bigr )
  & =
  \Functor^{\Bun,\catSeqName{M}}_{\Fib,\LCS}
    \bigl (
      \morphSeqName{f}^2 \morphCompose \morphSeqName{f}^1,
     o^1,
     o^3
    \bigr )
\end{align*}

\item[Associativity:]
Composition is associative by \pagecref{lem:labcomp}.

\item[Unit:]
$\Id[{{\seqname{L}^i}}]$ is an identity morphism by \pagecref{lem:labcomp}.
\end{description}
\end{proof}

$
\Functor^\Bun_{\Fib,\LCS}
  \bigl (
    \morphSeqName{f}^i,
    (\seqname{A}^i, \seqname{B}^i),
    (\seqname{A}^{i+1}, \seqname{B}^{i+1})
  \bigr )
$
is a morphism from $\LCS^\Bun \seqname{B}^i$ to \\
$\LCS^\Bun \seqname{B}^{i+1}$ and
$
\Functor^{\Bun,\catSeqName{M}}_{\Fib,\LCS}
  \bigl (
    \morphSeqName{f}^i,
    (\seqname{A}^i, \seqname{B}^i),
    (\seqname{A}^{i+1}, \seqname{B}^{i+1})
  \bigr )
$
is a strict morphism from
$\LCS^{\Bun,\catSeqName{M}} \seqname{B}^i$ to
$\LCS^{\Bun,\catSeqName{M}} \seqname{B}^{i+1}$.

\begin{proof}
\leavevmode
\begin{enumerate}
\item Prestructure morphism:
\newline
$\morphSeqName{f}^1$ $\Sigma$-commutes with
$\morphSeqName{F}^1$, $\morphSeqName{F}^2$.
\item m-atlas morphism:
\newline
$(\morphName{f}_0, \morphName{f}_1)$ is an m-atlas morphism from
$(\Triv[X^1-\pi^1-]{E^1}, \Triv[G^1-\rho^1-]{C^1})$ to \\
$(\Triv[X^2-\pi^2-]{E^2}, \Triv[G^2-\rho^2-]{C^2})$ by
\pagecref{lem:Bun-ATLmorph}.
\end{enumerate}
\end{proof}

$\Functor^\Bun_{\Fib,\LCS}$ is a functor from $\Bun \seqname{B}$ to
$\LCS^\Bun \seqname{B}$ and
$\Functor^{\Bun,\catSeqName{M}}_{\Fib,\LCS}$ is a functor from
$\Bun \seqname{B}$ to $\LCS^{\Bun,\catSeqName{M}} \seqname{B}$.

\begin{proof}
$\Functor^\Bun_{\Fib,\LCS}$ and
$\Functor^{\Bun,\catSeqName{M}}_{\Fib,\LCS}$ satisfy the definition of a
functor:
\begin{description}
\item[Preservation of endpoints:]
\leavevmode
\begin{align*}
  \Functor^\Bun_{\Fib,\LCS}
    \bigl (
      \morphSeqName{f}^i
      (\seqname{A}^i, \seqname{B}^i),
      (\seqname{A}^{i+1}, \seqname{B}^{i+1})
    \bigr )
  & =
    \bigl (
      \morphSeqName{f}^{i,LCS},
      \seqname{L}^i,
      \seqname{L}^{i+1}
    \bigr )
\\*
  \Functor^{\Bun,\catSeqName{M}}_{\Fib,\LCS}
    \bigl (
      \morphSeqName{f}^i
      (\seqname{A}^i, \seqname{B}^i),
      (\seqname{A}^{i+1}, \seqname{B}^{i+1})
    \bigr )
  & =
    \bigl (
      \morphSeqName{f}^{i,LCS},
      \seqname{L}^{i,\catSeqName{M}},
      \seqname{L}^{i+1,\catSeqName{M}}
    \bigr )
\\*
  \Functor^\Bun_{\Fib,\LCS} (\seqname{A}^i, \seqname{B}^i)i
  & =
  \seqname{L}^i
\\*
  \Functor^{\Bun,\catSeqName{M}}_{\Fib,\LCS}(\seqname{A}^i, \seqname{B}^i)
  & =
  \seqname{L}^{i,\catSeqName{M}}
\end{align*}

\item[Composition:]
This follows from the proof above that $\LCS^\Bun \seqname{B}$ and
$\LCS^{\Bun,\catSeqName{M}} \seqname{B}$ are categories.

\item[Identity:]
\leavevmode
\begin{align*}
  \Functor^\Bun_{\Fib,\LCS} (\Id[{{(\seqname{A}^i, \seqname{B}^i)}}])
  & =
  \Functor^\Bun_{\Fib,\LCS}
    \Bigl (
      \bigl ( \Id[E^i], \Id[X^i], \Id[Y^i], \Id[G^i] \bigr ),
      (\seqname{A}^i, \seqname{B}^i),
      (\seqname{A}^i, \seqname{B}^i)
    \Bigr )
\\
  & =
    \Bigl (
      \bigl (
        \Id[E^i],
        \Id[X^i] \times \Id[Y^i],
        \Triv{X^i},
        \Triv{Y^i},
        \Triv{G^i}
      \bigr ),
      \seqname{L}^i,
      \seqname{L}^i
    \Bigr )
\\
  & =
  \Id[{\seqname{L}^i}]
\\
  & =
  \Id[{\Functor^\Bun_{\Fib,\LCS} (\seqname{A}^i, \seqname{B}^i)}]
\\
  & =
  \Functor^{\Bun,\catSeqName{M}}_{\Fib,\LCS}
    (\Id[{{(\seqname{A}^i, \seqname{B}^i)}}])
\\
  & =
  \Functor^{\Bun,\catSeqName{M}}_{\Fib,\LCS}
    \Bigl (
      \bigl ( \Id[E^i], \Id[X^i], \Id[Y^i], \Id[G^i] \bigr ),
      (\seqname{A}^i, \seqname{B}^i),
      (\seqname{A}^i, \seqname{B}^i)
    \Bigr ) =
\\
  & =
    \Bigl (
      \bigl (
        \Id[E^i],
        \Id[X^i] \times \Id[Y^i],
        \Triv{X^i},
        \Triv{Y^i},
        \Triv{G^i}
      \bigr ),
      \seqname{L}^{i,\catSeqName{M}},
      \seqname{L}^{i,\catSeqName{M}}
    \Bigr )
\\
  & =
  \Id[{\seqname{L}^{i,\catSeqName{M}}}]
\\
  & =
  \Id[{{\Functor^{\Bun,\catSeqName{M}}_{\Fib,\LCS}
    (\seqname{A}^i, \seqname{B}^i)}}]
\end{align*}
\end{description}
\end{proof}

$\morphSeqName{f}^{i,LCS}$ is $(0, 2, 4)$-equivalent to
$\morphSeqName{f}^{i,LCS}$ iff $\morphSeqName{f}^i$ is bundle equivalent
to $\morphSeqName{g}^i$.

\begin{proof}
Both statments are true iff $\morphName{f}^i_E = \morphName{g}^i_E$,
$\morphName{f}^i_X = \morphName{g}^i_X$ and
$\morphName{f}^i_G = \morphName{g}^i_G$
\end{proof}
\end{theorem}

\begin{definition}[Functor from local coordinate spaces to fiber bundles]
\label{def:LCStoBun}
Let $\catName{E}$ be a model category,
$\catSeqName{XYG\rho}$ a $G$-$\rho$-model category,
$\catName{X}$, $\catName{Y}$ and $\catName{G}$ topological categories,
$\catSeqName{M}$ a sequence of categories,
$\seqname{M}^i$, $i=1,2$, a sequence,
$\seqname{E}^i \objin \catName{E}$,
$\seqname{C}^i = (C^i,\catName{C}^i) \objin \catSeqName{XYG\rho}$,
$X^i \objin \catName{X}$,
$Y^i \objin \catName{Y}$,
$G^i \objin \catName{G}$,
$\pi^i \maps E^i \morphEpi X^i$,
$\pi^i_X \maps C^i \morphEpi X^i$,
$\pi^i_Y \maps C^i \morphEpi Y^i$,
\nobreaks
  {{
    $
      \star^i \maps G^i \times G^i \morphEpi G^i
    $
  }}
$\rho^i \maps Y^i \times G^i \morphEpi Y^i$,
$\Singcat{\Triv{\catSeqName{M}^i}}$ a sequence of categories,
$
  \seqname{L}^i \defeq
  \Biggl (
    \Singcat{\Triv{\catSeqName{M}^i}},
    \seqname{M}^i,
    \seqname{A}^i,
    \seqname{F}^i,
    \Sigma
  \Biggr )
$
and
$
  \seqname{L}^{i,\catSeqName{M}} \defeq
  (
    \catSeqName{M}^i,
    \seqname{M}^i,
    \seqname{A}^i,
    \seqname{F}^i,
    \Sigma
  )
$
local coordinate spaces,
$\seqname{B}^i \defeq (E^i, X^i, Y^i, G^i, \pi^i, \rho^i)$
and
$\morphSeqName{f} \defeq (\morphName{f}_\beta, \beta \prec \length(\seqname{M}^1)$
a morphism from $\seqname{L}^1$ to $\seqname{L}^2$,
satisfying
\begin{enumerate}
\item
$
  \head(\catSeqName{M},5) =
  \Bigl  (
    \catName{E},
    \catSeqName{XYG\rho},
    \catName{X},
    \catName{Y},
    \catName{G}
  \Bigr  )
$
\item
$
  \head \Biggl ( \Singcat{\Triv{\catSeqName{M}^i}},5 \Biggr )  =
  \Biggl  (
    \singcat{\Triv[X^i-\pi^i]{E^i}},
    \singcat{\Triv[G^i-\rho^i-]{X^i,Y^i}},
    \singcat{X^i},
    \singcat{Y^i},
    \singcat{G^i}
  \Biggr  )
$
\item
$
  \head(\seqname{M}^i,5) =
  (
    \seqname{E}^i,
    \seqname{C}^i,
    X^i,
    Y^i,
    G^i
  )
$
\item $\seqname{M}^i \seqin \catSeqName{M}$
\item
$
  \head(\seqname{F}^i,5) =
  (
    \pi^i,
    \pi^i_X,
    \pi^i_Y,
    \star^i,
    \rho^i
  )
$
\item $C^i = X^i \times Y^i$
\item $G^i$ is a topological group with group operation $\star^i$.
Subsequent references to $G^i$ should be read as referring to
the group rather than just the underlying topological space.
\item $\pi^i$ is surjective.
\item $\pi^i_X = \pi_1 \maps X^i \times Y^i \morphEpi X^i$.
\item $\pi^i_Y = \pi_2 \maps X^i \times Y^i \morphEpi Y^i$.
\item $\rho^i$ is an effective right action of $G^i$ on $Y^i$.
\item
$
  \head(\Sigma,5) =
  \bigl ( (0,2), (1,2), (1,3), (4,4,4), (3,4,3) \bigr )
$
\end{enumerate}
\begin{remark}
Due to the commutation requirement, specifying
$\catName{M}_4 = \Trivcat[\mathrm{group-}]{\seqname{G}}$ is not
necessary in order to ensure that $\morphName{f}_G$ is a homomorphism.
\end{remark}
Then
\begin{equation}
\Functor^\Bun_{\LCS,\Fib} L^i \defeq
  (E^i, X^i, Y^i, \pi^i, G^i, \rho^i, \seqname{A}^i)
\end{equation}
\begin{equation}
\Functor^\Bun_{\LCS,\Fib}
  \bigl (
    \morphSeqName{f},
    (\seqname{A}^1, \seqname{B}^1),
    (\seqname{A}^2, \seqname{B}^2)
  \bigr )
\defeq
  \bigl (
    (\morphName{f}_0, \morphName{f}_2, \morphName{f}_3, \morphName{f}_4),
    \Functor^\Bun_{\LCS,\Fib} L^1,
    \Functor^\Bun_{\LCS,\Fib} L^2
  \bigr )
\end{equation}
\end{definition}

\begin{theorem}[Functor from local coordinate spaces to fiber bundles]
\label{the:LCStoBun}
Let
\begin{enumerate}
\item
$\catName{E}$ be a model category

\item
$\catSeqName{XYG\rho}$ a $G$-$\rho$-model category

\item
$\catName{X}$, $\catName{Y}$ and $\catName{G}$ topological categories

\item
$\catSeqName{M}$ a sequence of categories

\item
$
  \seqname{E}^\alpha \defeq (E^\alpha,\catName{E}^\alpha)
  \objin
  \catName{E}
$,
$\alpha \prec \Alpha$

\item
$
  \seqname{C}^\alpha \defeq (C^\alpha,\catName{C}^\alpha)
  \objin
  \catSeqName{XYG\rho}
$

\item
$X^\alpha \objin \catName{X}$

\item
$Y^\alpha \objin \catName{Y}$

\item
$G^\alpha \objin \catName{G}$ a topological group with group operation
$\star^\alpha$

\item
$\pi^\alpha \maps E^\alpha \morphEpi X^\alpha$ surjective

\item
$\pi^\alpha_X \maps C^\alpha \morphEpi X^\alpha$,

\item
$\pi^\alpha_Y \maps C^\alpha \morphEpi Y^\alpha$ surjective

\item
$\star^\alpha \maps G^\alpha \times G^\alpha \morphEpi G^\alpha$ surjective

\item
$\rho^\alpha \maps Y^\alpha \times G^\alpha \morphEpi Y^\alpha$ is
an effective right action of $G^\alpha$ on $Y^\alpha$

\item
$\Singcat{\Triv{\catSeqName{M}^\alpha}}$ a sequence of categories

\item
$\seqname{M}^\alpha$ a sequence

\item
$
  \seqname{L}^\alpha \defeq
  \Biggl (
    \Singcat{\Triv{\catSeqName{M}^\alpha}},
    \seqname{M}^\alpha,
    \seqname{A}^\alpha,
    \seqname{F}^\alpha,
    \Sigma
  \Biggr )
$
and
$
  \seqname{L}^{\alpha,\catSeqName{M}} \defeq
  (
    \catSeqName{M}^\alpha,
    \seqname{M}^\alpha,
    \seqname{A}^\alpha,
    \seqname{F}^\alpha,
    \Sigma
  )
$
local coordinate spaces

\item
$
  \seqname{B}^\alpha \defeq
  (
    E^\alpha,
    X^\alpha,
    Y^\alpha,
    \pi^\alpha,
    G^\alpha,
    \rho^\alpha
  )
$

\item
$\seqname{B} \defeq \set{{ \seqname{B}^\alpha}}[\alpha \prec \Alpha]$
\end{enumerate}
such that
\begin{enumerate}[resume]
\item
$
  \head(\catSeqName{M},5) =
  \Bigl  (
    \catName{E},
    \catSeqName{XYG\rho},
    \catName{X},
    \catName{Y},
    \catName{G}
  \Bigr  )
$
\item
$
  \head \Biggl ( \Singcat{\Triv{\catSeqName{M}^\alpha}},5 \Biggr ) =
  \Biggl  (
    \singcat{\Triv{E^\alpha}},
    \singcat{\Triv[G^\alpha-\rho^\alpha-]{X^\alpha,Y^\alpha}},
    \singcat{X^\alpha},
    \singcat{Y^\alpha},
    \singcat{G^\alpha}
  \Biggr  )
$
\item
$
  \head(\seqname{M}^\alpha,5) =
  (
    \seqname{E}^\alpha,
    \seqname{C}^\alpha,
    X^\alpha,
    Y^\alpha,
    G^\alpha
  )
$
\item $\seqname{M}^\alpha \seqin \catSeqName{M}$
\item
$
  \head(\seqname{F}^\alpha,5) =
  (
    \pi^\alpha,
    \pi^\alpha_X,
    \pi^\alpha_Y,
    \star^\alpha,
    \rho^\alpha
  )
$

\item $C^\alpha = X^\alpha \times Y^\alpha$
\item $G^\alpha$ is a topological group with group operation $\star^\alpha$
Subsequent references to $G^\alpha$ should be read as referring to
the group rather than just the underlying topological space.
\item $\pi^\alpha$ is surjective.
\item $\pi^\alpha_X = \pi_1 \maps X^\alpha \times Y^\alpha \morphEpi X^\alpha$.
\item $\pi^\alpha_Y = \pi_2 \maps X^\alpha \times Y^\alpha \morphEpi Y^\alpha$.
\item $\rho^\alpha$ an effective right action of $G^\alpha$ on $Y^\alpha$.
\item The coordinate functions in $\seqname{A}^\alpha$ preserve fibers.
\item
$
  \head(\Sigma,5) =
  \bigl ( (0,2), (1,2), (1,3), (4,4,4), (3,4,3) \bigr )
$
\end{enumerate}

Let
\begin{enumerate}
\item $\alpha^i \prec \Alpha$, $i=1,2$
\item
$\seqname{M}^i \defeq \seqname{M}^{\alpha^i}$
\item
$\seqname{E}^i = (E^i,\catSeqName{E}^i) \defeq \seqname{E}^{\alpha^i}$
\item
$\seqname{C}^i = (C^i,\catName{C}^i) \defeq \seqname{C}^{\alpha^i}$
\item
$\seqname{A}^i \defeq \seqname{A}^{\alpha^i}$
\item
$X^i \defeq X^{\alpha^i}$
\item
$Y^i \defeq Y^{\alpha^i}$
\item
$G^i \defeq G^{\alpha^i}$
\item
$\pi^i \defeq \pi^{\alpha^i}$
\item
$\pi^i_X \defeq \pi^{\alpha^i}_X$
\item
$\pi^i_Y \defeq \pi^{\alpha^i}_Y$
\item
$\star^i \defeq \star^{\alpha^i}$
\item
$\rho^i \defeq \rho^{\alpha^i}$
\item
$
  \Singcat{\Triv{\catSeqName{M}^i}} \defeq
  \Singcat{\Triv{\catSeqName{M}^{\alpha^i}}}
$
\item
$\seqname{L}^i \defeq \seqname{L}^{\alpha^i}$
\item
$
  \seqname{L}^{i,\catSeqName{M}} \defeq
  \seqname{L}^{\alpha^i,\catSeqName{M}}
$
\item
$\seqname{B}^i \defeq \seqname{B}^{\alpha^i}$
\item
$
  \morphSeqName{f} \defeq
  (\morphName{f}_\beta, \beta \prec \length(\catSeqName{M})
$
a morphism from $\seqname{L}^1$ to $\seqname{L}^2$
\item
$\morphName{f}_E \defeq \morphName{f}_0$
\item
$\morphName{f}_C \defeq \morphName{f}_1$
\item
$\morphName{f}_X \defeq \morphName{f}_2$
\item
$\morphName{f}_Y \defeq \morphName{f}_3$
\item
$\morphName{f}_G \defeq \morphName{f}_4$
\end{enumerate}
Then:

\begin{enumerate}
\item
$\Functor^\Bun_{\LCS,\Fib} \seqname{L}^i$ and
$\Functor^\Bun_{\LCS,\Fib} \seqname{L}^{\alpha^i,\catSeqName{M}}$ are
fiber bundles.

\begin{proof}
$
  \Functor^\Bun_{\LCS,\Fib} \seqname{L}^i =
  (E^i, X^i, Y^i, \pi^i, G^i, \rho^i, \seqname{A}^i)
  =
  \Functor^\Bun_{\LCS,\Fib} \seqname{L}^{i,\catSeqName{M}}
$
satisfy the conditions of \pagecref{def:Bun}:

\begin{description}[style=unboxed]
\item%
  [
    $\seqname{C}^i$ is a $G^\alpha$-$\rho^\alpha$ model space of
    $X^\alpha \times Y^\alpha$:
  ]
\leavevmode
$\catSeqName{XYG\rho}$ is a $G$-$\rho$-model category and
$\seqname{C}^i \objin \catSeqName{XYG\rho}$ by hypothesis, and thus
$\seqname{C}^i$ is a $G^\alpha$-$\rho^\alpha$ model space of
$X^\alpha \times Y^\alpha$:  by \pagecref{def:Grhomodcat}.

\item%
  [
    $\seqname{A}^i$ is an m-atlas of $\seqname{E}^i$ in the coordinate
    space $\seqname{C}^i$:
  ]
\leavevmode
By \pagecref{def:LCS}.

\item%
  [
    $\seqname{A}^i$ is an m-atlas of \ensuremath{\Triv[X^i-\pi^i]{E^i}}
    in the coordinate space \ensuremath{\Triv[G-\rho-]{X,Y}}:
  ]
\leavevmode
By \pagecref{def:Grhomodtriv}.

\item%
  [
    $\seqname{A}^i$ is a $\pi$-$G$-$\rho$-bundle atlas of $E^i$ in the
    coordinate space $X^i \times Y^i$:
  ]
\leavevmode
By \pagecref{lem:Bun-Atl}.
\end{description}
\end{proof}

\item
$\morphName{f}_C = \morphName{f}_X \times \morphName{f}_Y$ is a
$G^1$-$G^2$-$\rho^1$-$\rho^2$ morphism from
$\Triv[G^1-\rho^1-]{X^1,Y^1}$ to $\Triv[G^2-\rho^2-]{X^2,Y^2}$.

\begin{proof}
$\morphName{f}_C$ satisfies the conditions of \pagecref{def:Grhomorph}:

\begin{description}
\item[$\morphName{f}_C = \morphName{f}_1$ is a model function:]
\leavevmode
By
\pagecref{def:LCS}.

\item%
  [
    $\morphName{f}_C$ is a model function that preserves the group action:
  ]
\leavevmode
$\morphSeqName{f}$ $\Sigma$-commutes with
$\morphSeqName{F}^1,\morphSeqName{F}^2$, as shown in \pagecref{fig:fF1F2},
by \pagecref{def:LCS} and \pagecref{def:premorph}.
\end{description}
\end{proof}

\item
$
\Functor^\Bun_{\LCS,\Fib}
  (
    \morphSeqName{f},
    \seqname{L}^1,
    \seqname{L}^2
  )
$
and
$
\Functor^\Bun_{\LCS,\Fib}
  (
    \morphSeqName{f},
    \seqname{L}^{1,\catSeqName{M}},
    \seqname{L}^{2,\catSeqName{M}}
  )
$
are bundle maps from $(\seqname{A}^1,\seqname{B}^1)$ to
$(\seqname{A}^2,\seqname{B}^2)$.

\begin{proof}
$
\Functor^\Bun_{\LCS,\Fib}
  (
    \morphSeqName{f},
    \seqname{L}^1,
    \seqname{L}^2
  )
$
and
$
\Functor^\Bun_{\LCS,\Fib}
  (
    \morphSeqName{f},
    \seqname{L}^{1,\catSeqName{M}},
    \seqname{L}^{2,\catSeqName{M}}
  )
$
satisfy the conditions of \pagecref{lem:Bun-map}
\begin{description}[style=unboxed]
\item[$\seqname{B}^i$ is a protobundle:]
\leavevmode
By hypothesis

\item%
  [
    $\seqname{A}^i$ is a $\pi^i$-$G^i$-$\rho^i$-atlas of $E^i$ in the
    coordinate space $C^i = X^i \times Y^i$:
  ]
\leavevmode
Shown above.

\item%
  [
    $\morphName{f}_X \times \morphName{f}_Y$ is a
    $G^1$-$G^2$-$\rho^1$-$\rho^2$ morphism from
    \ensuremath{ \Triv[G^1-\rho^1-]{X^1,Y^1}} to
    \ensuremath{\Triv[G^2-\rho^2-]{X^2,Y^2}}:
  ]
\leavevmode
Shown above.

\item%
  [
    \ensuremath{(\morphName{f}_E, \morphName{f}_X \times \morphName{f}_Y)}
    is an
    \ensuremath{\Triv[X^1-\pi^1-]{E^1}}-%
    \ensuremath{\Triv[X^2-\pi^2-]{E^2}}
    m-atlas morphism from $\seqname{A}^1$ to $\seqname{A}^2$ in the
    coordinate spaces \ensuremath{\Triv[G^1-\rho^1-]{X^1,Y^1}},
    \ensuremath{\Triv[G^2-\rho^2-]{X^2,Y^2}}:
  ]
\leavevmode
By \pagecref{lem:Grhomodtriv}.
\end{description}
\end{proof}

\item
$\Functor^\Bun_{\LCS,\Fib}$ is a functor from $\LCS^\Bun \seqname{B}$ to
$\Bun \seqname{B}$ and is a functor from \\
$\LCS^{\Bun,\catSeqName{M}} \seqname{B}$ to $\Bun \seqname{B}$.

\begin{proof}
\leavevmode
\begin{description}
\item[Preservation of endpoints:]
\leavevmode
\begin{align*}
  \Functor^\Bun_{\LCS,\Fib}
    (
      \morphSeqName{f}^1,
      \seqname{L}^1,
      \seqname{L}^2
    \bigr )
  & =
    \bigl (
      (
        \morphName{f}^1_E,
        \morphName{f}^1_X,
        \morphName{f}^1_Y,
        \morphName{f}^1_G
      ),
      \seqname{B}^1,
      \seqname{B}^2
    \bigr )
\\*
  \Functor^\Bun_{\LCS,\Fib}
    (\catSeqName{M}^i, \seqname{M}^i, \seqname{A}^i, \seqname{F}^i, \Sigma)
  & =
    \seqname{B}^i
\end{align*}

\item[Composition:]
\leavevmode
\begin{align*}
  \Functor^\Bun_{\LCS,\Fib}
    \Bigl (
      \bigl (
        \morphSeqName{f}^2,
        \seqname{L}^2,
        \seqname{L}^3
      \bigr )
      \morphCompose[A]
      \bigl (
        \morphSeqName{f}^1,
        \seqname{L}^1,
        \seqname{L}^2
      \bigr )
    \Bigr )
  & =
  \Functor^\Bun_{\LCS,\Fib}
    \Bigl (
       \bigl (
         \morphSeqName{f}^2 \morphCompose[()] \morphSeqName{f}^1,
         \seqname{L}^1,
         \seqname{L}^3
       \bigr )
    \Bigr )
\\*
  & =
  \bigl (
    (
      \morphName{f}^2_E \morphCompose \morphName{f}^2_2,
      \morphName{f}^2_X \morphCompose \morphName{f}^1_X,
      \morphName{f}^2_Y \morphCompose \morphName{f}^1_Y,
      \morphName{f}^2_G \morphCompose\morphName{f}^1_G
    ),
    \seqname{B}^1,
    \seqname{B}^3
  \bigr )
\\*
  & =
  \Functor^\Bun_{\LCS,\Fib}
    \bigl (
      \morphSeqName{f}^2,
      \seqname{L}^2,
      \seqname{L}^3
    \bigr )
  \morphCompose[A]
  \Functor^\Bun_{\LCS,\Fib}
    \bigl (
      \morphSeqName{f}^1,
      \seqname{L}^1,
      \seqname{L}^2
    \bigr )
\end{align*}

\item[Identity:]
\leavevmode
\begin{align*}
  \Functor^\Bun_{\LCS,\Fib} (\Id[{\seqname{L}^i})]
  & =
  \Functor^\Bun_{\LCS,\Fib}
    \Bigl (
      \Id[{\seqname{M}^i}],
      \seqname{L}^i,
      \seqname{L}^i
    \Bigr )
\\*
  & =
    \biggl (
      \bigl (
        \Id[E^i],
        \Id[X^i],
        \Id[Y^i],
        \Id[G^i]
      \bigr ),
      (\seqname{A}^i,\seqname{B}^i),
      (\seqname{A}^i,\seqname{B}^i)
    \biggr )
\\*
  \Id[{{(\seqname{A}^i,\seqname{B}^i)}}]
  & =
  \Id[{\Functor^\Bun_{\LCS,\Fib} \seqname{L}^i}]
\end{align*}
\end{description}

The proof does not depend on the category, so it applies to
$\LCS^{\Bun,\catSeqName{M}} \seqname{B}$ as well.
\end{proof}

\item
$
  \Functor^\Bun_{\LCS,\Fib}
  \morphCompose
  \Functor^\Bun_{\Fib,\LCS}
$
and
$
  \Functor^\Bun_{\LCS,\Fib}
  \morphCompose
  \Functor^{\Bun,\catSeqName{M}}_{\Fib,\LCS}
$
are the identity functor on $\Bun \seqname{B}$.
\begin{proof}
Expanding the definitions, we have
\begin{enumerate}
\item
$
  \Functor^{\Bun,\catSeqName{M}}_{\Fib,\LCS}
    (\seqname{A}^i, \seqname{B}^i) =
  \seqname{L}^i
$
\item
$
  \Functor^\Bun_{\LCS,\Fib} \seqname{L}^i =
  (\seqname{A}^i, \seqname{B}^i)
$
\item
$
  \Functor^{\Bun,\catSeqName{M}}_{\Fib,\LCS}
    \bigl (
      (
        \morphName{f}_E,
        \morphName{f}_X,
        \morphName{f}_Y,
        \morphName{f}_G
      ),
      (\seqname{A}^1, \seqname{B}^1),
      (\seqname{A}^2, \seqname{B}^2)
    \bigr ) = \\
    \bigl (
      (
        \morphName{f}_E,
        \morphName{f}_X \times \morphName{f}_Y,
        \morphName{f}_X,
        \morphName{f}_Y,
        \morphName{f}_G
      ),
      \seqname{L}^1,
      \seqname{L}^2
    \bigr )
$
\item
$
  \Functor^\Bun_{\LCS,\Fib}
    \bigl (
      (
        \morphName{f}_E,
        \morphName{f}_X \times \morphName{f}_Y,
        \morphName{f}_X,
        \morphName{f}_Y,
        \morphName{f}_G
      ),
      \seqname{L}^1,
      \seqname{L}^2
    \bigr ) = \\
    \bigl (
      (
        \morphName{f}_E,
        \morphName{f}_X,
        \morphName{f}_Y,
        \morphName{f}_G
      ),
      (\seqname{A}^1, \seqname{B}^1),
      (\seqname{A}^2, \seqname{B}^2)
    \bigr )
$
\end{enumerate}

The proof does not depend on the category, so it applies to
$\LCS^{\Bun,\catSeqName{M}} \seqname{B}$ as well.
\end{proof}

\item
$
  \Functor^\Bun_{\Fib,\LCS}
  \morphCompose
  \Functor^\Bun_{\LCS,\Fib}
$
is the identity functor on $\LCS^\Bun \seqname{B}$ and
$
  \Functor^{\Bun,\catSeqName{M}}_{\Fib,\LCS}
  \morphCompose
  \Functor^\Bun_{\LCS,\Fib}
$
is the identity functor on $\LCS^{\Bun,\catSeqName{M}} \seqname{B}$.

\begin{proof}
Expanding the definitions, we have
\begin{enumerate}
\item
$
  \Functor^\Bun_{\LCS,\Fib} \seqname{L}^i =
  (\seqname{A}^i, \seqname{B}^i)
$
\item
$
  \Functor^\Bun_{\Fib,\LCS} (\seqname{A}^i, \seqname{B}^i) =
  \seqname{L}^i
$
\item
$
  \Functor^\Bun_{\Fib,\LCS}
    \bigl (
      (
        \morphName{f}^1_E,
        \morphName{f}^1_X,
        \morphName{f}^1_Y,
        \morphName{f}^1_G
      ),
      (\seqname{A}^1, \seqname{B}^1),
      (\seqname{A}^2, \seqname{B}^2)
    \bigr ) = \\
    \bigl (
      (
        \morphName{f}^1_E,
        \morphName{f}^1_X \times \morphName{f}^1_Y,
        \morphName{f}^1_X,
        \morphName{f}^1_Y,
        \morphName{f}^1_G
      ),
      \seqname{L}^1,
      \seqname{L}^2
    \bigr )
$
\item
$
  \Functor^\Bun_{\LCS,\Fib}
    \bigl (
      (
        \morphName{f}^1_E,
        \morphName{f}^1_X \times \morphName{f}^1_Y,
        \morphName{f}^1_X,
        \morphName{f}^1_Y,
        \morphName{f}^1_G
      ),
      \seqname{L}^1,
      \seqname{L}^2
    \bigr ) = \\
    \bigl (
      (
        \morphName{f}^1_E,
        \morphName{f}^1_X,
        \morphName{f}^1_Y,
        \morphName{f}^1_G
      ),
      (\seqname{A}^1, \seqname{B}^1),
      (\seqname{A}^2, \seqname{B}^2)
    \bigr )
$
\end{enumerate}

The same proof applies for $\Functor^{\Bun,\catSeqName{M}}_{\Fib,\LCS}$
with $\seqname{L}^{i,\catSeqName{M}}$ in place of $\seqname{L}^i$.
\end{proof}
\end{enumerate}
\end{theorem}

\part{Future directions}
\label{part:fut}
If this paradigm proves useful, it can be extended to include a set of
admissible functions on the model neighborhoods of the charts,
possibly using the language of sheaves. That might be desirable for
coordinate spaces more general than Fr\'echet spaces.

Further work is needed to determine whether it is productive to define
(semi-)strict m-atlas morphisms in terms of concrete categories over a
category of model spaces rather than directly in terms of a model
category.

Further work is needed to determine whether it is productive to
allow a local coordinate space to have more than one atlas, e.g., for
more than one bundle structure on the same base space.

The extension of paracompactness to model spaces is intended to be
useful for partitions of unity on fiber bundles. Further work is
needed to determine whether that is actually the case.

The definitions given here include some fairly strong conditions,
e.g., AOC. Further work is needed to determine whether they should be
relaxed for applications beyond manifolds and fiber bundles.

Further work is needed to determine whether the concepts of
category-based atlases\footnote{As opposed to pseudogroup based}
of model spaces and prestructures have general utility.

Further work is needed to determine where it is most fruitful to require
full subcategories and where just subcategories.

Further work is needed to devise a definition of constraints that
expresses global properties, e.g., compactness, and is both clear and
rigorous.

Further work is needed to determine conditions for mappings associated
with atlas morphisms to be model functions.

This paper uses the language of category theory as an organizing
principle, but defines various notions concretely with sets. It may be
desirable to abstract away some of the details, in the spirit of, e.g.,
topoi.

It may be desirable to define presequences in the language of natural
transformations.

\begin{thebibliography}{9}

\bibitem[Ad\'amek, Herrlich, Strecker, 1990] {JoyCat} Ji\v{r}\'i Ad\'amek,
Horst Herrlich,
George E. Strecker.
\textit{Abstract and Concrete Categories The Joy of Cats},
John Wiley and Sons, Inc., 1990.

\bibitem[Kelley, 1955] {GenTop} John L. Kelley,
\textit{General Topology},
D. Van Nostrand Company (first edition), 1955.

\bibitem[Kobayashi, 1996] {FoundDiffGeo1} Shoshichi Kobayashi,
Katsumi Nomizu,
\textit{Foundations of Differential Geometry, Volume I},
ISBN 0-471-15733-3, John Wiley and Sons, Inc., 1996.

%\bibitem[Lee, 2000] {DiffPhys} Jeffrey Marc lee,
%\textit{Differential Geometry},
%Analysis and Physics, 2000.

\bibitem[Mac Lane, 1998] {CftWM} Saunders Mac Lane,
\textit{Categories for the Working Mathemation, 2nd edition},
ISBN 0-387-98403-8, Springer-Verlag, 1998.

\bibitem[Metz, 2019] {Taxonomy} Shmuel (Seymour J.) Metz,
\textit%
{
  Towards a taxonomy of atlases and of morphisms between them
},
\href{https://arxiv.org/pdf/1906.11690}{{\tt arXiv:1906.11690}, 2019}

\bibitem[Steenrod, 1999] {TopFib} Norman Steenrod,
\textit{The Topology Of Fibre Bundles}, ISBN 0-691-00548-6,
Prineton University Press (seventh printing), 1999.

\end{thebibliography}

\end{document}
